\documentclass[11pt,a4paper]{article}
\usepackage[margin=25mm]{geometry}
\usepackage{fontspec}
\setmainfont[BoldFont=lmroman10-bold.otf,ItalicFont=lmroman10-italic.otf,
 BoldItalicFont=lmroman10-bolditalic.otf]{lmroman10-regular.otf}
\setsansfont{lmsans10-regular.otf}
\setmonofont[Scale=0.87]{lmmono10-regular.otf}
\usepackage{amsmath,amssymb,amsthm,mathtools,mathrsfs}
\usepackage{booktabs,longtable,array,enumitem,microtype}
\usepackage{xcolor}
\definecolor{inkblue}{HTML}{173B52}
\usepackage[colorlinks=true,linkcolor=inkblue,urlcolor=inkblue]{hyperref}
\usepackage{xurl}
\usepackage{bookmark}
\usepackage{fancyhdr}
\pagestyle{fancy}\fancyhf{}
\fancyhead[L]{\small GRH\_NIGHT: consequences and proof gates}
\fancyfoot[C]{\thepage}
\setlength{\headheight}{14pt}
\setlength{\emergencystretch}{3em}
\setlist{nosep}
\newtheorem{proposition}{Proposition}
\newcommand{\A}{A_D}
\newcommand{\M}{M_H}
\title{\textbf{GRH\_NIGHT}\\Consequences, methods, and the Grothendieck programme}
\author{GPT-6 Astra Pro research run\\Model effort setting not exposed}
\date{21 November 2026, 07:\allowbreak{}16 CET}
\begin{document}
\maketitle
\begin{abstract}
This finite catalogue records 362 consequences, refinements, counterexamples and explicit gates.\allowbreak{}The input $A_D$ is the Alp\"oge RH/Dirichlet-GRH theorem designated in the task. The input $M_H$ is Masullo\textquotesingle s claimed universal unitary-Hecke GRH and is always conditional. Classical proofs are located rather than reproduced; new formulations receive complete proofs. ``Proved here means established in this run, not a claim of publication novelty. The source formalization was not rebuilt. The full 940-page draft was not certified. Explicit source errors are distinguished from the conditional endpoint assumption. IN PROGRESS.
\end{abstract}
\tableofcontents
\newpage
\section*{How to read the evidence}
\addcontentsline{toc}{section}{How to read the evidence}
The six-column \path{CONSEQUENCES.tsv} is the authoritative statement and
dependency catalogue. \path{QUESTIONS.tsv} records the state of every
question treated. \path{CHECK_EVENTS.tsv} records actual check events,
including who performed them and whether a proof was read, independently
derived, checked by a finite computation, or carried from the supplied
record. A citation verified at the statement level is not a claim that the
whole classical proof was rechecked. All new proof texts in the final
checked part have a complete-reading event; computational examples are
not substituted for a general proof.

All critical-line assertions concern the stated nontrivial divisor, with
Tate shifts, trivial zeros, poles and omitted Euler factors handled as
specified. A signed Brauer factorization gives zero/pole locations; it
does not supply the nonnegative divisor orders required by Artin
holomorphy. The notation $A$, $A_D$ refers to the same adopted theorem,
and $M$, $M_H$ to the same conditional draft claim.

The original prompt asks for an open-ended exhaustive investigation.
This is a finite research record. Unresolved conjectures and unavailable
proofs are displayed as such rather than given invented resolutions.
Archive locators distinguish PDF page, repository footer and manuscript
numbering where they differ. The old papers absent from the supplied ZIP
are cited through their attached return records, without a claim of fresh
whole-proof verification.

\part{Checked findings and located consequences}

% BEGIN number_theory
\section{Number-theory consequences with the hypotheses separated}
\label{sec:number-theory}

Throughout this section $A$ denotes the user-designated Alp\"oge theorem
for the Riemann zeta function and Dirichlet L-functions.  It is an adopted
cited input, not a claim that this run rebuilt the unseen Lean dependencies.
$M$ denotes the conditional all-number-field unitary Hecke-GRH assertion
of Masullo, \S1.2.  Every use of $M$ below remains conditional.  The
accompanying catalogue is \path{work/number_theory/rows.json}; its
literature records include retrieval evidence and distinguish a checked
statement from a reread of an entire classical proof.

There are three different geometric translations.  Rational primes are
the closed points of $\operatorname{Spec}\mathbb Z$; prime ideals are the
closed points of $\operatorname{Spec}\mathcal O_K$; and a conjugacy condition
on Frobenius is a condition on a finite \'{e}tale cover of an open arithmetic
curve.  The last dictionary is the Galois-category equivalence in SGA~1,
Expos\'e~V, \S7.  The class group is
$\operatorname{Pic}(\operatorname{Spec}\mathcal O_K)$ and its class fields
are finite \'{e}tale covers.  These translations identify the objects to
which the estimates apply.  They do not turn number-field GRH into a
consequence of finite-field purity.  The inherited cote-1 trace programme
is a comparison of explicit formulas; no new archive scan is certified by
this section.  Exact archive work belongs to the geometry section.

\subsection{The classical arithmetic bounds being used}

\paragraph{NT-P01--NT-P04: rational primes and gaps.}
Schoenfeld, \emph{Sharper bounds for the Chebyshev functions
$\theta(x)$ and $\psi(x)$. II} (1976), Theorem~10, printed p.~337,
equations (6.2) and (6.3), gives, under RH,
\[
 |\psi(x)-x|<\frac{\sqrt{x}\log^2x}{8\pi}\quad(x\ge73.2),\qquad
 |\theta(x)-x|<\frac{\sqrt{x}\log^2x}{8\pi}\quad(x\ge599).
\]
Equation (6.1) replaces $\log^2x$ by $(\log x-2)\log x$ for
$x\ge23\cdot10^8$.  Corollary~1 gives
$|\pi(x)-\operatorname{li}(x)|<\sqrt{x}\log x/(8\pi)$ for
$x\ge2657$.  Here $\operatorname{li}$ has its usual normalization.
The constants in Theorem~10 were checked in the primary indexed text.
An independently retrieved primary proof locator for the same
prime-counting bound is Johnston (2022), Theorem~3.2, equation (3.6),
\emph{The Ramanujan Journal} \textbf{59}, 1307--1321.
It requires RH only to a height $T$ satisfying
$9.06\sqrt{x/\log x}/\log\log x\le T$; under $A$ one may choose
such a $T$ for each $x$.  Its \S2 explicitly says its main theorem applies
also to the ordinary, non-half-weight counting functions.  Equation (3.3)
there also improves the threshold for the displayed $\psi$ inequality
to $x\ge59$.  Thus the inaccessible original Corollary~1 page is no
longer a proof-source gate for NT-P03.

Dudek--Greni\'e--Molteni, \emph{Primes in explicit short intervals on RH}
(2016), Theorem~1.1, author PDF p.~1, gives a prime in
\[
 (x-c\sqrt{x}\log x,\ x+c\sqrt{x}\log x),
 \qquad c=\tfrac12+2/\log x,\qquad x\ge2.
\]
The same theorem gives at least $\sqrt{x}$ primes after replacing $c$ by
$c+1$.  These results imply the classical
$p_{n+1}-p_n=O(\sqrt{p_n}\log p_n)$ consequence of RH.  Neither this
bound nor the assertion of RH supplies a proof here of Cram\'er's
$O(\log^2p_n)$ conjecture, twin primes, or the binary Goldbach conjecture.
Those problems require additional information about prime correlations.

\paragraph{NT-P05--NT-P08: progressions and nonresidues.}
Greni\'e--Molteni--Perelli (2016), Theorem~1, author PDF p.~2, proves that
absolute $x_0,c_1,c_2>0$ work in
\[
 \pi(x+h;q,a)-\pi(x;q,a)\ge
 c_2\frac{h}{\varphi(q)\log x},\qquad
 c_1\varphi(q)\sqrt{x}\log x\le h\le x,\quad x\ge x_0,
\]
assuming the Dirichlet GRH modulo $q$ and $(a,q)=1$.  The following
sentence in that source gives the asymptotic if the ratio to the lower
scale tends to infinity.  We have not converted the unspecified absolute
constants into numerical constants.

Lamzouri--Li--Soundararajan (LLS), corrected
\path{arXiv:1309.3595v3}, is used for the sharp explicit conductor
statements.  Corollary~1.2, author p.~3, gives the least prime in a reduced
class modulo $q>3$ at most $(\varphi(q)\log q)^2$.
Corollary~1.1, p.~2, gives a quadratic nonresidue below $(\log q)^2$ for
prime $q\ge5$.  Theorem~1.1(1), p.~2, gives the proper-subgroup bound
recorded with its explicit functions $A(q),B(q)$ in NT-P08.
These are the already carried IA02 consequences, with the corrected
version and hypotheses rechecked; they are not new findings of this run.
Geometrically the progression is a Frobenius class in a cyclotomic cover,
and the nonresidue is detected by a finite rank-one character.

\paragraph{NT-P09: the least primitive root is a different quantifier.}
Shoup, \emph{Searching for primitive roots in finite fields} (1992),
Theorem~3, author PDF p.~4, gives under Dirichlet GRH
\[
 g(p)\ll r^4(\log r+1)^4(\log p)^2,
 \qquad r=\omega(p-1).
\]
The printed theorem was checked on the page image.  Here is the short
simplification to $g(p)\ll\log^6p$.  If $q_1<\cdots<q_r$ are the distinct
prime divisors of $p-1$, then $q_j\ge j+1$ and therefore
$p-1\ge\prod_{j=1}^r(j+1)$.  For $r\ge4$ the last at least $r/2$ factors
are at least $r/2$, so
$\log(p-1)\ge(r/2)\log(r/2)$.  Consequently
$r(\log r+1)\ll\log p$; enlarging the absolute constant handles $r<4$.
Substitution proves the assertion.  This finds a small generator for a
fixed $\mathbb F_p^\times$.  The density theorem for one fixed integer
as $p$ varies uses different, generally nonabelian, fields (NT-A01).

\subsection{A quantitative Frobenius trace transfer}

\paragraph{NT-C01 and NT-C02: the precise cited Chebotarev input.}
Let $L/K$ be a finite Galois extension, $G=\operatorname{Gal}(L/K)$,
$n_L=[L:\mathbb Q]$, and $\Delta_L=|\operatorname{disc}L|$.
For a conjugacy class $C$ put $\delta_C=|C|/|G|$.  For unramified
primes only, define
\[
 \psi_C(x)=\sum_{\substack{\mathfrak p\text{ unramified},\ m\ge1\\
                    (N\mathfrak p)^m\le x}}
 \mathbf1_C(\operatorname{Frob}_{\mathfrak p}^{m})\log N\mathfrak p,
 \qquad
 \pi_C(x)=\sum_{\substack{\mathfrak p\text{ unramified}\\N\mathfrak p\le x}}
             \mathbf1_C(\operatorname{Frob}_{\mathfrak p}).
\]
Set
\begin{align*}
 E_L(x)&=\sqrt{x}\left[
       \left(\frac{\log x}{2\pi}+2\right)\log\Delta_L+
       \left(\frac{\log^2x}{8\pi}+2\right)n_L\right],\qquad x\ge1,\\
 B_L(x)&=\sqrt{x}\left[
       \left(\frac1{2\pi}+\frac3{\log x}\right)\log\Delta_L+
       \left(\frac{\log x}{8\pi}+\frac1{4\pi}+\frac6{\log x}\right)n_L
                         \right],\qquad x\ge2,\\
 \operatorname{Li}_2(x)&=\int_2^x\frac{dt}{\log t}.
\end{align*}
Greni\'e--Molteni, \emph{An explicit Chebotarev density theorem under GRH}
(2019), Theorem~1.1 and Corollary~1.2, author PDF pp.~2--3, give
\[
 |\psi_C(x)-\delta_Cx|\le\delta_C E_L(x),\qquad
 |\pi_C(x)-\delta_C\operatorname{Li}_2(x)|\le\delta_C B_L(x),
\tag{NT-C}
\]
under RH for $\zeta_L$.  The theorem explicitly includes the unramified
version of $\psi_C$.  Its constants and ranges were checked on the page
images.  In particular Artin holomorphy is not an additional hypothesis.
$M$ supplies this input for every $L$; $A$ supplies it if $L/\mathbb Q$
is abelian by the factorization into primitive Dirichlet functions.

\paragraph{NT-F01. Statement and full proof.}
For a complex class function $f$ on $G$ put
\[
 \mu_f=\frac1{|G|}\sum_{g\in G}f(g),\qquad
 \|f\|_{1,G}=\frac1{|G|}\sum_{g\in G}|f(g)|,
\]
and define $\Psi_f$ and $\Pi_f$ by replacing $\mathbf1_C$ in the two
sums above by $f$.  Then
\[
 |\Psi_f(x)-\mu_fx|\le\|f\|_{1,G}E_L(x),\qquad
 |\Pi_f(x)-\mu_f\operatorname{Li}_2(x)|\le\|f\|_{1,G}B_L(x).
\tag{NT-F01}
\]
If $f=\chi_V$ is the character of a finite-dimensional complex
representation, then $\mu_f=\dim V^G$ and
$\|f\|_{1,G}\le\sqrt{\langle\chi_V,\chi_V\rangle_G}$.

To prove the estimates, choose one representative $g_C$ from each
conjugacy class.  The finite identities
\[
 f=\sum_C f(g_C)\mathbf1_C,
 \qquad \mu_f=\sum_C\delta_Cf(g_C)
\]
give $\Psi_f-\mu_fx=\sum_Cf(g_C)(\psi_C-\delta_Cx)$.
The triangle inequality and (NT-C) bound its modulus by
$E_L(x)\sum_C\delta_C|f(g_C)|=E_L(x)\|f\|_{1,G}$.
The same finite argument proves the assertion for $\Pi_f$.
The averaging operator $|G|^{-1}\sum_g\rho(g)$ is the projection to $V^G$,
so its trace is $\mu_{\chi_V}=\dim V^G$.  Cauchy--Schwarz for the
probability measure assigning mass $1/|G|$ to each group element yields
\[
 \left(\frac1{|G|}\sum_g|\chi_V(g)|\right)^2
 \le\frac1{|G|}\sum_g|\chi_V(g)|^2
 =\langle\chi_V,\chi_V\rangle_G.
\]
This proves every assertion.  This is a linear functional estimate on the
complexified representation ring; virtual differences need not be
nonnegative, so their error is controlled by the displayed norm.

\paragraph{NT-F02. Finite \texorpdfstring{\'{e}}{e}tale point counts, with the boundary.}
Let $U$ be a nonempty open subscheme of
$\operatorname{Spec}\mathcal O_K$, and let $X\to U$ be finite \'{e}tale
of degree $d$.  Let $S_X$ be the set of geometric points of its generic
fibre, let $L/K$ be its finite Galois splitting field, and let $r$ be the
number of $G$-orbits on $S_X$, equivalently the number of generic connected
components.  Let $S$ be the finite set of primes not in $U$ which are
unramified in $L$.  Under RH for $\zeta_L$,
\[
 \left|\sum_{\substack{\mathfrak p\in U,\ m\ge1\\(N\mathfrak p)^m\le x}}
       \#X(k_{\mathfrak p^m})\log N\mathfrak p-rx\right|
 \le rE_L(x)+d|S|\log x,\qquad x\ge1.
\tag{NT-F02}
\]
Here $k_{\mathfrak p^m}$ denotes the degree-$m$ extension of
$k_{\mathfrak p}$.  In particular, when no additional unramified prime is
removed, the error is $rE_L(x)$, independent of the degree $d$.

For the proof, finite \'{e}taleness implies that inertia at a prime of $U$
acts trivially on the generic geometric fibre.  The action of
$\operatorname{Gal}(L/K)$ on this fibre is faithful by the definition of
the splitting field, so $L/K$ is unramified at every prime of $U$.
Let $\chi(g)=\#S_X^g$ be the permutation character.  The elementary
Frobenius fixed-point formula for a finite \'{e}tale scheme over a finite
field is
$\#X(k_{\mathfrak p^m})=\chi(\operatorname{Frob}_{\mathfrak p}^{m})$:
its geometric points rational over that field are precisely the points
fixed by the $m$th Frobenius power.  Also $0\le\chi(g)\le d$.
Double counting pairs $(g,s)$ with $gs=s$ gives
\[
 \sum_{g\in G}\chi(g)=\sum_{s\in S_X}|G_s|=r|G|.
\]
Thus both $\mu_\chi$ and $\|\chi\|_{1,G}$ equal $r$.  NT-F01 bounds
the sum over every unramified prime by $rx$ with error $rE_L(x)$.
For each omitted prime, the nonnegative lost contribution is at most
\[
 d\left\lfloor\frac{\log x}{\log N\mathfrak p}\right\rfloor
          \log N\mathfrak p\le d\log x.
\]
If $N\mathfrak p>x$ the contribution is zero.  Summing this bound over
$S$ and applying the triangle inequality proves (NT-F02), including
$x=1$.  Notice that replacing the nonnegative character norm by the square
root of its second moment would give a weaker bound; the orbit count is
the relevant parameter here.

\paragraph{NT-F03. Splitting types.}
For a partition $\lambda$ of $d=|S_X|$, let $D_\lambda\subseteq G$ be
the elements with cycle partition $\lambda$ on $S_X$, and put
$\delta_\lambda=|D_\lambda|/|G|$.  This is a union of conjugacy classes.
Sum the second inequality (NT-C) over these classes and use the triangle
inequality.  Their densities sum to $\delta_\lambda$, proving
\[
 \left|\#\{\mathfrak p:\ N\mathfrak p\le x,\ \mathfrak p
 \text{ unramified},\ \operatorname{Frob}_{\mathfrak p}\in D_\lambda\}
 -\delta_\lambda\operatorname{Li}_2(x)\right|
 \le\delta_\lambda B_L(x).
\]
If $D_\lambda$ is empty both sides before the absolute value are zero.
For the canonical finite \'{e}tale model of the generic algebra over the
unramified locus, each Frobenius orbit is one closed point of its fibre,
and its orbit length is that point's residue degree.  Hence this cycle
partition is exactly its splitting type.  For the originally specified
model $X$ over $U$, restrict the counted primes to $U$; at most $|S|$
terms have been removed, so the same estimate holds with the additional
error $|S|$.  This proves both interpretations without assigning a fibre
of $X$ outside its specified base.

\paragraph{NT-F04. Avoiding a prescribed finite set.}
Let $D\subseteq G$ be a nonempty conjugacy-stable set of density
$\delta_D$, and let $S$ be any finite set of primes of $K$.  The same
finite-union argument gives
$\pi_D(x)\ge\delta_D(\operatorname{Li}_2(x)-B_L(x))$.
At most $|S|$ of the counted primes belong to $S$.  If
\[
 \delta_D\bigl(\operatorname{Li}_2(x)-B_L(x)\bigr)>|S|,
\]
the count outside $S$ is strictly positive.  It is an integer, so an
eligible prime exists.  This proves the avoidance assertion without adding
an unsupported arbitrary-exclusion clause to Bach--Sorenson's least-prime
theorem.  That theorem is NT-C03, Theorem~5.1, printed p.~1729, with
RH hypotheses specified in its \S2, p.~1720.

\subsection{Class numbers, Kummer towers, and Euclidean algorithms}

\paragraph{NT-H01--NT-H04: the regulator normalization.}
Write
\begin{align*}
 U(q)&=2e^\gamma\left(\log\log q-\log2+\tfrac12+
                                      \frac1{\log\log q}\right),\\
 V(q)&=\frac{12e^\gamma}{\pi^2}\left(\log\log q-\log2+\tfrac12+
                   \frac1{\log\log q}+\frac{14\log\log q}{\log q}\right).
\end{align*}
LLS Theorem~1.5, author pp.~3--4, gives
$V(q)^{-1}\le|L(1,\chi)|\le U(q)$ for primitive conductor
$q\ge10^{10}$.  Its Corollary~1.3 is the already carried imaginary
quadratic bound in NT-H02.  The inequality is not a two-sided estimate of
order $\log\log q$: the reciprocal lower bound must be retained.

Here is the full deduction of NT-H03, including its normalization.
For $D>0$ fundamental put $K=\mathbb Q(\sqrt D)$ and let $R_D$ be the
regulator of the full unit group modulo its torsion, so
$R_D=\log\varepsilon$ for a fundamental unit $\varepsilon>1$, regardless
of its norm's sign.  The analytic class-number formula is
\[
 \mathop{\rm Res}_{s=1}\zeta_K(s)=\frac{2h_DR_D}{\sqrt D}.
\]
An exact independent locator is Milne, \emph{Class Field Theory}, v4.03,
Chapter~VI, Corollary~2.12 and Example~2.13, printed p.~191.
Since $\zeta_K(s)=\zeta(s)L(s,\chi_D)$ and $\zeta$ has residue one,
$L(1,\chi_D)=2h_DR_D/\sqrt D>0$.
Multiplying the two inequalities in NT-H01 by $\sqrt D/2$ proves
\[
 \frac{\sqrt D}{2V(D)}\le h_DR_D\le\frac{\sqrt D}{2}U(D),
 \qquad D\ge10^{10}.
\]
This completes the short proof; no bound on $R_D$ was suppressed.

For completeness, the already inherited abelian-field formulation
NT-H04 is the following precise product calculation.  If $K/\mathbb Q$
is abelian with Dirichlet-character group $X_K$, then
\[
 h_KR_K=\frac{w_K\sqrt{\Delta_K}}{2^{r_1}(2\pi)^{r_2}}
              \prod_{\substack{\chi\in X_K\\\chi\ne1}}L(1,\chi).
\tag{NT-H04}
\]
Indeed Artin factorization gives the product for $\zeta_K/\zeta$,
and taking the residue at one gives this formula by the same class-number
formula.  Its left side and the prefactor are positive, so taking moduli
turns the product into $\prod|L(1,\chi)|$.  For any factors whose
conductors are at least $10^{10}$ insert $U$ and $1/V$; retain the finite
remaining product exactly or bound it separately.  This proves every
claimed product bound.  Dividing by an unestimated regulator does not
prove the corresponding bound on $h_K$.  This is the correction already
made in RETURN\_ZETA\_TWO IA02 and IA09, not a new result of this run.

\paragraph{NT-A01 and NT-A02: the nonabelian input in Artin's primitive-root problem.}
For a fixed integer $g\ne-1$ that is not a square, put
$K_n=\mathbb Q(\zeta_n,g^{1/n})$.  Hooley's classical theorem gives a
positive density
\[
 \delta(g)=\sum_{n\ge1}\frac{\mu(n)}{[K_n:\mathbb Q]}
\]
of primes where $g$ is a primitive root, assuming RH for these Kummer
fields.  We cite Hooley, \emph{On Artin's conjecture} (1967), the
unnumbered theorem on pp.~219--220.  The original internal page has not
been freshly obtained here; a fully inspected precise proof locator for
the density implication is Lenstra, \emph{On Artin's conjecture and
Euclid's algorithm in global fields} (1977), Theorem~(3.1), printed
p.~206, specialized to $K=F=\mathbb Q$, $W=\langle g\rangle$, $k=1$,
with density (2.15), p.~206, and positivity criterion (4.1), p.~207.
The series uses actual field degrees; replacing it by the universal
Artin constant for every $g$ loses the entanglement correction.

Lenstra's version also imposes a conjugacy-stable Frobenius condition
$C\subseteq\operatorname{Gal}(F/K)$ and asks that the residue image of
a finitely generated positive-rank $W\subseteq K^\times$ have index
dividing $k$.  Its auxiliary fields are formed from
$K(\zeta_{q(\ell)},W^{1/q(\ell)})$ with
$q(\ell)=\ell^{v_\ell(k)+1}$, taking the compositum for primes dividing
a squarefree $n$.  Theorem~(3.1) assumes RH for these fields' zeta
functions.  Theorem~(4.1) gives the finite obstruction criterion for
positive density.  $M$ supplies the hypotheses; $A$ alone does not
supply them because the Kummer fields need not be abelian over
$\mathbb Q$.  In geometric terms this is a tower of Kummer torsors of
$\mathbb G_m$, not merely the cyclotomic torsion cover.

\paragraph{NT-A03: a PID becomes Euclidean under the stated Kummer GRH.}
Lenstra's Theorem~(9.1), printed p.~218, states the result recorded in
NT-A03 for the ring of $S$-integers with $|S|\ge2$ and trivial ideal
class group.  Equation (9.2) gives its least Euclidean function.
The original page image confirms the weak inequality $|S|\ge2$;
the OCR's apparent strict inequality was not adopted.  The result is a
classical cited theorem with an exact proof locator, not a new proof of
the Euclidean algorithm.  The independent PID hypothesis and the
distinction from a norm-Euclidean algorithm remain in the catalogue.

\subsection{Growth, cancellation, and zero statistics}

\paragraph{NT-L02 and NT-L03. The fixed-character Lindel\"of deduction.}
The classical input is Simoni\v c, \emph{Estimates for L-functions in the
critical strip under GRH with effective applications} (2023), Theorem~1,
equations (4)--(6), proved in its \S3.  For a function $F$ in the Selberg
class with polynomial Euler product, zero-free to the right of the
critical line, it implies, for fixed $\delta>0$ and large $|t|$,
\[
 \log|F(\sigma+it)|,\qquad -\log|F(\sigma+it)|
 \ll_{F,\delta}(\log|t|)^{1-2\delta}\log\log|t|,
 \quad \tfrac12+\delta\le\sigma\le1.
\tag{NT-L-log}
\]
For $\delta\ge1/2$ one uses instead a smaller positive $\delta$ and
the absolutely convergent Euler product on $\sigma>1$.
The modulus bound for $\log F$ in the cited theorem gives both signs of
its real part.  Dividing the right side of (NT-L-log) by $\log|t|$
shows that it tends to zero for fixed $0<\delta<1/2$.  Hence, for every
$\epsilon>0$,
\[
 F(\sigma+it),\quad F(\sigma+it)^{-1}
 \ll_{F,\delta,\epsilon}(1+|t|)^\epsilon
 \qquad(\sigma\ge\tfrac12+\delta),
\tag{NT-L-right}
\]
with pole neighbourhoods omitted if $F$ has a pole.  The finite range of
$t$ is absorbed into the constant on any closed set avoiding poles and
zeros.  On $\sigma\ge1+\delta$ the Euler product gives the required
uniform bound directly; compact intermediate strips are already covered
by (NT-L-log) or the cited theorem's adjacent range.

Here is the passage to the critical line, often suppressed in invoking
Lindel\"of.  Fix $0<\delta<1/2$.  If $F$ has degree $d$, its functional
equation and Stirling's formula imply
\[
 |F(\tfrac12-\delta+it)|
 \ll_{F,\delta}(1+|t|)^{d\delta}
       |F^\vee(\tfrac12+\delta-it)|.
\]
The dual has the same zero-free hypothesis, so NT-L-right bounds this by
$O((1+|t|)^{d\delta+\eta})$ for every $\eta>0$.  The right boundary
$\sigma=1/2+\delta$ has exponent $\eta$.
The Phragm\'en--Lindel\"of principle on this strip gives at its midpoint
\[
 |F(\tfrac12+it)|\ll_{F,\delta,\eta}
                       (1+|t|)^{d\delta/2+\eta}.
\]
Its hypotheses hold because the standard functional equation and
finite-order analytic continuation give polynomial growth in each fixed
vertical strip; if there is a pole, apply the principle above and below
a fixed height containing it.  Choose $\delta$ and $\eta$ so that
$d\delta/2+\eta<\epsilon$.  The remaining bounded interval is compact,
and the critical line contains no pole for these L-functions.  This
proves $F(1/2+it)\ll_{F,\epsilon}(1+|t|)^\epsilon$.

For primitive Dirichlet functions the hypotheses are classical, and $A$
is the zero-free input; this proves NT-L02.  For a fixed non-pure-norm
unitary Hecke character, the Hecke L-function is entire of finite order,
its gamma shifts have nonnegative real parts, and its degree is
$[K:\mathbb Q]$.  On grouping its ideal Euler factors by rational primes,
one obtains a polynomial Euler product of that degree with local roots
of modulus at most one.  The Dirichlet coefficients are bounded by
$d_{[K:\mathbb Q]}(n)=O_\epsilon(n^\epsilon)$: each local coefficient
is bounded by the corresponding coefficient of $(1-p^{-s})^{-[K:\mathbb Q]}$.
Thus it meets the explicit axioms in Simoni\v c \S2.1, and $M$ supplies
zero-freeness.  If the character is a pure unitary norm power, its
L-function is a vertical translate of $\zeta_K$, to which the same
argument applies before translation.  Imprimitive Euler factors, when
included in the convention, form a fixed finite product; they and their
inverses are bounded on each $\sigma\ge1/2+\delta$, since their roots
have modulus at most one.  This proves NT-L03 in all the conventions just
specified.  Constants may depend on the field, character, infinity type,
and $\epsilon$; this proof makes no uniform conductor-aspect claim.

\paragraph{NT-L05 and NT-L06. Mertens, Liouville, and reciprocal tails.}
Titchmarsh, second edition, Theorem~14.25(C), printed p.~370, is the exact
classical equivalence
$\mathrm{RH}\Longleftrightarrow M(x)=O_\epsilon(x^{1/2+\epsilon})$
for every $\epsilon>0$.  We use that theorem rather than reproduce its
contour proof.  The following short deductions are included in full.
For a prime power $p^a$,
\[
 \sum_{d^2\mid p^a}\mu(p^a/d^2)=
 \begin{cases}1,&a\text{ even},\\-1,&a\text{ odd},\end{cases}
\]
because only the term with remaining exponent zero or one is nonzero.
Both sides are multiplicative, hence
$\lambda(n)=\sum_{d^2\mid n}\mu(n/d^2)$ for every $n$.
Consequently
\[
 \left|\sum_{n\le x}\lambda(n)\right|
 \le\sum_{d\le\sqrt x}|M(x/d^2)|
 \ll_\epsilon x^{1/2+\epsilon}
                \sum_{d\ge1}d^{-1-2\epsilon}
 \ll_\epsilon x^{1/2+\epsilon}.
\]
This proves Liouville cancellation, with the dependence on $\epsilon$
retained in the convergent sum.

Let $s=\sigma+it$ and choose $\eta>0$ with $1/2+\eta<\sigma$.
Partial summation for $y>x$ gives
\[
 \sum_{x<n\le y}\frac{\mu(n)}{n^s}
 =M(y)y^{-s}-M(x)x^{-s}
     +s\int_x^y M(u)u^{-s-1}\,du.
\]
The first term tends to zero as $y\to\infty$, and the integral converges
absolutely after using $|M(u)|\ll_\eta u^{1/2+\eta}$.  Its modulus is at
most a constant times
$|s|x^{1/2+\eta-\sigma}/(\sigma-1/2-\eta)$.
Together with the boundary term this proves the stated tail estimate.
It also proves local uniform convergence of the reciprocal Dirichlet
series on $\operatorname{Re}s>1/2$ by choosing $\eta$ below the distance
of a compact set from the boundary.  It agrees with $1/\zeta(s)$ on
$\operatorname{Re}s>1$ by the absolutely convergent Euler product and
therefore throughout that half-plane by analytic continuation.  At
$s=1$, $1/\zeta$ has value zero.  Hence $\sum\mu(n)/n=0$, as asserted.

\paragraph{NT-L04, NT-L07, NT-L08: what moments and zero density say.}
Harper (2013), Theorem~1, author p.~1, gives
$\int_T^{2T}|\zeta(1/2+it)|^{2k}dt\ll_kT(\log T)^{k^2}$ for fixed
$k\ge0$ under RH.  Its proof does not supply the conjectural leading
constant for every $k$.  Simoni\v c (2023), Theorem~2, equations (7),(8),
provides the explicit effective Mertens constants in NT-L07.
For NT-L08, if $\sigma>1/2$, every nontrivial zero covered by $A$ or
conditionally $M$ has real part $1/2<\sigma$, so the indicated counting
set is empty.  At $\sigma=1/2$ this reasoning says nothing about
multiplicity.  Trivial zeros from gamma factors and zeros of removed
finite Euler factors are not silently included in this definition.

\paragraph{NT-Z01 and NT-Z02: the exact pair-correlation range.}
For zeta zero ordinates $0<\gamma,\gamma'\le T$, counted with
multiplicity, put
\[
 F(\alpha,T)=\left(\frac{T}{2\pi}\log T\right)^{-1}
 \sum_{\gamma,\gamma'}T^{i\alpha(\gamma-\gamma')}
                    \frac4{4+(\gamma-\gamma')^2}.
\]
Montgomery (1973), \S1, unnumbered theorem, printed p.~181, proves under
RH the formula in NT-Z01 uniformly on each $0\le\alpha\le1-\epsilon$.
Corollary~2, p.~182, gives the proportion $2/3+o(1)$ of simple zeros.
The extension of the pair-correlation formula beyond that range is
explicitly a conjecture in that paper.  Adopting RH removes the theorem's
hypothesis; it does not replace its range by all $\alpha$.

\paragraph{NT-Z03--NT-Z05: low zeros of the prime-modulus Dirichlet family.}
For prime $q$ average over the $q-2$ nonprincipal characters, with scaled
ordinates $x_{\chi,j}=\gamma_{\chi,j}\log q/(2\pi)$.
Hughes--Rudnick, \emph{Linear statistics of low-lying zeros of
L-functions} (2003), Theorem~3.1, author p.~5, gives the unitary mean for
admissible even test functions with Fourier support in $[-2,2]$.
Theorem~4.1, p.~7, gives variance
$\int_{-1}^1|u|\widehat f(u)^2du$ with support in $[-1,1]$.
Theorem~5.1, p.~9, gives the Gaussian moments only for
$m<2/\alpha$ when the support is in $[-\alpha,\alpha]$.
These identities are unconditional with complex $\gamma$ allowed;
$A$ supplies the real-ordinate interpretation.  They are not an
unrestricted all-moment Gaussian theorem for one fixed test function.

In that source's ordering on p.~1, $\gamma_{\chi,1}$ is the first
\emph{nonnegative} ordinate, so a central zero is included with value zero.
With exactly this convention, Corollary~8.2, p.~20, gives
$\liminf_q\min_{\chi\ne\chi_0}x_{\chi,1}\le1/4$ under $A$.
Theorem~8.3, pp.~20--21, gives a positive proportion for
$x_{\chi,1}<\beta$ if $\beta\ge0.633$.
No strictly-positive-zero variant was inferred by discarding central
zeros.  The geometric comparison is with finite-field family monodromy
and its symmetry group; it does not identify the unrestricted
number-field limiting distribution.

\paragraph{NT-Z06 and NT-Z07: the additional gates.}
No argument in the cited sources proves full GUE, simplicity of every
zero, or central nonvanishing for every quadratic Dirichlet function
from zero location alone.  We record the specific missing correlation
or nonvanishing estimates, without claiming a logical independence
theorem.  For Iwaniec--Luo--Sarnak, \emph{Low lying zeros of families of
L-functions} (2000), the standing assumption on printed p.~56 explicitly
includes RH for the GL$_2$ functions in the family as well as Dirichlet
GRH.  Theorem~1.1, pp.~58--59, gives support $(-2,2)$ for fixed even
weight and squarefree level tending to infinity.  $A$ supplies its
Dirichlet part.  A universal Hecke statement does not identify a general
non-CM GL$_2$ newform with a rank-one Hecke character.  Thus the GL$_2$
hypothesis remains an additional input in this invocation.

\paragraph{NT-Z08 and NT-Z09: two exact classical equivalents.}
Rodgers--Tao, \path{arXiv:1801.05914v5}, Theorem~1, author p.~3,
proves $\Lambda\ge0$ for the classical de Bruijn--Newman constant.
The definition and RH equivalence in its \S1 give
$A\Rightarrow\Lambda\le0$, hence $\Lambda=0$.
For the deformation as normalized there, real-rootedness holds exactly
for $t\ge\Lambda$, yielding the two assertions in NT-Z08.
Lagarias, \path{arXiv:math/0008177v2}, Theorem~1.1, author p.~1,
is the exact locator for NT-Z09's divisor-sum inequality.
These are cited classical consequences, not additional new proofs of RH.

\subsection{Fixed-character cancellation and the local zero count}

\paragraph{NT-L09 and NT-L10. A full Perron transfer.}
Here is the elementary analytic form used for both rational and ideal
M\"obius sums.  Suppose
$G(s)=\sum_{n\ge1}a_nn^{-s}$ is absolutely convergent for
$\operatorname{Re}s>1$, $|a_n|\le d_D(n)$ for a fixed positive integer
$D$, and $G$ extends holomorphically to $\operatorname{Re}s>1/2$.
Assume that for every $\delta,\eta>0$,
\[
 G(\sigma+it)\ll_{\delta,\eta}(1+|t|)^\eta
 \quad\text{uniformly for }\tfrac12+\delta\le\sigma\le2.
\tag{NT-Perron-input}
\]
Then $\sum_{n\le x}a_n\ll_\epsilon x^{1/2+\epsilon}$ for every
$\epsilon>0$.

We prove the assertion, recording the truncation error as well as the
contour displacement.  It suffices to take $0<\epsilon<1/4$ and
half-integral $x\ge4$, because replacing any real $y$ by
$\lfloor y\rfloor+1/2$ leaves the partial sum unchanged and changes its
size by a bounded factor; the bounded remaining range is absorbed into
the constant.  Put
$c=1+1/\log x$, $T=x^2$, $\delta=\epsilon/4$ and
$\eta=\epsilon/8$.  The classical truncated Perron kernel gives
\[
 \sum_{n\le x}a_n=
 \frac1{2\pi i}\int_{c-iT}^{c+iT}G(s)\frac{x^s}{s}\,ds+R,
\]
with
\[
 |R|\ll\sum_{n\ge1}|a_n|(x/n)^c
           \min\left(1,\frac1{T|\log(x/n)|}\right).
\tag{NT-Perron-error}
\]
This is the usual Perron formula; an exact classical locator is
Titchmarsh, Lemma~3.12 and its proof, printed pp.~61--62.  The kernel
estimate follows by moving the integral of $y^s/s$ to the left or right
according to $y>1$ or $y<1$, taking the residue at zero in the first
case, and integrating the horizontal tails; the alternative trivial
bound handles $|\log y|\le T^{-1}$.

For $n\le x/2$ or $n\ge2x$ the denominator $|\log(x/n)|$ is bounded
below by $\log2$, so that part of (NT-Perron-error) is
\[
 \ll x^cT^{-1}\sum_{n\ge1}d_D(n)n^{-c}
 =x^cT^{-1}\zeta(c)^D\ll xT^{-1}(\log x)^D.
\]
For $x/2<n<2x$, use $d_D(n)\ll_{D,\eta}x^\eta$ and
$|\log(x/n)|\gg |x-n|/x$.  Since $x$ is half-integral,
$\sum_{x/2<n<2x}|x-n|^{-1}\ll\log x$.  Also $(x/n)^c\ll1$,
uniformly for these $x,n$.  This part of the error is therefore
$O(x^{1+\eta}T^{-1}\log x)$.  Both parts are smaller than
$x^{1/2+\epsilon}$.

Move the vertical line to $\sigma_0=1/2+\delta$.  The integrand has no
pole in the intervening rectangle: $G$ is holomorphic there and
$\sigma_0>0$.  On either horizontal segment,
(NT-Perron-input) bounds the integral by
$O(x^cT^{\eta-1})=O(x^{-1+2\eta})$.
On the new vertical segment it gives
\[
 \ll x^{\sigma_0}\left(1+\int_1^Tt^{\eta-1}\,dt\right)
 \ll_\eta x^{1/2+\delta}T^\eta
 =O(x^{1/2+\epsilon/2}).
\]
The bounded part $|t|\le1$ is harmless because $\sigma_0>0$ and $G$
is bounded on its compact segment.  This proves the asserted partial-sum
bound with room in the exponent.

Apply this lemma first to $G(s)=1/L(s,\chi)$.
Its coefficients are $a_n=\mu(n)\chi(n)$, including the usual zero
values at integers not coprime to the conductor.  Under $A$, the
reciprocal is holomorphic on $\operatorname{Re}s>1/2$; the principal
L-function's pole produces a zero of the reciprocal, not a pole.
NT-L02 supplies (NT-Perron-input), so the twisted bound follows.
Character orthogonality, for $(a,q)=1$, is the finite identity
\[
 \sum_{\substack{n\le x\\n\equiv a\ ({\rm mod}\ q)}}\mu(n)
 =\frac1{\varphi(q)}\sum_{\chi\ ({\rm mod}\ q)}
       \overline{\chi(a)}\sum_{n\le x}\mu(n)\chi(n).
\]
Imprimitive functions differ by finitely many Euler factors whose zeros
lie on $\operatorname{Re}s=0$, and their reciprocals are bounded in any
closed half-strip to the right of $1/2$.  Thus each term satisfies the
same fixed-character argument.  A finite sum proves the fixed-$q$
progression bound, with no uniformity in $q$ claimed.

For a unitary Hecke character, expand
\[
 \frac1{L(s,\eta)}
 =\sum_{(\mathfrak a,\mathfrak f)=1}
         \mu_K(\mathfrak a)\eta(\mathfrak a)(N\mathfrak a)^{-s}
 =\sum_{n\ge1}a_nn^{-s}.
\]
For each $n$, $|a_n|$ is at most the number of ideals of norm $n$.
If $[K:\mathbb Q]=D$, comparison of the rational-prime Euler factors
of $\zeta_K$ with $(1-p^{-s})^{-D}$ gives this number at most $d_D(n)$.
Indeed there are at most $D$ prime ideals above $p$, each of residue
degree at least one, and replacing every local factor by a degree-one
factor can only increase its nonnegative coefficients.  Holomorphy and
(NT-Perron-input) follow conditionally on $M$ from NT-L03, including
pure norm twists.  The lemma proves NT-L10.  Averaging over the finite
character group of any fixed ray class group then projects onto one ray
class exactly as above; all sums are over ideals prime to its modulus.
This completes every claimed cancellation statement.

\paragraph{NT-Z10. The fixed infinite-order Hecke zero count.}
Let $\eta$ be any fixed unitary Hecke character of a number field $K$.
Its number of nontrivial zeros, counted with multiplicity, in
\[
 0\le\operatorname{Re}s\le1,\qquad |\operatorname{Im}s-T|\le1
\]
is $O_{K,\eta}(\log(|T|+3))$, unconditionally.  This assertion includes
infinite-order characters; its constant is not asserted uniform in
their infinity types.

Here is a full local argument.  On any fixed vertical strip the finite
Hecke L-function has a polynomial upper bound in $|t|$, away from its
possible pole.  In the present argument this follows directly from the
unconditional first part of Simoni\v c's Theorem~1, equation (2), on
$\sigma\ge1/2$, and from the functional equation and Stirling's formula
on the reflected strip.  The Selberg-class hypotheses for a non-pure-norm
character were verified in NT-L03 without using GRH; a pure norm twist
is a translate of the Dedekind function.  Consequently constants $C,A>0$
exist such that, for large $|T|$,
\[
 \max_{|s-(2+iT)|\le4}|L(s,\eta)|\le C(|T|+3)^A.
\tag{NT-J-upper}
\]
For a pure norm twist the only pole has a fixed imaginary coordinate, so
choose $|T|$ sufficiently large that this disk does not contain it.
In all other cases the function is entire.

The Euler product at the centre gives a lower bound that a general
Schwartz Mellin transform need not have:
\[
 |L(2+iT,\eta)|
 \ge\prod_{\mathfrak p\nmid\mathfrak f}(1+(N\mathfrak p)^{-2})^{-1}
 \ge\prod_{\mathfrak p}(1-(N\mathfrak p)^{-2})
 =\zeta_K(2)^{-1}>0.
\tag{NT-J-lower}
\]
These comparisons follow term by term from $|\eta(\mathfrak p)|=1$
and $1/(1+u)\ge1-u$ for $0\le u<1$, and persist on taking limits of
the absolutely convergent products.

Let $n_T(3)$ count zeros in the disk of radius $3$ about $2+iT$.
Jensen's formula on the radius-$4$ disk, or its limiting-radius version
if a zero lies on the outer circle, implies
\[
 n_T(3)\log(4/3)
 \le\log\max_{|s-(2+iT)|\le4}|L(s,\eta)|
                -\log|L(2+iT,\eta)|.
\]
The right side is $O_{K,\eta}(\log(|T|+3))$ by (NT-J-upper) and
(NT-J-lower).  Every point with $0\le\operatorname{Re}s\le1$ and
$|\operatorname{Im}s-T|\le1$ is at distance at most $\sqrt5<3$
from the centre, so its zeros are included in $n_T(3)$.  The remaining
bounded interval of $T$ meets only finitely many zeros, by meromorphic
analyticity and the absence of an identically zero function; increasing
the constant covers it.  This proves the claim.

The finite-order conductor-explicit counterpart is
Greni\'e--Molteni2019, Appendix~A, Theorem~A.1, author p.~26.
NT-Z10 instead proves the fixed-character form needed for the local
arithmetic divisor, and supplies the infinite-order version of root
MT-M04.  The argument uses the lower bound at its Jensen centre and
polynomial vertical growth.  It does not repair a statement quantified
over every function in an arbitrary Schwartz Mellin-transform space.

\subsection{An explicit uniform cyclotomic specialization}

\paragraph{NT-P11 and NT-P12.}
For $q\ge3$ and $(a,q)=1$, define
\[
 \psi(x;q,a)=\sum_{\substack{p^m\le x\\p^m\equiv a\ ({\rm mod}\ q)}}\log p,
 \qquad
 \pi(x;q,a)=\#\{p\le x:p\equiv a\ ({\rm mod}\ q)\}.
\]
Put $e_q=1$ if $q\equiv2\pmod4$ and $e_q=0$ otherwise.
Then $A$ implies
\begin{align*}
 \left|\psi(x;q,a)-\frac{x}{\varphi(q)}\right|
 &\le\sqrt{x}\left[
 \left(\frac{\log x}{2\pi}+2\right)\log q+
 \frac{\log^2x}{8\pi}+2\right]+e_q\log x,\qquad x\ge1,\\
 \left|\pi(x;q,a)-\frac{\operatorname{Li}_2(x)}{\varphi(q)}\right|
 &\le\sqrt{x}\left[
 \left(\frac1{2\pi}+\frac3{\log x}\right)\log q+
 \frac{\log x}{8\pi}+\frac1{4\pi}+\frac6{\log x}\right]+e_q,
 \qquad x\ge2.
\end{align*}

For the full proof, take $L=\mathbb Q(\zeta_q)$ and $K=\mathbb Q$.
The Galois group is $(\mathbb Z/q\mathbb Z)^\times$, of order
$n_L=\varphi(q)$, and its conjugacy classes are singletons.
The cyclotomic discriminant formula gives
\[
 \Delta_L=q^{\varphi(q)}
    \prod_{p\mid q}p^{-\varphi(q)/(p-1)},\qquad
 \log\Delta_L\le\varphi(q)\log q.
\]
At every prime $p\nmid q$, arithmetic Frobenius acts on $\zeta_q$ by
$\zeta_q\mapsto\zeta_q^p$.  Thus its $m$th power is the element
represented by $p^m$ modulo $q$.  The use of arithmetic or geometric
Frobenius merely replaces the labelled class by its inverse; the estimate
is identical after that relabelling.  The Dirichlet factorization of
$\zeta_L$ makes RH for it a consequence of $A$.
Apply NT-C01 and NT-C02 with $\delta_C=1/\varphi(q)$ and insert the
discriminant upper bound; all displayed coefficients multiplying
$\log\Delta_L$ are positive on the stated ranges.

If $q$ is odd or $4\mid q$, the primes unramified in $L$ are exactly
those not dividing $q$, so the resulting sums are the progression sums
just defined.  If $q\equiv2\pmod4$, write $q=2m$ with $m$ odd.
Then $L=\mathbb Q(\zeta_m)$ and the prime $2$ is unramified, whereas it
is excluded from a reduced class modulo $q$.  For odd primes the class
condition modulo $m$ is equivalent to that modulo $q$, since reduction
$(\mathbb Z/q\mathbb Z)^\times\to(\mathbb Z/m\mathbb Z)^\times$ is
an isomorphism.  The unramified Chebotarev sum can therefore differ only
by the powers of $2$ whose Frobenius powers lie in the designated class.
Their total weight is at most
$\lfloor\log x/\log2\rfloor\log2\le\log x$.
For the prime-counting sum the discrepancy is at most one.  Applying the
triangle inequality proves both formulas, including this degeneration.

The small moduli are transparent separately:
\begin{align*}
 \psi(x;1,0)&=\psi(x),& \pi(x;1,0)&=\pi(x),\\
 \psi(x;2,1)&=\psi(x)-\lfloor\log x/\log2\rfloor\log2,&
 \pi(x;2,1)&=\pi(x)-\mathbf1_{x\ge2}.
\end{align*}
For $1\le x<2$ the floor is zero.  Hence NT-P01 and NT-P03 give the
corresponding bounds in their stated ranges, with the explicit powers-of-two
correction retained for modulus two.  These are estimates on the open
cyclotomic arithmetic curve, and the additional term is exactly its
removed-prime contribution.

\subsection{A specified quadratic family and central nonvanishing}

\paragraph{NT-Z11. The exact one-level-density input.}
Let $\mathcal D$ consist of the nonzero odd squarefree integers, of both
signs.  For $d\in\mathcal D$, $8d$ is a fundamental discriminant and
$\chi_{8d}$ is its primitive quadratic character.  Put
$\ell_X=\log(X/(2\pi e))$.  Let $w$ be an even, nonnegative, nonzero
Schwartz function and write
\[
 W_w(X)=\sum_{d\in\mathcal D}w(d/X),\qquad
 \mathcal Z_{w,\phi}(X)=\frac1{W_w(X)}
  \sum_{d\in\mathcal D}w(d/X)
       \sum_{\rho=1/2+i\gamma}\phi(\gamma\ell_X/(2\pi)),
\]
where the inner sum counts every nontrivial zero with multiplicity.
Use the Fourier convention
$\widehat\phi(u)=\int_{\mathbb R}\phi(t)e^{-2\pi iut}\,dt$.
For a real even Schwartz function $\phi$ with
$\operatorname{supp}\widehat\phi\subset(-2,2)$, $A$ implies
\[
 \lim_{X\to\infty}\mathcal Z_{w,\phi}(X)
 =\widehat\phi(0)-\frac12\int_{-1}^1\widehat\phi(u)\,du.
\tag{NT-Z11}
\]
This is the leading term of Fiorilli--Parks--S\"odergren,
\path{arXiv:1601.06833}, Theorem~1.1, author p.~2, with definitions
(1.1)--(1.2) on the same page.  Its proof uses Theorem~3.5, author p.~12, and is concluded on p.~14.
The GRH input in the prime-sum step is Dirichlet GRH
(Lemma~2.5, author p.~7), which $A$ supplies.  No unrestricted Fourier
support, individual central nonvanishing, or nonabelian GRH is being
asserted here.

\paragraph{NT-Z12. A full nonvanishing deduction, including deweighting.}
Under $A$, at least $15/16$ of this specified family have a nonzero
central value in the following precise sense:
\[
 \liminf_{X\to\infty}
 \frac{\#\{d\in\mathcal D:|d|\le X,\ L(1/2,\chi_{8d})\ne0\}}
      {\#\{d\in\mathcal D:|d|\le X\}}\ge\frac{15}{16}.
\tag{NT-Z12}
\]
The same lower bound holds for the weighted proportions defined by any
of the fixed weights $w$ in NT-Z11.  This is a deduction from the cited
classical density theorem, not a claim of a new nonvanishing theorem.
Here are all approximation and counting details.

Let $m_d$ be the order of vanishing at $1/2$.  The completed primitive
quadratic function has root number $+1$.  One way to see the sign without
choosing a Gauss-sum square root is to divide the symmetric completed
Dedekind functional equation of $\mathbb Q(\sqrt{8d})$ by that of
$\mathbb Q$: the factorization
$\zeta_{\mathbb Q(\sqrt{8d})}=\zeta L(s,\chi_{8d})$ and the real or
complex gamma factors identify the quotient, up to a nonzero constant,
with the completed quadratic function.  The gamma duplication formula
handles the imaginary quadratic case.  Thus this entire function is
even about $s=1/2$.  Its gamma factor is finite and nonzero at the
centre, so $m_d$ is an even nonnegative integer.  In particular
\[
 2\mathbf1_{L(1/2,\chi_{8d})=0}\le m_d.
\tag{NT-central-parity}
\]
These functional equations are the standard Hecke equations; a precise
source for the completed convention is Ramakrishnan--Valenza,
Theorem~7-19, printed p.~279.

Choose real even smooth functions $g_j$ compactly supported in $(-1,1)$
such that $\int g_j=1$ and
$g_j\to g=\tfrac12\mathbf1_{[-1,1]}$ in $L^2(\mathbb R)$.
For example, smooth even cutoffs of $[-1+1/j,1-1/j]$, normalized by
their integrals, have these properties.  Define
\[
 f_j(t)=\int g_j(u)e^{2\pi iut}\,du,\qquad
 \phi_j(t)=|f_j(t)|^2.
\]
Then $\phi_j$ is real, even, nonnegative and Schwartz,
$\phi_j(0)=1$, and
$\operatorname{supp}\widehat\phi_j\subset(-2,2)$.
The support assertion follows because
$\widehat\phi_j=g_j*\widetilde g_j$, where
$\widetilde g_j(u)=\overline{g_j(-u)}$.
Plancherel gives $f_j\to f$ in $L^2$, with
\[
 f(t)=\frac{\sin(2\pi t)}{2\pi t}.
\]
Moreover, Cauchy--Schwarz yields
\[
 \|\phi_j-|f|^2\|_1
 \le\|f_j-f\|_2(\|f_j\|_2+\|f\|_2)\longrightarrow0.
\]
Consequently their Fourier transforms converge uniformly.  Since
\[
 \widehat{|f|^2}(u)=
 \begin{cases}\tfrac12-\tfrac14|u|,&|u|\le2,\\0,&|u|>2,
 \end{cases}
\]
we obtain
\[
 \lim_{j\to\infty}
 \left(\widehat\phi_j(0)-\frac12
                   \int_{-1}^1\widehat\phi_j(u)\,du\right)
 =\frac12-\frac12\cdot\frac34=\frac18.
\tag{NT-test-limit}
\]
Under $A$ every ordinate in NT-Z11 is real.  The nonnegativity of
$\phi_j$ and $\phi_j(0)=1$ therefore imply
\[
 \frac1{W_w(X)}\sum_{d\in\mathcal D}w(d/X)m_d
 \le\mathcal Z_{w,\phi_j}(X).
\]
First take $\limsup_{X\to\infty}$ for a fixed $j$, use NT-Z11, and
then let $j\to\infty$.  The result and (NT-central-parity) give
\[
 \limsup_{X\to\infty}
 \frac{\sum_{d\in\mathcal D}w(d/X)
               \mathbf1_{L(1/2,\chi_{8d})=0}}{W_w(X)}\le\frac1{16}.
\tag{NT-weighted-vanishing}
\]
This proves the weighted assertion without applying the density theorem
directly to a nonsmooth Fourier transform or to support touching $2$.

For completeness, the ordinary count follows by a separate elementary
argument.  If $A_+(Y)$ counts positive odd squarefree integers up to $Y$,
the squarefree indicator gives
\begin{align*}
 A_+(Y)&=\sum_{\substack{k\le\sqrt Y\\k\ {\rm odd}}}\mu(k)
       \#\{n\le Y/k^2:n\ {\rm odd}\}\\
 &=\frac Y2\sum_{k\ {\rm odd}}\frac{\mu(k)}{k^2}+O(\sqrt Y)
 =\frac4{\pi^2}Y+O(\sqrt Y).
\end{align*}
Here the omitted tail is $O(Y\sum_{k>\sqrt Y}k^{-2})=O(\sqrt Y)$,
and the Euler product is
$\sum_{k\ {\rm odd}}\mu(k)/k^2=8/\pi^2$.
Thus $N(X):=\#\{d\in\mathcal D:|d|\le X\}
 =8X/\pi^2+O(\sqrt X)$.
Partial summation of this counting formula gives, for every fixed
smooth compactly supported $w$,
\[
 W_w(X)=\frac4{\pi^2}X\int_{\mathbb R}w(t)\,dt+O_w(\sqrt X).
\]
Indeed apply Stieltjes integration to the positive count $A_+$ and
its reflected negative count; the integrated error is bounded by a
constant depending on $\int_0^\infty\sqrt t\,|w'(t)|dt$.

Fix $\delta>0$.  Take such an even smooth $w$ equal to $1$ on
$[-1,1]$, between $0$ and $1$ everywhere, and supported in
$[-1-\delta,1+\delta]$.  The number of vanishing members with
$|d|\le X$ is bounded above by its weighted counterpart.  Dividing by
$N(X)$ and using (NT-weighted-vanishing) gives a $\limsup$ of at most
$\frac1{16}\cdot\frac12\int w\le(1+\delta)/16$.
Letting $\delta\downarrow0$ proves (NT-Z12).

The characters in this statement are quadratic rank-one Artin objects.
Their varying conductors describe a family of double covers of open
arithmetic curves.  The symplectic expression in NT-Z11 matches the
finite-field hyperelliptic-family comparison discussed immediately after
Theorem~1.1 of the cited paper.  This comparison does not enlarge the
analytic support range or turn the $15/16$ bound into nonvanishing for
each member.  The ordinary proportion proved here concerns the
$\chi_{8d}$ family exactly; it has not been silently relabelled as all
fundamental discriminants.

\subsection{Growing fields, class numbers, and restricted unramified towers}

\paragraph{NT-H05 and NT-H06. Exact growing-field inputs.}
For pairwise nonisomorphic number fields $K_i\ne\mathbb Q$, write
$g_i=\tfrac12\log\Delta_{K_i}$ and $n_i=[K_i:\mathbb Q]$.
Let $N_q(K_i)$ count prime ideals of norm $q$, where $q$ ranges over
rational prime powers.  The family is \emph{asymptotically exact} when
the following limits exist for every indicated index:
\[
 \phi_q=\lim_i\frac{N_q(K_i)}{g_i},\qquad
 \phi_{\mathbb R}=\lim_i\frac{r_1(K_i)}{g_i},\qquad
 \phi_{\mathbb C}=\lim_i\frac{r_2(K_i)}{g_i}.
\]
Let $h_i,R_i$ be the ordinary class number and full unit regulator, and
let $\kappa_i=\operatorname{Res}_{s=1}\zeta_{K_i}(s)$.
Tsfasman--Vl\u adu\c t, \path{arXiv:math/0205129},
GRH Theorem~7.2, author p.~44 (proof pp.~44--49), gives
\begin{align*}
 \lim_i\frac{\log(h_iR_i)}{g_i}
 &=1+\sum_q\phi_q\log\frac q{q-1}
      -\phi_{\mathbb R}\log2-\phi_{\mathbb C}\log(2\pi),\\
 \lim_i\frac{\log\kappa_i}{g_i}
 &=\sum_q\phi_q\log\frac q{q-1}.
\end{align*}
Only RH for the family of Dedekind functions is needed: $M$ supplies it;
$A$ suffices when every $K_i/\mathbb Q$ is abelian.  Neither normality
nor a tower is required in this theorem.  If $n_i/g_i\to0$, then
$N_q(K_i)\le n_i$ makes all the displayed invariants zero, so the
first limit is $1$.  This proves the usual Brauer--Siegel asymptotic
in this range without adding a normality hypothesis.

The same paper's GRH Theorem~3.1, author p.~19 (proof pp.~19--21),
supplies the basic inequality
\[
 \sum_q\frac{\phi_q\log q}{\sqrt q-1}
 +\alpha_{\mathbb R}\phi_{\mathbb R}
 +\alpha_{\mathbb C}\phi_{\mathbb C}\le1,
\tag{NT-basic}
\]
where
\[
 \alpha_{\mathbb R}=\tfrac12\log(8\pi)+\tfrac\pi4+\tfrac\gamma2,
 \qquad
 \alpha_{\mathbb C}=\log(8\pi)+\gamma,
\]
and $\gamma$ is Euler's constant.  Both original theorem pages were
inspected as images.  These are cited classical theorems.  The next
two consequences and their exact tower interpretations are proved below.

\paragraph{NT-H07. A root-discriminant bound with prescribed splitting.}
Suppose $n_i\to\infty$,
\[
 \Delta_{K_i}^{1/n_i}\longrightarrow\varrho\in(1,\infty),
 \qquad r_1(K_i)/n_i\longrightarrow\tau\in[0,1].
\]
Let $P$ be any set of rational primes such that every $p\in P$ splits
completely in every $K_i$.  Under RH for all $\zeta_{K_i}$,
\[
 \log\varrho\ge\log(8\pi)+\gamma+\frac\pi2\tau
          +2\sum_{p\in P}\frac{\log p}{\sqrt p-1}.
\tag{NT-H07}
\]
If $P$ is infinite, the inequality includes convergence of this
nonnegative series.  No normality assumption on $K_i/\mathbb Q$ is
needed.

To prove the claim, first take a subsequence with strictly increasing
degrees; its fields are pairwise nonisomorphic and its stated limits
are unchanged.  Note that
$g_i/n_i\to\tfrac12\log\varrho>0$.
For every fixed prime power $q$, the numbers $N_q(K_i)/g_i$ are
bounded since $N_q(K_i)\le n_i$.  Successive subsequence extraction,
followed by a diagonal subsequence over the countable prime powers,
makes all of them converge.  The archimedean ratios already converge;
on the resulting asymptotically exact subfamily they are
\[
 \phi_{\mathbb R}=\frac{2\tau}{\log\varrho},\qquad
 \phi_{\mathbb C}=\frac{1-\tau}{\log\varrho}.
\]
Complete splitting gives $N_p(K_i)=n_i$, hence
$\phi_p=2/\log\varrho$ for $p\in P$.
Apply (NT-basic) and discard the nonnegative terms from other norms.
For every finite $P_0\subset P$, multiplication by $\log\varrho$
yields
\[
 2\sum_{p\in P_0}\frac{\log p}{\sqrt p-1}
 +2\tau\alpha_{\mathbb R}+(1-\tau)\alpha_{\mathbb C}
 \le\log\varrho.
\]
The identity
$2\alpha_{\mathbb R}-\alpha_{\mathbb C}=\pi/2$
gives the claimed expression for finite $P_0$.  Taking the supremum
over finite subsets proves convergence and the same bound for all of
$P$.  The limiting quantities $\varrho,\tau$ are unchanged by the
subsequence, so this proves the original assertion.

\paragraph{NT-H08. Finiteness of a restricted arithmetic fundamental group.}
Fix a number field $K$ and a finite set $P$ of rational primes that
split completely in $K$.  Let $K^{\rm ur,P}$ be the compositum of all
finite extensions of $K$ unramified at every finite prime, in which
every prime of $K$ above $P$ splits completely.  Put
\[
 C(P)=8\pi\exp\left(\gamma+
                 2\sum_{p\in P}\frac{\log p}{\sqrt p-1}\right).
\]
If $\Delta_K^{1/[K:\mathbb Q]}<C(P)$, then $M$ implies
$[K^{\rm ur,P}:K]<\infty$.
If $K$ is totally real and only totally real extensions are allowed,
the corresponding conclusion holds under the weaker inequality
$\Delta_K^{1/[K:\mathbb Q]}<e^{\pi/2}C(P)$.
The inequalities are strict: no equality case is claimed.

Here is the full proof.  Being unramified at finite primes and splitting
completely at specified primes are both preserved by conjugation over
$K$ and by finite composita.  Thus the displayed compositum is Galois
over $K$, and all of its finite subextensions have these properties.
If it had infinite degree, it would contain a sequence of finite
Galois subextensions $K_i/K$ with degrees tending to infinity.
The relative discriminant formula and finite unramifiedness give
\[
 \Delta_{K_i}=\Delta_K^{[K_i:K]},\qquad
 \Delta_{K_i}^{1/[K_i:\mathbb Q]}=
                  \Delta_K^{1/[K:\mathbb Q]}=: \varrho.
\]
If $K=\mathbb Q$, this equality gives discriminant one for any such
$K_i$; Minkowski's discriminant bound already rules out a nontrivial
extension.  Otherwise $\varrho>1$.
After taking a subsequence, $r_1(K_i)/[K_i:\mathbb Q]$ tends to some
$\tau\in[0,1]$.  Every $p\in P$ splits completely in $K_i/\mathbb Q$
because it splits in $K$ and each prime above it splits in $K_i/K$.
$M$ supplies the Dedekind RH hypotheses of NT-H07.  That result yields
$\varrho\ge e^{\pi\tau/2}C(P)\ge C(P)$, a contradiction.
In the totally real version all $K_i$ are totally real, so $\tau=1$
and the same proof gives the second threshold.  This completes the proof.

For $P=\varnothing$, the thresholds are
$8\pi e^\gamma\simeq44.763232$ and
$8\pi e^{\gamma+\pi/2}\simeq215.332516$.
The group $\operatorname{Gal}(K^{\rm ur,P}/K)$ is the quotient of
$\pi_1^{\text{ét}}(\operatorname{Spec}\mathcal O_K)$ obtained by
killing the closed normal subgroup generated by Frobenius elements at
the primes above $P$.  Requiring total reality also kills the
archimedean complex-conjugation classes.  This is the finite-\'{e}tale
Galois-category dictionary, SGA~1, Expos\'e~V, \S7.  The discriminant
criterion therefore gives a concrete conditional finiteness statement
for this quotient.  It is a restriction on arithmetic towers in the
cote-149/GT comparison programme, not a claim of a new anabelian
reconstruction theorem.

\subsection{Prime races: zero location versus relations among ordinates}

\paragraph{NT-R01 and NT-R02. What $A$ alone supplies.}
Fix distinct reduced classes $a_1,\ldots,a_r$ modulo $q$, and set
\[
 E_{q;\boldsymbol a}(x)=\frac{\log x}{\sqrt x}
 \bigl(\varphi(q)\pi(x;q,a_1)-\pi(x),\ldots,
       \varphi(q)\pi(x;q,a_r)-\pi(x)\bigr).
\]
Rubinstein--Sarnak, \emph{Chebyshev's Bias} (1994), Theorem~1.1,
printed p.~174, proof \S2.1, gives under $A$ a probability measure
$\mu_{q;\boldsymbol a}$ such that, for every bounded continuous $f$,
\[
 \lim_{X\to\infty}\frac1{\log X}
       \int_2^X f(E_{q;\boldsymbol a}(x))\frac{dx}{x}
   =\int_{\mathbb R^r}f\,d\mu_{q;\boldsymbol a}.
\tag{NT-R01}
\]
Its normalization uses $\pi(x)$ in each coordinate.  Theorem~1.2,
printed p.~175, proof \S2.2, and the discussion following it imply
\[
 \underline\delta\{x:\pi(x;4,3)>\pi(x;4,1)\}>0,
 \qquad
 \underline\delta\{x:\pi(x;4,1)>\pi(x;4,3)\}>0,
\tag{NT-R02}
\]
where $\underline\delta S=\liminf_{X\to\infty}
(\log X)^{-1}\int_{S\cap[2,X]}dx/x$.
The same conclusion holds for the two classes modulo three.
These statements concern logarithmic sampling.  The existence of the
limiting distribution for continuous tests does not by itself justify
substituting an ordering-cone indicator: its boundary could have mass.

\paragraph{NT-R03 and NT-R04. The separate ordinate-independence input.}
Use the source's Grand Simplicity Hypothesis (GSH): the nonnegative
ordinates of zeros of all primitive Dirichlet functions, listed with
multiplicity, are linearly independent over $\mathbb Q$.  In particular
this excludes central zeros and repeated ordinates.  GSH remains an
additional hypothesis here.

Under $A$ and GSH, for $r\ge2$ every strict ordering race has a
logarithmic density in $(0,1)$; the source's formula (1.2) and
Remark~1.3, printed p.~176, give the requisite regularity and positive
mass.  For fixed $r$,
\[
 \max_{\substack{a_1,\ldots,a_r\pmod q\\
                   \text{distinct reduced classes}}}
 \left|\delta\{x:\pi(x;q,a_1)>\cdots>\pi(x;q,a_r)\}
                   -\frac1{r!}\right|\longrightarrow0
       \quad(q\to\infty).
\tag{NT-R04}
\]
This is Theorem~1.5, printed p.~177, proof \S3.2.
For the exact symmetry classification, put
$c(q,a)=-1+\#\{b\pmod q:b^2\equiv a\pmod q\}$.
Theorem~1.4, printed p.~177, proof Proposition~3.1, says that the
limiting law is invariant under every coordinate permutation precisely
in the following cases: $r=2$ and $c(q,a_1)=c(q,a_2)$; or $r=3$
and, after labelling, the classes are
$a_1,a_1\rho,a_1\rho^2$ with $\rho\ne1$ and $\rho^3=1\pmod q$.
This is symmetry of the full law.  It must not be conflated with a
converse assertion about one numerical race density.  The discussion
on p.~177 gives, in particular, the strict bias
$1/2<\delta\{x:\pi(x;4,3)>\pi(x;4,1)\}<1$ under these extra
hypotheses.  The relevant original pp.~174--178 were visually inspected
because their PDF text encoding is defective.

Geometrically these are races between Frobenius classes in cyclotomic
covers of $\operatorname{Spec}\mathbb Z[1/q]$.  Their limiting
distribution is an explicit-formula statistic of the corresponding
rank-one objects.  The additional GSH gate concerns relations among
zero ordinates, so the zero-location inputs $A$ and $M$ do not discharge
it in this argument.  A finite-field monodromy comparison likewise
does not remove that number-field hypothesis.

\subsection{Multiplicity, gaps, and the first zero: quantitative limits}

\paragraph{NT-Z13. An argument bound with its actual archimedean hypothesis.}
Carneiro--Chandee--Milinovich, \emph{A note on the zeros of zeta and
$L$-functions} (2015), author PDF Theorem~1, p.~2, proves under RH
\[
 |S(t)|\le \frac14\frac{\log t}{\log\log t}
       +O\left(\frac{\log t\log\log\log t}{(\log\log t)^2}\right).
\tag{NT-Z13a}
\]
Here $S(t)=\pi^{-1}\arg\zeta(1/2+it)$, with the argument continued
from $2$ through $2+it$ and the mean of its two one-sided values used
at a zero.  The proof is in \S3, pp.~3--5.

We record the precise Hecke specialization of their Theorem~5,
p.~7, equations (4.4)--(4.5), whose proof occupies pp.~7--12.
Let $\eta$ be a primitive finite-order Hecke character of a number
field $K$ of degree $n$.  Write its completion, up to a nonzero
constant independent of $s$, as
\[
 \Lambda(s,\eta)=G_\eta(s)L(s,\eta),\qquad
 G_\eta(s)=Q_\eta^{s/2}\prod_{j=1}^n\Gamma_{\mathbb R}(s+\mu_j),
 \qquad Q_\eta=\Delta_K\operatorname N\mathfrak f_\eta,
\]
where $\Gamma_{\mathbb R}(s)=\pi^{-s/2}\Gamma(s/2)$.
Each real place contributes $\mu_j=0$ or $1$ according to the sign
character.  Each complex place contributes the pair $0,1$.
Set
\[
 C_\eta=Q_\eta\prod_{j=1}^n(\mu_j+3),\quad
 C_\eta(t)=C_\eta(|t|+1)^n,\quad
 D_\eta(t)=\log\log\bigl(C_\eta(t)^{3/n}\bigr).
\]
For $t>0$ define
$S_\eta(t)=-\pi^{-1}\int_{1/2}^{\infty}
\operatorname{Im}(L'/L)(\sigma+it,\eta)\,d\sigma$ away from zero
ordinates, and use one-sided means at them.  Under RH for this
particular $L$-function,
\[
 |S_\eta(t)|\le\frac{\log C_\eta(t)}{4D_\eta(t)}
   +O\left(\frac{\log C_\eta(t)\log D_\eta(t)}{D_\eta(t)^2}\right).
\tag{NT-Z13b}
\]
The implied constant is absolute.  In particular, for fixed $K,\eta$,
\[
 |S_\eta(t)|\le\frac n4\frac{\log t}{\log\log t}
   +O_{K,\eta}\left(
       \frac{\log t\log\log\log t}{(\log\log t)^2}\right).
\tag{NT-Z13c}
\]

Here is the check of the source hypotheses, which matters for $M$.
After grouping prime ideals by their underlying rational prime, the
Hecke Euler factor is a product of at most $n$ linear reciprocal
factors with local parameters of modulus at most one: factor each
$1-\eta(\mathfrak p)p^{-fs}$ into its $f$ roots and pad missing
factors by zero parameters.  Thus the source's parameter $\vartheta$
is zero.  The completion has order one, the standard functional
equation with the conjugate character, and is entire except for the
trivial character's poles at $0,1$.  Its listed gamma parameters are
real and nonnegative.  These are precisely the analytic conditions
in \S4, p.~6, supplied by classical Hecke theory (RV1999,
Theorem~7-19).  $M$ supplies the remaining zero-location assumption;
$A$ suffices for Dirichlet characters over $\mathbb Q$.
In particular, the source explicitly requires gamma parameters to be
real or to occur in conjugate pairs.  We do not silently use it for
an arbitrary infinite-order unitary Hecke character.  The paired
$\pm it_k$ family in root MT-H01--H03 can be checked separately
against this same hypothesis.

\paragraph{NT-Z14. Full deduction of multiplicity and gap bounds.}
For each fixed $K,\eta$ as in NT-Z13, under its GRH, a zero
$1/2+i\gamma$ with $\gamma\to+\infty$ has multiplicity
\[
 m_\eta(\gamma)\le\frac n2\frac{\log\gamma}{\log\log\gamma}
   +O_{K,\eta}\left(
    \frac{\log\gamma\log\log\log\gamma}{(\log\log\gamma)^2}\right).
\tag{NT-mult}
\]
For consecutive distinct sufficiently large positive ordinates
$\gamma<\gamma'$,
\[
 \gamma'-\gamma\le\frac{\pi}{\log\log\gamma}
   +O_{K,\eta}\left(
    \frac{\log\log\log\gamma}{(\log\log\gamma)^2}\right).
\tag{NT-gap}
\]
Consequently every sufficiently large real $T$ lies within
\[
 \frac{\pi}{2\log\log T}
 +O_{K,\eta}\left(\frac{\log\log\log T}{(\log\log T)^2}\right)
\]
of a zero ordinate.  The analogous
negative-height statements follow by conjugating $\eta$.  These are
classical consequences in a fixed-character formulation; they assert
neither simplicity nor pair correlation.

We give the entire deduction, including the jump and endpoint issues.
In a sufficiently high horizontal region $L(s,\eta)$ has no pole;
GRH and the boundary nonvanishing theorem give no zero to the right
of the critical line.  Continue $\log L$ from its Euler-product
normalization there.  Near $\rho=1/2+i\gamma$ write
$L(s,\eta)=(s-\rho)^m U(s)$ with $U$ holomorphic and nonzero.
Approaching the critical line from its right, the argument of
$s-\rho$ changes from $-\pi/2$ to $\pi/2$ across the ordinate.
Hence
\[
 S_\eta(\gamma+)-S_\eta(\gamma-)=m.
\]
Taking one-sided limits in NT-Z13c bounds this jump by twice its
right-hand side and proves (NT-mult).

Let $\Theta_\eta(t)=\operatorname{Im}\log G_\eta(1/2+it)$, with
any continuous determination at large positive $t$.
The functional equation and conjugation give
$\Lambda(1/2+it,\eta)=\varepsilon_\eta
\overline{\Lambda(1/2+it,\eta)}$.
Therefore, on an interval containing no zero, the argument of the
completed function is constant modulo $\pi$, and
\[
 S_\eta'(t)=-\frac1\pi\Theta_\eta'(t),\qquad
 \frac1\pi\Theta_\eta'(t)
    =\frac n{2\pi}\log t+O_{K,\eta}(1).
\tag{NT-phase}
\]
The derivative formula is Stirling's formula applied to each of the
$n$ displayed $\Gamma_{\mathbb R}$ factors; the conductor term is
constant.  The same identity and NT-Z13c show that positive ordinates
are unbounded: if there were no zeros beyond a height, integration
would force $S_\eta(t)$ to have size comparable to $t\log t$.

For large $\gamma$, a zero-free gap cannot have length at least one.
Otherwise it contains an interval of length $1/2$ whose endpoints
are between $\gamma$ and $\gamma+1$.  Integrating (NT-phase) over it
gives a decrease of size $(n/(4\pi)+o(1))\log\gamma$, whereas
NT-Z13c bounds the difference of its endpoint values by
$O_{K,\eta}(\log\gamma/\log\log\gamma)$, a contradiction.
Thus $h=\gamma'-\gamma<1$ eventually.  Integrate (NT-phase) from
just above $\gamma$ to just below $\gamma'$ and then take endpoint
limits.  Uniformly on this interval,
\[
 \left(\frac n{2\pi}\log\gamma+O_{K,\eta}(1)\right)h
   =S_\eta(\gamma+)-S_\eta(\gamma'-)
   \le\frac n2\frac{\log\gamma}{\log\log\gamma}
      +O_{K,\eta}\left(
       \frac{\log\gamma\log\log\log\gamma}
            {(\log\log\gamma)^2}\right).
\]
Division proves (NT-gap).  Any $T$ between consecutive ordinates is
at distance at most half their gap from one endpoint.  Since that
gap is less than one, replacing $\gamma$ by $T$ in the slowly varying
bound changes only the error term.  This proves the last assertion
and completes the proof.  For $K=\mathbb Q$ and the trivial
character, $n=1$ gives the usual zeta multiplicity and gap bounds
under $A$ alone.

\paragraph{NT-Z15 and NT-Z16. Conductor aspect and the first ordinate.}
In the same finite-order Hecke setting put
$D_\eta=\log\log(C_\eta^{3/n})$.
Theorem~5, equation (4.5), gives
\[
 \operatorname{ord}_{s=1/2}L(s,\eta)
 \le\frac{\log C_\eta}{2D_\eta}
   +O\left(\frac{\log C_\eta\log D_\eta}{D_\eta^2}\right).
\tag{NT-Z15}
\]
There is no gamma pole or zero at $1/2$, so the order equals that
of the completion in the source.  Theorem~7, p.~13, with its full
proof on that page, gives at least one zero $1/2+i\gamma$ satisfying
\[
 |\gamma|\le\frac{\pi}{2D_\eta}
          +O\left(\frac{\log D_\eta}{D_\eta^2}\right)
 \quad\text{as } C_\eta^{3/n}\longrightarrow\infty.
\tag{NT-Z16}
\]
Both implied constants are universal.  A possible central zero is
included: this is not a bound for the first strictly positive
ordinate.  Keeping $n$ inside $D_\eta$ is essential when the degree
varies.  For primitive Dirichlet characters of conductor $q>2$, the
specializations are
\[
 \operatorname{ord}_{1/2}L(s,\chi)
       \le(1/2+o(1))\frac{\log q}{\log\log q},\qquad
 \min_\rho|\operatorname{Im}\rho|
       \le\frac{\pi/2+o(1)}{\log\log q}.
\]
Here $A$ supplies every required GRH hypothesis.  For finite-order
Hecke characters over general fields the corresponding input is
the finite-order part of $M$; it is not supplied by $A$ merely
because the base field happens to be abelian over $\mathbb Q$.

These bounds constrain the zero divisor of rank-one arithmetic
local systems and their Artin inductions.  In the Grothendieck
comparison language they are quantitative restrictions on the
archimedean explicit-formula spectrum.  They supply neither a
finite-field purity argument over $\operatorname{Spec}\mathbb Z$
nor a general geometric multiplicity-one theorem.  They refine the
specific open gates in NT-Z06 while leaving them open.

\subsection{Further prime constants and quadratic-family moments}

\paragraph{NT-P13. Explicit and asymptotic prime-gap constants.}
Under the RH part of $A$, every $x\ge4$ has a prime in
\[
 \left[x,x+\frac{22}{25}\sqrt x\log x\right].
\]
This is Carneiro--Milinovich--Soundararajan,
\emph{Fourier optimization and prime gaps} (2019), author v2
Theorem~5, p.~5; the explicit proof is \S6, pp.~21--26.
Consequently, for every consecutive pair with $p_n>3$,
$p_{n+1}-p_n\le(22/25)\sqrt{p_n}\log p_n$.
For clarity about the closed left endpoint, apply the interval
statement to $x=p_n+\varepsilon$ with
$0<\varepsilon<p_{n+1}-p_n$, then let $\varepsilon\downarrow0$.
Corollary~4 on the same page, following Theorem~3 and its proof in
\S5, strengthens the asymptotic constant:
\[
 \limsup_{n\to\infty}
       \frac{p_{n+1}-p_n}{\sqrt{p_n}\log p_n}<\frac{21}{25}.
\]
The explicit all-$x$ statement and the sharper asymptotic statement
have different quantifiers.  No claim is made that these are the
latest possible constants.

\paragraph{NT-P14--NT-P16. Fourier improvements in the conductor aspect.}
Carneiro--Milinovich--Quesada-Herrera--Ramos,
\emph{Fourier optimization, the least quadratic non-residue, and
the least prime in an arithmetic progression}, arXiv v2 dated
August 2025, Theorems~1 and 2, author p.~2, gives under $A$
\[
 \limsup_{p\to\infty}\frac{n_p}{(\log p)^2}<0.7615,
 \qquad
 \limsup_{p\to\infty}\frac{r_p}{(\log p)^2}<0.7615.
\tag{NT-P14}
\]
Here $p$ runs over primes, $n_p$ is the least positive quadratic
nonresidue modulo $p$, and $r_p$ is the least prime $r$ with
$(r/p)=1$.  The nonresidue is itself prime.  Theorem~1 follows from
Theorem~6, whose proof is \S4, pp.~13--17; the modification proving
Theorem~2 is \S4.4, pp.~16--17.  These bounds improve the
asymptotic constant; they do not replace NT-P07's explicit threshold
$p\ge5$ by an unstated numerical threshold.

More generally, for a nonprincipal character $\chi$ modulo $q$ of
fixed order $\ell\ge2$, define
$n_\chi=\min\{n\ge1:\chi(n)\notin\{0,1\}\}$.
Theorem~6, p.~5, and Corollary~7, p.~6, give
\[
 n_\chi\le(c_\ell+o_\ell(1))(\log q)^2,
 \qquad
 \begin{array}{c|rrrr}
   \ell&2&3&4&5\\ \hline
   c_\ell&0.7615&0.6707&0.6131&0.5765
 \end{array}
\tag{NT-P15}
\]
as $q\to\infty$.  For fixed $\ell\ge6$ the same corollary permits
\[
 c_\ell=\frac14\left(
  1-\frac{\frac{\ell}{\ell-1}\log\bigl((3\ell-4)/\ell\bigr)}
             {\log(\ell-1)}+\frac1{\ell-1}\right)^{-2}.
\]
The character need not be primitive.  Its values zero at primes
dividing $q$ are explicitly excluded in the definition.  If
$n_\chi$ were composite, every prime divisor would be smaller and
coprime to $q$, hence would have character value one by minimality;
multiplicativity would give $\chi(n_\chi)=1$, a contradiction.
Thus these are also bounds for the least prime outside the kernel
of the cyclotomic character.  The limit $c_\ell\to1/4$ is a
statement about these constants; we do not assert uniformity when
$\ell$ varies with $q$.

For the least prime $P(a,q)\equiv a\pmod q$, the same paper's
Theorem~3, author p.~3, proof \S5, pp.~18--22, gives under $A$
\[
 \limsup_{q\to\infty}
  \max_{(a,q)=1}\frac{P(a,q)}{(\varphi(q)\log q)^2}<\frac89.
\tag{NT-P16}
\]
The uniformity in the reduced class follows from the proof: its
errors in (5.15)--(5.16) depend on the fixed Fourier test function
and not on $a$.  The numerical optimization certification in the
paper is cited, not rerun here.  These results concern closed points
and selected Frobenius classes of cyclotomic covers.  They sharpen
the effective detection bounds in the SGA~1 Galois-category
dictionary; they do not enlarge the class of $L$-functions furnished
by $A$.

\paragraph{NT-L11. Quadratic central moments with the full family specified.}
For each fixed real $k\ge0$, $A$ and Harper,
\emph{Sharp conditional bounds for moments of the Riemann zeta
function}, Theorem~2, author p.~15, give
\[
 \sum_{\substack{X\le |d|\le2X\\d\ {
m fundamental}}}
       |L(1/2,\chi_d)|^k
     \ll_k X(\log X)^{k(k+1)/2}.
\tag{NT-L11}
\]
The family includes positive and negative fundamental discriminants;
$\chi_d$ is the associated primitive quadratic character.  The
argument is in \S5, pp.~15--17, using Propositions~3 and 4.
The input is only GRH for the quadratic functions in question.
This is an upper bound of the predicted order, not an asymptotic
formula with its leading constant.  In particular it asserts no
bound for negative moments or reciprocals of central values.

\paragraph{NT-L12. Full fixed-threshold large-value deductions.}
For every fixed $\alpha>0$, $A$ implies
\[
 \operatorname{meas}\{t\in[T,2T]:
            |\zeta(1/2+it)|\ge(\log T)^\alpha\}
       \ll_\alpha \frac{T}{(\log T)^{\alpha^2}}.
\tag{NT-L12a}
\]
For every fixed $\alpha>1/2$ it also implies
\[
 \#\{d\text{ fundamental}:X\le|d|\le2X,
          |L(1/2,\chi_d)|\ge(\log X)^\alpha\}
       \ll_\alpha
          \frac{X}{(\log X)^{(\alpha-1/2)^2/2}}.
\tag{NT-L12b}
\]
Here is the complete proof.  For any fixed $k>0$, Markov's
inequality and Harper's Theorem~1 (NT-L04) bound the first measure by
\[
 (\log T)^{-2k\alpha}
       \int_T^{2T}|\zeta(1/2+it)|^{2k}\,dt
       \ll_k T(\log T)^{k^2-2k\alpha}.
\]
Choose the fixed value $k=\alpha$ to obtain (NT-L12a).
Similarly NT-L11 bounds the second count by
$O_k(X(\log X)^{k(k+1)/2-k\alpha})$.
For $\alpha>1/2$, choose the positive fixed value
$k=\alpha-1/2$; completing the square gives
$k(k+1)/2-k\alpha=-(\alpha-1/2)^2/2$.
This proves (NT-L12b).  All constants may depend on the fixed
threshold exponent; no choice of $k$ growing with $T$ or $X$ has
been made.  Central zeros cause no difficulty because only positive
moments occur.  The result is a direct optimization of the cited
moment bounds, not a new moment theorem.

The quadratic family consists of rank-one Artin objects coming from
quadratic covers with growing discriminant.  Its central-value
moment exponent is an arithmetic family statement.  A geometric
symmetry-type comparison is useful context, but is not substituted
for Harper's analytic proof or the exact GRH input.

\subsection{Closed-point Frobenius relations in unramified towers}

\paragraph{NT-H09. A relative splitting-degree budget.}
Let $K\ne\mathbb Q$ be a number field of degree $n$, discriminant
$\Delta_K$, and root discriminant $\varrho=\Delta_K^{1/n}$.
Let $S$ be a finite set of prime ideals of $K$, with
$q_{\mathfrak p}=\operatorname N\mathfrak p$, and assign to each
$\mathfrak p\in S$ a positive integer $m_{\mathfrak p}$.
Suppose that $K_i/K$ are finite Galois extensions unramified at all
finite primes, $d_i=[K_i:K]\to\infty$, and the common residue
degree of primes above $\mathfrak p$ in $K_i/K$ is at most
$m_{\mathfrak p}$ for every $i$.  Assume
$r_1(K_i)/[K_i:\mathbb Q]\to\tau$ and RH for every
$\zeta_{K_i}$.  Then
\[
 \log\varrho\ge\log(8\pi)+\gamma+\frac\pi2\tau
  +\frac2n\sum_{\mathfrak p\in S}
       \frac{\log q_{\mathfrak p}}
            {q_{\mathfrak p}^{m_{\mathfrak p}/2}-1}.
\tag{NT-H09}
\]
Thus $M$ supplies the needed zero hypotheses.  No nesting of the
extensions is required.  Complete splitting is the case
$m_{\mathfrak p}=1$ and does not require the rational prime below
$\mathfrak p$ to split in $K$.

Here is the full proof.  Pass to strictly increasing degrees so that
the fields are pairwise nonisomorphic.  Finite unramifiedness gives
$\Delta_{K_i}=\Delta_K^{d_i}$ and
$g_i=\tfrac12 d_i\log\Delta_K$.
There are finitely many possible residue degrees at each point of
the finite set $S$, so pass to a subsequence on which these degrees
are fixed, say $f_{\mathfrak p}\le m_{\mathfrak p}$.
As in NT-H07, boundedness of $N_Q(K_i)/g_i$ and a diagonal
subsequence over rational prime powers $Q$ give an asymptotically
exact family.  Neither the limiting signature nor the displayed
residue degrees change.  Since the extension is Galois and
unramified, precisely $d_i/f_{\mathfrak p}$ primes lie above
$\mathfrak p$, each of norm $q_{\mathfrak p}^{f_{\mathfrak p}}$.
Different base prime ideals give disjoint sets of primes upstairs,
even when their absolute norms coincide.  Therefore
\[
 \phi_Q\ge\frac2{\log\Delta_K}
       \sum_{\substack{\mathfrak p\in S\\
                    q_{\mathfrak p}^{f_{\mathfrak p}}=Q}}
           \frac1{f_{\mathfrak p}}.
\]
The archimedean invariants are
$\phi_{\mathbb R}=2n\tau/\log\Delta_K$ and
$\phi_{\mathbb C}=n(1-\tau)/\log\Delta_K$.
Substitute these lower bounds in the TV inequality (NT-basic),
discard the other nonnegative terms, and multiply by
$\log\Delta_K/n$.  Since
$\log(q_{\mathfrak p}^{f_{\mathfrak p}})/f_{\mathfrak p}
=\log q_{\mathfrak p}$, the finite-place contribution is at least
\[
 \frac2n\sum_{\mathfrak p\in S}
      \frac{\log q_{\mathfrak p}}
           {q_{\mathfrak p}^{f_{\mathfrak p}/2}-1}
 \ge\frac2n\sum_{\mathfrak p\in S}
      \frac{\log q_{\mathfrak p}}
           {q_{\mathfrak p}^{m_{\mathfrak p}/2}-1}.
\]
The archimedean part simplifies using
$2\alpha_{\mathbb R}-\alpha_{\mathbb C}=\pi/2$ exactly as in
NT-H07.  This proves (NT-H09), including the possibility of
coincident prime-power norms.

\paragraph{NT-H10. A finite quotient from bounded-order Frobenius relations.}
Let
$G_K=\pi_1^{\text{ét}}(\operatorname{Spec}\mathcal O_K)$, and
for the data $S,\boldsymbol m$ above let $N$ be the closed normal
subgroup generated by
$\operatorname{Frob}_{\mathfrak p}^{m_{\mathfrak p}}$ for
$\mathfrak p\in S$.  The normal subgroup is independent of the
choices inside the Frobenius conjugacy classes.  Put
\[
 C_K(S,\boldsymbol m)=8\pi\exp\left(
       \gamma+\frac2n\sum_{\mathfrak p\in S}
       \frac{\log q_{\mathfrak p}}
            {q_{\mathfrak p}^{m_{\mathfrak p}/2}-1}\right).
\]
Under $M$, if $\varrho<C_K(S,\boldsymbol m)$, then $G_K/N$ is
finite; equivalently, $N$ is open.  For totally real $K$, first
quotient $G_K$ by the closed normal subgroup generated by complex
conjugations.  In that totally real quotient the weaker condition
$\varrho<e^{\pi/2}C_K(S,\boldsymbol m)$ suffices.  No equality
case is asserted.

To prove this, suppose the stated quotient is infinite.  Its finite
Galois subextensions $K_i/K$ can then be chosen with degrees tending
to infinity.  They are unramified at every finite prime.  In each
finite Galois group the order of
$\operatorname{Frob}_{\mathfrak p}$ divides
$m_{\mathfrak p}$, and that order equals its residue degree.
Choose a subsequence whose signature ratios tend to
$\tau\in[0,1]$.  NT-H09 gives
$\varrho\ge e^{\pi\tau/2}C_K(S,\boldsymbol m)
\ge C_K(S,\boldsymbol m)$, contradicting the strict inequality.
In the totally real quotient every $K_i$ is totally real, so
$\tau=1$ and the second conclusion follows.  If $K=\mathbb Q$,
Minkowski's discriminant bound already makes the whole unramified
group trivial, so this omitted case satisfies the conclusion
unconditionally.  This completes the proof.

For $m_{\mathfrak p}=1$ this is the unrestricted closed-point
version of NT-H08.  If $S$ consists of all primes above rational
primes that split completely in $K$, then each such rational prime
contributes $n$ equal summands and NT-H08's constant is recovered.
For larger $m_{\mathfrak p}$ the imposed relations only bound the
orders of Frobenius elements.  The conclusion concerns the index
of their closed normal closure; it does not assert that this normal
subgroup is generated by finitely many elements as an ordinary
topological group.  This is a quantitative Galois-category
consequence of the TV explicit-formula inequality, with its precise
arithmetic input recorded, rather than a general anabelian
reconstruction statement.

\subsection{Symmetrization removes the infinity-type restriction for zero bounds}

\paragraph{NT-Z17. Uniform zero bounds for arbitrary unitary Hecke infinity type.}
Let $\eta$ be a primitive unitary Hecke character of a number field
$K$, $n=[K:\mathbb Q]$, which is not a pure norm character.
Write its standard completion as
\[
 \Lambda(s,\eta)=Q_\eta^{s/2}
         \prod_{j=1}^n\Gamma_{\mathbb R}(s+\mu_j)L(s,\eta),
 \qquad Q_\eta=\Delta_K\operatorname N\mathfrak f_\eta,
\]
up to a harmless nonzero constant.  Unitarity gives
$\operatorname{Re}\mu_j\ge0$, but the $\mu_j$ need not occur
in conjugate pairs.  For any real $T$, put
\[
 \mathscr C_\eta(T)=Q_\eta\prod_{j=1}^n(|\mu_j+iT|+3),
 \qquad
 R_\eta(T)=\log\log\bigl(\mathscr C_\eta(T)^{3/n}\bigr).
\]
Under GRH for $L(s,\eta)$, in particular under $M$,
\[
 \operatorname{ord}_{s=1/2+iT}L(s,\eta)
 \le\frac{\log\mathscr C_\eta(T)}{2R_\eta(T)}
  +O\left(\frac{\log\mathscr C_\eta(T)\log R_\eta(T)}
                 {R_\eta(T)^2}\right),
\tag{NT-Z17a}
\]
with an absolute implied constant.  Also, as
$\mathscr C_\eta(T)^{3/n}\to\infty$, there is a nontrivial zero
$1/2+i\gamma$ of this same $L$-function such that
\[
 |\gamma-T|\le\frac{\pi}{2R_\eta(T)}
                 +O\left(\frac{\log R_\eta(T)}{R_\eta(T)^2}\right),
\tag{NT-Z17b}
\]
again with an absolute implied constant.  The degree-normalized
conductor is retained, so these are uniform formulations even when
$K,\eta,T$, and $n$ vary within the stated class.

Here is the complete transfer proof.  Define the product and its
completion by
\[
 F_T(s)=L(s+iT,\eta)L(s-iT,\overline\eta),\qquad
 \mathcal F_T(s)=\Lambda(s+iT,\eta)
                         \Lambda(s-iT,\overline\eta).
\]
Classical Hecke theory makes both completed factors entire of order
one precisely because the pure norm exception has been excluded.
The gamma parameters of the product are
$\mu_j+iT$ and $\overline\mu_j-iT$, in conjugate pairs.
The factors $Q_\eta^{iT/2}$ and $Q_\eta^{-iT/2}$ cancel, so its
finite conductor is $Q_\eta^2$ and its degree is $2n$.
After grouping prime ideals by rational primes, the Euler parameters
of either shifted factor have modulus at most one.  Thus the product
has the parameter $\vartheta=0$ in CCM2015, \S4.  Conjugation and
the Hecke functional equation show that it is self-dual with root
number one.  GRH for $\eta$ implies GRH for its conjugate and for
both shifts, hence for the product.  It therefore satisfies all
the hypotheses of CCM2015 Theorems~5 and 7, including the actual
gamma-pairing condition.

Its analytic conductor in that source is exactly
\[
 C(\mathcal F_T)=Q_\eta^2
       \prod_{j=1}^n(|\mu_j+iT|+3)^2
       =\mathscr C_\eta(T)^2.
\]
Consequently its denominator is
\[
 \log\log\bigl(C(\mathcal F_T)^{3/(2n)}\bigr)
                  =R_\eta(T),
\]
with no alteration of degree dependence.
Conjugation gives
$\operatorname{ord}_{1/2-iT}L(s,\overline\eta)
=\operatorname{ord}_{1/2+iT}L(s,\eta)$.
The gamma factors are finite and nonzero at the points under
consideration, since their arguments have positive real part.
Thus the central order of $\mathcal F_T$ is twice the order on
the left of (NT-Z17a).  Apply CCM Theorem~5, equation (4.5): the
factor two in $\log C(\mathcal F_T)$ is cancelled by division
by this doubled order.  This gives (NT-Z17a), with a universal
error constant as stated.

By CCM Theorem~7 there is a zero $1/2+iu$ of $\mathcal F_T$
with the bound on $|u|$ appearing in (NT-Z17b).  If its first
factor vanishes, $L(s,\eta)$ has ordinate $T+u$.
If its second factor vanishes, conjugating that zero gives an
ordinate $T-u$ for $L(s,\eta)$.  Either case proves (NT-Z17b).
This argument also permits $u=0$, as it should.

For each fixed non-pure-norm $K,\eta$,
$\log\mathscr C_\eta(T)=n\log|T|+O_{K,\eta}(1)$ as
$|T|\to\infty$.  Substitution in the two bounds gives exactly
the multiplicity and centered covering-radius estimates of
NT-Z14 for this arbitrary infinity type.  Applying the covering
bound at the midpoint of a large gap first shows that its length
is less than one, and then gives
\[
 \gamma'-\gamma\le\frac{\pi}{\log\log\gamma}
    +O_{K,\eta}\left(
        \frac{\log\log\log\gamma}{(\log\log\gamma)^2}\right)
\]
for consecutive large positive ordinates.  If $\eta$ is a pure
norm character, its $L$-function is $\zeta_K(s+i\tau)$ for a
fixed real $\tau$.  Translate the finite-order trivial-character
case of NT-Z14 to obtain the same fixed-character asymptotic
conclusions.  We do not apply CCM verbatim to its two shifted
completed zeta factors: they have poles away from $0,1$ and would
violate that theorem's stated hypotheses.

This proves a zero-divisor transfer, not an individual argument
bound for $S_\eta(T)$ with unpaired gamma parameters.  The product
is the conjugate-dual pair of a rank-one arithmetic object and its
norm twist.  It provides a precise analytic use of that symmetry
without presupposing an algebraic motive for the infinite-order
character; compare root MT-H01 and the geometry programme's
self-dual determinant constructions.
When $L(s,\eta)=L(s+i\tau,\chi)$ for a primitive Dirichlet
character $\chi$, $A$ already supplies the needed zero-location
input by translation.

\paragraph{NT-Z18. At least half of the prime-conductor character family is nonvanishing.}
Let $q\to\infty$ run over primes and average over all $q-2$
nonprincipal characters modulo $q$.  Under $A$,
\[
 \frac1{q-2}\sum_{\chi\ne\chi_0}
             \operatorname{ord}_{s=1/2}L(s,\chi)
       \le\frac12+O\left(\frac1{\log q}\right),
\]
and therefore
\[
 \#\{\chi\ne\chi_0\pmod q:L(1/2,\chi)\ne0\}
       \ge\left(\frac12-O\left(\frac1{\log q}\right)\right)(q-2).
\tag{NT-Z18}
\]
For the full proof take
$f(x)=(\sin(2\pi x)/(2\pi x))^2$, with $f(0)=1$.
This is real, even, nonnegative, entire, and decays as
$O((1+|x|)^{-2})$ on the real line.  Its Fourier transform is
$\widehat f(u)=\tfrac12(1-|u|/2)_+$, supported on $[-2,2]$,
and $\int f=1/2$.  It meets Hughes--Rudnick2003 Definition~2.1,
author p.~4, including the analytic-strip condition in \S2.1.
Their Theorem~3.1, p.~5, permits this closed support and gives
average $W_f=1/2+O(1/\log q)$.
Under $A$ all the scaled ordinates in $W_f$ are real, so every
summand is nonnegative; the central zeros alone contribute their
multiplicity because $f(0)=1$.  This proves the mean-order bound.
The indicator of central vanishing is at most that multiplicity,
which proves (NT-Z18) after summing.  No parity assumption on
central multiplicity is used for this non-self-dual family.

The exact family and the positivity use of $A$ distinguish this
unitary-family corollary from the quadratic-family statement
NT-Z12.  In geometric language it concerns the central zero divisor
of all nontrivial rank-one objects of the prime cyclotomic cover;
it does not assert nonvanishing for each object.

\subsection{Classwise prime counts and the average conductor}

\paragraph{NT-C06. Explicit prime counts in a ray class.}
Let $K$ have degree $n$ and absolute discriminant $\Delta_K$.
Let $\mathfrak m=\mathfrak m_0\mathfrak m_\infty$ be a modulus,
$G=\operatorname{Cl}_{\mathfrak m}(K)$ its ray class group, and
$h=|G|$.  Write $H/K$ for its ray class field.  For
$\chi\in\widehat G$, let $\mathfrak f_\chi$ be the finite part
of the conductor of its primitive Hecke character.  Define
\[
 a_{\mathfrak m}=\log\Delta_K+
           \frac1h\sum_{\chi\in\widehat G}
                    \log\operatorname N\mathfrak f_\chi.
\]
The following unconditional relations will be used:
\[
 a_{\mathfrak m}=\frac{\log\Delta_H}{h},\qquad
 0\le a_{\mathfrak m}\le
 \log\Delta_K+\left(1-\frac1h\right)
                     \log\operatorname N\mathfrak m_0.
\tag{NT-ray-conductor}
\]
The lower inequality is only a convenient bound; the exact expression
retains the mean of the individual conductor logarithms.

For a class $c\in G$, put
\begin{align*}
 \psi_c^{\mathfrak m}(x)
 &=\sum_{\substack{\mathfrak p\nmid\mathfrak m_0,\ j\ge1\\
           (\operatorname N\mathfrak p)^j\le x,\ [\mathfrak p]^j=c}}
                  \log\operatorname N\mathfrak p,\\
 \pi_c^{\mathfrak m}(x)
 &=\#\{\mathfrak p\nmid\mathfrak m_0:
             \operatorname N\mathfrak p\le x,\ [\mathfrak p]=c\}.
\end{align*}
Let $e_{\mathfrak m}$ be the number of prime ideals dividing
$\mathfrak m_0$ which are unramified in $H/K$; in particular,
$e_{\mathfrak m}\le\omega_K(\mathfrak m_0)$.
Under RH for $\zeta_H$, and hence under $M$, one has
\begin{align*}
 \left|\psi_c^{\mathfrak m}(x)-\frac{x}{h}\right|
 &\le\sqrt x\left[
       \left(\frac{\log x}{2\pi}+2\right)a_{\mathfrak m}
        +\left(\frac{\log^2x}{8\pi}+2\right)n\right]
       +e_{\mathfrak m}\log x,
       &&x\ge1,                                      \tag{NT-C06a}\\
 \left|\pi_c^{\mathfrak m}(x)-\frac{\operatorname{Li}_2(x)}h\right|
 &\le\sqrt x\left[
       \left(\frac1{2\pi}+\frac3{\log x}\right)a_{\mathfrak m}
       +\left(\frac{\log x}{8\pi}+\frac1{4\pi}
                         +\frac6{\log x}\right)n\right]
       +e_{\mathfrak m},
       &&x\ge2.                                      \tag{NT-C06b}
\end{align*}
Here $\operatorname{Li}_2(x)=\int_2^xdt/\log t$.
Both estimates are uniform in the class $c$.
Their errors contain no multiplicative ray class number: its role
is retained in the main term and in the average conductor.
Replacing $a_{\mathfrak m}$ by the upper bound in
(NT-ray-conductor) gives a version using only
$K,\mathfrak m_0,h$.

We give the full deduction, with the exact classical inputs identified.
Class-field reciprocity and existence identify $G$ with
$\operatorname{Gal}(H/K)$ and the class of each
$\mathfrak p\nmid\mathfrak m_0$ with its Frobenius;
see Milne, \emph{Class Field Theory}, v4.03, V, Theorems~3.5--3.6,
printed p.~158/PDF p.~167.  The finite conductors of the associated
characters divide $\mathfrak m_0$, and the trivial character has
conductor one.  The conductor--discriminant formula gives
\[
 \Delta_H=\Delta_K^h
              \prod_{\chi\in\widehat G}
                      \operatorname N\mathfrak f_\chi.
\]
Its proof locator is Neukirch, \emph{Algebraic Number Theory},
VII, (11.7)--(11.9), printed pp.~533--534/PDF pp.~541--542;
(11.10), printed p.~535/PDF p.~543, identifies the degree-one
Artin and class-field conductors.  The complete supplied proof of
the conductor--discriminant formula on those pages was read on
the original images.  Taking logarithms proves the equality in
(NT-ray-conductor).  Bounding each of the $h-1$ nontrivial
conductors by $\mathfrak m_0$ proves its upper bound.

Every conjugacy class of $\operatorname{Gal}(H/K)$ is a singleton,
so its density is $1/h$.  Moreover $[H:\mathbb Q]=nh$.
Apply NT-C01/C02, the precise Greni\'e--Molteni2019
Theorem~1.1 and Corollary~1.2, to that singleton.  Dividing their
errors by $h$ replaces $\log\Delta_H$ by
$a_{\mathfrak m}$ and $[H:\mathbb Q]$ by $n$, giving exactly
the bracketed expressions above.

Those source counts include every prime unramified in $H/K$.
All primes outside $\mathfrak m_0$ are unramified, but a prime
inside $\mathfrak m_0$ may be unramified as well.  Each such prime
contributes at most
\[
 \sum_{\substack{j\ge1\\(\operatorname N\mathfrak p)^j\le x}}
          \log\operatorname N\mathfrak p\le\log x
\]
to a prime-power count, and at most one to an ordinary prime
count.  Removing these $e_{\mathfrak m}$ possible primes and
using the triangle inequality proves (NT-C06a)--(NT-C06b), including
degenerate moduli.  This completes the proof.  For example, the
prime $2$ is unramified in the ray class field $\mathbb Q$ for
modulus $2\infty$ and is correctly removed by this boundary term.

The input is Dedekind RH for this one field $H$; no Artin
holomorphy conjecture is needed.  When $H/\mathbb Q$ is abelian,
$A$ suffices by Dirichlet factorization.  In particular, $A$
supplies every base-field-$\mathbb Q$ ray-class specialization.
For general $K$, abelianity of $K/\mathbb Q$ alone does not assert
abelianity of $H/\mathbb Q$ and is not substituted for this input.
Geometrically, the count takes place in the finite abelian
quotient of $\pi_1^{\text{ét}}(\operatorname{Spec}
\mathcal O_K[1/\mathfrak m_0])$ specified by the modulus and the
real-place conditions.  The quantity $a_{\mathfrak m}$ is the
average Artin conductor cost of its rank-one objects.

\paragraph{NT-C07. Ideal classes and degree-one prime representatives.}
Let $h_K=|\operatorname{Cl}(K)|$ and assume RH for
$\zeta_{\mathbb Q}$ and the Hilbert class-field Dedekind function.  Then every ordinary ideal class
contains a prime ideal of residue degree one over $\mathbb Q$
whose norm is at most
\[
 \left[h_K\bigl(4\log\Delta_K+2.5n\bigr)+5\right]^2.
\tag{NT-C07}
\]
For the narrow ideal class group the same conclusion holds with
$h_K^+$ in place of $h_K$, assuming RH for the corresponding
narrow class field.  $M$ supplies either hypothesis.
For either group, the classwise estimates (NT-C06a)--(NT-C06b)
hold with $a_{\mathfrak m}=\log\Delta_K$ and
$e_{\mathfrak m}=0$.

Indeed, use the modulus with finite part one and with empty real
part for the ordinary group, or all real places for the narrow
group.  Its class field $H/K$ is unramified at all finite primes;
consequently $\Delta_H=\Delta_K^{[H:K]}$.
The prime-count assertions follow at once from NT-C06.
For $K\ne\mathbb Q$ one has $H\ne\mathbb Q$, so
Bach--Sorenson1996 Theorem~5.1, printed p.~1729/PDF p.~13,
applies to each singleton of $\operatorname{Gal}(H/K)$.
Its degree-one prime has norm at most
$(4\log\Delta_H+2.5[H:\mathbb Q]+5)^2$.
Substitute $[H:K]=h_K$ or $h_K^+$ and the discriminant equality
to obtain (NT-C07).  The source requires RH for $\zeta_H$
and $\zeta_{\mathbb Q}$; $M$ provides both, and $A$ already
provides the latter.  If $K=\mathbb Q$, both groups are trivial
and the prime $2$ directly proves the displayed bound, without
applying the source theorem to its excluded field.
This proves the assertion in all cases.

The ordinary class group is
$\operatorname{Pic}(\operatorname{Spec}\mathcal O_K)$; the narrow
version adds the real sign conditions.  These are prime
representatives of each individual class, distinct from a small
set of primes generating the whole group.  No arbitrary extra
forbidden-prime set is included in (NT-C07); that additional
condition can instead be tested by the explicit count in NT-C06.

\paragraph{NT-C08. Certified ray-class Mertens constants and Euler products.}
Retain the notation and exact RH hypothesis of NT-C06, and let
$\delta=1/h$.  For each ray class $c$ there are real constants
$B_c$ and $A_c>0$ such that, for every $x\ge2$,
\begin{align*}
 \left|\sum_{\substack{\mathfrak p\nmid\mathfrak m_0\\
          [\mathfrak p]=c,\ \operatorname N\mathfrak p\le x}}
        \frac1{\operatorname N\mathfrak p}
        -\delta\log\log x-B_c\right|
       &\le U_{\mathfrak m}(x),                       \tag{NT-C08a}\\
 \left|\log\prod_{\substack{\mathfrak p\nmid\mathfrak m_0\\
          [\mathfrak p]=c,\ \operatorname N\mathfrak p\le x}}
       (1-(\operatorname N\mathfrak p)^{-1})^{-1}
       -\delta\log\log x-\log A_c\right|
       &\le U_{\mathfrak m}(x)+\frac{n}{2\lfloor x\rfloor},
                                                        \tag{NT-C08b}
\end{align*}
where the explicit error is
\[
 U_{\mathfrak m}(x)=\frac1{\sqrt x}
 \left[
   \left(\frac3{2\pi}+\frac9{\log x}\right)a_{\mathfrak m}
   +\left(\frac{3\log x}{8\pi}+\frac5{4\pi}
                    +\frac{18}{\log x}\right)n\right]
       +\frac{2e_{\mathfrak m}}x.
\tag{NT-C08-error}
\]
In particular the product divided by $(\log x)^{1/h}$ tends
to $A_c$.  The two finite expressions in (NT-C08a)--(NT-C08b),
with these errors, give certified intervals for $B_c$ and
$\log A_c$ once the relevant primes have been enumerated.
No uncomputed numerical value of either constant is asserted.

Here is the complete proof.  Abbreviate
$\pi_c^{\mathfrak m}(t)=\delta\operatorname{Li}_2(t)+R_c(t)$,
and let $S_c(x)$ be the reciprocal-prime sum in (NT-C08a).
Partial summation, including possible primes of norm $2$, gives
\[
 S_c(x)=\frac{\pi_c^{\mathfrak m}(x)}x+
                  \int_2^x\frac{\pi_c^{\mathfrak m}(t)}{t^2}\,dt.
\]
The lower boundary uses $\pi_c^{\mathfrak m}(2^-)=0$;
at $x=2$ the displayed formula is still valid.
Since $\operatorname{Li}_2(2)=0$, integration by parts yields
\[
 \frac{\operatorname{Li}_2(x)}x+
          \int_2^x\frac{\operatorname{Li}_2(t)}{t^2}\,dt
               =\log\log x-\log\log2.
\]
The estimate in NT-C06 makes $R_c(t)/t^2$ absolutely integrable
on $[2,\infty)$.  Define
\[
 B_c=-\delta\log\log2+
                       \int_2^\infty\frac{R_c(t)}{t^2}\,dt.
\]
These identities give the exact error formula
\[
 S_c(x)-\delta\log\log x-B_c
          =\frac{R_c(x)}x-
                        \int_x^\infty\frac{R_c(t)}{t^2}\,dt.
\tag{NT-C08-tail}
\]
Apply NT-C06b to both terms.  The elementary integrals needed are
\[
 \int_x^\infty t^{-3/2}\,dt=\frac2{\sqrt x},\qquad
 \int_x^\infty t^{-3/2}\log t\,dt
                      =\frac{2\log x+4}{\sqrt x},\qquad
 \int_x^\infty\frac{t^{-3/2}}{\log t}\,dt
                      \le\frac2{\sqrt x\log x}.
\]
For the $a_{\mathfrak m}$ terms the point and integral bounds
sum to $3/(2\pi)+9/\log x$ after division by $\sqrt x$.
For the $n$ terms they sum to
$3\log x/(8\pi)+5/(4\pi)+18/\log x$.
The boundary-prime contribution is
$e_{\mathfrak m}/x+\int_x^\infty e_{\mathfrak m}t^{-2}dt
=2e_{\mathfrak m}/x$.  This proves (NT-C08a) with exactly
the error displayed.

To pass to the product, define the absolutely convergent series
\[
 D_c=\sum_{\substack{\mathfrak p\nmid\mathfrak m_0\\
                         [\mathfrak p]=c}}
                 \sum_{k\ge2}\frac1{k(\operatorname N\mathfrak p)^k},
 \qquad A_c=\exp(B_c+D_c).
\]
For $q\ge2$,
$\sum_{k\ge2}1/(kq^k)\le1/(2q(q-1))$.
There are at most $n$ prime ideals of $K$ of any given
integer norm: if that norm is a rational prime power $p^f$,
there are at most $n/f\le n$ primes of that residue degree,
by the degree formula over $p$, and otherwise there are none.
Consequently
\[
 0\le\sum_{\substack{\mathfrak p\nmid\mathfrak m_0\\
              [\mathfrak p]=c,\ \operatorname N\mathfrak p>x}}
            \sum_{k\ge2}\frac1{k(\operatorname N\mathfrak p)^k}
 \le\frac n2\sum_{j>\lfloor x\rfloor}\frac1{j(j-1)}
 =\frac{n}{2\lfloor x\rfloor}.
\]
The logarithm of the finite Euler product is $S_c(x)$ plus
the corresponding truncated $k\ge2$ series.  Subtract
$\delta\log\log x+B_c+D_c$, use (NT-C08a) and this last
tail estimate, and obtain (NT-C08b).  Both errors tend to zero;
exponentiation proves the positive limiting constant and the
claimed interval certification.  This completes the proof.

The same deductions apply to ordinary and narrow ideal classes
with $a_{\mathfrak m}=\log\Delta_K$ and no boundary term.
For $K=\mathbb Q$ they are consequences of $A$ for every ray
class modulus.  The geometric object is the reciprocal-norm
measure on the closed points with a prescribed image in the
abelian fundamental-group quotient; its regularized Euler
product has a quantitatively controlled real constant.
This extends the finite prime counts to a convergent
regularization, without identifying that constant with a
particular motivic period.

\subsection{A conductor-uniform critical-line bound}

\paragraph{NT-L13. Uniform Lindel\"of bounds for non-pure-norm unitary Hecke characters.}
Let $\eta$ be a primitive unitary Hecke character of $K$ which is
not a pure norm character, and let $n=[K:\mathbb Q]$.
Use the completion and the sharp conductor already defined in
NT-Z17:
\[
 \mathscr C_\eta(T)=Q_\eta
                  \prod_{j=1}^n(|\mu_j+iT|+3),
 \qquad y=\log\log\mathscr C_\eta(T).
\]
Assume GRH for $L(s,\eta)$ and RH for $\zeta_{\mathbb Q}$.
Thus $M$ supplies both inputs; $A$ suffices for a norm shift
of a nonprincipal primitive Dirichlet character.
If $y\ge10$, then
\[
 |L(1/2+iT,\eta)|\le
  \exp\left[
    \left(\frac3{8y}+\frac{23n}{25y^2}\right)
                      \log\mathscr C_\eta(T)\right].
\tag{NT-L13}
\]
All constants are explicit and independent of $K,\eta,T,n$.
In particular, for every $\epsilon>0$ one has
$|L(1/2+iT,\eta)|\le\mathscr C_\eta(T)^\epsilon$ whenever
\[
 y\ge\max\left\{10,\frac3{4\epsilon},
                         \sqrt{\frac{46n}{25\epsilon}}\right\}.
\tag{NT-L13-threshold}
\]
This gives the conductor-uniform Lindel\"of statement for fixed
degree.  If the degree varies, the same conclusion in the form
$|L|\le\mathscr C^{o(1)}$ follows along any family for which
$y\to\infty$ and $n=o(y^2)$.  The dependence on degree is thus
specified rather than absorbed in an implied constant.

We verify the full specialization and the threshold deduction.
The classical input is Chandee2009, \emph{Explicit upper bounds
for $L$-functions on the critical line}, Corollary~1.1,
author p.~3, whose axioms and conductor are on p.~2,
equations (2)--(6).  Its proof is on p.~7, equations
(14)--(15); the underlying logarithmic inequality is proved
in Theorem~2.1, pp.~4--6.
Apply it to $s\mapsto L(s+iT,\eta)$.
The gamma parameters are $\mu_j+iT$ and have nonnegative
real part; this source does not impose the gamma-pairing
condition encountered in CCM2015.  Standard primitive Hecke
continuation and the functional equation give its analytic
axioms, and non-pure-norm unitary Hecke functions have no
completed pole or zero on the boundary lines.  The latter
nonvanishing is also part of the adopted GRH zero-location
input; its unconditional classical form is recorded in
Ramakrishnan--Valenza, Theorem~7-28, printed p.~289.

Grouping the prime-ideal Euler factors over each rational
prime produces $n$ parameters of modulus at most one,
padding with zero parameters at ramified primes when needed.
Indeed a factor $1-\eta(\mathfrak p)p^{-fs}$ factors into
$f$ linear factors in $p^{-s}$ whose roots have modulus one;
the degree formula gives total degree at most $n$.
The height shift multiplies those roots by $p^{-iT}$ and
preserves their moduli.  This proves the Ramanujan hypothesis
used in the corollary.  It also proves that the logarithmic
Euler coefficient at $p^k$ has modulus at most $n\log p$.
The source's proof of its explicit prime-power estimate uses
RH for $\zeta_{\mathbb Q}$, which is why that input has been
stated separately above.

In the source's normalization the conductor is
\[
 C_{\mathrm{Ch}}=
 \frac{Q_\eta}{\pi^n}
        \prod_{j=1}^n\left|\frac14+
                              \frac{\mu_j+iT}{2}\right|
 =\frac{Q_\eta}{(2\pi)^n}
        \prod_{j=1}^n|1/2+\mu_j+iT|
 \le\mathscr C_\eta(T).
\]
It is legitimate to use this larger conductor even if only
its double logarithm exceeds ten.  To see this directly in
the cited proof, retain the free parameter $x$ in equation
(14), replace $\log C_{\mathrm{Ch}}$ by
$\log\mathscr C_\eta(T)$ in its term with positive
coefficient $3/(4\log x)$, and then choose
$x=(\log\mathscr C_\eta(T))^2$.
The estimate (15) and the remaining terms depend only on
$x$ and $n$, and $\log x=2y\ge20$ meets their threshold.
The same collection of constants that proves Corollary~1.1
therefore gives exactly (NT-L13).  This explains the conductor
normalization and does not silently replace it by an
asymptotically equivalent one when the degree varies.
If the central value is zero, the nonlogarithmic bound is
immediate; otherwise the logarithmic argument applies.

Finally the two conditions
$y\ge3/(4\epsilon)$ and
$y^2\ge46n/(25\epsilon)$ make the two terms
$3/(8y)$ and $23n/(25y^2)$ at most $\epsilon/2$ each.
Exponentiation proves (NT-L13-threshold).
The varying-degree assertion follows from the same two
terms tending to zero.  This completes the deduction.

Pure norm characters have shifted completed zeta poles and
are not inserted into these pole-free axioms; their fixed-function
Lindel\"of consequences remain in NT-L03.  Geometrically,
(NT-L13) bounds the critical value of a rank-one unitary
arithmetic object uniformly in its finite and archimedean
conductors.  No algebraicity of an infinite-order infinity
type, negative-moment estimate, or nonzero lower bound at the
critical point is required or inferred.

\paragraph{NT-L14. Explicit two-sided bounds on the one-line.}
Let $\eta$ be a primitive unitary non-pure-norm Hecke character
over a degree-$n$ field, and set
\[
 C_1(\eta,T)=\frac{Q_\eta}{(2\pi)^n}
                   \prod_{j=1}^n|1+\mu_j+iT|.
\]
Here the exact normalization is retained.
For $\log C\ge23n$ put $Y=\log\log C-\log(2n)>0$ and define
\begin{align*}
 k_n&=2.31+\frac{22.59}{n}(e^{0.31n}-1-0.31n),\\
 j_n&=2+\frac{4.18}{n}(e^{0.69n}-1-0.69n),\\
 j_n'&=9+\frac{16.74}{n}(e^{0.69n}-1-0.69n),\\
 U_n(C)&=(2e^\gamma)^n
       \left(Y^n+\frac n2Y^{n-1}+\frac{nk_n}4Y^{n-2}\right),\\
 V_n(C)&=\left(\frac{12e^\gamma}{\pi^2}\right)^n
       \left(Y^n+\frac n2Y^{n-1}+\frac{nj_n}4Y^{n-2}
                         +\frac{n^2j_n'Y^n}{\log C}\right).
\end{align*}
Under GRH for $L(s,\eta)$ and RH for $\zeta_{\mathbb Q}$,
\[
  V_n(C_1(\eta,T))^{-1}
        \le |L(1+iT,\eta)|\le U_n(C_1(\eta,T))
        \qquad\text{if }\log C_1(\eta,T)\ge23n.
\tag{NT-L14}
\]
Thus $M$ supplies the hypotheses, while $A$ supplies them for
norm shifts of nonprincipal primitive Dirichlet characters.
For fixed $n$ the leading upper order is
$(2e^\gamma\log\log C)^n$, and the reciprocal leading order
is $(12e^\gamma\log\log C/\pi^2)^n$.
These are two-sided inequalities, not a claim that every value
has either extremal order.

This is Lumley2018, \emph{Explicit bounds for $L$-functions
on the edge of the critical strip}, Theorems~2.1--2.2,
author pp.~3--4, and Corollary~2.3, p.~5.
The definitions and axioms are in \S1.1, pp.~2--3;
the proofs are \S4.1, pp.~11--13, and \S4.2, pp.~14--16.
The exact constants (2.1)--(2.3) and both theorem formulas
were inspected on the original p.~4 image.
The rational-prime Euler-factor and analytic-axiom verification
is exactly the one given in NT-L13, with parameters shifted
by $iT$.  The source allows complex infinity parameters and
requires an entire completion, which explains the non-pure-norm
condition.  Its Lemmas~3.6--3.7, p.~10, explicitly use RH for
$\zeta_{\mathbb Q}$; that input is retained here.
Its quantity $J_1(n)$ is at most $j_n$, while $J_2(n)=j_n'$.
Since $Y>0$, replacing $J_1(n)$ by its displayed upper bound
preserves the reciprocal inequality.  This verifies every
specialization and normalization in (NT-L14).
The numerical constants remain the cited classical theorem;
no new numerical optimization is being claimed.

\paragraph{NT-H11. Relative class-number--regulator bounds over any base field.}
Let $E/K$ be a finite abelian extension with group $G$, and
write $n=[K:\mathbb Q]$.
For every $\chi\in\widehat G\setminus\{1\}$, let
$L(s,\chi)$ mean the primitive Hecke function attached to that
Artin character, including its correct ramified local factors.
Let $\mathfrak f_\chi$ be its finite conductor and let
$a_\chi$ count the real places of $K$ where its sign is odd.
Its conductor for NT-L14 at height zero is explicitly
\[
 C_\chi=
 \frac{\Delta_K\operatorname N\mathfrak f_\chi}
                  {(2\pi)^n}\,2^{a_\chi+r_2(K)}.
\tag{NT-relative-conductor}
\]
For a number field $F$, let $h_F$ be the ordinary ideal class
number, $R_F$ its full unit regulator, $w_F$ its number of roots
of unity, and $r_1(F),r_2(F)$ its signature.  Put
\[
 D(E/K)=\frac{w_E}{w_K}
       \sqrt{\frac{\Delta_E}{\Delta_K}}\,
       2^{r_1(K)-r_1(E)}(2\pi)^{r_2(K)-r_2(E)}.
\]
Assume RH for $\zeta_E$ and $\zeta_{\mathbb Q}$; in particular,
$M$ suffices.  If every nontrivial character satisfies
$\log C_\chi\ge23n$, then
\[
 D(E/K)\prod_{\chi\ne1}V_n(C_\chi)^{-1}
 \le\frac{h_ER_E}{h_KR_K}
 \le D(E/K)\prod_{\chi\ne1}U_n(C_\chi).
\tag{NT-H11}
\]
If some conductors fall below this threshold, let
$\mathcal S$ be precisely that finite set of characters and
put $F_{\mathrm{small}}=\prod_{\chi\in\mathcal S}|L(1,\chi)|$.
The valid bounds are then (NT-H11) with the products restricted
to $\chi\notin\mathcal S\cup\{1\}$ and both endpoints
multiplied by $F_{\mathrm{small}}$.
Thus no small-conductor value has been silently set to one.
A bound on $h_E/h_K$ alone additionally requires information
on $R_E/R_K$.

Here is the complete proof, including the ramified factors.
At a prime $\mathfrak p$ of $K$ put $q=\operatorname N\mathfrak p$.
Let $I\subset G$ be its inertia group, of order $e$, and let
$g\in G/I$ be Frobenius of order $f$.
There are $|G|/(ef)$ primes of $E$ above $\mathfrak p$,
all of norm $q^f$; the local Dedekind factor of $E$ is therefore
$(1-q^{-fs})^{-|G|/(ef)}$.
Only the characters trivial on $I$ contribute a local Hecke
factor, namely $(1-\chi(g)q^{-s})^{-1}$.
On restricting the characters of $G/I$ to its cyclic subgroup
$\langle g\rangle$, each of the $f$ roots of unity occurs
$|G|/(ef)$ times.  Consequently the product of these local
Hecke factors is the same Dedekind factor.
This proves on the absolutely convergent half-plane, and hence
by continuation,
\[
 \zeta_E(s)=\zeta_K(s)\prod_{\chi\ne1}L(s,\chi).
\tag{NT-relative-factorization}
\]
This uses the classical class-field identification of finite-order
Hecke characters with the characters of $G$, as in Milne V,
Theorems~3.5--3.6; the local factor computation above proves
the actual Euler-product identity used here.

All the nontrivial Hecke factors are entire, and the only pole
of $\zeta_K$ is at $1$.  A nontrivial zero off the critical
line of any factor would therefore be a nontrivial zero of
$\zeta_E$.  Thus RH for $\zeta_E$ supplies the Hecke GRH
hypotheses required by NT-L14.  The remaining RH for
$\zeta_{\mathbb Q}$ has been stated explicitly.
At each real place the gamma parameter of $\chi$ is $0$ or $1$;
at a complex place the two real-gamma parameters are $0,1$.
Substitution into the definition of $C_1$ proves
(NT-relative-conductor), with no asymptotic normalization.

Let $\kappa_F=\operatorname{Res}_{s=1}\zeta_F(s)>0$.
Taking residues in (NT-relative-factorization) and absolute
values gives
\[
 \frac{\kappa_E}{\kappa_K}
                  =\prod_{\chi\ne1}|L(1,\chi)|.
\]
The analytic class-number formula in Milne VI,
Corollary~2.12 and Example~2.13, printed p.~191/PDF p.~200,
is
\[
 \kappa_F=\frac{2^{r_1(F)}(2\pi)^{r_2(F)}h_FR_F}
                         {w_F\sqrt{\Delta_F}}.
\]
Its ratio therefore yields the exact identity
\[
 \frac{h_ER_E}{h_KR_K}
             =D(E/K)\prod_{\chi\ne1}|L(1,\chi)|.
\]
Apply NT-L14 to each factor above its threshold and leave the
remaining factors equal to their actual values.  This proves
both (NT-H11) and the small-conductor version.
Empty products give the correct identity when $E=K$.
The proof is complete.

For $E/\mathbb Q$ abelian, $A$ supplies RH for $\zeta_E$,
recovering and sharpening the characterwise framework of NT-H04.
For a general abelian extension $E/K$, the needed Dedekind
input is genuinely that of $E$, not just that of $K$.
The geometric comparison is between the Picard groups of the
two arithmetic curves and the rank-one character decomposition
of their finite abelian cover.  Its regulator and real-place
factors are integral parts of the comparison, not omitted
normalization terms.

\subsection*{NT-L15. A mean-square criterion for simple Hecke zeros and reciprocal derivatives}
Let $\eta$ be a fixed primitive unitary Hecke character of a number field $K$,
including a pure norm character. Define coefficients $b_\eta(\mathfrak a)$ by
\[
 L(s,\eta)^{-1}=\sum_{\mathfrak a\ne0}
       b_\eta(\mathfrak a)(N\mathfrak a)^{-s}
                    \qquad(\Re s>1),
 \qquad
 \mathcal M_\eta(x)=\sum_{N\mathfrak a\le x}b_\eta(\mathfrak a).
\]
Thus $b_\eta(\mathfrak a)=\mu_K(\mathfrak a)\eta(\mathfrak a)$
on ideals prime to the conductor and is zero on the other ideals.
Assume the additional weak mean-square hypothesis
\[
 B_\eta:=\limsup_{X\to\infty}\frac1{\log X}
       \int_1^X\left|\frac{\mathcal M_\eta(x)}x\right|^2dx<\infty.
                                                        \tag{NT-L15a}
\]
Then GRH holds for $L(s,\eta)$, every nontrivial zero is simple, and
\[
 \sum_{\rho=1/2+i\gamma}\frac1{|\rho L'(\rho,\eta)|^2}
                          \le B_\eta.                   \tag{NT-L15b}
\]
In particular $B_\eta>0$ and
\[
 |L'(\rho,\eta)|^{-1}\le\sqrt{B_\eta}\,|\rho|,
 \qquad |L'(\rho,\eta)|^{-1}=o(|\rho|),
                                                        \tag{NT-L15c}
\]
where the little-$o$ is along zeros of unbounded height. Also
\[
 \sum_{|\gamma|\le T}|L'(1/2+i\gamma,\eta)|^{-2}=o(T^2).
                                                        \tag{NT-L15d}
\]
No use of $A$ or $M$ is required once (NT-L15a) is assumed. Neither $A$ nor
$M$ is asserted to imply this additional mean-square hypothesis.
For $K=\mathbb Q$, $\eta=1$, the classical source is Titchmarsh,
\emph{The Theory of the Riemann Zeta-function}, 2nd ed.,
\S14.29: hypothesis (14.29.1) and Theorem14.29(A), pp.376--377;
reciprocal derivative estimate (14.29.3), p.377;
Theorem14.29(B)/(14.29.4), p.377, with proof ending p.380.
The following proof gives the Hecke formulation and its precise limsup constant.

\begin{proof}
Put $F(y)=e^{-y/2}\mathcal M_\eta(e^y)$ for $y\ge0$ and
$A(Y)=\int_0^Y|F(y)|^2dy$. Changing variables $x=e^y$ gives
$B_\eta=\limsup_{Y\to\infty}A(Y)/Y$.
For every $\epsilon>0$ there is $C_\epsilon$ such that
\[
 A(Y)\le (B_\eta+\epsilon)Y+C_\epsilon\qquad(Y\ge0).
                                                        \tag{NT-L15e}
\]
In particular the Laplace integral
\[
 H(w)=\int_0^\infty F(y)e^{-wy}\,dy
\]
is holomorphic for $\Re w>0$. To check this, integrate $|F|^2$ against
$e^{-\sigma y}$ by parts using (NT-L15e), then apply Cauchy--Schwarz to
$F(y)e^{-wy}$. The same argument with any fixed factor $y^j$ proves
uniform convergence of every derivative on each compact subset of that
half-plane. For $\Re w>1/2$, absolute convergence of the inverse Euler
series and partial summation give
\[
 H(w)=\int_1^\infty\mathcal M_\eta(x)x^{-(w+1/2)-1}\,dx
      =\frac1{(w+1/2)L(w+1/2,\eta)}.                    \tag{NT-L15f}
\]
The meromorphic identity theorem extends (NT-L15f) throughout $\Re w>0$.
Since $H$ is holomorphic there and $w+1/2\ne0$, $L(s,\eta)$ has no zeros
with $\Re s>1/2$. The conjugate character has coefficients
$b_{\bar\eta}(\mathfrak a)=\overline{b_\eta(\mathfrak a)}$, hence the
same mean-square hypothesis and the same zero-free half-plane.
The Hecke functional equation now excludes nontrivial zeros of
$L(s,\eta)$ with $\Re s<1/2$. This proves GRH, including non-real
characters; possible pure-norm poles only give zeros of $H$ and cause
no difficulty.

For $\sigma>0$ introduce the positive inner product
\[
 \langle f,g\rangle_\sigma
       =\sigma\int_0^\infty f(y)\overline{g(y)}e^{-\sigma y}dy.
\]
An integration by parts and (NT-L15e) show that
\[
 \|F\|_\sigma^2
   =\sigma^2\int_0^\infty A(y)e^{-\sigma y}dy,
 \qquad \limsup_{\sigma\downarrow0}\|F\|_\sigma^2\le B_\eta.
                                                        \tag{NT-L15g}
\]
Each $e_t(y)=e^{ity}$ has norm one. Consequently
\[
 |\sigma H(\sigma+it)|
       =|\langle F,e_t\rangle_\sigma|
       \le\|F\|_\sigma.
                                                        \tag{NT-L15h}
\]
At a zero $\rho=1/2+i\gamma$, (NT-L15f) has a pole of the same order
as that zero, since $\rho\ne0$. The bound (NT-L15h) as
$\sigma\downarrow0$ rules out order greater than one. Thus every zero
is simple and
\[
 \lim_{\sigma\downarrow0}\langle F,e_\gamma\rangle_\sigma
       =a_\gamma:=\frac1{\rho L'(\rho,\eta)}.
                                                        \tag{NT-L15i}
\]
For two distinct ordinates,
\[
 \langle e_\gamma,e_\nu\rangle_\sigma
       =\frac\sigma{\sigma-i(\gamma-\nu)}\longrightarrow0;
\]
the same inner product is one when $\gamma=\nu$.
Take an arbitrary finite set $S$ of zero ordinates and put
$P_S=\sum_{\gamma\in S}a_\gamma e_\gamma$.
Expanding the nonnegative number $\|F-P_S\|_\sigma^2$, using
(NT-L15g)--(NT-L15i), and taking the upper limit gives
\[
 0\le\limsup_{\sigma\downarrow0}\|F-P_S\|_\sigma^2
       \le B_\eta-\sum_{\gamma\in S}|a_\gamma|^2.
\]
The supremum over finite $S$ proves (NT-L15b). Nontrivial Hecke zeros
exist, by the classical zero-count asymptotic, so $B_\eta>0$.
Every summand is at most $B_\eta$, proving the first assertion of
(NT-L15c). Square summability makes $a_\gamma\to0$ along unbounded
ordinates and proves the second assertion. Finally, for $T\ge1$,
\[
 T^{-2}\sum_{|\gamma|\le T}|L'(\rho,\eta)|^{-2}
   =\sum_\gamma {\bf1}_{|\gamma|\le T}
           \frac{1/4+\gamma^2}{T^2}|a_\gamma|^2.
\]
Each factor before $|a_\gamma|^2$ tends to zero and is at most $5/4$.
Dominated convergence for the summable series proves (NT-L15d).
\end{proof}

The finite Bessel argument proves an inequality. It does not assert an
exact mean-square limit, completeness of an infinite Fourier expansion
of $F$, or a limiting distribution for $F$. Such assertions require
additional approximation estimates.
Geometrically, the coefficients are the inverse Euler weights of the
rank-one arithmetic object attached to $\eta$; squarefree ideals are
squarefree effective divisors on $\operatorname{Spec}\mathcal O_K$.
The new hypothesis is analytic control of these weighted divisors and
is not a consequence of a finite-field purity theorem established here.

\subsection*{NT-Z19. Weak mean square supplies the canonical arithmetic basis gate}
Use the canonical completed, pole-removed $\Xi_\eta$, inverse Mellin
function $\psi_\eta$, and closed quotient
\[
 Q_\eta=\mathcal E/\overline{\psi_\eta*\mathcal E},\qquad
 \mathcal E=\{f\in C^\infty(\mathbb R):
      \sup_x e^{N|x|}|f^{(j)}(x)|<\infty\text{ for all }N,j\ge0\}
\]
of MT-M12/MT-M16. Under (NT-L15a), all zeros are simple and
\[
 \overline{\psi_\eta*\mathcal E}
   =\{f:\mathcal Mf(\rho)=0\text{ for every nontrivial zero }\rho\},
 \qquad Q_\eta\simeq s(Z).
                                                        \tag{NT-Z19a}
\]
Here $\mathcal Mf(s)=\int_\mathbb Rf(x)e^{sx}dx$,
and $s(Z)$ is the space of sequences decreasing faster than every
power of $1+|\operatorname{Im}\rho|$. The quotient has an absolute
Schauder basis with these Mellin evaluations as coefficient functionals.
Convolution by every $h\in\mathcal E$ is nuclear. In the diagonal
Hilbert realization $\ell^2(Z)$ it is trace class, with Hilbert trace
\[
 \operatorname{tr}_{\ell^2(Z)}T_h=\sum_\rho\mathcal Mh(\rho).
\]
Writing $h^\#(x)=e^{-x}\overline{h(-x)}$, the operator
$T_{h*h^\#}$ is positive and
\[
 \operatorname{tr}_{\ell^2(Z)}T_{h*h^\#}
                     =\sum_\rho|\mathcal Mh(\rho)|^2\ge0.
                                                        \tag{NT-Z19b}
\]
\begin{proof}
NT-L15 proves GRH, simplicity, and
$|L'(\rho,\eta)|\ge B_\eta^{-1/2}|\rho|^{-1}$.
Since $|\rho|\le1+|\operatorname{Im}\rho|$, this is precisely the
polynomial derivative lower bound of MT-M12, or condition (1) of
MT-M16 with every multiplicity equal to one. The entire proofs of
those results give (NT-Z19a) and the absolute basis. MT-M13/MT-M17
then give nuclearity and the trace-class Hilbert realization. For
completeness, on each coordinate the convolution eigenvalue is
$\mathcal Mh(\rho)$, and direct substitution gives
$\mathcal M h^\#(s)=\overline{\mathcal Mh(1-\bar s)}$.
GRH implies $1-\bar\rho=\rho$, so the product eigenvalue is
$|\mathcal Mh(\rho)|^2$. Rapid Mellin decay and NT-Z10's local zero
count make both displayed trace series absolutely convergent. This
proves positivity and (NT-Z19b).
\end{proof}
This is a conditional analytic realization of the zero divisor with an
explicit arithmetic gate. It does not identify the quotient with an
independently constructed motivic or étale cohomology group.
For the Riemann zeta function the sole additional arithmetic condition
is the classical weak Mertens conjecture; $A$ is not needed separately.

\subsection*{NT-H12. Class-group torsion from the field's own Dedekind RH}
Fix integers $n\ge2$ and $\ell\ge2$. For every degree-$n$ number field
$K$ with absolute discriminant $D_K$ whose Dedekind zeta function satisfies
RH, and every $\epsilon>0$,
\[
 |\operatorname{Cl}_K[\ell]|
       \ll_{n,\ell,\epsilon}
 D_K^{\,1/2-1/(2\ell(n-1))+\epsilon}.                   \tag{NT-H12a}
\]
Only RH for $\zeta_K$ is needed. In particular $M$ supplies the hypothesis
conditionally, and adopted $A$ supplies it when $K/\mathbb Q$ is abelian.
No RH for a normal closure or Hilbert class field is required by the
argument below. The exponent is the classical Ellenberg--Venkatesh
bound, not a new class-group exponent. Our point is the exact
Dedekind-only input reduction.

The precise cited algebraic input is Ellenberg--Venkatesh,
\emph{Reflection principles and bounds for class group torsion},
author preprint, Lemma3 and its proof, p.6: if
$0<\delta<1/(2\ell(n-1))$ and there are $N$ prime ideals of $K$
of norm at most $D_K^\delta$, unramified over $\mathbb Q$ and not
extensions of prime ideals of any proper intermediate field, then
\[
 |\operatorname{Cl}_K[\ell]|
       \ll_{n,\ell,\epsilon,\delta}D_K^{1/2+\epsilon}/N.
                                                        \tag{NT-H12b}
\]
The source's Proposition1, p.7, applies this lemma using completely
split primes. We instead count degree-one prime ideals. The implicit
constant in (NT-H12a) is not assigned an unchecked numerical value.

\begin{proof}
Let $\pi_K(x)$ count all prime ideals of $K$ of norm at most $x$.
Apply NT-C02 to the trivial Galois extension $K/K$. Its conjugacy density
is one, there are no relative ramified primes to omit, and RH for
$\zeta_K$ gives, for $x\ge2$,
\[
 |\pi_K(x)-\operatorname{Li}_2(x)|\le B_K(x),
\]
where the completely explicit error is
\[
 B_K(x)=\sqrt x\left[
 \left(\frac1{2\pi}+\frac3{\log x}\right)\log D_K+
 \left(\frac{\log x}{8\pi}+\frac1{4\pi}+\frac6{\log x}\right)n
                  \right].                            \tag{NT-H12c}
\]
Every prime ideal of residue degree at least two above its rational
prime has underlying prime at most $\sqrt x$. There are at most $n$
prime ideals above each rational prime, so such ideals number at most
$n\pi(\sqrt x)\le n\sqrt x$. There are at most $n\omega(D_K)$
prime ideals above rational primes ramified in $K$. Therefore the
number $N_K^{(1),\mathrm{ur}}(x)$ of unramified degree-one prime ideals
satisfies the useful finite inequality
\[
 N_K^{(1),\mathrm{ur}}(x)
   \ge\operatorname{Li}_2(x)-B_K(x)-n\sqrt x-n\omega(D_K).
                                                        \tag{NT-H12d}
\]
Here $\omega$ counts distinct rational prime divisors.

These degree-one ideals meet the remaining condition in (NT-H12b).
Indeed, if a prime ideal $\mathfrak p$ of norm $p$ were
$\mathfrak q\mathcal O_K$ for a prime ideal of a proper subfield $F$,
then multiplicativity of the norm under ideal extension would give
\[
 p=N_K(\mathfrak p)=(N_F\mathfrak q)^{[K:F]}.
\]
The right side is a proper positive integral power of an integer
at least two, and cannot be the rational prime $p$.

Fix $\delta>0$ and put $x=D_K^\delta$. For large $D_K$, the bound
$\operatorname{Li}_2(x)\ge x/(2\log x)$ follows by integrating on
$[x/2,x]$. The error in (NT-H12c) is
$O_{n,\delta}(D_K^{\delta/2}\log D_K)$, and
$\omega(D_K)\le\log D_K/\log2$. Consequently
\[
 N_K^{(1),\mathrm{ur}}(D_K^\delta)
                  \gg_{n,\delta}D_K^\delta/\log D_K.   \tag{NT-H12e}
\]
Write $\alpha=1/(2\ell(n-1))$. It suffices to prove (NT-H12a) for
$0<\epsilon<\alpha$, since a result with any smaller positive epsilon
implies the assertion with a larger one. Choose
$\delta=\alpha-\epsilon/3$ and apply (NT-H12b) with its epsilon
parameter equal to $\epsilon/3$. Equations (NT-H12b) and (NT-H12e)
give
\[
 |\operatorname{Cl}_K[\ell]|
   \ll D_K^{1/2-\alpha+2\epsilon/3}\log D_K
   \ll_{n,\ell,\epsilon}D_K^{1/2-\alpha+\epsilon}.
\]
The bounded-discriminant cases are absorbed into the constant by
Hermite's finiteness theorem (Milne, \emph{Algebraic Number Theory}
v3.08, Theorem8.43, printed pp.151--152). This proves the claim.
\end{proof}
The geometry is the torsion subgroup of
$\operatorname{Pic}(\operatorname{Spec}\mathcal O_K)$.
The source's Arakelov argument incorporates all archimedean places;
our analytic step counts rational degree-one closed points without
requiring a completely split Frobenius class.

\subsection*{NT-H13. Bounds for unramified cyclic covers and Kummer torsors}
Retain $n\ge2$, $K$, and RH for $\zeta_K$ from NT-H12.
For a prime $\ell$, the number of $K$-isomorphism classes of cyclic
extensions $E/K$ of degree $\ell$, unramified at every finite place
and split at every real place, is exactly
\[
 c_\ell(K)=\frac{|\operatorname{Cl}_K[\ell]|-1}{\ell-1}
       \ll_{n,\ell,\epsilon}
       D_K^{\,1/2-1/(2\ell(n-1))+\epsilon}.             \tag{NT-H13a}
\]
For odd $\ell$ the real-place condition is automatic. For $\ell=2$
it is essential to this ordinary-class-group formula.

For any integer $\ell\ge2$, let $S$ be the set of primes of $K$
above divisors of $\ell$, let $U=\operatorname{Spec}\mathcal O_K[1/\ell]$,
and let $w_K$ be the number of roots of unity in $K$. The number of
isomorphism classes of $\mu_\ell$-torsors over $U$ is
\[
 \begin{split}
 |H^1_{\text{ét}}(U,\mu_\ell)|
  &=\ell^{r_1+r_2-1+|S|}\gcd(\ell,w_K)
                             |\operatorname{Pic}(U)[\ell]|\\
  &\le\ell^{r_1+r_2+|S|}|\operatorname{Cl}_K[\ell]|\\
  &\ll_{n,\ell,\epsilon}
                 D_K^{\,1/2-1/(2\ell(n-1))+\epsilon}.
 \end{split}                                           \tag{NT-H13b}
\]
The finite coefficient $\ell^{r_1+r_2+|S|}$ is at most
$\ell^{n(1+\omega(\ell))}$. These torsors need not be connected;
$\mu_\ell$ need not be a constant group scheme over $K$.

\begin{proof}
Global class field theory identifies the ordinary Hilbert class field
Galois group with $\operatorname{Cl}_K$, with real places split.
The convention is stated explicitly in Milne, \emph{Algebraic Number
Theory} v3.08, Remark4.11 and its footnote, printed p.72; the general
reciprocity/existence input is Milne, \emph{Class Field Theory}, V3.5--3.6.
Degree-$\ell$ cyclic subextensions correspond to kernels of surjective
homomorphisms $\operatorname{Cl}_K\to\mathbb Z/\ell\mathbb Z$.
For prime $\ell$ every nonzero homomorphism is surjective, and
\[
 |\operatorname{Hom}(\operatorname{Cl}_K,\mathbb Z/\ell\mathbb Z)|
      =|\operatorname{Cl}_K[\ell]|.
\]
Exactly $\ell-1$ such homomorphisms have each kernel, by postcomposition
with automorphisms of the target. This proves the equality in
(NT-H13a), and NT-H12 gives the bound. A real-place inertia group has
order at most two, so it is trivial in an extension of odd prime degree.

Put $E_{K,S}=\mathcal O_K[1/\ell]^\times$. Since $\ell$ is invertible
on $U$, the étale Kummer sequence and
$H^1_{\text{ét}}(U,\mathbb G_m)=\operatorname{Pic}(U)$ give
\[
 0\longrightarrow
 E_{K,S}/E_{K,S}^{\ell}
 \longrightarrow H^1_{\text{ét}}(U,\mu_\ell)
 \longrightarrow\operatorname{Pic}(U)[\ell]
 \longrightarrow0.                                    \tag{NT-H13c}
\]
A precise source is Stacks Project, Section59.28, Tag03PK,
Lemma59.28.1 and its following cohomology sequence; Lemma59.28.5,
Tag040Q, identifies the classes as line bundles with an $\ell$th-power
trivialization. Thus this cohomology group is not being confused with
Picard torsion alone.

The $S$-unit theorem gives
$\mathcal O_K[1/\ell]^\times\simeq\mu(K)\times
\mathbb Z^{r_1+r_2-1+|S|}$ as abelian groups. Here $\mu(K)$ is cyclic
of order $w_K$. See Milne, \emph{Algebraic Number Theory} v3.08,
Theorem5.11 and its complete proof, printed p.90. Taking the quotient
by $\ell$th powers has cardinality
$\ell^{r_1+r_2-1+|S|}\gcd(\ell,w_K)$, proving the first equality in
(NT-H13b).

Localization makes $\operatorname{Pic}(U)$ the quotient of
$\operatorname{Cl}_K$ by the classes of primes in $S$. For any
surjection $A\to B$ of finite abelian groups, the map
$A/\ell A\to B/\ell B$ is surjective and
$|A[\ell]|=|A/\ell A|$, $|B[\ell]|=|B/\ell B|$.
It follows that $|\operatorname{Pic}(U)[\ell]|
\le|\operatorname{Cl}_K[\ell]|$; no injection on torsion induced by
localization is asserted. Apply NT-H12 and
$\gcd(\ell,w_K)\le\ell$, $r_1+r_2\le n$,
$|S|\le n\omega(\ell)$ to finish the proof.
\end{proof}
This is a literal finite étale/Kummer interpretation of a conditional
class-group bound: cyclic quotients of the arithmetic fundamental
group with the stated real conditions, and torsors under the finite
étale group scheme $\mu_\ell$ over $U$. It gives no bound for arbitrary
nonabelian covers or for all Selmer groups without further reduction.

\subsection*{NT-L16. Zero multiplicity forces logarithmic mean-square growth}
For a fixed primitive unitary Hecke function, use the inverse coefficients
and summatory function $\mathcal M_\eta$ of NT-L15, and put
\[
 I_\eta(X)=\int_1^X|\mathcal M_\eta(x)/x|^2dx.
\]
If $\rho=1/2+i\gamma$ is a zero of multiplicity $m\ge1$, write
$c_\rho=L^{(m)}(\rho,\eta)/m!$. Then, without assuming GRH,
\[
 \limsup_{X\to\infty}\frac{I_\eta(X)}{(\log X)^{2m-1}}
       \ge\frac1{(2m-1)((m-1)!)^2\,|\rho c_\rho|^2}>0.
                                                        \tag{NT-L16a}
\]
More strongly, the reciprocal leading-coefficient terms may be summed
over \emph{all distinct critical-line zeros of exact multiplicity $m$}:
\[
 \limsup_{X\to\infty}\frac{I_\eta(X)}{(\log X)^{2m-1}}
 \ge\frac1{(2m-1)((m-1)!)^2}
   \sum_{\substack{\Re\rho=1/2\\ \operatorname{ord}_\rho L=m}}
            \frac1{|\rho c_\rho|^2}.                    \tag{NT-L16a'}
\]
This includes the possibility that the sum and the limsup are infinite.
Consequently, for every $0\le a<2m-1$,
\[
 \limsup_{X\to\infty}I_\eta(X)/(\log X)^a=+\infty.
                                                        \tag{NT-L16b}
\]
If instead there is a zero with real part greater than $1/2$, then
(NT-L16b) holds for every finite real $a\ge0$.
The critical-line location of the zero in (NT-L16a) is an explicit
hypothesis; $M$ supplies it for every nontrivial zero, but is not
used in the proof of this implication.

\begin{proof}
Set $F(y)=e^{-y/2}\mathcal M_\eta(e^y)$ and
$A(Y)=\int_0^Y|F(y)|^2dy=I_\eta(e^Y)$.
Let $b=2m-1$ and
$B=\limsup_{Y\to\infty}A(Y)/Y^b$.
If $B=+\infty$, both asserted lower bounds are immediate.
Otherwise, for every $\epsilon>0$,
$A(Y)\le(B+\epsilon)Y^b+C_\epsilon$.
As in NT-L15, Cauchy--Schwarz with an exponential weight gives a
holomorphic Laplace transform on $\Re w>0$, and partial summation
in its initial absolute-convergence half-plane gives
\[
 H(w)=\int_0^\infty F(y)e^{-wy}dy
      =\frac1{(w+1/2)L(w+1/2,\eta)}.
                                                        \tag{NT-L16c}
\]
The identity extends meromorphically through $\Re w>0$.
Use again the inner product
\[
 \langle f,g\rangle_\sigma=\sigma\int_0^\infty
       f(y)\overline{g(y)}e^{-\sigma y}dy.
\]
Integration by parts yields
\[
 \begin{split}
 \|F\|_\sigma^2
 &=\sigma^2\int_0^\infty A(y)e^{-\sigma y}dy\\
 &\le (B+\epsilon)\Gamma(b+1)\sigma^{1-b}
                               +C_\epsilon\sigma.
 \end{split}                                           \tag{NT-L16d}
\]
Put $k=m-1$ and use the normalized polynomial waves
\[
 u_{\gamma,\sigma}(y)=
       \frac{\sigma^k y^k e^{i\gamma y}}{\sqrt{(2k)!}},
 \qquad v_\sigma(y)=\sigma^k F(y).
\]
Their Gram matrix has entries
\[
 \langle u_{\gamma,\sigma},u_{\nu,\sigma}\rangle_\sigma
       =\left(\frac{\sigma}{\sigma-i(\gamma-\nu)}\right)^{2k+1}.
                                                        \tag{NT-L16e}
\]
In particular each vector has norm one, and for any fixed finite set
of distinct ordinates its Gram matrix tends to the identity as
$\sigma\downarrow0$. Polynomial growth of $A$ permits differentiation
under the Laplace integral on every compact subset of $\Re w>0$,
by weighted Cauchy--Schwarz. Hence
\[
 \langle v_\sigma,u_{\gamma,\sigma}\rangle_\sigma
 =\frac{\sigma^{2k+1}}{\sqrt{(2k)!}}
                 (-1)^k H^{(k)}(\sigma+i\gamma).
\]
At a zero $\rho=1/2+i\gamma$ of exact order $m=k+1$,
the meromorphic identity (NT-L16c) gives
$H(w)=(\rho c_\rho)^{-1}(w-i\gamma)^{-(k+1)}$
plus terms of smaller pole order. Differentiating this finite
principal part yields
\[
 \lim_{\sigma\downarrow0}
 \langle v_\sigma,u_{\gamma,\sigma}\rangle_\sigma
       =\frac{\sqrt{(2k)!}}{k!\,\rho c_\rho}.           \tag{NT-L16f}
\]
For a fixed finite set $S$ of the specified zeros, finite-dimensional
orthogonal projection onto the span of these waves gives the Gram
matrix version of Bessel's inequality. By (NT-L16e), letting
$\sigma\downarrow0$ in this inequality and using (NT-L16d) gives
\[
 \frac{(2k)!}{(k!)^2}\sum_{\rho\in S}|\rho c_\rho|^{-2}
   \le \limsup_{\sigma\downarrow0}\|v_\sigma\|_\sigma^2
   \le (B+\epsilon)(2k+1)!.
\]
Let $\epsilon\downarrow0$ and take the supremum over finite $S$.
Since the terms are nonnegative, this proves (NT-L16a') and
therefore (NT-L16a). For $m=1$ it is exactly the Abel--Bessel
normalization of NT-L15. The one-frequency model
$F(y)=a y^{m-1}e^{i\gamma y}/(m-1)!$ has
$A(Y)=|a|^2Y^{2m-1}/((2m-1)((m-1)!)^2)$, so this constant is
sharp for the single-pole Laplace argument; no attainment by an
arithmetic Hecke function is asserted.
The positive limsup in (NT-L16a) gives (NT-L16b) along an unbounded sequence
when $a<2m-1$.

Finally suppose a zero has real part greater than $1/2$.
A finite limsup in (NT-L16b), for any fixed $a\ge0$, would imply
$A(Y)\le C(1+Y^a)$. The same weighted Cauchy--Schwarz argument
makes $H$ holomorphic on $\Re w>0$. Equation (NT-L16c) would then
exclude every zero of $L$ in $\Re s>1/2$, a contradiction.
This proves the final assertion.
\end{proof}
This is a direct multiplicity-sensitive version of the weak-Mertens
Laplace argument. It is a necessary analytic growth condition, not
an asymptotic formula for $I_\eta(X)$.

\subsection*{NT-L17. Actual Kummer fields obstruct universal weak Mertens and full-jet positivity}
Let $r\ge3$ be an odd prime and let $p\ne r$ be a rational prime. Put
\[
 K=\mathbb Q(\zeta_r),\qquad E=\mathbb Q(\zeta_r,p^{1/r}).
\]
There is a nontrivial primitive finite-order Hecke character $\theta$
of $K$ whose entire finite-part $L$-function $H(s)=L_K(s,\theta)$
satisfies the all-prime factorization
\[
                  \zeta_E(s)=\zeta_K(s)H(s)^{r-1}.     \tag{NT-L17a}
\]
The function $H$ has infinitely many nontrivial zeros. Hence $\zeta_E$
has infinitely many nontrivial zeros of multiplicity at least $r-1$,
\emph{unconditionally}.

Let $\mathcal M_E(x)=\sum_{N\mathfrak a\le x}\mu_E(\mathfrak a)$
and $I_E(X)=\int_1^X|\mathcal M_E(x)/x|^2dx$. Then
\[
            \limsup_{X\to\infty}I_E(X)/\log X=+\infty.
                                                        \tag{NT-L17b}
\]
Thus a universal weak-Mertens hypothesis for all Dedekind zeta
functions is false. This statement makes no assertion about an
individual Dirichlet function's weak-Mertens conjecture.

If RH for $\zeta_E$ is additionally assumed, choose any zero $\rho$
of $H$. It is then on the critical line. With
$m=\operatorname{ord}_\rho\zeta_E\ge r-1$, NT-L16 gives the explicit
nonzero lower bound
\[
 \limsup_{X\to\infty}\frac{I_E(X)}{(\log X)^{2m-1}}
 \ge
 \frac1{(2m-1)((m-1)!)^2\,
             |\rho\,\zeta_E^{(m)}(\rho)/m!|^2}.
                                                        \tag{NT-L17c}
\]
In particular $\limsup I_E(X)/(\log X)^{2r-3}>0$; it is infinite
if the chosen multiplicity exceeds $r-1$. If instead $\zeta_E$
has an off-critical zero, its functional equation supplies a
right-half-strip zero, and every polynomial-logarithmic mean bound
fails by NT-L16. This gives the same positive-limsup conclusion
unconditionally by a dichotomy, while (NT-L17c) retains its stated
critical-zero/RH hypothesis.

There is also a concrete consequence for MT-M20: the canonical
\emph{full-jet} quotient for $\zeta_E$ has no positive-definite
continuous inner product for which every convolution extends boundedly
and every $C_{h*h^\sharp}$ is positive. Under $M$, the Hilbert traces
in MT-M19 are nevertheless nonnegative. Trace positivity and a faithful
positive operator realization are different assertions here.

\begin{proof}
The elementary Galois calculation in CX-ALG-024 gives
$[E:K]=r$ and
$\operatorname{Gal}(E/\mathbb Q)=\mathbb F_r\rtimes\mathbb F_r^\times$.
For clarity, $p$ is unramified in the cyclotomic field $K$, and
$X^r-p$ is Eisenstein at every prime above $p$, proving the degree.
The translations send $p^{1/r}$ to $\zeta_r^a p^{1/r}$.
The automorphism fixing $p^{1/r}$ and sending $\zeta_r$ to
$\zeta_r^b$ conjugates a translation by $a$ to one by $ba$.
Thus the $r-1$ nontrivial characters of
$\operatorname{Gal}(E/K)=\mathbb F_r$ form a single orbit under
$\operatorname{Gal}(K/\mathbb Q)$.

Class field theory identifies these characters with primitive
finite-order Hecke characters $\theta^a$ of $K$.
Conjugating ideals by an automorphism of $K/\mathbb Q$ preserves
absolute norms and carries Frobenius and inertia by conjugation.
Consequently it permutes every finite Euler factor, including the
ramified factors, and all $L_K(s,\theta^a)$ with $a\ne0$ are the
same function $H(s)$. The unconditional all-prime abelian factorization
proved within NT-H11 therefore gives (NT-L17a).
Equivalently, CX-ALG-024's permutation-character identity gives
$H=\zeta_{\mathbb Q(p^{1/r})}/\zeta_{\mathbb Q}$.
No general Artin-holomorphy conjecture is being invoked: $H$ is an
actual nontrivial Hecke $L$-function, so it is entire by the classical
Hecke theorem (Ramakrishnan--Valenza, Theorem7-19, printed p.279).

Greni\'e--Molteni, AppendixA TheoremA.1, author p.26, applies
unconditionally to primitive finite-order Hecke characters. Its
zero count has main term
$([K:\mathbb Q]/\pi)T\log T+O_{K,\theta}(T)$.
Hence $H$ has infinitely many nontrivial zeros in $0<\Re s<1$.
At each such zero $\zeta_K$ is holomorphic, because its only pole
is at $1$. Equation (NT-L17a) therefore gives multiplicity at
least $r-1$ in $\zeta_E$.
If the limsup in (NT-L17b) were finite, NT-L15 applied to the
trivial primitive Hecke character over $E$ would make all these
zeros simple, a contradiction. This proves (NT-L17b).
The quantitative claims follow exactly from NT-L16 and the
functional equation of $\zeta_E$, as described above.
Finally MT-M20 proves that a faithful positive completion with the
stated convolution properties requires simple zeros. The actual
multiple zeros just proved exclude it. MT-M19 supplies the separate
trace-positive conclusion under RH for $\zeta_E$, which $M$ supplies
conditionally.
\end{proof}
For example $r=3$, $p=2$ gives the degree-six $S_3$ field
$E=\mathbb Q(\sqrt[3]{2},\sqrt{-3})$ and at least double nontrivial
zeros. The argument given here concerns these explicitly monomial
Galois fields; this proof alone does not address every nonabelian
Galois field.
Geometrically this is the positive multiplicity of a single
rank-one Hecke realization in the zeta realization of an affine
finite Galois cover. It is compatible with GRH and with the
unconditional adapted-block synthesis of MT-M18, but prevents
a faithful positive Hilbert realization of the full jet algebra,
even when its traces are positive under $M$.


\subsection*{NT-E01. Cyclic reductions from precisely the division-field zeta functions}
Let $E/K$ be an elliptic curve over any number field, with or without
complex multiplication, and put $K_m=K(E[m])$ in a fixed algebraic closure.
Write $S_{E/K}(x)$ for the number of good prime ideals $\mathfrak p$ of
$K$ with $N\mathfrak p\le x$ for which $E_{\mathfrak p}(k_{\mathfrak p})$
is cyclic. The classical conditional density theorem gives
\[
 S_{E/K}(x)=\delta_{E/K}\frac{x}{\log x}
                 +o_{E,K}\!\left(\frac{x}{\log x}\right),\qquad
 \delta_{E/K}=\sum_{m\ge1}\frac{\mu(m)}{[K_m:K]}\ge0.       \tag{NT-E01a}
\]
The series is absolutely convergent. A sufficient analytic input is
\emph{RH for $\zeta_{K_m}$ for every squarefree $m$, including $m=1$}.
Thus $M$ supplies this consequence conditionally. Neither GRH for
$L(s,E)$ nor Artin holomorphy for the division-field Galois groups is an
additional hypothesis. No assertion of a positive density is implicit
in (NT-E01a).

The exact cited locators are Murty, \emph{On Artin's conjecture}
(1983), Theorem1, printed p.153, with its proof on pp.153--155;
and Campagna--Stevenhagen, \emph{Cyclic reduction densities for
elliptic curves} (2023), \S2, Lemma2.1, Corollary2.2 and
Propositions2.3--2.4 (author v2 pp.6--9). The latter verifies the
geometric degree growth, root-discriminant bound and large-prime
tail needed by Murty's theorem for arbitrary $E/K$.

\paragraph{Exact input identification.}
For a squarefree $m$, the compositum of the fields $K_\ell$ for
$\ell\mid m$ is $K_m$: the Chinese remainder decomposition of
$E[m]$ gives both inclusions. These are therefore exactly the fields
called $L_k$ in Murty's Theorem1. Its stated RH assumptions concern
their Dedekind zeta functions, rather than unspecified Artin factors.
The geometric criterion is that, at a good prime of characteristic
different from $\ell$, the full rank-two $\ell$-torsion is rational
over the residue field precisely when Frobenius acts trivially on
$E[\ell]$. Such full torsion prevents cyclicity. The characteristic
$p$ geometric $p$-torsion has rank at most one, so contributes no
additional rank-two obstruction. Corollary2.2 removes the resulting
finite set of residue-characteristic and model exceptions. This
identifies the inclusion--exclusion in (NT-E01a) with the stated
prime count. The limiting-density argument itself is the cited
classical theorem, not a new proof here.

Geometrically, each $K_m/K$ is the finite monodromy cover trivializing
the locally constant rank-two torsion sheaf of the elliptic scheme
over the appropriate arithmetic open. This uses finite quotients of
$\pi_1^{\text{ét}}$, and does not assert the Riemann hypothesis for
the infinite-image rank-two elliptic $L$-function. Work/run:
\path{GRH_NIGHT/number_theory}.

\subsection*{NT-E02. A finite Frobenius certificate and density-zero obstructions}
Suppose $E/K$ is non-CM. Let $T$ consist of the rational primes
dividing $30D_K$, those below a prime of bad reduction, and those for
which the image on $E[\ell]$ is not $\mathrm{GL}_2(\mathbb F_\ell)$.
This is a finite set, by the open image theorem. Put
$N=\prod_{\ell\in T}\ell$, $G_N=\operatorname{Gal}(K_N/K)$, and
\[
 C_N=\{g\in G_N:g|_{K_\ell}\ne1\text{ for every }\ell\mid N\}.
\]
Campagna--Stevenhagen Theorem1.1, author p.3, gives
\[
 \delta_{E/K}=\frac{|C_N|}{|G_N|}
   \prod_{\ell\nmid N}
       \left(1-\frac1{(\ell^2-1)(\ell^2-\ell)}\right).       \tag{NT-E02a}
\]
The product is strictly positive. Consequently $C_N=\varnothing$
is an unconditional finite certificate that all but finitely many
good reductions are noncyclic. If $C_N\ne\varnothing$, the hypothesis
of NT-E01 implies a positive density of cyclic reductions.

\paragraph{Complete finite-group deduction.}
For $m\mid N$, the event that a uniform $g\in G_N$ is trivial on
$K_m$ has probability $[K_m:K]^{-1}$. Intersections of these events
for primes dividing $m$ equal the $K_m$ event, by the compositum
identity proved above. Finite inclusion--exclusion therefore yields
\[
 \frac{|C_N|}{|G_N|}=\sum_{m\mid N}\frac{\mu(m)}{[K_m:K]}.
\]
The cited factorization gives (NT-E02a). For $\ell\ge2$, the
reciprocal of $(\ell^2-1)(\ell^2-\ell)$ is at most a constant
times $\ell^{-4}$ and is strictly less than one. The convergence
of its sum implies positivity of the product, by applying the
convergent logarithm series to all sufficiently large factors and
retaining the finitely many others.

If $C_N$ is empty, every element of $G_N$ lies in the kernel of at
least one restriction to $K_\ell$, $\ell\mid N$. Every good prime
outside the finite primes above $N$ and those ramified in $K_N$ has
such a Frobenius. Its reduction then contains
$(\mathbb Z/\ell\mathbb Z)^2$ for one such $\ell$, and is noncyclic.
This proves the finite-obstruction assertion without RH. Conversely,
$C_N\ne\varnothing$ makes the density expression positive, so NT-E01
gives infinitely many cyclic primes under its exact RH input.

For any $E/K$, including CM curves, the simpler obstruction
$E[\ell]\subset E(K)$ prevents cyclic reduction at every good prime
away from $\ell$, by the same reduction argument. Full rational
two-torsion is thus an immediate obstruction. Over a general $K$
the absence of this particular obstruction is not a positivity
criterion: the finite kernels in (NT-E02a) may jointly cover $G_N$
without any one being all of $G_N$. Actual elliptic examples of this
phenomenon are supplied by Campagna--Stevenhagen, Theorem5.1 and
\S5 (author pp.14--16); their construction uses three quadratic
division fields with Klein-four compositum. No independence of
division fields is silently assumed.

This is a finite monodromy-cover certificate, not a criterion for
a prescribed rational point to generate every cyclic reduction.
It also makes no claim about prime values of $\#E(\mathbb F_p)$.
Work/run: \path{GRH_NIGHT/number_theory}.

\subsection*{NT-E03. Lang--Trotter upper bounds without an extra holomorphy conjecture}
For a fixed non-CM elliptic curve $E/\mathbb Q$, set
$a_p=p+1-\#E(\mathbb F_p)$ at good primes, and
$P_{E,a}(x)=\#\{p\le x:p\nmid N_E,\ a_p=a\}$.
Zywina, \emph{Bounds for the Lang--Trotter conjectures},
Theorem1.2, author p.1, proves
\[
 P_{E,a}(x)\ll_E\frac{x^{4/5}}{(\log x)^{3/5}},\qquad
 P_{E,0}(x)\ll_E\frac{x^{3/4}}{(\log x)^{1/2}}.            \tag{NT-E03a}
\]
The first constant is uniform over integers $a$. These are upper
bounds, not the Lang--Trotter asymptotic. The source defines its GRH
on author p.3 as Dedekind RH for number fields. The proof permits the
following sufficient smaller input: adopted $A$, together with RH
for $\zeta_{\mathbb Q(E[\ell])}$ for every sufficiently large prime
$\ell$. Hence $M$ supplies the conclusion. Non-CM $\mathrm{GL}_2$
GRH is not being inferred.

\paragraph{Complete holomorphy/input audit.}
For large $\ell$ the mod-$\ell$ image is $G=\mathrm{GL}_2(\mathbb F_\ell)$.
Write $L=\mathbb Q(E[\ell])$, let $B$ be its upper triangular
subgroup, and let $U$ be its unipotent subgroup. Zywina's complete
\S4 proof, author pp.10--11, applies its Theorem2.3 to
$L^U/L^B$. Its group $B/U$ is abelian, so its Artin functions are
Hecke functions and the holomorphy assumption of Theorem2.3 is
known. Moreover $L/L^U$ has the cyclic group $U$, so the all-prime
abelian factorization NT-H11 gives
$\zeta_L=\zeta_{L^U}\prod_{\chi\ne1}L_{L^U}(s,\chi)$ with every
other factor entire. Thus RH for $\zeta_L$ implies RH for
$\zeta_{L^U}$, exactly what Theorem2.3 requires.

For $a=0$ the proof replaces $U$ by the subgroup $H$ of triangular
matrices with equal diagonal entries. This group is abelian:
writing a matrix as $aI+b\left(\begin{smallmatrix}0&1\\0&0\end{smallmatrix}\right)$
makes commutativity explicit. The same factorization transfers RH
from $\zeta_L$ to $\zeta_{L^H}$, while $B/H$ is again abelian.
The separate step choosing a prime $\ell$ split in
$\mathbb Q(\sqrt{a_p^2-4p})$ uses only RH for quadratic Dedekind
zeta functions and for $\zeta_{\mathbb Q}$: see Murty--Murty--Saradha
(1988), Lemma4.4 and its entire proof, printed pp.269--270.
These are supplied by $A$. This proves the stated sufficient input
reduction. The quantitative estimates in (NT-E03a) remain the
cited classical theorem.

The geometric data used here are finite mod-$\ell$ monodromy of
$H^1_{\text{ét}}(E_{\overline{\mathbb Q}},\mathbb F_\ell)$, with
the usual dual convention immaterial to the division field, and
abelian covers of intermediate arithmetic bases. An arbitrary
nonabelian use of Zywina's Theorem2.3 would still require its stated
Artin holomorphy hypothesis. Work/run:
\path{GRH_NIGHT/number_theory}.

\subsection*{NT-E04. Uniform counts with a prescribed Frobenius field}
Let $E/\mathbb Q$ be non-CM and let $k$ be an imaginary quadratic
field with class number $h_k$ and absolute discriminant $d_k$.
Put
\[
 P_{E,k}(x)=\#\{p\le x:p\nmid N_E,\
            \mathbb Q(\sqrt{a_p^2-4p})\cong k\}.
\]
Zywina Theorem1.3(i), author p.2, gives under Dedekind GRH
\[
 P_{E,k}(x)\ll_E
   h_k^{-3/5}\frac{x^{4/5}}{(\log x)^{3/5}}
                   +x^{1/2}(\log x)^3,                    \tag{NT-E04a}
\]
with a constant independent of $k$. Its proof is \S5.1,
author pp.13--14, after the algebraic reductions on pp.11--13.
Lemma5.5, author p.13, also gives $P_{E,k}(x)=0$ when $d_k>4x$.

For precision, this field family is larger than the division
fields in NT-E03. Let $H_k$ be the Hilbert class field of $k$,
and for split primes $\ell\ge5$ let $L_{k,\ell}/H_k$ be the
field of the joint representation formed from $E[\ell]$ and
the modulus-$\ell$ ray-class representation
$\psi_{k,\ell}:G_{H_k}\to
(\mathcal O_k/\ell\mathcal O_k)^*/\mathcal O_k^*$.
Zywina \S3.2--3.3, pp.7--10, defines it. The proof uses RH for
the Dedekind zeta functions of the auxiliary fields
$L_{k,\ell}^{\mathcal U}$ appearing in \S5, together with RH
for $\zeta_k$ and $\zeta_{\mathbb Q}$ in the split-prime choice.
Here $\mathcal B$ is the preimage of the triangular subgroup and
$\mathcal U$ is the image of the pairs $(aI+bN,a)$; the quotient
$\mathcal B/\mathcal U$ is abelian, so holomorphy is known.
This is an exact sufficient family specified by the cited proof;
$M$ supplies all of it. We do not replace it by Dirichlet GRH alone.

These counts distinguish endomorphism fields of reductions using
finite torsion and ray-class monodromy. They do not establish the
conjectured $x^{1/2}/\log x$ asymptotic. Work/run:
\path{GRH_NIGHT/number_theory}.

\subsection*{NT-E05. A refined lower bound for distinct Frobenius fields}
Let $\mathcal D_E(x)$ be the set of imaginary quadratic fields
arising as $\mathbb Q(\sqrt{a_p^2-4p})$ at good primes $p\le x$
for a fixed non-CM $E/\mathbb Q$. Under the exact input of NT-E04
for every imaginary quadratic $k$, together with adopted $A$,
\[
 |\mathcal D_E(x)|\gg_E
 \frac{x^{2/7}}
      {(\log x)^{4/7}(\log\log x)^{6/7}}
 \qquad(x\longrightarrow\infty).                         \tag{NT-E05a}
\]
In particular $M$ supplies these inputs. This is a newly proved
formulation in this run, not a claim of literature novelty.
Zywina's Corollary1.4, author p.2 with proof p.15, gives
$x^{2/7}/(\log x)^{10/7}$ using a weaker class-number estimate;
the proof below retains the sharper checked NT-H02 input.

\paragraph{Complete proof.}
Put $\Lambda(t)=\log\log(e^e t)$ for $t\ge1$. The explicit
NT-H02 bound, together with its finitely many excluded fundamental
discriminants, gives an absolute $c>0$ such that
\[
 h_k\ge c\,\frac{\sqrt{d_k}}{\Lambda(d_k)}
 \quad\hbox{for every imaginary quadratic }k.             \tag{NT-E05b}
\]
Indeed the displayed $V(d)$ in NT-H01 is $O(\log\log d)$ as
$d\to\infty$; lowering $c$ handles the finitely many smaller
discriminants, all of whose class numbers are positive.

Let $D=|\mathcal D_E(x)|$ and take $x$ sufficiently large that
$D\ge1$. Every such field has $d_k\le4x$ by NT-E04. Distinct
imaginary quadratic fields have distinct positive absolute
fundamental discriminants, so, ordering these discriminants,
the $j$th is at least $j$. It follows that
\begin{align*}
 \sum_{k\in\mathcal D_E(x)}h_k^{-3/5}
 &\ll \Lambda(4x)^{3/5}
                \sum_{k\in\mathcal D_E(x)}d_k^{-3/10}\\
 &\ll \Lambda(x)^{3/5}\sum_{j=1}^D j^{-3/10}
 \ll \Lambda(x)^{3/5}D^{7/10}.
\end{align*}
The last estimate uses
$\sum_{j=1}^D j^{-3/10}\le(10/7)D^{7/10}$, directly by integration of
$t^{-3/10}$ on $[1,D]$; its implicit constant is independent of
$E$, $x$ and the chosen fields.

Every good prime contributes to exactly one of these fields.
The prime number theorem, removal of the finitely many primes
dividing $N_E$, and (NT-E04a) therefore give
\[
 \frac{x}{\log x}\ll_E
    D^{7/10}\frac{x^{4/5}\Lambda(x)^{3/5}}{(\log x)^{3/5}}
                      +D x^{1/2}(\log x)^3.               \tag{NT-E05c}
\]
At least one right-hand term is at least a fixed positive
multiple of $x/\log x$. If it is the first, solving the
inequality gives
\[
 D\gg_E x^{2/7}(\log x)^{-4/7}\Lambda(x)^{-6/7}.
\]
If it is the second, then $D\gg_E x^{1/2}/(\log x)^4$.
The ratio of this latter lower bound to the former tends to
infinity, since it is
\[
 x^{3/14}(\log x)^{-24/7}\Lambda(x)^{6/7}\longrightarrow\infty.
\]
Thus the first lower bound follows in both cases for all
sufficiently large $x$. Finally
$\Lambda(x)\sim\log\log x$, which proves (NT-E05a).

This measures the diversity of Frobenius endomorphism fields
among closed fibres of an elliptic scheme. It does not prove
prime-order reductions or the Koblitz conjecture. Work/run:
\path{GRH_NIGHT/number_theory}.

\subsection*{NT-R05--NT-R07. Frobenius error laws with only Dedekind RH}

Fix a finite Galois extension $L/K$ of number fields, with group $G$.
For a real class function $t$ on $G$, put
\[
 \delta_t=\frac1{|G|}\sum_{g\in G}t(g),\qquad
 \pi_t(x)=\sum_{\substack{N\mathfrak p\le x\\\mathfrak p\ \text{unramified in }L}}
 t(\operatorname{Frob}_{\mathfrak p}),\qquad
 E_t(y)=ye^{-y/2}\bigl(\pi_t(e^y)-\delta_t\operatorname{Li}_2(e^y)\bigr).
\]
All assertions below require only RH for $\zeta_L$; hence $M$ suffices,
and $A$ suffices when $L/\mathbb Q$ is abelian. No assertion is uniform
as the extension varies.

\paragraph{The cited analytic input and its exact field requirement (NT-R05).}
Fiorilli--Jouve, \emph{Distribution of Frobenius elements in families of
Galois extensions}, author arXiv:2001.05428v2, Proposition3.18,
pp.32--34, proves the limiting-distribution theorem by a
Besicovitch $B^2$ approximation. Their Corollary3.17 and
Lemmata3.14--3.15, pp.30--32, supply the cyclic-induction version
which avoids Artin holomorphy. We use its specialization to
$\beta=1/2$, for real $t$. This gives the $B^2$ expansion below;
the full analytic almost-periodicity theorem is a cited classical
input. The precise specialization does not require $L/\mathbb Q$
to be Galois.

Here are all coefficients, with the conjugacy-class normalization
made explicit. For each conjugacy class $C$ choose $g_C\in C$,
put $H_C=\langle g_C\rangle$ and $F_C=L^{H_C}$, and define
\[
 a_t(\rho)=\sum_{C\subset G}\frac{|C|}{|G|}t(C)
       \sum_{\chi\in\widehat H_C}\overline{\chi(g_C)}
             \operatorname{ord}_{s=\rho}L(s,\chi,L/F_C).
 \tag{NT-R05a}
\]
The sum over $C$ means a sum over conjugacy classes. All functions
inside this formula are finite-order Hecke functions, with their
full local factors. The coefficient is independent of the choices
and is equivalently the residue at $\rho$ of the formal logarithmic
derivative attached to $t$.

Indeed, character orthogonality in the cyclic group and induction give
\[
 1_C=\frac{|C|}{|G|}
       \sum_{\chi\in\widehat H_C}\overline{\chi(g_C)}
                \operatorname{Ind}_{H_C}^G\chi.
\]
Artin induction for Euler factors, including inertia invariants,
then proves (NT-R05a). Moreover
$\zeta_L=\prod_{\chi\in\widehat H_C}L(s,\chi,L/F_C)$.
Every factor is holomorphic in the open critical strip. Thus RH for
$\zeta_L$ forces the nontrivial zeros of every factor in this formula
onto the critical line, without invoking holomorphy of any
nonabelian irreducible Artin function. This is the required analytic
input reduction in the cited cyclic argument.

Write $\Gamma$ for the set of distinct positive ordinates of zeros
of $\zeta_L$, and set
\[
 q_t=\frac1{|G|}\sum_{g\in G}t(g^2),\qquad
 b_t=-q_t-2a_t(1/2).
\]
Reality of $t$ gives $a_t(\overline\rho)=\overline{a_t(\rho)}$,
so $b_t$ is real. The cited expansion is
\[
 E_t(y)\ \sim_{B^2}\ b_t-
       2\Re\sum_{\gamma\in\Gamma}
        \frac{a_t(1/2+i\gamma)}{1/2+i\gamma}e^{i\gamma y}.
 \tag{NT-R05b}
\]
The notation means convergence of symmetric height truncations for
the seminorm
\[
 \|f\|_{B^2}^2=\limsup_{Y\to\infty}\frac1Y
                     \int_{\log2}^Y|f(y)|^2\,dy.
\]
The constant $-q_t$ is the prime-square term; the central zero
contributes $-2a_t(1/2)$. The local zero count NT-Z10 implies
\[
 \sum_{T<\gamma\le T+1}|a_t(1/2+i\gamma)|^2
       \ll_{L,K,t}\log^2(T+3),\qquad
 \sum_{\gamma\in\Gamma}
 \frac{|a_t(1/2+i\gamma)|^2}{1/4+\gamma^2}<\infty.
 \tag{NT-R05c}
\]
To check this directly, each nonnegative factor order in
(NT-R05a) is at most the corresponding order of $\zeta_L$.
The sum of these latter orders in a unit interval is $O_L(\log(T+3))$.
Taking absolute values in the finite class-character sum and then
squaring proves the first estimate; summing over integer intervals
proves the second. Multiplicities are retained throughout.

\paragraph{Complete joint-law and covariance deduction (NT-R06).}
For any finite list of real class functions $t_1,\ldots,t_d$, the
vector $(E_{t_1}(y),\ldots,E_{t_d}(y))$ has a limiting probability
law $\nu$ on $\mathbb R^d$. Precisely, the averages of every bounded
continuous function of this vector converge to its $\nu$ integral.
The mean vector is $(b_{t_j})_j$, and its covariance matrix is
\[
 \Sigma_{jk}=2\Re\sum_{\gamma\in\Gamma}
       \frac{a_{t_j}(1/2+i\gamma)
             \overline{a_{t_k}(1/2+i\gamma)}}{1/4+\gamma^2}.
 \tag{NT-R06a}
\]
Every series here converges absolutely. This is a full new
formulation of the joint deduction from the cited scalar expansion,
not a claim of literature novelty.

To prove it, take the closure $\mathcal T$ of the one-parameter
subgroup
$y\mapsto(e^{i\gamma y})_{\gamma\in\Gamma}$ in the countable
product of circles and give it Haar probability measure. Every
continuous character of a finite-dimensional projection has the
form $e^{i\lambda y}$ along this subgroup. Its Cesaro mean is zero
unless $\lambda=0$, in which case it is one. Stone--Weierstrass
therefore proves that the time averages of every continuous
function of finitely many coordinates equal the corresponding
Haar integrals. Distinct positive frequencies are orthogonal in
$L^2(\mathcal T)$; so are a positive frequency and the conjugate of
another positive frequency.

By (NT-R05c) the series
\[
 Z_j=b_{t_j}-2\Re\sum_{\gamma\in\Gamma}
       \frac{a_{t_j}(1/2+i\gamma)}{1/2+i\gamma}z_\gamma
\]
converges in $L^2(\mathcal T)$. Let $\nu$ be the law of this vector.
For any bounded Lipschitz test function, Cauchy--Schwarz bounds the
difference between its time average on $E$ and on a common finite
truncation by the Euclidean sum of the $B^2$ approximation errors.
The same bound holds for its Haar integrals on $Z$ and that
truncation. The finite-truncation averages converge by the preceding
character argument; letting the truncation grow proves convergence
for bounded Lipschitz tests. These tests determine weak convergence
of probability measures on $\mathbb R^d$, which gives the stated
bounded-continuous formulation. Taking Haar means and second moments
in the $L^2$ series proves (NT-R06a), including the factor two.
Cauchy--Schwarz and (NT-R05c) prove absolute convergence of its
mixed series.

No independence of the coordinates $z_\gamma$ was used: all rational
relations between the ordinates remain encoded in $\mathcal T$.
In particular this argument proves neither a Gaussian law nor
independence of zeros. A zero covariance in some direction means
exactly that its noncentral coefficients $a_t(1/2+i\gamma)$ all
vanish; its limiting law is then the point mass at $b_t$.

\paragraph{Finite etale fibres have a negative mean and a quantitative
one-sided lower density (NT-R07).}
Let $X\to U\subset\operatorname{Spec}\mathcal O_K$ be a nonempty
finite etale scheme, with generic algebra $\prod_{i=1}^r K_i$ and
Galois splitting field $L/K$. Put
\[
 \Pi_X(x)=\sum_{\substack{\mathfrak p\in U\\N\mathfrak p\le x}}
                         \#X(\mathcal O_K/\mathfrak p),\qquad
 E_X(y)=ye^{-y/2}\bigl(\Pi_X(e^y)-r\operatorname{Li}_2(e^y)\bigr).
\]
Apply NT-R05 to the permutation character $t$ of the generic
geometric fibre $S$. Omitting the finitely many additional primes
outside $U$ changes $\pi_t(x)$ by a bounded amount; its normalized
contribution tends to zero and has $B^2$ seminorm zero. Burnside's
lemma gives $\delta_t=r$. The full Artin permutation identity gives
\[
 a_t(1/2)=\sum_{i=1}^r\operatorname{ord}_{s=1/2}\zeta_{K_i}(s)
       \ge0.
\]
To evaluate the sign of $q_t$, decompose $S=\coprod_iG/H_i$.
Counting the pairs $(g,aH_i)$ fixed by $g^2$ gives
\[
 q_t=\sum_{i=1}^r
      \frac{\#\{g\in G:g^2\in H_i\}}{|H_i|}\ge r,
\]
because every element of $H_i$ belongs to the set being counted.
Thus the mean is $-B_X$, where
\[
 B_X=q_t+2\sum_i\operatorname{ord}_{1/2}\zeta_{K_i}\ge r>0,
 \qquad
 V_X=2\sum_{\gamma\in\Gamma}
          \frac{|a_t(1/2+i\gamma)|^2}{1/4+\gamma^2}<\infty.
\]
Consequently
\[
 \liminf_{X_0\to\infty}\frac1{\log X_0}
 \int_2^{X_0}
 1_{\{\Pi_X(x)<r\operatorname{Li}_2(x)\}}\,\frac{dx}{x}
       \ge\frac{B_X^2}{B_X^2+V_X}>0.
 \tag{NT-R07a}
\]
Here is the probability step in full. If $Z$ has the limiting law,
then $W=Z+B_X$ has mean zero and variance $V_X$. For $u\ge0$,
Markov's inequality gives
$\Pr(W\ge B_X)\le(V_X+u^2)/(B_X+u)^2$.
For $V_X>0$, choose $u=V_X/B_X$; for $V_X=0$ the variable
$Z=-B_X$ almost surely. In both cases
$\Pr(Z<0)\ge B_X^2/(B_X^2+V_X)$.
Weak convergence and the open set $(-\infty,0)$ imply the same
lower bound for the liminf of the time proportions. The substitution
$x=e^y$ gives (NT-R07a).

This is a positive lower logarithmic density. The existence of the
sign density, its being greater than one half, and densities of all
strict orderings are not asserted. They require additional control
of the limiting measure, notably its mass on the relevant boundary.
Geometrically these statements concern the trace of Frobenius on
finite-monodromy local systems and, for NT-R07, the permutation
realization of $H^0$ of a finite etale arithmetic cover. They use
ordinary Hecke/Dedekind realizations and do not invoke GRH for an
infinite-image geometric representation.

\subsection*{NT-R08. A discriminant-uniform lower logarithmic density}

Retain the finite etale setting of NT-R07 and assume only RH for
$\zeta_L$. There is an absolute constant $C>0$ such that
\[
 V_X\le C r^2\log^2(\Delta_L+2),\qquad
 \underline\delta\{\Pi_X(x)<r\operatorname{Li}_2(x)\}
 \ge \frac1{1+C\log^2(\Delta_L+2)}.
 \tag{NT-R08a}
\]
Here $\underline\delta$ denotes the lower logarithmic density
defined by the left side of (NT-R07a). The same constant works
for every number field $K$, every finite Galois splitting field
$L/K$, and every nonempty finite etale cover with this splitting
field. No lower bound on zero spacing or upper bound on individual
multiplicities independent of $\Delta_L$ is assumed. The
almost-periodic limiting process is still taken separately for
each fixed extension; no uniform rate of convergence to the law
is asserted.

\paragraph{Complete proof.}
Write $n=[L:\mathbb Q]$ and $D=\Delta_L$. Let
$m_L(\gamma)=\operatorname{ord}_{s=1/2+i\gamma}\zeta_L(s)$.
For every $\gamma>0$, the cyclic-character formula (NT-R05a)
and nonnegativity of the permutation character $t$ give
\[
 |a_t(1/2+i\gamma)|
 \le\sum_C\frac{|C|}{|G|}t(C)
       \sum_{\chi\in\widehat H_C}
          \operatorname{ord}_{1/2+i\gamma}L(s,\chi,L/F_C)
 =r\,m_L(\gamma).
 \tag{NT-R08b}
\]
This is an exact inequality: the inner sum equals the order
of $\zeta_L$, and the outer sum equals $r$ by Burnside.

The unconditional zero-count formula in Grenie--Molteni,
AppendixA TheoremA.1, author p.26, specialized to the trivial
character of $L$, implies with an absolute implied constant
\[
 U_j:=\sum_{\substack{\gamma\ {\rm distinct}\\j<\gamma\le j+1}}
                  m_L(\gamma)
       \ll \log D+n\log(j+3)\qquad(j=0,1,2,\ldots).
 \tag{NT-R08c}
\]
For clarity, the source counts both signs and includes
multiplicities. Its main term is
$T\{\log D+n\log(T/(2\pi e))\}/\pi$, with error
$O(\log D+n\log T+n)$ for $T\ge1$, and an additional
$O(n)$ constant term. For $j\ge2$ subtract its values at
$j-1$ and $j+1$; the difference of the main terms has the
displayed size by differentiation, and the sum of the two
error bounds has that size as well. This bounds the positive
ordinates in $(j,j+1]$. For $j=0,1$ use the source bound for
the total count at $T=3$. The absolute constants in its table
therefore give (NT-R08c) uniformly in the field.

Within a unit interval the orders are nonnegative, so the sum
of their squares is at most $U_j^2$. Equations (NT-R06a) and
(NT-R08b) yield
\[
 \begin{split}
 V_X
 &\le 2r^2\sum_{\gamma>0}
                  \frac{m_L(\gamma)^2}{1/4+\gamma^2}\\
 &\ll r^2\sum_{j=0}^{\infty}
          \frac{(\log D+n\log(j+3))^2}{1/4+j^2}\\
 &\ll r^2(\log D+n)^2.
 \end{split}
 \tag{NT-R08d}
\]
The final constant is absolute, since the series with numerator
$1$, $\log(j+3)$ or $\log^2(j+3)$ all converge against this
denominator.

We include the elementary removal of the degree parameter.
Minkowski's discriminant bound gives
\[
 \sqrt D\ge\frac{n^n}{n!}\left(\frac{\pi}{4}\right)^{r_2}.
\]
An exact reference is Milne, \emph{Algebraic Number Theory},
v3.08, Theorem4.3 printed p.70 and the proof of Theorem4.9
printed p.72. For $n\ge2$, the arithmetic-geometric mean
inequality gives $n!\le((n+1)/2)^n$, and $r_2\le n/2$.
Consequently
\[
 \sqrt D\ge
 \left(\frac{2n}{n+1}\sqrt{\frac{\pi}{4}}\right)^n
 \ge\left(\frac{2\sqrt\pi}{3}\right)^n.
\]
Since $2\sqrt\pi/3>1$, this proves
$n\ll\log D$ with an absolute constant for $n\ge2$.
The case $n=1,D=1$ is handled by $\log(D+2)$.
Substitution in (NT-R08d) proves the first assertion of
(NT-R08a). Finally $B_X\ge r$ in NT-R07 gives
\[
 \frac{B_X^2}{B_X^2+V_X}
 \ge \frac1{1+C\log^2(D+2)},
\]
which proves the lower-density assertion.

The variance estimate is consistent with the squared
$\|t\|_1^2$ bound in Fiorilli--Jouve Proposition4.6(87),
author p.40 and its proof p.41. The proof above derives the
coarse bound directly from the uniform zero count and gives
the finite etale consequence without invoking a bounded
multiplicity hypothesis. No numerical value for the absolute
constant is asserted.

In the Grothendieck language of NT-F02 and NT-R07, this is a
uniform lower-frequency statement for the Frobenius trace on
the permutation realization of $H^0$ of a finite etale cover.
The generic rank of that realization may grow; the final lower
bound depends on the splitting-field discriminant, not that
rank. It does not give a uniform threshold in $x$ or a uniform
error term for the approach to the limiting law.

\subsection*{NT-E06. Uniform CM angle counts from one Hecke-power family}
Let $E/k$ be an elliptic curve whose complex multiplication is defined
over the number field $k$. Put
\[
 a_{\mathfrak p}=N\mathfrak p+1-\#E_{\mathfrak p}(k_{\mathfrak p}),
 \qquad t_{\mathfrak p}=\frac{a_{\mathfrak p}}{\sqrt{N\mathfrak p}}.
\]
Only primes of good reduction are counted. The classical CM character
theorem supplies a unitary Hecke character $\eta$ of $k$ for which
\[
 t_{\mathfrak p}=\eta(\mathfrak p)+\overline{\eta(\mathfrak p)}
                =2\cos(2\pi\theta_{\mathfrak p}),
 \qquad \theta_{\mathfrak p}\in\mathbb R/\mathbb Z.       \tag{NT-E06a}
\]
Assume RH for $L_k(s,\eta^m)$ for every integer $m$, including
$m=0$, where the function is $\zeta_k$. This is a precisely specified
subfamily of $M$; Dirichlet GRH alone does not supply its nonzero
infinity types. Then, uniformly over all arcs $J\subset\mathbb R/\mathbb Z$
and all intervals $I\subset[-2,2]$, for $x\ge2$,
\begin{align}
 \#\{\mathfrak p:N\mathfrak p\le x,\ \theta_{\mathfrak p}\in J\}
   &= |J|\operatorname{Li}_2(x)
       +O_{E,k}\bigl(\sqrt x\log^2(2x)\bigr), \tag{NT-E06b}\\
 \#\{\mathfrak p:N\mathfrak p\le x,\ t_{\mathfrak p}\in I\}
   &= \mu(I)\operatorname{Li}_2(x)
       +O_{E,k}\bigl(\sqrt x\log^2(2x)\bigr), \tag{NT-E06c}\\
 d\mu(t)&=\frac{dt}{\pi\sqrt{4-t^2}}.\notag
\end{align}
The constants are independent of the interval and its endpoint choices.
In particular the intervals may vary with $x$. No numerical implicit
constant or uniformity in the curve or base field is asserted.

\paragraph{Classical input and complete deduction.}
We use the Hecke prime-sum estimate recorded by Bucur--Fit\'e--Kedlaya,
\emph{Effective Sato--Tate conjecture for abelian varieties and
applications}, Proposition2.5, equation(2.3), published p.1720
(author v2 p.6). Its CM interpretation is explicitly identified in
footnote2 to Proposition4.7, published p.1740 (author p.20).
For the fixed curve and field it gives, uniformly in $m\ne0$ and $x\ge2$,
\begin{align}
 S_m(x):=\sum_{N\mathfrak p\le x}\eta(\mathfrak p)^m
       &=O_{E,k}\bigl(\sqrt x\log(2x+2|m|)\bigr),\notag\\
 S_0(x)&=\operatorname{Li}_2(x)
       +O_{E,k}\bigl(\sqrt x\log(2x)\bigr).               \tag{NT-E06d}
\end{align}
Their sums omit the primes over the absolute conductor; restoring any
good primes omitted by that convention changes a sum by $O_{E,k}(1)$.
The representation dimension in this specialization is one, while the
motivic weight is $|m|$: the distinction is the point of their
Remark2.7. The finite conductor of $\eta^m$ divides that of $\eta$,
and its archimedean parameters grow at most linearly in $|m|$.
Every nonzero power has a nonzero angular infinity type, so is not
a pure norm twist. Consequently the required entire continuation,
functional equation, and boundary nonvanishing are classical Hecke
facts; the only conjectural analytic input in (NT-E06d) is the
stated RH family. The standard Hecke completion over $k$ has
degree $[k:\mathbb Q]$, which is fixed here.

For clarity, the elementary harmonic-analysis input is
Erd\H{o}s--Tur\'an, \emph{On a problem in the theory of uniform
distribution.I} (1948), TheoremIII, printed p.374 (parallel
proceedings pagination p.1150). For $P$ points on the circle and an
integer $H\ge1$ its discrepancy bound is
\[
 \left|\#\{\theta_j\in J\}-|J|P\right|
 \le C\left(\frac P{H+1}
       +\sum_{m=1}^{H}\frac1m
                    \left|\sum_{j=1}^{P}e^{2\pi i m\theta_j}\right|\right),
                                                               \tag{NT-E06e}
\]
where $C$ is absolute. The original statement includes closed
endpoints. The same bound, after changing $C$, holds with any endpoint
convention: shrink or enlarge an arc, use the same bound, and pass to
the limit for the finite set of points. Singletons are handled by this
argument as well.

Apply (NT-E06e) with $H=\lceil x\rceil$ and the angles in (NT-E06a).
The first term is $O_k(1)$, since the number of prime ideals of norm at
most $x$ is at most $[k:\mathbb Q]x$. By (NT-E06d) the sum is at most
\[
 C_{E,k}\sqrt x\log(4x+2)\sum_{m\le\lceil x\rceil}\frac1m
     =O_{E,k}\bigl(\sqrt x\log^2(2x)\bigr).
\]
Replacing $P$ by $\operatorname{Li}_2(x)$ using $m=0$ proves
(NT-E06b). The inverse image of a trace interval under
$\theta\mapsto2\cos(2\pi\theta)$ is the union of at most two arcs.
The pushforward of normalized Lebesgue measure is exactly $\mu$;
adding the two estimates proves (NT-E06c), with singleton endpoints
already accounted for. This proves all the stated uniformities.

\paragraph{Geometric meaning and locator.}
The local system $R^1f_*\mathbb Q_\ell$ of a smooth elliptic scheme over
an arithmetic open has rank two and weight one. CM splits its scalar
extension into two conjugate rank-one characters; their unitary
normalizations give the circle above. Its tensor powers are the
specific Hecke family used in the proof. Thus this is a Frobenius
distribution statement for a CM torus in the motivic monodromy group,
not a deduction of GRH for an arbitrary non-CM elliptic $L$-function.
Work/run: \path{GRH_NIGHT/number_theory},
\path{cm_prime_angles.tex}, NT-E06; sources S-NT-BFK and S-NT-ET.

\subsection*{NT-E07. A least-prime bound in any positive-measure CM interval}
Under the exact hypotheses of NT-E06, for any trace interval $I$ with
$\mu=\mu(I)>0$, there is a good prime ideal with $t_{\mathfrak p}\in I$
and
\[
 N\mathfrak p\le C_{E,k}\,\mu^{-2}\log^6(e/\mu).           \tag{NT-E07}
\]
The analogous conclusion holds for an arc of length $\mu$. This is a
newly proved formulation of the classical discrepancy consequence;
no claim of literature novelty or numerical value of $C_{E,k}$ is made.

\paragraph{Complete proof.}
Let $A_{E,k}$ be a common error constant in (NT-E06b)--(NT-E06c).
For $x\ge4$,
$\operatorname{Li}_2(x)\ge x/(2\log x)$, by integrating over $[x/2,x]$.
Write $L=\log(e/\mu)\ge1$ and set $x=C\mu^{-2}L^6$, with $C\ge4$.
Since $\log L\le L$ and $\log(1/\mu)\le L$, we have
\[
 \log(2x)\le(\log(2C)+8)L,\qquad
 \mu\sqrt x=\sqrt C\,L^3.
\]
The prime count is positive provided
$\mu\sqrt x>2A_{E,k}\log x\,\log^2(2x)$.
It is therefore enough to choose the fixed $C$ so that
\[
 \sqrt C>2A_{E,k}\bigl(\log(2C)+8\bigr)^3.
\]
Such a $C$ exists because $\sqrt C/\log^3(2C)\to\infty$.
This choice is independent of $I$, proves the bound, and also covers
intervals with endpoints included or excluded. Geometrically it
detects a prescribed positive-measure region in the compact CM
monodromy torus by one closed arithmetic fibre. Work/run:
\path{GRH_NIGHT/number_theory}, \path{cm_prime_angles.tex}, NT-E07.

\subsection*{NT-E08. An extremal CM trace asymptotic with a power-saving error}
Let now $E/k$ be any elliptic curve with potential CM by the imaginary
quadratic field $K$. Put $F=kK$ and $d=[F:k]\in\{1,2\}$. Let $\eta$
be a CM unitary Hecke character of $E_F$. Assume RH for the functions
$L_F(s,\eta^m)$ for every integer $m$, including $m=0$.
For either choice of sign, define
\[
 M^\pm_{E/k}(x)
  =\#\{\mathfrak p\text{ good}:N\mathfrak p\le x,\
               a_{\mathfrak p}=\pm\lfloor2\sqrt{N\mathfrak p}\rfloor\}.
\]
Then
\begin{align}
 M^\pm_{E/k}(x)
   &=\frac1{\pi d}\int_2^x\frac{t^{-1/4}}{\log t}\,dt
       +O_{E,k}\bigl(x^{5/8}\log^2(2x)\bigr),             \tag{NT-E08a}\\
 M^\pm_{E/k}(x)
   &\sim \frac4{3\pi d}\frac{x^{3/4}}{\log x}.
                                                               \tag{NT-E08b}
\end{align}
In particular $M$ supplies the result conditionally. The factor is
$4/(3\pi)$ when all CM is defined over $k$, and $2/(3\pi)$ in the
quadratic descent case. The integral in (NT-E08a) is retained because
replacing it by its first asymptotic term would introduce a larger
logarithmic error. This is a fully proved formulation in this run,
not an assertion that the asymptotic or its constant is new in the
literature.

\paragraph{Complete proof when $d=1$.}
It suffices first to consider the positive sign. If $a$ is an integer
with $a\le y$, then $a=\lfloor y\rfloor$ if and only if $a>y-1$.
Thus the Hasse bound gives the exact identity
\[
 a_{\mathfrak p}=\lfloor2\sqrt{N\mathfrak p}\rfloor
 \quad\Longleftrightarrow\quad
 t_{\mathfrak p}\in I_{N\mathfrak p},\qquad
 I_u=(2-u^{-1/2},2].                                    \tag{NT-E08c}
\]
The strict lower endpoint is essential and makes the identity valid
also when $N\mathfrak p$ is a square. No fractional-part factor occurs.
Direct integration of the arcsine density gives, uniformly for $u\ge2$,
\[
 \mu(I_u)=\frac1\pi\arccos\!\left(1-\frac1{2\sqrt u}\right)
         =\frac1\pi u^{-1/4}+O(u^{-3/4}).                 \tag{NT-E08d}
\]
For example the last estimate follows by setting $v=u^{-1/2}$ and
integrating $(4z-z^2)^{-1/2}=(2\sqrt z)^{-1}(1+O(z))$
from $0$ to $v$; the constants are absolute on $0\le v\le2^{-1/2}$.

For $x\ge4$ set $h=x^{-1/8}$ and partition $[2,x]$ into intervals
$(u_{j+1},u_j]$, with $u_0=x$, successive ratios $u_j/u_{j+1}=1+h$,
and the final lower endpoint truncated to $2$. Every ratio is at most
$1+h$. Since $I_u$ decreases with $u$, the contribution on this
interval is bracketed by the counts using the two fixed intervals
$I_{u_j}$ and $I_{u_{j+1}}$. Apply (NT-E06c) at both endpoints.
After summing, the total error in each bracket is at most
\[
 C_{E,k}\sum_j\sqrt{u_j}\log^2(2u_j)
       \ll_{E,k}h^{-1}\sqrt x\log^2(2x),                 \tag{NT-E08e}
\]
because $\sum_{j\ge0}(1+h)^{-j/2}\ll h^{-1}$ for $0<h\le1$.
Primes of norm exactly $2$ contribute $O_k(1)$.

The main brackets are the upper and lower Riemann sums of the
decreasing function $\mu(I_u)$ against the positive measure
$d\operatorname{Li}_2(u)=du/\log u$. Replacing $\mu(I_u)$ using
(NT-E08d) introduces an error bounded by a constant times
\[
 \int_2^x\frac{u^{-3/4}}{\log u}\,du+1
      \ll x^{1/4}+1.                                   \tag{NT-E08f}
\]
Indeed values of $u^{-3/4}$ at the two endpoints differ by a factor
at most $(1+h)^{3/4}$. For $f(u)=u^{-1/4}$ the corresponding
upper and lower sums differ by at most
$((1+h)^{1/4}-1)$ times the lower sum. They therefore bracket
$\int_2^x u^{-1/4}\,du/\log u$ with error
\[
 \ll h\int_2^x\frac{u^{-1/4}}{\log u}\,du
 \ll h x^{3/4}.                                       \tag{NT-E08g}
\]
Combining (NT-E08e)--(NT-E08g), then substituting $h=x^{-1/8}$,
proves (NT-E08a) when $d=1$. Applying exactly the same argument to
$-t_{\mathfrak p}$ proves the negative sign, since $\mu$ is symmetric.
Finally l'H\^opital's rule, or a single integration by parts, gives
\[
 \int_2^x\frac{u^{-1/4}}{\log u}\,du
      \sim \frac43\frac{x^{3/4}}{\log x},
\]
and $x^{5/8}\log^2(2x)=o(x^{3/4}/\log x)$.

\paragraph{Complete quadratic descent.}
The field defining all endomorphisms is $F=kK$. One can see this
directly on invariant differentials: in characteristic zero the
differential embeds $\operatorname{End}(E_{\bar k})\otimes\mathbb Q=K$
faithfully in $\bar k$, compatibly with $G_k$. Its kernel of Galois
action is therefore $G_{kK}$. If $d=2$, at an unramified good prime
inert in $F/k$, Frobenius interchanges the two CM eigenspaces of the
Tate module. Its matrix has zero diagonal, so $a_{\mathfrak p}=0$.
Such a prime is never extremal because $\lfloor2\sqrt{N\mathfrak p}\rfloor
\ge2$.

Every split good prime of $k$ has two primes of $F$ of the same norm
and the same Frobenius trace. A prime of $F$ above an inert prime has
norm $(N\mathfrak p)^2$; consequently the number of such primes of
norm at most $x$ is $O_k(\sqrt x)$, even by the elementary bound for
prime ideals above rational primes. Ramified primes and primes where
good reduction changes under base extension form a fixed finite set.
It follows that
\[
 M^\pm_{E/F}(x)=2M^\pm_{E/k}(x)+O_{E,k}(\sqrt x).
\]
Apply the already proved connected-CM case over $F$ and divide by
two. This proves both formulas for $d=2$, with the same stated
error scale.

\paragraph{Literature comparison, geometry, and proof location.}
Bucur--Fit\'e--Kedlaya Proposition4.9 and Corollary4.10, published
pp.1741--1742, give the order $x^{3/4}/\log x$ and the quadratic
base-change mechanism under this same Hecke-power input.
James--Pollack, \emph{Extremal primes for elliptic curves with
complex multiplication} (2017), Theorem1 (author p.2), states the
$2/(3\pi)$ asymptotic over $\mathbb Q$ unconditionally; its separate
erratum replaces the narrow-sector analytic input. That published
result and its erratum are comparison references, not inputs to
the proof above. This run does not certify the erratum's sector
estimate independently.

The arithmetic content is the number of closed fibres for which the
trace on $H^1$ reaches the integral Hasse endpoint. The CM torus and
its rank-one Hecke characters supply all the analytic hypotheses;
there is no implication here for extremal traces of a non-CM curve.
Work/run: \path{GRH_NIGHT/number_theory},
\path{cm_prime_angles.tex}, NT-E08; complete deduction from NT-E06,
with source inputs S-NT-BFK and S-NT-ET and comparison S-NT-JP.

\subsection*{NT-E09. Shrinking absolute deficits from the CM Hasse endpoint}
Retain $E/k$, its imaginary quadratic CM field $K$, $F=kK$, and
$d=[F:k]\in\{1,2\}$ from NT-E08. Fix real numbers
\[
 c>0,\qquad -\frac12<\lambda<\frac12,\qquad
 \alpha=\frac14-\frac{\lambda}{2}\in(0,\tfrac12).
\]
Set $u_0=\max\{2,c^{1/(1/2-\lambda)}\}$, so that
$c u_0^{\lambda-1/2}\le1$, and put, for $u\ge u_0$,
\[
 g_{c,\lambda}(u)=\frac1\pi
       \arccos\!\left(1-\frac c2u^{\lambda-1/2}\right).
                                                               \tag{NT-E09a}
\]
For either sign, let
\[
 D^\pm_{E/k;c,\lambda}(x)
 =\#\{\mathfrak p\text{ good}:N\mathfrak p\le x,\
              0\le2\sqrt{N\mathfrak p}\mp a_{\mathfrak p}
                           <c(N\mathfrak p)^\lambda\}.
                                                               \tag{NT-E09b}
\]
Assume RH for all integral powers of the unitary CM Hecke character
of $E_F$, including its zeroth power $\zeta_F$. Then, as $x\to\infty$,
\begin{align}
 D^\pm_{E/k;c,\lambda}(x)
   &=\frac1d\int_{u_0}^{x}\frac{g_{c,\lambda}(u)}{\log u}\,du
       +O_{E,k,c,\lambda}
           \bigl(x^{5/8+\lambda/4}\log^2(2x)\bigr),       \tag{NT-E09c}\\
 D^\pm_{E/k;c,\lambda}(x)
   &\sim\frac{\sqrt c}{\pi d(3/4+\lambda/2)}
                     \frac{x^{3/4+\lambda/2}}{\log x}.   \tag{NT-E09d}
\end{align}
The same statements hold if only prime ideals of absolute residue
degree one are counted. All parameters $E,k,c,\lambda$ are fixed;
the constants and threshold are not asserted uniform as
$\lambda$ approaches either endpoint.
For $-1/2<\lambda<0$ this gives infinitely many primes with an
absolute Hasse deficit tending to zero. This is a new complete
formulation in the run, without a claim of literature priority.

\paragraph{Complete proof of the moving-target estimate.}
First suppose $d=1$. For the positive sign the condition in
(NT-E09b), at norm $q\ge u_0$, is precisely
\[
 t_{\mathfrak p}\in
  J_q=(2-cq^{\lambda-1/2},2].
\]
The arcsine measure of $J_u$ is exactly $g_{c,\lambda}(u)$.
Let $v=cu^{-2\alpha}$, so $0<v\le1$. Elementary differentiation
and integration give
\[
 g_{c,\lambda}(u)\asymp_{c,\lambda}u^{-\alpha},
 \qquad
 -u\,g'_{c,\lambda}(u)
       =\frac{\alpha\sqrt v}{\pi\sqrt{1-v/4}}.            \tag{NT-E09e}
\]
In fact $g(u)\ge\sqrt v/\pi$, because
$g(u)=\pi^{-1}\int_0^v(4z-z^2)^{-1/2}\,dz$.
Since $v\le1$, (NT-E09e) consequently yields the explicit bound
\[
 0\le-\frac{u g'(u)}{g(u)}\le\frac{2\alpha}{\sqrt3}.
                                                               \tag{NT-E09f}
\]
Thus $g$ is decreasing and its values on a multiplicative interval
of ratio at most $1+h$ differ by a factor at most
$(1+h)^{2\alpha/\sqrt3}=1+O_\lambda(h)$ for $0<h\le1$.

Set $h=x^{-1/4+\alpha/2}$, which tends to zero because $\alpha<1/2$.
Partition $[u_0,x]$ into successive multiplicative intervals of
ratio $1+h$, truncating the last one at $u_0$. On each interval,
the moving interval $J_q$ lies between its two endpoint intervals.
Applying NT-E06 at both norm endpoints and adding the errors gives
\[
 O_{E,k}\!\left(\sqrt x\,h^{-1}\log^2(2x)\right),
                                                               \tag{NT-E09g}
\]
by the geometric series already evaluated in (NT-E08e).
The two main terms are Riemann sums for
$\int_{u_0}^x g(u)\,du/\log u$. By (NT-E09f) their difference is
bounded by a constant times
\[
 h\int_{u_0}^x\frac{g(u)}{\log u}\,du
    \ll_{c,\lambda}h x^{1-\alpha}.                      \tag{NT-E09h}
\]
The finite range below $u_0$ adds $O_{E,k,c,\lambda}(1)$.
Both powers in (NT-E09g)--(NT-E09h) equal
$x^{3/4-\alpha/2}=x^{5/8+\lambda/4}$. This proves (NT-E09c)
for $d=1$. Reflection of the trace interval proves the negative
sign with the identical mass and constants.

To obtain (NT-E09d), (NT-E09a) gives
$g(u)\sim(\sqrt c/\pi)u^{-\alpha}$. Comparison of the positive
integrals, followed by differentiation of
$x^{1-\alpha}/((1-\alpha)\log x)$, yields
\[
 \int_{u_0}^x\frac{g(u)}{\log u}\,du
     \sim\frac{\sqrt c}{\pi(1-\alpha)}
                           \frac{x^{1-\alpha}}{\log x}.
\]
The ratio of the error in (NT-E09c) to this main term is
$O(x^{-1/4+\alpha/2}\log^3(2x))=o(1)$, exactly because
$\alpha<1/2$.

If $d=2$, the nonsplit good primes have $a_{\mathfrak p}=0$ by
the CM eigenspace argument in NT-E08. They cannot satisfy
(NT-E09b) once $q$ exceeds a fixed threshold, since
$2\sqrt q>c q^\lambda$ for large $q$. Split primes again have two
lifts of equal norm and trace, while nonsplit primes upstairs have
square norms. The same $O_{E,k}(\sqrt x)$ correction therefore gives
the factor $1/d$. Finally, the number of prime ideals of any fixed
number field with absolute residue degree at least two and norm at
most $x$ is $O_k(\sqrt x)$: their underlying rational prime is at
most $\sqrt x$, and there are at most $[k:\mathbb Q]$ primes above it.
Because $5/8+\lambda/4>1/2$, deleting them does not change either
formula.

\paragraph{Exact range, analytic input, and geometric meaning.}
The small-angle expansion is used only to identify the asymptotic
constant. The exact arcsine integral is retained in the
power-saving statement. In particular an error of size
$\int^{x}u^{-3\alpha}\,du/\log u$ from replacing it by the first
term need not fit within (NT-E09c); for example it is larger when
$\lambda>3/10$. No such replacement is made.

At $\lambda=-1/2$ the predicted exponent becomes $1/2$ and this
argument no longer gives a smaller error. No endpoint assertion,
quadratic-polynomial prime theorem, or corresponding non-CM
statement is claimed. At $\lambda=0,c=1$ the strict inequality
recovers the exact integer extremal event in NT-E08. The result
describes shrinking neighbourhoods of a scalar Frobenius in the
compact CM monodromy torus of $H^1$, while retaining the actual
integer trace and absolute norm of the arithmetic fibre.

The only analytic hypothesis is the same power family as NT-E06;
$M$ supplies it conditionally. Sources S-NT-BFK and S-NT-ET supply
the classical prime-sum/discrepancy premises already audited there.
Work/run: \path{GRH_NIGHT/number_theory},
\path{cm_shrinking_deficits.tex}, NT-E09.

\subsection*{NT-A04. General Hecke angle discrepancy and its norm-twist exception}
Let $K$ be a number field and let $\eta$ be a unitary Hecke character
which is \emph{not} a finite-order character times a unitary power
of the norm. For every prime ideal unramified for $\eta$, define
$\theta_{\mathfrak p}\in\mathbb R/\mathbb Z$ by
$\eta(\mathfrak p)=e^{2\pi i\theta_{\mathfrak p}}$.
Assume RH for $L_K(s,\eta^m)$ for all positive integers $m$, and RH
for $\zeta_K$. Then, uniformly over all circle arcs $J$,
\[
 \#\{\mathfrak p\text{ unramified}:N\mathfrak p\le x,\
                          \theta_{\mathfrak p}\in J\}
 =|J|\operatorname{Li}_2(x)
       +O_{K,\eta}\bigl(\sqrt x\log^2(2x)\bigr)
 \quad(x\ge2).                                         \tag{NT-A04a}
\]
Any endpoint convention is allowed. In particular, uniformly for
$z\in\mathbb C$,
\[
 \#\{\mathfrak p\text{ unramified}:N\mathfrak p\le x,\
                                  \eta(\mathfrak p)=z\}
       \ll_{K,\eta}\sqrt x\log^2(2x).                    \tag{NT-A04b}
\]
The hypothesis on the character is essential for the angular law;
infinite order by itself is insufficient.

\paragraph{Exact classical locator and deductions.}
Rajan, \emph{Distribution of values of Hecke characters of infinite
order}, Acta Arithmetica85 (1998), Proposition4, printed
pp.289--290, gives (NT-A04a) with $\pi_K(x)$ in place of
$\operatorname{Li}_2(x)$. The character restriction is the one fixed
in Theorem1, printed p.286: in that paper a \emph{Dirichlet character}
means a finite-order Hecke character of $K$, and equivalence means
equality up to a power of the norm. The arguments use normalized
angles modulo one. Proposition4 has a constant depending only on
$K$ above a character-dependent threshold; allowing the constant
to depend on the fixed $\eta$ gives the displayed range $x\ge2$.
The finite ramification set is omitted throughout. Its analytic
premises are ordinary Hecke continuation and the precisely stated
power RH family; $M$ supplies them conditionally.

The prime-ideal estimate NT-C02 for the trivial extension $K/K$
replaces $\pi_K(x)$ by $\operatorname{Li}_2(x)$ with an error
$O_K(\sqrt x\log(2x))$. Shrinking or enlarging an arc and passing
to the limit for the finite collection of primes handles all
endpoint conventions, exactly as in NT-E06. For $|z|=1$, shrink
arcs to the singleton $z$; their main terms tend to zero while the
uniform error is fixed. This proves (NT-A04b). For $|z|\ne1$ the
set is empty.

\paragraph{Complete check of the norm-twist obstruction.}
Suppose instead that $\eta(\mathfrak p)=
\chi(\mathfrak p)(N\mathfrak p)^{it}$, where $\chi$ has finite order
$h$ and $t\ne0$. Set $\tau=ht\ne0$. Then the $h$th Fourier moment
of its prime angles is
\[
 \frac1{\pi_K(x)}\sum_{N\mathfrak p\le x}\eta(\mathfrak p)^h
       =\frac{x^{i\tau}}{1+i\tau}+o_{K,\chi,t}(1).       \tag{NT-A04c}
\]
Omitting the fixed ramification set does not affect this formula.
Here is a direct proof under the same available input $M$.
By NT-C02,
$\pi_K(u)=\operatorname{Li}_2(u)+O_K(\sqrt u\log(2u))$.
Abel summation therefore gives
\[
 \sum_{N\mathfrak p\le x}(N\mathfrak p)^{i\tau}
 =\int_2^x\frac{u^{i\tau}}{\log u}\,du
      +O_{K,\tau}(\sqrt x\log(2x)+1).
\]
Integration by parts with primitive
$u^{1+i\tau}/(1+i\tau)$ shows that the integral is
\[
 \frac{x^{1+i\tau}}{(1+i\tau)\log x}
           +O_\tau\!\left(\frac{x}{\log^2 x}+1\right).
\]
For the remainder, bound
$\int_2^xdu/\log^2u\ll x/\log^2x$, splitting at $\sqrt x$ if
desired. Since $\pi_K(x)\sim x/\log x$, (NT-A04c) follows.
Its absolute value tends to $(1+\tau^2)^{-1/2}>0$, whereas a
uniform angular distribution would force every nonzero Fourier
moment to tend to zero. Thus even under $M$ these infinite-order
characters do not obey (NT-A04a). The source exclusion prevents
the norm-twist poles from being silently treated as entire
nonprincipal factors.

\paragraph{Geometric dictionary and location.}
The general statement concerns rank-one unitary characters of the
idele class group. If the character is the unitary normalization
of an algebraic Hecke character, this is the Frobenius-value
distribution of the associated rank-one CM motive and its tensor
powers. A general unitary infinity type is not automatically
asserted to be an algebraic or \'{e}tale motive.
Work/run: \path{GRH_NIGHT/number_theory},
\path{hecke_shrinking_angles.tex}, NT-A04; primary source
S-NT-RAJAN, with the complete exception proof above.

\subsection*{NT-A05. Normalized shrinking arcs for infinite-order Hecke values}
Under the exact assumptions of NT-A04, fix
$c>0$ and $0<\alpha<1/2$, and set
$u_0=\max\{2,c^{1/\alpha}\}$.
For a center $\vartheta\in\mathbb R/\mathbb Z$ and $u\ge u_0$,
let $J_u(\vartheta)$ be the arc centered at $\vartheta$ of length
$cu^{-\alpha}$, with any fixed endpoint convention. Then, as $x\to\infty$,
\begin{align}
 &\#\{\mathfrak p\text{ unramified}:u_0<N\mathfrak p\le x,\
               \theta_{\mathfrak p}\in J_{N\mathfrak p}(\vartheta)\}
       \notag\\
 &\hspace{12mm}
   =c\int_{u_0}^x\frac{u^{-\alpha}}{\log u}\,du
       +O_{K,\eta,c,\alpha}
          \bigl(x^{3/4-\alpha/2}\log^2(2x)\bigr)
       \notag\\
 &\hspace{12mm}\sim\frac{c}{1-\alpha}\frac{x^{1-\alpha}}{\log x}.
                                                               \tag{NT-A05}
\end{align}
The error constant is uniform in the chosen center $\vartheta$.
All other parameters are fixed, and no assertion at $\alpha=1/2$
is made.

\paragraph{Complete proof.}
Set $h=x^{-1/4+\alpha/2}$. Partition $[u_0,x]$ into
successive multiplicative intervals of ratio $1+h$, with the
last endpoint truncated at $u_0$. The arcs are nested for the
fixed center. On each norm interval the moving prime count is
bounded above and below by the counts for its two endpoint arcs.
Apply (NT-A04a) at both norm endpoints, with a common constant
independent of $\vartheta$. The sum of all endpoint errors is
\[
 \ll_{K,\eta}\log^2(2x)
      \sum_{j\ge0}\sqrt{x(1+h)^{-j}}
 \ll_{K,\eta}\sqrt x\,h^{-1}\log^2(2x).
\]
The main brackets are the upper and lower Riemann sums for
$c\int_{u_0}^x u^{-\alpha}\,du/\log u$.
On an interval of ratio at most $1+h$, the endpoint values of
$u^{-\alpha}$ differ by a factor at most
$(1+h)^\alpha=1+O_\alpha(h)$. The difference of these Riemann
sums is consequently
$O_{c,\alpha}(h x^{1-\alpha})$.
With the chosen $h$, both power terms are $x^{3/4-\alpha/2}$.
This proves the equality in (NT-A05), with uniformity in the
center and the stated endpoint conventions.

Finally, integration by parts gives
$\int_{u_0}^x u^{-\alpha}\,du/\log u
\sim x^{1-\alpha}/((1-\alpha)\log x)$.
The ratio of the error to that main term is
$O(x^{-1/4+\alpha/2}\log^3(2x))=o(1)$.
This proves the asymptotic and the exact open range
$0<\alpha<1/2$. The window is normalized in the compact rank-one
torus; it is an arithmetic prime-counting assertion, and it does
not identify zeros of the Hecke function with independent random
angles. Work/run: \path{GRH_NIGHT/number_theory},
\path{hecke_shrinking_angles.tex}, NT-A05.

\subsection*{NT-E10. The degree-one endpoint reduces to finitely many CM norm forms}
Let $E/k$ have complex multiplication by the imaginary quadratic
field $K$, let $F=kK$, and put $D=|\operatorname{disc}K|$.
Fix $c>0$. This statement is unconditional. Consider good prime
ideals $\mathfrak p$ of $k$ of absolute residue degree one, so that
$N\mathfrak p=p$ is a rational prime, and fix either sign
$a=\pm a_{\mathfrak p}$. Write
\[
 \delta=2\sqrt p-a,\qquad B=4p-a^2 .
\]
Outside the finite ramification set of $F/k$, a prime satisfying
$0\le\delta<c/\sqrt p$ and $p>c/2$ must split in $F/k$.
At every such split prime there is a nonzero integer $b$ with
\[
 B=Db^2,\qquad
 \pi=\frac{a+b\sqrt{\operatorname{disc}K}}2\in\mathcal O_K,\qquad
 N_{K/\mathbb Q}\pi=p.                                  \tag{NT-E10a}
\]
For each fixed $c$, at all sufficiently large split good
degree-one primes the endpoint event is precisely
\[
 0\le2\sqrt p-a<\frac c{\sqrt p}
 \quad\Longleftrightarrow\quad
 a>0\ \hbox{and}\quad 0<Db^2<4c.                        \tag{NT-E10b}
\]
The equality $Db^2=4c$ never qualifies. Consequently, if $4c\le D$,
there are only finitely many degree-one good primes satisfying
either endpoint event. Outside the ramification set of $F/k$,
every possible exception then has $p\le c/2$.

\paragraph{Complete algebraic reduction.}
At a split good prime, all geometric CM endomorphisms specialize
to endomorphisms defined over $\mathbb F_p$. Their rational algebra
$K$ embeds in $\operatorname{End}^0(E_{\overline{\mathbb F}_p})$.
Here specialization is faithful: for $\ell\ne p$, good reduction
identifies the Tate modules, the action of an endomorphism on them
is compatible with specialization, and a nonzero endomorphism
cannot kill all $\ell$-power torsion. Indeed a nonzero endomorphism
of an elliptic curve is an isogeny with finite kernel.

The centralizer of this $K$ inside the rational endomorphism
algebra is $K$ itself. To check this without a classification of
the reduction, use the faithful action
\[
 \operatorname{End}^0(E_{\overline{\mathbb F}_p})
      \otimes_{\mathbb Q}\mathbb Q_\ell
       \hookrightarrow\operatorname{End}_{\mathbb Q_\ell}(V_\ell E)
\]
on the two-dimensional Tate module. The usual integral
injectivity follows from the fact that an endomorphism killing
all $\ell^n$-torsion factors through $[\ell^n]$; passing to the
$\ell$-adic completion and then inverting $\ell$ gives the displayed
map. The centralizer of
$K\otimes\mathbb Q_\ell$ there has dimension two: if the quadratic
algebra is a field the module has rank one over it, and if it is
split its two faithful eigenspaces each have dimension one.
Thus the rational centralizer has dimension at most two and
already contains the two-dimensional algebra $K$.

The $p$-power Frobenius endomorphism commutes with the specialized
CM endomorphisms, because they are defined over $\mathbb F_p$.
It therefore belongs to $K$. Its characteristic polynomial is
$X^2-a_{\mathfrak p}X+p$, so it is an algebraic integer of norm
$p$ and trace $a_{\mathfrak p}$. For the opposite sign replace
Frobenius by its negative, which still belongs to $\mathcal O_K$
and has norm $p$. Writing an algebraic integer in a quadratic
field in the usual integral basis gives
\[
 \pi=(a+b\sqrt{\operatorname{disc}K})/2,\qquad
 a\equiv b\,\operatorname{disc}K\pmod2,\qquad
 4p=a^2+Db^2.
                                                               \tag{NT-E10c}
\]
The integer $b$ is nonzero: otherwise $\pi$ would be a rational
algebraic integer and $p=\pi^2$ would not be prime.
This proves (NT-E10a). The argument uses only the Frobenius
characteristic polynomial, good-reduction specialization and
faithful Tate-module actions; no CM reduction-density theorem
or GRH is required.

\paragraph{Complete analytic inequality at the boundary.}
For a rational prime $p$, $2\sqrt p$ is not integral, so the Hasse
bound is strict and $B>0$. The identity
\[
 \delta\sqrt p=\frac{B}{2+a/\sqrt p}                    \tag{NT-E10d}
\]
holds, with a positive denominator smaller than $4$. Hence
$\delta\sqrt p<c$ implies $B<4c$ strictly.
If $p>c/2$, the same inequality forces $a>0$.
Conversely, suppose $a>0$, $p>c/2$, and $B<4c$.
Both $a$ and $2\sqrt p-c/\sqrt p$ are then positive.
Squaring the desired strict inequality shows the exact equivalence
\[
 \delta\sqrt p<c
 \quad\Longleftrightarrow\quad
 B<4c-\frac{c^2}{p}
 \quad\Longleftrightarrow\quad
 p>\frac{c^2}{4c-B}.                                   \tag{NT-E10e}
\]
There are only finitely many possibilities $B=Db^2<4c$.
If this set is nonempty, let $B_{\max}$ be its largest element and
take
\[
 p>\max\left\{\frac c2,\frac{c^2}{4c-B_{\max}}\right\}.
                                                               \tag{NT-E10f}
\]
Then the final inequality in (NT-E10e) holds simultaneously for
all those possibilities. If the set is empty, no split prime
qualifies. This proves (NT-E10b), including exclusion of equality.

For a good prime unramified and nonsplit in $F/k$, the two CM
eigenspaces are interchanged, so $a_{\mathfrak p}=0$, as proved in
NT-E08. Then $\delta\sqrt p=2p$, which fails the endpoint
inequality for $p>c/2$. The remaining ramified primes form a
fixed finite set. If $4c\le D$, (NT-E10a) gives $B\ge D\ge4c$
at every split prime, establishing the stated finiteness and
explicit restriction on the possible unramified exceptions.

\paragraph{What the endpoint gate says.}
The remaining possible rational primes belong to the finite
collection of quadratic norm families
\[
 p=\frac{a^2+Db^2}{4},\qquad
 0<|b|<2\sqrt{c/D},\qquad
 a>0,\qquad a\equiv b\,\operatorname{disc}K\pmod2.        \tag{NT-E10g}
\]
An actual prime of $k$ must still have the indicated
\emph{curve-specific Frobenius}: arbitrary prime values of these
forms are not asserted to have the desired trace for this
particular $E/k$. Thus the reduction neither discards
congruence/sign conditions nor proves infinitude of quadratic
polynomial prime values. It locates exactly the boundary left
open by NT-E09.

Geometrically, the endomorphism field forces the Frobenius norm
equation into the integral lattice $\mathcal O_K$ of the CM
torus. At this endpoint, the transverse lattice coordinate $b$
is bounded, rather than growing with the norm. Work/run:
\path{GRH_NIGHT/number_theory},
\path{cm_endpoint_gate.tex}, NT-E10; complete proof above,
using the CM eigenspace and good-reduction framework of NT-E06--08.

\subsection*{NT-A06--A07. Gaussian prime angles at almost every fixed center}

Write \(\mathbb T_G=\mathbf R/(\pi/2)\mathbf Z\), with probability
measure \(d\mu_G(\theta)=2\,d\theta/\pi\). For a nonzero ideal
\(\mathfrak a=(\alpha)\) of \(\mathbf Z[i]\), its angle is
\(\theta_{\mathfrak a}=\arg\alpha\pmod{\pi/2}\). Put
\[
 \Xi_m((\alpha))=(\alpha/\overline\alpha)^{2m},
 \qquad m\in\mathbf Z .
\]
These are well-defined conductor-one unitary Hecke characters.
For \(m\ne0\) their finite-part functions are entire, with completion
\[
 \pi^{-(s+2|m|)}\Gamma(s+2|m|)L(s,\Xi_m)
 =
 \pi^{-(1-s+2|m|)}\Gamma(1-s+2|m|)L(1-s,\Xi_m).
\]
The precise analytic assumption in this subsection is RH for
\(L(s,\Xi_m)\) for every positive integer \(m\); negative \(m\)
give the same functions. Conditional input \(M\) supplies this family.
The mean prime count below uses only the unconditional prime ideal
theorem for \(\mathbf Q(i)\). Input \(A\) alone does not supply the
nonzero angular infinity types.

\paragraph{Classical variance input.}
Rudnick--Waxman, \emph{Angles of Gaussian primes},
\path{arXiv:1705.07498v2}, Theorem~3.2, author p.~12,
proved in \S4.4, p.~20, gives the following. Fix nonzero nonnegative
functions \(f\in C_c^\infty((-1/2,1/2))\), even, and
\(\Phi\in C_c^\infty((1/2,1))\). Define
\[
 F_K(\theta)=\sum_{j\in\mathbf Z}
 f\left(\frac{K}{\pi/2}(\theta-j\pi/2)\right),\qquad
 P_{K,X}(\theta)=
 \sum_{\mathfrak p}\Phi(N\mathfrak p/X)\log N\mathfrak p\,
                         F_K(\theta_{\mathfrak p}-\theta).
\]
In the range used here, \(3\le K\le X\le K^2\), that theorem gives
\[
 \operatorname{Var}_{\mu_G}(P_{K,X})
       \ll_{f,\Phi}\frac{X}{K}(\log K)^2.             \tag{A06.1}
\]
Lemma~3.1, author pp.~11--12, gives
\[
 \mathbb E_{\mu_G}P_{K,X}
       \sim c_{f,\Phi}\frac XK,\qquad
 c_{f,\Phi}=\left(\int_{\mathbf R}f\right)
            \left(\int_0^\infty\Phi\right)>0.        \tag{A06.2}
\]
The asymptotic is uniform in \(K\): its entire \(K\)-dependence is
the exact factor \(1/K\) in the integral of \(F_K\).
The nonzero Fourier modes in \S4 use precisely the functions
\(L(s,\Xi_m)\) above; no zero-correlation or simplicity hypothesis
is part of the variance upper bound. The full asymptotic variance
conjecture of that paper is not used.

\paragraph{NT-A06: an explicit exceptional-measure formulation.}
Let \(J_K(\theta)\) be the open centered arc of length \(\pi/(2K)\).
For \(X\) sufficiently large and \(3\le K\le X\le K^2\),
let \(E_{K,X}\) be the set of centers whose arc \(J_K(\theta)\)
contains no prime ideal angle of norm in \((X/2,X)\), and let
\(E^{(1)}_{K,X}\) impose the same absence for absolute-degree-one
prime ideals. Then
\[
 \mu_G(E_{K,X})\ll\frac{K(\log K)^2}{X},\qquad
 \mu_G(E^{(1)}_{K,X})\ll\frac{K(\log K)^2}{X}+\frac1K .
                                                               \tag{A06.3}
\]
The implied constants are fixed once the two auxiliary functions
are fixed; no numerical constant is claimed.

\begin{proof}
For all sufficiently large \(X\), (A06.2) gives
\(\mathbb E P_{K,X}\ge c_{f,\Phi}X/(2K)>0\), uniformly in \(K\).
If \(E_{K,X}\) occurs, every summand in \(P_{K,X}\)
vanishes by its support. Chebyshev's inequality therefore gives
\[
 \mu_G\{P_{K,X}=0\}
 \le \frac{\operatorname{Var}(P_{K,X})}
           {(\mathbb E P_{K,X})^2}
 \ll \frac{K(\log K)^2}{X}.
 \]
The only prime ideals of absolute degree greater than one in
\(\mathbf Q(i)\) are those above inert rational primes, all with
angle zero. The set of centers whose arc \(J_K(\theta)\) contains
zero has measure \(1/K\). Outside that set any prime ideal found
by the first assertion has degree one. This proves the second
assertion, including the endpoint conventions (endpoints form a
measure-zero set).
\end{proof}

\paragraph{NT-A07: almost every center, at every sufficiently fine scale.}
Fix \(\epsilon>0\). Under the same angular RH family, for
\(\mu_G\)-almost every fixed \(\theta\), there is
\(K_0(\theta,\epsilon)\) such that for every real \(K\ge K_0\),
\(J_K(\theta)\) contains the angle of a degree-one Gaussian prime
ideal satisfying
\[
 N\mathfrak p\le C_\epsilon
              K(\log K)^3(\log\log K)^{1+\epsilon}.       \tag{A07.1}
\]
The constant \(C_\epsilon\) is independent of the center.
For almost every center this produces infinitely many such prime
ideals. The threshold may depend on the center. This is a
complete deduction from the cited variance bound, not a claim
of priority over other almost-sure formulations.

\begin{proof}
For integers \(j\) sufficiently large put
\[
 K_j=2^j,\qquad
 X_j=K_j(\log K_j)^3(\log\log K_j)^{1+\epsilon}.
\]
Then \(K_j\le X_j\le K_j^2\). Let \(E_j\) be the event that
no degree-one prime ideal of norm in \((X_j/2,X_j)\) has angle
in \(J_{K_j}(\theta)\). Formula (A06.3) gives
\[
 \mu_G(E_j)\ll_\epsilon
      \frac1{j(\log j)^{1+\epsilon}}+2^{-j}.
\]
This series converges. The first Borel--Cantelli lemma, which
requires no independence, says that almost every center belongs
to only finitely many \(E_j\). Consequently every sufficiently
large dyadic scale has such a prime.

For arbitrary sufficiently large real \(K\), choose \(j\) with
\(K\le 2^j<2K\). The centered arc \(J_{K_j}(\theta)\) is contained
in \(J_K(\theta)\). For all sufficiently large \(K\),
\(\log K_j\le2\log K\) and
\(\log\log K_j\le2\log\log K\); hence
\[
 X_j\le 2^{5+\epsilon}
           K(\log K)^3(\log\log K)^{1+\epsilon}.
\]
This proves (A07.1), for example with \(C_\epsilon=2^{5+\epsilon}\)
after increasing the threshold. This constant comes only from
the dyadic comparison, not from an explicit numerical variance
estimate. The selected dyadic primes have norm at least \(X_j/2\),
which tends to infinity, so infinitely many distinct prime ideals
occur.
\end{proof}

\paragraph{Quantifier boundary and geometric placement.}
Almost every fixed center cannot be replaced by a bound uniform
over all centers at this norm scale. If a nonreal Gaussian
integer \(a+ib\) has \(a,b>0\) and norm at most \(X\), then
\(\arg(a+ib)\ge\arctan(1/\sqrt X)\). Hence the open arc
\((0,\arctan(1/\sqrt X))\) contains no degree-one prime angle
of norm at most \(X\). A Gaussian prime may be represented in
the closed first quadrant after multiplication by a unit;
the degree-one primes on the axes do not give angles in that
open arc. This elementary gap explains why the exceptional set
and its order of quantifiers matter.

In the Grothendieck/Frobenius dictionary the prime ideals are
closed points of \(\operatorname{Spec}\mathbf Z[i]\), and the
\(\Xi_m\) are the integer weights of its compact CM torus.
The variance averages over the center of a torus neighborhood;
the first Borel--Cantelli step turns that averaged bound into
a statement about almost every fixed neighborhood system.
No independent random primes or independent zeros are assumed.
Run location: \path{GRH_NIGHT/number_theory};
the exact source is \path{S-NT-RW}, and the full new deductions
are (A06.3) and NT-A07 above.


\subsection*{NT-A08--A11. Prescribed ray class and Hecke angle, with uniform parameters}

Let \(K\) have degree \(n\) and absolute discriminant \(D\).
For a primitive unitary Hecke character \(\lambda\), use the standard
finite-part completion, with the complex factors split by duplication:
\[
 \Lambda_K(s,\lambda)=
 q_\lambda^{s/2}\prod_{j=1}^{n}\Gamma_{\mathbf R}(s+\kappa_{\lambda,j})
 L_K(s,\lambda),\qquad
 q_\lambda=D\,N\mathfrak f_\lambda,\qquad
 \Gamma_{\mathbf R}(s)=\pi^{-s/2}\Gamma(s/2).
                                                        \tag{A08.1}
\]
An irrelevant constant of absolute value one in the functional equation
is allowed. At a real place the parameter is
\(\epsilon+i\tau\), \(\epsilon\in\{0,1\}\); at a complex place the
two parameters are \(|k|/2+i\tau\) and \(|k|/2+i\tau+1\).
Put
\[
 {\cal C}(\lambda)=q_\lambda
                    \prod_{j=1}^{n}(3+|\kappa_{\lambda,j}|).
\]
These are the local Hecke conventions of Ramakrishnan--Valenza
Theorem~7-19 and its local-factor discussion, also checked in
MT-M28; Rajan, \S1, pp.~281--283, records the conductor and
infinity-type behavior under powers.

\paragraph{The exact classical prime-sum input.}
Iwaniec--Kowalski, \emph{Analytic Number Theory}, Theorem~5.15,
printed p.~114, equation~(5.56), states, under RH and the
Ramanujan bound for a function of degree \(d\), that
\[
 \psi_f(x)=r_fx+
 O\bigl(\sqrt x\log(2x)
             [\log{\cal C}(f)+d\log(2x)]\bigr),\qquad x\ge2,
                                                        \tag{A08.2}
\]
with an absolute constant. Here \(\psi_f\) uses the coefficients
of \(-L_f'/L_f\), \(r_f\) is the pole order at \(1\), and
\({\cal C}(f)\) is the analytic conductor at zero, defined in
\S5.1, (5.6)--(5.7), p.~95.
Its preceding proof sentence uses the truncated explicit formula
(5.53), pp.~112--113, and the local zero count (5.26).
This is a cited classical theorem; its preceding exercise and
full analytic development are not reproved here.

The book's definition in \S5.1, pp.~93--95, requires its infinity
parameters to be real or occur in conjugate pairs. An arbitrary
individual unitary Hecke function need not meet that convention.
The following complete reduction removes that restriction without
enlarging the order of the error term.

\paragraph{NT-A08: a uniform complex Hecke prime-sum bound.}
Suppose \(\lambda\) is not a pure unitary norm character and
RH holds for \(L_K(s,\lambda)\). Then, with an absolute constant,
\[
 \left|\sum_{N\mathfrak a\le x}
       \Lambda_K(\mathfrak a)\lambda(\mathfrak a)\right|
 \ll \sqrt x\log(2x)
          [\log{\cal C}(\lambda)+n\log(2x)].             \tag{A08.3}
\]
The value at an ideal not coprime to the conductor is zero.
Consequently, for any integral ideal \(\mathfrak f\) divisible by
\(\mathfrak f_\lambda\),
\[
 \sum_{\substack{N\mathfrak p\le x\\\mathfrak p\nmid\mathfrak f}}
            \lambda(\mathfrak p)
 \ll \sqrt x[\log{\cal C}(\lambda)+n\log(2x)]
                       +\log N\mathfrak f.             \tag{A08.4}
\]
For \(\lambda=1\), the same bound applies after subtracting
\(\operatorname{Li}_2(x)\) from the left side, using RH for
\(\zeta_K\). This principal case is handled separately.

\begin{proof}
Write
\[
 -\frac{L'}L(s,\lambda)=\sum_{r\ge1}b_r r^{-s}.
 \]
Grouping the ideal Euler factors by rational prime gives
\(|b_r|\le n\Lambda(r)\). In fact at \(r=p^a\) the coefficient
is a sum of terms \(f_{\mathfrak p}\log p\) times unit complex
numbers, over residue degrees \(f_{\mathfrak p}\mid a\), and
\(\sum_{\mathfrak p\mid p}f_{\mathfrak p}\le n\).
The local roots of the resulting degree-\(n\) rational-prime
Euler product have absolute value at most one.

For real \(t\) form
\[
 F_t(s)=L_K(s+it,\lambda)L_K(s-it,\overline\lambda).
 \]
Because \(\lambda\) is not pure norm, both factors are entire;
there are no shifted poles. Its degree is \(2n\), its conductor
is \(q_\lambda^2\), and its infinity parameters are
\(\kappa_{\lambda,j}+it\) and their conjugates. Its functional
equation has root number one. These statements include the
bad-prime Euler factors. The Hecke continuation, finite order
and boundary nonvanishing supply the remaining axioms of
Iwaniec--Kowalski. RH for the two translated factors follows
from RH for \(\lambda\) and complex conjugation.
For \(|t|\le2\),
\[
 \log{\cal C}(F_t)\le
           2\log{\cal C}(\lambda)+2n\log2 .
 \]
Its logarithmic derivative has coefficient
\(2\operatorname{Re}(b_r r^{-it})\).
Thus (A08.2), with pole order zero, bounds
\[
 R_t(y):=\sum_{r\le y}\operatorname{Re}(b_r r^{-it}),
 \qquad
 |R_t(y)|\le C\sqrt y\log(2y)
           [\log{\cal C}(\lambda)+n\log(2y)]              \tag{A08.5}
\]
uniformly in \(|t|\le2\).

Here are both steps recovering the imaginary part. Fix \(x\ge4\),
put \(a=\sqrt x\) and \(t=\pi/(2\log x)\), and hold this \(t\)
fixed while \(u\) ranges over \([a,x]\).
Denote the right side of (A08.5), with \(y=x\), by \(B_x\).
Write \(b_r=u_r+iv_r\). First partial summation gives
\[
 \begin{split}
 \sum_{a<r\le y}u_r\cos(t\log r)
  ={}&[R_0(y)-R_0(a)]\cos(t\log y)\\
    &+\int_a^y[R_0(u)-R_0(a)]
                       \frac{t\sin(t\log u)}u\,du .
 \end{split}
\]
The total variation of the cosine on this interval is at most
one, so this sum is \(O(B_x)\), uniformly for \(a\le y\le x\).
Since
\(\operatorname{Re}(b_r r^{-it})
   =u_r\cos(t\log r)+v_r\sin(t\log r)\),
it follows that
\[
 V(y):=\sum_{a<r\le y}v_r\sin(t\log r)=O(B_x).
 \]
Throughout \([a,x]\) the sine is at least \(1/\sqrt2\).
A second partial summation gives
\[
 \sum_{a<r\le x}v_r
   =\frac{V(x)}{\sin(t\log x)}
    +\int_a^x V(u)
                \frac{t\cos(t\log u)}
                     {u\sin^2(t\log u)}\,du
   =O(B_x),
 \]
because \(1/\sin(t\log u)\) has total variation at most
\(\sqrt2-1\). The initial segment has absolute sum at most
\(n\psi(\sqrt x)\ll n\sqrt x\), by the elementary Chebyshev
bound. It is absorbed in \(B_x\), since
\(\log{\cal C}(\lambda)\ge n\log3\).
The real part was already bounded by \(R_0(x)\).
This proves (A08.3); \(2\le x<4\) follows directly from the
coefficient bound.

Removing higher prime powers costs \(O(n\sqrt x\log(2x))\),
which is absorbed in (A08.3). Indeed, their underlying rational
primes are at most \(\sqrt x\), and the total weight above
each is at most \(n\log x\).
Abel summation removes the prime logarithm: if the weighted
error is \(O(\sqrt u\log(2u)E(u))\), with
\(E(u)=\log{\cal C}(\lambda)+n\log(2u)\), the endpoint error
is \(O(\sqrt xE(x))\), and its integral is bounded by
\[
 E(x)\int_2^x\frac{\log(2u)}
                     {\sqrt u(\log u)^2}\,du
                       \ll\sqrt x E(x).
 \]
Omitting additionally the primes dividing \(\mathfrak f\)
costs at most \(\omega(\mathfrak f)\le
\log N\mathfrak f/\log2\).
This proves (A08.4). For the principal character use the
Dedekind prime-count estimate NT-C02, with the same deletion.
In particular the principal factor is never put into \(F_t\).
\end{proof}

\paragraph{NT-A09: joint class and arc discrepancy.}
Let \(\mathfrak m=\mathfrak m_0\mathfrak m_\infty\) be a ray
modulus, \(G=\operatorname{Cl}_{\mathfrak m}(K)\) and \(h=|G|\).
Let \(\eta\) be unitary and not of the form finite order times
a unitary norm power. Put
\[
 \begin{aligned}
 \mathfrak f&=\operatorname{lcm}(\mathfrak m_0,\mathfrak f_\eta),
 & B_\eta&=4+\max_j|\kappa_{\eta,j}|,\\
 Q&=D\,N\mathfrak f\,B_\eta^n,
 & {\cal L}(x)&=\log Q+2n\log(3x).
 \end{aligned}                                                     \tag{A09.1}
\]
For \(a\in G\) and an arc \(J\subset\mathbf R/\mathbf Z\), let
\(N_{a,J}(x)\) count primes \(\mathfrak p\nmid\mathfrak f\)
of norm at most \(x\) with ray class \(a\) and
\(\eta(\mathfrak p)\in e^{2\pi iJ}\).
Assume RH for the primitive Hecke functions attached to
\[
             \chi\eta^m,\qquad
     \chi\in\widehat G,\quad m=0,1,2,\ldots .             \tag{A09.2}
\]
Conditional input \(M\) suffices. Then
\[
 N_{a,J}(x)=\frac{|J|}{h}\operatorname{Li}_2(x)
           +O\bigl(\sqrt x\log(2x){\cal L}(x)\bigr),
 \qquad x\ge2,                                         \tag{A09.3}
\]
with an absolute constant, uniformly in \(K,\eta,\mathfrak m,a,J\).
The same bound holds if only absolute-degree-one primes are counted.
In particular there is no factor \(h\) in the error.

\begin{proof}
For \(\chi\in\widehat G\), its finite conductor divides
\(\mathfrak m_0\); its infinity type has only real signs.
The finite conductor of \(\chi\eta^m\) divides \(\mathfrak f\).
At complex places its angular and imaginary parameters grow
by \(m\); at real places its parity remains in \(\{0,1\}\).
The conventions (A08.1) therefore give directly
\[
 {\cal C}(\chi\eta^m)\le Q(m+1)^n,\qquad m\ge0.           \tag{A09.4}
\]
For \(m>0\), none of these characters is pure norm. Otherwise,
if \(\chi\eta^m=|\cdot|^{it}\), the character
\(\eta|\cdot|^{-it/m}\) would have finite order: its \(m\)-th
power is the finite-order character \(\chi^{-1}\).
This contradicts the hypothesis on \(\eta\).
For \(m=0\), only \(\chi=1\) is principal.

Apply (A08.4), treating that principal character separately.
The additional \(\log N\mathfrak f\) is absorbed by
\(\sqrt x\log Q\). For each \(m\ge0\) we obtain
\[
 \sum_{\substack{N\mathfrak p\le x\\\mathfrak p\nmid\mathfrak f}}
       \chi(\mathfrak p)\eta(\mathfrak p)^m
 =\mathbf1_{\chi=1,m=0}\operatorname{Li}_2(x)
 +O\left(\sqrt x\{\log Q+n\log[(m+1)(2x)]\}\right).
                                                               \tag{A09.5}
\]
The ray-class indicator is
\[
 \mathbf1_{[\mathfrak p]=a}
       =\frac1h\sum_{\chi\in\widehat G}
                        \overline{\chi(a)}\chi(\mathfrak p).
 \]
Its \(h\) terms cancel the prefactor \(1/h\) when the errors
in (A09.5) are summed in absolute value. Thus the number
\(N_a(x)\) in class \(a\) has main term \(\operatorname{Li}_2(x)/h\),
and every nonzero angular Fourier sum in that class has
the same error bound without \(h\).

Apply the Erd\H{o}s--Tur\'an inequality S-NT-ET to the finite
list of angles in this single class, with \(H=\lceil x\rceil\).
Its zero-frequency truncation error is
\(O(N_a(x)/(H+1))=O(n)\): at most \(n\) prime ideals lie over
each rational prime of size at most \(x\).
For \(1\le m\le H\),
\(\log Q+n\log[(m+1)(2x)]\ll{\cal L}(x)\).
Summing the \(1/m\) weights gives (A09.3), first for half-open
arcs. Singleton bounds follow by decreasing arcs to an endpoint;
changing either endpoint convention changes the count by at
most the same error. This also covers arcs crossing zero.
Finally there are at most \(n\pi(\sqrt x)=O(n\sqrt x)\)
prime ideals of absolute residue degree at least two and norm
at most \(x\). Their deletion is absorbed in the error.
\end{proof}

\paragraph{NT-A10: shrinking arcs with separate class-number dependence.}
Fix \(c>0\) and \(0<\alpha<1/2\), and let
\(u_0=\max\{2,c^{1/\alpha}\}\).
For a fixed center let the arc at norm \(u>u_0\) have length
\(c u^{-\alpha}\). Let \(N_{a,c,\alpha}(x)\) count precisely
the primes \(\mathfrak p\nmid\mathfrak f\) with
\(u_0<N\mathfrak p\le x\), ray class \(a\), and angle in
that arc at \(u=N\mathfrak p\).
For \(x>\max\{4,u_0\}\), this joint count satisfies
\[
 \begin{split}
 N_{a,c,\alpha}(x)
  ={}&\frac c h\int_{u_0}^x\frac{u^{-\alpha}}{\log u}\,du\\
   &+O_{\alpha,c}\left(
       \sqrt x\log(2x){\cal L}(x)
       +h^{-1/2}x^{3/4-\alpha/2}\sqrt{{\cal L}(x)}
                  \right).                            
 \end{split} \tag{A10.1}
\]
The constant is uniform in the center and in the parameters
of (A09.1). For varying fields and ray moduli, a sufficient
uniform asymptotic regime is
\[
 h\,{\cal L}(x)(\log x)^2=o(x^{1/2-\alpha}).              \tag{A10.2}
\]
For fixed data this holds automatically. The same assertions
hold after restricting to degree-one primes.

\begin{proof}
Set \(T(x)=\log(2x){\cal L}(x)\). Partition \((u_0,x]\)
geometrically with ratio \(1+\delta\), \(0<\delta\le1/2\).
Nested centered arcs bound the moving count on each interval
between its two fixed-endpoint arc counts. Apply (A09.3) to
each difference of prime-counting functions. The geometric
sum of endpoint errors is
\[
 O\bigl(\delta^{-1}\sqrt x\,T(x)\bigr),
 \]
since \(\sum \sqrt{u_j}\ll\delta^{-1}\sqrt x\), and
\(T(u_j)\le T(x)\).
The two main terms bracket the displayed integral, and their
difference from it is
\[
 O_\alpha\left(\frac c h\delta
           \int_{u_0}^x\frac{u^{-\alpha}}{\log u}\,du\right)
 \ll_{\alpha,c}\frac c h\delta
                   \frac{x^{1-\alpha}}{\log x}
 \]
for \(x\) above a threshold depending only on \(c,\alpha\).
This follows from the ratio \((1+\delta)^\alpha=1+O_\alpha(\delta)\);
the last incomplete interval obeys the same bounds.

Write \(A=\sqrt x\,T(x)\) and
\(B=(c/h)x^{1-\alpha}/\log x\).
Taking \(\delta=\min\{1/2,\sqrt{A/B}\}\) gives
\[
 A/\delta+B\delta\ll A+\sqrt{AB}.
 \]
Indeed when \(\sqrt{A/B}\le1/2\) both terms equal \(\sqrt{AB}\);
otherwise \(B<4A\). Since
\[
 \sqrt{AB}
  =\sqrt{c/h}\,x^{3/4-\alpha/2}
       \sqrt{\frac{\log(2x)}{\log x}{\cal L}(x)},
 \]
this proves (A10.1). Its main term is asymptotic to
\(c\,x^{1-\alpha}/[h(1-\alpha)\log x]\).
The ratios of the two error terms to this main term are,
up to fixed constants,
\[
 h{\cal L}(x)(\log x)^2x^{\alpha-1/2},
 \qquad
 \left[h{\cal L}(x)(\log x)^2x^{\alpha-1/2}\right]^{1/2}.
 \]
They tend to zero under (A10.2).
Using the degree-one version of (A09.3) in the identical
partition proves the last assertion.
\end{proof}

\paragraph{NT-A11: a least prime in a prescribed class and arc.}
Let \(0<\mu\le1\) be the length of a fixed arc. Put
\[
 d=\frac h\mu,\qquad
 V=\log(eQ)+n\log(e\,n\,d).
 \]
Under (A09.2), there is an absolute constant \(C>0\) such
that every ray class and every arc of length \(\mu\) contains
a degree-one prime \(\mathfrak p\nmid\mathfrak f\) with
\[
                       N\mathfrak p\le C d^2 V^6 .      \tag{A11.1}
\]
No numerical value of \(C\) is claimed.

\begin{proof}
Use (A09.3) with its degree-one version. Let \(A_0\) be an
absolute error constant there, and put \(x=C d^2 V^6\).
We have \(d\ge1\), \(V\ge n\), \(n\log d\le V\) and
\(n\log V\le2V\). The last inequality follows by writing
\(V=nv\), since \(v\ge\log(e\,n)\) and
\(\log V=\log n+\log v\le2v\).
Consequently, for \(C\ge4\),
\[
 \log x\le(\log C+14)V/n,\qquad
 \log(2x)\le(\log C+15)V/n,\qquad
 {\cal L}(x)\le(32+2\log C)V .
 \]
Also \(\operatorname{Li}_2(x)\ge x/(2\log x)\) for \(x\ge4\).
Thus a sufficient condition for the main term
\(\operatorname{Li}_2(x)/d\) to exceed the error is
\[
 \sqrt C>
 2A_0(\log C+14)(\log C+15)(32+2\log C).
 \]
An absolute \(C\) satisfying this exists, because \(\sqrt C\)
eventually dominates that cubic polynomial in \(\log C\).
The resulting count is positive, proving (A11.1).
\end{proof}

\paragraph{Geometric and run placement.}
The finite condition is a Frobenius fiber in the ray-class
cover of \(\operatorname{Spec}\mathcal O_K[1/\mathfrak f]\);
the angular condition lies in a compact rank-one Hecke torus.
Their joint character algebra is
\(\widehat G\times\mathbf Z\), and (A09.2) is exactly its
Hecke \(L\)-function family. For normalized algebraic Hecke
characters this is the tensor-character description of a CM
torus with a finite Artin factor. For a general unitary
infinity type no algebraic-motive realization is asserted.
The formulas keep the finite mass \(1/h\), the finite
conductor, the infinity parameters and the field degree
separate. They strengthen the separate class and angle
formulations NT-C06/A04--A05. Run:
\path{GRH_NIGHT/number_theory}; full new proofs are
NT-A08--A11 above, with classical input S-NT-IK and S-NT-ET.



\subsection*{NT-C09--C13. Smaller class-group generators and exact computational input}

Let \(K\) be a number field of degree \(n\ge2\), discriminant
\(D=|\operatorname{disc}K|\), and \(r_1\) real places. All class
groups below are ordinary ideal class groups. Put
\[
 G=\operatorname{Cl}(K),\qquad
 G(T)=\langle[\mathfrak p]:N\mathfrak p\le T\rangle .
\]
The analytic input used for the uniform bounds is
\[
 {\cal R}_K:\quad
 \text{RH for }\zeta_{\mathbf Q}
 \text{ and for every }L_K(s,\chi),\quad \chi\in\widehat G,
                                                               \tag{C09.1}
\]
including \(\chi=1\), whose function is \(\zeta_K\).
Equivalently, it suffices to assume RH for \(\zeta_{\mathbf Q}\)
and \(\zeta_{H_K}\), where \(H_K\) is the ordinary Hilbert class
field: the all-prime abelian factorization NT-H11 has no
factor poles in the open critical strip. Thus \(M\) supplies
(C09.1). The designated input \(A\) suffices if \(H_K/\mathbf Q\)
is abelian; abelianness of \(K/\mathbf Q\) alone does not give
that conclusion.

\paragraph{NT-C09 and NT-C10: the precise published bounds.}
Greni\'e--Molteni, \emph{Breaking the 4 barrier for the bound of
a generating set of the class group}, Math.\ Comp.\ 94 (2025),
3145--3163, Theorems~1--4, author pp.~2--3, give the following
conclusions under (C09.1):
\[
 \begin{array}{ll}
 G(B_n)=G,&
 B_n=\left(4-\dfrac1{2n}+\dfrac1{2n^2}\right)(\log D)^2,\\[5pt]
 G(B'_n)\supseteq G^2,&
 B'_n=\left(4-\dfrac2{3n}+\dfrac3{8n^2}\right)(\log D)^2 .
 \end{array}                                                    \tag{C09.2}
\]
Without the respective \(n^{-2}\) corrections, these conclusions
hold for every \(n\ge5\). For \(n=2,3,4\), the verified ranges
are \(\log D\le l_-\) or \(\log D\ge l_+\), with
\[
 \begin{array}{c|rrr}
 n&2&3&4\\ \hline
 l_-&63&353&833\\
 l_+&12184&18455&27911\\ \hline
 l'_-&65&379&993\\
 l'_+&7968&9832&13108
 \end{array} .                                                  \tag{C09.3}
\]
The primed rows apply to \(G^2\).
The excluded middle intervals are intervals where the sharper
bound is not asserted, rather than a list of actual counterexamples.
For odd \(|G|\), squaring is an automorphism and both primed
bounds generate all of \(G\).

The proofs are in \S4, author pp.~10--12, and \S5,
pp.~12--13. Their analytic inputs are visible in \S2,
pp.~3--5 (class-character and Dedekind zero sums), and
Proposition~1, pp.~8--9 (RH for \(\zeta_{\mathbf Q}\)).
The small-discriminant input is their 2018
\emph{Explicit bounds for generators of the class group},
Theorem~3.5, equation~(3.10), author p.~6, using the
same class-character and Dedekind inputs.
The threshold computations, Odlyzko bounds and certified zero
tables remain credited parts of the published proof; they were
not numerically rerun here.

The containment in (C09.2) is the precise conclusion proved
at the start of the 2025 paper's \S5. Its introductory definition
of \(T'_K\) uses the phrase of generating \(G^2\); arbitrary
small prime classes need not themselves belong to \(G^2\).
We use the unambiguous subgroup-containment formulation throughout.

\paragraph{NT-C11: a finite certificate for a quotient, with exclusions.}
Here a smaller RH input suffices. Fix a subgroup \(H\le G\).
Assume RH only for the nontrivial class characters \(\chi\)
that are trivial on \(H\). No assumption on \(\zeta_K\) or
\(\zeta_{\mathbf Q}\) is needed in this paragraph.

Let \(L>0\), and let \(F:\mathbf R\to\mathbf R_{\ge0}\) be even, continuous,
compactly supported in \([-L,L]\), with nonnegative Fourier
transform. Assume the regularity needed in the Weil--Poitou
formula, for instance that \(F\) is piecewise \(C^1\) and
\((F(0)-F(x))/x\) is of bounded variation on \((0,\infty)\).
Define
\[
 \begin{split}
 I(F)&=\int_0^\infty
                  \frac{F(0)-F(x)}{2\sinh(x/2)}\,dx,\\
 J(F)&=\int_0^\infty\frac{F(x)}{2\cosh(x/2)}\,dx,\\
 W_F(\mathfrak p)&=\log N\mathfrak p
       \sum_{m\ge1}\frac{F(m\log N\mathfrak p)}
                              {(N\mathfrak p)^{m/2}},\\
 R_K(F)&=F(0)\{\log D-(\gamma+\log8\pi)n\}
                                      +nI(F)-r_1J(F).
 \end{split}                                                    \tag{C11.1}
\]
Let \(E\) be any set of prime ideals to be omitted. If
\[
       2\sum_{\mathfrak p}
             (-1)^{\mathbf1_E(\mathfrak p)}W_F(\mathfrak p)
                          >R_K(F),                            \tag{C11.2}
\]
then the classes of primes \(\mathfrak p\notin E\) with
\(N\mathfrak p<e^L\), together with \(H\), generate \(G\).
Only finitely many terms in (C11.2) are nonzero. This is a
quotient formulation of Greni\'e--Molteni (2025),
Theorem~5, author pp.~13--14, with the support stated in
the logarithmic variable in which the test function is evaluated.

\begin{proof}
Suppose the stated classes and \(H\) generate a proper
subgroup \(U<G\). A nontrivial character of \(G/U\), inflated
to \(G\), gives a character \(\chi\ne1\) trivial on \(H\)
and on every selected prime. Write
\[
       \phi(s)=\int_{\mathbf R}F(x)e^{(s-1/2)x}\,dx.
\]
The Weil--Poitou identity, in the normalization of the
2025 source (3), author p.~3, gives
\[
 2\sum_{\mathfrak p}\log N\mathfrak p
       \sum_{m\ge1}
       \frac{\operatorname{Re}(\chi(\mathfrak p)^m)
                          F(m\log N\mathfrak p)}
                              {(N\mathfrak p)^{m/2}}
       =R_K(F)-\sum_{\rho_\chi}\phi(\rho_\chi).
                                                               \tag{C11.3}
\]
The supplied RH hypothesis and Fourier positivity make every
zero term nonnegative. A prime outside \(E\) contributing to
the sum has norm below \(e^L\), hence character value one.
Indeed continuity and compact support force \(F(L)=0\),
so the endpoint creates no term. At a prime in \(E\),
\(\operatorname{Re}(\chi(\mathfrak p)^m)\ge-1\).
Since \(F\ge0\), the left side of (C11.3) is at least
the left side of (C11.2), whereas its right side is at most
\(R_K(F)\). This contradicts the strict inequality and
proves the assertion. The zero sum can equivalently be
read as the symmetric-height limit in the classical formula;
its terms are nonnegative here, so the comparison is valid.
\end{proof}

There is a directly computable sufficient version. Let \(T>1\),
\(L=\log T\), \(F(x)=(1-|x|/L)_+\), and write
\[
 {\mathfrak G}=\sum_{j\ge0}\frac{(-1)^j}{(2j+1)^2}.
\]
Then it suffices that
\[
 \begin{split}
 2\sum_{(N\mathfrak p)^m<T}
   (-1)^{\mathbf1_E(\mathfrak p)}
   \frac{\log N\mathfrak p}{(N\mathfrak p)^{m/2}}
        \left(1-\frac{m\log N\mathfrak p}{\log T}\right)
 >{}&\log D-n(\gamma+\log8\pi)-\frac{\pi r_1}{2}\\
    &+\frac{\pi^2 n/2+4{\mathfrak G}r_1}{\log T}.
 \end{split}                                                    \tag{C11.4}
\]
For \(H=1\) and \(E=\varnothing\), this is exactly the
Belabas--Diaz y Diaz--Friedman criterion, Corollary~2.2,
author pp.~6--7; its rounded Theorem~1.1 is on p.~3.
Here \(H=1\) denotes the trivial subgroup.

\begin{proof}
The Fourier transform of the triangular function is
\(L(\sin(tL/2)/(tL/2))^2\), with the value at zero taken
by continuity. It is nonnegative and meets all the
regularity assumptions. Also
\[
 1-F(x)\le x/L\quad(x\ge0),\qquad
 \int_0^\infty\frac{dx}{2\cosh(x/2)}=\frac\pi2.
\]
Thus
\[
 I(F)\le\frac{\pi^2}{2L},\qquad
 J(F)\ge\frac\pi2-\frac{4{\mathfrak G}}L.
 \]
For completeness, the constants follow by expanding
\[
 \frac1{2\sinh(x/2)}
       =\sum_{j\ge0}e^{-(j+1/2)x},\qquad
 \frac1{2\cosh(x/2)}
       =\sum_{j\ge0}(-1)^j e^{-(j+1/2)x}
\]
and integrating \(x\) times the series. The first integration
is monotone; the second is justified by the integrable
absolute sum of the \(x\)-weighted integrals. Their values
are \(\pi^2/2\) and \(4{\mathfrak G}\), respectively.
These inequalities bound \(R_K(F)\) above by the right
side of (C11.4); now apply (C11.2).
\end{proof}

\paragraph{NT-C12: what can actually be enumerated in polynomial time.}
Suppose an integral basis of the maximal order \(\mathcal O_K\)
and its full integer multiplication table are supplied.
Under (C09.1), one can deterministically enumerate, in
polynomial time in this input's bit length, prime ideals
whose classes generate \(G\), all of norm at most \(B_n+2\).
With the additional promise that \(|G|\) is odd, \(B'_n+2\)
may be used instead. The harmless additive two permits
certified rational rounding of the displayed logarithmic
bound; no unproved separation bound from an integer is needed.
The output is a list of integral ideal lattices.
This does not compute a relation lattice, the class number,
units, or a presentation of the class group.

\begin{proof}
First, the discriminant is computable by exact linear algebra
from the trace form. If all multiplication constants have
absolute value at most \(2^b\), the matrix of multiplication
by a basis element has that entry bound. Therefore each
trace-form entry has absolute value at most \(n^2 2^{2b}\).
Hadamard's inequality gives
\[
       \log_2 D\le 2bn+\frac{5n}{2}\log_2 n.                  \tag{C12.1}
\]
Consequently any integer \(T=O(1+(\log D)^2)\) is polynomial
in the supplied input length.

An integer \(T\) with \(B_n\le T\le B_n+2\) is obtainable
in polynomial time. Compute a rational upper approximation
to \(B_n\) with error less than \(1/2\), and take its ceiling.
For example, write \(D=2^m u\), \(1\le u<2\), and evaluate
\[
 \log u=2\sum_{j\ge0}
       \frac{1}{2j+1}\left(\frac{u-1}{u+1}\right)^{2j+1},
\]
as well as the same series for \(\log2\).
The ratios have absolute value at most \(1/3\); geometric
tail bounds give certified accuracy \(O((1+\log D)^{-1})\)
for \(\log D\), dividing the error allowance for \(\log2\)
by \(m+1\) before forming \(m\log2+\log u\).
This uses polynomially many rational bit operations. Squaring
and multiplying by the fixed rational coefficient gives
the required approximation. The same procedure applies
to \(B'_n\).

We describe an unconditional procedure to list all prime
ideals with norm at most this integer \(T\).
List rational primes \(p\le T\), even by trial division.
For each such \(p\), form the \(n\)-dimensional algebra
\[
                   A=\mathcal O_K/p\mathcal O_K
\]
over \(\mathbf F_p\), using the given table. Its Frobenius
\(F:a\mapsto a^p\) is \(\mathbf F_p\)-linear and is
computed by repeated squaring. Choose \(e\) with \(p^e\ge n\).
Then
\[
                   \operatorname{Nil}(A)=\ker F^e.           \tag{C12.2}
\]
Indeed a nilpotent in an \(n\)-dimensional algebra has
nilpotence index at most \(n\), whereas Frobenius is
injective on every finite residue field. Form the reduced
quotient \(R=A/\operatorname{Nil}(A)\). It is a product of
finite fields, and
\[
          B=\{b\in R:b^p=b\}\simeq\mathbf F_p^{\,g},
                                                               \tag{C12.3}
\]
where \(g\le n\) is the number of factors.

A basis of \(B\) is obtained by solving linear equations.
Use its elements successively to split \(R\) into the
simultaneous eigenspaces of their multiplication maps.
Every eigenvalue belongs to \(\mathbf F_p\), so it is enough
to scan \(a=0,\ldots,p-1\) and compute kernels of
\(m_b-aI\) on each current component. These maps are
diagonalizable because \(b^p=b\) and \(X^p-X\) has distinct
roots. The common eigenspaces are exactly the field factors:
a basis of \(\mathbf F_p^g\) separates every pair of
coordinates. All kernels are ideals, and the nonzero
resulting components give the maximal ideals of \(R\)
as kernels of the corresponding projections. Pulling
these back to \(A\) lists all maximal ideals \(W\subset A\).

If \(f=\dim_{\mathbf F_p}(A/W)\), the associated prime
ideal of \(\mathcal O_K\) has norm \(p^f\); retain it
exactly when \(p^f\le T\). It is explicitly the inverse
image of the subspace \(W\). To output an integer basis,
put a basis of \(W\) in reduced row echelon form, lift its
entries to \(\{0,\ldots,p-1\}\), and add \(p\) times each
coordinate vector at a nonpivot position. These \(n\)
vectors generate precisely the inverse-image lattice:
subtracting the lifted rows using the integral pivot
coordinates leaves zero pivot coordinates and multiples
of \(p\) elsewhere. Their determinant has absolute value
\(p^f\), as also follows from the quotient size.

All algebra dimensions are at most \(n\). The number of
scanned scalars is polynomial in \(p,n\), and \(p\le T\)
is polynomial in the input length by (C12.1). All modular
entries use \(O(\log T)\) bits. There are at most \(n\)
prime ideals above each \(p\), so the number of output
ideals is at most \(nT\), with polynomial total encoding
length. This proves the claimed deterministic bit
complexity. The generation conclusion follows from
(C09.2) because \(T\ge B_n\), or \(T\ge B'_n\) with
odd class number.
\end{proof}

The supplied maximal-order basis is essential to this
formulation. Starting only with a defining polynomial
does not supply it for free. The proof makes no complexity
claim about finding that basis, identifying relations,
or proving an odd-class-number promise.

\paragraph{NT-C13: finite Frobenius detection of ordinary class-field quotients.}
Under (C09.1), let \(E/K\) be abelian, unramified at every
finite place and split at every real place. If all primes
of norm at most \(B_n\) split completely in \(E\), then
\(E=K\). If the smaller bound \(B'_n\) is used instead,
then \(\operatorname{Gal}(E/K)\) has exponent at most two.
The correction-free bounds in (C09.3) have the same
respective conclusions.

\begin{proof}
Ordinary global class field theory gives a surjection
\[
       G\twoheadrightarrow\operatorname{Gal}(E/K)
\]
sending each prime class to its Frobenius. Complete
splitting kills the indicated small-prime subgroup.
At \(B_n\) that subgroup is \(G\), so the quotient is
trivial. At \(B'_n\) it contains \(G^2\); consequently
the quotient is a quotient of \(G/G^2\) and has exponent
at most two. The sharper-range versions follow identically.
\end{proof}

In geometric language, \(G=\operatorname{Pic}
(\operatorname{Spec}\mathcal O_K)\), and these tests concern
abelian finite \'etale arithmetic covers with the ordinary
real-place splitting condition. The finite prime-sum
certificate proves generation of a prescribed quotient
by its closed-point Frobenius elements. The enumeration
proof exposes the precise integral-algebra input needed
to make these points concrete. Run:
\path{GRH_NIGHT/number_theory}; primary inputs
\path{S-NT-GM2025-CG}, \path{S-NT-BDDF},
and \path{S-NT-GM2018-CG}.



\subsection*{NT-A12--A14. A nonabelian Frobenius class and a Hecke angle}

Let \(L/K\) be a finite Galois extension with group \(G\), and let
\(C\subset G\) be one conjugacy class. Choose \(g\in C\) and put
\[
 \begin{aligned}
 H&=\langle g\rangle,\quad h=|H|,\quad F=L^H,\quad d=[F:K],\\
 \delta&=\frac{|C|}{|G|},\quad n_F=[F:\mathbf Q],\quad
                              n_L=[L:\mathbf Q]=h n_F .
 \end{aligned}
                                                               \tag{A12.1}
\]
Let \(\eta\) be a unitary Hecke character of \(K\) which is not a
finite-order character times a unitary norm power. Write
\(\nu=\eta\circ N_{F/K}\), and use the primitive character attached
to every product below. Let \(S\) be a finite set of prime ideals of
\(K\) containing the primes ramified in \(L/K\) and the support of
\(\mathfrak f_\eta\). Extra prescribed omissions may be included.
Set
\[
 \begin{aligned}
 \mathfrak s_F&=\prod_{\mathfrak q\mid\mathfrak p,\,
                                      \mathfrak p\in S}\mathfrak q,\\
 B_\nu&=4+\max_j|\kappa_{\nu,j}|,\\
 Q_C&=D_L
       \bigl(N_F\mathfrak f_\nu\,N_F\mathfrak s_F\bigr)^h
       B_\nu^{n_L},\\
 {\cal L}_C(x)&=\log Q_C+2n_L\log(3x).
 \end{aligned}                                                 \tag{A12.2}
\]
Here \(D_L\) is the absolute discriminant; the gamma parameters
\(\kappa_{\nu,j}\) have precisely convention (A08.1).
All conductors and norms in (A12.2) are over \(F\).

The exact analytic input is RH for
\[
 L_F(s,\chi\nu^m),\qquad
       \chi\in\widehat H,\quad m=0,1,2,\ldots,                 \tag{A12.3}
\]
where class field theory regards \(\chi\) as a finite-order Hecke
character for the cyclic extension \(L/F\). In particular the
\((\chi,m)=(1,0)\) member is \(\zeta_F\).
Conditional input \(M\) supplies this family. No holomorphy or RH
for an irreducible nonabelian Artin function is assumed.

\paragraph{NT-A12: uniform joint discrepancy.}
For an arc \(J\subset\mathbf R/\mathbf Z\), count prime ideals
\(\mathfrak p\notin S\) with \(N_K\mathfrak p\le x\),
Frobenius class \(C\), and \(\eta(\mathfrak p)\in e^{2\pi iJ}\).
Then
\[
 N_{C,J}(x)
   =\delta |J|\operatorname{Li}_2(x)
      +O\bigl(\delta\sqrt x\log(2x){\cal L}_C(x)\bigr),
       \qquad x\ge2.                                      \tag{A12.4}
\]
The constant is absolute, uniformly in all the displayed data,
the arc and its endpoint conventions. The same bound holds when
only absolute-degree-one primes of \(K\) are counted.

\begin{proof}
We first check the possible norm degeneracy. The diagonal idele
map \(j:C_K\longrightarrow C_F\) satisfies
\[
 N_{F/K}(j(c))=c^d,\qquad |j(c)|_F=|c|_K^d.
\]
If \(\nu=\theta|\cdot|_F^{it}\) with \(\theta\) of finite order,
composition with \(j\) gives
\((\eta|\cdot|_K^{-it})^d=\theta\circ j\).
If the latter has order dividing \(r\), the former character
raised to \(dr\) is trivial. This contradicts the hypothesis on
\(\eta\). Thus \(\nu\) also is not finite order times norm.
For \(m>0\), an equality
\(\chi\nu^m=|\cdot|_F^{it}\) would make
\((\nu|\cdot|_F^{-it/m})^m=\chi^{-1}\), giving the same
contradiction. Hence every \(m>0\) member of (A12.3) is entire,
with no pure-norm pole. For \(m=0\) only \(\chi=1\) is principal.
The principal member is treated separately by NT-A08.

Next we retain the sum of conductors instead of replacing each
one by its largest value. Locally, the conductor of a product
divides the product of conductors, and taking a power cannot
increase the conductor. Therefore
\[
 \mathfrak f_{\chi\nu^m}\mid\mathfrak f_\chi\mathfrak f_\nu,
 \qquad
 {\cal C}(\chi\nu^m)
   \le D_F N_F\mathfrak f_\chi N_F\mathfrak f_\nu
                       [B_\nu(m+1)]^{n_F}.               \tag{A12.5}
\]
The gamma-factor inequality follows directly from (A08.1):
a finite-order character changes only the real parities, while
the imaginary and complex angular parameters of \(\nu^m\)
grow linearly in \(m\). It also holds at \(m=0\).
The conductor--discriminant formula for the cyclic extension
\(L/F\) gives the exact identity
\[
       D_L=D_F^h\prod_{\chi\in\widehat H}
                                      N_F\mathfrak f_\chi .
                                                               \tag{A12.6}
\]
This includes all ramified local factors. The precise classical
locator is Neukirch, Chapter VII, (11.9), printed p.~534,
with the supplied conductor/induction proof on pp.~533--534
and the one-dimensional conductor identification (11.10),
p.~535; see S-NT-NEUKIRCH. Summing the logarithms in (A12.5)
thus yields
\[
 \sum_{\chi\in\widehat H}\log{\cal C}(\chi\nu^m)
 \le \log D_L+h\log N_F\mathfrak f_\nu
          +n_L\log B_\nu+n_L\log(m+1).                 \tag{A12.7}
\]

Every character in (A12.3) is unramified outside the primes
dividing \(\mathfrak s_F\). Indeed a local norm sends units
to units, and the characters of \(L/F\) ramify only over the
ramification locus already in \(S\). The prime version of
NT-A08, deleting the additional primes in \(\mathfrak s_F\),
therefore gives
\[
 \begin{split}
 P_{\chi,m}(x)
  &:=
    \sum_{\substack{N_F\mathfrak q\le x\\
                      \mathfrak q\nmid\mathfrak s_F}}
                 \chi(\mathfrak q)\nu(\mathfrak q)^m\\
  &=\mathbf1_{\chi=1,m=0}\operatorname{Li}_2(x)
    +O\left(\sqrt x\{\log{\cal C}(\chi\nu^m)
                       +n_F\log(2x)\}
                   +\log N_F\mathfrak s_F\right).
 \end{split}                                                   \tag{A12.8}
\]
Its constant is absolute. The last error is simply the
number of additional omitted prime ideals, at most
\(\log N_F\mathfrak s_F/\log2\); a conductor exponent is
not required to count these omissions. Formula (A12.8)
also follows directly from the proof of (A08.4).

We now establish the induction identity used at primes:
\[
       \mathbf1_C
       =\delta\sum_{\chi\in\widehat H}
                         \overline{\chi(g)}\,
                            \operatorname{Ind}_H^G\chi .
                                                               \tag{A12.9}
\]
Orthogonality on \(H\) says
\(\sum_\chi\overline{\chi(g)}\chi(y)=h\mathbf1_{\{g\}}(y)\).
By the definition of induced class functions,
\(\operatorname{Ind}_H^G\mathbf1_{\{g\}}\) has value
\(|Z_G(g)|/h\) on \(C\) and zero elsewhere. Since
\(\delta |Z_G(g)|=1\), (A12.9) follows.

For a prime \(\mathfrak p\notin S\), the cosets in \(G/H\)
fixed by a Frobenius element correspond exactly to prime
ideals \(\mathfrak q\) of \(F\) over \(\mathfrak p\) with
relative residue degree one. The trace of
\(\operatorname{Ind}_H^G\chi\) is the sum of the corresponding
\(\chi(\operatorname{Frob}_{\mathfrak q})\).
For these primes,
\(\nu(\mathfrak q)=\eta(\mathfrak p)\), because the ideal
norm is \(\mathfrak p\). Use the same arithmetic Frobenius
convention in this identity and in class field reciprocity.
Multiplying (A12.9) by \(\eta(\mathfrak p)^m\) and summing
therefore gives
\[
 T_m^C(x):=
 \sum_{\substack{N_K\mathfrak p\le x,\ \mathfrak p\notin S\\
                  \operatorname{Frob}_{\mathfrak p}\in C}}
                        \eta(\mathfrak p)^m
 =\delta\sum_{\chi\in\widehat H}
              \overline{\chi(g)}P_{\chi,m}(x)
           +O(\delta n_L\sqrt x).                       \tag{A12.10}
\]
To verify the last error, the extra primes in the right-hand
sum have relative residue degree at least two, hence lie
over primes of \(K\) with norm at most \(\sqrt x\).
There are at most \(d[K:\mathbf Q]\sqrt x=n_F\sqrt x\)
such primes in \(F\). Each of the \(h\) character sums has
terms of modulus one, so their total contribution is at
most \(\delta h n_F\sqrt x\). Primes in \(S\) have been
deleted on both sides; no Frobenius value is assigned to
a ramified prime.

Combining (A12.7)--(A12.10), and using
\(\log Q_C\ge n_L\log4\), proves
\[
 T_m^C(x)=\delta\mathbf1_{m=0}\operatorname{Li}_2(x)
     +O\left(\delta\sqrt x
        \{\log Q_C+n_L\log[(m+1)(2x)]\}\right).
                                                               \tag{A12.11}
\]
All constants remain absolute. Notice that the \(h\)
summands have already been absorbed in the exact
discriminant and in \(n_L\); they are not replaced by
an additional factor \(|G|\).

Apply Erd\H{o}s--Tur\'an, S-NT-ET Theorem III, to the angles
of the primes with Frobenius class \(C\), with cutoff
\(\lceil x\rceil\). The zero-frequency truncation error
is \(O([K:\mathbf Q])\), since there are at most
\([K:\mathbf Q]x\) prime ideals of norm at most \(x\).
It is absorbed because
\(\delta n_L=|C|[K:\mathbf Q]\ge[K:\mathbf Q]\).
For the other terms, the braces in (A12.11) are
\(O({\cal L}_C(x))\) throughout \(1\le m\le\lceil x\rceil\);
their harmonic weights sum to \(O(\log(2x))\).
The \(m=0\) instance of (A12.11) supplies the main term
in (A12.4). Decreasing arcs to an endpoint bounds its
singleton count by the same error; thus all endpoint
conventions are allowed.
Finally the number of absolute-degree-at-least-two primes
of \(K\) with norm at most \(x\) is at most
\([K:\mathbf Q]\sqrt x\le\delta n_L\sqrt x\).
Deleting them proves the degree-one version.
\end{proof}

\paragraph{NT-A13: a shrinking angle in the specified class.}
Fix \(c>0\), \(0<\alpha<1/2\) and
\(u_0=\max\{2,c^{1/\alpha}\}\). Fix an angular center.
Let \(N_{C,c,\alpha}(x)\) count exactly the primes
\(\mathfrak p\notin S\) with \(u_0<N_K\mathfrak p\le x\),
Frobenius class \(C\), and angle in the centered arc
of length \(c(N_K\mathfrak p)^{-\alpha}\). Under (A12.3),
\[
 \begin{split}
 N_{C,c,\alpha}(x)
   ={}&\delta c\int_{u_0}^x\frac{u^{-\alpha}}{\log u}\,du\\
     &+O_{c,\alpha}\left(
       \delta\sqrt x\log(2x){\cal L}_C(x)
       +\delta x^{3/4-\alpha/2}\sqrt{{\cal L}_C(x)}
                     \right),\qquad x>\max\{4,u_0\}.
 \end{split}                                                   \tag{A13.1}
\]
The constant is uniform in the center and all the field and
character parameters. A sufficient relative asymptotic regime is
\[
         {\cal L}_C(x)(\log x)^2=o(x^{1/2-\alpha}).
                                                               \tag{A13.2}
\]
For fixed arithmetic data this holds automatically, and the
main term is asymptotic to
\(\delta c x^{1-\alpha}/[(1-\alpha)\log x]\).
The same assertions hold for absolute-degree-one primes.

\begin{proof}
Partition \((u_0,x]\) geometrically with ratio
\(1+\varepsilon\), \(0<\varepsilon\le1/2\).
The centered arcs are nested, so their endpoint arcs bracket
the moving count on each subinterval. Apply (A12.4) to both
endpoint prime counts and sum. Since
\(\sum_j\sqrt{u_j}\ll\sqrt x/\varepsilon\), the total
discrepancy is at most a constant times \(A/\varepsilon\),
where
\[
 A=\delta\sqrt x\log(2x){\cal L}_C(x),\qquad
 B=\frac{\delta c x^{1-\alpha}}{\log x}.
\]
The two main terms bracket the integral in (A13.1).
Their difference is \(O_{c,\alpha}(B\varepsilon)\):
on one mesh interval the ratio of the arc lengths is at most
\((1+\varepsilon)^\alpha=1+O_\alpha(\varepsilon)\), and
\(\int_{u_0}^x u^{-\alpha}du/\log u
 \ll_{c,\alpha}x^{1-\alpha}/\log x\).
The finite initial range in this integral bound depends only
on \(c,\alpha\). Its possible contribution to the count is
absorbed uniformly using
\([K:\mathbf Q]\le\delta n_L\le\delta\log Q_C/\log4\).

Choose \(\varepsilon=\min\{1/2,\sqrt{A/B}\}\).
Then \(A/\varepsilon+B\varepsilon\ll A+\sqrt{AB}\), and
\[
 \sqrt{AB}
  =\delta\sqrt c\,x^{3/4-\alpha/2}
       \sqrt{\frac{\log(2x)}{\log x}{\cal L}_C(x)}.
\]
This proves (A13.1). After division by the main term,
the two errors are bounded, up to fixed \(c,\alpha\)
constants, by
\[
 R(x)={\cal L}_C(x)(\log x)^2x^{\alpha-1/2},
 \qquad \sqrt{R(x)}.
\]
They tend to zero under (A13.2). Applying the identical
partition to the degree-one version of (A12.4) proves
the final assertion.
\end{proof}

\paragraph{NT-A14: a least prime with both conditions.}
For an arc of length \(0<\mu\le1\), put
\[
       V_C=\log(eQ_C)+n_L\log(e n_L/\mu).
\]
There is an absolute constant \(C_0>0\) such that (A12.3)
implies the existence of an absolute-degree-one prime
\(\mathfrak p\notin S\), with Frobenius class \(C\) and
angle in this arc, satisfying
\[
                  N_K\mathfrak p\le C_0\mu^{-2}V_C^6 .
                                                               \tag{A14.1}
\]
No numerical value of \(C_0\) is claimed.

\begin{proof}
Let \(A_0\) be an absolute error constant in (A12.4) for
degree-one primes. Set \(x=C_0\mu^{-2}V_C^6\), with
\(C_0\ge4\). We have \(V_C\ge n_L\),
\(n_L\log(1/\mu)\le V_C\), and
\(n_L\log V_C\le2V_C\).
For the last inequality write \(V_C=n_L v\), note
\(v\ge\log(e n_L)\), and use
\(\log(n_L v)\le2v\). Consequently
\[
 \begin{split}
 \log x&\le(\log C_0+14)V_C/n_L,\\
 \log(2x)&\le(\log C_0+15)V_C/n_L,\\
 {\cal L}_C(x)&\le(32+2\log C_0)V_C .
 \end{split}
\]
Since \(\operatorname{Li}_2(x)\ge x/(2\log x)\),
the main term \(\delta\mu\operatorname{Li}_2(x)\)
exceeds the error in (A12.4) as soon as
\[
 \sqrt{C_0}>
  2A_0(\log C_0+14)(\log C_0+15)(32+2\log C_0).
\]
Here the common factor \(\delta\) cancels; dropping the
remaining factor \(n_L^{-2}\le1\) makes this condition
only stronger. An absolute \(C_0\) satisfying it exists.
The count is then positive, proving (A14.1).
\end{proof}

\paragraph{Geometric meaning and exact work placement.}
The finite condition is a conjugacy-class Frobenius condition
for the finite \'{e}tale \(G\)-cover over the arithmetic open
\(\operatorname{Spec}\mathcal O_K\setminus S\). Passing to
\(F=L^{\langle g\rangle}\) is the corresponding intermediate
cover; its cyclic character decomposition is the explicit
induction (A12.9). The second condition is a rank-one
compact Hecke torus coordinate, pulled back by the norm.
For normalized algebraic Hecke characters this is a finite
Artin factor coupled to a CM torus representation. For
arbitrary unitary infinity type no algebraic-motive
realization is asserted. Formula (A12.3) gives the exact
rank-one analytic family replacing any unproved general
Artin-holomorphy input.
Full proofs: \path{frobenius_hecke_angles.tex},
NT-A12--A14, also in \path{section.tex}.
Run: \path{GRH_NIGHT/number_theory}.
Classical premises: NT-A08 (S-NT-IK), S-NT-ET and
S-NT-NEUKIRCH; the simultaneous formulations are proved
here, with no literature-priority claim.


\subsection*{NT-E11. CM Hasse deficits inside a nonabelian Frobenius class}

Let \(E/K\) be an elliptic curve all of whose complex
multiplication is defined over \(K\), and let \(\eta\) be its
normalized unitary CM Hecke character. Thus at good primes
\[
 a_{\mathfrak p}
     =\sqrt{N_K\mathfrak p}
          \{\eta(\mathfrak p)+\overline{\eta(\mathfrak p)}\}.
                                                               \tag{E11.1}
\]
The CM character premise and its precise classical locator
are those of NT-E06 (S-NT-BFK, Proposition 2.5 and the CM
identification in the footnote to Proposition 4.7).
Its nonzero angular infinity type excludes finite order
times a unitary norm power.
Choose any finite Galois \(L/K\), any conjugacy class \(C\),
and the notation \(\delta,Q_C,{\cal L}_C\) of NT-A12.
Include in \(S\) all bad-reduction primes, as well as the
ramification and conductor primes required there.
Assume exactly the cyclic rank-one RH family (A12.3).
Conditional input \(M\) supplies it.

Fix \(c>0\), \(-1/2<\lambda<1/2\), and put
\[
 \begin{aligned}
 \alpha&=\frac14-\frac{\lambda}{2},\qquad
 u_0=\max\{2,c^{1/(1/2-\lambda)}\},\\
 w(u)&=\frac1\pi
       \arccos\left(1-\frac c2 u^{\lambda-1/2}\right)
                  \quad(u\ge u_0).
 \end{aligned}                                             \tag{E11.2}
\]
For either sign, count the primes \(\mathfrak p\notin S\)
with \(u_0<N_K\mathfrak p\le x\), Frobenius class \(C\), and
\[
       0\le 2\sqrt{N_K\mathfrak p}\mp a_{\mathfrak p}
                      <c(N_K\mathfrak p)^\lambda .
\]
Their number \(D_{C;c,\lambda}^{\pm}(x)\) satisfies, for
\(x>\max\{4,u_0\}\),
\[
 \begin{split}
 D_{C;c,\lambda}^{\pm}(x)
   ={}&\delta\int_{u_0}^x\frac{w(u)}{\log u}\,du\\
     &+O_{c,\lambda}\left(
        \delta\sqrt x\log(2x){\cal L}_C(x)
        +\delta x^{5/8+\lambda/4}\sqrt{{\cal L}_C(x)}
                        \right).                         
 \end{split} \tag{E11.3}
\]
All other dependence is explicit in (A12.2). The same formula
holds when only absolute-degree-one primes are counted.
A sufficient uniform relative asymptotic regime is
\[
       {\cal L}_C(x)(\log x)^2=o(x^{1/4+\lambda/2}).
                                                               \tag{E11.4}
\]
For fixed arithmetic data it gives
\[
 D_{C;c,\lambda}^{\pm}(x)
    \sim \frac{\delta\sqrt c}{\pi(3/4+\lambda/2)}
                    \frac{x^{3/4+\lambda/2}}{\log x}.
                                                               \tag{E11.5}
\]

In particular let \(M_C^\pm(x)\) count the primes outside \(S\)
of norm at most \(x\), Frobenius class \(C\), and
\(a_{\mathfrak p}=\pm\lfloor2\sqrt{N_K\mathfrak p}\rfloor\).
For fixed \(E,K,L,C,S\),
\[
 \begin{split}
 M_C^\pm(x)
     &=\frac{\delta}{\pi}
          \int_2^x\frac{u^{-1/4}}{\log u}\,du
                +O_{E,K,L,C,S}\bigl(x^{5/8}\sqrt{\log(2x)}\bigr),\\
 M_C^\pm(x)&\sim
             \frac{4\delta}{3\pi}\frac{x^{3/4}}{\log x}.
 \end{split}                                                   \tag{E11.6}
\]
The integral, rather than only its leading asymptotic term, is
retained in the error-bearing formula. No general non-CM
elliptic \(L\)-function RH is used.

\begin{proof}
Write \(\eta(\mathfrak p)=e^{2\pi i\theta_{\mathfrak p}}\).
For the positive sign the displayed deficit condition says
\[
 \cos(2\pi\theta_{\mathfrak p})
                      >1-\frac c2 q^{\lambda-1/2},
                \qquad q=N_K\mathfrak p .
 \]
It is exactly the arc about zero of normalized length \(w(q)\),
with its outer endpoints omitted. The negative sign gives
the arc of the same length centered at \(1/2\). Equality
at the Hasse endpoint is allowed and is already in these
arcs. All endpoint conventions are covered by NT-A12.

Here is the required uniform control of the moving length.
Set \(v=cu^{-2\alpha}\), so \(0<v\le1\). Direct integration
and differentiation give
\[
 w(u)=\frac1\pi\int_0^v\frac{dz}{\sqrt{4z-z^2}},
 \qquad
 -u w'(u)=\frac{\alpha\sqrt v}{\pi\sqrt{1-v/4}},
 \qquad w(u)\ge\frac{\sqrt v}{\pi}.
 \]
Consequently
\[
 w(u)\asymp_{c,\lambda}u^{-\alpha},\qquad
    0\le-\frac{u w'(u)}{w(u)}
                    \le\frac{2\alpha}{\sqrt3}.            \tag{E11.7}
\]
On a multiplicative interval of ratio at most
\(1+\varepsilon\), with \(0<\varepsilon\le1/2\),
the endpoint lengths therefore have ratio
\(1+O_\lambda(\varepsilon)\).

Partition \((u_0,x]\) by such intervals. The arcs centered
at zero, or at \(1/2\), are nested. Their two endpoint
counts, using (A12.4), bracket the moving count. As in the
complete proof of NT-A13, the total error is
\[
 \begin{gathered}
 O_{c,\lambda}(A/\varepsilon+B\varepsilon),\\
 A=\delta\sqrt x\log(2x){\cal L}_C(x),\qquad
 B=\frac{\delta x^{1-\alpha}}{\log x}.
 \end{gathered}                                           \tag{E11.8}
\]
For completeness the first term follows from
\(\sum_j\sqrt{u_j}\ll\varepsilon^{-1}\sqrt x\).
For the second, (E11.7) bounds the difference of the
upper and lower integral sums by
\[
 O_\lambda\left(\delta\varepsilon
                    \int_{u_0}^x\frac{w(u)}{\log u}\,du\right)
             \ll_{c,\lambda} B\varepsilon .
 \]
Their common limiting integral is precisely the first
term in (E11.3), so no small-angle expansion has been
substituted into this comparison. The initial bounded
range is uniform after absorption by
\([K:\mathbf Q]\le\delta n_L\) and
\(\log Q_C\ge n_L\log4\).

Take \(\varepsilon=\min\{1/2,\sqrt{A/B}\}\).
Then (E11.8) is \(O_{c,\lambda}(A+\sqrt{AB})\), and
\[
 \sqrt{AB}\ll
          \delta x^{3/4-\alpha/2}\sqrt{{\cal L}_C(x)}
       =\delta x^{5/8+\lambda/4}\sqrt{{\cal L}_C(x)}.
 \]
This proves (E11.3). Using the degree-one form of (A12.4)
throughout proves its degree-one version.
Since \(w(u)\sim(\sqrt c/\pi)u^{-\alpha}\), the integral
has asymptotic
\[
 \frac{\delta\sqrt c}{\pi(1-\alpha)}
                          \frac{x^{1-\alpha}}{\log x}.
 \]
The relative errors are bounded by a constant times
\(R(x)+\sqrt{R(x)}\), where
\(R(x)={\cal L}_C(x)(\log x)^2x^{\alpha-1/2}\).
This proves (E11.4)--(E11.5).

For (E11.6), take \(c=1\), \(\lambda=0\). For any integer
\(a\le y\), \(a=\lfloor y\rfloor\) is equivalent to
\(a>y-1\); this includes integral \(y\).
Apply this with \(a=a_{\mathfrak p}\), or
\(-a_{\mathfrak p}\), and \(y=2\sqrt q\).
The floor event is exactly the strict deficit condition
in (E11.3), even when \(q\) is a square.
The primes of norm exactly two contribute at most
\([K:\mathbf Q]\le\delta n_L\), absorbed by its error.
The elementary expansion on \(u\ge2\) is
\[
 w(u)=\frac1\pi u^{-1/4}+O(u^{-3/4}).
 \]
Its integral remainder is \(O(x^{1/4})\) and is absorbed
by the first error in (E11.3).
For fixed arithmetic data
\({\cal L}_C(x)=O_{E,K,L,C,S}(\log(2x))\).
Thus the second error is \(O(x^{5/8}\sqrt{\log(2x)})\);
the first, \(O(\sqrt x\log^2(2x))\), is smaller for
sufficiently large \(x\), and is absorbed after adjusting
the fixed-data constant. Finally
\(\int_2^x u^{-1/4}du/\log u
 \sim(4/3)x^{3/4}/\log x\), which proves (E11.6).
\end{proof}

\paragraph{Scope and geometric placement.}
The finite étale cover may have an arbitrary nonabelian
Galois group; for example it may be a division-field
cover attached to another, non-CM elliptic curve.
The trace in this theorem is still the trace of the
specified CM curve \(E\). The proof couples the finite
Frobenius condition to its rank-one CM torus by the cyclic
intermediate covers of NT-A12. It does not infer GRH for
the standard \(L\)-function of that second curve, or a
Koblitz prime-order statement. The assumption that all
CM is defined over \(K\) is retained; no unproved descent
factor is inserted into this simultaneous formulation.
At \(\lambda=-1/2\), the positive exponent in (E11.4)
vanishes and the argument gives no asymptotic; the
finite norm-form gate in NT-E10 remains applicable.
For general \(\lambda\), the exact arcsine integral in
(E11.3) is essential to the stated error, as explained
in NT-E09. This is a newly proved formulation, with no
claim of literature priority.
Work/run: \path{GRH_NIGHT/number_theory},
\path{cm_frobenius_deficits.tex}, NT-E11, also
\path{section.tex}; full analytic input and its
classical locators are (A12.3), NT-A08 and NT-E06.


\subsection*{NT-A15--A17. Several Hecke coordinates and the exact independence gate}

Retain \(L/K,G,C,g,H,F,h,\delta,n_L,S,\mathfrak s_F\)
from NT-A12. Fix \(r\ge1\) unitary Hecke characters
\(\eta_1,\ldots,\eta_r\) of \(K\). Require that the character
\[
             \prod_{i=1}^r\eta_i^{m_i},
                    \qquad 0\ne\mathbf m\in\mathbf Z^r,  \tag{A15.1}
\]
is never finite order times a unitary norm power.
Enlarge \(S\) to contain all their conductor primes.
Write \(\nu_i=\eta_i\circ N_{F/K}\), and put
\[
 \begin{aligned}
 \mathfrak f_*&=\operatorname{lcm}_i\mathfrak f_{\nu_i},\\
 B_*&=4+\sum_{i=1}^r\max_j|\kappa_{\nu_i,j}|,\\
 Q_*&=D_L(N_F\mathfrak f_*\,N_F\mathfrak s_F)^h B_*^{n_L},\\
 {\cal L}_*(x)&=\log Q_*+2n_L\log(3x).
 \end{aligned}                                             \tag{A15.2}
\]
The exact RH family is
\[
 L_F\left(s,\chi\prod_{i=1}^r\nu_i^{m_i}\right),
       \qquad \chi\in\widehat H,\quad\mathbf m\in\mathbf Z^r .
                                                               \tag{A15.3}
\]
It is a subfamily of \(M\). The mixed products in this family
are part of the hypothesis; separate assertions only for
the individual coordinates are not substituted for them.

\paragraph{Classical discrepancy input, with a precise locator.}
For \(N\) points \(\boldsymbol\theta_j\in[0,1)^r\), the
Erd\H{o}s--Tur\'an--Koksma inequality bounds their normalized
rectangular discrepancy by
\[
 \begin{gathered}
 R(\mathbf m)=\prod_i\max(1,|m_i|),\\
 D_N\le(3/2)^r\left\{\frac2{H+1}
     +\sum_{0<\|\mathbf m\|_\infty\le H}
       \frac{1}{R(\mathbf m)}
       \left|\frac1N\sum_{j=1}^N
                  e^{2\pi i\mathbf m\cdot\boldsymbol\theta_j}
       \right|\right\}.
 \end{gathered}                                             \tag{A15.4}
\]
The exact statement is Dick--Pillichshammer,
\emph{Digital Nets and Sequences}, author preprint
Theorem 3.26, printed/PDF p.~84 (S-NT-DP).
That page refers its proof to Drmota--Tichy,
\emph{Sequences, Discrepancies and Applications},
\S1.2.2. The theorem page was image checked; the
referred proof is a cited premise, not a newly audited
argument in this run.

\paragraph{NT-A15: joint boxes and a least prime.}
For a product \(J=J_1\times\cdots\times J_r\) of circular
arcs, let \(\mu=\prod_i|J_i|\). Under (A15.1)--(A15.3),
the count of primes outside \(S\), Frobenius class \(C\),
norm at most \(x\), and coordinate angles in \(J\) is
\[
 N_{C,J}^{(r)}(x)=\delta\mu\operatorname{Li}_2(x)
       +O_r\bigl(\delta\sqrt x
                    (\log(2x))^r{\cal L}_*(x)\bigr).
                                                               \tag{A15.5}
\]
The constant depends only on \(r\); all other dependence
is displayed. Endpoint conventions are arbitrary, and
the same bound holds for absolute-degree-one primes.
For \(\mu>0\), this box contains such a degree-one prime
of norm at most
\[
 C_r\mu^{-2}V_*^{\,2r+4},\qquad
       V_*=\log(eQ_*)+n_L\log(e n_L/\mu),                 \tag{A15.6}
\]
where \(C_r\) depends only on \(r\).

\begin{proof}
Apply the diagonal-idele argument of NT-A12 to each
nonzero monomial in (A15.1). Its norm pullback is still
not finite order times norm. Multiplication by a
finite-order \(\chi\) therefore cannot make it pure norm.
Only \(\chi=1,\mathbf m=0\) is principal.
For negative powers the finite conductor is unchanged;
for any product its conductor divides the least common
multiple of the individual conductors, since all factors
are trivial on each sufficiently deep local unit group.
Thus
\[
 {\cal C}\left(\chi\prod_i\nu_i^{m_i}\right)
  \le D_FN_F\mathfrak f_\chi N_F\mathfrak f_*
       [B_*(1+\|\mathbf m\|_\infty)]^{n_F}.              \tag{A15.7}
\]
At a complex place the absolute value of the sum of
angular parameters is at most the sum of their absolute
values; imaginary parameters add. At real places the
parity remains zero or one. These facts prove the
gamma-factor part of (A15.7), including negative \(m_i\).
Summing logarithms over \(\chi\) uses exactly
\(D_L=D_F^h\prod_\chi N_F\mathfrak f_\chi\).

The whole cyclic induction and omitted-prime argument
(A12.8)--(A12.10) now applies to the monomial.
In particular its relative-degree correction remains
\(\delta n_L\sqrt x\), independently of \(\mathbf m\).
It gives the complete Fourier estimate
\[
 \begin{split}
 T_{\mathbf m}^C(x)
  &:=
    \sum_{\substack{N_K\mathfrak p\le x,\ \mathfrak p\notin S\\
                     \operatorname{Frob}_{\mathfrak p}\in C}}
                          \prod_i\eta_i(\mathfrak p)^{m_i}\\
  &=\delta\mathbf1_{\mathbf m=0}\operatorname{Li}_2(x)
      +O\left(\delta\sqrt x\{
         \log Q_*+n_L\log[(1+\|\mathbf m\|_\infty)(2x)]
                                  \}\right).
 \end{split}                                                   \tag{A15.8}
\]
The constant here is absolute: any dependence on \(r\)
has been put in \(B_*\), not in the Hecke prime estimate.
The principal factor is again handled separately.

Use (A15.4) in unnormalized form with \(H=\lceil x\rceil\).
The sum of the frequency weights is exactly
\[
 \sum_{0<\|\mathbf m\|_\infty\le H}\frac1{R(\mathbf m)}
       =\left(1+2\sum_{j=1}^H\frac1j\right)^r-1
                     \ll_r(\log(2H))^r .
 \]
The zero-frequency truncation contributes
\(O_r([K:\mathbf Q])\), absorbed as in NT-A12 by
\(\delta n_L\le\delta\log Q_*/\log4\).
All the braces in (A15.8) are \(O({\cal L}_*(x))\)
on this frequency cube. These facts prove (A15.5)
for half-open boxes not crossing the circle origins.
A circular box is a disjoint union of at most \(2^r\)
such boxes. Shrinking a coordinate interval onto a
specified endpoint bounds each coordinate-face count
by the same error. Summing the finitely many faces
permits all endpoint conventions.
The number of omitted absolute-degree-at-least-two
primes is at most \([K:\mathbf Q]\sqrt x\), again
absorbed. If the prime list is empty, the unnormalized
discrepancy assertion follows directly from the
zero-frequency estimate, so no division by zero is used.

For (A15.6), set \(V=V_*\), \(n=n_L\), and
\(x=C_r\mu^{-2}V^{2r+4}\). The inequalities
\(V\ge n\), \(n\log(1/\mu)\le V\), and
\(n\log V\le2V\), proved in NT-A14, imply
\[
 \begin{aligned}
 \log x&\le(\log C_r+4r+10)V/n,\\
 \log(2x)&\le(\log C_r+4r+11)V/n,\\
 {\cal L}_*(x)&\le(8r+24+2\log C_r)V.
 \end{aligned}
\]
Let \(A_r\) be a constant in the degree-one error
(A15.5). Since \(\operatorname{Li}_2(x)\ge x/(2\log x)\),
positivity follows, after cancellation of \(\delta\),
from
\[
 \sqrt{C_r}>
  2A_r(\log C_r+4r+10)(\log C_r+4r+11)^r
                    (8r+24+2\log C_r).
 \]
An appropriate \(C_r\ge4\) exists for every fixed \(r\).
The discarded factor \(n^{-(r+1)}\) is at most one.
This proves the least-prime assertion.
\end{proof}

\paragraph{NT-A16: simultaneously shrinking boxes.}
Fix \(c_i>0\), \(\alpha_i>0\), and put
\[
 A=\sum_i\alpha_i<\tfrac12,\qquad
 c_*=\prod_i c_i,\qquad
 u_0=\max\{2,c_1^{1/\alpha_1},\ldots,c_r^{1/\alpha_r}\}.
 \]
Fix every coordinate center and count \(u_0<N_K\mathfrak p\le x\)
in Frobenius class \(C\) whose \(i\)-th angular coordinate
lies in the arc of length \(c_i(N_K\mathfrak p)^{-\alpha_i}\).
Then
\[
 \begin{split}
 N_{C,\mathbf c,\boldsymbol\alpha}(x)
   ={}&\delta c_*\int_{u_0}^x\frac{u^{-A}}{\log u}\,du\\
     &+O_{\mathbf c,\boldsymbol\alpha,r}\left(
        \delta\sqrt x(\log(2x))^r{\cal L}_*(x)\right.\\
     &\hspace{38mm}\left.
        +\delta x^{3/4-A/2}
             (\log(2x))^{(r-1)/2}\sqrt{{\cal L}_*(x)}
                                    \right).
 \end{split}                                                   \tag{A16.1}
\]
The error is uniform in the centers and all arithmetic
parameters already displayed. The degree-one version
holds too. A sufficient relative asymptotic regime is
\[
          {\cal L}_*(x)(\log x)^{r+1}
                                  =o(x^{1/2-A}).           \tag{A16.2}
\]

\begin{proof}
The boxes are nested because every side length decreases.
Partition \((u_0,x]\) with multiplicative ratio
\(1+\varepsilon\), \(0<\varepsilon\le1/2\).
Applying (A15.5) to both endpoints gives a total error
at most a constant times \(A_0/\varepsilon\), where
\[
 A_0=\delta\sqrt x(\log(2x))^r{\cal L}_*(x).
 \]
Indeed the endpoint square roots form the same geometric
series as in NT-A13. The volume is exactly \(c_*u^{-A}\);
its ratio on a mesh interval is at most
\((1+\varepsilon)^A=1+O_A(\varepsilon)\).
Consequently the integral bracketing error is
\(O_{\mathbf c,\boldsymbol\alpha}(B_0\varepsilon)\), with
\[
                  B_0=\frac{\delta c_* x^{1-A}}{\log x}.
 \]
The canonical \(u_0\) bounds the finite initial range
using only the fixed \(\mathbf c,\boldsymbol\alpha\);
as before its field degree is absorbed by
\([K:\mathbf Q]\le\delta n_L\).
Choose \(\varepsilon=\min\{1/2,\sqrt{A_0/B_0}\}\).
The resulting error is \(O(A_0+\sqrt{A_0B_0})\), and
\[
 \sqrt{A_0B_0}
  =\delta\sqrt{c_*}\,x^{3/4-A/2}
          \sqrt{\frac{(\log(2x))^r}{\log x}{\cal L}_*(x)}.
 \]
This is the second term of (A16.1). Its main term is
asymptotic to
\(\delta c_*x^{1-A}/[(1-A)\log x]\).
The two relative errors are bounded by fixed constants
times \(R+\sqrt R\), where
\(R={\cal L}_*(x)(\log x)^{r+1}x^{A-1/2}\).
This proves (A16.2), and the degree-one argument is
identical. No assertion is made at \(A=1/2\).
\end{proof}

\paragraph{NT-A17: necessity of the monomial independence condition.}
For fixed \(L/K,C,S\), under \(M\) the joint prime angles
in this conjugacy class are Haar equidistributed on
\((\mathbf R/\mathbf Z)^r\) if and only if (A15.1) holds.
More precisely, if
\(\prod_i\eta_i^{m_i}=\theta|\cdot|^{it}\) for nonzero
\(\mathbf m\), with \(\theta\) of finite order dividing \(b\),
rename the real norm parameter so that the ideal value of
the \(b\)-th power is \(q^{i\tau}\). Then its normalized
prime Fourier moment is
\[
 \frac{1}{N_C(x)}
 \sum_{\substack{N_K\mathfrak p\le x,\ \mathfrak p\notin S\\
                  \operatorname{Frob}_{\mathfrak p}\in C}}
          \prod_i\eta_i(\mathfrak p)^{b m_i}
       =\frac{x^{i\tau}}{1+i\tau}+o(1).                  \tag{A17.1}
\]
Here \(N_C(x)\) is the number of primes in the sum.
For \(\tau=0\) the left side is exactly one once \(N_C(x)>0\).
Thus neither separate infinite order nor separate
coordinate equidistribution replaces (A15.1).

\begin{proof}
Sufficiency follows from (A15.5) for every fixed box:
its error is \(o(\delta\operatorname{Li}_2(x))\) for fixed
data, and the zero-frequency count has that denominator
asymptotic. Boxes characterize Haar equidistribution.

For necessity, raise the proposed relation to the \(b\)-th
power. Outside the conductor set, every summand in (A17.1)
is exactly \((N_K\mathfrak p)^{i\tau}\).
The fixed-field Chebotarev estimate under RH for \(\zeta_L\),
which \(M\) supplies, gives
\[
 N_C(u)=\delta\operatorname{Li}_2(u)
                       +O_{L,K,C,S}(\sqrt u\log(2u)).
 \]
Abel summation, with fixed \(\tau\), therefore yields
\[
 \sum_{\mathfrak p\text{ in the sum}}
                  (N_K\mathfrak p)^{i\tau}
    =\delta\int_2^x\frac{u^{i\tau}}{\log u}\,du
                +O_{L,K,C,S,\tau}(\sqrt x\log(2x)).
 \]
The lower endpoint can add only a fixed constant.
Differentiating
\(u^{1+i\tau}/[(1+i\tau)\log u]\) shows that this integral is
\[
       \frac{x^{1+i\tau}}{(1+i\tau)\log x}
                         +O_\tau(x/(\log x)^2).
 \]
Division by \(N_C(x)\sim\delta x/\log x\) proves (A17.1).
Its modulus tends to \(1/\sqrt{1+\tau^2}>0\).
The corresponding nonzero Fourier character has Haar
integral zero, so Haar equidistribution fails.
\end{proof}

\paragraph{Geometric meaning and work placement.}
The character lattice \(\mathbf Z^r\) describes the
compact torus coordinates; (A15.1) removes every
finite character or norm drift in that lattice.
For normalized algebraic Hecke characters this is the
torus part of the tensor category of the corresponding
CM objects, coupled to the finite \'{e}tale \(G\)-cover.
The mixed Hecke family (A15.3) records every tensor
character needed by the proof. This is not statistical
independence of zeros or a new assertion of functoriality
for general automorphic representations.
Full proofs: \path{joint_hecke_tori.tex}, NT-A15--A17,
also \path{section.tex}; run
\path{GRH_NIGHT/number_theory}.
Classical premises are NT-A08/NT-A12, S-NT-NEUKIRCH,
S-NT-DP and the explicit Chebotarev result NT-C02.


\subsection*{NT-E12. Distinct CM fields obstruct simultaneous integral extremes}

Let \(E_1,E_2\) be elliptic curves over a number field \(K\),
with all their complex multiplication defined over \(K\).
Suppose their imaginary quadratic CM fields \(M_1,M_2\)
are distinct. At a prime of common good reduction, put
\[
 q=N_K\mathfrak p,\qquad
 a_i=q+1-\#E_i(\mathcal O_K/\mathfrak p).
 \]
Then the following implication is unconditional:
\[
            a_1^2=a_2^2
       \quad\Longrightarrow\quad
       q\text{ is a square and }a_1^2=a_2^2=4q.
                                                               \tag{E12.1}
\]
In particular, at an absolute-degree-one prime the two
absolute traces are different. Simultaneous integral
extremal events
\[
            |a_1|=|a_2|=\lfloor2\sqrt q\rfloor
                                                               \tag{E12.2}
\]
can occur only at square norms and exact Hasse endpoints.
For all primes of norm at most \(x\) satisfying (E12.2),
the elementary bound
\[
        \#\{\mathfrak p:\text{(E12.2)},\ N_K\mathfrak p\le x\}
                              \le [K:\mathbf Q]\sqrt x
                                                               \tag{E12.3}
\]
holds. No RH assumption is used.

\begin{proof}
We give the centralizer argument, including the rational
Frobenius exception. Let \(E/K\) be either curve, with
CM field \(M\), and let \(\overline E\) be its good
reduction over \(\mathbf F_q\). Good-reduction
specialization injects \(\operatorname{End}_K(E)\) into
\(\operatorname{End}_{\mathbf F_q}(\overline E)\).
Thus \(M\) embeds in
\(A=\operatorname{End}_{\overline{\mathbf F}_q}^0(\overline E)\).
Every element in that embedded field commutes with the
\(q\)-power Frobenius \(\pi\), since its endomorphisms
are defined over \(\mathbf F_q\).

Choose any auxiliary prime \(\ell\ne\operatorname{char}\mathbf F_q\).
The faithful Tate-module action injects
\(A\otimes\mathbf Q_\ell\) into
\(\operatorname{End}_{\mathbf Q_\ell}(V_\ell\overline E)\),
whose underlying space has dimension two.
The action of the two-dimensional \'{e}tale algebra
\(M\otimes\mathbf Q_\ell\) is faithful and has rank one
over that algebra. If the algebra is a field, this
follows from the two dimensions; if it is
\(\mathbf Q_\ell\times\mathbf Q_\ell\), faithfulness
makes each of its two idempotent summands nonzero,
so each has dimension one.
Its centralizer in the two-by-two endomorphism algebra
is therefore exactly \(M\otimes\mathbf Q_\ell\).

Commutation with a basis of \(M\) is a finite system of
rational linear equations in \(A\). The centralizer of
\(M\) in \(A\) consequently has rational dimension at
most two, and already contains \(M\). It equals \(M\).
Since \(\pi\) commutes with \(M\), we conclude \(\pi\in M\).
These specialization and faithful-action facts are
the same classical good-reduction interface used in
the full proof of NT-E10; no assertion about a
Tate-conjecture surjectivity is needed.

The characteristic polynomial on \(V_\ell\) is
\[
                       X^2-aX+q .
 \]
If \(\pi\notin\mathbf Q\), then \(\mathbf Q(\pi)=M\);
its minimal polynomial has degree two and must be the
displayed characteristic polynomial. In particular
\[
                 M=\mathbf Q(\sqrt{a^2-4q}),
                         \qquad a^2-4q<0.               \tag{E12.4}
 \]
If instead \(\pi=t\in\mathbf Q\), its action is scalar.
Its trace and determinant give \(a=2t\) and \(q=t^2\).
As \(q\) is an integer, a rational square root \(t\)
is an integer. Thus \(q\) is a square and \(a^2=4q\).
Conversely a repeated root of the displayed polynomial
forces \(\pi=a/2\): inside the field \(M\), an element
whose square \((\pi-a/2)^2\) is zero is zero.

Now suppose \(a_1^2=a_2^2\). If this common value is
not \(4q\), both Frobenius elements are nonrational,
and (E12.4) identifies both \(M_1\) and \(M_2\) with
the same imaginary quadratic field. This contradicts
their distinctness. Hence the common value is \(4q\),
which proves (E12.1). If \(q=p\) is a rational prime
it is not a square, giving the degree-one conclusion.
More generally every odd absolute residue degree is
excluded.

Event (E12.2) implies equality of squared traces, so
it can occur only as stated. At square norm \(q=p^{2f}\le x\),
the underlying rational prime satisfies \(p\le\sqrt x\).
There are at most \([K:\mathbf Q]\) prime ideals of \(K\)
over each rational prime. Summing this bound over
the rational primes at most \(\sqrt x\) proves
(E12.3), even with the weaker count by all integers.
\end{proof}


\paragraph{Earlier rational special case.}
Bucur--Fit\'e--Kedlaya, Remark~4.11, published p.~1743
(author v2 p.~22), already gives the obstruction
$a_p(E_1)=a_p(E_2)=\lfloor2\sqrt p\rfloor$ over
$\mathbf Q$ for distinct CM fields and $p>3$,
crediting Andrew Sutherland. The formulation here
retains arbitrary base fields, equal absolute traces,
the precise square-norm exception, and the quantitative
count of higher-degree primes. The rational positive-extreme
special case is not claimed as a new literature result.

\paragraph{Meaning of the shrinking-boundary gate.}
Each separate extremal CM trace asks for an angle
of width comparable to \(q^{-1/4}\). Two such
requirements have total shrinking exponent \(1/2\),
exactly the boundary excluded from NT-A16.
Equations (E12.1)--(E12.3) exhibit a genuine integral
Frobenius obstruction there: continuous torus Haar
mass cannot be used to infer degree-one prime
infinitude at this scale. The obstruction concerns
the two integral traces and their quadratic
endomorphism algebras; it is compatible with
equidistribution in every fixed positive-volume box.
It persists after imposing any finite Frobenius
condition, since that only removes primes.
Geometrically the scalar commutant of each CM action
determines the nonrational Frobenius field, so two
distinct CM algebras cannot give the same noncentral
characteristic polynomial.
Full proof: \path{cm_joint_endpoint_obstruction.tex},
NT-E12, also \path{section.tex}; run
\path{GRH_NIGHT/number_theory}.


\subsection*{NT-E13. The endpoint obstruction descends to potential CM curves}

Let \(E_1,E_2\) be elliptic curves over a number field \(k\)
with distinct imaginary quadratic potential-CM fields
\(M_1,M_2\). All endomorphisms need not be defined over \(k\).
Put \(F_i=kM_i\), \(F=F_1F_2\), and let \(S\) contain
the bad-reduction primes of both curves and the finite
primes ramified in \(F/k\). At any prime
\(\mathfrak p\notin S\), write \(q=N_k\mathfrak p\)
and let \(a_i\) be the two Frobenius traces.
Then
\[
       a_1^2=a_2^2\ne0
       \quad\Longrightarrow\quad
       \begin{cases}
       \mathfrak p\text{ splits completely in }F/k,\\
       q\text{ is a square and }a_1^2=4q .
       \end{cases}
                                                               \tag{E13.1}
\]
Consequently the simultaneous integral extreme
\[
              |a_1|=|a_2|=\lfloor2\sqrt q\rfloor
                                                               \tag{E13.2}
\]
is impossible at every absolute-degree-one prime outside
\(S\). Among all prime ideals outside \(S\) of norm at most
\(x\), its count is at most \([k:\mathbf Q]\sqrt x\).
These assertions are unconditional.
The nonzero condition in (E13.1) is essential: both curves
may have trace zero in nonsplit CM components.

\begin{proof}
For a characteristic-zero CM elliptic curve the field
over which all endomorphisms are defined is \(kM_i\).
One way to see the exact quadratic action is to let an
endomorphism act on the one-dimensional tangent space.
In characteristic zero this is a faithful map from the
CM field into \(\overline k\), compatible with Galois
action. The action on endomorphisms therefore factors
through the two automorphisms of \(M_i\), with kernel
precisely the group fixing \(kM_i\).
This is the same endomorphism-field premise used in
NT-E08 and NT-E10.

Fix an auxiliary prime \(\ell\) different from the
residue characteristic. Over an algebraic closure of
\(\mathbf Q_\ell\), the faithful CM action decomposes
the two-dimensional Tate module into its two distinct
one-dimensional eigenspaces. For an arithmetic
Frobenius at \(\mathfrak p\), conjugation of
endomorphisms agrees with its action on \(M_i\).
If the restriction to \(F_i/k\) is nontrivial, Frobenius
interchanges those two eigenspaces. Its matrix in that
decomposition has zero diagonal, hence trace zero.
The Frobenius trace is \(a_i\), since the curve has
good reduction and \(\mathfrak p\) is unramified in
the indicated quadratic extension.

Thus \(a_i\ne0\) forces trivial Frobenius in \(F_i/k\).
If \(a_1^2=a_2^2\ne0\), the Frobenius is trivial on
both \(F_i\), hence on their compositum \(F\).
This proves the splitting assertion in (E13.1).
Choose a prime \(\mathfrak P\) of \(F\) over \(\mathfrak p\).
It has the same residue field, norm \(q\), and both
traces \(a_i\). The two base-changed curves now have
all CM defined over \(F\), with the same distinct
fields \(M_1,M_2\). NT-E12 applies and forces square
norm and exact Hasse equality, proving (E13.1).

Since \(q\ge2\), its extremal integer
\(\lfloor2\sqrt q\rfloor\) is positive. Event (E13.2)
therefore satisfies the nonzero premise of (E13.1).
It cannot occur at norm a rational prime.
For a square norm \(q=p^{2f}\le x\), the underlying
rational prime is at most \(\sqrt x\), and at most
\([k:\mathbf Q]\) prime ideals of \(k\) lie above each
one. This gives the stated count bound.
\end{proof}

\paragraph{An explicit pair over \(\mathbf Q\).}
For
\[
        E_1:y^2=x^3-x,\qquad E_2:y^2=x^3+1,
                                                               \tag{E13.3}
\]
there is no rational prime \(p>3\) at which both
\(|a_p(E_i)|=\lfloor2\sqrt p\rfloor\).
Indeed the automorphism
\((x,y)\mapsto(-x,iy)\) gives CM by \(\mathbf Z[i]\)
on \(E_1\), while
\((x,y)\mapsto(\zeta_3x,y)\) gives CM by
\(\mathbf Z[\zeta_3]\) on \(E_2\).
These are the maximal orders in the distinct
quadratic fields \(\mathbf Q(i)\) and
\(\mathbf Q(\sqrt{-3})\); their displayed generators
are defined over those fields.
The two integral equations have discriminants
\(64\) and \(-432\), respectively, so both curves
have good reduction at every \(p>3\).
The compositum is \(\mathbf Q(\zeta_{12})\), ramified
only at \(2,3\). Every \(p>3\) is therefore within
the hypotheses just proved.


\paragraph{Earlier rational special case.}
Bucur--Fit\'e--Kedlaya, Remark~4.11, published p.~1743
(author v2 p.~22), already gives the obstruction
$a_p(E_1)=a_p(E_2)=\lfloor2\sqrt p\rfloor$ over
$\mathbf Q$ for distinct CM fields and $p>3$,
crediting Andrew Sutherland. The formulation here
retains arbitrary base fields, equal absolute traces,
the precise square-norm exception, and the quantitative
count of higher-degree primes. The rational positive-extreme
special case is not claimed as a new literature result.

\paragraph{Input, scope and geometric placement.}
The proof combines the exact CM component action with
the rational-Frobenius exception of NT-E12. It does
not presume that a nonsplit component has a continuous
circle distribution; its trace is zero instead.
It gives a concrete obstruction over \(\mathbf Q\)
at the total shrinking exponent \(1/2\), and remains
true after any further finite Frobenius restriction.
The relevant geometry is the component group acting
on the two CM eigenlines, followed on the split
component by the two distinct quadratic endomorphism
algebras. No GRH or Artin-holomorphy input occurs.
Full proof: \path{potential_cm_joint_endpoint.tex},
NT-E13; run \path{GRH_NIGHT/number_theory}.

\subsection*{NT-A18--A19. Reciprocal-prime weighting removes the norm-drift obstruction}

Keep the finite Galois extension \(L/K\), conjugacy class
\(C\), cyclic subgroup \(H=\langle g\rangle\), field
\(F=L^H\), and \(\delta=|C|/|G|\) of NT-A12.
Fix unitary Hecke characters \(\eta_1,\ldots,\eta_r\)
of \(K\), and let \(S\) contain all ramification and
conductor primes. Write
\(\nu_i=\eta_i\circ N_{F/K}\).
For this result the independence condition is only
\[
 \prod_i\eta_i^{m_i}\text{ has infinite order}
              \qquad(0\ne\mathbf m\in\mathbf Z^r).         \tag{A18.1}
\]
Nontrivial unitary norm powers are allowed.
The sufficient analytic input is RH for \(\zeta_L\)
and for every primitive, non-pure-norm member of
\[
 L_F\left(s,\chi\prod_i\nu_i^{m_i}\right),
      \quad \chi\in\widehat H,\quad
                         0\ne\mathbf m\in\mathbf Z^r .
                                                               \tag{A18.2}
\]
Pure norm members require only RH for \(\zeta_F\),
already supplied by RH for \(\zeta_L\). Thus \(M\)
supplies exactly the stated rank-one input, without
general Artin holomorphy. All data in this subsection
are fixed; no uniformity as frequencies approach a
norm relation is asserted.

\paragraph{NT-A18: finite harmonic constants and the exact weighted torus gate.}
Let \(\mathcal P_C\) be the primes outside \(S\) with
Frobenius class \(C\). For \(q_{\mathfrak p}=N_K\mathfrak p\),
put
\[
 W_C(x)=\sum_{\substack{\mathfrak p\in\mathcal P_C\\
                                      q_{\mathfrak p}\le x}}
                       q_{\mathfrak p}^{-1},\qquad
 W_{\mathbf m}(x)=
 \sum_{\substack{\mathfrak p\in\mathcal P_C\\
                                      q_{\mathfrak p}\le x}}
       \frac{\prod_i\eta_i(\mathfrak p)^{m_i}}{q_{\mathfrak p}}.
 \]
There are finite constants \(B_C\in\mathbf R\) and
\(B_{\mathbf m}\in\mathbf C\), for every nonzero
\(\mathbf m\), such that
\[
 \begin{aligned}
 W_C(x)&=\delta\log\log x+B_C
                         +O_{L,K,S}\!\left(\frac{\log(2x)}{\sqrt x}\right),\\
 W_{\mathbf m}(x)&=B_{\mathbf m}
          +O_{\mathbf m,\boldsymbol\eta,L,K,S}\!
              \left(\frac1{\log x}+\frac{\log(2x)}{\sqrt x}\right).
 \end{aligned}                                             \tag{A18.3}
\]
If none of the factors for this \(\mathbf m\) in
(A18.2) is pure norm, its \(1/\log x\) term can be
omitted. The series defining \(B_{\mathbf m}\) is
ordered by prime norm; absolute convergence of its
first powers is not claimed.

The probability measures
\[
 \frac1{W_C(x)}
   \sum_{\substack{\mathfrak p\in\mathcal P_C\\q_{\mathfrak p}\le x}}
        \frac{\delta_{(\eta_1(\mathfrak p),\ldots,\eta_r(\mathfrak p))}}
             {q_{\mathfrak p}}
                                                               \tag{A18.4}
\]
therefore tend weakly to Haar measure on \((S^1)^r\).
Every fixed product of arcs has limiting mass equal
to its product length, with arbitrary endpoints.
For each fixed nonzero Fourier frequency the error
is \(O_{\mathbf m}(1/\log\log x)\).
The same conclusions, with possibly changed finite
constants, hold if only absolute-degree-one primes
are counted.

Conversely a finite-order nonzero monomial prevents
full Haar convergence for these weights. Thus, under
\(M\), (A18.1) is the exact condition for full
reciprocal-prime-weighted Haar distribution. It is
strictly weaker than the ordinary-prime condition
in NT-A17.

There is also a convergent Euler-product formulation.
For every nonzero \(\mathbf m\), the products
\[
 \prod_{\substack{\mathfrak p\in\mathcal P_C\\q_{\mathfrak p}\le x}}
 \left(1-\frac{z\prod_i\eta_i(\mathfrak p)^{m_i}}
                    {q_{\mathfrak p}}\right)^{-1}
                                                               \tag{A18.5}
\]
converge locally uniformly for \(|z|<2\) to a
nonvanishing holomorphic function. In particular
their values at \(z=1\) have finite nonzero limits.

\begin{proof}
First consider a single primitive unitary Hecke
character \(\lambda\) of \(F\), with all conductor
primes removed. Additional finitely many primes
may also be removed. If \(\lambda\) is not pure norm,
NT-A08 under its RH gives, for fixed data,
\[
 A_\lambda(u):=\sum_{N_F\mathfrak q\le u}\lambda(\mathfrak q)
                        =O_\lambda(\sqrt u\log(2u)).
 \]
Abel summation, including any primes of norm two,
gives
\[
 \sum_{N_F\mathfrak q\le x}
            \frac{\lambda(\mathfrak q)}{N_F\mathfrak q}
   =\frac{A_\lambda(x)}x+
                         \int_2^x\frac{A_\lambda(u)}{u^2}\,du .
 \]
The integral over \([2,\infty)\) is absolutely
convergent. Its tail and the endpoint term are
\(O_\lambda(\log(2x)/\sqrt x)\). This proves the
claimed finite limit in the non-pure case.

Now let the ideal values of a pure norm character
be \((N_F\mathfrak q)^{i\tau}\), with \(\tau\ne0\).
Under RH for \(\zeta_F\), the prime count after the
same finite deletions is
\(\operatorname{Li}_2(u)+E(u)\), where
\(E(u)=O_{F,S}(\sqrt u\log(2u))\).
Stieltjes integration by parts with
\(u^{i\tau-1}\) gives
\[
 \begin{split}
 \sum_{N_F\mathfrak q\le x}(N_F\mathfrak q)^{i\tau-1}
  ={}&\int_2^x\frac{u^{i\tau-1}}{\log u}\,du
       +x^{i\tau-1}E(x)\\
    &-(i\tau-1)\int_2^x E(u)u^{i\tau-2}\,du .
 \end{split}                                                   \tag{A18.6}
\]
The \(E\)-integral converges absolutely, and its
tail together with the endpoint is
\[
        O_{F,S}\left((1+|\tau|)
                               \frac{\log(2x)}{\sqrt x}\right).
 \]
The main integral becomes
\(\int_{\log2}^{\log x}e^{i\tau y}\,dy/y\).
It converges at infinity by integration by parts,
with the explicit tail bound
\[
       \left|\int_{\log x}^\infty
                      \frac{e^{i\tau y}}y\,dy\right|
                         \le\frac2{|\tau|\log x}.         \tag{A18.7}
\]
These computations prove actual convergence to a
finite constant in the pure norm case too.
There is no application of the shifted-pole
realification argument in this case.

Condition (A18.1) survives norm pullback:
if a pulled-back monomial had finite order \(b\),
composition with the diagonal idele map would
make its original monomial have order dividing
\(b[F:K]\). Thus, for nonzero \(\mathbf m\),
\(\chi\prod_i\nu_i^{m_i}\) is never principal.
If it is pure norm, its parameter \(\tau\) is
therefore nonzero. RH for \(\zeta_L\) supplies
RH for \(\zeta_F\) by the cyclic factorization
\(\zeta_L=\prod_{\chi\in\widehat H}L_F(s,\chi)\):
the nontrivial finite characters are entire, so
a nontrivial zero of any factor cannot be
cancelled in the open critical strip.

Apply the cyclic class-function identity (A12.9)
to the weighted prime sum. The part over primes
of \(F\) with relative residue degree one is
exactly \(W_{\mathbf m}(x)\), with coefficient
\(\delta\overline{\chi(g)}\) in the \(\chi\)-sum.
Each full prime sum over \(F\) has a finite limit
by (A18.6)--(A18.7) or the preceding non-pure
argument. The additional primes of relative
residue degree at least two give an absolutely
convergent correction, uniformly in their unit
character values.

Here is a quantitative check of that correction.
Every such prime has absolute residue degree
at least two. For \(n_F=[F:\mathbf Q]\),
\[
 \begin{split}
 \sum_{\substack{\mathfrak q:\ f_{\mathfrak q}\ge2\\
                                   N_F\mathfrak q>x}}
                \frac1{N_F\mathfrak q}
 &\le n_F\sum_{j=2}^{n_F}\sum_{p>x^{1/j}}p^{-j}\\
 &\le n_F\sum_{j=2}^{n_F}
       \left(x^{-1}+\frac{x^{-1+1/j}}{j-1}\right)
                       \ll_F x^{-1/2}.
 \end{split}                                                   \tag{A18.8}
\]
The inner bound follows by comparison with the
integral of \(t^{-j}\), and the rational primes
may be enlarged to all integers. The empty
sum for \(n_F=1\) causes no exception.
Combining these finitely many \(\chi\)-sums and
corrections proves the second line of (A18.3).

For the zeroth frequency, NT-C02 and finite
deletions give
\[
 \pi_C(u)=\delta\operatorname{Li}_2(u)+R_C(u),
          \qquad R_C(u)=O_{L,K,S}(\sqrt u\log(2u)).
 \]
The complete reciprocal-sum calculation of
NT-C08 applies verbatim: define
\[
 B_C=-\delta\log\log2+
                             \int_2^\infty\frac{R_C(u)}{u^2}\,du .
 \]
Then the error is exactly
\(R_C(x)/x-\int_x^\infty R_C(u)du/u^2\).
This proves the first line of (A18.3), and
\(W_C(x)\to\infty\).
Deleting absolute-degree-at-least-two primes
of \(K\) changes each Fourier sum by a convergent
series with the tail (A18.8), now over \(K\).
This proves every assertion about the degree-one
version.

Divide each nonzero Fourier sum by \(W_C(x)\).
Its limit is zero, with the stated fixed-frequency
rate. For completeness, trigonometric polynomials
are uniformly dense in continuous functions on
the torus, for instance by product Fej\'er means.
For a polynomial the limiting integral is Haar
integration by the just-proved Fourier limits.
For a continuous function approximate uniformly;
both the measures in (A18.4) and Haar measure
have mass one, so the approximation error is
at most twice the uniform error. This proves
weak convergence. Continuous upper and lower
approximations to a fixed box, with arbitrarily
small Haar integral gap at its faces, prove the
box assertion and every endpoint convention.

If instead a nonzero monomial is finite order
dividing \(b\), its \(b\)-th power is identically
one outside the conductor set. The corresponding
nonzero Fourier coefficient of (A18.4) is exactly
one whenever the denominator is positive.
Its Haar integral is zero, which proves necessity.

Finally set \(u_{\mathfrak p}=\prod_i\eta_i(\mathfrak p)^{m_i}\).
For \(|z|<2\), use the power-series logarithm of
each factor in (A18.5). Its first-degree sum is
\(zW_{\mathbf m}(x)\), converging locally uniformly.
For each fixed \(R<2\), all higher-degree sums
are uniformly absolutely convergent on \(|z|\le R\),
since
\[
 \sum_{\mathfrak p\in\mathcal P_C}\sum_{j\ge2}
       \frac{R^j}{j q_{\mathfrak p}^{\,j}}
 \le \frac{R^2}{2(1-R/2)}
                 \sum_{\mathfrak p\in\mathcal P_C}q_{\mathfrak p}^{-2}
                                                    <\infty .
 \]
The resulting logarithm is holomorphic. Its
exponential is the asserted locally uniform
Euler-product limit and never vanishes.
\end{proof}

\paragraph{NT-A19: a Dirichlet-GRH-only progression example.}
Fix primitive Dirichlet characters \(\chi_i\) of
conductors \(f_i\), and real numbers \(t_1,\ldots,t_r\)
which are linearly independent over \(\mathbf Q\).
For primes \(p\nmid\prod_i f_i\), put
\[
                    z_i(p)=\chi_i(p)p^{it_i}.
 \]
Fix a reduced arithmetic progression \(p\equiv a\pmod q\),
and put \(Q_0=\operatorname{lcm}(q,f_1,\ldots,f_r)\).
RH for the primitive Dirichlet characters with
conductor dividing \(Q_0\), including \(\zeta\) for
conductor one, suffices for the following assertions.
This is a finite subfamily of the user-adopted input
\(A_D\), and uses no additional \(M\) input.

The \(1/p\)-weighted points \((z_1(p),\ldots,z_r(p))\)
in that progression tend to full Haar measure on
\((S^1)^r\). Every nonzero mixed prime harmonic sum
converges to a finite constant. In contrast, the
ordinary equal-weight prime points fail to be
Haar equidistributed.

\begin{proof}
Partition the progression into its finitely many
reduced residue classes \(b\bmod Q_0\) satisfying
\(b\equiv a\bmod q\), omitting the finite conductor
primes. On each such class, a Fourier frequency
\(\mathbf m\ne0\) has value
\[
       \prod_i z_i(p)^{m_i}
        =\left(\prod_i\chi_i(b)^{m_i}\right)p^{i\tau_{\mathbf m}},
       \qquad \tau_{\mathbf m}=\sum_i m_i t_i\ne0.
 \]
The progression estimate NT-P12, with the
moduli-one-and-two convention there, gives
\[
 \pi(u;Q_0,b)
   =\frac{\operatorname{Li}_2(u)}{\varphi(Q_0)}
                              +O_{Q_0}(\sqrt u\log(2u)).
 \]
Apply exactly the Stieltjes calculation (A18.6),
with the main integral multiplied by
\(1/\varphi(Q_0)\). Equation (A18.7) proves
convergence of each Fourier harmonic sum on
each class, hence on their finite union.
The zeroth sum is
\(\varphi(q)^{-1}\log\log x+O_{q,\boldsymbol f}(1)\).
The same Fourier approximation argument therefore
proves full weighted Haar distribution.

For the ordinary weighting, choose any coordinate
\(i\) and an integer \(b_i\ge1\) with
\(\chi_i^{b_i}=1\). Its nonzero Fourier frequency
\(b_i\mathbf e_i\) has value \(p^{ib_i t_i}\).
The ordinary Abel calculation in NT-A17, now
using NT-P12 in the fixed progression, gives
\[
 \frac{\sum_{p\le x,\ p\equiv a\bmod q}p^{ib_i t_i}}
               {\pi(x;q,a)}
       =\frac{x^{ib_i t_i}}{1+ib_i t_i}+o(1).
 \]
Omitting conductor primes does not affect the
limit. Its modulus tends to the positive number
\((1+b_i^2t_i^2)^{-1/2}\), whereas the Haar
Fourier integral is zero. This proves the
ordinary-weight obstruction.
\end{proof}

\paragraph{Geometric and analytic meaning.}
The measures are reciprocal-norm measures on
closed points in a specified finite \'{e}tale
Frobenius class, with values in a compact torus
of rank-one Hecke coordinates. Their radial
norm drift is averaged by the logarithmic
measure \(dy/y\); a finite character relation
cannot be averaged away.
For general unitary infinity types this is an
idele-class formulation, not a claimed algebraic
motive realization. No zero-independence or
joint-zero-statistics hypothesis is used.
The constants in (A18.7) can grow when a nonzero
frequency approaches a norm relation, so no
uniform shrinking-box result is inferred here.
Full proofs: \path{reciprocal_prime_tori.tex},
NT-A18--A19; run \path{GRH_NIGHT/number_theory}.
Exact classical inputs are NT-A08 (S-NT-IK),
NT-C02 (S-NT-GM), the cyclic induction already
proved in NT-A12, and NT-P12 for the \(A_D\)
example. These are newly proved formulations
in this run, with no literature-priority claim.


\subsection*{NT-A20--A21. Euclidean ideal classes and one-point localization}

\paragraph{NT-A20: the exact Picard-group criterion.}
Let $K$ be a number field, let $S$ be a finite set of places
containing every archimedean place, and put $R=\mathcal O_{K,S}$.
Assume $|S|\ge2$. For every positive squarefree integer $m$, set
\[
 E_m=K\bigl(\zeta_m,(R^\times)^{1/m}\bigr).
 \tag{A20.1}
\]
These are finite number fields: the finitely generated group
$R^\times$ has only finitely many generators whose $m$th roots
need to be adjoined. Assume precisely RH for all $\zeta_{E_m}$.
Input $M$ supplies this family through its trivial Hecke characters.

A fractional ideal class $C\in\operatorname{Pic}(R)$ is
Euclidean in Lenstra's sense if and only if
\[
                 \operatorname{Pic}(R)=\langle C\rangle.
 \tag{A20.2}
\]
Thus, if $\operatorname{Pic}(R)$ is cyclic of order $h_S$,
it has exactly $\varphi(h_S)$ Euclidean ideal classes,
where $\varphi(1)=1$; otherwise it has none.
At most one of these classes can be Euclidean for the ideal norm.
In particular, when $h_S>2$ is cyclic, at least
$\varphi(h_S)-1$ classes admit a Euclidean reduction but
do not admit the norm reduction.

Here is the precise definition and the full elementary part
of the argument. Write
\[
 \mathcal E(R)=\{I:\ I\text{ is a nonzero fractional ideal of }R,\
                                      R\subseteq I\}.
\]
For a representative $\mathfrak c$ of $C$, a Euclidean
reduction is a map $\psi:\mathcal E(R)\to W$ to a well-ordered
set such that for every $I\in\mathcal E(R)$ and
$x\in I\mathfrak c\setminus\mathfrak c$, there is
$y\in\mathfrak c$ satisfying
\[
 J=(x-y)^{-1}I\mathfrak c\in\mathcal E(R),
                    \qquad \psi(J)<\psi(I).
 \tag{A20.3}
\]
The element $x-y$ is nonzero. The displayed membership
holds because $x-y\in I\mathfrak c$.
Replacing $\mathfrak c$ by a principal multiple leaves
the definition unchanged.

\begin{proof}
For necessity, suppose there is an ideal class outside
$\langle C\rangle$. Every fractional ideal class has a
representative in $\mathcal E(R)$: for a nonzero
$a$ in a fractional ideal $\mathfrak a$, the ideal
$a^{-1}\mathfrak a$ contains $R$.
Choose $I\in\mathcal E(R)$ outside that subgroup with
$\psi(I)$ smallest. Then $I\ne R$, and invertibility
of $\mathfrak c$ gives $I\mathfrak c\ne\mathfrak c$.
Choose $x$ in their difference and apply (A20.3).
Since $[J]=[I]C$, the class $[J]$ is still outside
$\langle C\rangle$, contradicting the minimal choice.
This proves necessity without RH.

For sufficiency, apply Lenstra's Proposition~(2.2),
printed p.~127 of \emph{Euclidean ideal classes},
Ast\'erisque~61 (1979), 121--131. Its hypotheses are
exactly $|S|\ge2$ and RH for (A20.1), and its conclusion
is that every class generating $\operatorname{Pic}(R)$
has such a reduction. The printed inequality is
$|S|\ge2$; the searchable OCR incorrectly renders it
as a strict inequality. This is a cited classical
analytic input. The paper refers its sufficiency
argument to the methods of Lenstra~1977, reference~[9];
it does not print a full new analytic proof at p.~127.
We do not represent the statement check as a fresh
verification of that omitted analytic argument.

The number of generators of a finite cyclic group of
order $h_S$ is $\varphi(h_S)$, giving the count.
For the norm assertion, use the single fixed map
\[
                       \psi_N(I)=N_R(I)^{-1}\in\mathbf N.
 \tag{A20.4}
\]
It is integral because $I^{-1}$ is integral and
$N_R(I)^{-1}=\#(R/I^{-1})$.
If two classes were Euclidean for this same map,
choose an $I\ne R$ of smallest $\psi_N(I)$.
For either class, (A20.3) produces a $J$ with smaller
value. Minimality among nontrivial overideals forces
$J=R$, so both classes must equal $[I]^{-1}$.
The set from which $I$ is chosen is nonempty:
there is a prime ideal outside the finite set $S$.
Finally,
\[
 \psi_N\bigl((x-y)^{-1}I\mathfrak c\bigr)<\psi_N(I)
                  \quad\Longleftrightarrow\quad
                     N_R(x-y)<N_R(\mathfrak c),
\]
which is exactly norm-Euclidean reduction. This proves
the claimed uniqueness and the lower count of classes
that are not norm-Euclidean.
\end{proof}

The necessity argument and fixed-function uniqueness
are elementary reformulations of Lenstra's
Lemma~(1.5), Theorem~(1.6), and the paragraph after
Exercise~(1.7), printed pp.~124--125; those complete
elementary proofs were read. The mathematical content
of the sufficiency is classical. The explicit family,
class count, and comparison with $M$ here make its
input and consequences reviewable. No finite-order
Dirichlet input alone supplies the general unit
Kummer family. No complexity bound for evaluating
the well-ordered reduction is claimed.

\paragraph{NT-A21: a bounded single closed point makes a Euclidean ring.}
Let $K$ be any number field whose ordinary class group
is cyclic, of order $h$. Write $n=[K:\mathbf Q]$ and
$D=|\operatorname{disc}K|$. Under $M$, for every
generator $C$ of $\operatorname{Cl}(K)$ there is a
prime ideal $\mathfrak p$ of absolute residue degree
one such that
\[
 [\mathfrak p]=C,\qquad
 N\mathfrak p\le
                \bigl[h(4\log D+2.5n)+5\bigr]^2,
 \tag{A21.1}
\]
and the ring
\[
             R=\mathcal O_{K,S_\infty\cup\{\mathfrak p\}}
 \tag{A21.2}
\]
is Euclidean. In particular, this holds for fields
whose original unit group is finite. The norm bound
controls the one closed point removed; it is not a
bound for the running time of the Euclidean reduction.

The exact analytic inputs, weaker than universal $M$,
are RH for $\zeta_{\mathbf Q}$ and for the ordinary
Hilbert class-field zeta function $\zeta_{H_K}$,
followed by RH for the fields (A20.1) formed from
the particular ring (A21.2). The latter family
depends on the selected prime. No claim that the
first two zeta functions imply that family is made.

\begin{proof}
NT-C07, the checked specialization of Bach--Sorenson
Theorem~5.1, gives (A21.1). Its discriminant
substitution is $D_{H_K}=D^h$ and $[H_K:\mathbf Q]=hn$.
For $K=\mathbf Q$ take $\mathfrak p=(2)$ directly;
it satisfies (A21.1).

Extension of ideals to $R$ gives the exact quotient
\[
 \operatorname{Pic}(R)=
       \operatorname{Cl}(K)/\langle[\mathfrak p]\rangle.
 \tag{A21.3}
\]
Indeed every ideal of $R$ is represented by a finite
product of primes different from $\mathfrak p$,
so extension is onto. If an ideal becomes principal,
its divisor differs from a principal divisor only at
$\mathfrak p$, showing that the kernel is the displayed
cyclic subgroup. Conversely every power of
$\mathfrak p$ becomes the unit ideal.
Since $[\mathfrak p]$ generates $\operatorname{Cl}(K)$,
the group in (A21.3) is zero.

There are at least two places in
$S_\infty\cup\{\mathfrak p\}$. Apply NT-A20 to its
principal ideal class; equivalently apply the
classical PID statement NT-A03. This makes $R$
Euclidean. For clarity, the equivalence between a
principal Euclidean ideal class and a Euclidean domain
can also be seen directly. For nonzero $b\in R$ put
$\phi(b)=\psi(b^{-1}R)$, with $\psi$ from (A20.3) and
$\mathfrak c=R$. If $a/b\notin R$, apply (A20.3) to
$I=b^{-1}R$ and $x=a/b$ to find $y\in R$. Then
$r=a-by\ne0$ and
\[
       (a/b-y)^{-1}b^{-1}R=r^{-1}R,\qquad
                                  \phi(r)<\phi(b).
 \tag{A21.4}
\]
If $a/b\in R$, the remainder is zero, so the usual
division property holds in all cases.

Finally $\mathfrak p^h=(a)$ for some $a\in\mathcal O_K$,
because $[\mathfrak p]$ has order $h$. Its divisor
is supported only at $\mathfrak p$, and hence
$R=\mathcal O_K[1/a]$. Thus (A21.2) is the actual
affine open obtained by deleting this single closed
point, even when $\mathfrak p$ is not principal.
\end{proof}

Geometrically, NT-A20 concerns reduction in an
invertible sheaf generating
$\operatorname{Pic}(\operatorname{Spec}R)$.
NT-A21 removes one quantitatively bounded closed point
whose divisor class kills that Picard group, and then
uses the unit-Kummer RH family to supply Euclidean
reduction. Cyclicity remains a hypothesis; the
argument neither asserts that all class groups are
cyclic nor that the ideal norm is a Euclidean function.


\subsection*{NT-E14--E16. Finite-field detection of characteristic-zero isogeny}

The analytic input here is Dedekind RH for certain finite
number fields. It is not RH for the standard $L$-function
of a non-CM elliptic curve.
For an even squarefree integer $m\ge2$, let
$\mathcal D_m$ be the collection of finite Galois
extensions $T/\mathbf Q$ satisfying
\[
 [T:\mathbf Q]\le384,\qquad
 T/\mathbf Q\text{ is unramified at every prime }p\nmid m.
 \tag{E14.1}
\]
The hypothesis RH for $\zeta_{\mathbf Q}$ and for
$\zeta_T$, $T\in\mathcal D_m$, is sufficient for the
applications below. This is a finite family: the
discriminant bound (E16.4) below and Hermite's theorem
give finiteness. A precise reference is Milne,
\emph{Algebraic Number Theory}, v3.08 (2020),
Theorems~8.42--8.43, printed pp.~151--152.
Input $M$ supplies this family.
The more restricted fields actually used in the
argument are specified explicitly below.

\paragraph{NT-E14: the published non-CM isogeny test.}
Let $E_1,E_2/\mathbf Q$ be elliptic curves without
complex multiplication, with conductors $N_1,N_2$,
and put $m=\operatorname{rad}(2N_1N_2)$.
If the two curves are not $\mathbf Q$-isogenous,
then there is a prime $p\nmid2N_1N_2$ with
\[
 a_p(E_1)\ne a_p(E_2),\qquad
 p\le \bigl(482\log m+2880\bigr)^2.
 \tag{E14.2}
\]
Consequently equality of the point counts at all
such primes up to this bound is equivalent to
$\mathbf Q$-isogeny.
This is Mayle--Wang, \emph{On the effective version
of Serre's open image theorem},
Bull.\ London Math.\ Soc.~56 (2024), 1399--1416,
Theorem~2, printed p.~1401; complete proof
pp.~1408--1410. Its proof omits $2$ as well as all
bad primes, although the theorem statement only
requires common good reduction.

The precise source dependency is useful. Proposition~12,
pp.~1407--1408, constructs a finite Galois quotient field
\[
 F\subset\mathbf Q(E_1[2^\infty],E_2[2^\infty]),
                      \qquad [F:\mathbf Q]\le192.
 \tag{E14.3}
\]
It detects a trace difference. Corollary~6,
pp.~1404--1405, then applies effective Chebotarev to
$T=F(\sqrt m)$ to exclude every prime dividing $m$.
Thus only RH for this actual $\zeta_T$ and
$\zeta_{\mathbf Q}$ is used analytically.
The field $F$ is a quotient of the full $2$-adic image;
it is not asserted to be the first $2$-division field.
The convenient fixed family (E14.1) contains every
possible such $T$. The improved numerical constants
in (E14.2) use the published Bach--Sorenson Table~3
bounds through the source's computation
(13)--(16). Those numerical inputs are credited,
not represented as a fresh rerun of that table.

\paragraph{NT-E15: the published effective open-image bound.}
For a non-CM elliptic curve $E/\mathbf Q$ with
conductor $N_E$, input $M$ implies that
$\overline\rho_{E,\ell}$ is surjective for every
prime $\ell>B_E$, where
\[
                 B_E=964\log\operatorname{rad}(2N_E)+5760.
 \tag{E15.1}
\]
This is Mayle--Wang Theorem~1, printed p.~1401,
proved in \S5.2, pp.~1412--1413.
It retains both the base field $\mathbf Q$ and
the non-CM hypothesis. It does not prove the
uniform proposed cutoff $37$ for all curves.

Here a sufficient analytic family is (E14.1) for
$m=\operatorname{rad}(2N_E)$, together with
RH for $\zeta_{\mathbf Q}$. Indeed the proof
associates to a nonsurjective prime $\ell>37$ a
nontrivial quadratic twist $E_D$ with
$\operatorname{rad}(N_{E_D})\mid m$.
It applies (E14.2) to $E,E_D$; the two curves
are not $\mathbf Q$-isogenous by its Lemma~14.
The witnessing prime has a nonzero trace divisible
by $\ell$, and the Hasse bound gives
$\ell\le2\sqrt p$, producing (E15.1).
All analytic fields are therefore the fields
$F_D(\sqrt m)$ from (E14.3) for these twists,
of degree at most $384$ and unramified outside $m$.
The finitely many relevant squarefree $D$ divide
$N_E$, with their signs allowed.
The source's modular-curve exclusions, local
inertia lemma, Faltings theorem, and twist facts
are unconditional classical inputs; this reduction
does not claim to reprove them.

\paragraph{NT-E16: a complete decision formulation from integral equations.}
Let $E_1,E_2$ be given by nonsingular integral
Weierstrass equations over $\mathbf Q$, of
discriminants $\Delta_1,\Delta_2\ne0$.
No minimality, factorization, conductor, or
non-CM promise is required in this paragraph.
Put
\[
 R=2|\Delta_1\Delta_2|,\qquad
 b=\operatorname{bitlength}(R),\qquad
                      B=(1532b+10181)^2.
 \tag{E16.1}
\]
Under $M$, the curves are $\mathbf Q$-isogenous
if and only if
\[
 \#E_1(\mathbf F_p)=\#E_2(\mathbf F_p)
                  \quad\hbox{for every prime }p\le B,\ p\nmid R.
 \tag{E16.2}
\]
Testing (E16.2) by direct point counting is a
deterministic algorithm polynomial in the total
bit length of the integral coefficients.
This decides isogeny over $\mathbf Q$; it does
not construct an isogeny, decide geometric
isogeny over arbitrary extensions, or compute an
endomorphism ring.

More precisely, the sufficient RH family is
RH for $\zeta_{\mathbf Q}$ and
all $\zeta_T$ in (E14.1) with
$m=\operatorname{rad}(R)$.
The algorithm never needs to compute $m$.
For any nonisogenous pair only one actual
quotient field $F$ and the field $F(\sqrt m)$
are needed in the proof.

\begin{proof}
We give the finite-quotient argument explicitly;
it is the $\ell=2$, rank-two case of the
Faltings--Serre construction in Mayle--Wang
Proposition~12.
Faltings's isogeny theorem says that two abelian
varieties are isogenous over their number field
if and only if their local Euler factors agree
at all but finitely many places: Invent.\ Math.~73
(1983), Corollary~2, equivalence (i)--(iii),
printed p.~361. For elliptic curves over
$\mathbf Q$, the good local factor is determined
by $1-a_pT+pT^2$. If $E_1,E_2$ are not
$\mathbf Q$-isogenous, we can therefore choose
a common good odd prime where their traces differ.
A Frobenius element at that prime has differing
rational $2$-adic Tate traces. This argument also
applies to CM curves and needs no separate
character-criterion citation. Mayle--Wang recall
the elliptic criterion in (11), printed p.~1407.

Write $\rho=(\rho_1,\rho_2)$ for the product
of the integral rank-two Tate representations.
Let $A$ be the $\mathbf Z_2$-linear span of
$\rho(G_{\mathbf Q})$ in
$M_2(\mathbf Z_2)\times M_2(\mathbf Z_2)$.
Products of image elements remain image elements,
so $A$ is a ring containing the identity.
It is a finite free $\mathbf Z_2$-module of rank
at most eight. Let $e$ be the least finite
$2$-adic valuation of a nonzero difference
$\operatorname{tr}\rho_1(\sigma)-
 \operatorname{tr}\rho_2(\sigma)$.
The linear functional
\[
 \lambda:A\longrightarrow\mathbf Z_2,\qquad
 (u,v)\longmapsto
             2^{-e}(\operatorname{tr}u-\operatorname{tr}v)
 \tag{E16.3}
\]
is integral on $A$ and induces a nonzero functional
on at least one image element in $A/2A$.
Let $G$ be the image of $G_{\mathbf Q}$ in
$(A/2A)^\times$, and let $C$ be the nonempty
subset on which this functional is nonzero.
Trace invariance makes $C$ conjugacy-stable.
Every lift of an element of $C$ has different
two Tate traces.

Here $|G|\le2^8-1=255$: zero is not a unit.
The homomorphism to $G$ is continuous.
Indeed $A$ is a closed finite lattice in the
ambient matrix module, and $2A$ is open in $A$.
The fixed field $F$ of its kernel is consequently
a finite Galois extension of $\mathbf Q$,
contained in the joint $2$-adic division tower.
It lies in some finite level
$\mathbf Q(E_1[2^j],E_2[2^j])$, with $j\ge1$, so its degree
divides
$|\mathrm{GL}_2(\mathbf Z/2^j\mathbf Z)|^2
 =36\cdot2^{8j-8}$.
Among divisors of such an integer not exceeding
$255$, the largest is $192$.
Thus $[F:\mathbf Q]\le192$.
This argument uses $A/2A$ itself; it does not
identify it with the reduction of the ambient
matrix algebra, an identification that need
not be injective for an unsaturated lattice.

N\'eron--Ogg--Shafarevich gives that $F/\mathbf Q$
is unramified outside the prime divisors of $R$.
Here the exact cited form is Mayle--Wang
Theorem~8 and Corollary~9, printed pp.~1406--1407,
with their Remark~11 for a product of two curves;
Theorem~8 cites Serre--Tate, \emph{Good reduction
of abelian varieties}, Ann.\ Math.~88 (1968),
Theorem~1.
Set $m=\operatorname{rad}(R)$ and
$T=F(\sqrt m)$. This is Galois of degree
$d\le384$ and unramified outside $m$.
As $m$ is even and squarefree, every prime
dividing $m$ ramifies in $\mathbf Q(\sqrt m)$,
and hence in $T/\mathbf Q$.
The inverse image of $C$ in
$\operatorname{Gal}(T/\mathbf Q)$ is nonempty
and conjugacy-stable.
NT-C03 applied to any one of its conjugacy
classes provides, under RH for $\zeta_T$
and $\zeta_{\mathbf Q}$, an unramified prime
$p$ with differing Tate traces and
\[
               p\le(4\log D_T+2.5d+5)^2.
\]
In particular $p\nmid R$, so both input equations
have good reduction and the differing traces
are their actual $a_p(E_i)$.

The standard different bound, in the exact
form of Mayle--Wang Lemma~7, printed p.~1405,
gives
\[
 \begin{split}
 \log D_T
   &\le(d-1)\log\operatorname{rad}(D_T)+d\log d\\
   &\le383\log m+384\log384.
 \end{split}
 \tag{E16.4}
\]
For completeness the local calculation is
$v_{\mathfrak q}(\mathfrak D_T)
 \le e_{\mathfrak q}-1+
             e_{\mathfrak q}v_p(e_{\mathfrak q})$.
For a Galois field, $e_{\mathfrak q}\mid d$;
summing residue degrees over primes above $p$
bounds the exponent of $p$ in $D_T$ by
$d-1+d\,v_p(d)$.
Sum after multiplying by $\log p$ over
ramified $p$. The first terms contribute
$(d-1)\log\operatorname{rad}(D_T)$ and the
remaining terms are at most $d\log d$.
This proves (E16.4), retaining the classical
local different inequality as its input.

Since $\log384<6$, the Chebotarev expression is
bounded by
\[
 \begin{split}
 4\log D_T+2.5d+5
    &\le1532\log m+1536\log384+965\\
    &<1532\log R+10181
     \le1532b+10181.
 \end{split}
 \tag{E16.5}
\]
The last inequality uses $R<2^b$ and $\log2<1$.
This proves the existence of a differing point
count within (E16.1) whenever the two curves
are not $\mathbf Q$-isogenous.
Conversely a $\mathbf Q$-isogeny identifies
the rational Tate representations, so the
point counts agree at every common good
prime. This proves the equivalence.

Finally let $L$ be the total input bit length.
Each discriminant is a fixed polynomial in
the five integral Weierstrass coefficients;
it is computable in polynomial time and
has bit length $O(L)$. Thus $b=O(L)$ and
$B=O(L^2)$. Trial division or a sieve
enumerates all primes up to the integer $B$
in polynomial time. Testing $p\nmid R$ only
needs a remainder calculation.
At each retained prime enumerate all
$(x,y)\in\mathbf F_p^2$ satisfying the reduced
Weierstrass equation and add the unique point
at infinity. This uses $O(p^2)$ elementary
modular operations; even the sum over all
integers $p\le B$ is $O(B^3)$.
The operands have polynomial bit length.
Consequently the entire finite test is
polynomial in $L$, with no factoring of $R$,
conductor calculation, or point-counting
oracle. A negative output includes a prime
and two different counts as an unconditional
certificate; correctness of every positive
output uses the stated RH input.
\end{proof}

NT-E14--E15 are precise published theorem
transfers with their analytic field families
identified from the full source proofs.
NT-E16 is a fully proved formulation of the
finite-quotient method with a weaker elementary
constant and a direct input model; no
literature-priority claim is made.
Geometrically, this detects isomorphism of
the semisimple rank-two Tate local systems
using closed points of the good-reduction
open of $\operatorname{Spec}\mathbf Z$,
then invokes Faltings to recover
$\mathbf Q$-isogeny. The finite quotients
have nonabelian images in general, but their
Dedekind functions are trivial-character
instances of $M$. No non-CM $\mathrm{GL}_2$
GRH or general Artin holomorphy has entered.
Work/run: \path{effective_elliptic_detection.tex},
\path{GRH_NIGHT/number_theory}.

\subsection*{NT-E17--E18. A quadratic-Dirichlet input for CM recognition}

Write $A_{\rm quad}$ for RH for $\zeta_{\mathbf Q}$ and
for every primitive quadratic Dirichlet $L$-function.
This is a subfamily of the user-designated input $A_D$.
No Hecke function over a number field of degree greater
than one needs to be assumed separately below.
For a nonzero integer polynomial $H$, put
\[
 h=\deg H,\qquad
 \mathcal H=\max\{1,\log\max_i|H_i|\}.
\]
The height is padded by one to include polynomials
whose nonzero coefficients are units.

\paragraph{Exact classical inputs and scope.}
We use Cremona--Sutherland,
\emph{Computing the endomorphism ring of an elliptic
curve over a number field}, Contemporary Mathematics~796
(2024), 75--102, in author version
arXiv:2301.11169v4 (15 December 2023).
All page numbers in this subsection refer to that
author version.

Its Proposition~19, pp.~9--10, proves under GRH that
every imaginary quadratic order discriminant $D<0$
of class number $h$ satisfies
\[
 |D|\le
 \left(\frac{12e^{2\gamma}}{\pi}\right)^2
 h^2\bigl(\log\log(h+1)+4\bigr)^4.
 \tag{E17.1}
\]
The analytic input in this proof is
Languasco--Trudgian, \emph{Uniform effective estimates
for $|L(1,\chi)|$}, Corollary~1, author p.~2,
applied to the primitive quadratic character of the
fundamental discriminant.
The remaining reduction from maximal to arbitrary
orders uses the class-number formula (2.1) and an
unconditional Euler-totient estimate.
Thus $A_{\rm quad}$ supplies the analytic hypothesis
of (E17.1). The cited small-discriminant
tabulations and numerical range checks in that
published proposition are retained as classical
inputs; they were not rerun here.

Let $\Phi_\ell(X,Y)\in\mathbf Z[X,Y]$ be the classical
modular polynomial for cyclic $\ell$-isogenies.
For monic irreducible $H\in\mathbf Z[X]$ and a root $j$,
define
\[
 m_\ell(j)=
 \operatorname{ord}_H
    \operatorname{Res}_Y\bigl(H(Y),\Phi_\ell(X,Y)\bigr).
 \tag{E17.2}
\]
Cremona--Sutherland equation (2.2), p.~8, gives
this equality with the number of conjugates joined
to $j$ by a cyclic $\ell$-isogeny, with multiplicity.
Their Corollary~14, p.~8, gives $m_\ell(j)=2$
for a CM $j$ of order discriminant $D$ whenever
$\ell\nmid D$ splits in $\mathbf Q(\sqrt D)$.
Their Proposition~16, p.~9, gives
$m_\ell(j)\le1$ for every prime $\ell$ if $j$ is not CM.
These are classical algebraic statements, independent
of GRH. Their supplied derivations are read here;
the upstream theory of isogeny volcanoes and
$K$-curves remains cited.

The published Algorithm~1 has conditional correctness
in Theorem~24, p.~14, and deterministic complexity
in Theorem~27 and its full proof, p.~14.
For that particular algorithm the published time
and space bounds are respectively
$(h^2\mathcal H)^{1+o(1)}$ and
$(h\mathcal H)^{1+o(1)}$.
For our modified test we assert the weaker
polynomial bound, with its exact-arithmetic
verification supplied below.
We use the deterministic polynomial-time computation
of $\Phi_\ell$ explicitly cited in the proof
of Theorem~27.
A direct primary locator is
Blake--Csirik--Rubinstein--Seroussi,
\emph{On the Computation of Modular Polynomials
for Elliptic Curves}, \S2, pp.~4--7:
the exact power-series and CRT algorithm is
given there, with the bit-complexity discussion
and coefficient bound (9) on p.~7.
The modified test below examines all the bounded
degrees, so it does not require the source's
least-ordinary-prime selection or its finite list
of abelian discriminants.

\paragraph{NT-E17: a bounded algebraic CM certificate under $A_D$.}
Let $H\in\mathbf Z[X]$ be monic and irreducible,
of degree $h\ge1$, and set
\[
 b_h=\operatorname{bitlength}(h+1),\qquad
                     P_h=(72b_h+264)^2.
 \tag{E17.3}
\]
Under $A_{\rm quad}$, $H$ is a Hilbert class
polynomial for an imaginary quadratic order if
and only if there is a prime $5\le\ell\le P_h$ with
\[
 H(X)^2\ \bigm|\ 
       \operatorname{Res}_Y\bigl(H(Y),\Phi_\ell(X,Y)\bigr)
                           \quad\hbox{in }\mathbf Z[X].
 \tag{E17.4}
\]
Scanning these primes and performing the exact
divisibility tests is deterministic and takes
polynomial time in $h\mathcal H$.
A positive answer, accompanied by $\ell$ and
the exact divisibility identity, is valid
unconditionally. The universal correctness of
a negative answer uses $A_{\rm quad}$.
Irreducibility is part of the stated input promise.

\begin{proof}
If (E17.4) holds, (E17.2) gives $m_\ell(j)\ge2$.
The cited non-CM bound therefore forces $j$ to be CM.
Its monic irreducible polynomial is then exactly
the Hilbert class polynomial of its endomorphism
order. This implication is unconditional.

For the converse suppose $H=H_D$.
First we replace (E17.1) by an elementary bound
whose constants can be used without computing
$\gamma$, $\pi$, or logarithms in the algorithm.
For $h\ge1$,
\[
 \log\log(h+1)+4
 \le\log(h+1)+4
 \le h+4\le3(h+1).
\]
Using $\gamma<1$, $e<3$, and $\pi>3$ in (E17.1),
we obtain
\[
 |D|<1296\cdot81\,h^2(h+1)^4
       <2^{17}(h+1)^6.
 \tag{E17.5}
\]
This bound includes $D=-3,-4$ and nonmaximal orders
because (E17.1) does.

Set $m=\operatorname{rad}(6|D|)$ and
\[
 K=\mathbf Q(\sqrt D),\qquad
 T=K(\sqrt m).
\]
The integer $m$ is positive, even, and squarefree.
The field $\mathbf Q(\sqrt m)$ is real quadratic
and $K$ is imaginary quadratic, so $T/\mathbf Q$
is biquadratic of degree four.
It is unramified outside the primes dividing $m$.
Every prime dividing $m$ ramifies already in
$\mathbf Q(\sqrt m)$, since its discriminant is $4m$.
Consequently every unramified prime of $T$ is
prime to $6D$.

The full Euler factorization of $\zeta_T$ is
$\zeta_{\mathbf Q}$ times the three quadratic
Dirichlet functions of its quadratic subfields.
Thus $A_{\rm quad}$ supplies RH for $\zeta_T$
and $\zeta_{\mathbf Q}$.
Apply NT-C03 to the identity class of
$\operatorname{Gal}(T/\mathbf Q)$.
There is a prime $\ell$ splitting completely in $T$,
with $\ell\nmid6D$ and
\[
                 \ell\le(4\log|\operatorname{disc}T|+15)^2.
\]
The Galois different bound proved in NT-E16,
with absolute degree four, gives
\[
 \log|\operatorname{disc}T|
       \le3\log m+4\log4
       \le3\log(6|D|)+4\log4.
 \tag{E17.6}
\]
It follows that
\[
 \begin{aligned}
 4\log|\operatorname{disc}T|+15
   &\le12\log|D|+12\log6+16\log4+15\\
   &<12\log|D|+60\\
   &<72\log(h+1)+204\log2+60\\
   &<72b_h+264.
 \end{aligned}
 \tag{E17.7}
\]
For the second line one can use
$\log6<9/5$ and $\log4<7/5$, which give the
stronger numerical sum $59$.
Both inequalities follow by summing the exponential
series through degree six:
the resulting positive margins over $6$ and $4$
are respectively $42629/1250000$ and
$592609/11250000$.
The last line uses $h+1<2^{b_h}$ and $\log2<1$.
We have proved $5\le\ell\le P_h$, $\ell\nmid D$,
and splitting in $K$.
The cited CM formula gives $m_\ell(j)=2$,
and hence (E17.4).

Finally $P_h=O(\log^2(h+2))$, so both the number
and the sizes of the tested primes are polynomial
in the input length.
The cited deterministic construction of
$\Phi_\ell$ takes polynomial time in $\ell$,
and its output has polynomial bit size.
Here is an elementary polynomial bound for
the rest of the exact computation.
The Sylvester matrix defining the resultant
has size $N=h+\ell+1$, and its entries are
integer polynomials in $X$ of degree at most
$\ell+1$ with polynomially bounded coefficient
bit lengths.
Its determinant has degree at most
$d=N(\ell+1)$.
Evaluate at $X=0,\ldots,d$.
Each resulting integer matrix has entries of
polynomial bit length, and its determinant can
be computed by exact rational elimination
in polynomial time.
The determinant-size bound
$N!(\max|a_{ij}|)^N$ bounds all these outputs.
Rational interpolation at the $d+1$ distinct
integers recovers the resultant in polynomial
time; the factorial denominators of the
Lagrange basis have polynomial bit length.
Finally division by the monic polynomial $H^2$
uses polynomially many exact integer operations.
At each division step the coefficient bit bound
increases by at most the coefficient bound for
$H^2$ plus a logarithm of its degree, so it
remains polynomial throughout.
The remainder vanishes exactly when (E17.4)
holds. This proves the stated complexity
and both directions of the test.
\end{proof}

The only analytic characters in this proof are the
quadratic character used in (E17.1), the three
quadratic characters of the displayed biquadratic
field, and the trivial character.
In particular $M$, ring-class Hecke RH, general
Artin holomorphy, and elliptic $\mathrm{GL}_2$ RH
are not additional inputs.
This is a new formulation and input reduction,
with no claim of literature priority.

\paragraph{NT-E18: translation to an elliptic curve over a varying number field.}
Suppose a number field
$F=\mathbf Q[t]/(f(t))$ is given by a monic
irreducible $f\in\mathbf Q[t]$, and a nonsingular
Weierstrass equation for $E/F$ is given with
its coefficients represented in the power basis.
The defining polynomial is supplied as a dense
coefficient list, and field elements as dense
power-basis vectors. All rational coefficients
are encoded in binary, so $n=\deg f$ is at most
the total input length;
the assertion that $f$ defines a field is the
input promise.
Under $A_{\rm quad}$, geometric CM of $E$
can be decided deterministically in polynomial
time in this total input length.
No maximal order, conductor, integer factorization,
or endomorphism basis is required.

\begin{proof}
Compute $\Delta\ne0$ and $j=c_4^3/\Delta$ using
exact arithmetic in $F$.
Multiplication is reduction modulo $f$;
inversion is rational linear algebra.
These operations have polynomial bit complexity.
Let $n=\deg f$. Compute the coordinate vectors
of $1,j,\ldots,j^n$ and find the first dependence.
It gives the monic minimal polynomial
$H_j\in\mathbf Q[X]$, of degree $h\le n$.

For clarity the coefficient-size assertion does
not assume a precomputed integral basis.
The multiplication matrix of $j$ has rational
entries of polynomial bit length in the input.
After a common denominator is cleared, its first
$n$ powers have polynomial bit length: multiplying
integer matrices increases the logarithmic
entry bound by at most the previous bound plus
the fixed matrix-entry bound and $\log n$.
Cramer's determinant bounds for the first
linear dependence therefore give coefficients
of $H_j$ of polynomial bit length.
Exact elimination consequently takes polynomial
time as well.

If $H_j$ has a nonintegral coefficient, then $j$
is not an algebraic integer and cannot be a
CM $j$-invariant; return non-CM.
Otherwise $H_j$ is a monic irreducible integer
polynomial. Apply NT-E17 to it.
The degree and padded logarithmic height of
$H_j$ are polynomially bounded in the original
input length, so the complete computation is
polynomial. The characterization of CM
$j$-invariants as the roots of the Hilbert class
polynomials proves correctness.
\end{proof}

This recognizes whether
$\operatorname{End}_{\overline F}(E)$ is $\mathbf Z$
or an imaginary quadratic order.
It does not output the order discriminant or
explicit endomorphisms.
Cremona--Sutherland Algorithm~2 does output
the discriminant, and its correctness and
termination are Theorem~30, p.~15.
Its near-quasilinear expected running time is
Theorem~35, p.~17, whose hypotheses include
the separate prime-distribution heuristic (H)
and the computational heuristics in its
reference [45, \S7.1], as well as GRH.
Neither $A_D$ nor $M$ by itself removes those
extra hypotheses.
Likewise Enge's floating-point class-polynomial
complexity, Corollary~1.3 of arXiv:cs/0601104v2,
pp.~2--3, explicitly retains a rounding-error
heuristic; it is not used in either proof above.
These distinctions concern the cited algorithms,
not an assertion that unconditional CM recognition
is impossible.

Geometrically (E17.4) detects the quadratic
endomorphism algebra of $H^1(E)$ through the
cyclic-isogeny correspondences on the moduli
line of elliptic curves. The conditional
existence of a bounded correspondence degree
comes from a finite biquadratic Frobenius cover
of an arithmetic open, in the finite-etale
dictionary of SGA~1, Expose~V, \S7.
The result distinguishes CM torus monodromy
from the non-CM case without proving a
rank-two automorphic GRH.
Work/run: \path{cm_recognition_dirichlet.tex},
\path{GRH_NIGHT/number_theory}.

% END number_theory

% BEGIN geometry_programme
\section{Geometry, motives, and the Grothendieck programme}
\label{sec:geometry}

The letters A and M have the dependency meaning fixed in the introduction.
In particular, accepting the stipulated endpoint A is not a certification of
every intermediate statement in the note, and assuming M does not certify
the functional analysis in its proposed proof. The findings below are new
formulations proved in this run; no claim of priority over the literature is
intended. Earlier findings are used with their original identifiers.

\subsection{The transfer results already established}
\label{ag:transfer}

\paragraph{AG-AT (the exact Artin--Tate domain).}
For a finite-type scheme over $\mathbb Z$ with exact all-prime point counts
\[
 \#X(\mathbb F_{p^r})=
 \sum_j p^{jr}\operatorname{Tr}(F_p^r\mid V_j^{I_p}),
\]
where the $V_j$ are finitely many virtual finite-image representations,
the identity
\[
 \zeta_X(s)=\prod_j L(V_j,s-j)
\]
is unconditional. This is ZETA\_TWO ZT01, in
\path{RETURN_ZETA_TWO_final_IMPLICATIONS.md}, Section 1,
extending BRIDGES F1. If the representations have abelian image, A gives
noninteger global divisor support on the lines $\Re s=j+1/2$.
AG-ARTIN records the arbitrary finite-image case: M gives the corresponding meromorphic
divisor-support assertion through integer Brauer induction (ZT02).
Negative exponents can produce poles; Artin holomorphy is an extra question.

\paragraph{AG-DED and AG-LOCAL (ramification and models).}
The permutation representation on the embeddings of a number field has
inertia invariants giving the Euler factors of $\operatorname{Spec}\mathcal
O_K$ at ramified primes as well. Thus every Dedekind zeta function is covered
by M, and abelian fields over $\mathbb Q$ are covered by A (ZT03).
If a factorization holds only away from a finite set $S$, the actual factors
at $S$ must be restored (ZT04). The counterexample ZT05, catalogued as AG-MODEL,
\[
 \zeta_{\operatorname{Spec}\mathbb Z[1/p]}(s)
       =(1-p^{-s})\zeta(s),
\]
has additional zeros $2\pi i k/\log p$, $k\ne0$. Such local divisor lines
are not shifted global critical lines. This is a statement about the
integral model, not a failure of the generic-fibre motivic normalization.

\paragraph{AG-CM (the potentially CM domain).}
ZETA\_TWO CM01 supplies the all-weight extension to pure objects in the
absolute-Hodge category generated by Artin motives, Tate motives, and
potentially CM abelian varieties. Its cited CM descent input gives a signed
product of Hecke functions of one weight; norm twisting makes these
characters unitary. The required M therefore includes infinite-order
unitary characters. The result is meromorphic divisor RH, and does not
declare absolute-Hodge projectors algebraic. The exact prior locator is
\path{RETURN_ZETA_TWO_final_IMPLICATIONS.md}, Section 2, CM01,
with its reported Anderson Section 5.7, Corollary 5.7.2 and Remark 5.7.3.
The original source has now been recovered and checked: G.~W.~Anderson,
\emph{Cyclotomy and an extension of the Taniyama group}, Compositio
Math.~57 (1986), 153--217. The category is defined in Section 5.3,
printed p.~178 (PDF p.~27); Theorem 6, Section 5.6, printed pp.~179--180
(PDF pp.~28--29), identifies its one-dimensional characters with algebraic
Hecke characters, including the inertia and Frobenius comparisons.
Corollary 5.7.2 and Remark 5.7.3 are both on printed p.~181 (PDF p.~30).
The former gives integer exponents, with no positivity claim; the latter
gives the common weight. Its internal reference to ``Corollary 5.8.2'' is
a printed cross-reference error, not a new result used here.

\paragraph{AG-CM-REC and AG-CM-PERIOD (classical conclusions already
stronger than zero location).}
In this precise potentially CM absolute-Hodge category, equality of the
whole $L$-functions implies isomorphism of objects: Anderson,
Proposition 5.7.5, printed pp.~181--182 (PDF pp.~30--31). This is
unconditional. The proof uses semisimplicity and Corollary 5.7.4's
identification, for a weight-zero object, of the pole order at $1$ with
the dimension of its invariant part. Equality of the local factors makes
the tensor products with a dual agree away from finitely many primes;
their local factors are nonzero and nonpolar at $1$. A pole of
$L(M\otimes M^\vee,s)$ therefore gives a nonzero morphism
$M\longrightarrow N$, and induction on the common rank proves the
claim. The proof was read in full, rather than inferred from GRH.
This is a recognition theorem for semisimple motives; it supplies no
effective finite bound for general infinite CM monodromy.

Anderson's Section 5.8, Theorem 7, printed p.~182 (PDF p.~31), also
proves Deligne's period conclusion
$L(M,0)/c^+(M)\in\mathbb Q$ for the stated \emph{tractible, critical}
objects: those for which the $L$-function is a single algebraic Hecke
$L$-function over a totally real or CM field. This is a separate
classical special-value theorem, using Siegel and Blasius, and is not
enlarged to arbitrary motives by A or M. It complements, rather than
replaces, the cycle/height and integral comparison gates below.

\paragraph{AG-WEIGHT and AG-NONCM (the residual geometric factors).}
With geometric Frobenius,
\[
 \operatorname{wt}N(n)=\operatorname{wt}N-2n,
 \qquad L(N(n),s)=L(N,s+n).
\]
The line for weight $w$ is $\Re s=(w+1)/2$. Positive weight is not the
remaining obstruction: CM objects of positive weight are covered by M,
whereas AG-NONCM-W0 records that $\operatorname{Sym}^2h^1(E)(1)$ for a non-CM elliptic curve still
has a non-CM factor despite having weight zero. AG-SCHEMES carries ZT06: it reduces a finite-type
scheme to finitely many smooth projective geometric factors and actual
exceptional fibres. ZT07 proves that, under A, an off-line elliptic zero
in $1/2<\Re s<3/2$ cannot be cancelled by the two Tate factors in the
good-model formula
\[
 \zeta_{\mathcal E}(s)
   =\zeta^S(s)\zeta^S(s-1)/L^S(E,s).
\]
These are retained reductions, not a proof of RH for the non-CM factor.
Framework locators carried from ZETA\_TWO are SGA 4 XVII Sections 5.1--5.3;
SGA 5 XV Section 3, no.~2, unnumbered theorem, equation (1);
EGA IV$_3$ (8.8.2); and Deligne, Weil I, Theorem 1.6, and Weil II,
Theorem 3.3.1.

\subsection{AG-FROB: a finite recognition theorem for Artin--Tate objects}
\label{ag:frob}

Let $L/\mathbb Q$ be a nontrivial finite Galois extension, with group $G$,
absolute discriminant $D_L$, and degree $n_L$. Put
\[
 B_L=(4\log D_L+\tfrac52 n_L+5)^2.
\]
Let
\[
 U=\bigoplus_j V_j(-j),\qquad W=\bigoplus_j W_j(-j),
\]
where all sums are finite and $V_j,W_j$ are actual finite-dimensional
characteristic-zero representations of $G$. Their Euler polynomials at
unramified primes are understood as those of the associated compatible
Artin--Tate systems; a single auxiliary prime may instead be chosen greater
than $B_L$.

\textbf{Statement.} Assume RH for $\zeta_L$ and $\zeta_{\mathbb Q}$.
If
\[
 \det(1-TF_p\mid U)=\det(1-TF_p\mid W)
 \quad\text{for every }p\le B_L\text{ unramified in }L,
\]
then $V_j\simeq W_j$ for every $j$, and hence $U\simeq W$.
For rational coefficients this is also an isomorphism in the rational
Artin--Tate category. A supplies the hypothesis when $L/\mathbb Q$ is
abelian; M supplies it in general. When $L=\mathbb Q$, the one prime $p=2$
suffices.

\begin{proof}
At an unramified prime, every eigenvalue in $V_j(-j)$ has the form
$p^j\omega$, where $\omega$ is a root of unity. Therefore its complex
absolute value is $p^j$. Distinct integers $j$ give distinct absolute
values, so equality of the total Euler polynomials separates uniquely into
equality of the polynomials in each weight. It follows that
$\operatorname{Tr}(F_p\mid V_j)=\operatorname{Tr}(F_p\mid W_j)$ for every
$j$ and every prime in the hypothesis.

Bach--Sorenson, \emph{Explicit bounds for primes in residue classes},
Math. Comp. 65 (1996), Theorem 5.1, printed p.~1729 (PDF p.~13), gives
an unramified prime $p\le B_L$ in every conjugacy class of $G$.
The precise analytic assumptions are stated in their Section 2,
printed p.~1720: zero-freeness of $\zeta_L$ and $\zeta_{\mathbb Q}$ in
$\Re s>1/2$. The reduction from cyclic to general Galois extensions is
Corollary 3.3, p.~1721. Thus the characters of $V_j$ and $W_j$ agree on
every conjugacy class. In characteristic zero finite-group representations
are semisimple, and equal characters give isomorphic representations.

For completeness, this last assertion also descends to a non-algebraically
closed characteristic-zero coefficient field $E$. The space of
$G$-equivariant maps is the solution space over $E$ of finitely many linear
equations. If it contains an invertible map after scalar extension, its
determinant polynomial is nonzero on that $E$-vector space. Since $E$ is
infinite it has a nonzero value at an $E$-point. That point is an invertible
$G$-equivariant map over $E$. For $L=\mathbb Q$, the eigenvalues $2^j$
already recover all multiplicities from the polynomial at $2$.
\end{proof}

In Grothendieck's language this is effective recognition inside a specified
semisimple Artin--Tate realization category. It extends
ALPOGE\_BRIDGES AQ-FINITE, whose sharper abelian conductor bound remains
available. It does not construct algebraic projectors for an arbitrary
geometric realization. Moreover, recognition of an Artin motive is weaker
than recognition of a finite cover as a scheme; the connected Gassmann
example in CX-GEO-004's companion finding makes this distinction explicit.
The forbidden-prime variant is handled there by a quadratic compositum.

\paragraph{The elliptic $h^1$ recognition extension (NT-E16).}
There is also an effective recognition statement beyond finite
monodromy, with the geometry supplied by Faltings's theorem.
For any two elliptic curves over $\mathbb Q$, including CM
curves, given by nonsingular integral Weierstrass equations,
put $b=\operatorname{bitlength}(2|\Delta_1\Delta_2|)$.
NT-E16 proves under M that equality of the point counts at
all primes $p\le(1532b+10181)^2$ not dividing
$2\Delta_1\Delta_2$ is equivalent to $\mathbb Q$-isogeny.
It uses only the Dedekind RH instances for the finite
Galois detection fields of degree at most $384$, together
with $\zeta_{\mathbb Q}$; it assumes no RH for the standard
elliptic $L$-function.  The complete proof also gives a
deterministic polynomial bit-time decision from those
equations without factoring their discriminants.
Thus the closed points of the actual good-reduction open
of $\operatorname{Spec}\mathbb Z$ recognize the rational
$h^1$ isogeny motive in this family.  An isogeny and its
dual give inverse rational correspondences after division
by their degree.  Neither isomorphism of the integral
curves nor an integral inverse correspondence is asserted.
This is the precisely scoped elliptic extension of the
recognition discussion in AG-FROB and AG-CM-REC.

\paragraph{AG-COVER: equal Artin motives need not give equal covers.}
Put $L=\mathbb Q(\sqrt2,\sqrt3)$ and let $K_1,K_2,K_3$ be its three
quadratic subfields. The finite schemes
\[
 X=\operatorname{Spec}L\amalg\operatorname{Spec}\mathbb Q
       \amalg\operatorname{Spec}\mathbb Q,
 \qquad
 Y=\coprod_{i=1}^3\operatorname{Spec}K_i
\]
are not isomorphic, but their rational Artin motives, and all their local
Artin Euler factors, are equal. Consequently
\[
 \zeta_L(s)\zeta_{\mathbb Q}(s)^2
       =\zeta_{K_1}(s)\zeta_{K_2}(s)\zeta_{K_3}(s).
\]
This identity includes the factors at ramified primes.

\begin{proof}
Write $G=\operatorname{Gal}(L/\mathbb Q)=C_2\times C_2$, with its three
subgroups $H_i$ of order two. The character of
$\mathbb Q[G]\oplus\mathbb Q^2$ takes value $6$ at the identity and
$2$ at every other element. Each nonidentity element belongs to exactly
one $H_i$, so the character of $\bigoplus_i\mathbb Q[G/H_i]$ has the
same values. Characteristic-zero semisimplicity and the descent argument
in AG-FROB give an isomorphism of rational representations. Isomorphic
representations have isomorphic inertia invariants with their residual
Frobenius action, proving equality of every local factor and hence the
zeta identity in the initial half-plane and by continuation. On the other
hand the degrees of the connected components are $4,1,1$ for $X$ and
$2,2,2$ for $Y$, precluding a scheme isomorphism.
\end{proof}

The same example over $\mathbb Z$ uses the respective rings of integers;
over $\mathbb Z[1/6]$ these are finite \emph{\'etale} covers. Thus even
complete Euler data forget the multiplication on an Artin algebra. A
applies to all fields here because they are abelian over $\mathbb Q$,
but does not restore the forgotten multiplication. The independent
transitive seven-element example CX-GEO-005 treats the connected
$G$-set distinction without asserting an arithmetic realization.

\subsection{AG-LAMBDA: the precise closure supplied by tensor operations}
\label{ag:lambda}

\textbf{Statement.} The finite direct sums of abelian Artin characters with
integer Tate twists are closed under tensor products, duals, exterior and
symmetric powers, and direct summands in characteristic zero. All their
$L$-functions have the A divisor-support conclusion. Replacing abelian
finite image by arbitrary finite image gives the corresponding M conclusion.
The same operations preserve the potentially CM absolute-Hodge category
used in CM01, so its M conclusion applies there as well.

\begin{proof}
Over an algebraic closure of the coefficient field an abelian Artin--Tate
object is a sum of lines $\chi_a(-j_a)$. A tensor product or a monomial
appearing in a symmetric or exterior power is a line
$(\prod_a\chi_a^{e_a})(-\sum_a e_aj_a)$. The character still has finite
abelian image and the Tate index is an integer. Duality replaces the
character and index by their inverses and negatives. Direct summands remain
in the same semisimple category, and descent preserves the asserted
factorization. The nonabelian assertion follows because all these
operations still factor through a finite group and commute with Tate
twists; apply ZT02 to their resulting Artin representations. Finally, the
potentially CM category is the tensor, dual, and summand closure in CM01.
Symmetric and exterior powers are the images of the usual idempotents in
the symmetric-group action on tensor powers, so belong to it too.
\end{proof}

The placement is SGA 6 V Section 2 and Section 7, especially (7.1.1),
and the determinant/$\lambda$-ring discussion carried in ZETA\_TWO ZT08
and SV01. These closures do not generate an arbitrary non-CM $GL_2$
motive from CM objects; enlarging a list of tensor operations is not the
missing analytic theorem.

\subsection{AG-SIGN: central signs with the necessary normalization}
\label{ag:sign}

\textbf{Statement.} Let $F$ be a nonzero holomorphic function on a
neighbourhood of $[c,\infty)$, real-valued there. Suppose $F$ has no zeros
on $(c,\infty)$ and is positive at one sufficiently large real point.
If $r=\operatorname{ord}_{s=c}F$, then
\[
 \frac{F^{(r)}(c)}{r!}>0.
\]
If, moreover, $F(c+z)=\varepsilon F(c-z)$ near $0$, with
$\varepsilon\in\{1,-1\}$, then $(-1)^r=\varepsilon$.
For a real, standardly normalized, holomorphic Hecke or CM $L$-function,
the absence of real zeros to the right of the centre follows from M
together with its usual analytic continuation and Euler-product theory.
The analogous Dirichlet statement uses A whenever no pole is present.

\begin{proof}
A continuous real nonvanishing function has constant sign on the connected
interval $(c,\infty)$. The normalization makes that sign positive.
Write $F(c+t)=a_rt^r+O(t^{r+1})$, with $a_r\ne0$. For $t>0$ the
left-hand side is positive, so $a_r>0$ after division by $t^r$ and passage
to the limit. Comparing the leading terms in the functional equation gives
$a_r=\varepsilon(-1)^r a_r$, hence the parity assertion.
\end{proof}

\paragraph{AG-HEIGHT (what this proves for a supplied height formula).}
If $F(c)=0$, then $F'(c)\ge0$: it is positive when $r=1$ and zero when
$r>1$. Consequently, an independently established identity
$H(Z)=C F'(c)$ with a specified constant $C>0$ yields $H(Z)\ge0$.
The elementary deduction is complete; the existence, domain, and
normalization of the height identity are extra inputs. It does not assert
that every homologically trivial cycle has such a formula, or that its
$L$-function has simple central vanishing. For an identity involving the
first nonzero derivative, the sign is strictly positive. For a formula with
an additional sign, that sign must be retained.

\paragraph{AG-SIGN-NORM (a counterexample to the unnormalized inference).}
Zero locations and functional-equation parity alone do not determine the
sign: $F(s)=s-c$ and $-F(s)$ have exactly the same zeros and the same odd
functional equation, but their central derivatives have opposite signs.
This is a counterexample to the stated analytic inference, not an example
of two standardly normalized arithmetic $L$-functions.

These findings refine \path{OURS_HODGE_BRIDGE_SKETCH_20261121.md},
Section 3.5. Their archive placement is cote 44, scans 2, 4 and 6,
Theorems 4--5, and scans 8--9, the letter to Serre of 12 August 1964,
as reported in that attached sketch. No original scan has been inferred
from this sign argument. The remaining arithmetic Hodge-index gate is
the cycle/height comparison for the stated class of cycles, not a sign
convention suppressed in a formula.

\subsection{AG-CRITICAL: why off-centre critical integers gain little}
\label{ag:critical}

\textbf{Statement.} Suppose $N$ is a pure, integral-weight, explicitly
Hecke-factored motive of weight $w$, with its usual functional equation
and archimedean factors. At an endpoint of the unitary strip, require each
factor in the specified signed Hecke factorization (or its dual, when the
functional equation is used) to be nonpolar. This factorwise requirement
prevents a denominator pole from producing a zero. Let $m$ be critical in
Deligne's sense: $L_\infty(N,s)$ and $L_\infty(N^\vee,1-s)$
are both regular at $s=m$. If $m\ne(w+1)/2$, then the standard
zero-free half-plane and boundary for unitary Hecke $L$-functions,
together with the functional equation, already give $L(N,m)\ne0$.
M is not needed for this conclusion.

The precise unitary boundary theorem used here is
D.~Ramakrishnan--R.~J.~Valenza, \emph{Fourier Analysis on Number Fields}
(1999), Section 7.7, Theorem 7-28, printed p.~289: for every number
field and every unitary idele class character, there is no zero at
$1+it$. Section 7.4, Theorem 7-19, printed p.~279, supplies the completed
continuation and functional equation, with its pure norm-character pole
exceptions retained. These are classical inputs independent of M.

\begin{proof}
After unitary normalization set $z=m-w/2$. If $w$ is even, $z$ is an
integer; if $w$ is odd, $z$ is a half-integer. An integer $m$ with
$0<z<1$ therefore has $z=1/2$, and is precisely the excluded centre.
For every other $m$ one has $z\ge1$ or $z\le0$. In the first case
standard unitary Hecke nonvanishing in $\Re z\ge1$, away from a pole,
gives the result for each factor and hence its finite signed product.
In the second case the functional equation relates the completed value
to its dual at $1-z\ge1$. Criticality ensures that both archimedean
factors are finite; gamma functions have no zeros. The nonzero dual
value therefore gives a nonzero original value. The factorwise nonpolar
condition excludes a Tate pole or its inverse being misread as a finite
nonzero special value. For a product with nonnegative exponents, the
usual no-pole condition implies this condition.
\end{proof}

Thus the strongest new real-axis information from M lies between the
off-centre integer points and at signs of central leading coefficients;
it does not evaluate regulators, periods, central ranks, or integral
torsion. This is the special-value distinction in ZETA\_TWO Section 4
and in the Hodge sketch's edge-versus-centre table. The relevant motivic
conventions are Deligne, \emph{Valeurs de fonctions $L$ et p\'eriodes
d'int\'egrales}, Sections 1.1--1.3; the period conjectures numbered 1.8
and 2.8 are separate inputs.

\subsection{AG-DET: Hecke versions of the square-variable determinant}
\label{ag:det}

The following formulation records exactly which analytic data are needed.
Let $\Lambda$ be a nonzero entire function of order at most one, with
\[
 \Lambda(c+z)=\varepsilon\Lambda(c-z),\qquad
 \varepsilon\in\{1,-1\},
\]
and let $r=\operatorname{ord}_{s=c}\Lambda$. Assume that every other zero
is $c+i\gamma$, with $\gamma$ real and nonzero. Choose one positive
$\gamma$ from each pair, with multiplicity, and write the sequence as
$\gamma_1,\gamma_2,\ldots$ in nondecreasing order, so $\gamma_1$
denotes the least positive zero when the sequence is nonempty.
Define the removable entire function
\[
 H(t)=t^{-r}\Lambda(c+it),\qquad
 D(w)=H(\sqrt w)/H(0).
\]
No choice of square-root branch affects $D$.

\textbf{Statement.} The diagonal operator
$B=\operatorname{diag}(\gamma_j^{-2})$ is positive and trace class, and
\[
 D(w)=\prod_j(1-w\gamma_j^{-2})=\det(I-wB).
\]
In particular this applies under M to standard entire selfdual completed
Hecke functions and to the entire completed $L$-function of a CM elliptic
curve, with its centre $c=1$. The usual order, functional equation,
and completion are additional classical inputs, not assertions supplied
by the word ``GRH''. Poles must first be removed in a separately specified
entire completion; arbitrary meromorphic Artin quotients are not included.

\begin{proof}
Comparison of the first nonzero Taylor coefficients shows
$\varepsilon=(-1)^r$. Hence $H(-t)=H(t)$, so $D$ is entire by the even
power series for $H$. Its order is at most one half: a bound of order
$1+\delta$ for $H(t)$ gives a bound of order $(1+\delta)/2$ for $D(w)$.
Its zeros, with multiplicity, are exactly $\gamma_j^2$.
The convergence exponent of the zeros of an entire function of order at
most one is at most one. In particular
$\sum_j|\gamma_j|^{-2}<\infty$. Equivalently the diagonal positive
operator $B$ has finite trace norm. Hadamard factorization for $D$, whose
order is strictly less than one, has no nonconstant exponential factor
and uses genus-zero factors. Since $D(0)=1$, its constant factor is one.
This gives the displayed product. The Fredholm determinant of a diagonal
trace-class operator is that same absolutely locally uniformly convergent
product, proving the identity. An exact analytic locator for the
factorization used here is C.~McMullen, \emph{Advanced Complex Analysis},
Math 213a course notes, version 16 November 2026, Theorem 3.16,
printed p.~92 (PDF p.~93); Theorems 3.10--3.12 give the zero-exponent
bound used just above. This is classical background, not a new claim
dependent on that revision date.
\end{proof}

This extends the prior zeta realization BR-TR-01 to a precisely specified
Hecke domain. Its relationship to cote 1 scans 79--80 and 111--112 is
the trace/Fredholm theorem reported in
\path{OURS_INTAKE_RESEARCH_SA23_20261120.md}, Section 4.1,
geo-H-lidskii and geo-H-genus. Those sheets are not themselves an
arithmetic RH claim. The diagonal presentation here is defined from
the zeros. The further comparison MT-M29 identifies this very operator
with the inverse square of the positive-frequency norm generator on
the associated graded of the canonical locally finite module in the
actual adelic summation quotient. Its image-power filtration retains
every zero multiplicity while removing the nilpotent extensions.
The resulting metric uses the fixed Mellin coordinate and its
positivity still uses the stated A or M zero-location input.
Thus the particular arithmetic-module identification is now supplied;
an independently positive cohomology or GGP-period realization remains
a separate construction. The ordinary trace-GNS completion has only
one dimension per distinct zero and does not replace that associated
graded when multiplicities occur.

\subsection{AG-MOMENT and AG-HANKEL: recoverable positive spectral data}
\label{ag:moment}

Assume the hypotheses of AG-DET and infinitely many distinct positive
zeros. Put $\lambda_j=\gamma_j^{-2}$ and
\[
 s_k=\operatorname{Tr}(B^k)=\sum_j\lambda_j^k\quad(k\ge1).
\]
The logarithmic determinant expansion is
\[
 -\log D(w)=\sum_{k\ge1}\frac{s_k}{k}w^k
 \qquad (|w|<\gamma_1^2).
\]

\textbf{AG-MOMENT.} For every $k\ge1$,
\[
 s_k^{-1/k}\le\gamma_1^2\le\frac{s_k}{s_{k+1}}.
\]
The left bounds increase to $\gamma_1^2$; the right bounds decrease to it.
If $m_1$ is the multiplicity of the least positive zero, then
$\lim_{k\to\infty}s_k\gamma_1^{2k}=m_1$.

\begin{proof}
Let $a=\max_j\lambda_j=\gamma_1^{-2}$. Then
$s_k\ge a^k$ and $s_{k+1}\le a s_k$, proving the two inequalities.
The $\ell^p$ norm monotonicity gives
$s_{k+1}^{1/(k+1)}\le s_k^{1/k}$. One direct proof is to divide by
$s_k^{1/k}$: the resulting nonnegative numbers have $k$-th powers
summing to one, are each at most one, and thus have $(k+1)$-th powers
summing to at most one. Hence the left bounds increase.
Cauchy--Schwarz applied to
$\lambda_j^{k/2}$ and $\lambda_j^{(k+2)/2}$ gives
$s_{k+1}^2\le s_ks_{k+2}$, so the right bounds decrease.

There are finitely many indices with $\lambda_j=a$, since $\sum\lambda_j$
is finite. For $k\ge1$,
$0\le(\lambda_j/a)^k\le\lambda_j/a$, whose sum is finite.
Dominated convergence shows
$s_k/a^k\to m_1$. It follows that both $s_k^{1/k}$ and
$s_{k+1}/s_k$ tend to $a$, which proves both limits and the
multiplicity formula.
\end{proof}

\textbf{AG-HANKEL.} For all integers $h\ge1$ and $d\ge0$, the matrix
\[
 (s_{h+i+j})_{0\le i,j\le d}
\]
is positive definite. Its determinant is therefore strictly positive.

\begin{proof}
For a nonzero vector $a=(a_0,\ldots,a_d)$ and
$P(X)=\sum_{i=0}^d a_iX^i$, its Hermitian quadratic form equals
\[
 \sum_{j\ge1}\lambda_j^h|P(\lambda_j)|^2.
\]
This converges because the $\lambda_j$ are bounded and
$\sum\lambda_j^h<\infty$. All terms are nonnegative. A nonzero
polynomial has only finitely many distinct zeros, whereas the positive
$\lambda_j$ have infinitely many distinct values. At least one term
is positive. This proves strict positive definiteness.
\end{proof}

These consequences strengthen the two-by-two moment inequalities carried
in BR-TR-02. They describe actual determinant coefficients available in
the cote 1 analytic framework, while leaving the arithmetic/geometric
realization gate unchanged. They assert no simplicity and no
random-matrix statistics.

\subsection{AG-SIGNED: a finite moment witness for a surviving pole}
\label{ag:signed}

Consider a finite signed product of specified entire selfdual functions
of order at most one, with a common centre $c$, each satisfying the
zero-location hypotheses of AG-DET. Let $r\in\mathbb Z$ be the net
central order. Divide out the central Laurent factor and normalize its
nonzero constant term. Its square-variable quotient has the form
\[
 D(w)=\prod_{\lambda\in\mathcal L}(1-w\lambda)^{m_\lambda},
 \qquad m_\lambda\in\mathbb Z\setminus\{0\},\quad
 \lambda>0,\quad \sum_\lambda |m_\lambda|\lambda<\infty,
 \qquad D(0)=1.                                      \tag{AG-S1}
\]
Here $\mathcal L$ has no positive accumulation point, and the signed
multiplicities are taken \emph{after} cancellation. The form (AG-S1)
follows by applying AG-DET to each factor. More generally the proposition
below applies whenever (AG-S1) is independently supplied. Put
\[
 s_k=\sum_\lambda m_\lambda\lambda^k\quad(k\ge1),\qquad
 H_d=(s_{1+i+j})_{0\le i,j\le d}.
\]
These moments are the local logarithmic coefficients
$-\log D(w)=\sum_{k\ge1}s_kw^k/k$ near $0$.

\textbf{Statement.} All the matrices $H_d$ are positive semidefinite
if and only if $m_\lambda>0$ for every $\lambda\in\mathcal L$.
Equivalently, the normalized quotient has no noncentral poles. The
central order condition $r\ge0$ must be checked separately. If a
negative multiplicity survives, a finite matrix $H_d$ already detects it,
with the explicit witness described below.

\begin{proof}
Set $a=\sup\mathcal L$ and introduce the finite signed measure
\[
 \mu=\sum_\lambda m_\lambda\lambda\,\delta_\lambda
 \quad\hbox{on }[0,a].
\]
The empty support case is immediate. For every complex polynomial
$P(X)=\sum_{i=0}^d b_iX^i$, absolute convergence gives
\[
 \sum_{i,j=0}^d\overline b_i s_{1+i+j}b_j
       =\int |P(x)|^2\,d\mu(x).
\]
Positive multiplicities therefore imply all the required inequalities.

Conversely, suppose $m_{\lambda_0}<0$. The point $\lambda_0>0$ is
isolated in the compact set $\{0\}\cup\mathcal L$. Let
\[
 \delta=\operatorname{dist}
   (\lambda_0,(\{0\}\cup\mathcal L)\setminus\{\lambda_0\})>0,
 \qquad R>a+\lambda_0,
 \qquad q=1-\delta^2/R^2<1.
\]
For each nonnegative integer $N$, the polynomial
\[
 P_N(x)=\left(1-\frac{(x-\lambda_0)^2}{R^2}\right)^N
\]
has degree at most $2N$, takes the value $1$ at $\lambda_0$, and has
$|P_N(\lambda)|\le q^N$ at every other point of $\mathcal L$.
Writing $W=\sum_{\lambda\ne\lambda_0}|m_\lambda|\lambda$, we obtain
\[
 \int |P_N|^2\,d\mu
   \le -|m_{\lambda_0}|\lambda_0+Wq^{2N}.
\]
Choose $N$ with $Wq^{2N}<|m_{\lambda_0}|\lambda_0$; take $N=0$
if $W=0$. Then the coefficient vector of $P_N$ is a negative
quadratic-form witness for the matrix $H_{2N}$, of size $2N+1$.
This proves the converse and its finite-witness assertion. In (AG-S1)
a negative net multiplicity is exactly a pole at $w=\lambda^{-1}$,
which proves the analytic equivalence. Removing the central Laurent
factor erased its order, explaining the separate condition on $r$.
\end{proof}

Geometrically this strengthens the cote 1 determinant bridge in the
meromorphic setting. A or M can first localize an identified divisor;
positivity of these finite moment matrices is then an exact additional
condition preventing cancellation from leaving poles. It is not supplied
by zero location alone. In particular, universal Hecke GRH plus Brauer
induction does not prove this positivity, and hence does not prove Artin
holomorphy. The criterion is for the specified quotient: symmetrizing
an arbitrary nonselfdual function with its dual can cancel a pole between
the two factors, so holomorphy of that product must not be attributed to
each unsymmetrized factor. The local character-matrix test in MT-ART02
is a complementary criterion at an individual Artin divisor point.

\subsection{AG-EXT: the precise information that the determinant forgets}
\label{ag:ext}

The earlier findings BR-MIX-01 and ZETA\_TWO SV01--SV02 are retained.
At a good prime an adapted matrix for a mixed object is block triangular,
so its determinant equals that of its weight-graded object. The nonsplit
Kummer extension from the unit $2$ over $\mathbb Z[1/2]$ and its split
Tate object therefore have the same incomplete $L$-function. The omission
of ramified factors is essential: inertia invariants need not preserve a
short exact sequence. These findings refute recovery of extensions from
semisimplified good Euler data; they do not assert that all richer
arithmetic data agree.

AG-EXT-K records the independent rational mixed-Tate background from Deligne--Goncharov,
\emph{Groupes fondamentaux motiviques de Tate mixte},
arXiv:math/0302267v2 (27 February 2003), (2.1.3), printed p.~14
(PDF p.~16):
\[
 \operatorname{Ext}^1(\mathbb Q(0),\mathbb Q(n))=
 \begin{cases}
 0&n\le0,\\
 \mathcal O_S^\times\otimes\mathbb Q&n=1,\\
 K_{2n-1}(k)\otimes\mathbb Q&n\ge2.
 \end{cases}
\]
Propositions 2.2--2.3 on the same page identify the corresponding graded
free motivic Lie algebra, while Section 2.4 explicitly distinguishes a
free presentation from a canonical set of generators. These statements
are unconditional background. They do not calculate integral torsion or
the pro-$\ell$ kernels in the prompt's cote 149 programme.

\subsection{An indexed integral $K$-group statement and its exact Tate-twist consequences}
\label{ag:k8}

\paragraph{AG-K8: the inherited theorem now has a primary proof locator.}
The assertion $K_8(\mathbb Z)=0$ in RETURN\_ALPOGE\_BRIDGES,
R-KGROUP/F-SM14--16, is a separate unconditional theorem.
It is Theorem~6.4 of Dutour Sikiri\'c--Elbaz-Vincent--Kupers--Martinet,
\emph{Voronoi complexes in higher dimensions, cohomology of
$GL_N(\mathbb Z)$ for $N\ge8$ and the triviality of $K_8(\mathbb Z)$},
JIMJ~25 (2026), 565--590; the primary author version
arXiv:1910.11598v1 has Theorem~6.4 and its proof on p.~19.
The result and its short proof were already available in 2019;
2026 is the journal issue year, not a consequence of the adopted A.

The exact dependency chain also explains the integral input.
With $R=\mathbb Z[1/7!]$, their Proposition~6.1 (author p.~18)
gives $H_9(BQ;R)=0$, and Proposition~6.3 (p.~19) makes the
Hurewicz map $\pi_9(BQ)\otimes R\to H_9(BQ;R)$ injective.
Since $\pi_9(BQ)=K_8(\mathbb Z)$, this proves vanishing after
inverting $2,3,5,7$.  Weibel VI, Corollary~9.8 (chapter p.~55)
removes the 2-primary part; Proposition~10.5 (p.~63) removes
the $3,5,7$ parts because those primes are regular.
These steps prove the stated integral vanishing from the cited
homological computations and standard integral comparisons.
Neither a rational special value nor a zero-location assertion
supplies the missing integral information by itself.

Kupers's independent \emph{A short proof that $K_8(\mathbb Z)=0$},
dated 21 October 2019, Theorem~A and pp.~1--4, supplies a
second short route.  Its rank-filtration spectral sequence has
zero entries on total degree nine after inverting $7!$;
the rest of the proof uses the low $K$-groups and the sphere
spectrum to pass from homology to homotopy.  The complete
four-page proof was read, as were the complete final reductions
in the joint paper.  The underlying earlier Voronoi enumeration,
Church--Putman proof and Quillen--Lichtenbaum theorem remain
cited classical inputs; no fresh audit of their computations
or full proofs is represented here.  In the usual homological
Serre spectral-sequence convention, the short proof's displayed
transgression $(9,0)\to(0,8)$ is $d_9$, not its printed $d_8$;
this indexing correction does not affect the argument.

\begin{proposition}[AG-K8-TWIST: exact integral and finite-twist vanishing]
For every odd prime $p$, let $U_p=\operatorname{Spec}\mathbb Z[1/p]$,
let $C_p$ be the $p$-primary part of
$\operatorname{Cl}(\mathbb Q(\zeta_p))$, and let
$\omega:\operatorname{Gal}(\mathbb Q(\zeta_p)/\mathbb Q)
\to\mathbb Z_p^\times$ be the Teichmüller character.
Then, unconditionally,
\begin{align*}
 H^2_{\text{ét}}(U_p,\mathbb Z_p(5))&=0,
 &H^2_{\text{ét}}(U_p,\mu_p^{\otimes5})&=0,\\
 C_p^{(-4)}&=0,
 &K_{8+2a(p-1)}(\mathbb Z)\{p\}&=0\qquad(a\ge0).
\end{align*}
Here the integral cohomology is the standard continuous
$p$-adic group in the cited comparison, and braces denote
the entire $p$-primary subgroup.  The class-group assertion
is the component $C_p^{(p-5)}$ in the joint paper's
Corollary~7.1, author p.~20, for $p>3$; the same statement
at $p=3$ follows from the proof below.  These are precisely
indexed classical consequences, not new applications of A or M.
\end{proposition}
\begin{proof}
Weibel VI, Theorem~8.2 and Corollary~8.3 (chapter pp.~47--48)
give, for odd $p$ and $i>0$,
\begin{align*}
 K_{2i}(\mathbb Z)\{p\}&\simeq
 H^2_{\text{ét}}(U_p,\mathbb Z_p(i+1)),\\
 K_{2i}(\mathbb Z)/p&\simeq
 H^2_{\text{ét}}(U_p,\mu_p^{\otimes(i+1)}).
\end{align*}
Putting $i=4$ and using AG-K8 proves the first two assertions.
We give the class-group comparison explicitly to fix its
character and the passage from mod-$p$ to integral vanishing.
Put $R_p=\mathbb Z[\zeta_p,1/p]$ and
$G=\operatorname{Gal}(\mathbb Q(\zeta_p)/\mathbb Q)$.
The finite etale cover $\operatorname{Spec}R_p\to U_p$
has degree $p-1$, so its Hochschild--Serre sequence,
with coefficients of exponent $p$, reduces to taking
$G$-invariants.  There is one prime of $\mathbb Q(\zeta_p)$
over $p$, and it is principal.  Consequently
$\operatorname{Pic}(R_p)=\operatorname{Cl}(\mathbb Q(\zeta_p))$.
The field is totally imaginary and only that one finite
place is removed.  The global Brauer invariant sequence
therefore gives $\operatorname{Br}(R_p)=0$; this is exactly
Weibel (8.1.1) and Example~8.3.2.  The Kummer sequence now
identifies $H^2_{\text{ét}}(R_p,\mu_p)$ with
$\operatorname{Pic}(R_p)/p$.
Tensoring by the rank-one module $\mu_p^{\otimes4}$,
and then taking invariants, yields
\[
 H^2_{\text{ét}}(U_p,\mu_p^{\otimes5})
 \simeq\bigl((C_p/p)\otimes\mu_p^{\otimes4}\bigr)^G
 \simeq (C_p/p)^{(-4)}.
\]
Thus the character is $\omega^{-4}$, not $\omega^{-5}$.
The order $p-1$ of $G$ is a $p$-adic unit, so the
Teichmüller idempotents split $C_p$ integrally and commute
with reduction modulo $p$.  Hence
$C_p^{(-4)}/p=0$.  A finite abelian $p$-group equal to its
$p$-multiple is zero, proving the full-component vanishing.

Finally set $i=4+a(p-1)$.  The mod-$p$ cyclotomic
character has order dividing $p-1$, whence
$\mu_p^{\otimes(i+1)}\simeq\mu_p^{\otimes5}$.
The same comparison gives
$K_{8+2a(p-1)}(\mathbb Z)/p=0$.
Positive even $K$-groups of $\mathbb Z$ are finite;
thus their $p$-primary parts vanish.  This proves the
entire displayed family.  It is also the $k=2$
instance of the periodic criterion in Weibel VI,
Theorem~10.9 (p.~65), with the regular-prime case
provided by Proposition~10.5.
\end{proof}

Geometrically, this kills one precisely specified
Tate-twisted $H^2$ group on each $U_p$ and the matching
cyclotomic unramified abelian class-group component.
It does not kill all positive-weight relation modules,
choose normalized Soule generators, or identify the
full pro-$p$ kernel in the LM6/Ihara programme.
The attached Alpoge footnote and step~17 still do not
export an indexed theorem on their own; the theorem
and comparison just supplied are separate proved inputs.

\subsection{AG-KUMMER: an integral obstruction on the actual cote 149 sheets}
\label{ag:kummer}

The original Montpellier PDF was inspected here, not merely the intake.
It has a cover page: PDF p.~80 has repository footer 79, and so on.
At PDF p.~80/footer 79, equations (13)--(17) relate an inertia-normalizer
extension, its $H^2$ class, and Kummer cohomology. PDF p.~81/footer 80,
(18)--(19), identifies the relevant line bundle as the normal bundle of
the section. PDF p.~82/footer 81, (20)--(21), and PDF p.~85/footer 84,
(22), describe the $H^1$ and completed-unit ambiguity. These are the
actual geometric objects on the cited sheets. Their link to LM6
A.44--A.45 is supplied by ORDER item 31; no theorem about an irregular
prime or $K_8(\mathbb Z)$ is printed on these pages.

Here is a precise integral consequence, independent of A and M. Let $S$
be a scheme on which $\ell$ is invertible, put $P=\operatorname{Pic}(S)$,
and let $\operatorname{Br}'(S)=H^2_{\text{ét}}(S,\mathbb G_m)$.
Write $\widehat P_\ell=\varprojlim_nP/\ell^nP$ and
$T_\ell A=\varprojlim_n A[\ell^n]$, with transition maps multiplication
by $\ell$. There are exact sequences
\begin{align*}
 0&\longrightarrow\widehat P_\ell
   \longrightarrow\varprojlim_nH^2_{\text{ét}}(S,\mu_{\ell^n})
   \longrightarrow T_\ell\operatorname{Br}'(S)\longrightarrow0,\\
 0&\longrightarrow\Gamma(S,\mathcal O_S^*)^{\wedge}_\ell
   \longrightarrow\varprojlim_nH^1_{\text{ét}}(S,\mu_{\ell^n})
   \longrightarrow T_\ell P\longrightarrow0.
\end{align*}
We deliberately use inverse-limit cohomology: identifying it with a
chosen definition of continuous cohomology requires its own
$\varprojlim{}^1$ check. If $P$ is finitely generated, the second right
term is zero. A line bundle $\mathcal L$ has zero compatible integral
$\ell$-adic Chern class if and only if its class in $P$ is torsion of
order prime to $\ell$.

\begin{proof}
The Kummer exact sequence gives
\[
 0\longrightarrow P/\ell^nP
 \longrightarrow H^2(S,\mu_{\ell^n})
 \longrightarrow\operatorname{Br}'(S)[\ell^n]\longrightarrow0,
\]
and the analogous degree-one sequence with units modulo $\ell^n$th
powers and $P[\ell^n]$. The exact source is Stacks Project,
Lemma 59.28.1, tag 03PK, and its immediately following cohomology
sequence. The transition maps on the left terms are surjective.
Consequently a compatible sequence on the right lifts compatibly in
the middle: after choosing a lift at level $n$, correct any lift at
level $n+1$ by a lift of its discrepancy in the left term. This recursive
construction proves right-surjectivity after taking inverse limits;
left-exactness is immediate. It proves both displayed sequences.

Now write the finitely generated group $P$ as
$\mathbb Z^a\oplus P[\ell^\infty]\oplus P_{\ell'}$, with finite torsion
summands. Its $\ell$-adic completion is
$\mathbb Z_\ell^a\oplus P[\ell^\infty]$, and the kernel of completion
is exactly $P_{\ell'}$. Also no nonzero compatible infinitely
$\ell$-divisible torsion sequence can occur in a finitely generated
group, so $T_\ell P=0$. The injection in degree two then proves the
Chern-class criterion.
\end{proof}

\paragraph{AG-KUMMER-2: a fully explicit integral/rational separation.}
Let $K=\mathbb Q(\sqrt{-35})$, $\omega=(1+\sqrt{-35})/2$, and
$S=\operatorname{Spec}\mathbb Z[\omega,1/2]$. The invertible ideal
$\mathfrak p=(3,\omega)$ determines a line bundle $\mathcal L$ whose
compatible integral $2$-adic Chern class is nonzero, although its image
after tensoring with $\mathbb Q_2$ is zero.

\begin{proof}
The ring of integers is $\mathbb Z[\omega]$ and its discriminant is
$-35$. The minimal polynomial $X^2-X+9$ is irreducible modulo $2$,
so $(2)$ is the only prime over $2$, and is principal. Modulo $3$ it is
$X(X-1)$; thus the primes over $3$ are
$\mathfrak p=(3,\omega)$ and $\overline{\mathfrak p}=(3,\omega-1)$.
Since $N(\omega)=9$ and $\omega$ is a unit modulo
$\overline{\mathfrak p}$, unique ideal factorization gives
$\mathfrak p^2=(\omega)$. The ideal $\mathfrak p$ is not principal:
an integral generator would have norm $3$, whereas
\[
 N(a+b\omega)=a^2+ab+9b^2=(a+b/2)^2+35b^2/4
\]
cannot be $3$ for integers $a,b$.

Minkowski's bound is $2\sqrt{35}/\pi<4$, so every ideal class has a
representative of norm $1$, $2$, or $3$. There is no ideal of norm $2$;
the two of norm $3$ have the same class of order two. Hence the class
group is $\mathbb Z/2$. The precise classical bound is Milne,
\emph{Algebraic Number Theory}, v3.08 (19 July 2020), Theorem 4.3,
PDF p.~72. Inverting $2$ quotients the class group by the classes of
the primes over $2$, which here are principal, so
$\operatorname{Pic}(S)=\mathbb Z/2$ as well. To see this last assertion
directly, fractional ideals of the localization are obtained by forgetting
the valuation at $(2)$; a relation introduced in the localization differs
from a principal ideal only by a power of that principal prime.

AG-KUMMER now injects this nonzero class into the inverse limit of
$H^2(S,\mu_{2^n})$. It has order two there, so rationalization kills it.
The total space of $\mathcal L$ has a zero section with normal bundle
$\mathcal L$, making the example a literal normal-bundle obstruction of
the kind appearing in cote 149. No claim identifying this particular
scheme's group-cohomology extension class is needed.
\end{proof}

This proves a concrete limit of rational information: a rationally zero
normal Chern class need not be integrally zero. A supplies Dirichlet GRH
for the quadratic field here, but changes none of this calculation.
The finite-level Kummer test is therefore a genuine extra gate for an
integral inertia-splitting argument. It does not refute the original
LM6 statements, whose comparison maps must still be supplied.

\subsection{AG-C34-TYPE and AG-C34-REPAIR: resolving the commutator notation}
\label{ag:c34repair}

The official cote 34 PDF has a cover: PDF p.~53 is manuscript/footer
p.~52, and PDF p.~64 is p.~63. Theorem 1.1 on the former page fixes a
field $k$ finitely generated over its prime field, a normal geometrically
connected $k$-scheme $X$ locally of finite type, a prime
$\ell\ne\operatorname{char}k$, and an open subgroup
$\pi\subseteq\pi_1(X)$. It defines $\pi_0$ to be the image in
$G_k$ and $I=\pi\cap\pi_1(X_{\bar k})$, in
\[
 1\longrightarrow I\longrightarrow\pi\longrightarrow\pi_0
 \longrightarrow1.
\]
Its $I(\ell)$ is the maximal abelian pro-$\ell$ quotient of $I$.
The theorem states that $I(\ell)$ is finitely generated over
$\mathbb Z_\ell$ and that its $\pi_0$-coinvariants are finite.
The proof was located and read in Sections 2--3, PDF pp.~70--78
(manuscript pp.~69--77), rather than guessed from the intervening
applications. Section 2.4, PDF p.~72, makes the dual $H^1$-invariant
reduction; Section 2.5, PDF pp.~73--74, uses a curve compactification,
its cyclotomic boundary, and the finite-field curve theorem after
spreading out. Sections 3.1--3.3, PDF pp.~75--78, reduce normal $X$ to
a generic hyperplane curve using surjectivity on fundamental groups.
These are classical geometric inputs; A and M do not enter them.

In the crossed-out Remark 1.17 on PDF p.~64, $H$ is a quotient of $I$
whose kernel is $\pi$-stable, and the readable final hypothesis is that
$H$ is a solvable pro-$\ell$ group. The preceding proposed
$\ell$-adic Lie-group language is itself struck out. The remark writes
\[
 1\longrightarrow H\longrightarrow\pi'\longrightarrow\pi_0
 \longrightarrow1
\]
and places the closure of $[\pi_0,H]$ as an open subgroup of $\pi'$.
This is a tentative, struck manuscript assertion, not a theorem.

\paragraph{AG-C34-TYPE (the literal ambient-group obstruction).}
For an infinite $\pi_0$, every commutator of a lift of an element of
$\pi_0$ with an element of $H$ lies in $H$, and its closure cannot be
open in $\pi'$. Indeed the quotient by $H$ is the infinite group
$\pi_0$, whereas a subgroup open in a profinite group has finite index.
For a concrete example satisfying the manuscript's geometric hypotheses,
take $X=\mathbb G_m/\mathbb F_q$ and $\ell\nmid q$. Its geometric
pro-$\ell$ quotient $H=\mathbb Z_\ell(1)$ is abelian and hence
solvable; the rational point $1$ gives a section. Arithmetic Frobenius
acts as multiplication by $q$, so the action commutators generate
$(q-1)\mathbb Z_\ell(1)$. This is open in $H$ but is not open in
$\pi'$, whose quotient is $G_{\mathbb F_q}\simeq\widehat{\mathbb Z}$.
Thus the literal ambient group in the inherited Q03 cannot be retained.
For nonabelian $H$ there is a second issue: $\pi_0$ acts canonically
only by outer automorphisms unless a section or lifts are specified.

\paragraph{AG-C34-REPAIR (the intrinsic corrected statement).}
Let $H$ be \emph{any} pro-$\ell$ quotient of $I$ with $\pi$-stable
kernel, and let $P$ be the resulting quotient extension of $\pi_0$ by
$H$. Define $K$ to be the closed subgroup generated by all
$[p,h]$, $p\in P$, $h\in H$. Then $K$ is open in $H$, and
\[
 H/K\simeq (H^{\mathrm{ab}})_{\pi_0},
 \qquad [H:K]\le |I(\ell)_{\pi_0}|<\infty.       \tag{AG-C34.1}
\]
In particular the solvability and Lie-group hypotheses are unnecessary
for this precise corrected conclusion.

\begin{proof}
Because $H$ is normal in $P$, conjugating a commutator $[p,h]$ by
$a\in P$ gives $[apa^{-1},aha^{-1}]$. Hence $K$ is normal in $P$.
It contains the closed commutator subgroup of $H$, so $H/K$ is abelian;
all conjugation by $P$ is trivial on it. Conjugation on
$H^{\mathrm{ab}}$ factors canonically through $\pi_0$, irrespective
of any section. Killing its action differences is exactly the additional
quotient by the images of the commutators $[p,h]$. This proves the
isomorphism in (AG-C34.1).

The surjection $I\twoheadrightarrow H$ induces an equivariant
surjection $I(\ell)\twoheadrightarrow H^{\mathrm{ab}}$ by the
universal property of the maximal abelian pro-$\ell$ quotient.
Taking coinvariants remains surjective: it is simply quotienting both
compact abelian groups by the closed subgroups generated by their
action differences. Theorem 1.1 therefore supplies a surjection from
the finite group $I(\ell)_{\pi_0}$ to $H/K$. This proves the bound.
Since $K$ is closed, finite index in the compact group $H$ implies
openness (the finite Hausdorff quotient is discrete).
\end{proof}

This repairs the archived Q03 as a theorem already implied by its
stated abelian predecessor. It does not identify the unqualified
symbol $[\pi_0,H]$ in every possible intended meaning. A smaller
commutator subgroup formed from a chosen section, or openness in a
different ambient group, needs a separately stated question. In
particular the repair is not described as an automatic proof of all
interpretations of the crossed-out remark.

\paragraph{AG-C34-INDEX (an explicit Frobenius refinement).}
More generally, let $M=H^{\mathrm{ab}}$ be a finitely generated
$\mathbb Z_\ell$-module with continuous $\Gamma$-action, let $T$ be
its finite torsion subgroup, and let $\Lambda=M/T$. Suppose some
$F\in\Gamma$ has no eigenvalue $1$ on
$\Lambda\otimes\mathbb Q_\ell$. Then
\[
 |\operatorname{coker}(F-1\mid M)|
   =|T^F|\ell^{v_\ell\det(1-F\mid\Lambda\otimes\mathbb Q_\ell)},
 \qquad
 |M_\Gamma|\le
 |T^F|\ell^{v_\ell\det(1-F\mid\Lambda\otimes\mathbb Q_\ell)}.
                                                               \tag{AG-C34.2}
\]
Equality in the second assertion holds if $\Gamma$ is topologically
generated by $F$.

\begin{proof}
The assumption gives $\Lambda^F=0$. Apply the snake lemma to
$0\to T\to M\to\Lambda\to0$ with vertical maps $F-1$.
It yields
\[
 0\longrightarrow T/(F-1)T\longrightarrow M/(F-1)M
 \longrightarrow\Lambda/(F-1)\Lambda\longrightarrow0.
\]
For an endomorphism of the finite group $T$, kernel and cokernel have
the same cardinality. Smith normal form over $\mathbb Z_\ell$ gives
the order of the last group as the displayed power of $\ell$.
This proves the equality. The continuous image $(F-1)M$ is compact
and hence closed; the full coinvariant quotient is a quotient of it.
If the powers of $F$ are dense in $\Gamma$, every action difference
is a limit of differences of powers, all of which belong to
$(F-1)M$, so the two quotients agree.
\end{proof}

For a smooth proper geometrically connected curve $C/\mathbb F_q$,
let $H$ be its full geometric pro-$\ell$ fundamental group and $J$ its
Jacobian. Then $H^{\mathrm{ab}}=T_\ell J$, compatibly with Galois
action, and the corrected commutator quotient is
\[
 H/\overline{[P,H]}
       \simeq T_\ell J/(F-1)T_\ell J
       \simeq J(\mathbb F_q)[\ell^\infty].             \tag{AG-C34.3}
\]
For the fundamental-group comparison, an exact modern primary locator
is J.~S.~Milne, \emph{Jacobian Varieties}, Proposition 9.1,
pp.~30--32, together with \emph{Abelian Varieties}, Remarks 15.3
and 15.5, pp.~28--29. To check the last isomorphism directly, use
$0\to T_\ell J\to V_\ell J\to J(\overline{\mathbb F}_q)
[\ell^\infty]\to0$. The fixed subgroup of $T_\ell J$ is zero:
$J(\mathbb F_q)$ is finite, so a compatible indefinitely divisible
sequence of its $\ell$-power torsion points vanishes. Thus $F-1$ is
invertible on $V_\ell J$, and the snake lemma gives (AG-C34.3).
Consequently the index is exactly $\ell^{v_\ell\#J(\mathbb F_q)}$.
This finite-field calculation even uses less than the full Weil bound.
In higher-dimensional smooth proper situations purity supplies the
corresponding no-eigenvalue-$1$ test, but finite-field purity is already
classical and no new global GRH input is involved.

\subsection{AG-C34-SECTION: the solvable quotient with a chosen section}
\label{ag:c34section}

There is a second precise interpretation of the struck Remark 1.17 that
can now be settled. Keep the exact geometric hypotheses and notation of
AG-C34-REPAIR. Suppose the quotient extension
$1\to H\to P\to\pi_0\to1$ has a specified continuous section, so
$P=H\rtimes\pi_0$. Assume that $H$ is solvable of finite derived
length, using closed derived subgroups. This is stronger than merely
being an inverse limit of finite solvable groups; every pro-$\ell$
group already satisfies that latter condition. For the action $\alpha$
given by the section put
\[
 N=\overline{\langle h^{-1}\alpha_g(h):h\in H,
                                      \ g\in\pi_0\rangle}.
\]
\textbf{Statement.} The subgroup $N$ is normal in $P$ and open in
$H$. Thus the chosen-section action commutators have finite index in
the geometric group. They need not equal the intrinsic subgroup
$\overline{[P,H]}$ of AG-C34-REPAIR.

\begin{proof}
For a fixed automorphism $\alpha$, write
$c_\alpha(h)=h^{-1}\alpha(h)$. The identity
\[
 c_\alpha(hx)=x^{-1}c_\alpha(h)x\,c_\alpha(x)
\]
shows that the subgroup generated by all the $c_\alpha(h)$ is normal
in $H$. Applying $\alpha_g$ to $c_{\alpha_f}(h)$ gives
$c_{\alpha_{gfg^{-1}}}(\alpha_g(h))$, so it is also stable under
$\pi_0$. Its closure $N$ is consequently normal in $P$. Set
$Q=H/N$. The induced action of $\pi_0$ on $Q$ is trivial by
construction, whence $P/N=Q\times\pi_0$.

We claim that every open subgroup $U$ of $Q$ has finite abelianization.
Take the inverse image of $U\times\pi_0$ first in $P$ and then in
the original open arithmetic subgroup $\pi$. This gives an open
subgroup $\pi_U$ of $\pi_1(X)$ with the same arithmetic image
$\pi_0$. Its geometric kernel $I_U$ maps onto $U$. Conjugation on
$U^{\mathrm{ab}}$ is trivial: the $U$-part acts by inner
automorphisms and the $\pi_0$-part acts trivially. Therefore
\[
 I_U(\ell)_{\pi_0}\twoheadrightarrow U^{\mathrm{ab}}.
\]
The left group is finite by Theorem 1.1 applied to the open subgroup
$\pi_U$, proving the claim. This use of the theorem on \emph{all
the relevant finite covers}, rather than on $H$ alone, is essential.

Let $Q^{(0)}=Q$ and
$Q^{(i+1)}=\overline{[Q^{(i)},Q^{(i)}]}$. Starting with the open
subgroup $Q$, the claim implies inductively that $Q^{(i+1)}$ has
finite index in the open subgroup $Q^{(i)}$, and hence is open in
$Q$. Since $Q$ is a quotient of the solvable group $H$, this finite
derived series ends at $1$. Thus $1$ is open in $Q$, so the compact
group $Q$ is finite. Equivalently $N$ is open in $H$.

For the non-equality assertion, take a finite nonabelian $\ell$-group
$H$ and let $\pi_0$ act trivially. Then $N=1$ whereas
$\overline{[P,H]}=[H,H]\ne1$. Both are open in $H$. This last
example concerns the subgroup identity, and asserts no realization of
that abstract example by a particular arithmetic variety.
\end{proof}

This is additional movement on BRIDGES Q03: under the explicitly stated
section hypothesis, its solvability condition is sufficient for the
action-commutator conclusion in the corrected ambient group $H$.
No section is inferred from purity or from A/M. Without a section the
intrinsic assertion AG-C34-REPAIR remains the canonical formulation.

\subsection{AG-C34-NONSOLV: why compact analytic groups alone do not suffice}
\label{ag:c34nonsolv}

The finite-derived-length condition just used cannot be replaced, at
the level of this group argument, by finite generation or compact
$\ell$-adic analyticity. The following counterexample is deliberately
abstract; no geometric realization is asserted.

Fix an odd prime $\ell$ and put
\[
 H_n=\ker\bigl(\mathrm{SL}_2(\mathbb Z_\ell)
            \longrightarrow\mathrm{SL}_2(\mathbb Z/\ell^n)\bigr),
 \qquad H=H_1.
\]
Then $H$ is an infinite, topologically finitely generated compact
pro-$\ell$ group, and every open subgroup of $H$ has finite
abelianization. Nevertheless, for the trivial action of any profinite
group $\Gamma$, the action-commutator subgroup is $1$, which is not
open in $H$. In particular the absence of eigenvalue $1$ on
$H^{\mathrm{ab}}\otimes\mathbb Q_\ell$ does not repair that
inference: here this vector space is zero.

\begin{proof}
For each $n\ge1$, the map
$1+\ell^n A\mapsto A\bmod\ell$ identifies
$H_n/H_{n+1}$ with the additive group
$\mathfrak{sl}_2(\mathbb F_\ell)$. The determinant equation gives
trace zero. Surjectivity follows from the upper and lower unipotent
matrices and the diagonal matrix
$\operatorname{diag}(1+\ell^n,(1+\ell^n)^{-1})$, whose leading terms
are the basis
\[
 E=\begin{pmatrix}0&1\\0&0\end{pmatrix},\quad
 F=\begin{pmatrix}0&0\\1&0\end{pmatrix},\quad
 T=\begin{pmatrix}1&0\\0&-1\end{pmatrix}.
\]
Thus $H$ is an inverse limit of finite $\ell$-groups and is compact;
the upper unipotent matrices with entry in $\ell\mathbb Z_\ell$
already show that it is infinite.
It is explicitly a three-dimensional $\ell$-adic analytic group:
write its matrices as
$\bigl(\begin{smallmatrix}a&b\\c&(1+bc)/a\end{smallmatrix}\bigr)$
with $a\in1+\ell\mathbb Z_\ell$ and
$b,c\in\ell\mathbb Z_\ell$. These three parameters are analytic
coordinates, and multiplication and inversion are analytic in them.

The three displayed lifts at level $1$ topologically generate $H$.
Indeed, for $n\ge1$ and odd $\ell$,
\[
 (1+\ell^n A)^\ell\equiv1+\ell^{n+1}A
                       \pmod{\ell^{n+2}}.
\]
Their successive $\ell$-power iterates therefore give all three basis
vectors in every $H_n/H_{n+1}$. Any element of $H$ can be matched
successively modulo $H_2,H_3,\ldots$ by products of elements in the
closed group they generate. The products converge to the given element,
so that closed group is all of $H$.

For $n,m\ge1$, direct matrix multiplication gives
\[
 [1+\ell^n A,1+\ell^m B]
       \equiv1+\ell^{n+m}(AB-BA)
                      \pmod{\ell^{n+m+1}}.
\]
The Lie brackets $[E,F]=T$, $[T,E]=2E$, and $[T,F]=-2F$
span $\mathfrak{sl}_2(\mathbb F_\ell)$ because $\ell$ is odd.
Let $C_n=\overline{[H_n,H_n]}$. For every $k\ge2n$, commutators
of $H_n$ with $H_{k-n}\subseteq H_n$ therefore show that
$C_n\cap H_k$ surjects onto $H_k/H_{k+1}$. Successively lifting a
given element of $H_{2n}$ through these quotients and taking the
convergent product proves
\[
 H_{2n}\subseteq C_n.
\]
Every open subgroup $U\subseteq H$ contains some $H_n$. Its closed
commutator subgroup contains $C_n$, and hence $H_{2n}$, so
$U^{\mathrm{ab}}$ is finite. Taking $U=H$ also proves
$H^{\mathrm{ab}}\otimes\mathbb Q_\ell=0$.
For a trivial $\Gamma$-action, each action commutator is $1$,
which has infinite index in $H$. This proves the claimed failure.
\end{proof}

The example pinpoints what AG-C34-SECTION uses beyond abelian
coinvariant control: its induction must terminate. It does not challenge
the intrinsic commutator theorem, since $\overline{[H,H]}$ is open
here, nor any additional geometric constraints that could exclude this
trivial-action example in a proposed realization.

\subsection{Cote 149: the fundamental-group class and the section obstruction}
\label{ag:c149-comparison}

The comparison omitted by a bare Kummer calculation is available under
precise homotopy hypotheses. Borne--Emsalem--Stix,
\emph{Lifting Galois sections along torsors}, Math. Nachr.~288 (2015),
1939--1951, arXiv:1301.4429v2, Proposition 3.1 and Proposition 3.4
(PDF pp.~5--6), and Proposition 4.4 (PDF pp.~8--9), supply the
following classical statement. The complete proofs and the continuous
cohomology argument in Proposition 4.2 (PDF p.~7) have been read.

\paragraph{AG-C149-COMP (the exact comparison).}
Let $S$ be a geometrically connected, geometrically unibranch variety
over a characteristic-zero field $k$, with
$\pi^{\text{ét}}_2(S,\bar s)=0$. Let $\mathcal L$ be a line
bundle and $\mathcal L^\times$ its complement of the zero section.
For compatible base points there is an exact sequence
\[
 1\longrightarrow\widehat{\mathbb Z}(1)
 \longrightarrow\pi_1(\mathcal L^\times)
 \longrightarrow\pi_1(S)\longrightarrow1.
\]
The conjugation action on its kernel factors through $G_k$ and is
cyclotomic. Quotienting the kernel by its prime-to-$\ell$ factors gives
an extension
\begin{equation}\label{eq:ag-c149-extension}
 1\longrightarrow\mathbb Z_\ell(1)
 \longrightarrow\mathcal E_\ell(\mathcal L)
 \longrightarrow\pi_1(S)\longrightarrow1.
\end{equation}
The comparison isomorphism
\[
 H^2_{\mathrm{cont}}(\pi_1(S),\mathbb Z_\ell(1))
 \simeq H^2_{\mathrm{cont}}(S,\mathbb Z_\ell(1))
\]
sends its extension class $e_\ell(\mathcal L)$ to the arithmetic
Kummer class $c_{1,\ell}(\mathcal L)$. Continuous cohomology here is
the cohomology of the inverse system $(\mu_{\ell^n})$; it is not
silently identified with an inverse limit in arbitrary situations.

The source proves exactness from the homotopy sequence of the torus
torsor and proves the class comparison by compatible transgressions:
the identity of the Tate module maps to the negative of each of the
two classes. The hypothesis on $\pi_2$ is substantive: over an
algebraically closed field, the torsor
$\mathbb A^{d+1}\setminus\{0\}\to\mathbb P^d$, $d\ge1$, has trivial
total-space fundamental group, so its fibre inertia does not inject
(source Remark 3.2, PDF p.~5). Smooth geometrically connected curves
other than the geometric projective line are algebraic $K(\pi,1)$
spaces (source Section 3.1, PDF p.~4).

Geometrically, $\mathcal L$ is the normal bundle of the zero section in
$\operatorname{Tot}(\mathcal L)$, so this is an actual punctured
normal-bundle extension. It instantiates the $K(\pi,1)$ comparison
explicitly contemplated on cote 149, original PDF p.~80, manuscript
p.~79, equations (13)--(17), followed by the normal-bundle discussion
on PDF p.~81, manuscript p.~80. For the torsion-line-bundle hyperbolic family below,
AG-C149-NORMAL supplies the corresponding inertia-normalizer map.
A different family still requires its own comparison; none of these
statements identifies the global tripod kernel in LM6 A.44--A.45.

\begin{proposition}[AG-C149-LIFT: torsion classes vanish on every section]
Under the hypotheses of AG-C149-COMP, a continuous section
$s:G_k\to\pi_1(S)$ of the arithmetic projection lifts to
$\mathcal E_\ell(\mathcal L)$ if and only if
$s^*e_\ell(\mathcal L)=0$. If $\mathcal L$ is torsion in
$\operatorname{Pic}(S)$, every such section lifts, although the
extension \eqref{eq:ag-c149-extension} need not split over
$\pi_1(S)$. When lifts exist, their classes modulo conjugation by
$\mathbb Z_\ell(1)$ form a torsor under
\[
 H^1_{\mathrm{cont}}(k,\mathbb Z_\ell(1))
   \simeq\varprojlim_n k^\times/(k^\times)^{\ell^n}.
\]
\end{proposition}
\begin{proof}
Pulling back \eqref{eq:ag-c149-extension} by $s$ gives an extension of
$G_k$ with class $s^*e_\ell(\mathcal L)$. A lift is precisely a
continuous splitting of this pullback, so the class criterion follows
from extension classification, equivalently source Lemma 4.5
(PDF p.~9) and its preceding comparison.

We verify the continuous-cohomology issue directly for the field.
Hilbert 90 and the Kummer sequence give
\[
 H^1(k,\mu_{\ell^n})=k^\times/(k^\times)^{\ell^n},\qquad
 H^2(k,\mu_{\ell^n})=\operatorname{Br}(k)[\ell^n].
\]
The maps on the first groups are the surjective quotient maps, and
those on the second groups are multiplication by $\ell$. Thus
$\varprojlim^1 H^1(k,\mu_{\ell^n})=0$, and the defining exact
sequence for continuous cohomology gives
\[
 H^2_{\mathrm{cont}}(k,\mathbb Z_\ell(1))
   \simeq T_\ell\operatorname{Br}(k).
\]
For any abelian group $A$, its Tate module
$T_\ell A=\varprojlim(A[\ell^n],\ell)$ is torsion-free over
$\mathbb Z_\ell$: if $\ell t=0$, then compatibility gives
$t_n=\ell t_{n+1}=0$ for every $n$, and induction excludes all
$\ell$-power torsion. Multiplication by an integer prime to $\ell$
is invertible. If $m[\mathcal L]=0$, additivity of the Kummer class
therefore gives
$m s^*e_\ell(\mathcal L)=0$, whence $s^*e_\ell(\mathcal L)=0$.
This is the $\mathbb G_m$, $\ell$-primary specialization of the
source's Proposition 4.2 and Lemma 4.6.

Finally, once one splitting is fixed, the difference of a second
splitting from it is a continuous $1$-cocycle. Conjugation by the
kernel changes that cocycle by a coboundary. This proves the torsor
claim. The finite groups $H^0(k,\mu_{\ell^n})$ form a
Mittag--Leffler inverse system; hence their $\varprojlim^1$ vanishes
and the degree-one continuous-cohomology sequence gives the displayed
Kummer identification.
\end{proof}

\begin{proposition}[AG-C149-ELLIPTIC: non-splitting does not obstruct all sections]
Let $E/\mathbb Q$ be the elliptic curve
\[
 y^2=x^3-x,
 \qquad O=\text{the point at infinity},\qquad P=(0,0),
\]
and let $\mathcal L=\mathcal O_E(P-O)$.
The class of the extension $\mathcal E_2(\mathcal L)\to\pi_1(E)$
is nonzero of order $2$, so this extension does not split. Its image
with $\mathbb Q_2(1)$ coefficients is zero. Nevertheless, every
Galois section $G_{\mathbb Q}\to\pi_1(E)$ lifts through this
extension. Thus vanishing of the obstruction after pullback by
every Galois section does not force the global integral class to
vanish.
\end{proposition}
\begin{proof}
The point $P$ is a nonidentity point of order $2$, since inversion on
this Weierstrass curve sends $(x,y)$ to $(x,-y)$. The standard
identification $E(\mathbb Q)=\operatorname{Pic}^0(E)$ sends $P$ to
$[\mathcal L]$, so $\mathcal L$ has order $2$.

Its discriminant is $64$, so the displayed integral model has good
reduction at $3$. Direct enumeration gives
\[
 E(\mathbb F_3)=\{O,(0,0),(1,0),(2,0)\}
                 \simeq(\mathbb Z/2\mathbb Z)^2.
\]
Properness and good reduction give a group homomorphism
$E(\mathbb Q)\to E(\mathbb F_3)$. If $P=2Q$ for a rational point
$Q$, reduction would yield $\bar P=2\bar Q=0$, whereas
$\bar P=(0,0)\ne O$. Therefore $P\notin2E(\mathbb Q)$.
If $\mathcal L=\mathcal M^{\otimes2}$ in $\operatorname{Pic}(E)$,
degrees give $\deg\mathcal M=0$, and the same Picard identification
would express $P$ as twice a rational point. Thus
$[\mathcal L]\ne0$ in $\operatorname{Pic}(E)/2$.

The Kummer injection
\[
 \operatorname{Pic}(E)/2\hookrightarrow H^2(E,\mu_2)
\]
shows that the reduction of $c_{1,2}(\mathcal L)$ modulo $2$ is
nonzero. Consequently the continuous integral class is nonzero;
its order is exactly $2$ by additivity and $\mathcal L^{\otimes2}
\simeq\mathcal O_E$. The curve has genus $1$ and is an algebraic
$K(\pi,1)$, so AG-C149-COMP identifies this with the nonzero
extension class. A torsion class maps to zero in cohomology with
$\mathbb Q_2(1)$ coefficients. Finally AG-C149-LIFT applies to
every Galois section, and such sections do exist because
$O\in E(\mathbb Q)$. This proves all the claims.
\end{proof}

This supplies the group/cohomology comparison missing from the purely
arithmetic example AG-KUMMER-2, under explicit hypotheses and for a
literal normal bundle. It also separates three questions on the cote:
nonvanishing of a global integral extension class, its rational image,
and the obstruction for a specified section. None of these distinctions
uses A or M, and none identifies an irregular-prime class-group
component or proves the remaining LM6 tripod-kernel assertion.

\subsection{An actual hyperbolic-family cusp normalizer}
\label{ag:c149-normalizer}

We now make the remaining normalizer map explicit for a family
constructed from a torsion line bundle. The only local normalizer
input is Stix, \emph{On cuspidal sections of algebraic fundamental
groups}, author manuscript dated 24 February 2012, Section 3,
Lemma 14 (PDF p.~7) and Proposition 16(1)--(2) (PDF p.~8).
For a characteristic-zero hyperbolic curve, two distinct prolonged
cusps have inertia groups with trivial intersection, and the
arithmetic normalizer of a cusp inertia group is its decomposition
group. The complete proofs of these statements have been read.
The product input is SGA~1, Expos\'e X, Corollary 1.7, recomposed
pp.~203--204 (arXiv PDF pp.~219--220), whose proof and Corollary 1.8
have also been read. Its proper-factor hypothesis will be retained.

\begin{proposition}[AG-C149-NORMAL: a cyclic-cover normalizer comparison]
Let $k$ have characteristic zero, let $S/k$ be smooth, proper and
geometrically connected, and fix a line bundle $\mathcal L$ with
an isomorphism $\tau:\mathcal L^{\otimes n}\simeq\mathcal O_S$,
where $n\ge2$. Set
\[
 B=\mathbb P^1_k\setminus\{0,1,\infty\},\qquad
 U=\{v\in\mathcal L^\times:\tau(v^{\otimes n})\ne1\}.
\]
Then $U\to S$ is a smooth family of hyperbolic curves of type
$(0,n+2)$. Compactify the line bundle by adjoining its infinity
section, so that the boundary consists of the zero section, the
infinity section, and the degree-$n$ finite \'etale divisor
$\tau(v^{\otimes n})=1$. Choose compatible geometric base points,
and let $I\simeq\widehat{\mathbb Z}(1)$ be the zero-cusp inertia
subgroup of a geometric fibre, viewed inside $\pi_1(U)$.
The open immersion $U\to\mathcal L^\times$ restricts to an
isomorphism
\begin{equation}\label{eq:ag-cusp-normalizer}
 N_{\pi_1(U)}(I)\ \xrightarrow{\ \sim\ }
 \pi_1(\mathcal L^\times)
\end{equation}
over $\pi_1(S)$, with the usual compatibility up to choices of
paths. In particular, it gives an exact sequence
\[
 1\longrightarrow I\longrightarrow N_{\pi_1(U)}(I)
 \longrightarrow\pi_1(S)\longrightarrow1,
\]
whose action on $I$ is cyclotomic. No $\pi_2(S)=0$ hypothesis is
needed for this explicit torsion-line-bundle construction.
\end{proposition}
\begin{proof}
The maps $v\mapsto\tau(v^{\otimes n})$ give finite \'etale
$\mu_n$-torsors
\[
 U\longrightarrow S\times B,\qquad
 \mathcal L^\times\longrightarrow S\times\mathbb G_m.
\]
This follows locally by writing the map as $z\mapsto u z^n$ with
$u$ a unit: its derivative is invertible when $z\ne0$.
After the finite \'etale cover
$T=\{e\in\mathcal L^\times:\tau(e^{\otimes n})=1\}\to S$,
the family is
$T\times(\mathbb P^1\setminus(\{0,\infty\}\cup\mu_n))$.
Thus it has the stated smooth hyperbolic fibres, and $U$ and
$\mathcal L^\times$ are geometrically connected.

Write $\Pi_X=\pi_1(X)$ and $G=G_k$.
Proper K\"unneth over $\bar k$, followed by the geometric/arithmetic
exact sequences, gives
\[
 \Pi_{S\times B}=\Pi_S\times_G\Pi_B,\qquad
 \Pi_{S\times\mathbb G_m}=\Pi_S\times_G\Pi_{\mathbb G_m}.
\]
Indeed the maps induced by projections are isomorphisms on their
geometric kernels and on the quotient $G$, and hence are
isomorphisms. This verifies the arithmetic product assertion even
when $S$ has no rational point.

Use additive notation for $\mu_n$. Choose a point in the geometric
fibre of $T$, so that $T$ is represented by a continuous crossed
cocycle $a:\Pi_S\to\mu_n$. Let
$b:\Pi_B\to\mu_n$ represent the Kummer torsor $z^n=t$, where
$t$ is the coordinate on $B$. The two actions on $\mu_n$ factor
through the same cyclotomic action of $G$. The first displayed
cover has cocycle $a+b$: locally write $v=e z$ with
$\tau(e^{\otimes n})=1$ and $z^n=t$. Consequently its fundamental
group is the stabilizer
\begin{equation}\label{eq:ag-cyclic-stabilizer}
 H=\{(\sigma,\beta)\in\Pi_S\times_G\Pi_B:
                 a(\sigma)+b(\beta)=0\}.
\end{equation}
The zero set of a crossed cocycle is a subgroup. Unless the
$\mu_n$ action is trivial, it need not be a normal subgroup; no
ordinary homomorphism-kernel assertion is being made here.
The same description gives a subgroup $H'$ for
$\Pi_{\mathcal L^\times}$, using the Kummer cocycle on
$\Pi_{\mathbb G_m}$.

Choose the zero cusp of $B$, its inertia $I_0$ and its
decomposition group $D_0$. In the product group the relative
inertia is $\{1\}\times I_0$. Since $t$ has ramification index
$n$ in the local parameter $v$ upstairs, the inertia in $H$ is
\[
 I=\{1\}\times nI_0.
\]
This also verifies its injection: multiplication by $n$ on the
free profinite procyclic group $I_0$ is injective. Equivalently,
the Kummer cocycle restricted to $I_0$ is its reduction modulo
$n$, so its intersection with $H$ is exactly $nI_0$.

If $\beta\in\Pi_B$ normalizes $nI_0$, then
$I_0\cap\beta I_0\beta^{-1}$ contains the nontrivial subgroup
$nI_0$. Stix's Proposition 16(1) forces the two prolonged cusps
to coincide, so $\beta\in D_0$. Conversely, $D_0$ normalizes
$I_0$ and its characteristic subgroup $nI_0$. Therefore
\[
 N_H(I)=H\cap(\Pi_S\times_G D_0).
\]
The filling map $B\to\mathbb G_m$ identifies
$D_0\xrightarrow{\sim}\Pi_{\mathbb G_m}$.
To see this without a choice of splitting, compare their exact
sequences with quotient $G$: on the quotient the map is the
identity, and on geometric inertia it is an isomorphism, since
both groups are $\widehat{\mathbb Z}(1)$ with the same local
parameter $t$. The injections of the punctured local group for
$B$ and for $\mathbb G_m$ are supplied by Stix's Lemma 14,
whose hypothesis is $\chi\le0$ and whose proof explicitly includes
$\mathbb G_m$. For the latter curve this also follows directly
from the Kummer covers $t\mapsto t^a$: geometrically, its entire
fundamental group is this inertia group.

The resulting isomorphism
$\Pi_S\times_GD_0\simeq\Pi_S\times_G\Pi_{\mathbb G_m}$
preserves the Kummer cocycles. It therefore identifies $N_H(I)$
with $H'$, proving \eqref{eq:ag-cusp-normalizer}.
Finally the rational point $1\in\mathbb G_m(k)$ gives
$\Pi_{\mathbb G_m}=\widehat{\mathbb Z}(1)\rtimes G$; its Kummer
cocycle is reduction of the first coordinate modulo $n$.
Thus the projection $H'\to\Pi_S$ is surjective, and its kernel
is $n\widehat{\mathbb Z}(1)\simeq\widehat{\mathbb Z}(1)$.
This proves exactness and the cyclotomic action, with the kernel
identified with $I$ by the same local parameter computation.
\end{proof}

\begin{proposition}[AG-C149-FULLLIFT: every arithmetic section lifts]
Under the hypotheses of AG-C149-NORMAL, every continuous section
$s:G_k\to\Pi_S$ lifts to $N_{\Pi_U}(I)$, and hence to $\Pi_U$.
This is a statement for the full profinite inertia group,
not just one $\ell$-primary quotient.
\end{proposition}
\begin{proof}
Using the last coordinates of the previous proof, the group in
question consists of pairs $(\sigma,z)$ with
\[
 \sigma\in\Pi_S,\quad z\in\widehat{\mathbb Z}(1),\qquad
 a(\sigma)+z\bmod n=0,
\]
and multiplication
$(\sigma,z)(\sigma',z')=(\sigma\sigma',z+\chi(\sigma)z')$,
where $\chi$ is the cyclotomic action. It suffices to lift the
crossed cocycle $-a\circ s:G_k\to\mu_n$ to a continuous
crossed cocycle with values in $\widehat{\mathbb Z}(1)$.

Kummer theory gives $H^1(k,\mu_n)=k^\times/(k^\times)^n$.
Choose $c\in k^\times$ representing $[-a\circ s]$ and a
compatible collection of roots $c^{1/m}\in\bar k$ for all
positive $m$, ordered by divisibility. Such a collection exists
because the finite nonempty root sets have surjective transition
maps. Then
\[
 \kappa(g)=\bigl(g(c^{1/m})/c^{1/m}\bigr)_m
       \in\varprojlim_m\mu_m=\widehat{\mathbb Z}(1)
\]
is a continuous crossed cocycle, each finite coordinate being
continuous. Its reduction modulo $n$ has the desired cohomology
class. If this reduction differs from $-a\circ s$ by the
coboundary $du$ of $u\in\mu_n$, lift $u$ to
$\tilde u\in\widehat{\mathbb Z}(1)$ and replace $\kappa$ by
$\kappa-d\tilde u$. Its reduction is then exactly $-a\circ s$.
The map $g\mapsto(s(g),\kappa(g))$ is the required continuous
group homomorphism and is a lift of $s$. Its projection to $G_k$
is the identity, so its composition with the inclusion in
$\Pi_U$ is an arithmetic section as well.
\end{proof}

\begin{proposition}[AG-C149-CUSP: a nonsplit normalizer with every section liftable]
Take $E/\mathbb Q$ and $\mathcal L=\mathcal O_E((0,0)-O)$
from AG-C149-ELLIPTIC, fix any rational trivialization of
$\mathcal L^{\otimes2}$, and form $U$ as above with $n=2$.
Then $U\to E$ is a family of four-punctured projective lines,
with two boundary sections and a degree-two \'etale boundary
divisor. Its zero-cusp normalizer extension over $\pi_1(E)$
does not split. Its $2$-primary class is nonzero of order $2$
and has zero rational image. Nevertheless every section
$G_{\mathbb Q}\to\pi_1(E)$ lifts through the full cusp normalizer.
\end{proposition}
\begin{proof}
The normal bundle of the zero section in the ruled compactification
is $\mathcal L$, because its neighbourhood is the total space of
$\mathcal L$. AG-C149-NORMAL identifies the actual cusp
normalizer with $\pi_1(\mathcal L^\times)$, over $\pi_1(E)$
and on inertia. The curve $E$ is an algebraic $K(\pi,1)$, so
AG-C149-COMP identifies the $2$-primary extension class with
$c_{1,2}(\mathcal L)$; AG-C149-ELLIPTIC proved that this class
is nonzero of order $2$. Thus its $2$-primary extension is
nonsplit, which also excludes a splitting of the full extension.
The same earlier computation gives its zero rational image.
AG-C149-FULLLIFT proves the assertion for every arithmetic
section, and the section associated with $O\in E(\mathbb Q)$
makes the assertion nonvacuous.
\end{proof}

Thus the normalizer comparison has been supplied for an explicit
hyperbolic family, rather than merely posited for a punctured line
bundle. This advances the actual integral-versus-rational and
section-lifting questions at cote 149, PDF pp.~80--82,
equations (13)--(21). The proof uses neither A nor M. Its scope is
the specified isotrivial family; identifying a different inertia
normalizer still needs its own map, and the global tripod-kernel
assertion LM6 A.45 remains separate.

\subsection{LM6 A.44--A.45: the exact integral generator gate}
\label{ag:lm6}

\paragraph{AG-LM6-REG (the known kernel-field theorem).}
For a prime $p$, let $\Omega_p^*$ be the fixed field of the kernel
of the outer Galois action on
$\pi_1^{\mathrm{pro}\text{-}p}
(\mathbb P^1_{\overline{\mathbb Q}}\setminus\{0,1,\infty\})$.
Let $\Omega_p$ be the maximal pro-$p$ extension of
$\mathbb Q(\mu_p)$ unramified away from $p$. Then
$\Omega_p^*\subseteq\Omega_p$, and equality is classical for odd
regular $p$. An exact primary statement is Hoshi,
\emph{On the Kernels of the Pro-$l$ Outer Galois Representations
Associated to Hyperbolic Curves over Number Fields}, revised March
2014 manuscript, Definition 1.8, Remark 1.8.1, and Theorem 1.9,
PDF p.~11. Its proof cites Sharifi's Theorem 1.1 and Brown's theorem.
This formulation corrects the informal phrase ``the kernel is the
Galois group of $\Omega_p$'': the kernel is
$\operatorname{Gal}(\overline{\mathbb Q}/\Omega_p^*)$.

Brown, \emph{Mixed Tate motives over $\mathbb Z$}, Ann. Math.~175
(2012), 949--976, Theorem 1.1 and the following paragraph, printed
p.~950, prove the motivic-generation assertion formulated as
Conjecture 1 on p.~949. Its Deligne--Ihara consequence concerns the
graded Lie algebra after tensoring with $\mathbb Q_p$. It does not
discard integral indices. Brown's theorem is a classical input here;
its entire proof has not been re-refereed in this run.

\paragraph{AG-LM6-IRREG (where the regular-prime argument uses more).}
Write $K=\mathbb Q(\mu_{p^\infty})$ and
$G=\operatorname{Gal}(\Omega_p^*/K)$, with its outer-action
filtration $F^mG$ and graded Lie algebra $\mathfrak g$.
Choose the usual maximal-image Soulé elements $\sigma_m\in F^mG$
for odd $m\ge3$. They define maps $\Psi:S\to G$ from the free
pro-$p$ group on these generators and
$\psi:\mathfrak s\to\mathfrak g$ from the free graded
$\mathbb Z_p$-Lie algebra. Sharifi,
\emph{Relationships between conjectures on the structure of pro-$p$
Galois groups unramified outside $p$}, Proc. Sympos. Pure Math.~70
(2002), 275--284, author PDF Theorem 1.1 p.~2 and proof pp.~8--9,
uses the freeness of
$\operatorname{Gal}(\Omega_p/\mathbb Q(\mu_p))$ at regular primes.
The full proof and its preceding group lemmas have been read.

For irregular $p$, his Theorem 1.3 (author PDF p.~3; proof p.~10),
together with Brown's rational freeness, says that $\Psi$ and
$\psi$ are not surjective if Greenberg's generalized conjecture
holds for $\mathbb Q(\mu_p)$. Here that conjecture says that the
unramified abelian pro-$p$ Iwasawa module over the compositum of
all $\mathbb Z_p$-extensions is pseudo-null. This is a conditional
obstruction to the proposed generator argument, not a disproof of
$\Omega_p^*=\Omega_p$: the latter could hold with additional
integral generators and relations.

\begin{proposition}[AG-LM6-NORM: the filtered Soulé normalization]
For odd $p$ and odd $m\ge3$, in the integral normalization of
Sharifi's equation (2.10), author PDF p.~7,
\[
 \kappa_m(F^mG)\subseteq
 p^{\sum_{j\ge1}\lfloor(m-1)/p^j\rfloor}\mathbb Z_p(m).
\]
In particular $\kappa_5(F^5G)\subseteq3\mathbb Z_3(5)$,
although $3$ is regular. Thus surjectivity of these filtered
characters onto the full displayed lattice is false in every
degree with $m-1\ge p$.
\end{proposition}
\begin{proof}
The cited identity is
\[
 \kappa_m=(p^{m-1}-1)(m-1)!\lambda_m^{(p)}
       \quad\hbox{on }F^mG,
\]
where $\lambda_m^{(p)}$ is integral, obtained from the integral
stable-derivation algebra. The first factor is a $p$-adic unit.
Counting the multiples of $p^j$ among $1,\ldots,m-1$ gives
$v_p((m-1)!)=\sum_{j\ge1}\lfloor(m-1)/p^j\rfloor$.
The inclusion follows, and for $(p,m)=(3,5)$ this valuation is $1$.
\end{proof}

This is a precise test for the suggestion in
\path{OURS_FABLE_PROGRAM_SUGGESTIONS_20261120_EXPLAINED.md},
Section 5, linked there to LM6 A.44--A.45, bridge B3, and cote 149
scans 79--85. It does not assert nonsurjectivity on a different,
larger Galois domain, and it does not replace a normalized character
by an unnormalized one. A proposed extension of the regular-prime
proof must identify the exact integral generator quotient, its
relations, and the cohomological comparison producing them.
Neither endpoint A nor endpoint M supplies these data. The
$p=2$ and irregular-prime kernel-field questions in the intake
therefore remain unresolved by this argument; they are not listed
as refuted.

\subsection{The formal-height mechanism in Alpöge step 2}
\label{ag:formal-height}

The supplied Fable suggestion, Section 2, gives the height-dependent
derivative estimate but passes through a Lubin--Tate normal form.
The following proof works directly with the given formal group.
It strengthens the estimate to an explicit bound and requires no
identification of a general deformation with a particular
Lubin--Tate group.

\begin{proposition}[AG-FG-EXP: an exact exponential coefficient bound]
Let $R$ be a complete unramified discrete valuation ring of mixed
characteristic $(0,p)$, with fraction field $K$, perfect residue
field, and valuation $v(p)=1$, so $v(K^\times)=\mathbb Z$;
set $v(0)=+\infty$.
Let $F(X,Y)$ be a one-dimensional commutative formal group over
$R$. Suppose its reduction has finite height $h\ge1$, and put
$q=p^h$. Write
\[
 [p]_F(t)=pt+\sum_{j\ge2}a_jt^j.
\]
Thus $a_j\in pR$ for $2\le j<q$, $a_q\in R^\times$, and
all $a_j\in R$. Normalize the formal exponential by
\[
 e(z)=\exp_F(z)=z+\sum_{n\ge2}e_nz^n.
\]
Then
\[
 v(e_n)\ge-\frac{n-1}{q-1},\qquad
 v(e_n)\ge-\left\lfloor\frac{n-1}{q-1}\right\rfloor.
\]
The second inequality records that the valuation group is integral.
\end{proposition}
\begin{proof}
We include the logarithm construction to specify all normalizations.
Put $b(t)=\partial_2F(t,0)\in1+tR[[t]]$ and
\[
 \omega_F=\frac{dt}{b(t)},\qquad D=b(t)\frac{d}{dt}.
\]
The differential is integral, and $D$ preserves $R[[t]]$ with
$(Dt)(0)=1$. Associativity, differentiated in the third variable
at zero, gives
\[
 b(F(x,y))=\partial_2F(x,y)b(y).
\]
Define $\lambda(t)=\int_0^t b(u)^{-1}\,du\in K[[t]]$.
The displayed identity implies
$\partial_y\lambda(F(x,y))=\lambda'(y)$.
Evaluating the resulting constant at $y=0$ yields
$\lambda(F(x,y))=\lambda(x)+\lambda(y)$.
Its linear coefficient is $1$, so it has a formal inverse $e$,
with $e(0)=0$ and $e'(0)=1$. In particular
\[
 e(pz)=[p]_F(e(z)).
\]

We prove the first inequality by induction; it is equality for
$n=1$. Comparing the coefficient of $z^n$ in this identity gives
\[
 (p^n-p)e_n=
 \sum_{j=2}^n a_j
  \sum_{\substack{n_1+\cdots+n_j=n\\n_i\ge1}}
       e_{n_1}\cdots e_{n_j}.
\]
Every $n_i<n$, and the induction hypothesis bounds the valuation
of each product below by $-(n-j)/(q-1)$.
For $2\le j<q$, multiplication by $a_j$ adds at least $1$;
division by $p^n-p$, whose valuation is $1$, therefore leaves
a bound $-(n-j)/(q-1)\ge-(n-1)/(q-1)$.
For $j\ge q$, multiplication adds at least $0$, and division
leaves
\[
 -1-\frac{n-j}{q-1}\ge-\frac{n-1}{q-1}.
\]
The ultrametric inequality proves the induction. Since each
$e_n\in K$ has integral valuation, rounding up gives the second
bound. No convergence assertion or special normal form was used.
\end{proof}

\begin{proposition}[AG-FG-DER: integral derivatives with the exact error]
With $R,F,D,h,q$ as above, for every $f\in R[[t]]$ and integer
$a\ge1$,
\begin{equation}\label{eq:ag-height-derivative}
 v\bigl((D^af)(0)\bigr)\ge
 \max\left\{0,\ v(a!)-
                 \left\lfloor\frac{a-1}{q-1}\right\rfloor\right\}.
\end{equation}
In particular, writing $s_p(a)$ for the base-$p$ digit sum,
\[
 v\bigl((D^af)(0)\bigr)\ge
 a\left(\frac1{p-1}-\frac1{p^h-1}\right)
 -\frac{s_p(a)}{p-1}+\frac1{p^h-1}.
\]
For height $2$ the coefficient of $a$ is
\[
 \frac{p}{p^2-1}=\frac1p+\frac1{p^3}+\frac1{p^5}+\cdots.
\]
\end{proposition}
\begin{proof}
Differentiating $\lambda(e(z))=z$ gives $e'(z)=b(e(z))$.
Consequently, for each $a$,
\[
 (D^af)(0)=a!\,[z^a]f(e(z)).
\]
In a term of $[z^a]e(z)^j$, with $1\le j\le a$, the
coefficient is a product of $j$ coefficients whose indices sum
to $a$. AG-FG-EXP bounds its valuation by $-(a-j)/(q-1)$,
which is at least $-(a-1)/(q-1)$. The integral coefficients of
$f$ and the ultrametric inequality give the same bound for
$[z^a]f(e(z))$. Rounding to integral valuations and multiplying
by $a!$ gives the second entry in the maximum. The first follows
because $D$ preserves $R[[t]]$.

Finally, counting multiples of $p^j$ among $1,\ldots,a$ gives
$v(a!)=\sum_{j\ge1}\lfloor a/p^j\rfloor
=(a-s_p(a))/(p-1)$. Subtracting $(a-1)/(q-1)$ gives the
displayed digit-sum bound. The height-$2$ identity is the
geometric-series sum. Since
$s_p(a)\le(p-1)(1+\lfloor\log_pa\rfloor)$, it also gives the
claimed linear gain with an explicit logarithmic error.
\end{proof}

\begin{proposition}[AG-FG-PPOWER: the exact restricted derivation]
With the same coordinate and normalization, reduction modulo $p$
satisfies
\[
 \overline D^{p}=\overline{a_p}\,\overline D
             \quad\hbox{on }(R/p)[[t]].
\]
Thus the $p$-th power is zero for $h\ge2$ and nonzero for $h=1$.
\end{proposition}
\begin{proof}
In characteristic $p$, the iterated Leibniz formula shows that
the $p$-th power of a derivation is a derivation. It is invariant
under the formal-group translations because $\overline D$ is.
Every invariant derivation is a scalar multiple of
$\overline D$: an invariant vector field $c(t)\partial_t$
satisfies $c(F(t,u))=\partial_1F(t,u)c(t)$, and setting $t=0$
gives $c(u)=b(u)c(0)$. It remains to compute the scalar
$(D^pt)(0)$ modulo $p$.

The coefficients $e_i$ are integral for $1\le i<p$ by
AG-FG-EXP, since $q\ge p$. In the coefficient recursion for
$n=p$, every summand with $2\le j<p$ lies in $pR$, and the
$j=p$ term is $a_p$. Hence
\[
 pe_p\equiv-a_p\pmod p.
\]
As $(D^pt)(0)=p!e_p$ and $(p-1)!\equiv-1\pmod p$,
the scalar is $a_p$ modulo $p$. The factorial congruence follows
by pairing inverse elements in $\mathbb F_p^\times$ and holds
also at $p=2$. The height definition gives the last assertions.
\end{proof}

\begin{proposition}[AG-FG-PRODUCT: independent formal factors]
Suppose the integral formal group is a product of one-dimensional
groups as above, of heights $h_1,\ldots,h_r$, with commuting
normalized invariant derivations $D_1,\ldots,D_r$. For
$f\in R[[t_1,\ldots,t_r]]$ and integers $a_i\ge0$, put
\[
 B_i(0)=0,\qquad B_i(a)=
 \max\left\{0,v(a!)-
       \left\lfloor\frac{a-1}{p^{h_i}-1}\right\rfloor\right\}
          \quad(a\ge1).
\]
Then
\[
 v\bigl((D_1^{a_1}\cdots D_r^{a_r}f)(0)\bigr)
       \ge\sum_{i=1}^r B_i(a_i).
\]
\end{proposition}
\begin{proof}
Regard $f$ as a power series in all variables other than $t_1$,
with coefficients in $R[[t_1]]$. Apply AG-FG-DER to each
coefficient and evaluate at $t_1=0$. The result lies in
$p^{B_1(a_1)}R[[t_2,\ldots,t_r]]$. Divide out this scalar and
repeat for each remaining coordinate. Evaluation and the other
coordinate derivations commute, proving the bound. If $a_i=0$,
the corresponding operation is simply evaluation and preserves
integrality. An isogeny decomposition over $K$ alone is not an
integral product decomposition and does not meet this hypothesis.
\end{proof}

\paragraph{AG-FG-THETA (what transfers to a theta calculation).}
Alpöge's note, step 2, PDF p.~3, asserts the height-$2$ exponent
for the normalized derivatives of a CM theta value. AG-FG-DER
supplies that exponent whenever the theta section, in a specified
local trivialization and integral parameter, is an element of
$R[[t]]$, and the operator is the integral invariant derivation
above. The period and trivialization comparison must be supplied
for that theta object. More generally, for constants
$u,b\in K^\times$ and $\widetilde D=uD$,
\[
 v\bigl((\widetilde D^{a}(bf))(0)\bigr)
  =v(b)+a\,v(u)+v\bigl((D^af)(0)\bigr),
\]
so their effect on the bound is explicit. A nonconstant change of
trivialization or a connection requires its own comparison.
The supplied Lean map, step-2 row, reports only a low-confidence
match to CM-curve/invariant-derivation declarations; no new exact
Lean lemma identifier is asserted here.

This proves the formal-group mechanism in Fable Section 2 and
shrinks its general-height gate to the actual theta integrality
and normalization data. Its archive map is inherited explicitly
from that record: cote 54 scans 5--7, the restricted $p$-Lie
algebra and Cartier-duality discussion; cotes 4--9 for crystals;
cotes 136--137, cote 16 scan 14, and cote 47 scan 9 for the
divided-power line. These additional sheets have not been newly
image-checked here. The theorem does not by itself place a theta
section on the crystalline site or replace a theta-lift chain.
It uses neither A nor M.

\subsection{Integral elliptic functions and actual special-value divisibility}
\label{ag:elliptic-integrality}

There is a concrete sector in which the formal integrality required
by AG-FG-THETA can be proved directly. The normalizations below agree
with the regularized Weierstrass functions used by Bannai--Kobayashi,
\emph{Integral structures on $p$-adic Fourier theory}, Ann. Inst.
Fourier 66 (2016), 521--550, Section 5, especially printed pp.~543--546
(Numdam PDF pp.~24--27). Their Theorems 5.1 and 5.8 locate the stronger
classical Katz congruences. The following explicit floor bound is a
direct application of the general formal-height argument proved here;
it does not assume a Lubin--Tate normal-form coordinate.

\begin{proposition}[AG-THETA-REG: integral regularizations]
Let $R$ be a complete unramified discrete valuation ring of mixed
characteristic $(0,p)$, $p\ge5$, with $v(p)=1$. Let $E/R$ have good
reduction and equation
\[
 y^2=4x^3-g_2x-g_3,\qquad g_2,g_3\in R,
\]
with unit discriminant. Use the parameter $t=-2x/y$, invariant
differential $\omega=dx/y$, and invariant derivation $D$ dual to
$\omega$. Let the formal group have height $h\in\{1,2\}$, and put
\[
 d_h(a)=\max\left\{0,\ v(a!)-
                   \left\lfloor\frac{a-1}{p^h-1}\right\rfloor\right\}
 \quad(a\ge1).
\]
Suppose $[b]\in\operatorname{End}_R(E)$ satisfies
$[b]^*\omega=b\omega$ with $b\in R^\times$. Then
\[
 f_b(t)=x(t)-b^2x([b](t))\in R[[t]],\qquad
 v(D^a f_b(0))\ge d_h(a).
\]
Define the normalized Weierstrass coefficients $\beta_m$ by
\[
 x(\exp_F z)=z^{-2}+\sum_{a\ge0}\beta_{a+2}\frac{z^a}{a!}.
\]
For $a\ge1$ one has the exact identity and bound
\begin{equation}\label{eq:ag-bh-bound}
 D^a f_b(0)=(1-b^{a+2})\beta_{a+2},\qquad
 v((1-b^{a+2})\beta_{a+2})\ge d_h(a).
\end{equation}
In particular
$v(\beta_m)\ge d_h(m-2)-1-v_p(m)$ for $m\ge3$.
If $E$ has an action of the full ring of integers of an imaginary
quadratic CM field over $R$, $p$ is
inert in that CM field, and $h=2$, then the stronger conclusion
$v(\beta_m)\ge d_2(m-2)$ holds whenever $p^2-1\nmid m$.
\end{proposition}
\begin{proof}
Write $x=t^{-2}u$. The equation of the curve gives
\[
 u=1+\frac{g_2t^4}{4u}+\frac{g_3t^6}{4u^2}.
\]
Its formal solution has $u\in1+t^4R[[t^2]]$, by coefficient
recursion, since the derivative of the defining equation with
respect to $u$ at $(t,u)=(0,1)$ is a unit. Thus $x(t)$ has integral
Laurent coefficients, leading term $t^{-2}$, and no other negative
power. The endomorphism has integral odd expansion
$[b](t)=bt+O(t^3)$: oddness follows from its compatibility with
inversion, which is $t\mapsto-t$. Since $b$ is a unit,
$b^2x([b](t))$ has the same sole negative term $t^{-2}$.
Their difference is therefore integral and regular in the formal
parameter. The derivation is normalized by $\omega/dt=1$ at zero,
so AG-FG-DER applies and gives the displayed derivative bound.

In characteristic zero the formal logarithm linearizes every
endomorphism: $[b](\exp_F z)=\exp_F(bz)$. Substitution into the
Laurent expansion cancels the pole and gives
\eqref{eq:ag-bh-bound} coefficient by coefficient.
Take now the integer multiplication endomorphism $b=1+p$.
For odd $p$, binomial expansion gives
$v_p((1+p)^m-1)=1+v_p(m)$: for $p\nmid m$ the first term
has valuation one and all later terms have larger valuation;
each further raising to the $p$th power increases the valuation
by one. This proves the uniform bound for $\beta_m$.

In the inert CM case, the residues of the CM integers contain
$\mathbb F_{p^2}$, so choose $b$ whose residue generates
$\mathbb F_{p^2}^{\times}$. It acts as the stipulated
endomorphism and is a unit. When $p^2-1\nmid m$, the factor
$1-b^m$ is a unit, proving the stronger bound.
\end{proof}

\begin{proposition}[AG-THETA-TORS: torsion translations close another integrality gate]
Keep the elliptic model and normalization above. If $P\in E(R)$
reduces away from the origin, then
\[
 x(P+_E t)\in R[[t]],\qquad
 v(D^a x(P))\ge d_h(a)\quad(a\ge1).
\]
Every nonzero torsion point of order prime to $p$ satisfies the
reduction condition after a finite unramified extension defining
it. For a model defined over a number field and a compatible
complex uniformization $E(\mathbb C)=\mathbb C/\Lambda$ with
$\omega=dz$, write $z_P$ for a lift of $P$. Then
\begin{equation}\label{eq:ag-eisenstein-derivative}
 D^a x(P)=(-1)^a(a+1)!
     \sum_{\gamma\in\Lambda}(z_P+\gamma)^{-a-2}.
\end{equation}
The valuation bound therefore applies to these actual algebraic
normalized Eisenstein values.
\end{proposition}
\begin{proof}
The rational function $x$ is regular on $E\setminus\{O\}$.
Translation by $P$ takes the formal neighbourhood of $O$ into
this open subscheme when $\bar P\ne O$. Pullback of $x$ is
therefore an element of $R[[t]]$. Translation commutes with the
invariant derivation, so AG-FG-DER proves the asserted bound at $P$.
For an integer $n$ prime to $p$, multiplication by $n$ on the
formal group has unit linear term and hence an inverse formal
series over $R$. It has no nonzero $n$-torsion in that
neighbourhood. Thus a nonzero prime-to-$p$ torsion point cannot
reduce to $O$. Moreover $E[n]/R$ is finite \'etale, so its
points become rational over a finite unramified extension;
the unramified hypothesis needed in AG-FG-DER is preserved.

In complex coordinates $x=\wp_\Lambda(z)$ and $D=d/dz$.
Differentiating the defining series for $\wp$ $a\ge1$ times
gives \eqref{eq:ag-eisenstein-derivative}; the resulting series
converges absolutely since $a+2>2$. The left side is algebraic
and fixes the period and differential normalization of the
right side. Its valuation is thus well defined through the
chosen $p$-adic embedding; no valuation of a free-standing
complex period is being asserted.
\end{proof}

\begin{proposition}[AG-THETA-HECKE: an explicitly normalized Hecke-value family]
Let the global elliptic curve have CM by $\mathcal O_K$, where
$K$ is imaginary quadratic of class number one. Suppose its
period lattice for the chosen algebraic differential is
$\Lambda=\Omega\mathcal O_K$. Put
$w_K=|\mathcal O_K^\times|$. For $m\ge3$ divisible by $w_K$,
the conductor-one algebraic Hecke character
$\psi_m((\alpha))=\alpha^m$ is well defined, and
\[
 \beta_m=w_K(m-1)!\,\Omega^{-m} L(m,\psi_m).
\]
Consequently all bounds for $\beta_m$ in AG-THETA-REG apply
to this displayed algebraic normalization. These values are
already in the region of absolute convergence, so the
divisibility conclusion uses neither A nor M.
\end{proposition}
\begin{proof}
Every nonzero integral ideal of $\mathcal O_K$ is principal,
and changing its generator multiplies $\alpha^m$ by $1$.
Thus the formula defines the stated Hecke character. Its value
on an ideal has complex absolute value $N\mathfrak a^{m/2}$,
so the series at $s=m$ is absolutely convergent for $m>2$.
Counting the $w_K$ generators of each ideal gives
\[
 \begin{aligned}
 L(m,\psi_m)&=\frac1{w_K}
   \sum_{\alpha\in\mathcal O_K\setminus\{0\}}
          \frac{\alpha^m}{(\alpha\bar\alpha)^m}\\
  &=\frac1{w_K}
   \sum_{\alpha\in\mathcal O_K\setminus\{0\}}\alpha^{-m}.
 \end{aligned}
\]
where the final equality uses the bijection
$\alpha\mapsto\bar\alpha$ of $\mathcal O_K$.
The Laurent expansion of $\wp_\Lambda$ has coefficient
$(m-1)\sum_{\gamma\in\Lambda\setminus\{0\}}\gamma^{-m}$
at $z^{m-2}$; this follows by expanding each summand in its
defining series, absolutely for $m\ge3$.
Comparison with the definition of $\beta_m$ gives
$\beta_m=(m-1)!\sum_{\gamma\ne0}\gamma^{-m}$, and scaling
the lattice by $\Omega$ proves the identity.
\end{proof}

For example, for $y^2=4x^3-4x$ over $\mathbb Q(i)$ at $p=7$,
take $m=52$. Then $48\nmid52$ and
$d_2(50)=v_7(50!)-\lfloor49/48\rfloor=8-1=7$.
Thus the normalized value
$4\cdot51!\,\Omega^{-52}L(52,\psi_{52})$ is divisible by
$7^7$ in the specified algebraic normalization. This is a
concrete special-value consequence of the formal method.
The larger theta-lift programme still needs its own section,
period and operator comparison; these proofs close the
regularized elliptic-function and prime-to-$p$ torsion sectors.

\begin{proposition}[AG-CM-INTEGRAL-MODEL: existence of the required CM models]
For every imaginary quadratic field $K$, with class number $h_K$
and $w_K$ roots of unity, there are a number field $F\supset K$
and an elliptic scheme with full $\mathcal O_K$-action over
$\mathcal O_F[1/3]$ such that
\[
 [F:K]\le\frac{48h_K}{w_K},\qquad
 F/K\text{ is unramified away from }3.
\]
A complex embedding can be chosen so that its generic curve is
isomorphic to $\mathbb C/\mathcal O_K$. At any prescribed finite
set of places above rational primes $p\ge5$ unramified in $K$,
one may choose one algebraic invariant differential which is
an integral generator at every chosen place. Its complex period
lattice is then $\Omega\mathcal O_K$.

In particular the model and differential hypotheses of
AG-THETA-CLASS/RAY/PCOND can be realized for every $K$ and every
finite set of inert primes $p\ge5$; at those places the completed
base is unramified over $\mathbb Q_p$ and the formal height is two.
This is an existence statement with a specified integral
normalization, not a claim about every twist or every preselected
differential of a given CM curve.
\end{proposition}
\begin{proof}
Use the CM moduli stack $\mathcal M$ of pairs $(E,\iota)$ with
$\iota:\mathcal O_K\to\operatorname{End}(E)$ inducing the
structure map on the Lie algebra. Kudla--Rapoport--Yang,
\emph{A peculiar modular form of weight one}, arXiv:math/9808143,
Section 5, Proposition 5.1 and its two deformation cases,
PDF pp.~23--25, prove the requisite finite étale CM moduli
statement over $\mathcal O_K$. This section explicitly allows
arbitrary imaginary quadratic $K$. Proposition 5.3 and
Corollary 5.4 identify its geometric characteristic-zero
isomorphism classes with $\operatorname{Cl}(K)$ and its coarse
scheme with the Hilbert class field's integer scheme.

Over $\mathcal O_K[1/3]$, add an ordered basis of $E[3]$.
The choices form a finite étale torsor with group
$\operatorname{GL}_2(\mathbb F_3)$, of order $48$.
No nontrivial automorphism of a CM pair fixes such a basis.
Indeed that automorphism is a unit $u\in\mathcal O_K^\times$;
if $u\ne1$, the isogeny $[u-1]$ has degree
$N_{K/\mathbb Q}(u-1)=|u-1|^2\le4$ and cannot kill the
subgroup $E[3]$ of degree $9$. The rigidified moduli object
therefore has no stabilizers and is a finite étale scheme
$Y$ over $\mathcal O_K[1/3]$, with its universal elliptic curve
and CM action. There are exactly $48/w_K$ level classes over
each characteristic-zero CM isomorphism class, so
$\deg(Y/\mathcal O_K[1/3])=48h_K/w_K$.

Choose the component containing the complex point
$\mathbb C/\mathcal O_K$ with any ordered basis of its
$3$-torsion. This component is
$\operatorname{Spec}\mathcal O_F[1/3]$ for a finite extension
$F/K$, since a connected finite étale scheme over this
normal Dedekind base is the corresponding normalized integer
scheme. Its degree is at most the degree of $Y$ and it is
unramified away from $3$. Pull back the universal curve.
At a prime unramified in $K/\mathbb Q$, its local base is
therefore unramified over $\mathbb Q_p$. The height-two
assertion at inert primes is exactly Case 2 of the cited
deformation analysis.

The invariant differential module is an invertible module on
$\mathcal O_F[1/3]$. Choose any nonzero differential in its
one-dimensional generic fibre. At the finitely many chosen
places it differs from a local integral generator by a scalar
of specified integral valuation. Weak approximation supplies
a single element of $F^\times$ with the opposite valuations.
Multiplying the differential by that element makes it a
generator at every chosen place. The chosen complex curve
has principal CM lattice, so scaling its differential changes
the common period $\Omega$ and leaves the lattice in the form
$\Omega\mathcal O_K$. Any two local integral generators
differ by a local unit; the valuation conclusions are invariant
under that remaining unit ambiguity.
\end{proof}

\begin{proposition}[AG-CM-MODEL-DEGREE: an A-only quantitative CM model]
Let $-D$ be a negative fundamental discriminant with
$D\ge10^{10}$, and put $K=\mathbb Q(\sqrt{-D})$.
Under adopted A, the field and CM model in
AG-CM-INTEGRAL-MODEL can be chosen with
\[
 [F:K]\le\frac{48e^\gamma}{\pi}\sqrt D
    \left(\log\log D-\log2+\frac12+
                         \frac1{\log\log D}\right).
\]
The extension remains unramified away from $3$ over $K$,
and the simultaneous local differential normalizations in
that proposition require no increase of its degree.
\end{proposition}
\begin{proof}
Here $w_K=2$. NT-H02, from Lamzouri--Li--Soundararajan,
corrected arXiv:1309.3595v3, Corollary 1.3 and (1.3),
author p.~4, gives
\[
 h_K\le\frac{2e^\gamma}{\pi}\sqrt D
    \left(\log\log D-\log2+\frac12+
                         \frac1{\log\log D}\right)
\]
from Dirichlet GRH, now supplied by A.
Multiply this by $48/w_K=24$ in the proved model-degree
bound. The model construction and weak approximation do
not otherwise change. This is a quantitative statement about
a field carrying a universal level-three CM elliptic curve,
or equivalently a point of this zero-dimensional CM moduli
scheme; it asserts neither a small defining polynomial nor
an algorithm for constructing it within the degree bound.
\end{proof}

\begin{proposition}[AG-THETA-CLASS: the class-number restriction is unnecessary]
Let $K$ be any imaginary quadratic field, with ring of integers
$\mathcal O$ and $w=|\mathcal O^\times|$. Fix a CM elliptic curve
$E_0$ over a number field, with its full $\mathcal O$-action and an
algebraic differential $\omega$ whose complex period lattice is
$\Omega\mathcal O$. At the chosen $p$-adic embedding, suppose that
$p\ge5$, the curve has a good model over the integers $R$ of a
finite unramified extension of $\mathbb Q_p$, the action extends
over $R$, and $\omega$ generates its invariant differentials.
Write $h$ for its formal height and normalize $v(p)=1$.
Use short Weierstrass models normalized by these differentials;
such models exist integrally since $6\in R^\times$.

Let $\psi$ be any conductor-one algebraic Hecke character of
infinity type $(m,0)$, where $m\ge3$; necessarily $w\mid m$.
Choose integral ideal-class representatives $\mathfrak a$ prime
to $p$. The quotient $E_{\mathfrak a}=E_0/E_0[\mathfrak a]$ has
a normalized differential $\omega_{\mathfrak a}$ with pullback
$\omega$ and lattice $\Omega\mathfrak a^{-1}$. With the
coefficients $\beta_m$ defined in AG-THETA-REG for these models,
\begin{equation}\label{eq:ag-cm-class-sum}
 w(m-1)!\,\Omega^{-m}L(m,\psi)
   =\sum_{[\mathfrak a]\in\operatorname{Cl}(K)}
      \psi(\mathfrak a)^{-1}
      \beta_m(E_{\mathfrak a},\omega_{\mathfrak a}).
\end{equation}
All the weights $\psi(\mathfrak a)^{-1}$ are $p$-adic units.
Consequently the valuation of the displayed normalized Hecke
value is at least
\[
 d_h(m-2)-1-v_p(m),
\]
and it is at least $d_2(m-2)$ when $p$ is inert, $h=2$, and
$p^2-1\nmid m$. There is no denominator equal to the class number.
\end{proposition}
\begin{proof}
First we check the integral normalization under ideal isogenies.
Since $N\mathfrak a$ is a unit in $R$, the subgroup
$E_0[\mathfrak a]\subset E_0[N\mathfrak a]$ is finite
\'etale. After a finite unramified extension it is constant;
its nonzero points have integral coordinates and remain distinct
from the origin after reduction, by AG-THETA-TORS. Apply the
normalized quotient construction of V\'elu, \emph{Isog\'enies
entre courbes elliptiques}, C.~R.~Acad.~Sci.~Paris 273 (1971),
A238--A241, formulas (7), (11), and Remark 2; in Ghitza's
translation these occur on PDF pp.~2, 4. Its coefficient formulas
are integral polynomials in the kernel coordinates. Their reductions
construct the quotient by the reduced, still separable kernel,
so the resulting cubic remains nonsingular. The construction
descends to $R$. Its pullback on invariant differentials is the
identity. Thus the quotient has good reduction, its differential
is an integral generator, and its formal group is isomorphic to
that of $E_0$, so has height $h$. The $\mathcal O$-action descends
because it preserves the kernel. Complex uniformization identifies
this quotient with $\mathbb C/(\Omega\mathfrak a^{-1})$ and
identifies its normalized differential with $dz$.

Every nonzero integral ideal in the class $[\mathfrak a]^{-1}$
has the form $(\alpha)\mathfrak a^{-1}$ for exactly $w$
elements $\alpha\in\mathfrak a\setminus\{0\}$, modulo this
unit multiplicity. Its character value and norm are
\[
 \psi((\alpha)\mathfrak a^{-1})
       =\alpha^m\psi(\mathfrak a)^{-1},\qquad
 N((\alpha)\mathfrak a^{-1})
       =\alpha\bar\alpha/N\mathfrak a.
\]
Since $|\psi(\mathfrak b)|=(N\mathfrak b)^{m/2}$, all series
below converge absolutely when $m>2$. Partitioning the Hecke
series by ideal classes gives
\[
 L(m,\psi)=\frac1w\sum_{[\mathfrak a]}
  \frac{(N\mathfrak a)^m}{\psi(\mathfrak a)}
      \sum_{0\ne\alpha\in\mathfrak a}\bar\alpha^{-m}
 =\frac1w\sum_{[\mathfrak a]}\psi(\mathfrak a)^{-1}
      \sum_{0\ne\gamma\in\mathfrak a^{-1}}\gamma^{-m}.
\]
The second equality uses the exact ideal identity
$\bar{\mathfrak a}=(N\mathfrak a)\mathfrak a^{-1}$.
The Laurent-coefficient calculation in AG-THETA-HECKE now proves
\eqref{eq:ag-cm-class-sum} for each lattice separately.

If the order of $[\mathfrak a]$ is $e$ and
$\mathfrak a^e=(\alpha)$, then
$\psi(\mathfrak a)^e=\alpha^m$. Because $\mathfrak a$ is prime
to $p$, $\alpha$ is a unit at every prime above $p$, hence
$v(\psi(\mathfrak a))=0$. Apply AG-THETA-REG to each normalized
quotient model and then use the ultrametric inequality in
\eqref{eq:ag-cm-class-sum}. The field containing the character
values may be ramified at $p$: only their unit valuations are
used. The formal derivative estimates were made on the unramified
models before extending the valuation. Neither averaging over
ideal classes nor dividing by their number occurs.
\end{proof}

\begin{proposition}[AG-THETA-RAY: proper conductor gives a torsion-value formula]
Keep the global CM model, embeddings, differential, and $p$-adic
hypotheses of AG-THETA-CLASS. Let $\psi$ have infinity type
$(m,0)$, $m\ge3$, and a proper conductor ideal
$\mathfrak f\ne\mathcal O$ prime to $p$. Fix the convention
\[
 \psi((\alpha))=\chi_{\mathfrak f}(\alpha)\alpha^m
 \quad ((\alpha,\mathfrak f)=1),
\]
where $\chi_{\mathfrak f}$ is a finite character of
$(\mathcal O/\mathfrak f)^\times$. In particular
$\chi_{\mathfrak f}(u)u^m=1$ for every unit $u$.
Choose integral class representatives $\mathfrak a$ prime to
$p\mathfrak f$, and for every
$c\in(\mathcal O/\mathfrak f)^\times$ choose
$\alpha_{\mathfrak a,c}\in\mathfrak a$ with residue $c$.
Let $E_{\mathfrak a,\mathfrak f}$ have differential-normalized
lattice and point
\[
 \Lambda_{\mathfrak a,\mathfrak f}
    =\Omega\bar{\mathfrak f}\mathfrak a^{-1},\qquad
 P_{\mathfrak a,c}
    =\frac{\Omega\bar\alpha_{\mathfrak a,c}}{N\mathfrak a}
          \pmod{\Lambda_{\mathfrak a,\mathfrak f}}.
\]
These curves have good unramified models with normalized integral
differentials and height $h$; the points are nonzero torsion of
order prime to $p$. Writing $D_{\mathfrak a,\mathfrak f}$ for
the corresponding invariant derivation, one has
\begin{equation}\label{eq:ag-cm-ray-sum}
 \begin{aligned}
 w(m-1)!\,\Omega^{-m}L(m,\psi)
   ={}&(-1)^{m-2}\sum_{[\mathfrak a]}
         \psi(\mathfrak a)^{-1}
       \sum_{c\in(\mathcal O/\mathfrak f)^\times}
          \chi_{\mathfrak f}(c)\\[-2pt]
     &\hspace{25mm}\cdot
       D_{\mathfrak a,\mathfrak f}^{m-2}
          x_{\mathfrak a,\mathfrak f}(P_{\mathfrak a,c}).
 \end{aligned}
\end{equation}
It follows that
\[
 v\bigl(w(m-1)!\,\Omega^{-m}L(m,\psi)\bigr)
       \ge d_h(m-2).
\]
No condition $w\mid m$ or $p^2-1\nmid m$ is needed here.
The expression whose valuation is taken is algebraic by the
displayed identity, with the period fixed by $\omega$.
\end{proposition}
\begin{proof}
Because $\mathfrak a$ is prime to $\mathfrak f$, its residue
map onto $\mathcal O/\mathfrak f$ is surjective, with kernel
$\mathfrak f\mathfrak a$. Thus the choices in the statement
exist, and their cosets are precisely the elements of
$\mathfrak a$ prime to $\mathfrak f$.

The normalized quotient attached to the ideal
$\mathfrak f\mathfrak a$ in the preceding proof has lattice
$\Omega\mathfrak f^{-1}\mathfrak a^{-1}$ and good model of
height $h$. Scale its lattice coordinate by the integer
$N\mathfrak f$. Since
$(N\mathfrak f)\mathfrak f^{-1}=\bar{\mathfrak f}$, this
gives the asserted lattice. On the Weierstrass model the change
is $x'=(N\mathfrak f)^{-2}x$,
$y'=(N\mathfrak f)^{-3}y$, and
$\omega'=(N\mathfrak f)\omega_{\mathfrak f\mathfrak a}$.
The scaling factor is a $p$-adic unit, so integrality, good
reduction, height and the generator property of the differential
are preserved.

Changing $\alpha_{\mathfrak a,c}$ by
$\mathfrak f\mathfrak a$ changes its displayed lift by
$\Omega\bar{\mathfrak f}\mathfrak a^{-1}$, so the point is
well defined. Multiplication by any element of
$\bar{\mathfrak f}$ takes it into that lattice; in particular
its order divides $N\mathfrak f$. It is nonzero, since being
zero would imply
$\alpha_{\mathfrak a,c}\in\mathfrak f\mathfrak a$, contrary
to $c$ being a unit modulo the proper ideal $\mathfrak f$.
These prime-to-$p$ points are defined after a finite unramified
extension and satisfy AG-THETA-TORS.

Now partition the absolutely convergent ideal series first by
class and then by residue. The same $w$-to-one generator count
as before gives
\[
 L(m,\psi)=\frac1w\sum_{[\mathfrak a]}
   \frac{(N\mathfrak a)^m}{\psi(\mathfrak a)}
   \sum_c\chi_{\mathfrak f}(c)
      \sum_{\alpha\in\alpha_{\mathfrak a,c}
                       +\mathfrak f\mathfrak a}\bar\alpha^{-m}.
\]
The unit compatibility in the statement guarantees that each
of the $w$ generators contributes the same summand even when
$w\nmid m$. For the lattice and point just constructed,
AG-THETA-TORS gives
\[
 D_{\mathfrak a,\mathfrak f}^{m-2}
     x_{\mathfrak a,\mathfrak f}(P_{\mathfrak a,c})
  =(-1)^{m-2}(m-1)!\,\Omega^{-m}(N\mathfrak a)^m
       \sum_{\alpha\in\alpha_{\mathfrak a,c}
                        +\mathfrak f\mathfrak a}\bar\alpha^{-m}.
\]
This proves \eqref{eq:ag-cm-ray-sum}. All coefficients in its
finite sum are $p$-adic units. For $\psi(\mathfrak a)$ this
follows by raising to the class order as before; the additional
factor $\chi_{\mathfrak f}(\alpha)$ is a root of unity and
has valuation zero. Each derivative has valuation at least
$d_h(m-2)$, proving the bound. A ramified field of finite-character
values causes no loss: the estimates already hold over unramified
fields of definition of the models and torsion points, and
the scalar weights remain units after any extension.
\end{proof}

These two formulas close the class-number and prime-to-$p$
conductor gates for the elliptic-function sector of Fable
Section 2 and Alp\"oge Step 2. Geometrically the class sum
runs over the prime-to-$p$ CM isogeny orbit, while a proper
ray conductor replaces an origin expansion by evaluation on
finite \'etale torsion sections. The resulting exact periods
belong to these algebraic differentials; no comparison with
an arbitrary higher-dimensional theta section or a theta-lift
lattice is asserted. All special values used here lie in the
region of absolute convergence and require neither A nor M.

\begin{proposition}[AG-FG-POLE: formal torsion and an explicit pole loss]
Let $R,F,D,h$ satisfy AG-FG-EXP and put $q=p^h$, $Q=q-1$.
Extend $v$, still normalized by $v(p)=1$, to an algebraic closure
of the fraction field. The following estimates are valid over
every finite, possibly ramified, extension $R'/R$.
For $g\in R'[[t]]$ one still has
\[
 v((D^ag)(0))\ge d_h(a)\qquad(a\ge1).
\]
Let $P$ be a nonzero formal point with $\delta=v(t(P))>0$,
and let $f(t)\in t^{-d}R[[t]]$, $d\ge1$. If
$\delta\le1/Q$, write $a-1=\kappa Q+\eta$,
$0\le\eta<Q$. Then
\begin{equation}\label{eq:ag-formal-pole-bound}
 v(D^af(P))\ge
 \max\left\{-(a+d)\delta,\quad
       v_p(a!)-\kappa-(\eta+d+1)\delta\right\}.
\end{equation}
If $P$ is formal torsion of exact order $p^r$, $r\ge1$,
then its parameter has the exact valuation
\[
 \delta_r=\frac1{q^{r-1}(q-1)}.
\]
Thus \eqref{eq:ag-formal-pole-bound} applies with $\delta_r$.
\end{proposition}
\begin{proof}
For the first assertion, each coefficient of $[z^a]\exp_F(z)^j$
is a sum of products with indices $n_1+\cdots+n_j=a$.
AG-FG-EXP gives the integer lower bound
\[
 -\sum_i\left\lfloor\frac{n_i-1}{Q}\right\rfloor
 \ge-\left\lfloor\frac{a-j}{Q}\right\rfloor
 \ge-\left\lfloor\frac{a-1}{Q}\right\rfloor.
\]
Multiplication by an integral coefficient of $g$ cannot lower
this bound. Multiplying by $a!$ and also using that $D$ preserves
$R'[[t]]$ proves $d_h(a)$. This proof rounds the exponential
coefficients over the original unramified ring; it does not round
arbitrary valuations in $R'$.

Put $u=t(P)$. The translated series satisfies
$F(u,T)=u+T A(T)$ with $A(T)\in R'[[T]]$.
All coefficients are well defined by convergence in $u$.
Expansion of $(u+TA(T))^{-d}$ shows that the coefficient of
$T^j$ has valuation at least $-(j+d)\delta$.
The remaining factor in $f(F(u,T))$ is integral, so if
$f(F(u,T))=\sum_{j\ge0}g_jT^j$, then
\[
 v(g_j)\ge-(j+d)\delta.
\]
Translation commutes with the invariant derivation. Therefore
\[
 D^af(P)=a!\,[z^a]\sum_{j\ge0}g_j\exp_F(z)^j,
\]
and the product estimate above yields
\[
 v(D^af(P))\ge v_p(a!)+
   \min_{1\le j\le a}
       \left\{-(j+d)\delta-
           \left\lfloor\frac{a-j}{Q}\right\rfloor\right\}.
\]
To compute this minimum, put $a-j=kQ+b$, $0\le b<Q$.
The expression equals
$-(a+d)\delta+k(Q\delta-1)+b\delta$.
Because $Q\delta\le1$, its minimum occurs at
$k=\kappa$, $b=0$, equivalently $j=\eta+1$.
It is $-\kappa-(\eta+d+1)\delta$.
Independently, $D$ preserves integral Laurent series and
increases their pole order by at most one. Thus
$D^af\in t^{-d-a}R[[t]]$ and its value has valuation at least
$-(a+d)\delta$. Combining the two bounds proves
\eqref{eq:ag-formal-pole-bound}.

Finally write
$[p](t)=pt+\sum_{j\ge2}a_jt^j$ as in AG-FG-EXP.
For a point of positive parameter valuation $\delta$, the
terms with $2\le j<q$ have valuation strictly greater than
$1+\delta$, and those with $j>q$ have valuation strictly
greater than $q\delta$. Unless $\delta=1/(q-1)$, the
valuation of $[p](t(P))$ is therefore
$\min\{1+\delta,q\delta\}$. A nonzero point killed by
$p$ must have $\delta=1/(q-1)$, since otherwise a unique
term of least valuation cannot sum to zero.
For exact order $p^r$, $r>1$, induction gives the parameter
valuation of $[p]P$ as $\delta_{r-1}\le1/(q-1)$.
If $\delta\ge1/(q-1)$, the parameter of $[p]P$ would have
valuation at least $1+1/(q-1)$, a contradiction.
Hence $\delta<1/(q-1)$ and $q\delta=\delta_{r-1}$.
This proves the formula for $\delta_r$.
\end{proof}

\begin{proposition}[AG-THETA-PCOND: the inert $p$-conductor case with its denominator]
Keep the hypotheses and normalization of AG-THETA-CLASS, and
suppose $p$ is inert and the formal height is $2$.
Let $\psi$ have infinity type $(m,0)$, $m\ge3$, and conductor
\[
 \mathfrak f=(p)^r\mathfrak f_0,\qquad r\ge1,
 \quad(\mathfrak f_0,p)=1.
\]
Write $a=m-2$, $a-1=\kappa(p^2-1)+\eta$,
$0\le\eta<p^2-1$, and
\[
 \delta_r=\frac1{p^{2(r-1)}(p^2-1)},\qquad
 B_2(a,r)=\max\{-(a+2)\delta_r,\,
       v_p(a!)-\kappa-(\eta+3)\delta_r\}.
\]
Then the exactly normalized algebraic value satisfies
\begin{equation}\label{eq:ag-cm-p-conductor-bound}
 v\bigl(p^{rm}w(m-1)!\,\Omega^{-m}L(m,\psi)\bigr)
 \ge
 \begin{cases}
 d_2(m-2),&\mathfrak f_0\ne\mathcal O,\\
 B_2(m-2,r),&\mathfrak f_0=\mathcal O.
 \end{cases}
\end{equation}
Equivalently, remove $p^{rm}$ from the left by subtracting
$rm$ from the corresponding lower bound. This is an explicit
denominator estimate, not an assertion of unrestricted integrality.
\end{proposition}
\begin{proof}
Choose class representatives $\mathfrak a$ prime to
$p\mathfrak f_0$, and representatives
$\alpha_{\mathfrak a,c}\in\mathfrak a$ of
$c\in(\mathcal O/\mathfrak f)^\times$ as in AG-THETA-RAY.
Use the good normalized models with lattices
$\Omega\bar{\mathfrak f}_0\mathfrak a^{-1}$; for
$\mathfrak f_0=\mathcal O$ these are $E_{\mathfrak a}$.
On them put
\[
 P^{(r)}_{\mathfrak a,c}
   =\frac{\Omega\bar\alpha_{\mathfrak a,c}}
                    {(N\mathfrak a)p^r}
       \pmod{\Omega\bar{\mathfrak f}_0\mathfrak a^{-1}}.
\]
The residue decomposition in AG-THETA-RAY is an identity of
absolutely convergent lattice sums even when the conductor has
a $p$-part. Its lattice there is now $p^r$ times the good
lattice just displayed. Under a lattice scaling by $p^r$,
$x$ scales by $p^{-2r}$ and the invariant derivation by
$p^{-r}$. The derivative of order $m-2$ therefore scales by
$p^{-rm}$, giving the exact identity
\[
 \begin{aligned}
 p^{rm}w(m-1)!\,\Omega^{-m}L(m,\psi)
 ={}&(-1)^{m-2}\sum_{[\mathfrak a]}\psi(\mathfrak a)^{-1}
    \sum_{c\in(\mathcal O/\mathfrak f)^\times}
       \chi_{\mathfrak f}(c)\\[-2pt]
 &\hspace{22mm}\cdot
       D_{\mathfrak a,\mathfrak f_0}^{m-2}
          x_{\mathfrak a,\mathfrak f_0}
                (P^{(r)}_{\mathfrak a,c}).
 \end{aligned}
\]
Each scalar coefficient is a $p$-adic unit, by the class-order
argument already proved. The curve models are still unramified;
the torsion points and finite-character values may require
ramified extensions.

If $\mathfrak f_0$ is proper, then
$[p^r]P^{(r)}_{\mathfrak a,c}$ is the nonzero prime-to-$p$
torsion point attached to the residue of
$\alpha_{\mathfrak a,c}$ modulo $\mathfrak f_0$.
It reduces away from zero. Consequently
$P^{(r)}_{\mathfrak a,c}$ itself reduces away from zero.
Translation of $x$ is integral over the possibly ramified
field defining this point, and the first part of AG-FG-POLE
gives $d_2(m-2)$ without changing that bound.

If $\mathfrak f_0=\mathcal O$, each point has exact order
$p^r$: it is killed by $p^r$, and being killed by $p^{r-1}$
would mean $\bar\alpha_{\mathfrak a,c}\in p\bar{\mathfrak a}$,
impossible because the residue $c$ is a unit modulo $(p)^r$.
The height-two special fibre has no nonzero geometric
$p$-torsion point. Indeed, the formal multiplicity at the
origin of $[p]$ is $p^2$, the full degree of this elliptic
isogeny; hence its geometric kernel has only that point.
Iteration gives the same assertion for $p^r$.
Thus all the displayed points are formal, and AG-FG-POLE gives
their exact parameter valuation $\delta_r$ and the bound
$B_2(m-2,r)$ for $x$, whose pole order is $2$.
Applying the ultrametric inequality to the finite identity
proves both cases of \eqref{eq:ag-cm-p-conductor-bound}.
\end{proof}

The comparison now covers every finite conductor in the inert
CM elliptic sector, with the denominator displayed when the
conductor meets $p$. The ramification affects the torsion points
and scalar values, while the exponential estimate is still
computed on the original unramified formal group. General theta
sections retain their separate integral model, connection and
period comparison; these explicit functions do not supply them.

\begin{proposition}[AG-THETA-NONINT: a precise reduced-theta normalization obstruction]
On the good CM model
\[
 E:\ y^2=4x^3-4x,\qquad t=-2x/y,\qquad\omega=dx/y,
\]
over $R=\mathbb Z_7[i]$, let $\lambda(t)$ be the formal logarithm
for $\omega$. For its complex period lattice, let $\theta(z)$
be the reduced theta function for the divisor $[0]$, normalized
by $\theta'(0)=1$. Then
\begin{equation}\label{eq:ag-theta-nonintegral}
 \frac{\theta(\lambda(t))}{t}
   =1-\frac5{12}t^4+\frac{1411}{2016}t^8+O(t^{12}),
 \qquad v_7\left(\frac{1411}{2016}\right)=-1.
\end{equation}
In particular this normalized formal series is not in $R[[t]]$.
Multiplication by an integral formal unit or an integral change
of local parameter with unit linear coefficient cannot make it
integral.
\end{proposition}
\begin{proof}
The theta convention is the one in Bannai--Kobayashi,
\emph{Algebraic theta functions and the $p$-adic interpolation
of Eisenstein--Kronecker numbers}, arXiv:math/0610163v4,
Example 1.9, PDF pp.~10--11:
\[
 \theta(z)=\exp(-e_2^*z^2/2)\sigma(z).
\]
Here the discriminant $64$ is a $7$-adic unit, and $R$ is
unramified quadratic since $-1$ is not a square modulo $7$.
The automorphism $(x,y)\mapsto(-x,iy)$ acts on $\omega$ by
$i$, giving the asserted CM action and lattice symmetry.
The lattice is stable under multiplication by $i$.
For positive regularization parameter, the weight-two lattice
sum defining $e_2^*$ changes sign under $\gamma\mapsto i\gamma$
and converges absolutely, so is zero. Its continuation is
therefore zero also. Thus $\theta=\sigma$ for this model.
This uses the canonical reduced theta convention, not an
arbitrary exponential modification of a sigma function.

Write $x=t^{-2}u$. The curve equation gives
$u^2-u=t^4$, with $u(0)=1$. Since $y=-2t^{-3}u$,
\[
 \lambda'(t)=\frac{\omega}{dt}
       =1-\frac{tu'}{2u}
       =\frac1{2u-1}
       =(1+4t^4)^{-1/2}.
\]
Integrating the first three binomial terms gives
\[
 \lambda(t)=t-\frac25t^5+\frac23t^9+O(t^{13}).
\]
For completeness the sigma coefficients can be computed from
the curve itself. Put $x(z)=\wp(z)$ in the coordinate with
$\omega=dz$. The differential equation
$x''=6x^2-2$, and the $i$-symmetry, give
\[
 \wp(z)=z^{-2}+\frac15z^2+\frac1{75}z^6+O(z^{10}).
\]
For instance, comparison of the constant and $z^4$ terms
gives $2c_2=12c_2-2$ and
$30c_6=12c_6+6c_2^2$. The defining product for sigma gives
$-(\log\sigma)''=\wp$, with
$\sigma(z)/z=1+O(z^4)$. Consequently
\[
 \log\frac{\sigma(z)}z
      =-\frac{z^4}{60}-\frac{z^8}{4200}+O(z^{12}),
 \qquad
 \sigma(z)=z-\frac{z^5}{60}-\frac{z^9}{10080}+O(z^{13}).
\]
Substitution of the computed logarithm yields the coefficient
of $t^4$ as $-2/5-1/60=-5/12$ and that of $t^8$ as
\[
 \frac23+5\left(-\frac25\right)\left(-\frac1{60}\right)
           -\frac1{10080}=\frac{1411}{2016}.
\]
The denominator is $7\cdot288$ and the numerator is $4$
modulo $7$, proving \eqref{eq:ag-theta-nonintegral}.
All earlier coefficients are $7$-integral.

If this series is multiplied by $u(t)\in R[[t]]^\times$,
the first nonintegral coefficient becomes
$u(0)(1411/2016)$ plus an integral element, and still has
valuation $-1$. Under $t=a_1T+a_2T^2+\cdots$, with
$a_1\in R^\times$ and all $a_j\in R$, its first
nonintegral coefficient becomes
$a_1^8(1411/2016)$ plus an integral element. This also has
valuation $-1$. If the divisor equation is changed from $t$
to $T$, their ratio is itself an integral unit, so the same
conclusion applies to that normalization as well.
\end{proof}

This refutes the blanket transfer that a canonically normalized
reduced theta section on every good inert CM model is automatically
an integral formal function. It identifies the missing integral
trivialization in that transfer. It neither refutes A nor asserts
failure of any stated derivative-divisibility theorem in its
specified normalization. The stronger classical theta theory
may supply different normalizations or estimates. In particular,
the ordinary-prime hypothesis in Bannai--Kobayashi,
Propositions 2.14--2.15, PDF pp.~28--29, is an explicit extra
hypothesis for their formal integrality theorem; it is absent
at the inert prime in this example.

\subsection{A proved Fourier refinement of the supersingular divisibility}

The following argument uses the proved integral Fourier estimates of
Bannai--Kobayashi, not the announcement of a future inert-prime paper in
their earlier theta manuscript.  Its input is the Lubin--Tate distribution
estimate in \textbf{S-AG-BK-FOURIER}, Proposition 5.3(i),(ii), printed
pp.~544--545, together with the Fourier moment and torsion-averaging
identities used in Proposition 5.5, pp.~546--547.  These statements and
their complete displayed proofs have been read.  The exact period
valuation is given in their \S2, printed pp.~526--527.  No assertion
about the full sigma-function radius is substituted for these results.

\begin{proposition}[AG-FG-FOURIER: divisibility from an integral trace family]
Let $p\geq3$, $q=p^2$, $Q=q-1$ and $\alpha=p/Q$.  Let $G$ be a
Lubin--Tate formal group for the unramified quadratic extension of
$\mathbf Q_p$, with an integral invariant differential whose value at the
origin is a unit.  Normalize its derivation $D$ by the formal logarithm.
Let $R\subset\mathcal O_{\mathbf C_p}$ be a complete discrete valuation
ring containing the coefficients of $G$, with $v_p(R\setminus\{0\})
\subset e^{-1}\mathbf Z$.  Suppose that $f_j\in R[[t]]$, $j\geq0$,
satisfy the convergent formal identities
\[
 \frac1q\sum_{\tau\in G[p]} f_j(t\mathbin\oplus\tau)
       =f_{j+1}([p]t).                                      \tag{AG-Fourier-trace}
\]
For every $n\geq1$ put
\[
 b_e(n)=
 \begin{cases}
 \max\{0,e^{-1}\lceil e(n\alpha-1)\rceil\},&Q\nmid n,\\
 \max\{0,e^{-1}\lceil e(n-1)\alpha\rceil\},&Q\mid n.
 \end{cases}
\]
Then $v_p(D^nf_j(0))\geq b_e(n)$ for every $j$.  In particular, for
an unramified coefficient ring,
\[
             v_p(D^nf_j(0))\ \geq\ \left\lfloor\frac{pn}{p^2-1}\right\rfloor.
                                                               \tag{AG-Fourier-floor}
\]
An integral change of parameter or multiplication of the invariant
differential by a unit does not change these valuation assertions.
\end{proposition}

\begin{proof}
Choose a Lubin--Tate period $\varpi$ in the Fourier construction.  Its
valuation is
\[
 v_p(\varpi)=\frac1{p-1}-\frac1{p^2-1}=\alpha.
\]
For an integral series $f$, denote its Fourier distribution on
$\mathcal O_{\mathbf Q_{p^2}}$ by $\mu_f$.  This is a distribution;
it need not be a bounded measure.  Proposition 5.3 gives
\[
 \left|\int_{\mathcal O^\times}x^n\,d\mu_f\right|_p\leq p,
 \qquad
 \left|\int_{\mathcal O^\times}(x^m-x^n)\,d\mu_f\right|_p
       \leq p^{-l+\alpha}
 \quad(m\equiv n\bmod p^lQ).                              \tag{AG-Fourier-est}
\]
The indicator of $\mathcal O^\times$ corresponds under the transform to
subtracting the average over $G[p]$.  Differentiating its transform and
using (AG-Fourier-trace) therefore gives, for $a_{j,n}=D^nf_j(0)$,
\[
 c_{j,n}:=a_{j,n}-p^n a_{j+1,n}
       =\varpi^n\int_{\mathcal O^\times}x^n\,d\mu_{f_j}.    \tag{AG-Fourier-moment}
\]
In degree zero its mass is $f_j(0)-f_{j+1}(0)\in R$.  Consequently,
when $Q\mid n$, the second estimate with exponents $n,0$ and $l=0$
shows that the $n$th moment has valuation at least $-\alpha$:
the constant mass is integral, and the difference has that valuation
bound.  For $Q\nmid n$ the first estimate instead gives valuation
at least $-1$.  After multiplication by $\varpi^n$, these are respectively
$(n-1)\alpha$ and $n\alpha-1$.  Each $c_{j,n}$ belongs to $R$, so rounding
up in its actual value group gives the corresponding lower bound in the
definition of $b_e(n)$.  The additional bound zero follows directly from
$a_{j,n},a_{j+1,n}\in R$, since $D$ is an integral derivation.

Iterate (AG-Fourier-moment):
\[
 a_{j,n}=\sum_{i=0}^{N-1}p^{in}c_{j+i,n}+p^{Nn}a_{j+N,n}.
\]
The last term has valuation at least $Nn$, uniformly in $j$, while all
terms of the sum have valuation at least $b_e(n)$.  Letting $N$ tend to
infinity proves the claim.  For $e=1$ and $Q\nmid n$, $n\alpha$ is
nonintegral because $(p,Q)=1$, hence $\lceil n\alpha-1\rceil
=\lfloor n\alpha\rfloor$.  For $Q\mid n$, $n\alpha$ is an integer
and $0<\alpha<1$, hence $\lceil(n-1)\alpha\rceil=n\alpha$.
Finally an integral parameter isomorphism preserves the coefficient
lattice, and rescaling the differential by a unit multiplies an $n$th
derivative by a unit to the power $n$.
\end{proof}

\begin{proposition}[AG-THETA-FOURIER: the exact floor bound for the CM values]
Use the good full-CM models and the algebraic differential in
AG-CM-INTEGRAL-MODEL, at an inert prime $p\geq5$.  Write
$B_p(n)=\lfloor pn/(p^2-1)\rfloor$.
For a CM endomorphism $b$ with $p\nmid N(b)$ and every $m\geq3$,
\[
 v_p\bigl((1-b^m)\beta_m(E,\omega)\bigr)\geq B_p(m-2).
                                                               \tag{AG-Fourier-BH}
\]
For $P$ defined over an unramified extension with $\bar P\ne O$,
\[
                  v_p(D^n x(P))\geq B_p(n)\qquad(n\geq1).   \tag{AG-Fourier-point}
\]
Thus the conductor-one identity of AG-THETA-CLASS has the improved
lower bounds
\[
 v_p\bigl(w(m-1)!\Omega^{-m}L(m,\psi)\bigr)\geq
 \begin{cases}
 B_p(m-2),&p^2-1\nmid m,\\
 B_p(m-2)-1-v_p(m),&p^2-1\mid m.
 \end{cases}                                               \tag{AG-Fourier-class}
\]
For the proper prime-to-$p$ conductor identity of AG-THETA-RAY the
lower bound is $B_p(m-2)$, with no weight congruence exception.
For the mixed conductor $\mathfrak f=(p)^r\mathfrak f_0$ of
AG-THETA-PCOND, $\mathfrak f_0\ne\mathcal O_K$, let $e$ be the
ramification index of a field containing the displayed points.  The
valuation of its left side
$p^{rm}w(m-1)!\Omega^{-m}L(m,\psi)$ is at least
\[
                         \max\{d_2(m-2),b_e(m-2)\}.         \tag{AG-Fourier-mixed}
\]
The pure $p$-conductor pole loss remains the one in AG-THETA-PCOND.
\end{proposition}

\begin{proof}
Full CM is relevant here.  At an inert prime the formal
$\mathcal O_{K,p}$-module has relative height one.  Over the completed
maximal unramified extension it is isomorphic to a Lubin--Tate module,
with unit linear coefficient.  This is the unique-lifting statement
used in \textbf{S-AG-KRY-CM}, Proposition 5.1, Case 2: a fixed
height-one formal $\mathcal O_{K,p}$-module and its fixed action on the
Lie algebra have a unique lift.  Apply it to the special fibre of a
Lubin--Tate module and to that of our full-CM formal group after their
geometric identification.  The differential is a generator, so the
normalization discrepancy is a unit.  This step uses the actual CM
action; it asserts no such classification for arbitrary height-two
deformations.  The unramified extension still has integral value group,
so the rounding in AG-FG-FOURIER is preserved.

The classical multiplication identity, also Lemma 5.4 of
\textbf{S-AG-BK-FOURIER}, is
\[
          \sum_{T\in E[p]}x(Q+T)=p^2x([p]Q).
\]
It holds as an identity of rational functions and hence after every
formal substitution in question.  Since $[b]$ permutes $E[p]$, it
implies that the integral regularization
$f_b=x-b^2x\circ[b]$ of AG-THETA-REG satisfies
(AG-Fourier-trace) with the constant family $f_j=f_b$.
Its $m-2$ derivative is $(1-b^m)\beta_m$, proving
(AG-Fourier-BH).

For (AG-Fourier-point), use
$f_j(t)=x([p]^jP+_E t)$.  The supersingular multiplication map $[p]$
is purely inseparable on the special fibre, so its only geometric
point above $O$ is $O$.  Thus $[p]^j\bar P\ne O$ for all $j$.
Each $f_j$ is an integral formal function because $x$ is regular in
that formal neighbourhood, and the same multiplication identity is
exactly (AG-Fourier-trace).  Apply the proposition with $e=1$.
The argument over a ramified coefficient extension instead gives
$b_e(n)$; no rounding to an integer is made there.

If $Q\nmid m$, choose a CM element $b$ whose residue generates
$\mathbf F_{p^2}^{\times}$.  Then $1-b^m$ is a unit.  If $Q\mid m$,
take $b=1+p$ and use
$v_p((1+p)^m-1)=1+v_p(m)$.  These choices in (AG-Fourier-BH),
followed by the exact finite ideal-class identity already proved in
AG-THETA-CLASS, give (AG-Fourier-class).  All character coefficients
in that identity are units.  They may enlarge the coefficient field,
but do not weaken a valuation bound proved before forming the sum.

The points in AG-THETA-RAY have prime-to-$p$ order, are defined over
unramified extensions, and reduce away from $O$.  Apply
(AG-Fourier-point) to each term of its exact identity.  In the mixed
conductor case of AG-THETA-PCOND the points are possibly ramified,
but still reduce away from $O$ because $[p]^rP$ is nonzero
prime-to-$p$ torsion.  The ramified version just proved and the earlier
independent bound $d_2(m-2)$ both apply.  The exact lattice scaling
$p^{-rm}$ has already been accounted for on the displayed left side.
For pure $p$-conductor, the points reduce to the pole; the integral
trace-family hypothesis is then absent, which is why its earlier
explicit Laurent bound is retained.
\end{proof}

\paragraph{Geometric and method translation.}
The operator in (AG-Fourier-trace) is the normalized trace for the finite
flat CM isogeny $[p]$.  The Fourier identity turns its subtraction into
restriction of a locally analytic distribution to the units.  The
integral function lattice and the unramified value group then provide
the exact divisibility.  This is a concrete trace-to-arithmetic-lattice
comparison relevant to Alpöge step 2 and the Fable \S2 bridge; the
finite-flat isogeny and the integral formal functions are the additional
geometric inputs.  It supplies neither an integral canonical theta
trivialization contradicted by AG-THETA-NONINT nor a theorem about
every automorphic theta lift.  The statement is unconditional and
does not use A or M.  In the cote 1 framework it exemplifies a
geometric trace with a specified integral lattice, in addition to the
analytic trace/determinant statements already catalogued there.

\subsection{A geometric period at nonsplit CM primes}

Here is a proved nonsplit-prime result with its actual period, rather
than an unverified use of the supersingular sigma-radius announcement.
\textbf{S-AG-AI-CM} is Andreatta--Iovita, \emph{Katz type $p$-adic
$L$-functions for primes $p$ non-split in the CM field},
\emph{Commentarii Mathematici Helvetici} 99 (2024), 641--716.
We use its Proposition 5.9, printed pp.~689--690, and the analyticity
proved in Proposition 5.6 and Remark 5.8.  The algebraic normalization
below is the one in its \S2.4; its comparison with complex $L$-values
is the Damerell comparison for the Eisenstein case and the stated
Waldspurger comparison for the cuspidal case.

\paragraph{AG-SHIMURA-PERIOD: the precise classical input.}
Fix an imaginary quadratic field $K$, a level $N\geq5$ satisfying the
Heegner condition $\mathcal O_K/\mathfrak N\simeq\mathbf Z/N\mathbf Z$,
and a nebentype $\epsilon$.  Let $p\geq5$ be inert in $K$ and
$p\nmid N$.  Put $c=p^n d$, $n\geq2$, with $d$ odd and prime to
$pN\operatorname{disc}(K)$.  Let $F$ and $\chi$ satisfy the hypotheses
of \S2.4 of the source: $\chi$ is of its type $(c,\mathfrak N,\epsilon)$,
has infinity type $(k+m,-m)$, $m\geq0$, and has exact $p$-part of
the finite conductor $p^n$.  In the Eisenstein case $F=E_{k,\epsilon}$
is the series of \S2.3.1, $\epsilon$ has parity $k$ and is nontrivial
when $k=2$.  In the cuspidal case $k\geq2$, the discriminant of $K$
is odd, $\chi|_{\mathbf A_{\mathbf Q}^{\times}}=\epsilon\mathbf N^k$,
and the finite local epsilon conditions of its Theorem 2.4 are all
$+1$.  These are retained hypotheses, not supplied by A or M.

Choose the CM elliptic curve of conductor $d$, a generically cyclic
subgroup $C^{(n)}\subset E[p^n]$ of order $p^n$, and
$E^{(n)}=E/C^{(n)}$, as in \S5.1.  Its endomorphism order has
conductor $p^n d$.  Write $\omega^{(n)}$ for an integral generator of
its invariant differentials, and use the isogeny-normalized
differentials on the ideal-class translates.  The algebraic number is
\[
 L_{\rm alg}(F,\chi)=
 \sum_{[\mathfrak a]\in\operatorname{Pic}(\mathcal O_c)}
   \chi_m(\mathfrak a)^{-1}\,
   \delta_k^mF\bigl(\mathfrak a*(E^{(n)},\psi_N^{(n)},
                                      \omega^{(n)})\bigr),
 \qquad \chi_m=\chi\mathbf N^m.                         \tag{AG-AI-algebraic}
\]
The operator $\delta_k^m$ is the weight-$k+2m$ operator obtained from
the Gauss--Manin connection and the CM splitting of the Hodge
filtration.  The source constructs a locally analytic function
$L_p(F,-)$ on its character domain and a period $\Omega_{p,n}$ such
that
\[
 L_{\rm alg}(F,\chi)=\Omega_{p,n}^{k+2m}L_p(F,\chi),
 \qquad
 v_p(\Omega_{p,n})=\frac1{p^{n-1}(p^2-1)}.               \tag{AG-AI-period}
\]
For a varying modular weight, the source requires $n\geq3$ at inert
$p\geq5$; the threshold $n\geq2$ above is for a fixed classical
weight.  This distinction is Definition 4.24, pp.~681--682.

\begin{proof}[Normalization and the source proof]
Lemma 3.1 of the source gives
\[
 v_p\bigl(\operatorname{Hdg}(E^{(n)})\bigr)
       =\frac1{p^{n-1}(p+1)}.
\]
The canonical subgroup supplies a distinguished differential
$\omega_n$ in the canonical Hodge sublattice, with
$\omega_n=\Omega_{p,n}\omega^{(n)}$.  By the definition of that
sublattice, its valuation relative to an integral generator is
$v_p(\operatorname{Hdg})/(p-1)$; this proves the asserted period
valuation.  Lemma 5.1 and Definition 5.2 use the isogeny through
$E/C^{(n)}[p]$ to give the required integral Hodge projection before
evaluation.  Thus the period specifies a lattice comparison, not
merely a nonzero complex scaling constant.

The $p$-depleted modular form, with the source's precise operator
convention in Corollary 4.23, is
\[
 F^{[p]}=F-VUF
       =F-a_pVF+\epsilon(p)p^{k-1}V^2F,
\]
where $UV=1$.  Definition 5.7 defines $L_p$ by the finite CM sum
for $\delta_k^mF^{[p]}$ evaluated with $\omega_n$, divided by
$| (\mathcal O_K/\mathfrak N)^\times |$ when its larger adelic
class set $H(c,\mathfrak N)$ is used.  Keep this factor: it is what
identifies the adelic sum with (AG-AI-algebraic).
The $V$ and $V^2$ terms factor through conductor $p^{n-1}d$ or
$p^{n-2}d$.  Their $\chi$-weighted sums vanish, because the restriction
of the primitive conductor-$p^n$ character to the order-$p$ kernel
of the preceding class-group map is nontrivial.  This is the last
paragraph of the complete proof of Proposition 5.9.  The remaining
sum is the one for $F$.  Replacing $\omega_n$ by
$\omega^{(n)}$ multiplies that weight-$k+2m$ evaluation by
$\Omega_{p,n}^{k+2m}$, giving (AG-AI-period).
Finally, Proposition 5.6 proves analyticity in local weight
coordinates by applying the Hodge projection to the constant
coefficient in the marked-section expansion; Remark 5.8(1) gives
the fixed-weight version used here.  The underlying overconvergent
Gauss--Manin iteration is the source's cited construction, not a
new assertion proved from GRH.
\end{proof}

\begin{proposition}[AG-SHIMURA-SLOPE: a uniform divisibility and interpolation test]
Fix the data of AG-SHIMURA-PERIOD and a closed analytic weight disk
$\mathcal D$ contained in one of its character branches on which
$L_p(F,-)$ is represented by a convergent power series.  There is a
finite constant $C_{F,\mathcal D}\geq0$, independent of $m$, such
that every classical point $\chi\in\mathcal D$ of type $(k+m,-m)$
satisfies
\[
 v_p L_{\rm alg}(F,\chi)
       \geq\frac{k+2m}{p^{n-1}(p^2-1)}-C_{F,\mathcal D}.  \tag{AG-AI-slope}
\]
For any sequence of such points converging in the weight topology,
with ordinary integer weights $m\longrightarrow+\infty$, the
unnormalized values tend to zero $p$-adically.  If the limiting
classical value is nonzero, these unnormalized values have no
continuous interpolant agreeing at that limit.  In the Eisenstein
case every classical point has a sequence with unchanged local-unit character data
and exact finite conductor of this form.
\end{proposition}

\begin{proof}
Write the convergent series on the chosen closed disk as
$L_p(F,T)=\sum a_j(T-T_0)^j$, of radius $r$.  Its Gauss bound
$M=\sup_j|a_j|_p r^j$ is finite.  If it is nonzero, take
$C_{F,\mathcal D}=\max\{0,\log_p M\}$; if it is identically zero,
take $C_{F,\mathcal D}=0$ and interpret $v_p(0)=+\infty$.
Then $v_pL_p(F,\chi)\geq-C_{F,\mathcal D}$ everywhere on the
disk.  Taking valuations in (AG-AI-period) proves (AG-AI-slope).
Its right side tends to infinity with $m$.  Continuity at the
limit would force the limiting unnormalized value to be zero,
which proves the claimed obstruction when that value is nonzero.

For the last assertion let $h$ be the class number and $w$ the
number of roots of unity of $K$, and put $H=hw$.  For a fractional
ideal $\mathfrak a$ choose a generator $\alpha$ of $\mathfrak a^h$
and put $\psi(\mathfrak a)=\alpha^{-w}$.  Changing the generator
multiplies it by a $w$th root of unity, so this is well defined;
multiplication of generators also proves multiplicativity.
Moreover $\psi((z))=z^{-H}$.  In the source's adelic infinity-type
convention this is a conductor-one algebraic Hecke character of
type $(H,0)$.  The character
$\eta=\psi/\overline\psi$ has type $(H,-H)$ and conductor one.
On ideals prime to $p$ its values are units of the unramified field
$K_p$.  Since its residue field has order $p^2$,
\[
      \eta^{(p^2-1)p^j}\longrightarrow1
\]
uniformly in the unit-valued $p$-adic character coordinates.  This
also follows directly from $v_p(u^{p^2-1}-1)\geq1$ for
$u\in\mathcal O_{K_p}^{\times}$ and successive $p$th powers.
For a fixed classical $\chi_0$, set
$\chi_j=\chi_0\eta^{(p^2-1)p^j}$.  Its restrictions to the local unit groups and its exact
finite conductor are unchanged, while its weight is
$m_j=m_0+H(p^2-1)p^j$.  Hence $\chi_j\to\chi_0$ and
$m_j\to+\infty$.  In the Eisenstein case these characters retain
all the stated admissibility conditions.  They eventually belong
to any sufficiently small closed analytic weight disk about
$\chi_0$, proving the final assertion.  We impose no assertion here
that a cuspidal local epsilon condition is preserved by an
arbitrary deformation without checking that condition.
\end{proof}

\paragraph{Place in the programme.}
The geometric object is the Hodge bundle of a modular curve at a CM
point of nonmaximal order, with its canonical subgroup and the
Gauss--Manin connection.  The relevant comparison is between an
integral differential and a canonical Hodge sublattice after a cyclic
$p$-power isogeny.  Its positive, explicitly computed valuation
explains both the divisibility and the normalization required for
$p$-adic interpolation.  This gives an actual Shimura-geometric
continuation of the theta/divisibility discussion in Alpöge step 2
and Fable \S2.  It is compatible with the formal-trivialization
obstruction in AG-THETA-NONINT.  Neither the interpolation theorem
nor its slope consequence uses A or M.  In the cuspidal case the
complex value is the specific Rankin value in the source's
Waldspurger theorem; it is not thereby assigned Hecke GRH or
Artin holomorphy.

\begin{proposition}[AG-SHIMURA-RAM: the ramified period and its exact slope]
Keep the Heegner, nebentype, parity and, where applicable, cuspidal
epsilon hypotheses of AG-SHIMURA-PERIOD, but let $p\geq5$ be
ramified in $K$, with unique prime $\mathfrak p$ above $p$.
Let $c=p^n d$, where $n\geq2$ and $d$ is odd and prime to
$pN\operatorname{disc}(K)$.  Restrict to characters whose exact
$p$-part of the conductor ideal is
$(p)^n=\mathfrak p^{2n}$.  No assertion about an odd
$\mathfrak p$-exponent is implicit here.

Start with $E$ having CM by $\mathcal O_d$ over the integers of
a finite extension of $\mathbb Q_p$.  Choose a subgroup
$C^{(n)}\subset E[p^n]$ which is cyclic of order $p^n$ on the
generic fibre, and whose generic-fibre $p$-torsion line differs
from $E[\mathfrak p]$.  Put $E^{(n)}=E/C^{(n)}$.
Its CM order has conductor $p^n d$.  With an integral invariant
differential $\omega^{(n)}$, the canonical differential is
$\omega_n=\Omega_{p,n}\omega^{(n)}$, and
\[
 v_p(\Omega_{p,n})=\frac1{2p^n(p-1)},\qquad
 L_{\rm alg}(F,\chi)=\Omega_{p,n}^{k+2m}L_p(F,\chi).
                                                        \tag{AG-AI-ram}
\]
Here $L_{\rm alg}$ uses precisely (AG-AI-algebraic), and $L_p$
uses the source's ramified Hodge projection.  For every closed
analytic character disk $\mathcal D$ in its domain there is a
finite $C_{F,\mathcal D}\geq0$ such that
\[
 v_p L_{\rm alg}(F,\chi)
       \geq\frac{k+2m}{2p^n(p-1)}-C_{F,\mathcal D}.       \tag{AG-AI-ram-slope}
\]
The threshold $n\geq2$ works both for a fixed classical modular
weight and for the source's analytic finite-slope modular families
at $p\geq5$.  For a fixed Eisenstein series, every classical point
has an explicit sequence of admissible characters converging to
it while $m\to+\infty$, preserving its local-unit character data
and exact conductor.  Consequently a nonzero algebraic value at
the limit again obstructs continuous interpolation without the
period normalization.
\end{proposition}

\begin{proof}
We first spell out the transversality condition and the CM order.
Over an algebraic closure the $p$-adic CM Tate lattice of $E$ is
free of rank one over $\mathcal O_{K_{\mathfrak p}}$.
Thus identify its $p^n$-torsion with
$p^{-n}\mathcal O_{K_{\mathfrak p}}/\mathcal O_{K_{\mathfrak p}}$.
Write a generator of the chosen cyclic group as
$p^{-n}\beta$.  The line $E[\mathfrak p]$ in $E[p]$ is
the nonunit line $\mathfrak p/(p)$, after this identification.
Our transversality condition is therefore exactly that $\beta$
is a unit.  It can be achieved after a finite extension by
choosing any such generator and taking its finite flat closure.
The quotient lattice is
\[
 \mathcal O_{K_{\mathfrak p}}+\mathbb Z_p p^{-n}\beta
 \sim \mathcal O_{K_{\mathfrak p}}+p^{-n}\mathbb Z_p
 =p^{-n}(\mathbb Z_p+p^n\mathcal O_{K_{\mathfrak p}}),
\]
where the homothety uses $\beta^{-1}$.  The multiplier ring of
the last lattice is exactly the order
$\mathbb Z_p+p^n\mathcal O_{K_{\mathfrak p}}$: membership is
necessary by applying a multiplier to $1$, and sufficient since
this is a ring.  The prime-to-$p$ local orders are unchanged by
the isogeny.  This proves the claimed global conductor.
The transversality condition is on generic fibres; it must not
be read as disjointness of the special fibres of finite flat
group schemes at this supersingular point.

Andreatta--Iovita, Lemma 3.2, printed pp.~659--660, gives
\[
 E[\mathfrak p]\text{ canonical},\qquad
 v_p\operatorname{Hdg}(E)=\frac12,\qquad
 v_p\operatorname{Hdg}(E^{(n)})=\frac1{2p^n}.
\]
Its proof uses the $p$-divisible CM endomorphism given by a
uniformizer of $K_{\mathfrak p}$ and the Katz--Lubin canonical
subgroup theorem.  Lemma 6.1, printed p.~692, supplies the
integral Hodge projection needed for evaluation.  Unlike the
inert case it includes the different, whose $p$-valuation is
$1/2$.  Its reduced numerical inequality is
$3p^2-p-1\leq p^3$, which holds for the present $p\geq5$.
The canonical Hodge sublattice has valuation
$v_p\operatorname{Hdg}(E^{(n)})/(p-1)$ relative to an integral
generator.  This proves the period valuation.

Propositions 6.4 and 6.6, printed pp.~693--694, establish
analyticity and (AG-AI-ram), respectively; their proofs reduce
to the full arguments of Propositions 5.6 and 5.9, with the
ramified projection just described.  Retain the precise
depletion $F^{[p]}=F-VUF$, the adelic class-sum denominator,
and primitive-character cancellation as in our preceding
normalization proof.  Definition 4.24 gives $n\geq2$ in the
ramified case even for the varying modular weight.
The finite Gauss bound for $L_p$ on $\mathcal D$, exactly as
in AG-SHIMURA-SLOPE, now gives (AG-AI-ram-slope).

For completeness use the same conductor-one character
$\eta=\psi/\overline\psi$ of type $(H,-H)$, where $H=hw$,
constructed in AG-SHIMURA-SLOPE.  Its unit values now lie in
the ramified quadratic field $K_{\mathfrak p}$, whose residue
field is $\mathbb F_p$.  Every unit $u$ satisfies
$v_p(u^{p-1}-1)\geq1/2$.  If $v_p(a)\geq1/2$ and $p\geq5$,
the binomial expansion gives
\[
 v_p((1+a)^p-1)\geq v_p(a)+1:
\]
the intermediate binomial coefficients are divisible by $p$,
and $p v_p(a)\geq v_p(a)+1$.  Induction proves uniformly that
$v_p(u^{(p-1)p^j}-1)\geq j+1/2$.
Consequently
\[
 \chi_j=\chi_0\eta^{(p-1)p^j}\longrightarrow\chi_0,
 \qquad m_j=m_0+H(p-1)p^j\longrightarrow+\infty.
\]
Conductor-one twisting changes no local-unit restrictions or
finite conductor, and in the Eisenstein case preserves all the
stated admissibility conditions.  The sequence eventually lies
in the chosen disk, so (AG-AI-ram-slope) proves the final claim.
No integrality rounding is used: the differential and period
may lie over ramified extensions and have fractional valuation.
\end{proof}

The ramified construction concerns the same Hodge-bundle and
Gauss--Manin comparison as the inert construction, but the
different and the canonical CM $\mathfrak p$-subgroup change
its valuation.  It is an unconditional geometric contribution
to the same Alpöge-step-2 and Fable-\S2 normalization gate,
not a consequence of A or M.

\subsection{Arthur blocks, the Lefschetz action, and the integral lattice}
\label{ag:arthur}

The supplied Hodge bridge sketch, Route 1.4 and Route 4, compares
Alpöge's added blocks $\mu[t]$ with Lefschetz strings and with the
integral ring of cote 16, scans 14--18, no.~7. The following
statements locate the classical comparison and separate the
additional representation and lattice calculations. The archive
locators here are inherited from that sketch and from
\path{RETURN_DP_MULT_cp01_IMPLICATIONS.md}, Section 6,
DP15, rather than newly checked manuscript images.

\paragraph{AG-ARTHUR-ARCH (the exact classical comparison).}
Arthur, \emph{Unipotent automorphic representations: conjectures},
Astérisque 171--172 (1989), 13--71, Section 9,
Proposition 9.1, printed pp.~62--66 (Clay PDF pp.~50--54),
compares the complex cohomology of the cohomological real packet
of Section 5 with the representation $r_\mu\circ\psi_\infty$.
Here $r_\mu$ is the minuscule representation specified by the
Shimura cocharacter. The equivalence respects the component-group
action, $\mathbb C^\times$, and the Lefschetz $SL_2$.
The Lefschetz action is constructed on printed pp.~60--61.
This is a statement about the sum of relative Lie algebra
cohomologies in the specified packet, after applying $r_\mu$;
it is not a comparison of integral lattices. The global
decomposition (9.1), printed p.~58, explicitly assumes
Conjecture 8.1 in that source. These exact local and global
scopes replace the sketch's unqualified identification.
The entire proof of Proposition 9.1 has been read: it compares
the $\mathbb C^\times$ and component-group characters, then the
Lefschetz weights using Vogan--Zuckerman Proposition 6.19.
That underlying theorem is cited, not re-proved here.

\begin{proposition}[AG-ARTHUR-PARAM: the two $SL_2$ factors]
Let $k$ be nonarchimedean with residue cardinality $q$, and choose
geometric Frobenius $\Phi$ with $|\Phi|=q^{-1}$. Let $\mu$ be an
unramified unitary character with $\mu(\Phi)=\alpha$, and let
$t\ge1$. An Arthur parameter trivial on the Deligne $SL_2$
and equal to $\mu\boxtimes\operatorname{Sym}^{t-1}$ on the
Arthur factor has associated Weil--Deligne representation
\[
 r(\Phi)e_j=\alpha q^{j-(t-1)/2}e_j\quad(0\le j<t),
 \qquad N=0.
\]
Its local factor is
\begin{equation}\label{ag:arthur-factor}
 L(s,r,0)=\prod_{j=0}^{t-1}
       (1-\alpha q^{-(s+(t-1)/2-j)})^{-1}.
\end{equation}
On the same underlying Weil representation, the monodromy chain
$Ne_0=0$, $Ne_j=e_{j-1}$ for $j\ge1$ is a different
Weil--Deligne representation when $t>1$, and has factor
\[
 L(s,r,N)=(1-\alpha q^{-(s+(t-1)/2)})^{-1}.
\]
\end{proposition}
\begin{proof}
Gan--Gross--Prasad, \emph{Branching laws for classical groups:
the non-tempered case}, arXiv:1911.02783v2, Section 2,
PDF p.~8, define the separate Deligne and Arthur factors.
The associated Langlands parameter inserts
$\operatorname{diag}(|w|^{1/2},|w|^{-1/2})$ into the Arthur
factor. On the monomials $x^{t-1-j}y^j$ its eigencharacters
are $\mu|\cdot|^{(t-1)/2-j}$. Triviality of the Deligne
factor gives $N=0$, proving the first description.
For the second description a direct calculation gives
$r(\Phi)Nr(\Phi)^{-1}=q^{-1}N$, the required
Weil--Deligne relation; inertia is trivial.
The local factor is, by definition,
\[
 \det(1-q^{-s}r(\Phi)\mid(\ker N)^I)^{-1}.
\]
For $N=0$ the kernel is the whole space; for the chain it is
$\mathbb C e_0$. This gives the two displayed answers.
The definition is Deligne, \emph{Les constantes des équations
fonctionnelles des fonctions $L$}, LNM 349 (1973),
Section 8.12, printed p.~572 (IAS PDF p.~72);
Wedhorn, arXiv:math/0011210v2, (3.2.1), PDF p.~37,
states the same convention explicitly.
\end{proof}

If a global unitary Hecke character $\mu$ is given, the formal
product of the factors in \eqref{ag:arthur-factor}, at the places
where these hypotheses hold, is
\[
 \prod_{j=0}^{t-1}L^S(s+(t-1)/2-j,\mu).
\]
M therefore places the zeros inherited from nontrivial zeros of
these constituents on the shifted lines $\Re s=1-t/2+j$;
for a Dirichlet character over $\mathbb Q$, A supplies this
deduction. Known poles and trivial zeros are kept separate.
Omitted Euler factors have their explicit divisors as in AG-LOCAL.
This is a product of shifted constituents, not a new pure
weight assertion. In geometric terms, weights of an associated
graded object do not determine its monodromy operator or the
$\ker N$ factor at a degeneration. No occurrence of an arbitrary
formal parameter in a theta lift is asserted.

\begin{proposition}[AG-ARTHUR-MIN: a block can change after $r_\mu$]
In characteristic zero,
\[
 \bigwedge\nolimits^2\operatorname{Sym}^3(\mathbb Q^2)
       \simeq\operatorname{Sym}^4(\mathbb Q^2)\oplus\mathbb Q
\]
as $SL_2$ representations. Thus the second fundamental
representation of $GL_4$, applied to a standard block of length
$4$, gives strings of lengths $5$ and $1$.
\end{proposition}
\begin{proof}
Set $v_j=x^{3-j}y^j$ for $0\le j\le3$ and use
$E=x\partial_y$, $F=y\partial_x$,
$H=x\partial_x-y\partial_y$. Let $w=v_0\wedge v_1$.
It has $H$-weight $4$ and $Ew=0$. The successive nonzero
vectors generated by $F$ are
\[
 \begin{split}
 w,&\qquad 2v_0\wedge v_2,\\
 2v_0\wedge v_3+6v_1\wedge v_2,&\qquad
 12v_1\wedge v_3,\qquad24v_2\wedge v_3.
 \end{split}
\]
Their weights are $4,2,0,-2,-4$, so they are independent;
their span is stable under $E,F,H$ and is the irreducible
highest-weight-$4$ module. For example the commutator identities
give $EF^jw=j(5-j)F^{j-1}w$, and $F^5w=0$.
The vector
$u=v_0\wedge v_3-3v_1\wedge v_2$ is killed by $E,F,H$.
It is independent of the preceding weight-zero vector.
The two submodules therefore have dimensions $5$ and $1$,
with direct sum the six-dimensional exterior square. Finite
rational $SL_2$ representations are determined here by their
Lie algebra action, or directly by exponentiating the nilpotent
operators, so this is the asserted group decomposition.
\end{proof}

The example is an explicit representation calculation; it does
not claim a new automorphic occurrence. Its geometric use is
precise: a Shimura-cohomology calculation must first apply the
representation $r_\mu$ appropriate to that Shimura datum.
Even a correct equality of standard Arthur parameters does not
allow one to omit this step. For a signature where $r_\mu$ is
the standard representation, the obstruction illustrated here
is absent, but the integral comparison below is still required.

\begin{proposition}[AG-LEFSCH-LATTICE: the rational string loses indices]
There are two graded, integral $\mathfrak{sl}_2$-stable lattices
in the same rational irreducible module, with identical
Frobenius characteristic polynomials in every degree and every
power, for which the extreme Lefschetz square has cokernel
$\mathbb Z/2$ and $0$, respectively.
\end{proposition}
\begin{proof}
On $V=\operatorname{Sym}^2(\mathbb Q^2)$ use the operators
$E,F,H$ just defined and set $m_j=x^{2-j}y^j$. Consider
\[
 M=\mathbb Zm_0\oplus\mathbb Zm_1\oplus\mathbb Zm_2,
 \qquad
 N=\mathbb Zm_0\oplus\mathbb Z(2m_1)\oplus\mathbb Z(2m_2).
\]
Both are stable: on $M$,
$Fm_0=2m_1$, $Fm_1=m_2$,
$Em_1=m_0$, $Em_2=2m_1$.
Writing $n_0=m_0,n_1=2m_1,n_2=2m_2$, on $N$ one has
$Fn_0=n_1$, $Fn_1=n_2$,
$En_1=2n_0$, $En_2=2n_1$.
The $H$-weights in both bases are $2,0,-2$.
Consequently $F^2:M_2\to M_{-2}$ multiplies by $2$,
whereas $F^2:N_2\to N_{-2}$ is an isomorphism.
Here subscripts denote $H$-weight, not the basis index.
For any positive integer $q$, the endomorphism
$\Phi_qm_j=q^jm_j$ preserves both lattices, satisfies
$\Phi_qF\Phi_q^{-1}=qF$, and has the same eigenvalues
$1,q,q^2$ in both. The same holds for all powers and for
restrictions to the graded pieces. Finally $N$ is not stable
under $F^{(2)}=F^2/2!$, since $F^{(2)}n_0=n_2/2$.
Thus even integral Lie operators do not automatically supply
an integral divided-power structure.
\end{proof}

This is an abstract integral-lattice counterexample, not a claim
that either lattice has been realized as the cohomology of a
specified variety. It establishes why rational parameter
recognition and Euler factors cannot determine the theta-series
index or the integral correspondence ring sought in Route 4.

\begin{proposition}[AG-LEFSCH-SPLIT: an exact sufficient lattice criterion]
Let $V$ be a finite rational $\mathfrak{sl}_2$ module with
\[
 V=\bigoplus_{m\ge0}
       \operatorname{Sym}^m(\mathbb Q^2)\otimes W_m.
\]
Suppose a lattice $M\subset V$ splits as
$\bigoplus_m M^{(m)}$ in these isotypic summands, each
$M^{(m)}$ splits into its $H$-weight lattices, and all
$E^{(a)}=E^a/a!$, $F^{(a)}=F^a/a!$ preserve each
isotypic lattice $M^{(m)}$.
For $r\ge0$ write $M_{\pm r}$ for the weight lattices.
Then
\begin{equation}\label{ag:integral-lefschetz-test}
 \operatorname{coker}
  (E^{(r)}:M_{-r}\longrightarrow M_r)
 \simeq
 \bigoplus_{\substack{m\ge r\\m\equiv r\ (2)}}
 \left(\mathbb Z/
     \binom{(m+r)/2}{r}\mathbb Z\right)^{\dim W_m}.
\end{equation}
For $E^r$ the cyclic orders are multiplied by $r!$.
Thus the inverse of $E^{(r)}$ has denominator dividing the
least common multiple of the displayed binomial coefficients.
\end{proposition}
\begin{proof}
The finite sums $\exp(E)=\sum_aE^{(a)}$ and
$\exp(-F)=\sum_a(-1)^aF^{(a)}$ preserve the lattice,
as do their inverses. Hence
$w=\exp(E)\exp(-F)\exp(E)$ is an integral automorphism.
On the standard two-dimensional module it sends $x$ to $-y$
and $y$ to $x$. It therefore carries weight $-r$ to weight
$r$ isomorphically on each $M^{(m)}$.
For $j=(m+r)/2$ the weight-$-r$ subspace of
$\operatorname{Sym}^m\otimes W_m$ is
$x^{m-j}y^j\otimes W_m$, and
\[
 E^{(r)}(x^{m-j}y^j)
   =\binom{j}{r}x^jy^{m-j},\qquad
 w(x^{m-j}y^j)=(-1)^{m-j}x^jy^{m-j}.
\]
Thus $E^{(r)}$ on this subspace is the displayed binomial
coefficient times an integral isomorphism, up to sign.
Its cokernel is the corresponding cyclic group repeated
$\dim W_m$ times. The assumed integral splitting permits the
direct sum over $m$. Multiplication by $r!$ gives the ordinary
power statement. The denominator assertion follows by applying
the integer annihilating this finite cokernel.
\end{proof}

For an extreme string ($m=r$), the divided power is unimodular
under these hypotheses, while the ordinary power contributes
$r!$. For intermediate weights the binomial factor can remain
nontrivial. This agrees with the inherited DP15 obstruction:
in dimension $3$, exterior divided powers are integral, but
the $H^2\to H^4$ paired block $J-I$ has cokernel
$\mathbb Z/2$. The general H8/Wilson diagonal result and
H8(iv)'s separate order/index claim retain their reported
status; they are not re-proved or silently identified with an
automorphic lattice here.

\begin{proposition}[AG-LEFSCH-INDEX: the index does not require integral splitting]
Keep the rational module $V=\bigoplus_m\operatorname{Sym}^m
(\mathbb Q^2)\otimes W_m$ and the operators above. Let
$M\subset V$ be any full lattice and put $M_j=M\cap V_j$,
without assuming that $M$ splits into its weight or isotypic
lattices. Fix $r\ge0$, and suppose that
\[
 f=E^{(r)}\text{ sends }M_{-r}\text{ into }M_r,
 \qquad w(M_{-r})=M_r,
 \quad w=\exp(E)\exp(-F)\exp(E).
\]
These two conditions hold, in particular, when all divided
powers of $E$ and $F$ preserve $M$.
Write
\[
 \begin{gathered}
 \mathcal D_r=\{m\ge r:m\equiv r\pmod2,\ W_m\ne0\},\\
 a_m=\binom{(m+r)/2}{r},\qquad d_m=\dim W_m,
 \qquad h_r=\sum_{m\in\mathcal D_r}d_m.
 \end{gathered}
\]
Then
\begin{equation}\label{eq:ag-lefschetz-index}
 [M_r:fM_{-r}]=\prod_{m\in\mathcal D_r}a_m^{d_m}.
\end{equation}
For the ordinary power $E^r$ the index is multiplied by
$(r!)^{h_r}$. Moreover, the exponent of $\operatorname{coker}f$,
equivalently the least denominator of $f^{-1}$ between these
lattices, divides $\prod_{m\in\mathcal D_r}a_m$.
Empty products are $1$.
\end{proposition}
\begin{proof}
The intersection of a full lattice with a rational subspace
is a full lattice in that subspace. Thus the two weight lattices
have rank $h_r$. The stated Weyl condition makes
$A=w^{-1}f$ an integral endomorphism of $M_{-r}$.
The monomial calculation in AG-LEFSCH-SPLIT gives its action on
the weight-$-r$ part of the $m$th rational isotypic summand:
it is multiplication by the nonzero integer
\[
 b_m=(-1)^{(m-r)/2}a_m.
\]
Consequently $f$ is injective with finite cokernel, and the
absolute determinant of $A$ is the right side of
\eqref{eq:ag-lefschetz-index}. The determinant-index formula
for an injective endomorphism of a free abelian group proves
the assertion, because $w$ is a lattice isomorphism.
Replacing $f$ by $r!f$ proves the ordinary-power formula.

For the exponent claim, the integral polynomial
$P(T)=\prod_{m\in\mathcal D_r}(T-b_m)$ annihilates $A$:
this is checked on each rational isotypic summand, where $A$
is scalar. Write $P(T)=P(0)+TQ(T)$ with $Q\in\mathbb Z[T]$.
Then
$P(0)M_{-r}=-A Q(A)M_{-r}\subseteq A M_{-r}$.
The absolute value of $P(0)$ is $\prod_m a_m$, proving the
annihilator and denominator statements. If the weight space
is zero the claims hold directly.
\end{proof}

\begin{proposition}[AG-LEFSCH-DEFECT: exact gluing and its possible primes]
Under the hypotheses of AG-LEFSCH-INDEX, put
\[
 N_{\pm r}=\bigoplus_{m\in\mathcal D_r}
       (M_{\pm r}\cap V^{(m)}),\qquad
 Q_{\pm r}=M_{\pm r}/N_{\pm r},
\]
where $V^{(m)}=\operatorname{Sym}^m(\mathbb Q^2)\otimes W_m$.
The groups $Q_{\pm r}$ are finite. Let
$\bar f:Q_{-r}\to Q_r$ be the induced map. There is an exact
sequence of finite groups
\begin{equation}\label{eq:ag-lefschetz-gluing}
 0\longrightarrow\ker\bar f
 \longrightarrow\bigoplus_{m\in\mathcal D_r}
        (\mathbb Z/a_m\mathbb Z)^{d_m}
 \longrightarrow\operatorname{coker}f
 \longrightarrow\operatorname{coker}\bar f
 \longrightarrow0.
\end{equation}
The first and last groups have equal orders. If a prime $\ell$
divides at most one of the integers $a_m$, the $\ell$-primary
cokernel is already forced: it is zero if no $a_m$ is divisible
by $\ell$, and otherwise it is
\[
 (\mathbb Z/\ell^{v_\ell(a_m)}\mathbb Z)^{d_m}
\]
for that unique $m$. Thus ambiguity beyond the split cyclic
formula can occur only at primes dividing two of its binomial
factors.
\end{proposition}
\begin{proof}
Each isotypic intersection is a full lattice in its rational
weight subspace, so their direct sum has finite index in the
whole weight lattice. The Weyl element preserves the rational
isotypic summands and identifies $N_{-r}$ with $N_r$, as well
as $M_{-r}$ with $M_r$. The operator $f$ acts on the $m$th
intersection as $b_m w$. It follows that it preserves these
two exact lattice sequences and has the split cokernel on
$N_{-r}\to N_r$ displayed in
\eqref{eq:ag-lefschetz-gluing}. The snake lemma, using
injectivity on both $N_{-r}$ and $M_{-r}$, gives that exact
sequence. Since $w$ identifies $Q_{-r}$ and $Q_r$, they have
the same finite order, so $\ker\bar f$ and
$\operatorname{coker}\bar f$ have the same order as well.

For the primary assertion use $A=w^{-1}f$ on
$M_{-r}\otimes\mathbb Z_\ell$. If all $b_m$ are units,
its determinant is a unit and its cokernel is zero.
If only $b_{m_0}$ is a nonunit, then every
$b_{m_0}-b_m$ for $m\ne m_0$ is a unit in $\mathbb Z_\ell$.
The polynomial projector
\[
 P_{m_0}(A)=\prod_{\substack{m\in\mathcal D_r\\m\ne m_0}}
       \frac{A-b_m}{b_{m_0}-b_m}
\]
therefore preserves this lattice. On the rational weight
space it projects onto $V^{(m_0)}_{-r}$, so it splits the
lattice into that intersection and a complementary summand.
On the first summand $A=b_{m_0}$ and the rank is $d_{m_0}$;
on the complement the determinant is a unit, so $A$ is an
integral automorphism. This gives exactly the asserted
$\ell$-primary group. For $r=0$ all $a_m=1$ and the first
case applies. This also avoids any issue with repeated unit
eigenvalues in that case.
\end{proof}

\paragraph{Sharpness of the denominator under the stated weaker hypotheses.}
Take $V=\operatorname{Sym}^3(\mathbb Q^2)\oplus
\operatorname{Sym}^7(\mathbb Q^2)$ and $r=1$.
Let $u=xy^2$ and $v=x^3y^4$ be the two weight-$-1$ monomials,
one in each summand. Then $A=w^{-1}E$ acts by $-2$ on $u$
and $-4$ on $v$. Choose
\[
 M_{-1}=\mathbb Zu+\mathbb Z\frac{u+v}{2},\qquad
 M_1=wM_{-1},
\]
and complete these to a full lattice using monomial lattices
in all other weights. In the displayed basis of $M_{-1}$,
\[
 A=\begin{pmatrix}-2&1\\0&-4\end{pmatrix}.
\]
Thus the required raising map is integral and the required
Weyl map is a lattice isomorphism. The cokernel is
$\mathbb Z/8\mathbb Z$, because the matrix has determinant
$8$ and its entries have greatest common divisor $1$.
With the split lattice $\mathbb Zu\oplus\mathbb Zv$ and its
Weyl image, the same rational map instead has cokernel
$\mathbb Z/2\oplus\mathbb Z/4$. Both indices are $8$,
but the least inverse denominator changes from $4$ to $8$.
This proves that the product annihilator in AG-LEFSCH-INDEX
cannot always be replaced by the least common multiple under
its stated hypotheses. This example asserts the local
raising/Weyl lattice conditions only; it supplies no claim
of stability under all divided powers or of a geometric
realization.

The theta-to-Lefschetz gate for a total index is now smaller:
identify the actual integral theta/period lattice with the
chosen cohomological lattice, and verify that the normalized
raising map is integral and the Weyl (or correctly normalized
Fourier) operator identifies the two weight lattices.
Integral isotypic splitting is unnecessary for the index.
It matters for the full Smith group only through the explicit
finite defect sequence; even there, the possible exceptional
primes are restricted as above. This does not assert the
missing theta-lattice comparison or identify it with H8(iv).
The rational Arthur comparison and AG-ARTHUR-MIN supply
different parts of the construction, and A and M do not by
themselves provide these integral maps.

\subsection{The geometric input to the finite-field factoring interface}

\paragraph{AG-CYCLIC-DIV (recent source theorem, with its full geometric proof checked).}
The additional manuscript PF26, OpenAI,
\emph{Deterministic Polynomial Factorization over Prime Fields},
4 October 2026, Section 4, pp.~13--19, gives the following
geometric mechanism. Let $q\ne\operatorname{char}K$ be an odd
prime, let the finite field $K$ contain $\mu_q$, and let
$F\in K[X]$ be monic, square-free and split, of degree
$N\ge q$ divisible by $q$. On the smooth projective curve
$C:Y^q=F(X)$, the deck transformation $\sigma(Y)=\zeta_qY$
acts on its Jacobian. Put $\lambda=1-\sigma$, let $\Lambda$
be the degree-zero divisor lattice on the $q$ points at infinity,
and write $e=(\sigma-1)\infty_0$.
For $r\ge0$ and $m=1+(q-1)r$, every solution
\[
 \lambda^t d=[W],\qquad d\in J(\bar K),\quad
 W\in\Lambda,\quad 1\le t\le m,
\]
is rational over the extension of $K$ of degree $q^r$, and
has a rational divisor representative there. This is exactly
Lemma 4.3 and Lemma 4.4; their definitions and complete proofs,
including Lemmas 4.1--4.2, were independently read and checked.
For $r=v_q(N)$ the extension degree is at most $N$.
Corollary 4.5 supplies the required ramification-class chains,
and Proposition 4.6 proves that the associated norm labels
must separate two ramification points by step $m$.

The decisive information is integral: $\Lambda$ is a free
rank-one $\mathbb Z[\zeta_q]$ module and
$Ne\in\lambda^{m-1}\Lambda\setminus\lambda^m\Lambda$.
The source retains a Jacobian class together with its actual
infinity divisor. Its filtration proof bounds the field of
definition, while its final lattice obstruction forces
separation for every sequence of permitted choices. This
connects a finite \'etale algebra's idempotent problem to
Jacobians and rational divisors, as used by CX-ALG-027/028.
Here $\lambda$ is the cyclotomic operator $1-\sigma$, not
a $\lambda$-ring operation. The result supplies a specific
integral geometric mechanism; it does not identify an
automorphic theta lattice or an arbitrary inertia extension.
The complete factoring transfers retain PF26's algebraic
algorithm as an additional recent-manuscript input. The
review recorded here covers its entire Section 4 only, not
the norm-solving implementation or a formal certificate.

\subsection{Central Artin orders and an abstract cohomological functor}

\paragraph{AG-DENINGER: a classical existence theorem with an extra analytic comparison.}
Deninger, \emph{Is there a Birch and Swinnerton-Dyer conjecture for
Dedekind zeta functions?}, arXiv:2504.15767v3 (1 April 2026),
Theorems 2.15--2.16, pp.~10--14, constructs an abstract functor
on number fields with symplectic forms, positive compatible star
operators and Galois descent.  On a Galois field its representation
is the sum, once each, of the irreducible symplectic Artin
representations with root number $-1$.  All such functors are
isomorphic, with automorphism group
$\operatorname{Map}(\Pi_-,\{\pm1\})$, where $\Pi_-$ denotes
these finite-image representations of $G_{\mathbb Q}$.
This abstract theorem is unconditional.  The additional
identification
$\dim V_K=\operatorname{ord}_{s=1/2}\zeta_K(s)$ uses the
Serre hypothesis stated on pp.~1--2: the central Artin order
is one for precisely those symplectic representations with
root number $-1$, and zero for every other irreducible.
Neither A nor M supplies that multiplicity assertion.

The following extension retains arbitrary central multiplicities.
Its proof also supplies the invariant-space argument for fields
which are not Galois over $\mathbb Q$.  In all statements a
positive star structure means a nondegenerate complex alternating
form $B$ and an antilinear, equivariant operator $J$ such that
$J^2=-1$ and $h(v,w)=B(v,Jw)$ is positive definite hermitian,
linear in $v$.  Equivariance of $J$ is essential.

\begin{proposition}[AG-CENTRAL-STRUCT: the exact finite positive-star criterion]
Let $G$ be finite and let
$\Theta=\sum_{\pi\in\operatorname{Irr}(G)}m_\pi\chi_\pi$,
with $m_\pi\in\mathbb Z$.  Suppose
$m_{\bar\pi}=m_\pi$, and that $m_\pi$ is even whenever
$\pi$ is orthogonal.  The following are equivalent:
\begin{enumerate}
\item Every $m_\pi$ is nonnegative.
\item The matrix $(\Theta(gh^{-1}))_{g,h\in G}$ is positive
semidefinite.
\item There is a complex $G$-representation of character $\Theta$
with a positive star structure.
\end{enumerate}
When these conditions hold, the structured representation is
unique up to equivariant isomorphism.  Its group of equivariant
automorphisms preserving $B$ and $J$ is
\[
 \prod_{\pi\ {
m symplectic}}\mathrm O(m_\pi)
 \ \times\!
 \prod_{\pi\ {
m orthogonal}}\mathrm{Sp}(m_\pi/2)
 \ \times\!
 \prod_{\{\pi,\bar\pi\},\,\pi\not\simeq\bar\pi}
                              \mathrm U(m_\pi).          \tag{AG-central-groups}
\]
Here $\mathrm O(m)$ is the real orthogonal group and
$\mathrm{Sp}(r)$ is the compact group on a quaternionic space
of dimension $r$; factors of size zero are omitted.
\end{proposition}

\begin{proof}
The equivalence of the first two conditions is the finite
convolution calculation of MT-ART02: on the matrix coefficients
of $\pi$, the normalized convolution matrix has eigenvalue
$m_\pi/\dim\pi$.  The third condition implies the first because
multiplicities of an actual representation are nonnegative.
It remains to construct and classify the positive star structures.

Fix invariant positive hermitian forms on irreducible
representations, obtained by averaging.  An orthogonal irreducible
has an antiunitary equivariant real structure $C_\pi$ with
$C_\pi^2=1$; a symplectic irreducible has an antiunitary
equivariant structure $J_\pi$ with $J_\pi^2=-1$.
For example, in the symplectic case define $A$ by
$h(v,w)=B_\pi(v,Aw)$ for a chosen invariant alternating form.
It is antilinear and equivariant, and $A^2=\lambda I$ by
Schur's lemma.  For $v\ne0$,
\[
 0<h(Av,Av)=B_\pi(Av,A^2v)
             =-\lambda B_\pi(v,Av)=-\lambda h(v,v).
\]
Thus $\lambda$ is negative real.  Moreover
$h(Av,Aw)=-\lambda h(w,v)$, so
$J_\pi=(-\lambda)^{-1/2}A$ is antiunitary and has square
$-1$.  For a symmetric invariant form the same calculation
has a plus sign: $\lambda>0$, and
$C_\pi=\lambda^{-1/2}A$ is antiunitary with square $1$.

For orthogonal $\pi$, equip its multiplicity space
$M_\pi=\mathbb C^{m_\pi}$ with its standard antiunitary
quaternionic structure, which exists because $m_\pi$ is even,
and use $C_\pi\otimes J_{M_\pi}$.  For symplectic $\pi$,
equip $M_\pi$ with ordinary conjugation and use
$J_\pi\otimes C_{M_\pi}$.  Tensor products here are antilinear
on both factors and therefore define a well-defined antilinear
operator on the complex tensor product.  Both constructions
square to $-1$ and are antiunitary for the product hermitian form.

For a non-selfdual pair set $W=\pi\otimes\mathbb C^{m_\pi}$
and use $W\oplus\bar W$ with its orthogonal sum hermitian form.
The map
\[
             J(w,\bar v)=(-v,\bar w)
\]
is equivariant, antiunitary and has square $-1$.
Take the orthogonal direct sum of these constructions.
If $h$ is its hermitian form, put
$B(v,w)=-h(v,Jw)$.  This is complex bilinear and nondegenerate.
Antiunitarity and $J^2=-1$ give
$h(w,Jv)=-h(v,Jw)$, so $B$ is alternating, and
$B(v,Jw)=h(v,w)$.  This proves existence.

For classification, Schur's lemma makes the hermitian form on
each isotypic component a tensor product with a positive form
on the multiplicity space.  Equivariance makes $J$ exchange
dual isotypic components.  On a selfdual component it therefore
has the form just used, with an antiunitary real structure on
the multiplicity space of a symplectic irreducible and an
antiunitary quaternionic structure on that of an orthogonal
irreducible.  These structures have orthonormal real or
quaternionic bases.  In the latter case one chooses a unit
vector $e$, then $Je$; they are orthogonal, and their
orthogonal complement is $J$-stable, giving induction.
For a non-selfdual pair choose an orthonormal basis on one
multiplicity space and let $J$ determine the other.
This proves uniqueness up to isomorphism.
Preserving $B$ and $J$ is equivalent to being unitary and
commuting with $J$.  On the respective multiplicity spaces
these are precisely $\mathrm O(m)$, $\mathrm{Sp}(m/2)$ and
$\mathrm U(m)$, proving (AG-central-groups).
\end{proof}

\begin{proposition}[AG-CENTRAL-FUNCTOR: the remaining central gate is exact]
For every irreducible finite-image complex representation $\pi$
of $G_{\mathbb Q}$ put
\[
             m_\pi=\operatorname{ord}_{s=1/2}L(s,\pi).
\]
The following central, rather than global, holomorphy condition
is equivalent to existence of an abstract positive-star functor
with the specified central Artin characters:
\[
 \text{every finite-image Artin function is holomorphic at }1/2.
                                                        \tag{AG-Central-H}
\]
More precisely, under (AG-Central-H) there is a functor $K\mapsto V_K$
on number fields and embeddings, with isometric pullbacks $i^*$,
adjoints $(i^*)^\dagger$, and
\[
 (i^*)^\dagger i^*=1,\qquad
 i^*(i^*)^\dagger=\frac1{[L:K]}
                 \sum_{\sigma\in\operatorname{Gal}(L/K)}\sigma
 \quad(L/K\text{ Galois}).                              \tag{AG-central-descent}
\]
For $L/\mathbb Q$ Galois its character is
$\sum_\pi m_\pi\chi_\pi$, with $\pi$ ranging over
$\operatorname{Irr}\operatorname{Gal}(L/\mathbb Q)$; for every
number field $K$ one has
\[
             \dim V_K=\operatorname{ord}_{s=1/2}\zeta_K(s).
                                                        \tag{AG-central-dimension}
\]
Such functors with those characters and descent are unique up
to isomorphism.  Their common automorphism group is the product
in (AG-central-groups) over all finite-image irreducibles of
$G_{\mathbb Q}$.  Conversely the prescribed character at every
Galois field forces (AG-Central-H).  An equality of dimensions
alone is not claimed to force this condition.
\end{proposition}

\begin{proof}
Brauer meromorphy makes the orders $m_\pi$ integers.
Complex conjugation of the Artin Dirichlet coefficients gives
$m_{\bar\pi}=m_\pi$.  At $s=1/2$ all archimedean factors are
finite and nonzero.  For a selfdual $\pi$ its functional
equation therefore gives
\[
                        (-1)^{m_\pi}=W(\pi),             \tag{AG-central-parity}
\]
even if the order is negative.  The Fröhlich--Queyrut theorem
gives $W(\pi)=1$ for orthogonal $\pi$.  An exact primary
statement is Queyrut, \emph{Sur l'équation fonctionnelle des
séries $L$ d'Artin pour des caractères réels}, Mémoires SMF 37
(1974), theorem on p.~133: it applies to representations
realizable over $\mathbb R$, not to all real-valued characters.
Thus the parity hypotheses of AG-CENTRAL-STRUCT hold.
Condition (AG-Central-H) gives the remaining nonnegativity.

Choose one representative of each irreducible and form the
algebraic direct sum
\[
 \mathcal V=\bigoplus_\pi \pi^{\oplus m_\pi},\qquad
                V_K=\mathcal V^{G_K},
 \quad G_K=\operatorname{Gal}(\overline{\mathbb Q}/K),  \tag{AG-central-construction}
\]
equipped with the blockwise forms and star operators of
AG-CENTRAL-STRUCT.  Every vector has finite orbit.  For a
finite Galois extension $L/\mathbb Q$, its invariant subspace
contains only the finitely many irreducibles factoring through
$\operatorname{Gal}(L/\mathbb Q)$, hence is finite dimensional.
For arbitrary $K$ embed it in such an $L$; then $V_K\subset V_L$
is also finite dimensional.  The restriction of $h$ to $V_K$
is positive, $J$ preserves $V_K$, and the restriction of $B$
is nondegenerate because $B(v,Jv)=h(v,v)>0$ for $v\ne0$.

An embedding $\alpha:K\hookrightarrow L$ extends to
$\sigma\in G_{\mathbb Q}$; use the restriction of $\sigma$
to $V_K$.  Two extensions differ on the right by an element
of $G_K$, so this map is independent of the choice.  It is
isometric and respects $B,J$, and composition is functorial.
For an inclusion $K\subset L$ with $L/K$ Galois, the average
of $\operatorname{Gal}(L/K)$ on $V_L$ is the orthogonal
projector to $V_L^{\operatorname{Gal}(L/K)}=V_K$.
This proves (AG-central-descent), including the case that
$K/\mathbb Q$ is not Galois.  Unnormalized transfer is
$[L:K](i^*)^\dagger$, so its composite with pullback is
multiplication by the extension degree.

Choose a finite Galois $L/\mathbb Q$ containing $K$ and put
$H=\operatorname{Gal}(L/K)$.  Artin formalism for the
permutation representation $\operatorname{Ind}_H^G1$ gives
\[
 \operatorname{ord}_{1/2}\zeta_K
       =\sum_{\pi\in\operatorname{Irr}(G)}
                    m_\pi\dim\pi^H
       =\dim V_K.
\]
This proves the dimension formula.  Conversely an actual
representation with the prescribed character has nonnegative
irreducible multiplicities, which is precisely central
holomorphy.  Representations over other base number fields
are covered as well by induction to $G_{\mathbb Q}$.

Finally consider any functor with the stated character and
descent properties.  Take its directed union over finite
Galois extensions of $\mathbb Q$; isometric pullbacks are
injective.  This union carries a locally finite $G_{\mathbb Q}$
action, and the averaging identity recovers its $G_K$-invariants
as $V_K$.  On each finite quotient the prescribed character
gives exactly the multiplicities $m_\pi$.  The union therefore
decomposes into the blocks of (AG-central-construction).
The classification in AG-CENTRAL-STRUCT gives a blockwise
equivariant isometry to that construction and hence a natural
isomorphism of functors.  Natural automorphisms are precisely
the compatible equivariant block automorphisms, which gives
the asserted product of compact classical groups.
\end{proof}

Under Serre's extra multiplicity conjecture all the surviving
factors in (AG-central-groups) are $\mathrm O(1)=\{\pm1\}$,
recovering Deninger's result.  With higher multiplicities the
same central positivity condition allows larger automorphism
groups.  MT-ART02 identifies the exact finite positive matrix
which must hold at each Galois level.  M confines the possible
Artin divisor to the critical line, but it does not supply
(AG-Central-H) at its centre.  These are abstract representation
spaces attached to arithmetic curves $\operatorname{Spec}\mathcal O_K$;
no geometric cohomology, integral lattice, cycle map, regulator
or formula for the leading coefficient has been constructed.

\begin{proposition}[AG-CENTRAL-REAL: the real-coefficient obstruction survives higher order]
For a Galois field $L/\mathbb Q$ satisfying central Artin
holomorphy, its representation $V_L$ in AG-CENTRAL-FUNCTOR
has an equivariant real form if and only if every symplectic
irreducible has even central multiplicity.  In particular a
symplectic irreducible with root number $-1$ forbids such a
real form, regardless of whether its central order is one,
three or larger.  The same obstruction applies unconditionally
to Deninger's sign-selected functor whenever that negative-root
representation occurs.
\end{proposition}
\begin{proof}
A real form is equivalent to an equivariant antilinear operator
$C$ with $C^2=1$.  On a symplectic isotypic component write
$C=J_\pi\otimes A$ with $J_\pi^2=-1$ as above.  Then $A$
is antilinear on its multiplicity space and must satisfy
$A^2=-1$.  If $A=M\circ\mathrm{conj}$ on $\mathbb C^m$,
then $M\bar M=-I$, so
$|\det M|^2=(-1)^m$; this requires $m$ even.
Conversely the standard quaternionic $A$ exists in every even
dimension.  Orthogonal blocks always have a real form, and
nonselfdual pairs have one by exchanging the two conjugate
summands.  This proves the equivalence.  Formula
(AG-central-parity) makes $m_\pi$ odd when $W(\pi)=-1$.
For the sign-selected functor its multiplicity is already one.
\end{proof}

Thus the root-number class studied in MT-ART10--12 can obstruct
a functorial real or rational coefficient realization, as well
as force central vanishing.  This statement concerns the
equivariant coefficient structure; it does not rule out the
positive complex quaternionic realization just constructed.

\begin{proposition}[AG-CENTRAL-Q8: an unconditional arithmetic real-form obstruction]
For the explicit tame quaternionic field
\[
 K=\mathbb Q(\sqrt5,\sqrt{41}),\qquad
 N=K\left(\sqrt{\frac{5+\sqrt5}{2}
                         \frac{41+\sqrt{205}}{2}}\right),
\]
the actual central-order character
$\sum_\pi\operatorname{ord}_{s=1/2}L(s,\pi)\,\pi$ is the
character of a complex representation of $\operatorname{Gal}(N/\mathbb Q)$,
but that representation has no equivariant real form.
This assertion needs neither M, general central Artin holomorphy,
nor Serre's conjecture on the exact central multiplicities.
\end{proposition}
\begin{proof}
MT-ART10 records the classical Martinet example (1971,
Theorem IV.1, p.~406, Example 3, p.~407), and Fr\"ohlich's
root-number criterion: its group is $Q_8$ and the unique
degree-two symplectic character $\chi$ has $W(\chi)=-1$.
Every irreducible is either linear or this $\chi$.
The linear factors are Dirichlet functions (including $\zeta$),
all holomorphic at $1/2$.  If $C_4\subset Q_8$ and $\lambda$
is a faithful linear character, then
$\chi=\operatorname{Ind}_{C_4}^{Q_8}\lambda$; hence its
Artin function is an entire nontrivial Hecke function.
Thus all central orders are nonnegative integers, so their
sum is an actual character.  The functional equation with
$W(\chi)=-1$ makes its multiplicity positive and odd.
AG-CENTRAL-REAL now forbids a real form.  The field and its
root sign are the located classical input from MT-ART10;
this deduction does not claim a new computation of their values.
\end{proof}

This is a concrete obstruction for a proposed real-coefficient
central cohomology of the arithmetic curve
$\operatorname{Spec}\mathcal O_N$, despite positivity of the
complex central character.  The complex symplectic/star
structure of AG-CENTRAL-STRUCT remains available.

\subsection{A precise obstruction to the exponential half-Tate candidate}

\paragraph{AG-EXP-VANISH: the located classical obstruction.}
Fresán--Jossen's theorem, reproduced with its complete proof as
Deninger's Theorem 4.1, pp.~16--17, states
\[
 \operatorname{Ext}^1_{\mathcal M^{\exp}(K)}
             (\mathbb Q(0),\mathbb Q(1/2))=0
                       \qquad(K\subset\mathbb C).
\]
Here $\mathbb Q(1/2)$ is the dual of the exponential motive
of $(\mathbb A^1,x^2)$, in the specified exponential-motive
category.  This is a proved obstruction to that particular
candidate for the central cohomology, not a vanishing theorem
for all proposed arithmetic cohomology theories.
The following formulation isolates the mechanism and extends it
to any monodromy with no eigenvalue one.

\begin{proposition}[AG-EXP-SEPARATE: monodromy separation kills the extension]
Let $N\in\mathcal M^{\exp}(K)$ have perverse realization
$j_!\mathcal L[1]$ on $\mathbb A^1_{\mathbb C}$, where
$j:\mathbb G_m\hookrightarrow\mathbb A^1$, and let $T$
be the rational monodromy of the finite-rank local system
$\mathcal L$ about zero.  If $T-I$ is invertible, then
\[
 \operatorname{Ext}^1_{\mathcal M^{\exp}(K)}
                         (\mathbb Q(0),N)=0.
\]
In particular $T=-1$ recovers AG-EXP-VANISH.
\end{proposition}
\begin{proof}
Take an extension $0\to N\to M\to\mathbb Q(0)\to0$.
The perverse realization of $\mathbb Q(0)$ is
$j_!\mathbb Q[1]$.  Exactness of realization gives an
extension of these two perverse sheaves.  Its stalk complex
at zero vanishes, since those of both end terms vanish.
On $\mathbb G_m$ it is an extension of local systems, so in
a basis adapted to the extension its monodromy is
\[
                   \begin{pmatrix}T&b\\0&1\end{pmatrix}.
\]
Changing the lift of the last basis vector by
$u=-(T-I)^{-1}b$ makes this matrix block diagonal.
The vanishing stalk at zero identifies the perverse sheaf
with the extension by zero of its restriction.  Hence the
perverse extension splits.  Its largest trivial subobject
is exactly $j_!\mathbb Q[1]$, because $T-I$ is invertible
on the other summand.

Apply the theorem of the fixed part for exponential motives,
Fresán--Jossen, \emph{Exponential motives}, Theorem 6.5.1 in
the version cited as [8] in Deninger.  The exact statement
quoted and used in his proof of Theorem 4.1 says that the
largest trivial perverse subobject is realized
by a classical submotive $M_0\subset M$.  The composite
$M_0\to M\to\mathbb Q(0)$ realizes an isomorphism.
The realization is faithful and exact, so its kernel and
cokernel are zero; the composite is an isomorphism of motives.
Its inverse supplies a section of $M\to\mathbb Q(0)$,
proving that the extension splits.
\end{proof}

The fixed-part theorem is an explicit classical input here;
its complete source proof is not re-established by the matrix
calculation.  The exact monodromy hypothesis also matters:
when $1$ is an eigenvalue, $(T-I)^{-1}$ is unavailable and
this argument says nothing about the extension group.
The obstruction narrows the central cohomology construction
suggested by the trace/Fredholm programme of cote 1.  It does
not evaluate the distinct integral Kummer extensions of cote
149 or the general special-value $K$-groups in AG-EXT-K.

\subsection{The exact local exponential and determinant hypotheses}

The local regulator step in a theta or special-value construction
has a condition distinct from the global zero-location inputs.
We fix $G_p=\operatorname{Gal}(\overline{\mathbb Q}_p/\mathbb Q_p)$,
a finite extension $L/\mathbb Q_p$, and a de Rham
$L$-representation $V$ of $G_p$.  Put
\[
 D=D_{\mathrm{cris}}(V),\qquad
 t_V=D_{\mathrm{dR}}(V)/\operatorname{Fil}^0D_{\mathrm{dR}}(V).
\]
Frobenius on $D$ is $L$-linear here.  We use the convention
that $\mathbb Q_p(1)$ has Hodge--Tate weight $+1$,
crystalline Frobenius $p^{-1}$ and filtration jump $-1$.
The classical inputs are Bloch--Kato, \emph{L-functions and
Tamagawa numbers of motives}, Proposition 3.8 and its full
proof, pp.~354--359, Corollary 3.8.4, pp.~355--356,
and Definition 3.10, p.~359.  In particular
\[
\begin{split}
0\longrightarrow H^0(G_p,V)&\longrightarrow D^{\varphi=1}
 \longrightarrow t_V\xrightarrow{\exp_V}H_e^1(G_p,V)
 \longrightarrow0,\\
0\longrightarrow H^0(G_p,V)&\longrightarrow D
 \xrightarrow{\ d\mapsto((1-\varphi)d,\bar d)\ }
 D\oplus t_V\longrightarrow H_f^1(G_p,V)\longrightarrow0.
\end{split}
\tag{AG-local-exact}
\]
These are statements about rational local cohomology.  An
integral lattice or its Tamagawa index is additional data.

\begin{proposition}[AG-LOCAL-EXP: the exact exponential gate]
The kernel of $\exp_V:t_V\to H^1(G_p,V)$ is
\[
 \operatorname{im}\bigl(D^{\varphi=1}\longrightarrow t_V\bigr),
 \qquad
 H_f^1(G_p,V)/H_e^1(G_p,V)
       \simeq D/(1-\varphi)D.
\tag{AG-local-gate}
\]
Consequently the exponential is injective precisely when that
image is zero, and it is an isomorphism
$t_V\xrightarrow{\sim}H_f^1(G_p,V)$ precisely when
$D^{\varphi=1}=0$.
If $V\simeq V^*(1)$ and $H^0(G_p,V)=0$, then the following
are equivalent:
\[
 D^{\varphi=1}=0,\quad \exp_V\text{ is injective},\quad
 H_e^1=H_f^1,\quad H_g^1=H_f^1.
\]
Under these equivalent conditions the dual exponential induces
an isomorphism
\[
 H^1(G_p,V)/H_f^1(G_p,V)\xrightarrow{\sim}
            \operatorname{Fil}^0D_{\mathrm{dR}}(V).
\]
\end{proposition}
\begin{proof}
The first equality in (AG-local-gate) is exactness of the
first row of (AG-local-exact).  In the second row, quotient
the target $D\oplus t_V$ by the image of its second summand.
The resulting cokernel is $D/(1-\varphi)D$, proving the
second equality.  Since $1-\varphi$ is an endomorphism of
one finite-dimensional $L$-space, its cokernel vanishes if
and only if its kernel does.  If that kernel vanishes, both
injectivity and surjectivity of $\exp_V$ follow; conversely
surjectivity onto $H_f^1$ already forces it to vanish.
When $H^0=0$, the map $D^{\varphi=1}\to t_V$ is injective,
giving the additional injectivity equivalence.  Bloch--Kato
duality states that $H_e^1(V)$ and $H_g^1(V^*(1))$ are exact
annihilators, and $H_f^1(V)$ is the exact annihilator of
$H_f^1(V^*(1))$.  Selfduality and nondegeneracy therefore
give the last equivalence.  Finally the dual exponential is
adjoint to $\exp_{V^*(1)}$ under the perfect de Rham and
Tate pairings.  Dualizing the isomorphism onto $H_f^1$ gives
the asserted quotient isomorphism, with its stated target.
\end{proof}

\begin{proposition}[AG-LOCAL-EXCEPTION: generic selfduality does not suffice]
For every odd prime $p$, there is an irreducible crystalline
two-dimensional $\mathbb Q_p$-representation $V$ such that
\[
 \begin{gathered}
 \det V=\mathbb Q_p(1),\qquad
 H^0(G_p,V)=H^2(G_p,V)=0,\\
 \dim H_f^1(G_p,V)=1,\qquad
 \exp_V=0,\qquad\exp_V^*=0.
 \end{gathered}
\]
Its Hodge--Tate weights are $(2,-1)$.
\end{proposition}
\begin{proof}
Take $D=\mathbb Q_pe_1\oplus\mathbb Q_pe_2$, monodromy zero,
and
\[
 \varphi e_1=e_1,\qquad \varphi e_2=p^{-1}e_2,
 \qquad
 \operatorname{Fil}^iD=
 \begin{cases}
 D&i\le-2,\\
 \mathbb Q_p(e_1+e_2)&-1\le i\le1,\\
 0&i\ge2.
 \end{cases}
\tag{AG-local-example}
\]
The total Hodge and Newton numbers are both $-1$.
The only proper nonzero Frobenius-stable subspaces are
$\mathbb Q_pe_1$ and $\mathbb Q_pe_2$.  Each has induced
Hodge number $-2$, while its Newton number is respectively
$0$ or $-1$.  Thus $t_H\le t_N$ for every such subspace,
and $D$ is weakly admissible.  Colmez--Fontaine,
\emph{Construction des représentations $p$-adiques
semi-stables}, Theorem A (author version, p.~2), supplies
the crystalline representation $V$.  No proper nonzero
Frobenius-stable subspace has equal Hodge and Newton numbers,
so no such subspace is admissible.  The equivalence of
categories therefore also proves that $V$ is irreducible.

Its determinant has Frobenius $p^{-1}$ and its unique
filtration jump at $-1$, so it is $\mathbb Q_p(1)$.
The wedge pairing gives the required symplectic selfduality.
Moreover
\[
 H^0(G_p,V)=D^{\varphi=1}\cap\operatorname{Fil}^0D
   =\mathbb Q_pe_1\cap\mathbb Q_p(e_1+e_2)=0.
\]
Local duality gives $H^2=0$.  The map
$D^{\varphi=1}=\mathbb Q_pe_1\to
t_V=D/\mathbb Q_p(e_1+e_2)$ is an isomorphism.
The first row of (AG-local-exact) gives $\exp_V=0$ and
$H_e^1=0$; the second row gives $\dim H_f^1=\dim t_V=1$.
The local Euler characteristic gives $\dim H^1=2$.
Finally, by symplectic selfduality,
$\exp_{V^*(1)}=0$ as well.  Adjointness of exponential and
dual exponential under perfect pairings forces
$\exp_V^*=0$ on all of $H^1$.  Equivalently,
Bloch--Kato duality gives $H_g^1=(H_e^1)^\perp=H^1$.
\end{proof}

\paragraph{Exact scope of the recent-manuscript correction.}
Burungale--Kobayashi--Nakamura--Ota,
\emph{A local sign decomposition for symplectic self-dual
Galois representations of rank two}, arXiv:2508.17776v1,
allows $R=\mathbb Q_p$ in \S2.2(ii), p.~12.  Its
Definition 2.11, p.~16, defines genericity in this case
by $H^0=H^2=0$.  Thus (AG-local-example) satisfies exactly
those hypotheses, and the identity is a de Rham specialization
in the sense of \S5.1.1, p.~35.  Corollary 3.2, p.~19,
asserts nonvanishing of the ordinary dual exponential on a
basis of one of the signed lines.  Here it vanishes on the
whole cohomology, so that corollary is false in its stated
generality.  There is no purity or exclusion of Frobenius
fixed vectors in those definitions.

Separately, the displayed proof of Proposition 4.17,
pp.~28--29, uses $\exp^{-1}:H_f^1\to t_V$ as an
isomorphism.  AG-LOCAL-EXP gives its exact missing hypothesis.
This identifies a gap in that comparison as written; it does
not refute the main sign-decomposition theorem.
AG-LOCAL-SPLICE below repairs this precise determinant
comparison by a polarization-preserving specialization of
the full exact complex. It does not repair the separate
ordinary-exponential nonvanishing assertion. Applications for which the
crystalline period space vanishes, or which satisfy the
purity condition below, are not counterexamples.

\begin{proposition}[AG-LOCAL-DET: the finite determinant survives the exception]
Suppose $H^0(G_p,V)=0$.  With the ordering $D\oplus t_V$
as written in (AG-local-exact), that exact sequence gives
a canonical isomorphism of ordinary determinant lines
\[
 \delta_f:\det_L H_f^1(G_p,V)\xrightarrow{\sim}\det_L t_V.
\]
If $1-\varphi$ is invertible and $u$ is a top exterior
vector of $t_V$, its relation with the ordinary exponential is
\[
 \delta_f\bigl((\det\exp_V)(u)\bigr)
                      =\det(1-\varphi\mid D)\,u.
\tag{AG-local-det}
\]
The determinant isomorphism remains defined when the ordinary
exponential is zero.
\end{proposition}
\begin{proof}
The second row of (AG-local-exact) is now a short exact
sequence $0\to D\to D\oplus t_V\to H_f^1\to0$.
Use its exterior-power isomorphism, with the subspace basis
preceding lifted quotient vectors, and cancel the identical
source and target $\det D$ factors.  This defines $\delta_f$
without dividing by a Frobenius eigenvalue.  When $F=1-\varphi$
is invertible, choose a basis of $D$ and lift the exponential
of a basis of $t_V$ by vectors $(0,u_i)$.  The resulting
block matrix has diagonal blocks $F$ and $I$, proving
(AG-local-det).

In (AG-local-example), put $u=\bar e_1=-\bar e_2$ and
$a=1-p^{-1}$.  The image of $e_1$ in $D\oplus t_V$ is
$(0,u)$ and that of $e_2$ is $(ae_2,-u)$.
The quotient is generated by $h=[(e_1,0)]$.
In the ordered target basis $(e_1,e_2,u)$, the columns
consisting of those two image vectors and the lift of $h$
have determinant $-a$.  Hence
$\delta_f(h)=-a u\ne0$, although $\exp_V(u)=0$.
This also verifies directly that the two maps cannot be
interchanged at an exceptional period.
\end{proof}

\begin{proposition}[AG-LOCAL-PURITY: a geometric condition that removes the exception]
Assume all eigenvalues of crystalline Frobenius on $D$ are
algebraic numbers of a fixed Weil weight $w\ne0$, meaning
that every complex embedding has absolute value $p^{w/2}$.
Then $D^{\varphi=1}=0$, so $\exp_V:t_V\to H_f^1$ is an
isomorphism.  In particular this holds for
\[
 V=H^{2n+1}_{\text{ét}}
       (X_{\overline{\mathbb Q}_p},\mathbb Q_p)(n+1),
\]
when $X/\mathbb Q_p$ admits a smooth projective model over
$\mathbb Z_p$.  Crystalline comparison and the already-known
finite-field purity theorem give
$w=(2n+1)-2(n+1)=-1$.
\end{proposition}
\begin{proof}
The algebraic number $1$ has absolute value $1$ in every
complex embedding.  Since $p^{w/2}\ne1$, it cannot be an
eigenvalue.  Apply AG-LOCAL-EXP.  In the geometric case,
comparison identifies Frobenius eigenvalues with those of
the smooth projective reduction in degree $2n+1$, and the
Tate twist multiplies them by $p^{-(n+1)}$, giving the
displayed weight.  Exact classical locators are
Bhatt--Morrow--Scholze, \emph{Integral $p$-adic Hodge theory},
Theorem 1.1(i), author version p.~2, repeated as
Theorem 14.6(i), pp.~121--122; and Katz--Messing,
\emph{Some consequences of the Riemann hypothesis for varieties
over finite fields}, Theorem 1, pp.~74--75, and
Corollary 1, pp.~75--76.  If the reduction has several
geometric components, apply the latter theorem after a
finite residue extension to each component; eigenvalues
of a Frobenius power then have the claimed absolute
values, hence so do the original eigenvalues.
These are classical inputs, as in AG-PURITY.  No new purity
theorem or consequence of A or M is asserted.
\end{proof}

For motivic realizations this identifies the exact local
condition needed before turning a de Rham tangent vector
or a normalized theta derivative into a finite local
cohomology class by an invertible exponential.  At bad
reduction, or for an abstract local representation such
as (AG-local-example), purity of that precise crystalline
space must be checked separately.  Neither global GRH
input supplies it.  Even where the rational exponential
is invertible, the integral image, period normalization
and global cycle map remain the separate gates already
recorded in AG-THETA-FOURIER and AG-SHIMURA-PERIOD.

\subsection{Repairing the full local determinant comparison}

We retain the normalization $\varphi=p^{-1}$ on $L(1)$.
The classical potentially semistable theorem supplies the
Weil--Deligne representation of a de Rham representation.
In the unnormalized convention used here,
$N\varphi=p\varphi N$ and
$D_{\mathrm{cris}}$ is the inertia-invariant kernel of $N$.
Exact primary locators are Berger, \emph{Équations
différentielles $p$-adiques et $(\varphi,N)$-modules filtrés},
\S I.1, p.~16, for the conventions, and \S V.1, p.~31,
for de Rham $\Longleftrightarrow$ potentially semistable;
Nakamura, \emph{A generalization of Kato's local
$\varepsilon$-conjecture}, author version arXiv:1305.0880v3,
\S3.3, p.~35, constructs the associated Weil--Deligne object.
These substantial classical existence inputs are not new
consequences of A or M.

\begin{proposition}[AG-LOCAL-WD: the exceptional rank-two case is crystalline]
Let $V$ be a two-dimensional de Rham $L$-representation
of $G_p$ with $\det V=L(1)$.  If
$D_{\mathrm{cris}}(V)^{\varphi=1}\ne0$, then $V$ is
crystalline.  More precisely, if $V$ is not crystalline,
its crystalline period space is either zero or a line
on which $\varphi$ acts by $p^{-1}$ or $-p^{-1}$.
Consequently the ordinary exponential of every such
noncrystalline $V$ is an isomorphism $t_V\to H_f^1$.
\end{proposition}
\begin{proof}
Extend the coefficient field to split the finite inertia
representation on the two-dimensional Weil--Deligne
space $W$.  This does not change the relevant dimensions
or the presence of the eigenvalue one.  Inertia is
semisimple in characteristic zero, and its determinant
is trivial because the Weil--Deligne determinant of
$L(1)$ is unramified.  Thus $\dim W^I$ cannot be one:
a complementary inertia character would also have to be
trivial.  The only possibilities are zero and two.

If $W^I=0$, its subspace $(\ker N)^I$ is zero, hence
$D_{\mathrm{cris}}=0$.  If $W^I=W$, the representation
is semistable.  When $N=0$ it is crystalline.  Otherwise
$\ker N=\operatorname{im}N$ is a line.  If Frobenius
acts on that line by $\lambda$, the relation
$N\varphi=p\varphi N$ makes its action on the quotient
$p\lambda$.  Its determinant is $p^{-1}$, so
\[
                  p\lambda^2=p^{-1},\qquad
                  \lambda=\pm p^{-1}.
\]
Neither value is one.  This proves both assertions
about the crystalline space, and AG-LOCAL-EXP proves
the exponential assertion.  The line and its scalar
descend to $L$, since these two scalars already lie in
$\mathbb Q_p$.
\end{proof}

We now use the actual de Rham epsilon determinant,
including its graded ordering.  Let $p$ be odd and let
$V$ be symplectic selfdual of rank two, de Rham and
$H^0(G_p,V)=0$.  Write
\[
 \begin{gathered}
 H=H^1(G_p,V),\quad U=H_f^1(G_p,V),\quad E=D_{\mathrm{dR}}(V),\\
 F^0=\operatorname{Fil}^0E,\qquad t=E/F^0.
 \end{gathered}
\]
The symmetric Tate pairing $q$ on $H$ is perfect, and
$U$ is a one-dimensional Lagrangian.  The alternating
polarization $\beta$ on $E$ takes values in
$D_{\mathrm{dR}}(L(1))$.  Fix a basis $\zeta$ of
$L(1)$ and its Fontaine period $\tau$; the corresponding
de Rham basis is $\zeta/\tau$.  We identify the second
fundamental line with $L\zeta$ by the polarization.
Let $B_\epsilon:H\times H\to L\zeta$ be the alternating
pairing induced by the de Rham epsilon isomorphism.
Equivalently, define $w$ by
\[
                  B_\epsilon(x,y)=q(wx,y)\zeta.
\]
We use Nakamura's Weil--Deligne epsilon convention,
not the epsilon constant of the underlying Weil
representation alone.  His Remark 3.6, p.~38, records
their possible difference by
$\det(-\varphi\mid D_{\mathrm{st}}/D_{\mathrm{cris}})$.

\begin{proposition}[AG-LOCAL-SPLICE: the exceptional epsilon comparison survives]
With these conventions, even when
$D_{\mathrm{cris}}(V)^{\varphi=1}\ne0$, one has
\[
 B_\epsilon(x,y)
   =-\Gamma(V)\epsilon(W(V))q(x,y)\zeta
       \qquad(x\in U,\ y\in H).
\tag{AG-local-splice}
\]
In this exceptional case $V$ is crystalline and
$\epsilon(W(V))=1$.  Thus $U$ is exactly the
$-\Gamma(V)\epsilon(W(V))$ eigenspace of $w$.
Together with the inverse-exponential calculation on
the locus $D_{\mathrm{cris}}^{\varphi=1}=0$, this repairs
the stated conclusion of BKNO Proposition 4.17 without
restoring the false nonvanishing assertion in its
Corollary 3.2.
\end{proposition}
\begin{proof}
We first specify the determinant map that must be
specialized.  Put $V^*=V^\vee(1)$ and
$D^*=D_{\mathrm{cris}}(V^*)$.  The finite exact sequence
for $V$ and the Tate dual of the one for $V^*$ splice
to give
\[
\begin{aligned}
0\longrightarrow D&\longrightarrow D\oplus t
 \longrightarrow H\\
 &\longrightarrow (D^*)^\vee\oplus F^0
 \longrightarrow (D^*)^\vee\longrightarrow0.
\end{aligned}
\tag{AG-local-spliced-complex}
\]
The first map is $d\mapsto((1-\varphi)d,\bar d)$.
The last map uses the transpose of $1-\varphi^*$,
where $\varphi^*$ is the Frobenius on the actual Tate
dual.  The other maps use the finite local subgroup,
its exact annihilator, and the perfect de Rham pairing.
This is exactly Nakamura (22)--(24), pp.~36--37.
Apply his acyclic determinant trivialization (19),
cancel the two copies of each crystalline determinant,
and use
\[
        \det F^0\otimes\det t\longrightarrow\det E,
                  \qquad v\otimes u\longmapsto v\wedge u.
\]
The result is the canonical map
$\Theta:\det H\to\det E$ occurring in that definition.
Multiplying by $\Gamma(V)$ and then applying
$\theta_{\mathrm{dR}}$ gives $B_\epsilon$;
the latter map sends $\zeta/\tau$ to
$\epsilon(W(V))\zeta$ by Lemma 3.4, pp.~37--38.
We will specialize precisely this map, rather than
an unrelated determinant of $1-\varphi$.

By AG-LOCAL-WD, the exceptional representation is
crystalline.  Its determinant is $p^{-1}$ and it has
an eigenvalue one, so its two distinct Frobenius
eigenvalues are $1,p^{-1}$.  Choose the corresponding
basis of $D=E$.  Over $R=L[[z]]$ set
\[
 \varphi_z=
 \begin{pmatrix}1+z&0\\0&p^{-1}(1+z)^{-1}\end{pmatrix},
       \qquad \varphi_z^*=p^{-1}\varphi_z^{-t}.
\tag{AG-local-polarized-deformation}
\]
Keep $F^0$, the entire filtration and $\beta$ fixed.
Every alternating form in dimension two is multiplied
by the determinant of a matrix.  Consequently
$\beta(\varphi_zx,\varphi_zy)=p^{-1}\beta(x,y)$.
The same fixed polarization identifies $\varphi_z$
with $\varphi_z^*$ for every $z$.  This checks the
actual Tate twist and dual Frobenius; the two operators
$1-\varphi_z$ and $1-\varphi_z^*$ are conjugate and
have identical determinants.

The map
\[
 a_z:D_R\longrightarrow D_R\oplus t_R,
       \qquad d\longmapsto((1-\varphi_z)d,\bar d)
\]
is a split injection with free rank-one cokernel $U_R$.
Indeed $a_0$ is injective by $H^0=0$, so one of its
maximal minors is nonzero in $L$ and hence a unit in
$R$.  The same statement holds for the actual dual
map, and the fixed polarization identifies its
cokernel with $U_R$.  Form
$H_R=U_R\oplus U_R^\vee$ with its hyperbolic symmetric
pairing.  At $z=0$ identify it with $(H,q)$ preserving
$U$: any complement can be made isotropic by replacing
a vector $c$ by
$c-q(c,c)h/(2q(h,c))$, where $h$ spans $U$.
Using the fixed de Rham pairing and this Tate pairing,
splice the finite sequence over $R$ with its actual
dual, in exactly the order of
(AG-local-spliced-complex).  The sequence is exact and
all its modules are finite free.  The determinant
construction therefore gives a regular isomorphism
\[
                   \Theta_R:\det H_R\xrightarrow{\sim}\det E_R
\]
which specializes to the original $\Theta$.  This uses
only finite free linear algebra and base change of
determinants; no family of Galois representations is
being asserted.

For completeness, compute that map over the fraction
field $K=L((z))$.  Both $A=1-\varphi_z$ and
$A^*=1-\varphi_z^*$ are invertible there.  The finite
cokernel has the identification
\[
 [(d,u)]\longmapsto u-\overline{A^{-1}d},
\]
which makes the formal exponential $t_K\to U_K$ an
isomorphism.  In the target basis $D_K\oplus t_K$,
the columns of $a_z$ followed by the lift $(0,u)$ of
an exponential basis have determinant $\det A$.
In the dual sequence, put a basis of its subspace
first and lifts of a basis of its quotient last.
Solving the quotient map uses $(A^{*t})^{-1}$.
Moving its rank-one subspace past the two crystalline
coordinates has sign $(-1)^2=1$, so its corresponding
determinant is $(\det A^*)^{-1}$.
Their ratio is one, by the conjugacy just proved.
Moreover each crystalline determinant has graded
degree two.  Moving or cancelling these factors
introduces no further Koszul sign.  After these
cancellations the map is precisely
\[
 \begin{gathered}
 \det H_K\longrightarrow U_K\otimes(H_K/U_K)
 \xrightarrow{\ \exp^{-1}\otimes\exp^*\ }
 t_K\otimes F^0_K,\\
 t_K\otimes F^0_K\longrightarrow\det E_K,
       \qquad u\otimes v\longmapsto v\wedge u.
 \end{gathered}
\tag{AG-local-generic-splice}
\]
This verifies both the Euler-factor cancellation and
the remaining order; in particular its last arrow
is $-u\wedge v$ in the opposite order.

By construction the formal dual exponential is the
adjoint of the exponential for the actual polarized
dual finite sequence.  Hence
$q(\exp(u),y)=\beta(u,\exp^*(y))/(\zeta/\tau)$.
Equation (AG-local-generic-splice) consequently gives
\[
            \Theta_K(x\wedge y)
                  =-q(x,y)\,\zeta/\tau
                    \qquad(x\in U_K).
\]
Both sides come from regular $R$-linear maps
$U_R\otimes H_R\to\det E_R$.  Their equality over
$K$ implies equality over the domain $R$, and then
at $z=0$.  For the actual crystalline representation,
the Weil--Deligne representation is unramified with
$N=0$, so its epsilon constant is one by Nakamura
\S3.3, p.~35, property (4).  The fixed filtration
gives the fixed factor
$\Gamma(V)=(-1)^{k-1}$ for weights $(k,1-k)$.
Multiplication by this factor and $\theta_{\mathrm{dR}}$
proves (AG-local-splice).

Nondegeneracy of $q$ gives $wx=-\Gamma(V)x$ on $U$.
Since $B_\epsilon$ is alternating, $w$ is skew-adjoint
for the symmetric form $q$, so it has trace zero.
Its other eigenvalue is therefore $\Gamma(V)$;
the two eigenvalues are distinct.  Thus the indicated
eigenspace is exactly $U$.  On the complementary locus
$D^{\varphi=1}=0$, AG-LOCAL-EXP makes both maps used
in BKNO Proposition 4.17 actual isomorphisms.  The
displayed calculation in its complete proof,
pp.~28--29, then gives the same formula with the
stated Weil--Deligne epsilon factor.  The new argument
above fills the exceptional case where that proof's
inverse exponential is unavailable.
\end{proof}

This repairs the rational de Rham determinant comparison
for the local signed line.  It does not make the ordinary
exponential or dual exponential nonzero in
AG-LOCAL-EXCEPTION, and it does not supply an integral
theta-period comparison or a global arithmetic cycle.
For a noncrystalline rank-two CM realization with
determinant $L(1)$, AG-LOCAL-WD already removes this
particular exponential obstruction.  The universal
epsilon-isomorphism and its global applications remain
the separate results cited by BKNO; the proof here
concerns their precise de Rham determinant specialization.

\subsection{Effective integral denominators for the Frobenius grading}

Let $X/\mathbb F_q$ be smooth, projective and geometrically
connected of dimension $d$.  Use $\ell$-adic cohomology for
$\ell\ne p$, crystalline cohomology after inverting $p$,
or a rational cohomology satisfying Katz--Messing's stated
Poincaré duality, weak Lefschetz and zeta-function hypotheses.
Let $F$ act by Frobenius pullback and set
\[
 G_i(T)=\det(T-F\mid H^i(X)),\qquad b_i=\deg G_i,
 \qquad P(T)=\prod_{i=0}^{2d}G_i(T).
\]
These monic polynomials lie in $\mathbb Z[T]$, and the
roots of $G_i$ have every complex absolute value $q^{i/2}$.
Thus distinct nonconstant $G_i$ are relatively prime.
Katz--Messing, Theorem 1 and Theorem 2(1), pp.~74--77,
show that the same $G_i$ work in crystalline cohomology
and that the rational grading projectors are algebraic
polynomials in Frobenius.  The following formulation gives
an explicit denominator and its exact integral meaning.

We use the correspondence action
$C_*(x)=p_{2*}(p_1^*x\cup\operatorname{cl}(C))$.
Accordingly $\Phi$ denotes the transpose graph of the
$q$-power Frobenius morphism, so its action is $F$;
its powers are again integral graph correspondences.
All assertions of idempotence below concern their actions
on cohomology, equivalently the specified homological
quotient of the correspondence ring.

\begin{proposition}[AG-FROB-PROJECTOR: an explicit resultant denominator]
For an index $i$ with $b_i>0$, put
\[
 Q_i=\prod_{j\ne i}G_j,\qquad
 R_i=\operatorname{Res}(G_i,Q_i)\in\mathbb Z\setminus\{0\}.
\]
There is an integral cycle $C_i$, an integer linear
combination of transpose graphs of Frobenius iterates and the diagonal,
whose action on $H^j(X)$ is $R_i\delta_{ij}$.
Hence $\pi_i=C_i/R_i$ is a rational algebraic K\"unneth
projector.  Its displayed denominator satisfies
\[
 0<|R_i|\le
 \prod_{j\ne i}
       (q^{i/2}+q^{j/2})^{b_i b_j}.
\tag{AG-frob-resultant}
\]
One common denominator for all these projectors is
$\mathcal D=\prod_{i<j}|\operatorname{Res}(G_i,G_j)|$,
where factors involving $G_i=1$ are one.
\end{proposition}
\begin{proof}
The adjugate of the integral Sylvester matrix gives polynomials
$A_i,B_i\in\mathbb Z[T]$ satisfying
\[
                  A_iG_i+B_iQ_i=R_i.
\]
Indeed the Sylvester map sends the coefficient vectors of
polynomials of degrees less than $\deg Q_i$ and $\deg G_i$
to those of $A_iG_i+B_iQ_i$; multiplying its adjugate by
the coefficient vector of the constant polynomial produces
the displayed identity, up to the harmless common sign
fixed by the convention for the resultant.
Set $C_i=B_i(\Phi)Q_i(\Phi)$.  Cayley--Hamilton gives
$G_j(F)=0$ on $H^j$.  For $j\ne i$ the factor $G_j$
in $Q_i$ makes the action zero.  For $j=i$ the Bezout
identity makes it $R_iI$.  This proves the cycle assertion
and the cohomological projector identities.

Writing $\alpha$ and $\beta$ for the roots of $G_i$ and
$G_j$ in a common complex algebraic closure gives
\[
 |\operatorname{Res}(G_i,G_j)|
   =\prod_{\alpha,\beta}|\alpha-\beta|
   \le(q^{i/2}+q^{j/2})^{b_i b_j}.
\]
Multiplicativity of resultants proves (AG-frob-resultant).
Each $|R_i|$ divides $\mathcal D$, proving the common
denominator statement.  Zero cohomology groups simply
contribute the polynomial one.
\end{proof}

\begin{proposition}[AG-FROB-ORDER: the exact arithmetic order behind the certificate]
The natural injection of finite free abelian groups
\[
 A=\mathbb Z[T]/(P)\longrightarrow
 B=\prod_i\mathbb Z[T]/(G_i)
\]
has cokernel of order $\mathcal D$.  Let $e_i\in B$
be one in its $i$th factor and zero elsewhere.  The least
positive integer $a_i$ for which $a_ie_i$ belongs to $A$
is exactly the order of the class of $e_i$ in $B/A$;
it divides $|R_i|$.  It is computable by the Smith normal
form of the displayed integral map.  This is the exact
denominator in the abstract Frobenius polynomial order;
an action on a particular cohomology representation can
have further relations and a smaller denominator.
\end{proposition}
\begin{proof}
For two relatively prime monic polynomials $G,Q$ there
is an exact sequence
\[
 0\longrightarrow\mathbb Z[T]/(GQ)
 \longrightarrow\mathbb Z[T]/(G)\oplus\mathbb Z[T]/(Q)
 \xrightarrow{\ (a,b)\mapsto a-b\ }
 \mathbb Z[T]/(G,Q)\longrightarrow0.
\tag{AG-frob-crt}
\]
For injectivity, a polynomial divisible by both $G$ and
$Q$ is divisible by $GQ$ over $\mathbb Q[T]$, and monic
division makes the quotient integral.  If
$a-b=Gu+Qv$, the polynomial $a-Gu=b+Qv$ lifts the
pair, proving exactness in the middle.  The last map is
surjective.  Its target is the cokernel of multiplication
by $Q$ on the free group $\mathbb Z[T]/(G)$.
The determinant of that multiplication map is
$\operatorname{Res}(G,Q)$, so the target has its absolute
value as order.  Induct on the number of factors and use
the multiplicativity of the resultant.  The indices
multiply to give $\mathcal D$.

The assertion about $a_i$ is the definition of the order
of an element in the finite quotient.  The Bezout identity
in AG-FROB-PROJECTOR proves $R_ie_i\in A$ and hence
$a_i\mid |R_i|$.  Choosing the monomial bases in all
factors gives an explicit integral matrix for $A\to B$;
its Smith normal form computes the order of each $e_i$
class.  This also explains the last qualification: a
representation annihilated by $P$ need not be faithful
as an $A$-module.
\end{proof}

\begin{proposition}[AG-FROB-COLLISION: exact bad primes for Frobenius-only projectors]
Fix a prime $\ell\ne p$ and put
$\Lambda_j=H^j_{\text{ét}}(X_{\overline{\mathbb F}_q},
\mathbb Z_\ell)/\mathrm{torsion}$.
A polynomial $f\in\mathbb Z_\ell[T]$ acting as the identity
on $\Lambda_i$ and as zero on every $\Lambda_j$, $j\ne i$,
exists if and only if $\ell\nmid R_i$.
This exact obstruction concerns the algebra generated by
Frobenius, and need not obstruct an integral algebraic
K\"unneth projector supplied by other cycles.
\end{proposition}
\begin{proof}
If $R_i$ is a unit, the Bezout polynomial
$f=B_iQ_i/R_i$ is integral over $\mathbb Z_\ell$.
The characteristic polynomial of $F$ on the free lattice
$\Lambda_j$ is $G_j$, because it becomes that polynomial
after tensoring with $\mathbb Q_\ell$.  Cayley--Hamilton
therefore proves the required identities integrally.

Conversely, $\ell\mid R_i$ means that the reductions of
$G_i$ and some $G_j$, $j\ne i$, have a common root
$\alpha\in\overline{\mathbb F}_\ell$.
On both nonzero spaces
$\Lambda_i\otimes\overline{\mathbb F}_\ell$ and
$\Lambda_j\otimes\overline{\mathbb F}_\ell$, that root
is an eigenvalue of the reduced Frobenius.  Applying a
putative $f$ to eigenvectors forces simultaneously
$\bar f(\alpha)=1$ and $\bar f(\alpha)=0$, a contradiction.
No semisimplicity assumption is needed.

For the distinction in the last sentence, take
$X=\mathbb P^1_{\mathbb F_q}$.  Here
$G_0=T-1$, $G_2=T-q$, and
\[
 \pi_0=\frac{q-\Phi}{q-1},\qquad
 \pi_2=\frac{\Phi-1}{q-1}.
\]
The least denominator in $\mathbb Z[\Phi]$ is exactly
$q-1$: evaluation of any integer polynomial at $q$ and
$1$ gives values congruent modulo $q-1$, and the displayed
formulas attain that bound.  Yet, for a rational point
$x\in\mathbb P^1(\mathbb F_q)$, the two integral cycles
$[x\times\mathbb P^1]$ and $[\mathbb P^1\times x]$
are precisely the two K\"unneth projectors in our convention.
Thus every prime dividing $q-1$ is a genuine obstruction
within the Frobenius polynomial algebra, even though these
integral cycle projectors already exist.
\end{proof}

For cote 16, this separates three pieces of the problem:
finite-field purity supplies the rational algebraic grading;
the exact resultants control its Frobenius-generated integral
order; and an integral inverse Lefschetz correspondence still
requires its own cycle construction.  Neither the displayed
projectors nor their denominators prove idempotence modulo
rational equivalence in the full Chow ring.  The latter
would require control of the homologically trivial error,
for example an appropriate nilpotence theorem.  The
corresponding analytic inputs A and M do not enter these
already-pure finite-field deductions.  The next subsection
supplies the exact nilpotence lift and, in the Kimura-finite
sector, proves the stronger rational Chow identity directly.

\subsection{Chow projectors: nilpotence, parity and exact denominator primes}

Keep $G_i,P,R_i,C_i,\Phi$ from AG-FROB-PROJECTOR, omitting
the indices with $G_i=1$, and put
$\mathcal D=\prod_{i<j}|\operatorname{Res}(G_i,G_j)|$.
All products of correspondences in this subsection mean
composition.  The following elementary correction states
precisely which ring must satisfy the nilpotence hypothesis.

\begin{proposition}[AG-FROB-NIL: an integral polynomial correction]
Let $A$ be a unital algebra over
$R=\mathbb Q$ or $R=\mathbb Z[1/\mathcal D]$, containing
the specified element $\Phi$.  Suppose that, in this ring,
\[
                         P(\Phi)^N=0\qquad(N\ge1).
\tag{AG-chow-nil}
\]
Put $e_i=C_i/R_i$ and define
\[
 s_N(x)=\sum_{j=N}^{2N-1}
             {2N-1\choose j}x^j(1-x)^{2N-1-j},
                   \qquad \widetilde e_i=s_N(e_i).
\tag{AG-chow-correction}
\]
Then the $\widetilde e_i$ are pairwise orthogonal
idempotents summing to one, and
\[
                 G_i(\Phi)^N\widetilde e_i=0.
\tag{AG-chow-primary}
\]
They are the unique idempotents in the commutative
algebra $R[\Phi]$ congruent to the respective $e_i$
modulo its nilpotent ideal $(P(\Phi))$.
As rational combinations of integral graph cycles,
$\widetilde e_i$ have denominator dividing
$R_i^{\,2N-1}$; one common denominator divides
$\mathcal D^{\,2N-1}$.

For $A=\operatorname{CH}^d(X\times X)_{\mathbb Q}$,
the assertions are rational Chow identities.  They
are identities in
$\operatorname{CH}^d(X\times X)_{\mathbb Z[1/\mathcal D]}$
if (AG-chow-nil) is known in that localized ring.
Knowing it only after tensoring with $\mathbb Q$
does not assert those integral identities.
\end{proposition}
\begin{proof}
The binomial theorem gives
\[
 s_N(x)+s_N(1-x)=1,\qquad
 x^N\mid s_N(x),\qquad (1-x)^N\mid1-s_N(x)
                         \quad\hbox{in }\mathbb Z[x].
\]
Thus $x^N(1-x)^N$ divides $s_N(x)(1-s_N(x))$.
The degree is exactly $2N-1$: its leading coefficient
is $(-1)^{N-1}{2N-2\choose N-1}\ne0$, by the
alternating binomial-sum identity obtained from
Pascal's rule.  This also covers $N=1$.

Write $e_i=E_i(\Phi)$, where
$E_i(T)=B_i(T)Q_i(T)/R_i$.  The Bezout construction
gives $E_i\equiv\delta_{ij}\pmod{G_j}$.
Consequently
\[
 s_N(E_i(T))\equiv\delta_{ij}\pmod{G_j(T)^N}.
\]
The polynomials $G_j^N$ are pairwise comaximal over
$R$, because their pairwise resultants are units.
The Chinese remainder theorem therefore makes
\[
 \begin{gathered}
 s_N(E_i)^2-s_N(E_i),\quad
 s_N(E_i)s_N(E_j)\ (i\ne j),\\
 \sum_i s_N(E_i)-1,
 \quad G_i^N s_N(E_i)
 \end{gathered}
\]
all divisible by $P^N$.  Evaluation at $\Phi$ and
(AG-chow-nil) prove every stated identity.

The same congruences show that $s_N(E_i)-E_i$ is
divisible by $P$, so these idempotents lift the
claimed classes modulo $(P(\Phi))$.

For uniqueness, any two commuting idempotents $a,b$
with $a-b$ in a nilpotent ideal satisfy
$a(1-b)=(a-b)(1-b)$ in that ideal.  The left side
is itself idempotent, hence is zero.  The same
argument gives $b(1-a)=0$, so $a=b$.
The elements in question commute because they are
polynomials in $\Phi$.  This argument also explains
why commuting idempotent products that are nilpotent
must vanish; it does not assume a semisimple ring.

Finally (AG-chow-correction) can be written as
\[
 \widetilde e_i=
 \frac{1}{R_i^{\,2N-1}}
 \sum_{j=N}^{2N-1}{2N-1\choose j}
                  C_i^j(R_i-C_i)^{2N-1-j}.
\]
Every numerator is an integral polynomial in graph
correspondences, proving the denominator claim.
These calculations take place in the chosen ring
$A$, which proves the last distinction as well.
\end{proof}

There is a sharper answer for Kimura-finite Chow motives
over the finite base field itself.  We use the following
located classical tensor facts from André,
\emph{Motifs de dimension finie}, Bourbaki no.~929,
Astérisque 299 (2005): Theorem 2.8, p.~124, gives
Kimura finiteness for motives of abelian type;
Corollaries 3.19--3.22, p.~132, identify the even/odd
dimension and give categorical Cayley--Hamilton;
\S4.2, p.~138, states that Frobenius is a tensor
automorphism of the identity on motives over
$\mathbb F_q$ for every adequate equivalence.
The coefficient field is $\mathbb Q$; the geometric
base field is allowed to have positive characteristic.

\begin{proposition}[AG-FROB-CHOW: exact Cayley--Hamilton in a finite-dimensional motive]
If $h(X)$ is Kimura finite over $\mathbb F_q$, then
\[
                     P(\Phi)=0
       \quad\text{in }\operatorname{CH}^d(X\times X)_{\mathbb Q}.
\tag{AG-chow-cayley}
\]
Consequently the original $e_i=C_i/R_i$, with no
nilpotence correction, are pairwise orthogonal
rational Chow projectors summing to the diagonal,
and $G_i(\Phi)e_i=0$ in that ring.  They are
represented by cycles whose denominators divide
$R_i$.  This applies in particular when $h(X)$ is
of abelian type over $\mathbb F_q$.

For every prime $\ell\ne p$, the rational degree-$i$ Chow
projector can be represented by a polynomial in
$\Phi$ with coefficients in $\mathbb Z_{(\ell)}$
if and only if $\ell\nmid R_i$.  Thus the set of
necessary denominator primes away from $p$ is exact
inside the Frobenius polynomial algebra.  These
rational Chow identities do not by themselves give
idempotence in the localized integral Chow ring.
\end{proposition}
\begin{proof}
Choose a Kimura decomposition defined over the base
field,
\[
                   h(X)=M_+\oplus M_-.
\]
An exterior power of $M_+$ vanishes and a symmetric
power of $M_-$ vanishes.  Realization in super vector
spaces shows that $H(M_+)$ is purely even and
$H(M_-)$ purely odd: a nonzero odd vector has
nonzero exterior powers of all orders in the super
convention, and a nonzero even vector has nonzero
symmetric powers of all orders.  Thus these two
realizations are exactly $H^{\mathrm{even}}(X)$ and
$H^{\mathrm{odd}}(X)$.  Frobenius is a natural
endomorphism of the identity category, so it
commutes with the two splitting idempotents and
preserves $M_+$ and $M_-$.

We make the characteristic-polynomial normalization
explicit.  For an even object $M$ of dimension $m$
and an endomorphism $f$, categorical Cayley--Hamilton
uses the monic polynomial
\[
 c_f(T)=\sum_{j=0}^m(-1)^j
                 \operatorname{tr}(\wedge^j f)T^{m-j}.
\tag{AG-categorical-char}
\]
It annihilates $f$.  In terms of the reciprocal
polynomial $\det(1-tf)$ used in the cited discussion,
this is $T^m\det(1-T^{-1}f)$.
One can check this normalization directly: partially
trace the vanishing antisymmetrizer on $M^{\otimes(m+1)}$
after applying $1\otimes f^{\otimes m}$.  Group
permutations by the length $j+1$ of the cycle through
the first position.  There are
${m\choose j}j!$ choices for that cycle, and summing
the remaining signed permutations contributes
$(m-j)!\operatorname{tr}(\wedge^{m-j}f)$.
Division by $(m+1)!$ gives
\[
 0=\frac1{m+1}\sum_{j=0}^m(-1)^j
                  \operatorname{tr}(\wedge^{m-j}f)f^j,
\]
which is $(-1)^mc_f(f)/(m+1)=0$.
The cited dimension result ensures
$\wedge^{m+1}M=0$ for this $m$; an even object of
dimension zero is zero.  This is the usual classical
tensor proof, included to fix the reciprocal convention.

Apply (AG-categorical-char) to $\Phi|M_+$.
Trace compatibility with realization identifies its
polynomial with
\[
                         P_+(T)=\prod_{i\ {
m even}}G_i(T).
\]
For the odd part, pass to the category of formally
$\mathbb Z/2$-graded objects in the same rational
tensor category.  The original category embeds
fully faithfully as the degree-zero objects.
Let $\Pi$ be the degree-one unit, whose self-braiding
is $-1$ and whose square is the unit.
Then $M_-\otimes\Pi$ is even, because
\[
 \wedge^j(M_-\otimes\Pi)
              \simeq \operatorname{Sym}^j(M_-)\otimes\Pi^j.
\]
Give $\Pi$ the identity endomorphism.  Tensoring by
it changes the categorical trace sign of $M_-$.
The realization of $M_-$ was purely odd, so the
new trace and all exterior coefficients are exactly
the ordinary traces and exterior coefficients of
$H^{\mathrm{odd}}(X)$.  Hence the monic polynomial
for $\Phi|M_-\otimes1_\Pi$ is
\[
                         P_-(T)=\prod_{i\ {
m odd}}G_i(T).
\]
Cayley--Hamilton makes it zero.  Tensoring back by
the invertible $\Pi$ and using the faithful embedding
gives $P_-(\Phi)|M_-=0$ in the original Chow category.
Likewise $P_+(\Phi)|M_+=0$.  Their product
$P=P_+P_-$ therefore annihilates all of $h(X)$,
proving (AG-chow-cayley).

Apply AG-FROB-NIL with $N=1$, for which $s_1(x)=x$.
This proves the Chow projector and denominator
claims.  If $\ell\nmid R_i$, the displayed polynomial
$C_i/R_i$ has coefficients in $\mathbb Z_{(\ell)}$.
If $\ell\mid R_i$ and $\ell\ne p$, any polynomial
with such coefficients would act integrally on the
free $\ell$-adic lattices and give the same degree
projector.  AG-FROB-COLLISION rules this out by the
common reduced Frobenius eigenvalue.  This proves
the exact prime-set assertion.  The tensor proof
was in the category with rational coefficients;
it has not proved that its zero correspondence
remains zero before torsion is removed.
\end{proof}

\begin{proposition}[AG-FROB-CYCLES: cycle degrees and a conditional integral annihilator]
Under the rational hypothesis of AG-FROB-NIL,
\[
      \widetilde e_i\operatorname{CH}^r(X)_{\mathbb Q}=0
                         \qquad(i\ne2r).
\tag{AG-chow-degrees}
\]
This holds in particular for the original $e_i$ of
AG-FROB-CHOW.  It places all rational codimension-$r$
cycles in the degree-$2r$ Chow summand, but does not
assert that the cycle map on that summand is injective
or surjective.

If (AG-chow-nil) holds in the localized integral
Chow correspondence ring, then, for $0\le r\le d$
and $i\ne2r$, the integral image
$\widetilde e_i\operatorname{CH}^r(X)_{\mathbb Z[1/\mathcal D]}$
is annihilated by the nonzero integer
$G_i(q^r)^N$.  Moreover
\[
       0<|G_i(q^r)^N|
                      \le(q^r+q^{i/2})^{Nb_i}.
\tag{AG-chow-annihilator}
\]
\end{proposition}
\begin{proof}
The $q$-power Frobenius pulls a codimension-$r$
cycle back to $q^r$ times itself.  This can be seen
integrally at the generic point of each irreducible
cycle: the smooth ambient variety has a regular
local ring of dimension $r$ there, and Frobenius
replaces a regular parameter ideal by its $q$th
powers, whose colength is $q^r$.  Frobenius is
finite flat on a smooth variety over the perfect
base field (étale locally this is the free
monomial-basis map on affine space), so this is its
flat pullback on Chow groups.  Our transpose-graph
convention is precisely that pullback.

Apply (AG-chow-primary) to
$\widetilde e_i\operatorname{CH}^r$.
Frobenius acts there by $q^r$, so the relation gives
$G_i(q^r)^N\widetilde e_i z=0$ for every cycle $z$.
All complex absolute values of roots of $G_i$ are
$q^{i/2}$; if $i\ne2r$, none can equal $q^r$.
Thus $G_i(q^r)$ is a nonzero integer.  Dividing
over $\mathbb Q$ proves (AG-chow-degrees), and
multiplying the bounds
$|q^r-\alpha|\le q^r+q^{i/2}$ over its roots proves
(AG-chow-annihilator).  When the relation was assumed
only rationally, this argument gives the rational
vanishing only; no integral annihilator is then claimed.
\end{proof}

For clarity, denominators and equalities are separate
issues.  A cycle in
$\operatorname{CH}^d(X\times X)_{\mathbb Z[1/\mathcal D]}$
can become zero over $\mathbb Q$ while remaining
nonzero torsion over that localization.  If the map
from this localized Chow group to its rationalization
is injective, the rational projector identities
above do hold integrally; alternatively the explicit
localized nilpotence hypothesis of AG-FROB-NIL suffices.
Without either input the rational tensor theorem alone
does not decide the integral identities.  Algebraically,
$\mathbb Z[t]/(\ell t,t^2-t)$ has a nonzero torsion
idempotent $t$ invisible after tensoring with $\mathbb Q$;
this illustrates the issue and is not offered as a
geometric Chow-ring example.

Thus in the finite-field Artin--Tate and abelian/CM
motivic sector, classical parity and Frobenius already
promote the grading to rational Chow correspondences
with the explicit resultants above.  This is a genuine
refinement of the cote 16 grading subproblem, without
a new GRH premise.  The torsion-sensitive integral
correspondence ring, the inverse Lefschetz cycle and
the general non-finite-dimensional motive retain their
own hypotheses; neither A nor M supplies them here.

\begin{proposition}[AG-FROB-Z: an effective obstruction to one coefficient localization]
For every odd prime $\ell$, the adopted input A supplies
a rational prime $p$ with
\[
                p\equiv1\pmod\ell,
                     \qquad p<2(\ell\log\ell)^2,
\tag{AG-frob-Z-bound}
\]
such that neither cohomological degree projector of
$\mathbb P^1_{\mathbb F_p}$ is a polynomial in
Frobenius with coefficients in $\mathbb Z_\ell$.
Nevertheless both projectors are restrictions of
integral product-cycle projectors on
$\mathbb P^1_{\mathbb Z}$.

More generally, for any two finite sets of rational
primes $S,T$, there is a prime $p\notin T$ where
these degree projectors cannot be represented by
polynomials in Frobenius over $\mathbb Z[S^{-1}]$.
An explicit witness is provided by (AG-frob-Z-bound)
for any odd prime $\ell\notin S$ greater than every
element of $T$.  Thus shrinking the arithmetic base
at finitely many primes and inverting finitely many
coefficient primes do not make this Frobenius-only
grading construction integral at every remaining
reduction.
\end{proposition}
\begin{proof}
Bach--Sorenson, \emph{Explicit bounds for primes in
residue classes}, Theorem 5.3, printed pp.~1731--1732,
gives a prime $p\equiv m\pmod q$ with
$p<2(q\log q)^2$ under ERH, for the coprime residue
classes in question.  Apply it with $m=1,q=\ell$.
The analytic functions in this cyclotomic case are
the primitive Dirichlet factors, so A supplies the
required hypothesis.  This is the residue-class
specialization of the same explicit-prime source
used in CX-GEO-004, with its sharper Theorem 5.3
constant.  The classical published numerical
estimate is an input here; its computational
verification is not being re-created.

Frobenius on $H^0$ and $H^2$ of $\mathbb P^1$ has
eigenvalues $1,p$.  A polynomial producing one of
the two degree projectors would have values $1,0$
at these two arguments, in one order or the other.
But every $f\in\mathbb Z_\ell[T]$ satisfies
$f(p)\equiv f(1)\pmod\ell$ when $p\equiv1\pmod\ell$.
This is a contradiction, also the special case
$R_i=p-1$ of AG-FROB-COLLISION.  Here $p>\ell$,
so the distinct-coefficient-prime requirement
$\ell\ne p$ is automatic.

The section at infinity of $\mathbb P^1_{\mathbb Z}$
gives the two integral relative product cycles
\[
 [\{\infty\}\times\mathbb P^1],\qquad
 [\mathbb P^1\times\{\infty\}].
\]
Their restrictions to every fiber are the two
integral projectors already computed in
AG-FROB-COLLISION.  Thus the obstruction belongs
to the polynomial algebra generated by the
individual fiber's Frobenius, not to the availability
of integral algebraic cycles.

Finally choose $\ell\notin S$ greater than every
prime in $T$.  The prime from (AG-frob-Z-bound)
satisfies $p>\ell$, so it is outside $T$.
The coefficient ring $\mathbb Z[S^{-1}]$ embeds
in $\mathbb Z_\ell$, and any proposed polynomial
there would contradict the argument above.
This proves the precise two-localization statement.
Existence of such obstructing fibers without an
explicit bound already follows from Dirichlet's
theorem; A supplies the displayed effective witness.
\end{proof}


\subsection{Actual inverse Lefschetz cycles on products of elliptic curves}
\label{ag:lefschetz-geometric}

The inherited statement H8 (run 20261112), as used in the Hodge
sketch Route~4, identifies the exterior-algebra map with a subset
inclusion matrix.  RETURN\_DP\_MULT, DP15 and Q4e, already check
its three-dimensional $J-I$ obstruction at cote~16, scan~14,
section~7.  We retain that attribution.  The following formulation
adds an actual cycle realization and the exact least denominator;
it does not identify the separate H8(iv) order or an automorphic
lifting lattice.

\begin{proposition}[AG-LEFSCH-SMITH: the full geometric lattice]
Let $g\ge1$ and $X=\prod_{i=1}^{g}E_i$ be a product of elliptic curves over a
subfield $F\subset\mathbb C$, with their origins $o_i$.  Put
$\theta=\sum_i\operatorname{pr}_i^*[o_i]$, and let $L$ denote
cup product with $\theta$.  Fix $0\le k\le g$ and put $a=g-k$.
The map $L^a:H^k(X(\mathbb C),\mathbb Z)\longrightarrow
H^{2g-k}(X(\mathbb C),\mathbb Z)$ is injective, and its finite
cokernel is the following direct sum of cyclic groups:
\begin{equation}
 \begin{gathered}
 \displaystyle\bigoplus_{r=0}^{\lfloor k/2\rfloor}\ 
 \bigoplus_{i=0}^{r}
 \left(\mathbb Z\Big/
 a!\binom{r+a-i}{r-i}\mathbb Z\right)^{m_{r,i}},\\
 m_{r,i}=\binom{g}{k-2r}2^{k-2r}
 \left[\binom{g-k+2r}{i}-\binom{g-k+2r}{i-1}\right].
 \end{gathered}
 \tag{AG-lefschetz-coker}
\end{equation}
Here $\binom n{-1}=0$, and a cyclic group of order one is zero.
The displayed list is a diagonal form; it need not be in
invariant-factor order.  In particular, the exact least positive
integer clearing the cohomological inverse is
\begin{equation}
 D_{g,k}=a!\,d_{g,k},\qquad
 d_{g,k}=\operatorname{lcm}_{0\le j\le\lfloor k/2\rfloor}
                         \binom{a+j}{j}.
 \tag{AG-lefschetz-denominator}
\end{equation}
For the integral divided-power map $L^{[a]}=L^a/a!$, remove
$a!$ from every cyclic order in (AG-lefschetz-coker); its exact
inverse denominator is $d_{g,k}$.
\end{proposition}
\begin{proof}
Integral K\"unneth gives a direct sum indexed by the subset $S$
of factors contributing $H^1$.  Write $s=|S|$ and $k=s+2r$.
The $H^1$ tensor has rank $2^s$; each remaining factor contributes
its generator in degree zero or two.  If $n=g-s=2r+a$, the
source basis in this block is indexed by the $r$-subsets $J$
of an $n$-set, and the target by its $(r+a)$-subsets $K$.
All raising operators have even degree.  Their squares vanish,
so the matrix is
\[
 a!\bigl(\mathbf 1_{J\subseteq K}\bigr)_{K,J}
       =a!\,W_{r,r+a}(2r+a)^{\mathsf t}.
\]
Wilson's integral diagonal-form theorem applies because
$r\le r+a\le(2r+a)-r$.  Its diagonal entries are
$\binom{r+a-i}{r-i}$ with multiplicity
$\binom ni-\binom n{i-1}$, for $0\le i\le r$.
A primary proof is Ghorbani--Khosrovshahi--Maysoori--Mohammad-Noori,
\emph{Inclusion Matrices and Chains}, Corollary~3, author p.~10;
Theorem~1 and its complete induction occupy pp.~8--9,
with the unimodular equivalence defined on p.~7.  Thus this
is equivalence over $\mathbb Z$, not just a rational eigenvalue
calculation.  Multiplication by $a!$ preserves the same two
unimodular changes of basis.  There are $\binom gs$ subsets
$S$ and $2^s$ copies for each, proving the asserted direct sum
and injectivity.

For an injective map of free lattices of the same rank,
$N(L^a)^{-1}$ is integral if and only if $N$ annihilates its
cokernel.  Its least such $N$ is therefore the least common
multiple of the diagonal entries.  Replacing $i$ by $r-j$
gives exactly the binomials $\binom{a+j}{j}$ with
$0\le j\le r$.  The block $r=\lfloor k/2\rfloor$ occurs
with $s=0$ or $1$ and nonzero $H^1$ tensor, so every entry
needed for (AG-lefschetz-denominator) really occurs.  Finally,
the same subset expansion gives an integral matrix for
$L^{[a]}$, and removes precisely the common factorial.
\end{proof}

\begin{proposition}[AG-LEFSCH-CYCLES: the sharp denominator is attained by cycles]
In AG-LEFSCH-SMITH there are integral Chow projectors
$\Pi_j\in\mathrm{CH}^{g}(X\times_F X)$ for every degree,
and a correspondence
\[
 C_k\in\mathrm{CH}^{g-a}(X\times_F X)
                   \otimes\mathbb Z[1/D_{g,k}]
\]
whose coefficients can in fact all be taken in
$D_{g,k}^{-1}\mathbb Z$, such that, in the Chow
correspondence ring with coefficients $\mathbb Z[1/D_{g,k}]$,
\begin{align*}
 C_k&=\Pi_k C_k\Pi_{2g-k},\\
 C_k L^a\Pi_k&=\Pi_k,&
 L^a C_k\Pi_{2g-k}&=\Pi_{2g-k}.
\end{align*}
Here $L=\Delta_{X*}\theta$ denotes the algebraic correspondence.
There is likewise an inverse to $L^{[a]}$ with coefficients
in $d_{g,k}^{-1}\mathbb Z$.  These denominator bounds are sharp
even if arbitrary algebraic correspondences are allowed:
there is an integral correspondence inducing
$N(L^a)^{-1}$ on $H^{2g-k}$ if and only if $D_{g,k}\mid N$;
replace $D_{g,k}$ by $d_{g,k}$ for the divided-power map.
No A, M, Hodge conjecture, or Chow-nilpotence hypothesis is used.
\end{proposition}
\begin{proof}
Use the action convention
$Z_*\alpha=(p_2)_*(p_1^*\alpha\cdot Z)$.
On a pointed elliptic curve $(E,o)$ the following are integral
cycles, with the indicated correspondence degrees:
\[
 \begin{gathered}
 \pi_0=[o\times E],\quad \pi_2=[E\times o],\quad
 \pi_1=[\Delta_E]-\pi_0-\pi_2,\\
 U=[o\times o],\qquad V=[E\times E].
 \end{gathered}
\]
The $\pi_j$ have degree zero, $U$ degree $1$, and $V$ degree
$-1$.  Direct composition of product cycles gives
\[
 VU=\pi_0,\qquad UV=\pi_2,\qquad U^2=V^2=0.
\]
Moreover $\pi_0,\pi_1,\pi_2$ are orthogonal idempotents
summing to the diagonal; $U$ and $V$ kill the $\pi_1$ summand.
For example $VU$ is the pushforward of $[o,o,E]$ from
$E^3$, and $UV$ that of $[E,o,o]$.  The other relations follow
by the same product intersections, using $[o]^2=0$ in the
Chow ring of a curve.  Consequently
$\pi_0,\pi_2,U,V$ are exact matrix units on the two Tate
summands, already modulo rational equivalence with integral
coefficients.  These are identities over the field $F$.

Take external products of the three projectors and sum all
patterns of total degree $j$ to obtain $\Pi_j$.
Fix a subset $S$ carrying $\pi_1$, and set
$N_S=\bigotimes_{i\in S}(E_i,\pi_1)$.
The other factors are copies of the two Tate summands.
External products of $U,V,\pi_0,\pi_2$ give matrix units
between every two patterns $J,K$ on those factors, tensored
with the identity on $N_S$.  This needs no algebraic basis
or splitting of the $H^1$ of any elliptic curve.
Composition of these cycles is exactly composition of their
integer matrices.

On this actual Chow summand, $L$ is the sum of the raising
matrix units $U_i$.  Indeed $\Delta_{E_i*}[o_i]=[o_i\times o_i]$,
and $U_i$ kills the $\pi_1$ factor.  The $U_i$ commute and
square to zero, so the matrix of $L^a$ is the very same
$a!W_{r,r+a}(2r+a)^{\mathsf t}$ used above, now as a
matrix of Chow-motive morphisms.  Invert this square matrix
over $\mathbb Q$.  Its inverse entries belong to
$D_{g,k}^{-1}\mathbb Z$ by its integral diagonal form.
Replace each entry by the corresponding product-cycle
matrix unit and sum over $S$.  A unit taking a pattern of
size $r+a$ to one of size $r$ has degree $-a$, so the result
has codimension $g-a$ on $X\times X$.
The exact matrix-unit relations prove all three displayed
Chow identities over $\mathbb Z[1/D_{g,k}]$ itself:
they arise from identities between integer matrices after
clearing denominators, so no hidden Chow torsion is discarded.
In particular this construction does not
lift merely homological identities by an unspecified
nilpotence theorem.  The divided-power construction is the
same, with the factorial removed; its raising correspondence
is already integral by the subset expansion.

Multiplying the constructed correspondence by $D_{g,k}$
produces an integral cycle.  Conversely, an integral
correspondence acts integrally on the torsion-free Betti
cohomology.  If its action on $H^{2g-k}$ is
$N(L^a)^{-1}$, the lattice criterion in AG-LEFSCH-SMITH
forces $D_{g,k}\mid N$.  This proves the asserted exact
optimality among all cycles, not only the product cycles.
\end{proof}

For $k=2$, the divided-power cokernel is exactly
$\mathbb Z/(g-1)$, recovering H8 and DP15 rather than
reclassifying them as new findings.  For $(g,k)=(5,4)$,
(AG-lefschetz-coker) gives
$\operatorname{coker} L\simeq(\mathbb Z/2)^{43}\oplus
\mathbb Z/6$, and the least denominator is $6$.
The product-intersection construction and its localized Chow
identities also work over any field, by the very same matrix-unit
relations; the global sharpness assertion above uses the stated
Betti lattice.  Thus one coefficient denominator works for this
construction on every good reduction of a product over a number
field.  The conclusion applies in particular to products
of CM elliptic curves used as test objects in the Hodge
sketch.  It closes the cycle-realization and denominator
comparison for this product polarization.  General
polarizations, nonsplit abelian varieties, the integral
order in H8(iv), and the specific theta-lift lattice still
require their own comparisons.


\subsection{Integral Fourier transforms: the exact cycle gate and a CM torsion obstruction}
\label{ag:fourier-cycle-gate}

HodgeSketch Route~4 cites Beckmann--de Gaay Fortman from memory,
and compares the theta-lift divisibilities with algebraic divided
powers and Fourier transforms.  We now locate the precise theorem
and distinguish its cohomological conclusion from identities of
integral Chow correspondences.  The latter distinction already
matters on the CM curve used in AG-C149-ELLIPTIC.

\begin{proposition}[AG-FOURIER-MIN: the classical minimal-class gate]
Let $(A,\theta)$ be a principally polarized complex abelian variety
of dimension $g\ge2$, and identify $A$ with its dual by the
polarization.  Write $\mathcal A^{2*}(A)$ for the image of the
integral cycle-class map, and put
$\gamma_c=\theta^c/c!\in H^{2c}(A,\mathbb Z)$.
The following assertions are equivalent:
\begin{enumerate}
\item $\gamma_{g-1}$ is the class of an integral one-cycle.
\item There is an integral, possibly inhomogeneous, cycle
$\Gamma\in\mathrm{CH}(A\times\widehat A)$ with
$\operatorname{cl}(\Gamma)=\operatorname{ch}(\mathcal P_A)$;
thus it induces the full integral cohomological Fourier transform.
\item Every integral Hodge class of degree $2g-2$ on $A$ is
algebraic.
\item For each positive-degree algebraic integral cohomology
class $x$ and every $n\ge1$, the integral class $x^n/n!$ is
algebraic.
\end{enumerate}
These are Beckmann--de Gaay Fortman, Theorem~1.1,
with the constructive proof in their Lemmas~3.4--3.5,
Theorem~3.8, Propositions~3.11--3.12 and section~4.1
(author pp.~8--13).  They imply the existence of an integral
Fourier transform \emph{up to homology}.  A weak integral
Fourier transform agreeing with the rational Chow transform
requires the stronger lift of
$\Theta^{g-1}/(g-1)!\in\mathrm{CH}_1(A)_{\mathbb Q}$ to
$\mathrm{CH}_1(A)$, where $\Theta$ is the symmetric polarization
class.  That stronger equivalence is their Theorem~3.8;
one must not replace this Chow class by its cohomology class.
\end{proposition}
\begin{proof}
Here is the construction, with its exact integral input.
The Pontryagin ring of an abelian variety has canonical divided
powers on its positive-dimensional cycles by Moonen--Polishchuk,
Theorem~1.6 and Corollary~1.7 (author pp.~6--7).
Suppose $z\in\mathrm{CH}_1(A)$ has class $\gamma_{g-1}$.
Let $j_1,j_2:A\to A\times A$ be the two coordinate inclusions,
and let $\Delta$ be the diagonal.  The integral cycle
\[
 R=(-1)^{g+1}(j_{1*}z+j_{2*}z-\Delta_*z)
\]
has cohomology class
$\rho=c_1(\mathcal P_A)^{2g-1}/(2g-1)!$ by
Beckmann--de Gaay Fortman's Lemma~3.5.  Define
\[
 \Gamma=(-1)^g\sum_{n=0}^{2g}\bigl((-1)^gR\bigr)^{[n]},
\]
using those Pontryagin divided powers; the zeroth term is the
origin.  Their Lemma~3.4 gives
$\operatorname{cl}(\Gamma)=\operatorname{ch}(\mathcal P_A)$.
Indeed $n!\operatorname{cl}(R^{[n]})=
\operatorname{cl}(R)^{\star n}$, and integral cohomology of
$A\times A$ is torsion-free, so these divided classes agree
with the rational exponential identity.  Terms of dimension
larger than $2g$ vanish.  This proves $1\Rightarrow2$ without
assuming that $z_{\mathbb Q}=\Theta^{g-1}/(g-1)!$ in Chow.

The cohomological Fourier transform is an integral isomorphism
$H^2(A,\mathbb Z)\to H^{2g-2}(\widehat A,\mathbb Z)$,
with the appropriate Tate twist.  Lefschetz $(1,1)$ and the
correspondence $\Gamma$ therefore give $2\Rightarrow3$.
Assertion~3 implies~1 because $\gamma_{g-1}$ is an integral
Hodge class.  The transpose Fourier correspondence, followed
by inversion and its usual sign, gives an inverse on
cohomology and preserves algebraic classes integrally.
Consequently $(-1)^g\mathcal F$ identifies the intersection
and Pontryagin rings on their algebraic cohomology images.
It sends positive codimension to positive dimension.  The
Pontryagin divided powers just cited transport to the unique
operations $x\mapsto x^n/n!$ in the torsion-free cohomology
ring.  Thus $2\Rightarrow4$.  Taking $x=\theta$ and
$n=g-1$ gives $4\Rightarrow1$.

For the last, stronger statement use the same construction
with a lift of $\Theta^{g-1}/(g-1)!$ in rational Chow.
Now the exponential identity holds in rational Chow, so
$\Gamma_{\mathbb Q}=\operatorname{ch}(\mathcal P_A)$.
Conversely a weak integral transform sends a symmetric
polarization divisor to an integral cycle whose rational
class is $(-1)^{g-1}\Theta^{g-1}/(g-1)!$.
This is exactly the additional condition in Theorem~3.8.
The cohomological construction alone supplied no such
rational-Chow equality.
\end{proof}

\begin{proposition}[AG-DP-ALG: the exact algebraic denominator of a divided power]
For a principally polarized complex abelian variety define
\[
 \kappa_c(A)=\min\{N\ge1:N\gamma_c\in\mathcal A^{2c}(A)\},
 \qquad 0\le c\le g.
\]
This integer exists and divides $c!$.
There is an integral correspondence in
$\mathrm{CH}^{g+c}(A\times A)$ inducing
$N L^{[c]}:x\mapsto N\gamma_c\cup x$ on \emph{all}
cohomology if and only if $\kappa_c(A)\mid N$.
Thus the exact algebraic denominator of this whole
raising operator is $\kappa_c(A)$, even though its
cohomological matrix is integral.
\end{proposition}
\begin{proof}
The integers $N$ with $N\gamma_c$ algebraic form a nonzero
ideal of $\mathbb Z$: it contains $c!$, since
$c!\gamma_c=\theta^c$ is algebraic.  Its positive generator
is $\kappa_c(A)$.  If $C$ is the stipulated integral
correspondence, apply it to $1=\operatorname{cl}([A])$.
The result is the class of the integral cycle $(p_2)_*C$,
and equals $N\gamma_c$.  This proves necessity.
Conversely choose an integral cycle $Z$ with class
$N\gamma_c$.  The integral correspondence $\Delta_*Z$
acts by cup product with that class, by the projection
formula.  This proves sufficiency.  The statement concerns
the operator on all degrees; it does not infer an obstruction
to an inverse restricted to a particular pair of degrees.
\end{proof}

\paragraph{AG-DP-EGS: a proved obstruction to the blanket geometric extension.}
We import the precise primary theorem of Engel--de Gaay
Fortman--Schreieder, \emph{Matroids and the integral Hodge
conjecture for abelian varieties}, arXiv:2507.15704v3
(27 March 2026; manuscript dated 31 March): Theorem~1.1
on p.~2, with the final proof in section~8.3, pp.~79--80.
For a very general principally polarized complex abelian
variety of dimension $g\ge4$, every integral algebraic
class in codimension $2\le c\le g-1$ is an even multiple
of $\gamma_c$.  Their Corollary~1.2, on p.~2 with proof
on p.~80, says the image is exactly $2\mathbb Z\gamma_c$
when $g=4$ or $5$.  This source theorem is a cited input:
the final reductions and their stated antecedents were read,
but the full 85-page degeneration/combinatorial proof and
its computational certificates were not independently audited.

It follows that $\kappa_c(A)$ is even in this range.
In particular
\[
 \kappa_2(A)=2\quad(g\ge4),\qquad
 \kappa_c(A)=2\quad(g\in\{4,5\},\ 2\le c\le g-1).
\]
The first equality also uses $\kappa_2(A)\mid2$;
the second is precisely their exact cycle-image result.
AG-DP-ALG therefore gives the exact factor two obstruction
to algebraic divided-power raising in these cases.
AG-FOURIER-MIN rules out an integral Fourier transform even
up to homology on these very general varieties.  In particular
it also rules out a weak integral Chow Fourier transform.
These are formal consequences with their full proofs just
given; they do not reprove the imported theorem.

Every principally polarized complex abelian variety has the
same integral polarized exterior cohomology model
$\bigwedge^*\mathbb Z^{2g}$: a symplectic basis writes
$\theta=\sum_i u_i\wedge v_i$.  Accordingly the Wilson
matrices and all the integral groups in AG-LEFSCH-SMITH
also apply to its cohomology, whether or not the pairs
come from elliptic factors.  Yet $\kappa_2=1$ on a product
of elliptic curves, where the subset cycles are integral,
and $\kappa_2=2$ on a very general variety of the same
dimension $g\ge4$.  This is a geometric obstruction to
transporting integral cohomological divided powers to
algebraic correspondences solely from that lattice.
It addresses the blanket extension of HodgeSketch Route~4,
not H8's correct exterior-algebra calculation.  The
very-general complex statement supplies no counterexample
on the CM locus or over $\overline{\mathbb Q}$.

\begin{proposition}[AG-FOURIER-ISOGENY: isogeny degrees annihilate the obstruction]
Let $A$ be a complex abelian variety of dimension $g$,
and suppose there is an isogeny
$f:A\times B\to J$ of degree $d$, where $J$ is a product
of Jacobians of smooth projective curves.  Then
\[
 d\,\mathcal Z_1(A)=0,\qquad
 \mathcal Z_1(A)=\frac{H^{2g-2}(A,\mathbb Z)\cap H^{g-1,g-1}(A)}
                         {\operatorname{cl}\mathrm{CH}_1(A)}.
\]
For principally polarized $A$, in particular
$\kappa_{g-1}(A)\mid d$.  Several such isogenies, with
possibly different $B$ and $J$, give the greatest common
divisor of their degrees as an annihilator.  In particular,
if that greatest common divisor is one and $A$ is principally
polarized, all the equivalent conditions of AG-FOURIER-MIN
hold.  This isolates the finite-prime degree gate in attempts to obtain the
integral Fourier property by adjoining Jacobian factors
and passing through isogenies.
\end{proposition}
\begin{proof}
Products of Jacobians satisfy the integral Hodge conjecture
for one-cycles by Beckmann--de Gaay Fortman, Theorem~1.2
and Corollary~4.1, whose proof uses the actual Abel--Jacobi
curve class and the minimal-class criterion.
Let $i:A\hookrightarrow A\times B$ use the origin of $B$,
and let $q:A\times B\to A$ be projection.  For an integral
Hodge class $\alpha$ of degree $2g-2$, the class
$f_*i_*\alpha$ is an integral one-cycle Hodge class on $J$.
Choose an integral cycle $Z$ representing it.
The cycle $q_*f^*Z$ is integral, since the isogeny is finite
flat.  On integral cohomology
$f^*f_*\beta=d\beta$: the left side is the sum of
translations by the kernel, and translations act trivially
on cohomology.  Also $q_*i_*=1$.  Hence
$\operatorname{cl}(q_*f^*Z)=d\alpha$.
The asserted divisibility and greatest-common-divisor
claim follow, the latter by an integer Bezout combination.
For the very general varieties in AG-DP-EGS, all these
degrees must be even.  This recovers the precisely scoped
obstruction in Engel--de Gaay Fortman--Schreieder,
Corollary~1.5, p.~3, whose proof is on p.~81; it does not
exclude the separate product/CM families.
\end{proof}

\begin{proposition}[AG-FOURIER-CM2: the full motivic identities need factor two]
Let $E:y^2=x^3-x$ over $\mathbb Q$, with origin $o=\infty$.
The modified diagonal
\[
 \Delta_e=\Delta_{123}-\Delta_{12o}-\Delta_{1o3}-\Delta_{o23}
                +\Delta_{1oo}+\Delta_{o2o}+\Delta_{oo3}
       \in\mathrm{CH}_1(E^3)
\]
has exact order two, even after base change to $\mathbb C$.
Here $\Delta_{12o}$ denotes the image $t\mapsto(t,t,o)$,
and the other terms use the indicated coordinates.

There is a weak integral Fourier correspondence $F_0$
for the principal polarization, with the exact square
identity $F_0^2=-[-1]^*$ in integral Chow correspondences.
However, no choice of integral Fourier correspondence can
satisfy both that square identity and
\[
 F\circ m_* = \Delta^*\circ(F\otimes F)
\]
exactly while inducing the usual rational Chow Fourier
transform.  The least positive factor $d$ for which
both identities can hold after multiplication by $d$
is $2$, over each of $\mathbb Q$, $\overline{\mathbb Q}$
and $\mathbb C$.  The canonical $F_0$ already realizes
this factor over $\mathbb Q$; its product-compatibility
defect has exact order two in the group of integral
Chow correspondences.
\end{proposition}
\begin{proof}
This is an explicit CM specialization of Moonen--Polishchuk,
Proposition~2.5, Remark~2.6 and Theorem~3.11
(author pp.~10--11 and 17--19).  We give the full cycle
detection because it records information invisible in
cohomology.

Use the integral projectors $\pi_0,\pi_1,\pi_2$ from
AG-LEFSCH-CYCLES.  The double cover $x:E\to\mathbb P^1$
is branched at $o$.  Pulling back the diagonal relation
on $\mathbb P^1\times\mathbb P^1$ gives
\[
 [\Delta_E]+[\Gamma_{[-1]}]=2\pi_0+2\pi_2.
\]
Thus $[-1]$ acts as $-1$ on the $\pi_1$ motive, integrally.
Applying $\pi_1^{\otimes3}$ to the small diagonal gives
exactly $\Delta_e$.  To check this directly, each
$\pi_2$ collapses its coordinate to $o$ and gives the
seven displayed terms; collapsing all three coordinates
pushes a curve to a point and gives zero.  A $\pi_0$
in one coordinate gives the corresponding coordinate
axis.  Applying $1-\pi_2$ in another coordinate kills
that axis, while two different $\pi_0$'s already give
zero.  Hence there are no additional terms.
Simultaneous inversion fixes the small diagonal and
commutes with these projectors, but acts as $(-1)^3$
on their tensor product.  Consequently $2\Delta_e=0$
over $\mathbb Q$.

For nonvanishing work over $\mathbb Q(i)$ or $\mathbb C$.
Put $P=(0,0)$ and let $\alpha(x,y)=(-x,iy)$.
Define $f:E^2\to E^3$ by $f(a,b)=(a,\alpha(a),b)$,
and set $T=(\operatorname{pr}_2)_*f^!$ on these
codimension-two cycles.  The map $f$ is a regular
immersion, the zero fiber of the smooth homomorphism
$(x,y,z)\mapsto y-\alpha(x)$; its normal line bundle
is trivial.  The kernel of $1-\alpha$ consists of
$o,P$, each with multiplicity one.  The seven
intersections therefore give
\begin{align*}
 T(\Delta_{123})&=[o]+[P],&T(\Delta_{12o})&=2[o],\\
 T(\Delta_{1o3})&=[o],&T(\Delta_{o23})&=[o],\\
 T(\Delta_{1oo})&=[o],&T(\Delta_{o2o})&=[o],\\
 T(\Delta_{oo3})&=0.
\end{align*}
For the final term the curve is contained in the
immersion, so the refined intersection is the first
Chern class of the trivial excess normal line and
is zero.  The others are the simple kernels of
$1-\alpha$, $\alpha$, or the identity just described.
It follows that $T(\Delta_e)=[P]-[o]$.
This is a nonzero element of exact order two in
$\mathrm{Pic}^0(E)\simeq E$: $P\ne o$ and
$\operatorname{div}(x)=2[P]-2[o]$.
Thus $\Delta_e$ is nonzero even over $\mathbb C$;
together with the preceding identity this proves
its exact order, also over $\mathbb Q$.

The weak transform and its square can be seen directly.
With $U=[o\times o]$ and $V=[E\times E]$, put
\[
 F_0=V-U-\pi_1\in\mathrm{CH}(E^2).
\]
The normalized Poincare bundle under the principal
polarization has $c_1(\mathcal P_E)=-\pi_1$;
the same product intersections give
$c_1(\mathcal P_E)^2=-2U$.
Thus $F_{0,\mathbb Q}=\operatorname{ch}(\mathcal P_E)$.
The exact matrix-unit relations from AG-LEFSCH-CYCLES
yield
\[
 F_0^2=-\pi_0+\pi_1-\pi_2=-[-1]^*.
\]
To compute the product defect over $\mathbb Q$ itself,
put $\epsilon=U+\pi_0+\pi_2-[\Delta_E]-V$ and let
$j_{12},j_{23}:E^2\to E^3$ be the coordinate inclusions.
Let $\delta(x,y)=(x,y,x)$.  We claim the exact identity
\begin{equation}
 j_{12*}\epsilon\star j_{23*}\epsilon-\delta_*\epsilon
       =\Delta_e.                                      \label{ag:eq:cm-cube-defect}
\end{equation}
Here $\star$ is Pontryagin product on $E^3$.
For completeness, the calculation is integral and
componentwise.  Write $A_i$ for the three coordinate
axes, $D_{ij}$ for the two-coordinate diagonals, and
$S_{ij}$ for the coordinate surfaces.  Set
$\alpha=A_1+A_2-D_{12}$ and
$\beta=A_2+A_3-D_{23}$.  Then
\[
 j_{12*}\epsilon=[o,o,o]+\alpha-S_{12},\qquad
 j_{23*}\epsilon=[o,o,o]+\beta-S_{23}.
\]
The zero-dimensional components in \eqref{ag:eq:cm-cube-defect}
cancel, and the one-dimensional difference is
\[
 \alpha+\beta-(D_{13}+A_2-\Delta_{123})=\Delta_e.
\]
To check dimension two, for an integral linear form
$u$ in the three coordinates write $K_u=u^*[o]$.
The only nonzero terms in $\alpha\star\beta$ after
cancellation are
\[
 \alpha\star\beta=K_y-K_{y-z}-K_{x-y}+K_{x-y+z}.
\]
For example $A_1\star A_3=K_y$,
$A_1\star D_{23}=K_{y-z}$,
$D_{12}\star A_3=K_{x-y}$, and
$D_{12}\star D_{23}=K_{x-y+z}$; each map onto
the displayed surface has degree one.  The other
nonzero terms $S_{12},S_{23}$ cancel in pairs.
The dimension-two difference is therefore
\[
 K_{x-y+z}-K_{x-y}-K_{y-z}+K_{x-z}+K_y-K_x-K_z=0.
\]
This is an equality in $\mathrm{Pic}(E^3)$ over
$\mathbb Q$: apply the theorem of the cube to the
symmetric line bundle $\mathcal O_E(o)$ and the
three maps $x,-y,z$, and use
$K_{x+z}+K_{x-z}=2K_x+2K_z$ from the theorem of the
square.  Thus no division by two has been made.
In dimension three,
$\alpha\star S_{23}=0=S_{12}\star\beta$:
in each expression two degree-one maps onto $E^3$
cancel, while the remaining term has a smaller image
and pushes forward to zero.  The dimension-four
term is zero by dimension.  This proves
\eqref{ag:eq:cm-cube-defect} over $\mathbb Q$.

Here is how this exact cycle identity measures the
motivic identity, including its sign.  For abelian
varieties $X,Y$ and $\eta\in\mathrm{CH}(X\times Y)$,
write
\[
 G_\eta=j_2^*\circ(\eta\star-)\circ j_{1*}.
\]
Moonen--Polishchuk Lemma~3.6, whose proof is valid
over an arbitrary field, identifies $G_\eta$ with
the integral correspondence $([-1]_X\times1_Y)^*\eta$;
thus $\eta\mapsto G_\eta$ is an isomorphism of
abelian groups.  Direct substitution gives
$F_0=-G_\epsilon$.  Permuting the second and third
coordinates in \eqref{ag:eq:cm-cube-defect}, using the
symmetry of $\epsilon$ and $\Delta_e$, gives
\[
 j_{13*}\epsilon\star j_{23*}\epsilon
       -(\Delta\times1)_*\epsilon=\Delta_e.
\]
The projection formula for Pontryagin product then
computes the entire dual product defect as
\begin{equation}
 B:=F_0\circ\Delta^*+
       m_*\circ(F_0\otimes F_0)=G_{\Delta_e},       \label{ag:eq:cm-dual-defect}
\end{equation}
where on the right $X=E^2$ and $Y=E$.
Indeed the first term on the left is
$-j_3^*\circ(((\Delta\times1)_*\epsilon)\star-)
\circ j_{12*}$, and the second is the same expression
with positive sign and kernel
$j_{13*}\epsilon\star j_{23*}\epsilon$.
Consequently $B$ has exact order two, over $\mathbb Q$
and after the stated base changes.

Finally put $I=[-1]^*$ and
$P=F_0\circ m_*-\Delta^*\circ(F_0\otimes F_0)$.
The already proved equality $F_0^2=-I$ makes $F_0$
invertible integrally, with inverse $-IF_0$.
Inversion commutes with $F_0$ and is compatible
with the multiplication and diagonal maps, so
\[
 P=I F_0\,B\,(F_0\otimes F_0).
\]
The identity follows by substituting
$B=F_0\Delta^*+m_*(F_0\otimes F_0)$ and
$F_0^2=-I$.  The outer compositions are integral
isomorphisms.  Thus $P$ too has exact order two:
$F_0$ satisfies the required product identity
up to factor two already over $\mathbb Q$.

To exclude another integral lift with factor one,
we use the precise classical criterion:
Moonen--Polishchuk, Definition~3.4 and Theorem~3.11,
say that, for an elliptic curve over an algebraically
closed field and any isomorphism to its dual, a motivic
integral Fourier transform up to factor $d$ exists
if and only if $d\Delta_e=0$.  The two identities
in the statement are exactly their conditions (ii)
and (iii), together with their rational-Chow
condition (i).  Since $\Delta_e$ has exact order two,
this proves that no other integral lift or isomorphism
to the dual can make $d=1$ work over $\mathbb C$.
Such a lift over $\overline{\mathbb Q}$ or $\mathbb Q$
would give one over $\mathbb C$ by base change, so
it is impossible there as well.  The explicit
calculation above provides factor two over $\mathbb Q$,
without extending the necessity theorem beyond its
stated field hypothesis.  This proves all the asserted
minimal factors; the general criterion remains
properly credited to its source.
\end{proof}

The last result does not contradict AG-LEFSCH-CYCLES:
those inverse-map identities were proved using an exact
matrix subring, while the Fourier assertion also requires
compatibility with the multiplication correspondence.
Nor does it obstruct the Fourier--Mukai derived equivalence.
It locates a nonzero integral Chow torsion class that must
be retained if HodgeSketch Route~4 or cote~16 is interpreted
as a strict motivic Fourier comparison.  Rational motives,
integral Betti cohomology, and the associated $L$-functions
do not detect this class; the relevant integral cohomology
is torsion-free.  The zero-location inputs A and M do not
enter this calculation of the exact compatibility factor.

\subsection{Archive and run consequences, with the original scope retained}
\label{ag:archive}

\paragraph{AG-C1.} Cote 1/79--80 and 111--112 gives the trace/Fredholm
framework, and SA23 geo-H-lidskii/geo-H-genus is its reported modern
reconstruction. A yields the earlier zeta positivity realization;
M yields AG-DET for the stated Hecke functions. AG-MOMENT and AG-HANKEL
give further consequences in that framework. MT-M29 now derives the
AG-DET operator from the canonical image-power grading of the locally
finite part of the actual adelic summation quotient, retaining the
zero multiplicities. Its positivity still uses A or M; an independent
geometric positivity or GGP-period comparison remains a separate
construction.

\paragraph{AG-C16.} Cote 16, scans 14--18, no.~7, concerns the integral
Lefschetz correspondence ring. Scans 44--50 contain the Coates letter's
conditions A, C and index positivity; the attached Hodge sketch uses a scan
numbering one higher than the GCT\_NSRH intake. These are algebraicity and
integrality problems. Neither A nor M constructs the required
correspondence or bounds its denominator. The sketch's explicit route is
the rank comparison on $X\times X$, not a claim that zero location
reverses Grothendieck's implication to the Weil conjectures.
AG-FROB-PROJECTOR/ORDER/COLLISION now give exact denominators
for the grading inside the Frobenius algebra; AG-FROB-CHOW
promotes it to rational Chow projectors for Kimura-finite
motives, and AG-FROB-NIL/CYCLES state the precise integral
nilpotence and torsion conditions. Those constructions do not by themselves give inverse Lefschetz.
AG-LEFSCH-SMITH/CYCLES now construct actual inverse
correspondences, with the exact least denominator, for
products of pointed elliptic curves with the product
polarization.  AG-FOURIER-MIN now locates the exact
minimal-class gate for an integral cohomological Fourier
kernel, while AG-DP-ALG isolates algebraicity of each
whole divided-power operator.  The primary theorem
imported in AG-DP-EGS disproves the blanket integral
extension on very general complex principally polarized
abelian varieties, without asserting a counterexample on
the CM locus or over $\overline{\mathbb Q}$.
AG-FOURIER-CM2 gives a separate arithmetic obstruction:
on $y^2=x^3-x$ a full integral motivic Fourier product
identity requires exactly factor two.  Neither obstruction
invalidates the constructed product inverse Lefschetz maps,
and the particular theta-chain comparison remains open.

\paragraph{AG-C34-74 and AG-C34-64.} Original PDF p.~74, manuscript p.~73, (2.5.5),
has now been visually checked: its multiplicative $1$ denotes the trivial
invariant group, equivalently additive $0$. The finite-field curve theorem
is the stated input, so inherited Q04 is unchanged. The original
definitions and full proof route for Theorem 1.1 have also been read.
AG-C34-TYPE identifies the literal obstruction in Q03, and
AG-C34-REPAIR proves the coherent intrinsic correction. AG-C34-INDEX
gives its finite Frobenius/Jacobian index. AG-C34-SECTION also proves
the formulation using a specified section and solvable $H$;
AG-C34-NONSOLV gives its abstract group-theoretic boundary.
No assertion of new purity or
of a consequence of A/M is made.

\paragraph{AG-C44.} Cote 44 Theorems 4--5 and the Serre letter, at the
locators in AG-SIGN, concern polarization and heights. AG-SIGN corrects
the central-derivative bridge but leaves cycle existence and the
height formula explicit. The supplied Hodge sketch also cites a separate
2026 CM-Hodge theorem for its abelian positivity consequences; those
claims are not consequences of either A or M and are carried with
their original dependency rather than independently certified here.

\paragraph{AG-C149 and AG-GT.} ORDER\_tonight, item 31, points to cote
149 scans 79--85, LM6 A.44--A.45, integral bridge B3, and the
irregular-prime/$\ell=2$ questions. GRH does not remove the integral
class-group and generator-relation gates in AG-LM6-IRREG/NORM.
For an actual punctured normal bundle, AG-C149-COMP now gives the
classical comparison under the explicit $\pi_2=0$ hypothesis, and
AG-C149-ELLIPTIC separates its global class from every section's
obstruction. AG-C149-NORMAL/CUSP now supply the actual normalizer
comparison in a family of four-punctured projective lines; every
arithmetic section even lifts through the full profinite normalizer.
A different specified family still requires its own comparison map.
The semi-graded Euler information in AG-EXT explains why the endpoint
zero theorems do not evaluate that extension data. In SA22 run 1,
A6-node-character (AG-NODE) is the quadratic Artin character times the cyclotomic
character; A gives effective detection of that finite quadratic sign.
M additionally covers finite nonabelian splitting fields in effective
Frobenius tests such as AG-FROB. These improvements do not reverse the
recorded refutation AG-INERTIA of run 2 A:inertia-maximality: finite-index inertia
subgroups have the same weights. Nor do they identify the profinite GT
kernel from the cyclotomic quotient. The existing supplied claim about
rational Deligne--Drinfeld freeness (CC-GT03, external 2026 input) is
carried separately from the full profinite Ihara question R-GTHAT.

\paragraph{AG-PURITY and AG-WM.} Finite-field purity and the usual
Grothendieck--Lefschetz trace formula are already classical. The full
geometric weight--monodromy statement remains represented in these
materials by BR-WM-02's finite exterior-power trace criterion; its
remaining inequalities are geometric hypotheses, not consequences of
global Dirichlet or Hecke zero locations. Likewise AG-STANDARD records that the supplied archive
questions about standard conjectures require algebraic cycles or
correspondences. No formal logical independence from A or M is asserted;
the precise point is that the displayed deductions do not supply those
inputs.

% END geometry_programme

% BEGIN complexity
\section{Complexity, effective algebra, and finite geometric tests}
\label{sec:complexity}

Here $A_D$ denotes the rational Dirichlet-GRH input designated by the task as
Alp\"oge's theorem. The letter $M$ denotes Masullo's all-number-field Hecke-GRH
claim, and every use of $M$ remains conditional. The input size of an arithmetic
problem includes the circuits, explicitly listed variables, rational constants
in binary, and any promised field representation. A geometric object without
an encoding has no associated complexity bound.

\subsection{Classical consequences and inherited results}

The following catalogue entries preserve the original finding identifiers of
the earlier runs. They are not advertised as new findings of this session.
An exact locator to an earlier proof is distinguished from a fresh reading of
that proof. In particular, the earlier complete file \path{CX_GRH.tex} is
not in tonight's attachment: its \path{IMPLICATIONS}, \path{LEDGER},
\path{QUESTIONS}, and targeted intake are available. Statements that depend
on its elaborate new finite-field constructions are carried from those records;
this session does not manufacture an independent verification of the absent
proof text.

\paragraph{CX-ALG-001 and CX-ALG-002: primality.}
Bach, \emph{Explicit bounds for primality testing and related problems}
(1990), Theorem 2, p.~372, and its application on p.~373 give the cutoff
$2(\log N)^2$ for a compositeness witness. Schoof, \emph{Four primality testing
algorithms} (2008), Theorem 2.2, pp.~105--106, states and proves the Miller test
with this cutoff using only quadratic Dirichlet characters. Therefore $A_D$
discharges its analytic input. Writing $N-1=2^s m$ with $m$ odd, test the usual
conditions $a^m=1$ or $a^{2^i m}=-1$ modulo $N$ for every base through the
cutoff, also checking nontrivial gcds. Even inputs are handled separately.
There are $O((\log N)^2)$ bases and each modular exponentiation takes a
polynomial number of bit operations. With $b=\lceil\log_2N\rceil$ and
multiplication cost $\mathcal M(b)$, the elementary implementation uses
$O(b^3\mathcal M(b))$ bit operations. This is the transfer already recorded
as CX\_GRH CX-N01 and ALPOGE\_GCT C-GRH-02.

The statement $\mathrm{PRIMES}\in\mathrm P$ was already unconditional:
Agrawal--Kayal--Saxena, \emph{PRIMES is in P}, Annals 160 (2004), Theorems 4.1
and 5.1, with their 2019 erratum. The new input validates this particular
short-witness test; it does not first place primality in $\mathrm P$.
Geometrically these are computations in $\mathbf G_m$ over $\mathbf Z/N\mathbf Z$
and, for AKS, finite cyclotomic algebras. The archival companion is cote 28,
scan 3 (group-action axioms), as mapped in the CX\_GRH return, not freshly
inspected from an archive scan in this session.

\paragraph{CX-ALG-003--CX-ALG-004: generators and nonresidues.}
Bach's Theorem 3, p.~372, gives units below $3(\log N)^2$ generating
$(\mathbf Z/N\mathbf Z)^\times$. Applied to the proper subgroup of $r$th
powers in $\mathbf F_p^\times$, it gives a polynomial search for a specified
prime $r\mid p-1$. Membership is tested by exponentiation to $(p-1)/r$;
the other factors of $p-1$ need not be known. This is CX\_GRH CX-N02.
The original Bach PDF was not freshly recovered here (the publisher returned
403); its exact locator is carried from the preceding checked run. Schoof's
primary source independently supplies the quadratic special case used above.
In the Kummer sequence this is finding a nonzero class in
$\mathbf F_p^\times/(\mathbf F_p^\times)^r$, or a nontrivial
$\boldsymbol\mu_r$-torsor. The established companion locator is SGA 4 IX 3.2.

\paragraph{CX-ALG-006--CX-ALG-008: primitive roots and extension fields.}
CX\_GRH CX-N03 constructs a certified generator of $\mathbf F_p^\times$ when
the full factorization of $p-1$ is supplied, by constructing each primary
component and multiplying them. Shoup, \emph{Searching for Primitive Roots
in Finite Fields} (1992), author version Theorem 3, p.~4, and Section 6,
p.~13, gives a least primitive root of order
$O(r^4(\log r+1)^4(\log p)^2)$, where $r=\omega(p-1)$,
and hence $O((\log p)^6)$ under the
rational Dirichlet hypothesis. This is a guaranteed short candidate list
without a supplied factorization, as CX\_GRH CX-N04 explicitly records.
A candidate list does not identify or certify a particular primitive root.

The stronger extension-field construction CX\_GRH CX-N13 supplies a specified
$r$th nonresidue in a supplied $\mathbf F_{p^m}$ in
$\operatorname{poly}(m\log p)$ time when $r\mid p-1$ is prime. Its quadratic
case with even $m$ already has an unconditional construction. The theorem is
inherited from the full targeted check in
\path{OURS_INTAKE_CX_GRH_20261121.md}, Section 1. A fresh source audit
found a nonvanishing gap in one cited auxiliary formula; CX-ALG-016 gives
the exact counterexample, and CX-ALG-017--CX-ALG-018 give a fully checked
repair and re-proof of the needed construction. The defining
irreducible polynomial is part of the input. These are computations with
Kummer covers inside an explicitly represented finite \`etale algebra.

\paragraph{CX-ALG-009--CX-ALG-010: polynomial factorization over finite fields.}
The inherited theorem CX\_GRH CX-N15 has bound
\[
                 (d^{\log d}\log q)^{O(1)}
\]
for a degree-$d$ polynomial over a supplied $\mathbf F_q$. Under $M$ its
analytic input is satisfied by $\zeta_{\mathbf Q}$ and the particular
$\zeta_{E_\nu}$ used in the construction, where
$E_\nu=\mathbf Q(\zeta_{2\nu},p^{1/\nu})$ and $\nu\leq d^2$.
Its exact component locators are Adleman--Lenstra (1986), Theorem 1, p.~350,
Algorithm A, pp.~350--351; Bach--von zur Gathen--Lenstra (2001), Theorem 2,
p.~7; and the 2008 Ivanyos--Karpinski--R\'onyai--Saxena version, Main Theorem,
p.~3, Lemma 4.2, p.~14, Proposition 4.6, pp.~16--17.
The fresh primary published IKRS article confirms the quasipolynomial
factorization-versus-automorphism distinction (Theorems 1.1 and 1.3), and
its unconditional noncommutative zero-divisor result (Theorem 1.7).

The displayed bound remains quasipolynomial in $d$. The residual $A_D$-only
problem in CX\_GRH CX-U02 is the high-valuation branch
$p>d$, $r\nmid p(p-1)$ and
$v_r(p^{\operatorname{ord}_r(p)}-1)>1$: construct the maximal coefficient
Teichm\"uller unit needed by that descent, or replace the descent. The new
CX-ALG-023--CX-ALG-026 below replace the degree-$sr$ field construction by
a prime-degree field, cover the additional case $r\mid m$ for the supplied
base field $\mathbf F_{p^m}$, and reduce the remaining analytic family to
pure prime-degree number fields. They do not solve the unrestricted
$A_D$-only gate. Geometrically factorization computes
primitive idempotents, hence connected components, of a finite \`etale
algebra (SGA 1 V Section 7). It is a different computational problem from
factoring an integer with unknown prime divisors.

\paragraph{CX-ALG-011--CX-ALG-013: arithmetic Picard groups.}
Bach (1990), Theorem 4, p.~375, and its corollary, p.~376, give the familiar
$12(\log|D_K|)^2$ class-group generator bound, using class-character Hecke
$L$-functions over $K$. Thus $M$ supplies the premise. This is CX\_GRH
CX-N07 and ALPOGE\_BRIDGES CC-F03/F03b. If the Hilbert class field is
abelian over $\mathbf Q$, $A_D$ suffices by Dirichlet factorization; the
weaker condition that $K/\mathbf Q$ be abelian does not establish this.
A factor base is not the full lattice of relations between its ideal
classes. Relation collection, saturation, regulator information and their
running times remain separate inputs in CX\_GRH CX-N08 and CX-U03.
In Grothendieck's language the class group is
$\operatorname{Pic}(\operatorname{Spec}\mathcal O_K)$ and principal divisors
supply the relation lattice. The established companion locators are SGA 4
IX Theorem 3.3, SGA 6 V 3.10 and VI 3.2--3.3, and cote 2 scan 5 as recorded
by CX\_GRH.

\paragraph{CX-ALG-019: a complete class-group algorithm retains randomness.}
For a supplied negative fundamental discriminant $-d$, Biasse--Erukulangara,
\emph{A proof of the conjectured run time of the Hafner--McCurley class
group algorithm}, Main Theorem, author PDF p.~4, Algorithm 4, p.~15,
and Corollary 1, p.~16, give a Las Vegas computation of
$\operatorname{Cl}(\mathcal O_{\mathbf Q(\sqrt{-d})})$ with expected bound
\[
 \exp\!\left(\left(\frac3{\sqrt8}+o(1)\right)
                         \sqrt{\log d\log\log d}\right).
\]
Their Conjecture 1, p.~4, explicitly states number-field Hecke-GRH; hence
$M$ supplies the hypothesis. The proof also uses the class-group Cayley
graph expansion of Jao--Miller--Venkatesan (2009), Theorem 1.1 and
Corollary 1.3. Rational quadratic Dirichlet-GRH alone does not supply all
these class-character estimates. This is the complete algorithmic transfer
already recorded as CX\_GRH CX-N12, not a new algorithm of this session.

The source's bounded trial returns a certified answer or FAILURE, with
success probability $1-d^{-1+o(1)}$ and the displayed time bound. Repeating
on detected failure gives the stated expected Las Vegas bound. The
certification in Algorithm 4 compares the order $h'$ of the group presented
by the collected relations with $2h^*$, where $h^*\le h<2h^*$ and $h$
is the actual class number. Since $h\mid h'$ and $h'>0$, the condition
$h'<2h^*$ forces $h'=h$. Thus the restart does not accept an incorrect
relation lattice. The full earlier proof and its detailed source check are
credited to CX-N12; the primary theorem, analytic-input definition,
Algorithm 4 and Corollary 1 were freshly inspected here.
The runtime is subexponential in $\log d$, not polynomial in that input
length, and the algorithm remains randomized. Its geometric output is the
finite Picard group, with actual principal-divisor relations rather than
only a factor base. The prior companions are SGA 4 IX Theorem 3.3,
SGA 6 V 3.10 and VI 3.2--3.3, and cote 2 scan 5.

\paragraph{CX-ALG-020: small-prime walks in ray class groups.}
Jao--Miller--Venkatesan (2009), Theorem 1.1, author PDF p.~3,
Remark 1.2(a) and Corollary 1.3, p.~4, give the following exact transfer.
For degree-$n$ $K$, modulus $\mathfrak m$, and
$Q=|D_K|N\mathfrak m$, take the narrow ray class group $G$ and the
multiset of degree-one prime ideals of norm at most $x$, coprime to
$\mathfrak m$, together with their inverses. Conditional on $M$, its
Cayley adjacency eigenvalues satisfy
\[
 \lambda_1=2\operatorname{li}(x)+O(n\sqrt x\log(xQ)),\qquad
 |\lambda_\chi|=O(n\sqrt x\log(xQ))\quad(\chi\ne1).
\]
The constants here are absolute. For fixed $n$ and $B>2$,
$x\ge(\log Q)^B$, and sufficiently large $Q$, a walk of length
$t\ge C_{n,B}\log|G|/\log\log Q$ reaches any fixed subset $T$ with
probability at least $|T|/(2|G|)$. Quotients inherit the bound.
The precise hypothesis is GRH for the Hecke characters actually occurring;
for $K=\mathbf Q$, $A_D$ suffices. Random walks remain randomized.
Geometrically these are bounded-prime correspondences on arithmetic
Picard/ray-class quotients, complementing CX\_GRH CX-N12; the Picard
companions are SGA 4 IX Theorem 3.3 and cote 2 scan 5.

\paragraph{CX-ALG-022: the ordinary elliptic isogeny graph.}
The same paper's Theorem 1.5, p.~5, applies conditionally on $M$ to ordinary
elliptic curves over a fixed $\mathbf F_q$, with fixed point count and
the same geometric endomorphism order. The graph of isogenies of prime
degree less than $(\log(4q))^B$, $B>2$ fixed, has the stated expansion
bound. Theorem 1.6, pp.~6--7, transfers a discrete-logarithm solver working
on a fraction $\mu$ of that level to every curve on it with an expected
$\mu^{-1}\operatorname{poly}(\log q)$ overhead. These are conditional
classical transfers; they do not supply such a solver. The Hecke input is
for imaginary-quadratic ring-class characters, rather than an assertion
about a non-CM elliptic curve's global L-function. The geometric link is
the CM ideal-class action (Proposition 3.1, p.~11) and Deuring lifting
(Theorem 3.3, p.~12); ORDER item 20 supplies the abelian-scheme/isogeny
programme, alongside the inherited Picard comparison of cote 2 scan 5.

\paragraph{CX-HN-001--CX-HN-004: the exact Nullstellensatz input.}
Koiran's DIMACS 96-27, Section 2, Theorem 2, pp.~5--6, proves
$\mathrm{HN}\in\mathrm{AM}\subseteq\Pi_2^{\mathrm P}$ under its GRH input.
The August 8, 1997 erratum corrects the lower estimate in Theorem 8 and a
prime-product expression in Theorem 10; it explicitly preserves the main
statements. The primary author-hosted report and the entire erratum were
retrieved in this session. Section 5, Theorem 9 and Corollary 1, printed p.~12 (PDF page~13),
show that the analytic estimate is for the bounded primitive solution
field itself. A Galois closure is not required. These are the earlier
CX\_GRH CX-K01--K02; the circuit and homogeneous nonzero forms are also
stated in GCT V Theorem 2.10, p.~25.

Consequently $M$ gives the unrestricted transfer, whereas $A_D$ gives the
bounded abelian-witness promise of CX\_GRH CX-K03. All the arithmetic size
bounds in that promise matter: field degree, logarithmic discriminant and
denominators have the stated single-exponential bounds in the input length.
The universal replacement by an abelian solution field was already refuted
by $x^3-2$ in CX\_GRH CX-K04 and ALPOGE\_GCT C-GRH-03. This refutes that
specific replacement, not every possible future $A_D$-only proof.
The geometric object is a finite \`etale spread of the solution field over
an open of $\operatorname{Spec}\mathbf Z$; degree-one Frobenius orbits give
rational reductions. The companion locators are SGA 1 V Section 7 and
cote 98 scans 11--12 as carried in CX\_GRH.

\paragraph{CX-NN-001--CX-NN-004: normalization maps.}
The preceding CX\_GRH CX-K05 refines GCT V Theorem 3.4, pp.~27--28, to
$\mathrm{FP}^{\Sigma_2^{\mathrm P}}$ under $M$ for equation-given,
standard-graded algebras with explicit ambient variables and polynomial
output bits. The source proof's prefix-existence queries already have
$\Sigma_2$ form. Its unconditional polynomial-space statement remains
available. The resulting map
$\operatorname{Spec}R\to\mathbf A^{\dim R}$ is finite; the companion is
EGA II 6.1.4--6.1.5, and cote 98 scan 44 in the prior crosswalk.

The historical EXPH transfers for succinct small separating systems are
GCT V Theorems 4.10, p.~34, and 5.12, carried in CX\_GRH CX-K06. The
optimal-cardinality h.s.o.p. bounds of Theorems 4.1, p.~30, and 5.5 retain
their efficient defining-equation premise. Their output conventions cannot
be silently changed. Earlier unconditional improvements already exist:
Forbes--Shpilka, ECCC TR17-163, Theorem 7.1, p.~21, constructs the relevant
hitting set in polynomial space for unary $n,s,r$; its introductory
Theorems 1.2--1.3 record normalization consequences. Guo--Saxena--Sinhababu
(2019), Theorems 1.5--1.6, are recorded in CX\_GRH CX-K07. These are not
new consequences of tonight's analytic inputs.

\subsection{New formulations with complete proofs}

\paragraph{CX-ALG-005: deterministic prime-field square roots.}
\emph{Statement.} Under $A_D$, given an odd prime $p$ and $a\in\mathbf F_p$,
one can find all roots of $X^2-a$ or certify their absence in deterministic
time polynomial in $\log p$.

\emph{Proof.} The case $a=0$ is immediate. Euler's criterion decides whether
$a$ is a square. If it is not, return no roots. Otherwise find a quadratic
nonresidue by the deterministic search of CX-ALG-004. Apply the classical
Tonelli--Shanks algorithm with this nonresidue. More specifically, Algorithm
3.34 of Menezes--van Oorschot--Vanstone, \emph{Handbook of Applied Cryptography},
p.~100, has only one random step, its choice of the nonresidue in Step 2.
All its remaining steps are modular exponentiation, gcd inversion and a
loop of length at most $v_2(p-1)\leq\log_2p$; Fact 3.35 gives a polynomial
bit bound for those operations. Replacing Step 2 by the bounded deterministic
search preserves polynomial time and the algorithm's proved correctness.
The resulting roots are $r,-r$. Completing the square factors every quadratic
with nonzero leading coefficient; in characteristic two the separate field
case is trivial for $\mathbf F_2$. This does not provide a square-root oracle
modulo an integer whose factors are unknown.\hfill$\square$

\paragraph{CX-ALG-014: a compositeness witness need not give the usual gcd factor.}
\emph{Statement.} The rule ``take a Miller witness and compute
$\gcd(a,N)$ or $\gcd(a^m\pm1,N)$'' does not always output a nontrivial divisor.

\emph{Proof.} Take $N=35$, so $N-1=2\cdot17$, and $a=2$. Since
$2^{12}=1\pmod{35}$, one has $2^{17}=32\pmod{35}$. This is neither $1$ nor
$-1$ modulo $35$, so $2$ is a Miller witness. Yet
\[
 \gcd(2,35)=\gcd(32-1,35)=\gcd(32+1,35)=1.
\]
This is a counterexample to precisely the displayed extraction rule.
It is not a lower bound for general integer factorization and does not
rule out other algorithms.\hfill$\square$

\paragraph{CX-ALG-015: integer factorization with a supplied exponent multiple.}
\emph{Statement.} Given integers $N\ge2$ and $T\ge1$, with the promise
\[
       \lambda'(N):=\operatorname{lcm}_{p\mid N}(p-1)\mid T,
\]
the complete factorization of $N$ is computable deterministically in time
polynomial in $\log N+\log T$ under $A_D$. This is the input-length form
of Miller (1976), Lemma 5 and Theorem 3, p.~310. In particular, a supplied
value of $\varphi(N)$ or $\lambda(N)$ meets the promise. Neither value is
asserted to be available for an arbitrary $N$.

\emph{Proof of the stated input-length form.} Remove powers of two, test
primality, and detect perfect powers by integer-root computations. These
are polynomial-time operations. For a perfect power, recurse on its base
and multiply the resulting exponents. It remains to split an odd composite
with at least two distinct prime divisors. Write $T=2^s t$ with $t$ odd.
For every $a\leq\lceil3(\log N)^2\rceil$, first compute $\gcd(a,N)$;
return a proper nontrivial gcd if found, and otherwise test, for units $a$,
\[
             \gcd\bigl(a^{2^j t}-1,N\bigr),\qquad 0\leq j<s.
\]
Powers are computed modulo $N$ by repeated squaring.

To prove that some test splits $N$, put
$e_p=v_2(p-1)$ and $e=\max_{p\mid N}e_p$. If $e_p=e$ and $e_q<e$ for
two divisors, the quadratic character $a\mapsto(a/p)$ is a nontrivial
character of $(\mathbf Z/N\mathbf Z)^\times$. If all $e_p=e$, use
$a\mapsto(a/p)(a/q)$ for distinct divisors $p,q$ instead. This too is
nontrivial, by the Chinese remainder theorem. Bach's generating bound
CX-ALG-003 implies that some tested unit has character $-1$.
In the first case its reduction modulo $p$ has multiplicative order with
2-part $2^e$, whereas the 2-part modulo $q$ is smaller. In the second
case one Legendre symbol is $-1$ and the other is $1$, with the same
consequence after interchanging $p,q$. Since $t$ is a multiple of the odd
part of each $p-1$, one obtains
\[
       a^{2^{e-1}t}=-1\pmod p,\qquad
       a^{2^{e-1}t}=1\pmod q.
\]
The corresponding gcd is divisible by $q$ and not by $p$, so is proper.
The promise gives $s\ge e$, hence this test occurs in the list. This
argument uses only reductions modulo the distinct primes; repeated prime
factors in $N$ cause no problem.

Every divisor $D$ encountered in the recursion still satisfies
$\lambda'(D)\mid T$. A recursion tree has $O(\log N)$ leaves and nodes
after perfect-power compression. Each node uses polynomially many bases
and at most $\log_2T$ squarings per base, so the total bit cost is
polynomial in the displayed input length. Geometrically the known exponent
annihilates the multiplicative groups of the reduced residue fibers;
quadratic characters separate two unknown fibers and a gcd recovers their
different supports. This is a precise promise case of the integer-factoring
gate in CX-ALG-014, not an algorithm for producing the exponent promise.
\hfill$\square$

\paragraph{CX-ALG-016: the fixed-generator resolvent needs a nonzero test.}
\emph{Statement.} The literal fixed-generator formula in
Bhargava--Ivanyos--Mittal--Saxena (2017), Theorem 3.1, author PDF p.~3,
can be zero. This does not refute their nonresidue-construction conclusion;
CX-ALG-017 supplies the missing nonzero choice.

\emph{Proof.} Let $q=7$, $r=3$, $d=3$, and
$E=\mathbf F_7[X]/(X^3-2)$, with $\alpha=X\bmod(X^3-2)$.
The nonzero cubes in $\mathbf F_7$ are $1,-1$, so this cubic has no root
and is irreducible. Take the primitive cube root $\zeta=4$. Then
$\alpha^7=4\alpha$ and $\alpha^{49}=2\alpha$, whence
\[
 L=\alpha+4\alpha^7+2\alpha^{49}
   =(1+16+4)\alpha=0.
\]
Consequently $L^{(7^3-1)/3}=0$, whereas the stated right-hand side is
$\zeta^{-1}=2$. The displayed formula and its hypotheses were checked
against the rendered original PDF page, not merely its text extraction.
The invalid step is division by $L$ when inferring
$L^{q^k-1}=\zeta^{-1}$ from $L^{q^k}=\zeta^{-1}L$.
\hfill$\square$

\paragraph{CX-ALG-017: a basis-search repair for the resolvent.}
\emph{Statement.} Let $E=\mathbf F_q[X]/(f)$ be a supplied field of degree
$d=rk$, where $r$ is prime. Suppose a primitive $r$th root of unity
$\zeta\in\mathbf F_q$ is supplied, and assume $r\mid q-1$ for odd $r$,
or $4\mid q-1$ when $r=2$. Then an $r$th nonresidue in $E$, and its
nonresidue norm in $\mathbf F_q$, can be found deterministically in time
polynomial in $d\log q$. No analytic hypothesis is needed for this statement.

\emph{Proof.} Write $Q=q^k$ and define the $\mathbf F_q$-linear map
\[
       P(z)=\sum_{i=0}^{r-1}\zeta^i z^{Q^i}.
\]
As a polynomial in $z$, it is nonzero and has degree
$Q^{r-1}<Q^r=|E|$. Thus it cannot vanish at every element of $E$.
Since it is linear, it cannot vanish on every element of the power basis
$1,\alpha,\ldots,\alpha^{d-1}$. Scan this basis and choose
$\beta=P(\alpha^j)\ne0$. Cyclically shifting the sum and using
$z^{Q^r}=z$ gives $\beta^Q=\zeta^{-1}\beta$, hence
$\beta^{Q-1}=\zeta^{-1}$.

Put $S=1+Q+\cdots+Q^{r-1}$. For odd $r$, write $Q=1+rc$ and expand
modulo $r^2$ to get
\[
 S\equiv r+rc\frac{r(r-1)}2\equiv r\pmod{r^2}.
\]
For $r=2$, the hypothesis gives $S=1+Q\equiv2\pmod4$.
In both cases $S/r\equiv1\pmod r$. Therefore
\[
 \beta^{(q^d-1)/r}
   =\left(\beta^{Q-1}\right)^{S/r}=\zeta^{-1}\ne1.
\]
This proves that $\beta$ is an $r$th nonresidue. Its norm satisfies
\[
 \left(N_{E/\mathbf F_q}\beta\right)^{(q-1)/r}
      =\beta^{(q^d-1)/r}=\zeta^{-1},
\]
so is also a nonresidue. There are $d$ basis candidates and $r\le d$
terms per test; all exponent bit lengths are $O(d\log q)$.
Field arithmetic in the supplied representation and these exponentiations
therefore take polynomial time.\hfill$\square$

For the inherited CX\_GRH CX-N13, obtain $\zeta$ under $A_D$ from a
prime-field nonresidue. If odd $r\mid m$, apply the theorem with
$q=p,d=m$. If $r\nmid m$, a prime-field nonresidue $a$ remains a
nonresidue upstairs, since
\[
 a^{(p^m-1)/r}
   =\left(a^{(p-1)/r}\right)^{1+p+\cdots+p^{m-1}}
   =\left(a^{(p-1)/r}\right)^m\ne1.
\]
The remaining quadratic-even case is proved next. Thus the prior
construction is preserved with an explicit nonvanishing justification.
In the Kummer-cover language, the repair chooses a nonzero vector in the
specified Frobenius eigenspace before passing to its torsor class; a
particular chosen generator need not have a nonzero projection.

\paragraph{CX-ALG-018: the even-degree quadratic case, with no GRH.}
\emph{Statement.} For odd $p$ and even $m$, a supplied
$E=\mathbf F_p[X]/(f)$ of degree $m$ admits a deterministic
$\operatorname{poly}(m\log p)$ construction of a quadratic nonresidue.
This is a fresh complete re-proof of the even-degree clause of inherited
CX\_GRH CX-N13, not a new claim of priority over BIMS17 Corollary 1.3.

\emph{Proof.} Put $Q=p^{m/2}$ and $\alpha=X\bmod f$.
If $p\equiv1\pmod4$ or $4\mid m$, then $Q\equiv1\pmod4$.
The element $u=\alpha^Q-\alpha$ is nonzero, since $\alpha$ has degree
$m$ and is not in the subfield of degree $m/2$. It satisfies $u^Q=-u$.
Consequently
\[
       u^{(p^m-1)/2}=(-1)^{(Q+1)/2}=-1,
\]
so it is the required nonresidue.

It remains that $p\equiv3\pmod4$ and $m\equiv2\pmod4$.
Compute the subfield $K=\mathbf F_{p^2}\subseteq E$ as the kernel of
the $\mathbf F_p$-linear map $z\mapsto z^{p^2}-z$.
Choose $v\in K\setminus\mathbf F_p$ from a computed basis and set
$u=v^p-v\ne0$. Then $u^p=-u$ and $c=u^2\in\mathbf F_p$ is a
nonresidue: a square $c$ would force $u\in\mathbf F_p$.
Because $-c$ is a square, the exponent formula
$t=(-c)^{(p+1)/4}$ gives $t^2=-c$. Thus $i=u/t$ satisfies $i^2=-1$,
and $K=\mathbf F_p(i)$ is explicitly represented.

Start with $z=i$, of order four. Its norm to $\mathbf F_p$ is one.
Whenever the current $z=a+bi$ has norm one, its order is a power of two
at least four; hence $b\ne0$, as the 2-part of $\mathbf F_p^\times$ has
order two. The nonzero numbers
\[
                 \delta_+=(a+1)/2,\qquad \delta_-=(a-1)/2
\]
have product $-b^2/4$, a nonsquare. Exactly one is a square in
$\mathbf F_p$. Determine which by Euler's criterion, choose its square
root $x=\delta^{(p+1)/4}$, and set $y=b/(2x)$.
Using $a^2+b^2=1$ verifies directly that
$(x+yi)^2=a+bi=z$. The norm of this root is $1$ if
$\delta=\delta_+$ and $-1$ if $\delta=\delta_-$.
Replace $z$ by this root, thereby doubling its order, until the norm is
$-1$.

This loop terminates after $O(\log p)$ steps. Indeed the norm-one subgroup
has order $p+1$, whereas the 2-part of $K^\times$ has order
$2^{v_2(p^2-1)}=2^{1+v_2(p+1)}$. An element's order cannot continue
doubling inside the former subgroup. At termination
\[
        z^{(p^2-1)/2}
          =N_{K/\mathbf F_p}(z)^{(p-1)/2}=-1.
\]
The extension $E/K$ has odd degree $m/2$, so this $z$ remains a nonresidue
in $E$. All operations are modular exponentiation or linear algebra in
the supplied $m$-dimensional field, with polynomial bit cost.
\hfill$\square$

The quadratic square-root algorithm then applies in $E$ with this known
nonresidue. In characteristic two the separate square-root problem is
inverse Frobenius, $a\mapsto a^{2^{m-1}}$, and there is no quadratic
nonresidue. These are finite-\'etale and Kummer computations; SGA 1 V
Section 7 and SGA 4 IX 3.2 are the same companions as for CX-ALG-008.

\paragraph{CX-ALG-021: the genus quotient only needs rational Dirichlet-GRH.}
\emph{Statement.} For an imaginary quadratic field $K$, put
$C=\operatorname{Cl}(\mathcal O_K)$ and $G=C/C^2$. The small-prime
spectral and mixing bounds in CX-ALG-020 hold on $G$ under $A_D$, with
$n=2$, $\mathfrak m=(1)$ and $Q=|D_K|$. No Hecke-GRH assumption for
the remaining, non-genus characters of $C$ is needed. If $C$ has exponent
two, this covers the whole class-group graph.
Moreover, the images of prime ideals of prime norm at most
\[
                         (8\log|D_K|+15)^2
\]
generate $C/C^2$ under $A_D$.

\emph{Proof.} Let $H$ be the Hilbert class field of $K$.
By class field theory $\operatorname{Gal}(H/K)=C$. Complex conjugation
preserves $H$: the conjugate of an unramified abelian extension of $K$
is again such an extension, and the maximal one is unique. Its action on
$C$ is inversion, since
$\mathfrak a\overline{\mathfrak a}=(N\mathfrak a)$ is principal.
Actual complex conjugation has order two and restricts nontrivially to
$K$, so
\[
       \operatorname{Gal}(H/\mathbf Q)=C\rtimes\langle c\rangle,
                    \qquad cac^{-1}=a^{-1}.
\]
The subgroup $C^2$ is normal. Hence $F=H^{C^2}$ is Galois over
$\mathbf Q$, and its Galois group is
$(C/C^2)\times\langle c\rangle$, since inversion is trivial on a group
of exponent two. In particular $F/\mathbf Q$ is abelian.

Every character $\chi$ of $C/C^2$ extends to two one-dimensional
characters $\chi_+,\chi_-$ of $\operatorname{Gal}(F/\mathbf Q)$, by
assigning $c$ the values $1,-1$. Directly on the two cosets of
$\operatorname{Gal}(F/K)$,
\[
       \operatorname{Ind}_{\operatorname{Gal}(F/K)}^{
                       \operatorname{Gal}(F/\mathbf Q)}\chi
                         =\chi_+\oplus\chi_-.
\]
Artin induction of L-functions and abelian class field theory over
$\mathbf Q$ therefore express the class-character Hecke function
$L_K(s,\chi)$ as a product of the two associated Dirichlet L-functions.
This includes $\chi=1$, for which the product is $\zeta_K$.
The input $A_D$ gives precisely the zero-location premises for all these
functions.

For the explicit generator bound, fix any nontrivial character $\chi$ of
$C/C^2$. Its kernel in $C$ defines a quadratic unramified extension
$F_\chi/K$ inside $H$. The preceding conjugation argument shows that
$F_\chi/\mathbf Q$ is biquadratic. Its absolute discriminant is
$|D_K|^2$ by the relative-discriminant tower formula and unramifiedness.
Apply Bach--Sorenson Theorem 5.1 to the nonidentity element of
$\operatorname{Gal}(F_\chi/K)$, regarded inside
$\operatorname{Gal}(F_\chi/\mathbf Q)$. Its zeta hypothesis follows from
$A_D$, since this degree-four field is abelian over $\mathbf Q$.
The resulting unramified rational prime satisfies
\[
            p\le(4\log(|D_K|^2)+2.5\cdot4+5)^2
                    =(8\log|D_K|+15)^2.
\]
Its Frobenius is trivial on $K$, so it splits there, and either prime
ideal above it has $\chi$-value $-1$. If the prime ideals in the displayed
range generated a proper subgroup of $C/C^2$, a nontrivial character of
that elementary two-group would be trivial on the subgroup. The prime just
constructed contradicts this. This proves the generator assertion without
paying for the degree of the entire genus field.

The character-by-character proof of JMV Theorem 1.1, author p.~9,
uses only the L-function of the individual character being estimated.
The spectrum of the quotient graph consists exactly of characters
factoring through $C/C^2$ (Remark 1.2(a), p.~4). Thus all its eigenvalue
bounds hold under the premises just established. Lemma 2.1, pp.~7--8,
and the proof of Corollary 1.3, p.~10, give the claimed mixing bound
with $|G|=|C/C^2|$. If $G$ is trivial, the statement is immediate.
\hfill$\square$

This is the genus-level part of the arithmetic Picard action. It sharpens
the input split in CX-ALG-013: the full Hilbert class field may be
nonabelian over $\mathbf Q$, while its genus quotient is abelian and
already covered by $A_D$. It does not give the missing estimates for the
other class characters or a general deterministic class-group algorithm.

\paragraph{CX-GEO-001: existential constructible feasibility.}
\emph{Statement.} Conditional on $M$, deciding whether a quantifier-free
Boolean circuit of rational polynomial equalities and inequations has a
complex solution belongs to $\mathrm{AM}$. Its emptiness problem belongs to
$\mathrm{coAM}\subseteq\Sigma_2^{\mathrm P}$. Polynomial expressions are
given by arithmetic circuits and all original variables are explicitly listed.

\emph{Proof.} Represent each atomic equality $f(x)=0$ by a new Boolean
variable $b$ and a new variable $u$, with equations
\[
 b(b-1)=0,\qquad bf(x)=0,\qquad (1-b)(u f(x)-1)=0.
\]
If $b=1$, these require $f(x)=0$. If $b=0$, they require $f(x)\ne0$
and an inverse $u=f(x)^{-1}$. Conversely any assignment to $x$ extends
with $b$ equal to the truth value of $f(x)=0$; when $b=1$, any $u$ works.
An inequation is the negation of its equality atom.
For each Boolean gate, introduce its output variable $t$ and use
\[
 t-uv=0\quad(\mathrm{AND}),\qquad
 t-u-v+uv=0\quad(\mathrm{OR}),\qquad
 t+u-1=0\quad(\mathrm{NOT}),
\]
where $u,v$ here denote that gate's Boolean inputs. Impose $t(t-1)=0$
for each output if desired, and impose that the final output equals $1$.
Induction through the acyclic Boolean circuit proves that the resulting
polynomial system is feasible exactly when the original formula is true
at some $x$. The construction shares each original arithmetic circuit;
it does not expand a Boolean formula into disjunctive normal form.
Its size and rational-constant bit lengths increase polynomially.
Apply CX-HN-001 to the resulting system and take complements for emptiness.
The reduction concerns equations and inequations over $\mathbf C$, not
ordered-field inequalities. Geometrically it realizes the given constructible
locus as the image of an explicitly presented affine scheme and preserves
nonemptiness.\hfill$\square$

\paragraph{CX-GEO-002: radical membership.}
\emph{Statement.} Conditional on $M$, for
$I=(f_1,\ldots,f_s)\subseteq\mathbf Q[x_1,\ldots,x_n]$ and $g$ given
by rational arithmetic circuits, radical membership $g\in\sqrt I$ is in
$\mathrm{coAM}$ and nonmembership is in $\mathrm{AM}$.

\emph{Proof.} By the Hilbert Nullstellensatz over $\mathbf C$,
$g\notin\sqrt{I\mathbf C[x]}$ exactly when there is a point $x$ with
$f_i(x)=0$ for every $i$ and $g(x)\ne0$. This is equivalent to feasibility of
\[
               f_i(x)=0\quad(1\leq i\leq s),\qquad y g(x)-1=0.
\]
For completeness, the rational and complex membership tests agree. If
$g^m\in I\mathbf C[x]$, a representation by a finite list of polynomial
multiples gives a finite linear system with rational coefficients for those
multipliers' coefficients. A rational linear system having a complex solution
has a rational solution, by row reduction. Thus $g^m\in I$. The converse
follows by extension of scalars. The new system has polynomial size, so
CX-HN-001 gives the asserted classes. In scheme language the question is
whether the principal open $D(g)$ in $\operatorname{Spec}(\mathbf Q[x]/I)$
has a geometric point. Nilpotents matter only through the radical in this
particular problem.\hfill$\square$

\paragraph{CX-GEO-003: reduced containment and equality.}
\emph{Statement.} Conditional on $M$, inclusion and equality of affine
algebraic sets specified by rational arithmetic-circuit equations are in
$\mathrm{coAM}$. This statement does not assert the same bound for
scheme-theoretic ideal inclusion.

\emph{Proof.} Put $X=V(f_1,\ldots,f_s)$ and
$Y=V(g_1,\ldots,g_t)$. Then $X\nsubseteq Y$ exactly when
\[
        f_i(x)=0\quad(1\leq i\leq s),\qquad
        \sum_{j=1}^t u_jg_j(x)=1                              \tag{CX.1}
\]
is feasible. Indeed, the last equation is impossible if all $g_j(x)$ vanish;
if one does not vanish, its inverse supplies a solution with all other
$u_j=0$. For $t=0$ the sum is zero, making (CX.1) impossible, correctly
expressing $Y=\mathbf A^n$. For an empty $f$-list there is no restriction
on $x$, also correctly. Thus noninclusion is in $\mathrm{AM}$ and
inclusion is in $\mathrm{coAM}$.

Equality can be handled by the same single-reduction method. Introduce
$b$ with $b(b-1)=0$. Impose
\[
\begin{aligned}
 b f_i(x)&=0,& (1-b)g_j(x)&=0,\\
 b\left(\sum_j u_jg_j(x)-1\right)&=0,&
 (1-b)\left(\sum_i v_if_i(x)-1\right)&=0.
\end{aligned}
\]
For $b=1$ this is a point of $X\setminus Y$; for $b=0$ it is a point
of $Y\setminus X$. Hence inequality is one HN instance of polynomial size.
Finally, the ideals $(x)$ and $(x^2)$ have the same complex zero set but
are unequal. They show exactly why this reduction tests underlying reduced
sets rather than the structure sheaves or ideal inclusions.
For GCT this test applies when defining equations of the orbit closures
are actually supplied; producing such equations remains a separate gate.
\hfill$\square$

\paragraph{CX-GEO-004: a finite Frobenius test with prescribed excluded primes.}
\emph{Statement.} Let $L/\mathbf Q$ be finite Galois, let $G$ be its Galois
group, and let $E_i=L^{H_i}$ for subgroups $H_i\leq G$ ($i=1,2$).
Let $\chi_i$ be the permutation character of $G/H_i$. Set $n=[L:\mathbf Q]$
and $D=|D_L|$. Suppose the right halves of the critical strips for
$\zeta_L$ and $\zeta_{\mathbf Q}$ are zero-free, as supplied by $M$, or by
$A_D$ if $L/\mathbf Q$ is abelian. If $\chi_1\ne\chi_2$, there is a prime
unramified in $L$ such that
\[
       p\leq B_L=(4\log D+\tfrac52 n+5)^2                  \tag{CX.2}
\]
and the numbers of degree-one prime ideals of $E_1$ and $E_2$ above $p$
differ. Equivalently the numbers of $\mathbf F_p$-points of the two integer
models differ.

For any nonempty finite set $S$ of rational primes put
$d=\prod_{p\in S}p$, $F=\mathbf Q(\sqrt d)$, and $L'=LF$.
The same assertion holds with $p\notin S$ and bound $B_{L'}$, under the
corresponding hypotheses for $\zeta_{L'}$ and $\zeta_{\mathbf Q}$.
If $n'=[L':\mathbf Q]$ and $D'=|D_{L'}|$, an explicit upper estimate is
\[
 n'\leq2n,\qquad \log D'\leq2\log D+n\log|D_F|
                     \leq2\log D+n\log(4d),               \tag{CX.3}
\]
so the entirely explicit larger bound
\[
        \bigl(8\log D+4n\log(4d)+5n+5\bigr)^2            \tag{CX.4}
\]
also suffices. If $S=\varnothing$, use (CX.2) directly.

\emph{Proof.} Choose $\sigma\in G$ such that
$\chi_1(\sigma)\ne\chi_2(\sigma)$. Bach--Sorenson,
\emph{Explicit Bounds for Primes in Residue Classes} (1996), Theorem 5.1,
p.~1729, gives an unramified prime whose Frobenius conjugacy class is that
of $\sigma$, with bound (CX.2). Their Section 2, p.~1720, specifies the
needed zero-free right-half assumption for the top-field zeta function and
$\zeta_{\mathbf Q}$; no Artin-holomorphy conjecture is required.
For an unramified prime the Frobenius orbits on $G/H_i$ correspond to the
prime ideals above $p$, and the orbit lengths are their residue degrees.
Its fixed points therefore count precisely the degree-one primes.
This proves the first assertion.

For the exclusion refinement, every prime of $S$ ramifies in $F$. For odd
$p\mid d$ this follows from the quadratic discriminant formula; if $2\in S$,
then $d$ is even and $D_F=4d$, so 2 ramifies as well. Consequently every
prime in $S$ ramifies in $L'$, since its ramification index in the subextension
$F/\mathbf Q$ divides its index upstairs. Choose any lift of $\sigma$ to
$\operatorname{Gal}(L'/\mathbf Q)$ and apply the same theorem there. The
prime obtained is unramified in $L'$, hence avoids $S$, and its restriction
to $L$ has the required character difference.

We justify (CX.3) without an unstated discriminant estimate. If $F\subseteq L$,
then $L'=L$ and the inequalities are immediate. Otherwise $L$ and $F$ are
linearly disjoint because $F$ is quadratic. If $\omega$ is an integral basis
generator for $\mathcal O_F$, the order $\mathcal O_L[\omega]\subseteq
\mathcal O_{L'}$ has relative discriminant generated by $D_F$. The relative
field discriminant divides the order discriminant. The tower formula gives
$D'\leq D^2|D_F|^n$. Also $n'=2n$ and $|D_F|\leq4d$.
Substitution in (CX.2) yields (CX.4). If $L/\mathbf Q$ is abelian, $LF$ is
also abelian, so the stated $A_D$ alternative survives the exclusion step.
The exceptional source case $L=\mathbf Q$ cannot have unequal characters
for the two subfields in this statement, so causes no gap.\hfill$\square$

\emph{Consequences and limits.} Agreement at all the tested primes implies
$\chi_1=\chi_2$. Hence the corresponding characteristic-zero permutation
representations are isomorphic and their Artin $L$-functions, including
ramified factors formed on inertia invariants, agree. In particular
$\zeta_{E_1}=\zeta_{E_2}$. This is a finite trace certificate, not necessarily
an efficient algorithm in a short presentation of $E_i$: a common Galois
closure may have large degree. Nor is equality of permutation characters
an isomorphism theorem for finite covers, as the next result proves.
The Grothendieck framework is SGA 1 V Section 7 and the Frobenius trace
formalism; cote 98 scans 11--12 are the finite-$G$-set archival companion
carried from the CX\_GRH crosswalk.

\paragraph{CX-GEO-005: equal characters do not identify transitive $G$-sets.}
\emph{Statement.} Let $V=\mathbf F_2^3$, $G=\operatorname{GL}(V)$,
$X=V\setminus\{0\}$ and $Y=V^*\setminus\{0\}$ with the natural and dual
actions. Then $X$ and $Y$ are transitive seven-element $G$-sets with identical
permutation characters but are not isomorphic as $G$-sets.

\emph{Proof.} For $g\in G$, the number of fixed nonzero vectors is
$2^{\dim\ker(g-I)}-1$. The corresponding number of fixed nonzero covectors
is $2^{\dim\ker((g^{-1})^{\mathsf T}-I)}-1$. The matrices
$(g^{-1})^{\mathsf T}-I$ and $g^{\mathsf T}-I$ have the same rank, since
they differ by multiplication by an invertible matrix and sign. Transpose
preserves rank, so both fixed-point counts are equal.

Suppose the two $G$-sets were isomorphic. The stabilizer $H$ of a nonzero
vector, say $e_1$, would stabilize some nonzero covector and therefore its
kernel, a plane $W\subset V$. Such an invariant plane does not exist.
If $e_1\in W$, its image $W/\langle e_1\rangle$ is a line in
$V/\langle e_1\rangle$, invariant under all of
$\operatorname{GL}_2(\mathbf F_2)$ because $H$ induces that full group.
No line is invariant under that group. If $e_1\notin W$, projection identifies
$W$ with the two-dimensional quotient, so $W$ is a complement to
$\langle e_1\rangle$. But $H$ contains all transvections
$(a,v)\mapsto(a+u(v),v)$ with $u$ a linear functional on that quotient.
Choosing $u(v)=1$ for a nonzero projected vector in $W$ moves its lift by
$e_1$, which cannot stay in $W$. Both cases contradict invariance.
Thus no $G$-set isomorphism exists.\hfill$\square$

This is a representation-theoretic counterexample, needing no GRH and no
inverse-Galois assertion. For any connected finite \`etale $G$-torsor that
realizes this group, its two associated connected covers exhibit the same
limitation. No particular unverified realization over $\mathbf Q$ is used.
The example distinguishes the Galois category from its linearized
representation category; the latter is what Frobenius traces recognize.

\paragraph{CX-GEO-006: orbit membership with an explicit action.}
\emph{Statement.} Conditional on $M$, suppose an algebraic group $G$ and its
action on an affine space have polynomial-length rational circuit descriptions,
with every inverted denominator explicitly listed. Given rational points
$v,w$, deciding whether $w\in G(\mathbf C)v$ belongs to $\mathrm{AM}$.

\emph{Proof.} Introduce coordinates for $g$, impose the supplied equations
and inequations defining $G$, and impose the coordinate equations
$g\cdot v=w$. For example, in $\operatorname{GL}_n$ an additional variable
$z$ and equation $z\det(g)-1=0$ enforce invertibility; the determinant has
a polynomial-size arithmetic circuit. Inverted denominators of a rational
action are handled by the same inverse-variable equations. The resulting
constructible-feasibility instance has polynomial size and is equivalent to
orbit membership by definition of the action. Apply CX-GEO-001.
Orbit-closure membership is not represented by this system: points in the
boundary need not equal $g\cdot v$ for any $g$. If closure equations are
provided, CX-HN-001 and CX-GEO-003 apply to those equations; without them
the equation-access problem persists.\hfill$\square$

\paragraph{CX-GEO-007: polynomially parameterized orbit closures already admit a PSPACE test.}
\emph{Statement.} Let the ambient coordinates of $V$ be explicitly listed.
Suppose the action on a supplied rational vector $v$ is given by
division-free arithmetic circuits with rational coefficients for all its
coordinates. In either of the following
cases, membership of a rational point $w$ in the complex orbit closure is
in $\mathrm{PSPACE}$, without either GRH input: (a) the group is
$\operatorname{GL}_n$ and the map $g\mapsto g\cdot v$ extends by the supplied
circuits to a polynomial map on all $\operatorname{Mat}_n$; (b) the group is
$\operatorname{SL}_n$ and the supplied action coordinates are polynomial
in its matrix entries. This statement charges the explicit ambient
coordinate count in the input size.

\emph{Proof.} Guo--Saxena--Sinhababu (2019), Theorem 1.3, p.~5, together
with Lemma 4.1 and Theorem 4.2, pp.~15--16, gives a polynomial-space test
for whether zero lies in the Zariski closure of the image of a polynomial
map specified by arithmetic circuits: this is their approximate polynomial
satisfiability problem. In case (a), $\operatorname{GL}_n$ is dense in
$\operatorname{Mat}_n$. For a polynomial map $F$ on the latter, every closed
set containing $F(\operatorname{GL}_n)$ has closed inverse image containing
a dense subset, hence contains $F(\operatorname{Mat}_n)$. Their image
closures are equal. Apply the cited test to $F-w$.

For case (b), construct a polynomial map from an affine space into
$\operatorname{SL}_n$ whose image contains its Gaussian-elimination big
cell. Products of elementary unipotent matrices parameterize arbitrary
lower and upper unitriangular matrices. A diagonal matrix of determinant
one is a product of adjacent diagonal blocks
$\operatorname{diag}(t,t^{-1})$: for diagonal entries $d_1,\ldots,d_n$,
use $t_i=d_1\cdots d_i$ in positions $(i,i+1)$ for $1\leq i<n$.
Such a block belongs to the image of a
six-factor elementary-unipotent word, because, in its $2$ by $2$ block,
\[
 E_{12}(t)E_{21}(-t^{-1})E_{12}(t)
 E_{12}(-1)E_{21}(1)E_{12}(-1)
       =\operatorname{diag}(t,t^{-1}).
\]
Use independent parameters in each of the six positions of this fixed
word. The resulting map is polynomial and always has determinant one;
the displayed specializations show that it contains every such diagonal
block in its image. Place these words between the lower and upper
unitriangular words. Its image contains every $LDU$ decomposition in the
big cell, which is a dense open subset of $\operatorname{SL}_n$.
The number of elementary factors and parameters is $O(n^2)$ and the
matrix product has polynomial-size circuits. Composing with the supplied
action on $v$ gives a polynomial map with the same image closure as the
orbit map, by the same closed-inverse-image argument. Subtract $w$ and
apply the cited PSPACE algorithm. For $n=1$, the special-linear action is
trivial and membership reduces to $w=v$.\hfill$\square$

This gives an equation-free route for this exact orbit-closure encoding.
It does not give an AM bound from $M$ for approximate satisfiability, and
it does not turn an exponentially long list of ambient coefficients into
a polynomial-length input. Those are separate questions from the theorem.

\paragraph{CX-GEO-008: a finite Frobenius generating certificate.}
\emph{Statement.} Let $L/\mathbf Q$ be finite Galois with group $G$ and
assume the zeta hypotheses of CX-GEO-004. For each unramified rational
prime $p\leq B_L$, choose any representative $g_p$ of its Frobenius class.
The elements $g_p$ generate $G$. More strongly, a subgroup meeting every
Frobenius class observed in this range equals $G$. A sublist of at most
$\lfloor\log_2|G|\rfloor$ representatives suffices to generate $G$.
For a prescribed excluded set $S$, replace the bound by $B_{L'}$ from
CX-GEO-004 and use only primes outside $S$. As before, $A_D$ suffices for
abelian $L/\mathbf Q$, and $M$ supplies the general hypothesis.

\emph{Proof.} Let $H<G$ be proper and let $X=G/H$. Double counting gives
\[
 \sum_{g\in G}|\operatorname{Fix}_X(g)|
      =\sum_{x\in X}|\operatorname{Stab}_G(x)|=|G|.
\]
The identity fixes $|X|>1$ points. If every element fixed at least one,
the left side would exceed $|G|$. Hence there is $g$ with no fixed coset.
Since $g$ fixes $aH$ exactly when $a^{-1}ga\in H$, its conjugacy class
misses $H$. The explicit Chebotarev theorem used in CX-GEO-004 produces
that class at an unramified prime $p\leq B_L$. With excluded primes,
choose a lift to $\operatorname{Gal}(L'/\mathbf Q)$ and apply the same
argument at bound $B_{L'}$; every prime of $S$ ramifies in $L'$ and is
therefore avoided. Thus no proper subgroup can meet every observed class,
and in particular no proper subgroup contains all the chosen $g_p$.

Starting from the trivial subgroup, choose a listed $g_p$ outside the
subgroup generated so far until it equals $G$. Each proper enlargement
has index at least two, so after $k$ selections the generated subgroup has
order at least $2^k$. Consequently $k\leq\lfloor\log_2|G|\rfloor$.
For $G=1$ the empty list suffices.\hfill$\square$

In SGA 1 V Section 7 this is a finite quotient of the arithmetic
fundamental group, with a finite list of Frobenius conjugacy classes
certifying that no proper subgroup contains the image. The bound is
relative to the specified extension and its discriminant; it is not a
uniform presentation of the entire profinite fundamental group. The
archival companion is the finite-$G$-set discussion of cote 98 scans
11--12 recorded in the preceding CX\_GRH crosswalk, and the work target
is ORDER item 30's effective finite Galois-image question.

\paragraph{CX-GEO-009: bounded Laurent degenerations give an AM promise.}
\emph{Statement.} Let $F:\mathbf A^n\to\mathbf A^m$ and
$w\in\mathbf Q^m$ be specified by division-free rational arithmetic
circuits of total size $s$. Supply in unary an upper bound $D\ge1$ for
their syntactic degrees and an integer $B\ge0$.
Conditional on $M$, it is in $\mathrm{AM}$ to decide whether there are
Laurent polynomials
\[
 \beta_i(\varepsilon)=\sum_{j=-B}^{B}c_{i,j}\varepsilon^j,
 \qquad c_{i,j}\in\mathbf C,
\]
such that every coordinate of $F(\beta)-w$ belongs to
$\varepsilon\mathbf C[\varepsilon]$.
Consequently image-closure membership is in $\mathrm{AM}$ on the promise
that every positive instance admits such an arc with $D$ and $B$ bounded
polynomially in the original encoding length. This includes the polynomial
orbit parameterizations of CX-GEO-007 with the same arc promise.

\emph{Proof.} Use the coefficients $c_{i,j}$ as unknowns. At each circuit
gate $g$ of syntactic degree $d_g$, maintain arithmetic circuits for the
coefficients of its substituted Laurent polynomial in the range
$[-Bd_g,Bd_g]$. Input-variable gates have their prescribed coefficient
unknowns; a rational constant has only a constant coefficient. At an
addition or subtraction gate, add or subtract coefficient circuits.
At a multiplication gate, take the usual coefficient convolutions.
Syntactic degrees add at multiplication, so no exponent outside the stated
range can occur. Discard gates unused by the outputs; then each remaining
degree is at most $D$.

There are at most $2BD+1$ coefficient positions per gate, and each product
uses at most $(2BD+1)^2$ pairs. Sharing all earlier coefficient circuits
therefore gives a polynomial-size arithmetic-circuit construction in
$s,n,m,B,D$. For every output impose that all coefficients of exponents
$-BD,\ldots,0$ in $F(\beta)-w$ vanish. These are polynomial equations in
the $n(2B+1)$ unknown coefficients. The resulting HN instance has polynomial
encoding size and is feasible exactly when the displayed bounded Laurent
arc exists. Koiran's theorem gives the claimed AM bound under $M$.
The construction also covers $B=0$, when it is ordinary exact feasibility.

Any such arc puts $w$ in the image closure by GSS Lemma 4.1 and Theorem
4.2, pp.~15--16. Conversely the promise provides a bounded arc for every
point of that closure under consideration, proving the promise corollary.
All oracle queries concern this explicit coefficient system; its bounded
solution field is the exact number-field zeta input. An $A_D$ alternative
requires the bounded abelian-witness promise for this coefficient system,
not merely for the original rational points $v,w$.
\hfill$\square$

The generic GSS bound on Laurent precision can be exponential; writing
that bound in binary would not justify the polynomial-size coefficient
construction. This result therefore isolates a concrete smaller gate for
the promised families: a polynomially bounded degeneration, rather than
access to all equations of the orbit closure. It does not assert that
arbitrary GCT orbit closures meet that promise. The geometric objects are
finite Laurent jets of a parameterized degeneration; the group-action
framework is SGA 3 I 4.6.3.1 and 4.7.1--4.7.2, with cote 28 scan 3 as the
archival companion inherited from CX\_GRH. The work target is ORDER item
28's Koiran/GCT route.

\subsection{Gates retained rather than silently discharged}

\paragraph{CX-LIM-001: general derandomization and GCT.}
The exact remaining premises of the cited routes are uniform hitting sets
or independent circuit-hardness hypotheses. GCT V Theorem 4.7, p.~33,
assumes a specific permanent border-hardness bound. Its Lemma 4.8, on the
same page, extracts a hitting set from a sufficiently explicit separating
system. Neither construction is furnished by a statement about zeros of
Dirichlet or Hecke $L$-functions. This records the boundary of these proofs;
it is not a proof that no logical implication from GRH to general
derandomization could ever exist. No $\mathrm{BPP}=\mathrm P$, integer
factorization in $\mathrm P$, or $\mathrm{VP}\ne\mathrm{VNP}$ is claimed.

\paragraph{CX-LIM-002: positivity and purity.}
The earlier ALPOGE\_GCT P-F7 and P-F10 and PLETH PLE-F5/F7 already show
why pure geometry, nonnegative Hilbert data, efficient counting and actual
orbit-closure multiplicities are different inputs. Their established
counterexamples and corrected positivity-index implications are carried,
not rebranded. The new tests CX-GEO-002--CX-GEO-003 can verify specified
reduced-locus properties once equations are given. They do not construct a
uniform positive multiplicity formula or a strict asymptotic quotient-rank
inequality. Those are the surviving GCT gates.

\paragraph{CX-LIM-003: the separate counting premise.}
The attached CX\_GRH return also assumed an independent theorem denoted
$E$ and drew ETH/counting consequences in CX-E01--CX-E05. Tonight's two
analytic inputs $A_D$ and $M$ do not contain $E$. Those rows therefore do
not become consequences of these two inputs by being present in the same
archive. This is an input classification, not a new assertion about ETH.

\subsection{A smaller field-construction gate for deterministic factorization}

\paragraph{CX-ALG-023: prime-degree witnesses and maximal coefficient units.}
\emph{Statement.} Let $p,r$ be distinct primes, put
$s=\operatorname{ord}_r(p)$, $R=r^{v_r(p^s-1)}$, and
\[
 D_r=\mathbf F_p[Z]/(1+Z+\cdots+Z^{r-1}),\qquad \zeta=Z\bmod\Phi_r.
\]
Write $\rho_a(\zeta)=\zeta^a$ for $a\in\mathbf F_r^\times$, and
$\omega(a)$ for its multiplicative Teichm\"uller lift to the relevant
$r$-power modulus. An explicitly supplied field $F=\mathbf F_{p^r}$
permits, unconditionally in $\operatorname{poly}(r\log p)$ time,
construction of $\eta\in D_r^\times$ satisfying
\[
          \operatorname{ord}(\eta)=R,\qquad
          \rho_a(\eta)=\eta^{\omega(a)}\quad(a\in\mathbf F_r^\times).
\]
Conversely, such an $\eta$ permits construction of an explicit
$\mathbf F_{p^r}$ in the same polynomial bound. In particular, a supplied
$\mathbf F_{p^m}$ already supplies this coefficient unit whenever
$r\mid m$, regardless of the value of $v_r(p^s-1)$.

The equivalence is classical: Lenstra, \emph{Finding isomorphisms between
finite fields} (1991), Theorems (5.1)--(5.2), pp.~338--340, with its
normal-basis algorithm, Theorem (3.1), pp.~334--336. The following full
proof makes the exact constructive input and the new application to the
remaining CX\_GRH gate explicit; no priority for the equivalence is claimed.

\emph{Proof.} Compute a normal-basis element $\alpha$ for
$F/\mathbf F_p$ by the cited unconditional linear-algebra algorithm.
Form $C=F\otimes_{\mathbf F_p}D_r$, let $\sigma$ act by $p^s$-power
Frobenius on $F$ and trivially on $D_r$, and set
\[
             \beta=\sum_{i=0}^{r-1}\zeta^i\sigma^i(\alpha).
\]
The factors of $\Phi_r$ over $\mathbf F_p$ all have degree $s$.
Since $s\mid r-1$, the fields $F$ and $\mathbf F_{p^s}$ are linearly
disjoint. Thus each component of $C$ is a field, and the elements
$\sigma^i(\alpha)$ remain linearly independent over its coefficient
field $\mathbf F_{p^s}$. The displayed sum is consequently nonzero in
every component: $\beta$ is a unit. This use of a normal basis prevents
the vanishing defect in CX-ALG-016.

On $C$, $\sigma$ equals $p^s$-power Frobenius. Shifting the sum gives
$\sigma(\beta)=\zeta^{-1}\beta$. With $u=(p^s-1)/R$ and
$\delta=\beta^u$ it follows that
\[
                  \delta^R=\zeta^{-1},\qquad \delta^{rR}=1.
\]
Compute the lifts $\omega(a)$ modulo $rR$ and their inverses, and put
\[
          y=\prod_{a\in\mathbf F_r^\times}
                        \rho_a(\delta)^{\omega(a)^{-1}}.
\]
Reindexing the product by $a\mapsto ba$ proves
$\rho_b(y)=y^{\omega(b)}$. Moreover,
\[
          y^R=\prod_a\zeta^{-a\omega(a)^{-1}}
                   =\zeta^{-(r-1)}=\zeta.
\]
Hence $y$ has order $rR$. The element $\eta=y^r$ has order $R$,
and $\eta^{p^s-1}=1$. Therefore $\sigma(\eta)=\eta$.
Since $\gcd(s,r)=1$, the fixed field of $\sigma$ on $F$ is
$\mathbf F_p$, so $C^\sigma=D_r$. This proves $\eta\in D_r$ and
the required equivariance.

For the converse, equivariance and transitivity of the $\rho_a$ on the
components of $D_r$ imply that the image of $\eta$ in every component
$\mathbf F_{p^s}$ has order $R$. It is an $r$th nonresidue there.
The polynomial $Y^r-\eta$ is irreducible in each component: in a field
containing $\boldsymbol\mu_r$, its splitting group is a subgroup of
the prime-order group $\boldsymbol\mu_r$, and reducibility would force
a root and hence make $\eta$ an $r$th power. Form
\[
                       C'=D_r[Y]/(Y^r-\eta).
\]
Extend $\rho_a$ by $\rho_a(Y)=Y^{\omega(a)}$, taking lifts modulo
$rR$. The relations are preserved, and uniqueness and multiplicativity
of the lifts give an action of $\mathbf F_r^\times$.
This group acts transitively on the field components of $C'$, so its
fixed algebra is a field. It has degree $r$ over $\mathbf F_p$:
after extending scalars to an algebraic closure, $C'$ has $r(r-1)$
points, the action on them is free because its action on $\zeta$ is
free, and there are exactly $r$ orbits. Invariants commute with this
scalar extension, being kernels of explicit linear maps. The fixed
field is therefore an explicit $\mathbf F_{p^r}$.

All algebras in this proof have dimension at most $r(r-1)$.
The integer $R$ has $O(r\log p)$ bits; the lifts can be calculated by
Hensel lifting the simple roots of $X^{r-1}-1$ modulo $r$, or by
modular exponentiation. Fixed algebras and their multiplication tables
are obtained by linear algebra. These operations give the stated bound.
If $r\mid m$, the subfield $F\subseteq\mathbf F_{p^m}$ is the kernel
of $z\mapsto z^{p^r}-z$, computable in time polynomial in the supplied
field size in bits. Applying the first construction finishes the proof.
\hfill$\square$

Geometrically $D_r$ is a finite \`etale coefficient algebra; adjoining
an $r$th root is a Kummer cover, and taking the displayed invariants
is descent of a finite \`etale algebra. The companions are SGA~1,
Expos\'e~V, \S7 and SGA~4, Expos\'e~IX, 3.2, as in CX\_GRH
CX-N13--N15. The new use is the additional $r\mid m$ branch of
CX\_GRH CX-U02, not a claim that its older proof was absent or wrong.

\paragraph{CX-ALG-024: a pure prime-degree zeta function supplies the auxiliary prime.}
\emph{Statement.} Let $r$ be an odd prime and $p\ne r$ any prime, and write
\[
 F_r=\mathbf Q(p^{1/r}),\quad K_r=\mathbf Q(\zeta_r),\quad
 E_r=\mathbf Q(\zeta_r,p^{1/r}).
\]
Under $A_D$, RH for $\zeta_{F_r}$ is equivalent to RH for
$\zeta_{E_r}$. It implies a rational prime $q\ne p,r$ with
\[
 q\equiv1\pmod r,\qquad p^{(q-1)/r}\not\equiv1\pmod q,
 \qquad q\le B_{p,r},
\]
where the explicit sufficient bound is
\[
 B_{p,r}=
 \left(4r(2r-3)\log r+4(r-1)^2\log p
                       +\tfrac52r(r-1)+5\right)^2.
\]
Consequently an explicit $\mathbf F_{p^r}$ is constructible in
$\operatorname{poly}(r\log p)$ time from this analytic input.
Under the task convention $M$ supplies RH for the degree-$r$ field
$F_r$; its role here is conditional.

\emph{Proof.} The discriminant of the order $\mathbf Z[\zeta_r]$ is
$\pm r^{r-2}$, as follows by taking the product of the derivatives
of $\Phi_r$ at its roots. Thus $p$ is unramified in $K_r$.
The polynomial $X^r-p$ is Eisenstein at every prime of $K_r$ above $p$,
so $[E_r:K_r]=r$. The group of the splitting field $E_r/\mathbf Q$
is the affine group
\[
             G=\mathbf F_r\rtimes\mathbf F_r^\times,
       \qquad (a,b):t\longmapsto bt+a
\]
acting on the $r$ roots of $X^r-p$.
Let $P$ be this permutation character, and let $Q$ be the regular
character of the quotient $G/\mathbf F_r$, inflated to $G$.
There is an identity of characters
\[
                    \operatorname{reg}_G=(r-1)(P-1)+Q.
\]
At the identity both sides equal $r(r-1)$; at a nonzero translation,
$P=0$ and $Q=r-1$; and at an affine transformation with $b\ne1$,
$P=1$ and $Q=0$. These cases exhaust $G$, proving the identity.
Artin formalism, including its ramified inertia-invariant local factors,
now gives
\[
          \zeta_{E_r}(s)=\zeta_{K_r}(s)
                    \left(\frac{\zeta_{F_r}(s)}{\zeta_{\mathbf Q}(s)}
                    \right)^{r-1}.
\]
The abelian field $K_r$ and $\mathbf Q$ have their zero-location
assertions under $A_D$. Thus in the open critical strip off the middle
line the displayed identity proves the asserted equivalence. This uses
no general Artin-holomorphy conjecture.

For a discriminant bound, the relative order
$\mathcal O_{K_r}[p^{1/r}]$ has discriminant generated, up to sign, by
$r^r p^{r-1}$. The field discriminant divides this order discriminant.
The discriminant tower formula therefore yields
\[
      |D_{E_r}|\le
                  r^{\,r(2r-3)}p^{(r-1)^2},
      \qquad [E_r:\mathbf Q]=r(r-1).
\]
Apply Bach--Sorenson (1996), Theorem~5.1, p.~1729, to the conjugacy
class of the automorphism fixing $K_r$ and sending $p^{1/r}$ to
$\zeta_r p^{1/r}$. Its conjugates send that root to
$\zeta_r^a p^{1/r}$ for $a\in\mathbf F_r^\times$.
The theorem's precise zeta input (\S2, p.~1720) is supplied by the
preceding factorization and $A_D$. Substitution of the degree and
discriminant bounds into $(4\log|D_{E_r}|+2.5[E_r:\mathbf Q]+5)^2$
gives $B_{p,r}$.

The resulting unramified prime splits in $K_r$, hence $q\equiv1\pmod r$.
At a prime of $E_r$ over $q$, its Frobenius relation is
\[
       (p^{1/r})^q\equiv\zeta_r^a p^{1/r},\qquad
              p^{(q-1)/r}\equiv\zeta_r^a\ne1.
\]
Unramifiedness excludes $p$ and $r$, and the root remains nonzero.
For prime $r$, the last inequality is equivalent to
$r\nmid(q-1)/\operatorname{ord}_q(p)$.
The image of $p$ consequently generates the order-$r$ quotient of
$(\mathbf Z/q\mathbf Z)^\times$. Equivalently, $p$ is inert in the
unique degree-$r$ subfield of $\mathbf Q(\zeta_q)$.

This is exactly the algebraic input to Adleman--Lenstra (1986),
Algorithm~A, steps (2)--(3), pp.~350--351, and Lemma~6, p.~352.
That algorithm constructs an irreducible degree-$r$ polynomial by
Gaussian periods using $\operatorname{poly}(q\log p)$ operations;
the primality and power-congruence tests for $q$ use the same bound.
Enumerating primes until the first suitable $q$ halts by $B_{p,r}$,
which is $O(r^4\log^2(rp))$. No additional Hecke or radical oracle
enters this construction. This proves the polynomial bound.
\hfill$\square$

The geometric interpretation is a relation among the permutation
objects of an affine finite Galois cover: the degree-$r$ cover and its
cyclotomic quotient determine the required zeta input of the full cover.
It belongs to the Artin-motive and finite-monodromy framework of SGA~1
V \S7 and to the representation-ring operations of SGA~6 V.
It sharpens the auxiliary-prime input of CX\_GRH CX-N11--N15;
the earlier sufficient family remains valid.

\paragraph{CX-ALG-025: a complete factoring composition with a smaller analytic family.}
\emph{Statement.} Let $f\in\mathbf F_{p^m}[X]$ have degree $d\ge2$,
with the base field explicitly supplied. If $p\le d$, deterministic
factorization is unconditional in $\operatorname{poly}(d\log(p^m))$
time. Suppose $p>d$, and define
\[
 \mathcal B(p,d,m)=\{r\le d:\ r\text{ prime},\ r\nmid(p-1)m,
                 \ v_r(p^{\operatorname{ord}_r(p)}-1)>1\}.
\]
Under $A_D$ and RH for the finite family
\[
             \zeta_{\mathbf Q(p^{1/r})},\qquad r\in\mathcal B(p,d,m),
\]
one can factor $f$ deterministically in
$(d^{\log d}\log(p^m))^{O(1)}$ time. Only the primes from this family
actually encountered need their zeta hypothesis. If $\mathcal B$ is
empty, the result is $A_D$-only. In general $M$ supplies the remaining
input, conditionally. Every number field in the displayed new family
has prime degree at most $d$. If every prime $r\le d$ either divides
$m$ or satisfies $v_r(p^{\operatorname{ord}_r(p)}-1)=1$, the same
quasipolynomial bound is unconditional.

\emph{Proof.} First reduce to a square-free polynomial by gcds and
inverse Frobenius; repeated-factor multiplicities are retained.
For a square-free factor $h$, put
$A=\mathbf F_{p^m}[X]/(h)$ and compute the $\mathbf F_p$-linear fixed
algebra
\[
                         B=\{x\in A:x^p=x\}.
\]
If $h$ has $t$ irreducible factors, its Chinese-remainder decomposition
shows that $B\simeq\mathbf F_p^t$. This identity is used to prove
correctness, not as supplied factorization data. A basis and a
multiplication table for $B$ are found by linear algebra in dimension
at most $m\deg h$. If $t=1$, $h$ is irreducible and no split is needed.

If $p\le d$, take a nonscalar basis element $b\in B$ and test
$b-a$ for $a\in\mathbf F_p$. One of these is a nonzero zero divisor:
in the unknown split coordinates $b$ has at least two different values.
Thus at most $p$ tests of multiplication matrices supply a proper
idempotent, and recursive splitting has polynomial cost in this case.

Assume henceforth $p>d$. Apply the unconditional IKRS algorithm,
published Theorem~5.8, author PDF pp.~28--29, to the explicit algebra
$B$ and its scalar subalgebra $\mathbf F_p$. It either returns a
nonzero zero divisor or a semiregular automorphism $\sigma$ of order
$t$. Choose a prime $r\mid t$ and put $\tau=\sigma^{t/r}$;
this is a semiregular automorphism of prime order $r\le d<p$.
The group in question acts freely on the split components; in this
setting its powers have the required semiregularity directly.

There are four ways to obtain the coefficient unit $\eta$ of
CX-ALG-023 for this $r$.
If $r\mid m$, use the degree-$r$ subfield of the supplied base field.
If $r\mid p-1$, CX-ALG-004 finds an $r$th nonresidue $a\in\mathbf F_p$
under $A_D$; the prime-degree Kummer polynomial $X^r-a$ is irreducible,
so CX-ALG-023 applies. If $v_r(p^{\operatorname{ord}_r(p)}-1)=1$,
the formal root $\zeta\in D_r$ itself already has the required order
and equivariance. Otherwise $r\in\mathcal B$, and CX-ALG-024 followed
by CX-ALG-023 supplies $\eta$ from its one additional zeta hypothesis.
All these operations take polynomial time in $d\log(p^m)$.

For completeness, here is why that coefficient unit closes the
automorphism step. IKRS Proposition~4.6, published author PDF
pp.~18--19, constructs a Teichm\"uller element
$x\in(B\otimes D_r)^\times$ with
$\tau(x)=\zeta x$. Every component of $B\otimes D_r$ is
$\mathbf F_{p^s}$, where $s=\operatorname{ord}_r(p)$, so its
$r$-primary unit group has exponent $R$.
If $\langle x,\eta\rangle$ were cyclic, the element $\eta$ of
order $R$ would generate it. Then $x$ would be a power of the
coefficient element $\eta$ and hence fixed by $\tau$, contradicting
$\tau(x)=\zeta x$ and invertibility of $x$.
IKRS Lemma~4.2, published author PDF p.~16, therefore supplies a
nonzero zero divisor in $B$ in polynomial time. Its proof also gives
the descended idempotent: the annihilator of the difference of two
powers of these Teichm\"uller elements is invariant under every
$\rho_a$, so its identity lies in $B$.

From any nonzero zero divisor $z\in B$, compute the identity $e$ of
the nonzero proper ideal $Bz$ by linear algebra; it is a proper
idempotent. Such an $e\in B\subseteq A$ has values zero or one in
the factor fields. Hence $\gcd(h,e)$, with $e$ represented by its
polynomial residue, is a proper factor of $h$. Recurse on the two
factors. There are at most $2d-1$ nodes; at each node the IKRS cost
is bounded by $(d^{\log d}\log(p^m))^{O(1)}$, and all the added
operations are polynomial. If each $r\le d$ is covered by the
supplied-field or valuation-one branch, choose those branches first;
then no analytic input is used. This proves both stated bounds.
\hfill$\square$

This is a new composition and finite-input refinement using classical
algorithms. It gives a strictly smaller sufficient analytic family than
the earlier CX\_GRH CX-N15 construction through degrees $sr$ and
auxiliary primes dividing $r-1$. It also strictly enlarges the earlier
$A_D$ cases by exploiting $r\mid m$. The unrestricted $A_D$ gate is
now to construct a degree-$r$ field, equivalently its maximal coefficient
unit, only when
\[
       p>d,\qquad r\nmid p(p-1)m,\qquad
             v_r(p^{\operatorname{ord}_r(p)}-1)>1.
\]
This is a sufficient construction problem, not a claim that those RH
assumptions are necessary. The geometric task is still computation of
primitive idempotents of a finite \`etale algebra, SGA~1 V \S7;
the Kummer descent is SGA~4 IX 3.2. It changes CX\_GRH CX-U02 and
ORDER~28 at their actual coefficient-field gate.

\paragraph{CX-ALG-026: exact field extensions have the same prime-degree gate.}
\emph{Statement.} Given $\mathbf F_{p^m}$ and a positive integer $n$
charged in unary, an explicit extension of degree $n$ is constructible
in $\operatorname{poly}(n\log(p^m))$ time under $A_D$ and RH for
$\zeta_{\mathbf Q(p^{1/r})}$ only for prime divisors $r\mid n$ such that
\[
       r\ne p,\qquad r\nmid(p-1)m,\qquad
               v_r(p^{\operatorname{ord}_r(p)}-1)>1.
\]
If every prime divisor $r$ of $n$ either divides $pm$ or satisfies
$r\ne p$ and $v_r(p^{\operatorname{ord}_r(p)}-1)=1$, no analytic
hypothesis is needed. In particular, this holds if every such $r$ divides
$pm$. The output size, and hence the unary or equivalent dense-field
convention on $n$, is essential.

\emph{Proof.} Use Lenstra (1991), Theorem~(9.1), pp.~344--346.
It requires an irreducible degree-$r$ polynomial over the supplied base
field only for primes $r\mid n$ not dividing $m$.
For $r=p$, the Artin--Schreier construction of its Theorem~(7.1),
p.~342, is unconditional and polynomial in $p\log(p^m)$; here
$p\le n$ is charged to the output size. For $r\ne p$, obtain an
explicit $\mathbf F_{p^r}$ by the appropriate branch of
CX-ALG-023--CX-ALG-024. In the valuation-one branch use the converse
of CX-ALG-023 with $\eta=\zeta$; in the $r\mid p-1$ branch use
the prime-field nonresidue under $A_D$.
Convert this explicit field to an irreducible polynomial by Lenstra's
Theorem~(1.1), p.~330. Since $r\nmid m$, that polynomial stays
irreducible over $\mathbf F_{p^m}$: the two finite extensions of
$\mathbf F_p$ have coprime degrees. These are all inputs to
Theorem~(9.1), which supplies the claimed extension and polynomial
bound. Under the stated unconditional promise, use the supplied-field,
Artin--Schreier, or valuation-one branch before considering the
$r\mid p-1$ nonresidue branch. Thus no analytic branch is called.\hfill$\square$

The output is a connected degree-$n$ finite \`etale cover of the
supplied finite-field point. This is the explicit construction counterpart
to the isomorphism/uniqueness theorem for finite fields, not a replacement
for general polynomial factorization. Its run locators are CX\_GRH
CX-N15/CX-U02 and the construction interface in ORDER~28; the same
SGA~1 V \S7 finite-\'etale companion applies.

\subsection{Binary cocharacters remove a Laurent-precision promise}

\paragraph{CX-GEO-010: a polynomial-size equation system for one-parameter limits.}
\emph{Input convention.} Let $G/\mathbf Q$ be a connected split reductive
group acting linearly on $V=\mathbf Q^N$, with $N\ge1$.
The input supplies equations and inverse coordinates defining $G$,
polynomial arithmetic circuits over $\mathbf Q$ for its action in those
affine coordinates, and a weight basis for a
split maximal torus $T\simeq\mathbf G_m^r$. Its cocharacter lattice is
identified with $\mathbf Z^r$, and the complete list of coordinate
weights
\[
                     a_1,\ldots,a_N\in\mathbf Z^r
\]
is supplied in binary. These data are promised to define the indicated
group, action, and full cocharacter lattice. All dimensions, circuits,
equations, denominators, and coefficient bit lengths are charged to the
input size. The points $v,w$ have rational coordinates. Rank zero is the
trivial connected group and can be handled directly; assume $r\ge1$ below.

\emph{Statement.} The relation
\[
 \mathcal R(v,w):\quad
 \exists g,h\in G(\mathbf C),\ b\in\mathbf Z^r:\quad
       \lim_{t\longrightarrow0}\lambda_b(t)\,g v=h w,
 \qquad \lambda_b(t)=(t^{b_1},\ldots,t^{b_r}),
\]
has a polynomial-time many-one reduction to Hilbert Nullstellensatz over
$\mathbf Q$. In particular it is in AM under $M$. No bound on the
absolute exponents is supplied in unary. The equations remain a
polynomial-size system when $w$ is an additional unknown vector instead
of a fixed input.

\emph{Proof: bounded sign patterns.} Put
$H=\max(1,|a_{ji}|)$ and $B=N!H^N$.
Every sign pattern of the numbers $a_j\cdot b$ attained by an integer
$b$ is attained by another integer vector $b'$ with $|b'_i|\le B$.
Here is a complete size argument. Orient each row having nonzero sign
so that its dot product is positive, and require that dot product to be
at least one; retain a zero row as an equality. Write the unknown vector
as $x^+-x^-$, with both vectors nonnegative, and introduce a nonnegative
slack variable for each inequality. This gives a feasible system
\[
                       Cx=c,\qquad x\ge0,
\]
having $N$ rows, integral coefficients of absolute value at most $H$,
and entries of $c$ equal to zero or one. It is feasible because the
original integer vector supplies one such solution.

Choose a solution with the fewest positive coordinates. Its supported
columns of $C$ are linearly independent: otherwise a linear dependence,
oriented to have a positive coefficient, could be subtracted up to the
first vanishing coordinate, preserving nonnegativity and all equations.
Let $k\le N$ be the number of supported columns. Choose $k$ rows making
their square submatrix nonsingular. Cramer's rule gives the supported
coordinates with common positive denominator the absolute value of that
integer determinant. Multiplying by this denominator gives nonnegative
integer coordinates at most $k!H^k\le B$, since both the denominator and
each replaced-column determinant have that bound. The resulting
$x^+-x^-$ has coordinates of absolute value at most $B$. The zero
equalities remain zero, and every positive inequality remains positive,
because its right hand side has been multiplied by a positive integer.
Thus the entire prescribed sign pattern is preserved. If $k=0$, the
zero vector supplies the same conclusion.

\emph{Proof: algebraic encoding.} Choose
$L=\lceil\log_2(B+1)\rceil$. Represent $b_i$ as
\[
           b_i=\sum_{\ell=0}^{L-1}2^\ell
                         (u_{i\ell}-v_{i\ell}),
        \qquad u_{i\ell}^2=u_{i\ell},\quad
               v_{i\ell}^2=v_{i\ell}.
\]
Every complex solution of these equations assigns actual zero-or-one
values to the bits, so the represented vector is an integer vector.
Write $d_j=a_j\cdot b$ and give each coordinate three Boolean selectors
$s_j^+,s_j^0,s_j^-$ whose sum is one. Choose
\[
       J=L+\lceil\log_2(rH)\rceil+1.
\]
Let $P_j,Q_j$ be nonnegative integers represented by $J$ Boolean bits
each, and impose
\[
 s_j^+(d_j-1-P_j)=0,\qquad
 s_j^0d_j=0,\qquad s_j^-(d_j+1+Q_j)=0.
\]
The selectors now describe exactly the signs of $d_j$. All necessary
values of $|d_j|-1$ fit, since $|d_j|\le rH(2^L-1)$.
No ordering predicate over $\mathbf C$ has been inserted: the signs
are enforced by bounded integers represented by Boolean variables.

Impose the supplied equations for $g,h\in G$, including inverse
variables for all inverted denominators. Write $u=gv$ and $z=hw$ using
the supplied action circuits. For each coordinate impose
\[
              s_j^-u_j=0,\qquad
              (s_j^++s_j^-)z_j=0,\qquad
              s_j^0(u_j-z_j)=0.
\]
These are exactly the conditions that $t^{d_j}u_j$ have finite limit
$z_j$: a negative weight requires $u_j=z_j=0$, a zero weight requires
$u_j=z_j$, and a positive weight requires $z_j=0$.

Any solution therefore proves $\mathcal R(v,w)$. Conversely, from any
one-parameter limit, replace $b$ by the bounded vector having the same
signs. The limit has precisely the same coordinates, and all Boolean
and slack variables can be filled in. Hence the equation system is
feasible exactly when $\mathcal R(v,w)$ holds.
Finally $\log B=O(N\log N+N\log H)$, so $L,J$ and every coefficient
bit length are polynomial in the declared input size. Sharing the action
circuits gives a polynomial-size HN system. If $w$ is unknown, the action
circuits remain polynomial-size since the representation is linear in
that vector. Koiran's theorem gives the claimed AM consequence under
$M$.\hfill$\square$

This is an explicit encoding lemma proved here, not a claim of priority
for bounded integral witnesses in rational polyhedra. The geometric
object is a cocharacter filtration and its associated one-parameter
specialization. The inherited companions are SGA~3, Expos\'e~I,
4.6.3.1 and 4.7.1--4.7.2 for actions, and Expos\'e~VIII, 1.1--1.2 for
diagonalizable groups and their character lattices. Its run target is
the equation-access gate in CX\_GRH CX-K02/CX-K05 and ORDER~28; it
refines the bounded Laurent-arc route CX-GEO-009.

\paragraph{CX-GEO-011: orbit-closure intersection and the closed-target promise.}
\emph{Statement.} Under the explicit representation convention of
CX-GEO-010, deciding
\[
             \overline{Gv}\cap\overline{Gw}\ne\varnothing
\]
reduces in polynomial time to HN over $\mathbf Q$, and is consequently
in AM under $M$; disjointness is in coAM. If $Gw$ is promised closed,
the same upper bound holds for $w\in\overline{Gv}$. The nullcone test
$0\in\overline{Gv}$ is included without a separate promise, since the
orbit of zero is closed. These assertions concern complex Zariski
closures. They require neither defining equations for those closures nor
a unary bound on degeneration exponents.

\emph{Classical geometric input.} Every orbit closure of a reductive
group on an affine variety contains a unique closed orbit, and a
one-parameter subgroup has a limit in that orbit. An exact modern primary
locator is Bate--Herpel--Martin--R\"ohrle,
\emph{Cocharacter-closure and the rational Hilbert--Mumford Theorem},
arXiv:1411.7849v5, introduction p.~1 and Theorem~4.3, pp.~12--13,
specialized to the algebraically closed field $\mathbf C$.
Their introduction identifies the classical closed-subset formulation
as Kempf's Theorem~1.4 (1978). The current source check uses the displayed
primary theorem and its full proof, not a fresh audit of Kempf's
unretrieved original proof. Every cocharacter of a connected reductive
group is conjugate into the chosen maximal torus.

\emph{Proof.} First suppose $Gw$ is closed and
$w\in\overline{Gv}$. The unique closed orbit in that closure is then
$Gw$. The cited theorem gives a cocharacter with a limit in $Gw$.
Conjugate this cocharacter into $T$, moving the source point by a group
element and moving the limit within $Gw$. Thus $\mathcal R(v,w)$
holds. Conversely, if $\mathcal R(v,w)$ holds, its limit is in
$\overline{Gv}$; applying $h^{-1}$ and using invariance of the closure
gives $w\in\overline{Gv}$. This converse needs no closedness promise.
Applying CX-GEO-010 proves the promised membership statement.

For the intersection assertion, a nonempty intersection of the two
closures is a closed $G$-stable affine subvariety. It contains a closed
orbit: choose an orbit of minimal dimension in it, whose boundary, if
nonempty, would consist of smaller-dimensional orbits. That orbit is
closed also in $V$, since the intersection is closed. Uniqueness shows
that it is the closed orbit of both $\overline{Gv}$ and
$\overline{Gw}$. Choose a point $z$ on it. The preceding paragraph
then gives both $\mathcal R(v,z)$ and $\mathcal R(w,z)$.
Conversely, these two relations imply that $z$ belongs to both closures.
We have proved the exact existential equivalence
\[
 \overline{Gv}\cap\overline{Gw}\ne\varnothing
       \quad\Longleftrightarrow\quad
       \exists z\in\mathbf C^N:\quad
                  \mathcal R(v,z)\ \text{and}\ \mathcal R(w,z).
\]
Take two copies of the polynomial-size HN system in CX-GEO-010 and share
the unknown vector $z$. This is still one polynomial-size HN instance.
Koiran's theorem and complementation give the stated AM and coAM bounds.
\hfill$\square$

These are full reductions written here from classical geometry; no claim
that this is the first complexity formulation is made. The exact analytic
input is the Dedekind zeta function of a bounded algebraic solution field
of the constructed HN system, as in CX-HN-002. $M$ supplies it
conditionally. An $A_D$ alternative needs the bounded abelian-witness
promise for this new system. Rational input points and integral
cocharacters do not force the group elements or the shared limit point
to be defined over an abelian field.

The closed-target promise is substantive for membership. This argument
does not assert that an arbitrary point in an orbit closure is reachable
by a single cocharacter. Nor does deciding intersection construct a
separating invariant, a homogeneous system of parameters, or a
representation-theoretic multiplicity inequality. Such search and
lower-bound inputs remain the separate gates already recorded in
CX-NN-001--004 and CX-LIM-001--002. The geometric companion is the
reductive group action and its torus filtrations in SGA~3; the work
locators are CX\_GRH CX-K02/CX-K05, the GCT V orbit-intersection
interface (Theorem~9.8, p.~68), and ORDER~28.
If a GCT point is supplied by a small circuit in a representation having
exponentially many coordinates, expanding the complete coordinate and
weight lists may already be exponential. The explicit-input bound above
does not claim an efficient conversion from such a succinct encoding.

\paragraph{CX-GEO-012: sparse values at circuit gates suffice, even with large exponents.}
\emph{Statement.} Let $F:\mathbf A^n\to\mathbf A^m$ be given by
polynomial arithmetic circuits over $\mathbf Q$, using addition,
subtraction and multiplication, and let $w\in\mathbf Q^m$.
Append subtraction gates so that the outputs are $F_j-w_j$.
Given a positive integer $T$ in unary, the following existence problem
has a polynomial-time many-one reduction to HN:
\begin{quote}
Is there a Laurent-polynomial arc
$\beta(t)\in\mathbf C[t,t^{-1}]^n$ such that the value at every gate
of the augmented circuit has at most $T$ nonzero terms, and every output
belongs to $t\mathbf C[t]$?
\end{quote}
The polynomial bound is in the circuit/point input length and $T$;
there is no supplied degree bound or exponent bound. Thus the problem
is in AM under $M$. If a family is promised to admit such an arc with
polynomially bounded $T$ whenever $w\in\overline{F(\mathbf C^n)}$,
image-closure membership for that promised family is in AM under $M$.
This is a promise about values at intermediate gates, not only about the
support of the input coordinates or final outputs.

\emph{Proof.} Merge repeated occurrences of each input variable into a
single input node. For every gate $g$, introduce $T$ coefficient
variables $c_{g,i}$ and $T$ integer exponent variables $e_{g,i}$,
representing the Laurent polynomial
\[
                       A_g(t)=\sum_{i=1}^T c_{g,i}t^{e_{g,i}}.
\]
Zero coefficients and repeated exponents are allowed, so this represents
every Laurent polynomial with at most $T$ terms. At a constant gate
of value $a$, set all exponents to zero, the first coefficient to $a$,
and all other coefficients to zero. At an input-variable gate this list
defines the corresponding coordinate of $\beta$.

For an addition gate with predecessors $g_1,g_2$, form the list of
candidate terms in $A_{g_1}+A_{g_2}-A_g$; for subtraction use
$A_{g_1}-A_{g_2}-A_g$. At a multiplication gate use the $T^2+T$ terms
in
\[
              \sum_{i,j}c_{g_1,i}c_{g_2,j}
                         t^{e_{g_1,i}+e_{g_2,j}}-A_g(t).
\]
Write each candidate list as
$\sum_{a=1}^{M_g}C_{g,a}t^{E_{g,a}}$, where $M_g=O(T^2)$.
Each coefficient $C_{g,a}$ is a signed coefficient variable or a
product of two coefficient variables, and each exponent $E_{g,a}$ is
a sum of at most two exponent variables.

Collect the following integral linear forms in the vector $e$ of all
gate exponents: every difference $E_{g,a}-E_{g,b}$ within each candidate
list; each exponent fixed to zero at a constant gate; all pairwise
differences of slots at each output; and each individual output-slot
exponent. There are $R=O(sT^4)$ such forms, where $s$ is the augmented
circuit size; replace $R$ by one if necessary. Their coefficients have
absolute value at most two, even when both inputs of a multiplication
gate coincide.

Apply the bounded-sign-pattern argument of CX-GEO-010 to this list of
linear forms. Any assignment of integer exponents may be replaced by
one with
\[
                            |e_{g,i}|\le R!2^R
\]
preserving every one of the collected signs. In particular, this
preserves all equality classes of candidate exponents, the zero
exponents at constant gates, and the sign and equality classes at the
outputs. Keep all coefficient values unchanged. Every Laurent identity
at a gate is still true, because it says precisely that the sum of the
coefficients in each exponent class vanishes. The required output
condition is also preserved. Consequently a polynomial bit bound on all
exponents loses no solution.

Encode these bounded integers and the signs of all the collected forms
by the Boolean-variable construction in CX-GEO-010. Let
$q_{g,ab}$ denote the zero-sign selector for
$E_{g,a}-E_{g,b}$. For each candidate term $a$, impose
\[
                       \sum_{b=1}^{M_g}q_{g,ab}C_{g,b}=0.
\]
Since the selectors express actual equality of integer exponents,
these equations hold if and only if the candidate Laurent polynomial
is identically zero. They enforce exactly the relevant addition,
subtraction or multiplication identity, including every cancellation.
Induction over the circuit therefore shows that all the $A_g$ are its
actual values on the Laurent arc defined at the input nodes.

At an output $g$, let $q_{ij}$ indicate equality of $e_{g,i},e_{g,j}$,
and let $s_i^-,s_i^0$ be the negative and zero selectors for $e_{g,i}$.
Impose, for each $i$,
\[
                 (s_i^-+s_i^0)\sum_{j=1}^T q_{ij}c_{g,j}=0.
\]
These conditions say exactly that every exponent class of value at
most zero has coefficient zero. Hence $A_g(t)\in t\mathbf C[t]$.
Conversely any arc in the statement supplies lists and coefficients
satisfying all these equations; use the bounded replacement just proved
to fill in the Boolean variables.

There are polynomially many candidate terms, comparisons and equations.
The exponent bit bound is $O(R\log(R+1))$; the positive slack integers
needed in the sign encoding also have polynomial bit length. Coefficient
constraints use only the supplied rational constants, coefficient
variables, their pairwise products and Boolean selectors. Thus the HN
instance has polynomial encoding size, independently of the circuit's
syntactic degree. This proves the exact reduction and its AM consequence.

Finally, a feasible arc gives a map from the punctured affine line to
the source whose image under $F$ extends over zero with value $w$.
It therefore puts $w$ in the image closure. Under the stated promise,
every positive image-closure instance has such a sparse arc, proving the
last assertion.\hfill$\square$

For the bounded-degree, bounded-precision circuit setting of CX-GEO-009,
all gate values have at most $2BD+1$ possible exponents. Thus that route
is included when $BD$ is polynomial. The next example shows why sparse
gate values can remain useful when no polynomial unary precision bound
exists. The geometric language is that of maps from a punctured formal
disc whose images specialize at its closed point. For orbit-map
applications the inherited SGA~3 I 4.7.1--4.7.2 and cote 28 scan 3
companions apply. This moves the explicit-degeneration gate in
CX\_GRH CX-K02/CX-K05 and ORDER~28; it does not assert that generic
degenerations have sparse values at every circuit gate.

\paragraph{CX-GEO-013: exponential Laurent precision can be necessary at width one.}
For $s\ge1$ set $D=2^s$ and consider
\[
       F_s(X,Y)=(X,XY^D),\qquad w=(0,1).
\]
The map has an arithmetic circuit of size $O(s)$, using repeated
squaring to compute $Y^D$. The point $w$ lies in the image closure but
not in the image. The arc
\[
                         X=t^D,\qquad Y=t^{-1}
\]
has at most one term at every gate, including the zero value at the
second augmented output $XY^D-1$. It satisfies
$F_s(\beta(t))-w=(t^D,0)$.

Nevertheless every Laurent-polynomial arc with limit $w$ requires
an exponent of magnitude at least $D$ in its input coordinate $X$.
Indeed write $\operatorname{val}$ for the smallest exponent with
nonzero coefficient. A limit of zero in the first coordinate gives
$\operatorname{val}(X)>0$, while the second-coordinate limit of one
gives
\[
                \operatorname{val}(X)+D\operatorname{val}(Y)=0.
\]
Both coordinates are nonzero Laurent polynomials, so valuation is
additive on their product. The integer $\operatorname{val}(Y)$ must be
negative, whence $\operatorname{val}(X)\ge D$.
Thus the example refutes a universal polynomial unary Laurent-precision
bound in circuit size, even for two input coordinates and a width-one
degeneration. It does not refute polynomial bit bounds, which are what
CX-GEO-012 uses. The assertion that $w$ is not in the actual image
follows at once from $X=0\Rightarrow XY^D=0$; the displayed arc proves
closure membership.\hfill$\square$

This is a complete elementary counterexample proved here. Its run
locator is the bounded-precision gate of CX-GEO-009 and the refined
sparse-gate formulation CX-GEO-012, against the image-closure interface
in CX\_GRH CX-K02 and ORDER~28. It is an explicit specialization over
the punctured affine line, in the same geometric framework as that
refined statement.

\subsection{A separately isolated recent algebraic factoring input}

The following two transfers use an additional, explicitly identified
algebraic input. Write $\mathrm{PF26}$ for the algebraic assertions of
OpenAI, \emph{Deterministic Polynomial Factorization over Prime Fields},
4 October 2026: Theorem~1.2, printed/PDF p.~3, and the primary-table
interface of Proposition~8.1, pp.~35--37, with the uniform implementation
proved in Sections~2--10, pp.~6--44. Theorem~1.2 factors over a prime
field given suitable auxiliary primes, with cost polynomial in their
\emph{values}; Proposition~8.1 supplies the odd-degree splitter from a
field containing a specified root of unity and a generator of the
corresponding full primary multiplicative subgroup. Proposition~10.1,
pp.~41--44, gives the bound
\[
                  O(B^{500000}+B^{100}E^{20}),
\qquad B=20+(d+1)(\lceil\log_2p\rceil+1),
\]
for its full prime-field algorithm, where $E$ is the largest auxiliary
prime supplied. These are recent-manuscript inputs, separate from $A$
and $M$. The companion analytic assertion quoted in its Section~11 is
not used below. The complete invoked algebraic proof in Sections~2--10
has been read in this run; independent readers have also checked the
geometric Section~4 and the promised norm solver in Section~6 with its
Lemma~2.2 input. No formal certification or implementation is claimed.

\paragraph{CX-ALG-027: a polynomial transfer under $M+\mathrm{PF26}$.}
Assume $\mathrm{PF26}$. Conditional on $M$, a nonzero degree-$d$ polynomial
over any supplied finite field $\mathbb F_{p^m}$ can be completely
factored deterministically in time polynomial in $d+1$ and
$m\log p$. For each prime $r\le d$ in the large-characteristic
branch, put
\[
 E_{p,r}=\mathbb Q(\zeta_{12r},p^{1/r}).
\]
The precise sufficient analytic input is RH for all those
$\zeta_{E_{p,r}}$ and for $\zeta_{\mathbb Q}$.
There is an auxiliary prime $\ell\notin\{2,3,p,r\}$ with
\[
 \ell\equiv1\pmod{12r},\qquad
 p^{(\ell-1)/r}\not\equiv1\pmod\ell,
\]
satisfying the fully explicit sufficient bound
\[
 \ell\le
 \left[16r^2\log(12r^2)+16r(r-1)\log p+10r^2+5\right]^2.
 \tag{CX-PF26-prime}
\]
The assertion concerns finite-field polynomial factorization; it has
no integer-factorization conclusion.

\emph{Full proof of the transfer.}
Constants and degree-one polynomials are immediate. Put
$B=20+(d+1)(\lceil\log_2p\rceil+1)$ and first suppose
$p>B^{200000}$. In particular $p>\max(d,3)$.
For a prime $r\le d$ set $K=\mathbb Q(\zeta_{12r})$ and
$k=[K:\mathbb Q]\le4r$. Since $p\nmid12r$, the prime $p$ is
unramified in $K$. At any prime of $K$ above $p$, the polynomial
$X^r-p$ is Eisenstein. Thus $[E_{p,r}:K]=r$, and this extension is
cyclic because $K$ contains the $r$th roots of unity. It is Galois
over $\mathbb Q$ as well. Choose a nonidentity element of
$\operatorname{Gal}(E_{p,r}/K)$ and its conjugacy class in
$\operatorname{Gal}(E_{p,r}/\mathbb Q)$.

The cyclotomic discriminant bound and the relative power-basis
discriminant give
\begin{align*}
 \log|D_K|&\le k\log(12r),\\
 \log|D_{E_{p,r}}|
 &\le r\log|D_K|+
        k\bigl(r\log r+(r-1)\log p\bigr)\\
 &\le4r^2\log(12r^2)+4r(r-1)\log p,\qquad
 [E_{p,r}:\mathbb Q]\le4r^2.
\end{align*}
For completeness, the first inequality follows because the
discriminant of $\Phi_{12r}$ divides $(12r)^{\varphi(12r)}$:
differentiate $X^{12r}-1=\Phi_{12r}(X)H(X)$ at its primitive
roots and multiply. The field discriminant divides this polynomial
discriminant. For the second inequality, the discriminant of
$X^r-p$ is, up to sign, $r^rp^{r-1}$; the relative field
discriminant divides its ideal, and the discriminant tower formula
gives the displayed estimate.

Apply Bach--Sorenson Theorem~5.1 to $E_{p,r}/\mathbb Q$.
The exact source premise is RH for $\zeta_{E_{p,r}}$ and
$\zeta_{\mathbb Q}$; $M$ supplies it. Its bound
$(4\log|D_{E_{p,r}}|+2.5[E_{p,r}:\mathbb Q]+5)^2$
is precisely bounded by (CX-PF26-prime). The resulting unramified
prime has trivial Frobenius on $K$ and nontrivial Frobenius on the
cyclic Kummer extension over $K$. The first condition says
$\ell\equiv1\pmod{12r}$. At a degree-one prime of $K$ above
$\ell$, the second says that $X^r-p$ does not split; since
$\mathbb F_\ell$ contains $\mu_r$, it is exactly
$p^{(\ell-1)/r}\ne1$ in $\mathbb F_\ell$.
Each of $2,3,r$ ramifies in $K/\mathbb Q$, and $p$ ramifies
in $E_{p,r}/K$ by Eisenstein. Unramifiedness therefore excludes
every prime forbidden by the manuscript's auxiliary-prime definition.

The right side of (CX-PF26-prime) is polynomial in $d$ and
$\log p$; indeed it is $O(B^4)$ in this branch, since
$r\le d<p$ and $d\log p=O(B)$. Upward integer search,
trial-division primality tests, and modular powering find each
required prime with polynomial total cost in $B$. No unknown
integer factorization is needed. Theorem~1.2 and Proposition~10.1
of $\mathrm{PF26}$ now give polynomial-time factoring over
$\mathbb F_p$. In the branch $p\le B^{200000}$, ordinary
Berlekamp linear algebra and enumeration of the $p$ prime-field
scalars already give a polynomial bound in $B$.

Here is the complete passage to the supplied extension field; it
does not require a new analytic input. For a square-free polynomial
$h\in\mathbb F_{p^m}[X]$, form
\[
 C=\mathbb F_{p^m}[X]/(h),\qquad
 C_0=\ker(z\mapsto z^p-z).
\]
The kernel is computed by $\mathbb F_p$-linear algebra on a space
of dimension $m\deg h$. If $h$ has $t$ distinct irreducible
factors over $\mathbb F_{p^m}$, its quotient is a product of
$t$ finite fields, so $C_0\simeq\mathbb F_p^t$. Thus $t=1$
certifies irreducibility. If $t>1$, a computed basis supplies
$b\notin\mathbb F_p1$. Its minimal polynomial over
$\mathbb F_p$, obtained from the first dependence among
$1,b,b^2,\ldots$, is the product of the distinct linear factors
$Z-a$ corresponding to its component values. Its degree lies
between $2$ and $t\le\deg h$. Factor it by the preceding
prime-field procedure. For every root $a$, the polynomial
$\gcd(h,b-a)$ selects exactly those irreducible factors whose
component value is $a$. At least two resulting groups are nonempty,
and each gives a proper factor. Recursion decreases degrees and
therefore uses only polynomially many prime-field queries, all of
degree at most $d$.

General multiplicities are handled by derivative gcds. When the
derivative is zero, every exponent is a multiple of $p$; take the
$p$th root of each coefficient by inverse Frobenius
$a\mapsto a^{p^{m-1}}$, and recurse on the polynomial of smaller
degree, multiplying the output multiplicities by $p$. Otherwise,
split off a nontrivial derivative gcd or apply the square-free
procedure just given. At every branch the degrees decrease; at a
binary branch they sum to the parent degree. Coefficient arithmetic,
all matrices, and all powering exponents have sizes polynomial in
$d+1$ and $m\log p$. This proves the claimed complete
factorization bound.\hfill$\square$

The geometric content is a computation of primitive idempotents of a
finite \mbox{\'etale} algebra. The extra source's odd-degree mechanism
uses the Jacobian of $Y^r=F(X)$, the operator $1-\sigma$, and the
integral lattice of degree-zero divisors on its infinity fiber; its
norm implementation uses splitting of vector bundles on
$\mathbb P^1$. These are explicit geometric structures, not an
assumption that an arbitrary geometric existence argument is effective.
The attached-run comparison is CX\_GRH CX-N15/CX-U02, now refined by
CX-ALG-023--026. Those earlier quasipolynomial conclusions remain
available without the additional $\mathrm{PF26}$ input.

\paragraph{CX-ALG-028: the same polynomial transfer with the smaller finite family.}
Retain the supplied $\mathbb F_{p^m}$ and degree bound $d$.
Assume the primary-table odd-splitter assertion of
$\mathrm{PF26}$, and $A$. In the large-characteristic branch it
is enough to add RH for
\[
 \zeta_{\mathbb Q(p^{1/r})}
 \quad\text{only for primes }r\le d\text{ with }
 r\nmid(p-1)m,\qquad
 v_r\bigl(p^{\operatorname{ord}_r(p)}-1\bigr)>1.
 \tag{CX-PF26-small-family}
\]
One may restrict this family further to primes actually encountered
in the table construction. Under these inputs the same factoring
problem is deterministic polynomial time. In particular $M$ supplies
the family; and under $A+\mathrm{PF26}$ no additional analytic input
is needed when the displayed set is empty. This is a sufficient
promise and finite input list, not a necessity assertion.
Furthermore $\mathrm{PF26}$ alone suffices if every prime $r\le d$
either divides $m$ or satisfies
$v_r(p^{\operatorname{ord}_r(p)}-1)=1$ in that branch.

\emph{Full proof.}
The small-characteristic and extension-field reductions in
CX-ALG-027 are unconditional once a prime-field factorer is
available. In the large-characteristic branch, construct the
primary table in increasing prime order $r\le d$. The quadratic
entry is a nonsquare of $\mathbb F_p$, raised to the odd part of
$p-1$. If $v_2(p-1)=1$, use $-1$. If $2\mid m$, one can
instead obtain a nonsquare unconditionally by squaring a nonzero
trace-zero element of the supplied quadratic subfield. In the
remaining case $A$ supplies a nonsquare. All subsequent primes
are odd.

At an odd prime $r$, factor $\Phi_r$ over $\mathbb F_p$ using
only previously constructed entries: its degree is $r-1$.
Choose one irreducible factor and obtain
$K_r=\mathbb F_p(\zeta_r)$, of degree
$s=\operatorname{ord}_r(p)\mid r-1$. Put
$R=r^{v_r(p^s-1)}$.
If $R=r$, then $\zeta_r$ already generates the full
$r$-primary subgroup of $K_r^\times$. If $r\mid p-1$,
$A$ supplies an $r$th nonresidue in $\mathbb F_p=K_r$;
raising it to $(p-1)/R$ supplies the required primary generator.
These branches use no extra Dedekind function.

In every remaining case it suffices to construct a field
$U_r\simeq\mathbb F_{p^r}$. If $r\mid m$, it is already
available as the kernel of $z\mapsto z^{p^r}-z$ in the supplied
field. Otherwise the residual hypotheses are exactly those in
(CX-PF26-small-family), and CX-ALG-024 constructs $U_r$ in
polynomial time from $A$ and RH for the indicated pure
prime-degree field. That construction needs only its own
auxiliary-prime congruence modulo $r$; no twist by
$\mathbb Q(\zeta_{12})$ is introduced here.

For clarity, the remaining transfer from $U_r$ to a primary
generator is purely linear algebra. Since $\gcd(r,s)=1$,
$V_r=U_r\otimes_{\mathbb F_p}K_r$ is a field of degree $r$
over $K_r$. Write $Q=p^s$ and solve
\[
                      u^Q=\zeta_r u
\]
for a nonzero $u\in V_r$. This equation is $K_r$-linear.
It has a nonzero solution because $\zeta_r$ has norm one in
$V_r/K_r$, and in a cyclic finite multiplicative group the image
of $v\mapsto v^{Q-1}$ is exactly that norm-one subgroup.
The identity
$v_r(Q^r-1)=v_r(Q-1)+1$ shows that $u$ has $r$-primary
order $rR$: indeed $u^{Q-1}=\zeta_r$ has order exactly $r$.
Consequently
\[
             w=u^{(Q^r-1)/(rR)},\qquad\omega_r=w^r
\]
have orders $rR$ and $R$. The unique subgroup of $V_r^\times$
of order $R$ lies in $K_r^\times$, so $\omega_r\in K_r$.
Its coordinates there are recovered by the known inclusion and
linear algebra. All field dimensions are at most $r(r-1)$,
and all exponent bit lengths are polynomial in $r\log p$.

The induction defining the table is well founded. The even-degree
splitter is the manuscript's Proposition~2.3, pp.~9--10; it
uses only the quadratic entry. Its odd-degree splitter,
Proposition~8.1, pp.~35--37, consumes precisely a field of
degree at most $r-1$, a primitive $r$th root of unity, and the
full $r$-primary generator just constructed, for a prime
$r$ dividing the split polynomial's degree. Its bound is a fixed
polynomial in the original $B$. Berlekamp's fixed algebra and
the separating-element construction of Lemma~9.1, p.~38,
therefore factor every polynomial of degree below the next table
prime using the entries already present. Equivalently one can
repeat the recursion proof of CX-ALG-027, using the manuscript's
even and odd splitters on its split minimal polynomials. The
table is constructed once and is not recursively rebuilt. There
are at most $d$ entries and polynomially many smaller-degree
factoring calls, so all costs remain polynomial. This completes
the transfer proof. For the final promise, use $\zeta_r$ whenever
$R=r$, and otherwise use the supplied degree-$r$ subfield when
$r\mid m$. These choices never invoke the $A$ or extra Dedekind
branches. The quadratic choices just given have the same property.
Thus $\mathrm{PF26}$ alone suffices in that promised case.
\hfill$\square$

This row isolates both changes: the newer algebraic splitter supplies
the polynomial bound, while CX-ALG-023/024 supplies the smaller
analytic family and the $r\mid m$ exemption. The absence of an
extra field in a promised instance does not prove that Dirichlet GRH
controls the residual nonabelian fields in other instances. The
Kummer-cover and finite-\'etale-component dictionary is the same as
CX\_GRH CX-N13--N15/CX-U02 and CX-ALG-023--026; the recent source
adds a Jacobian and a vector-bundle computation with the exact
interfaces stated above.

\subsection{Two precise limits of the geometric promise refinements}

\paragraph{CX-GEO-014: width-one Laurent-arc feasibility is HN-complete.}
Consider the exact decision problem of CX-GEO-012 for the
\emph{supplied circuit}, including its intermediate gates. Even
when $T=1$, the target is zero, and every output polynomial has
degree at most two, this problem is polynomial-time many-one
equivalent to complex Hilbert-Nullstellensatz feasibility for
rational circuit systems. In particular, the width-one restriction
does not by itself supply a deterministic polynomial algorithm
for HN. This is a reduction statement, not a lower bound against
such an algorithm.

\emph{Full proof.}
The reduction to HN is CX-GEO-012, which is unconditional as a
reduction; $M$ is used only for its final AM classification. For
the converse, start with an HN system. Introduce one variable for
each arithmetic gate and impose the corresponding relation
$z=x+y$, $z=x-y$, or $z=xy$, together with the constant-gate
and output equations. This produces, in polynomial time, an
equivalent system
\[
                    f_1(X)=\cdots=f_s(X)=0
\]
whose polynomials have degree at most two and whose coefficient
bit lengths are polynomially bounded in the original input. Let
$X_1,\ldots,X_n$ be all its variables, introduce new variables
$Y_1,\ldots,Y_n$, and specify the circuit map
\[
 F(X,Y)=\bigl(f_1(X),\ldots,f_s(X),
            (Y_i)_{i=1}^n,(Y_i(1+X_i))_{i=1}^n\bigr).
 \tag{CX-width-one-guard}
\]
The supplied circuit explicitly computes every gate $1+X_i$
before multiplying it by $Y_i$. Thus each guard is on a path to
an output; it is not an unused gate. Every output polynomial has
degree at most two. Take the target to be zero and set $T=1$.

A solution $a$ of the HN system gives the constant arc
$X=a$, $Y=0$. Every intermediate value is a constant, hence
has support size at most one, and every output is zero. Conversely,
suppose a width-one Laurent-polynomial arc satisfies the condition
for this supplied circuit. Each input $X_i(t)$ is either zero or
$c_it^{e_i}$ with $c_i\ne0$. If $e_i\ne0$, the intermediate
value $1+c_it^{e_i}$ has two distinct nonzero terms, contradicting
the width-one requirement. Therefore all the $X_i$ are constants.
The outputs $f_j(X)$ are then constants; membership in
$t\mathbb C[t]$ forces each to be zero. They give a solution of
the original system. Both directions and the polynomial encoding
bound are now proved.\hfill$\square$

The dependence on the specified circuit is part of the problem:
one is not free to remove its guard gates by replacing it with
another presentation of the same map. The result concerns the
exact sparse-arc decision problem. It does not identify arbitrary
image-closure membership with that problem when the sparse-arc
promise is absent. Geometrically, the guards turn permitted maps
from the punctured affine line into constant maps in the original
coordinates. The run locators are CX-GEO-012, CX\_GRH CX-K02,
and the explicit-degeneration gate in ORDER~28.

\paragraph{CX-GEO-015: padded targets and affine intersection need distinct promises.}
Let $V$ be the space of homogeneous forms of positive degree $a$
in $D$ variables over $\mathbb C$, with action
$(g\cdot f)(x)=f(g^{-1}x)$. A nonzero form depending on at most
$k<D$ of those variables has $0$ in the closure of its
$\operatorname{SL}_D$-orbit. Its orbit is therefore not closed.
In particular, for $m>n\ge1$, the padded permanent
\[
                    \ell^{m-n}\operatorname{perm}_n
\]
in the usual $m^2$-variable ambient space fails the closed-target
promise of CX-GEO-011. Independently, for the
$\operatorname{GL}_D$ action on forms of any fixed positive degree,
\emph{every} two affine orbit closures intersect, since both
contain zero. Neither observation concerns projective orbit closures.

\emph{Full proof.}
After relabeling, suppose the form uses only the first $k$
coordinates. Consider the cocharacter
\[
 \lambda(t)=\operatorname{diag}
   \bigl(\underbrace{t^{-(D-k)},\ldots,t^{-(D-k)}}_{k},
         \underbrace{t^k,\ldots,t^k}_{D-k}\bigr).
\]
Its determinant is
$t^{-k(D-k)+(D-k)k}=1$. The inverse action scales each used
variable by $t^{D-k}$, so homogeneity gives
\[
                 \lambda(t)\cdot f=t^{a(D-k)}f\longrightarrow0.
\]
An invertible linear change of variables cannot carry a nonzero
form to zero. Hence zero is in the orbit closure and outside the
orbit, proving nonclosedness. The padded permanent has degree $m$
and uses at most $n^2+1$ variables. For $m>n\ge1$,
$n^2+1<m^2$, so the general argument applies with $D=m^2$.

For the general linear group, use the scalar cocharacter
$t\mapsto t^{-1}I_D$. It sends every degree-$a$ form to
$t^af$, which tends to zero. The assertion about intersections
follows for any pair, including the zero form. Projectivization
removes the zero vector, and a scalar multiple represents the
same projective point; thus this argument makes no assertion
about the corresponding projective closures.\hfill$\square$

These elementary facts explain the exact interfaces of
CX-GEO-010/011 in the GCT setting. General orbit-closure
intersection and closed-target containment remain different tasks
from containment of the standard padded permanent in a determinant
orbit closure. The attached locators are CX\_GRH CX-K02/CX-K05,
GCT\_NSRH and ORDER~28, with the action-map framework
SGA~3 I 4.7.1--4.7.2 and the inherited cote~28 scan~3 companion.
No new priority is claimed for these elementary instability facts.

% END complexity

% BEGIN root
\section{The two methods: precise transfers and repairs}
\label{sec:methods}

\subsection*{MT-A01. What the adopted Dirichlet theorem says adelically}
The supplied declarations \path{riemannHypothesis},
\path{LFunction_ne_zero_of_half_lt_re}, and
\path{LFunction_eq_zero_re} in \path{ALPOGE_Main.lean} have the
definitions printed in \path{ALPOGE_Statements.lean}.  We use these as the
cited input $A_D$, as requested.  The supplied archive contains these exported
declarations and the scan report, but not the imported tree
\path{Supplements.Kernels.MainTheoremsAll}; a kernel build of that tree is
not a check performed in this run.  This distinction concerns our evidence,
not the stipulated status of the input theorem.

Every unitary idele-class character of $\mathbb Q$ has the form
$\chi_0|\cdot|^{it}$, where $\chi_0$ is a finite-order Dirichlet character
and $t\in\mathbb R$.  Consequently $A_D$ gives the rank-one Hecke GRH over
$\mathbb Q$ for these characters, not just its finite-order specialization.
Indeed the norm map has a continuous section and compact kernel, which for
$\mathbb Q$ is $\widehat{\mathbb Z}^{\times}$ after the sign convention at
infinity.  A continuous character of this profinite group into the circle has
finite image: its image is a compact quotient of a profinite group, hence
totally disconnected, whereas an infinite closed subgroup of the circle is
the circle.  It therefore factors through some
$(\mathbb Z/N\mathbb Z)^\times$.  The remaining unitary character of
$\mathbb R_{>0}$ is $x^{it}$.  In the convention
$L(s,\chi_0|\cdot|^{it})=L(s+it,\chi_0)$, zeros are translated only in
their imaginary coordinate.  This proves the assertion.  Geometrically,
the finite-order part is a rank-one representation of the arithmetic
fundamental group; the continuous norm twist is an analytic normalization,
not an additional finite-monodromy motive.  The result does not cover
arbitrary characters over a nontrivial number field.

\subsection*{MT-M01. A nonzero counterexample to the draft's general zero-count lemma}
Masullo, Lemma 9.3.2.4, PDF p.~425, claims a local $O(\log |T|)$ zero count
for the Mellin transform of every $\varphi\in\mathcal S_E(G)$.
We checked the image of the original PDF page, including the displayed
Jensen inequality and its last inference.  The quantifier is present in the
PDF; this is not an inference from the model transcription.

Take $G=\mathbb R_{>0}$, trivial character, Haar measure $dt/t$, and set
\[
 g_a(u)=\frac{1}{\sqrt{4\pi a}}\exp\left(-\frac{u^2}{4a}\right),
 \qquad a=1+i,\qquad
 \varphi(e^u)=\frac{g_{1+i}(u)-g_{1-i}(u)}{2i}.
\]
The square roots use the branch with positive real part.  This is a real,
nonzero smooth function.  For every real $b$ and every nonnegative integer
$j$, the $j$th derivative of $e^{bu}g_{1\pm i}(u)$ is a polynomial in $u$
times $e^{bu-u^2/(4(1\pm i))}$.  Its absolute value is bounded by a
polynomial times $e^{bu-u^2/8}$, hence is rapidly decreasing.  Thus
$t^b\varphi(t)$ is Schwartz in the logarithmic coordinate for every $b$;
this is exactly the condition of Definition D.3.3.1, PDF p.~804 (also
\S1.1.2, PDF p.~11).  In particular $\varphi\in\mathcal S_E(G)$.

For $\operatorname{Re}a>0$, the Gaussian integral and analytic continuation
in $a$ give
\[
 \int_{\mathbb R}g_a(u)e^{su}\,du=e^{as^2},\qquad s\in\mathbb C.
\]
For completeness, the identity first follows by completing the square for
positive real $a$ and real $s$; both sides are entire in $s$ by Gaussian
domination on compact sets, and holomorphic in $a$ on
$\operatorname{Re}a>0$, where $\operatorname{Re}(1/a)>0$.  The identity
theorem proves the claimed continuation.  Therefore
\[
 \mathcal M\varphi(s)=e^{s^2}\sin(s^2).
\]
Its zeros on the positive imaginary axis are
$i\sqrt{k\pi}$, $k=1,2,\ldots$, and are simple.  For $I=[-1,1]$ and
$T\ge2$, their number in $I+i[T,T+1]$ is
\[
 \#\{k:T^2/\pi\le k\le(T+1)^2/\pi\}
   =\frac{2T+1}{\pi}+O(1).
\]
This is not $O(\log T)$.  The function is nonzero, so removing the zero
function from the lemma does not repair it.  The proof on p.~425 needs a
\emph{lower} bound for the modulus at a Jensen centre, whereas rapid decay
supplies upper bounds.  No statement about the zeros of a Hecke L-function
is contradicted by this example.

\subsection*{MT-M02. The displayed approximate identities are proper kernel projections}
Theorem 9.3.2.6, equation (9.3.2.10), PDF p.~426, defines subspaces by
intersecting kernels of all jet evaluations up to height $T$, then claims
their orthogonal projections tend strongly to the identity.  The printed
formula and the asserted inclusion in its proof have both been checked on
the page image.

Here is a counterexample even with a single zero, so the zero-count issue
in MT-M01 is irrelevant.  Put
\[
 g(u)=(4\pi)^{-1/2}e^{-u^2/4},\qquad
 \psi(e^u)=-g'(u).
\]
Then $\psi\in\mathcal S_E(\mathbb R_{>0})$, and integration by parts gives
$\mathcal M\psi(z)=z e^{z^2}$.  Its only zero is $0$, with multiplicity one,
strictly inside the strip for $I=[-1,1]$.  Choose Sobolev index $0$.
The draft's Hilbert norm is then
\[
 \|f\|_H^2=\int_{\mathbb R}|f(e^u)|^2e^{2|u|}\,du.
\]
The functional $\delta_0(f)=\int_{\mathbb R}f(e^u)\,du$ is continuous,
because Cauchy--Schwarz bounds its modulus by
$\|f\|_H(\int e^{-2|u|}\,du)^{1/2}=\|f\|_H$.
It is nonzero, since $\delta_0(g)=1$.  For every $T>0$ the subspace displayed
in (9.3.2.10) is the same proper closed subspace $\ker\delta_0$.
Its orthogonal projector cannot converge to the identity.  In particular,
the Riesz vector $r(e^u)=e^{-2|u|}$ has $\|r\|_H=1$ and is killed by every
one of these projectors.  It belongs to the indicated completion because
$C_c^\infty(\mathbb R)$ is contained in $\mathcal S_E$ in logarithmic
coordinates and is dense in this weighted $L^2$ space: truncate first and
then smooth on compact intervals, where the positive weight is bounded
above and below.  The ranges decrease when more evaluation
conditions are added, whereas the proof on p.~426 uses the opposite
inclusion.

\subsection*{MT-M03. The exact projection repair}
Let $H$ be a Hilbert space, let $(\ell_j)_{j\ge1}$ be continuous linear
functionals, and put
\[
 V_N=\bigcap_{j\le N}\ker\ell_j,\qquad V=\bigcap_{j\ge1}\ker\ell_j.
\]
Then $P_{V_N}\to P_V$ strongly.  The complementary projectors converge
strongly to $P_{V^\perp}$, the projector onto the closed span of the Riesz
vectors of the $\ell_j$.  Thus they give an approximate identity on the
spectral quotient $H/V\simeq V^\perp$.

Here is a full proof.  For fixed $x\in H$, put $y_N=P_{V_N}x$.  For
$M\ge N$, nestedness gives $P_{V_M}y_N=y_M$ and hence
\[
 \|y_N-y_M\|^2=\|y_N\|^2-\|y_M\|^2.
\]
The nonincreasing nonnegative sequence $\|y_N\|^2$ converges, so $(y_N)$ is
Cauchy.  Write its limit as $y$.  For each fixed $N$, all sufficiently late
$y_M$ lie in the closed space $V_N$, so $y\in V$.  Also
$x-y_N\perp V$ for every $N$; passing to the limit gives $x-y\perp V$.
Therefore $y=P_Vx$.  Taking complements proves the second assertion.
Finally the orthogonal complement of the closed span of the Riesz vectors
is precisely $V$, which proves the last identification.  Exhaustion by
height instead of an integer enumeration gives the same proof.

This repairs a functional-analytic direction error.  It does not identify
the common kernel with the closed adelic theta image, prove a Schauder
basis, establish nuclearity of the actual convolution operators, or justify
the geometric-kernel decomposition.  Those are separate statements and
must be checked in the quotient on which they are used.  In Grothendieck's
trace language this is the distinction between an annihilator and its
spectral quotient, before any trace or Fredholm determinant is taken.

\subsection*{MT-M04. An arithmetic replacement for the false general lemma}
For every fixed unitary Hecke character, the number of nontrivial zeros
in a unit-height window is $O_{K,\eta}(\log(|T|+3))$, unconditionally.
The full proof for arbitrary infinity type is NT-Z10 above: polynomial
vertical growth and a uniform lower bound at the Euler-product centre
give Jensen's estimate.  It includes pure norm twists, whose single
translated pole is excluded before applying Jensen at large height.
For a fixed primitive finite-order character a conductor-explicit
primary locator is Greni\'e--Molteni,
\emph{An explicit Chebotarev density theorem under GRH}, Appendix A,
Theorem A.1, author PDF p.~26: that appendix's zero count is unconditional.
Subtract the counts at consecutive heights.  The main term has derivative
$O(\log(|T|+3))$ for a fixed field and character and each remainder has the
same bound.  The arithmetic restriction therefore closes the local
zero-count gate, including infinite-order characters, without using GRH.
It supplies a route around MT-M01 for arithmetic divisors, but does not
establish Lemma 9.3.2.5 for arbitrary Schwartz convolution ideals.  The
constant in the infinite-order statement is fixed-character, not a
claimed uniform archimedean-conductor constant.

\subsection*{MT-ART01. A nonmonomial GL\texorpdfstring{$_2$}{2} consequence: odd Artin forms}
Assume $M_H$.  Let $\rho:G_{\mathbb Q}\to\operatorname{GL}_2(\mathbb C)$
be continuous, irreducible, odd and of finite image.  Then its Artin
L-function is entire and all of its nontrivial zeros lie on
$\operatorname{Re}s=1/2$.

The holomorphy input is already a theorem: Khare--Wintenberger,
\emph{Serre's modularity conjecture (I)}, Corollary 10.2(ii), author PDF
p.~21 (with Theorem 10.1(ii), pp.~20--21), associates $\rho$ to a weight-one
cuspidal newform.  Its standard L-function is entire.  Let $L$ be the finite
Galois extension cut out by $\rho$.  Brauer induction writes its character
as an integer linear combination of characters induced from
one-dimensional characters of subgroups of $\operatorname{Gal}(L/\mathbb Q)$.
Artin induction and class field theory therefore express $L(s,\rho)$ as
a finite product of integer powers of finite-order Hecke L-functions over
the corresponding fixed fields.  In the open critical strip, $M_H$ makes
each factor holomorphic and nonzero away from the critical line.  The
product, including its negative exponents, is consequently holomorphic and
nonzero there.  Combining this with the independent holomorphy theorem
proves GRH for $L(s,\rho)$.

This includes projective image $A_5$: no monomial expression with only
nonnegative exponents is required.  It is a new transfer in this catalogue,
not a new proof of the classical modularity theorem.  In motivic terms it
concerns odd rank-two Artin motives and their weight-one automorphic
realizations.  It gives no GRH for the weight-two L-function of an arbitrary
non-CM elliptic curve.  Nor does it prove holomorphy of arbitrary even
icosahedral or higher-dimensional Artin representations.  The inherited
Brauer divisor-support statements in RETURN\_ZETA\_TWO are the general
input behind this specialization.

\subsection*{MT-ART02. The missing positivity is a finite matrix at each zero}
Let $L/K$ be a finite Galois extension with group $G$, and let
$0<\operatorname{Re}s_0<1$.  By Brauer meromorphy define
\[
 n_\chi=\operatorname{ord}_{s=s_0}L(s,\chi,L/K)\in\mathbb Z,
 \qquad\Theta_{s_0}=\sum_{\chi\in\operatorname{Irr}(G)}n_\chi\chi.
\]
Form the $|G|$ by $|G|$ matrix
\[
 K_{g,h}=\Theta_{s_0}(g h^{-1}).
\]
Then all Artin L-functions for $L/K$ are holomorphic at $s_0$ if and only
if $K$ is positive semidefinite.  The matrix is Hermitian because
$\chi(g^{-1})=\overline{\chi(g)}$ and $n_\chi$ is real.

To prove the equivalence, consider the central convolution operator
$T=|G|^{-1}K$ on $\mathbb C[G]$.  On the matrix-coefficient space of an
irreducible character $\chi$, Schur orthogonality makes $T$ scalar with
eigenvalue $n_\chi/\chi(1)$ (interchanging $\chi$ and its dual under the
opposite convolution convention only permutes the list).  Each such space
has dimension $\chi(1)^2$ and together they exhaust the regular
representation.  Thus $K\succeq0$ is equivalent to $n_\chi\ge0$ for every
$\chi$, which is exactly absence of a pole of every factor.

The restrictions of $\Theta_{s_0}$ to cyclic subgroups are already actual
characters.  Indeed, for a cyclic subgroup $C$ and a one-dimensional
character $\eta$ of $C$, Frobenius reciprocity and Artin induction give
\[
 \langle\Theta_{s_0}|_C,\eta\rangle_C
   =\operatorname{ord}_{s=s_0}L(s,\eta,L/L^C)\ge0,
\]
because the right-hand side is an order of a holomorphic Hecke function in
the strip.  This proves positivity on every such restriction; it does not
prove the full matrix is positive.  Also
\[
 \Theta_{s_0}(1)=\sum_\chi n_\chi\chi(1)
   =\operatorname{ord}_{s=s_0}\zeta_L(s).
\]
The point of the criterion is the exact remaining object: a central
positive kernel in the Artin character ring.  $M_H$ removes all tests off
the critical line, while neither zero location nor cyclic positivity
supplies the full critical-line kernel.  An Alp\"oge-type positivity
argument would have to produce this kernel or an equivalent positive
decomposition, with the actual Artin orders as coefficients.

\subsection*{MT-ART04. All induced-linear positivity tests still leave a gap}
There is a virtual character which passes every positivity inequality
obtained by inducing a one-dimensional character of a subgroup, but is not
an actual character.  This refutes an abstract reconstruction rule; we do
not claim that this virtual character occurs as the orders of actual
Artin L-functions at a point.

Let $G=A_5$, and let $\chi$ be its four-dimensional augmentation character
for the natural action on five letters.  Define
\[
             \Theta=\operatorname{reg}_G-5\chi.
\]
The natural permutation action is doubly transitive, so its permutation
character has inner-product norm squared two.  Its invariant subspace is
one-dimensional; the orthogonal complement therefore has character
$\chi$ with norm squared one, proving that $\chi$ is irreducible.
The coefficient of $\chi$ in $\Theta$ is $4-5=-1$, so $\Theta$ is not an
actual character and its full matrix from MT-ART02 is not positive.

Nevertheless, for every $H\le G$ and every linear character $\eta$ of $H$,
\[
       \left\langle\Theta,\operatorname{Ind}_H^G\eta\right\rangle_G
                    \ge0.                                      \tag{MT.1}
\]
Here are all the details.  The conjugacy classes of $A_5$ have sizes
$1,15,20,12,12$, as follows by counting the double transpositions,
three-cycles and the two five-cycle classes.  A normal subgroup is a union
of these classes containing the identity, and its size divides $60$.
No proper nontrivial sum of these class sizes containing $1$ divides $60$;
thus $A_5$ is simple.  An action on cosets of a proper subgroup of index
$2,3,$ or $4$ would be an injective homomorphism into the corresponding
symmetric group, which is impossible by order.  Lagrange's theorem leaves
only the possible proper indices
\[
                    k\in\{5,6,10,12,15,20,30,60\}.
\]
Put $m=\langle\chi|_H,\eta\rangle_H$.  Since $\chi|_H$ has dimension four,
$0\le m\le4$.  Frobenius reciprocity says that
$\operatorname{Ind}_H^G\eta$, of dimension $k$, contains $m$ copies of
$\chi$, so $4m\le k$.  These two bounds imply $k-5m\ge0$ for every index
in the list except possibly $k=12$.  At $k=12$, $|H|=5$, and $H$ acts
regularly on the five letters.  The restriction of $\chi$ is the sum of
the four distinct nontrivial characters of $C_5$, so $m\le1$ and again
$k-5m\ge0$.  Finally
\[
 \left\langle\operatorname{reg}_G,
       \operatorname{Ind}_H^G\eta\right\rangle=k,
\]
proving (MT.1) for proper $H$.  For $H=G$, simplicity and nonabelianity
make the trivial character the only linear character, and its coefficient
in $\Theta$ is one.  This proves (MT.1) in all cases.

For comparison with the methods, Artin induction and Hecke holomorphy
give exactly inequalities of the form (MT.1) for an actual local Artin
order character.  The example proves that these inequalities alone do not
force a positive character, even before any question of zero location.
Thus a transfer of determinant positivity from the Alp\"oge method would
need additional arithmetic information, not just all the existing
induced-linear nonnegative orders.  The elementary exact computation
\path{checks/artin_cone_certificate.py} checks the negative matrix
quadratic form: with $v(g)=\chi(g)$ one has $v^*Kv=-900$.

\subsection*{MT-ART03. The classical simple-zero reduction}
Stark, \emph{Some effective cases of the Brauer--Siegel theorem},
Inventiones Mathematicae \textbf{23} (1974), Theorem 3, p.~144, proves
that if $\operatorname{ord}_{s_0}\zeta_L\le1$ and $s_0\ne1$, then all
Artin factors for $L/K$ are holomorphic at $s_0$.  We cite this classical
result rather than count another proof as a finding.  Applied with the
divisor-support theorem under $M_H$, it confines any still-possible pole
inside the strip to a multiple zero of the top-field zeta on the critical
line.  No simplicity assertion for these Dedekind zetas is being made.

\subsection*{MT-Q01. Exact path quadrics in every characteristic}
For $k\ge2$ let
\[
 V_k=\left\{(x_1,\ldots,x_k):\sum_{i=1}^{k-1}x_ix_{i+1}=0\right\}
       \subset\mathbb A^k_{\mathbb Z},\qquad
 r=\lfloor k/2\rfloor,\quad d=\lceil k/2\rceil.
\]
Then for every finite field $\mathbb F_q$, including characteristic two,
\[
 \#V_k(\mathbb F_q)=q^{k-1}+(q-1)q^{d-1}.                 \tag{MT.2}
\]
In the relative Grothendieck ring of varieties over $\mathbb Z$ one has
\[
                 [V_k]=\mathbb L^{k-1}+\mathbb L^d-\mathbb L^{d-1}.
                                                               \tag{MT.3}
\]
Here $\mathbb L=[\mathbb A^1_{\mathbb Z}]$; the identity is an identity of
classes, not a claim of a direct-sum splitting of all integral realizations.

To prove these formulas, group the odd coordinates into $x\in\mathbb A^d$
and the even coordinates into $y\in\mathbb A^r$.  The quadratic form
becomes $x^{\mathsf T}Ay$, where the integral $d$ by $r$ matrix $A$ has an
$r$ by $r$ triangular minor with diagonal entries one.  It has a left
inverse over $\mathbb Z$ and its columns extend to a unimodular basis of
$\mathbb Z^d$.  Thus after an invertible integral change of the $x$
coordinates the equation is $\sum_{j=1}^r x_jy_j=0$, with $d-r$ free
$x$ coordinates.  Over $y=0$ its fibre is $\mathbb A^d$.  Over
$\mathbb A^r\setminus\{0\}$, the map of vector bundles
$\mathcal O^d\to\mathcal O$ given by $(y_1,\ldots,y_r,0,\ldots,0)$ is
surjective, and its kernel is a vector bundle of rank $d-1$.
It is Zariski locally trivial on the opens where a coordinate $y_j$ is
invertible.  Additivity and the vector-bundle relation give
\[
 [V_k]=\mathbb L^d+(\mathbb L^r-1)\mathbb L^{d-1},
\]
which is (MT.3).  Counting the same fibres over $\mathbb F_q$ proves
(MT.2) directly, without a purity theorem.

Consequently the exact error term in the unweighted orthogonality identity
of the supplied Fable sketch is
\[
             q\,\#V_k(\mathbb F_q)=q^k+(q-1)q^d.
\]
For a fixed nonzero additive-character parameter its quadratic sum is
$q^d$, by summing first over $x$.  This gives square-root size when $k$ is
even; for odd $k$ the extra factor $q^{1/2}$ reflects the one-dimensional
radical.  It explains the sign and the exact constant for this family.
The general weighted operator estimate in
\path{OURS_FABLE_PROGRAM_SUGGESTIONS_20261120_EXPLAINED.md},
\S3, remains a separate inherited result.  This calculation does not prove
the sign or norm assertions for arbitrary sheaves in Alp\"oge's steps 3
and 5, nor identify this elementary model with all the constructions in
the omitted Lean dependency tree.

\subsection*{MT-Q02. Its arithmetic zeta function and shifted zero lines}
The exact all-prime formula in MT-Q01 gives
\[
 Z(V_k/\mathbb F_q,T)
     =\frac{1-q^{d-1}T}{(1-q^{k-1}T)(1-q^dT)},
 \qquad
 \zeta_{V_k}(z)
     =\frac{\zeta(z-k+1)\zeta(z-d)}{\zeta(z-d+1)}.       \tag{MT.4}
\]
Indeed apply (MT.2) over every $\mathbb F_{q^m}$ in
$Z=\exp(\sum_{m\ge1}\#V_k(\mathbb F_{q^m})T^m/m)$ and use
$-\log(1-aT)=\sum_{m\ge1}a^mT^m/m$.  Multiplying the resulting factors
over rational primes proves the second identity first in a right
half-plane of absolute convergence, and then meromorphically.
Under $A_D$ the nontrivial divisor of this function is supported on
\[
    \operatorname{Re}z=k-\tfrac12,\qquad
    \operatorname{Re}z=d+\tfrac12,\qquad
    \operatorname{Re}z=d-\tfrac12.
\]
Divisor orders are signed and cancellations are kept exactly as in
(MT.4); the denominator's zeros are possible poles, not additional zeros.
The shifted pole of $\zeta$ at $1$ and shifted trivial zeros are recorded
by the same explicit formula, separately from this nontrivial-divisor
assertion.  This is an all-characteristic, concrete Artin--Tate example
linking the path model of the method with arithmetic zeta functions and
the Grothendieck ring.  It illustrates why a mixed-weight scheme has
several shifted lines rather than a single universal critical line.

\subsection*{MT-L01. Damped Euler products and Adams operations}
Let $\rho$ be a finite-image complex representation of $G_K$ with no
invariant vectors, of dimension $d$.  At each finite place $v$ put
$A_v=\rho(\operatorname{Frob}_v)|_{V^{I_v}}$.  Its eigenvalues have
absolute value one.  For $|c|\le1$ define, initially for $\operatorname{Re}s>1$,
\[
 D_c(s,\rho)=\prod_v\det(1-cA_v(Nv)^{-s})^{-1}.
\]
If $L(s,\rho)$ is holomorphic and nonzero on
$H=\{\operatorname{Re}s>1/2\}$, then $D_c$ has a holomorphic nonvanishing
continuation to $H$, given by
\begin{align*}
 D_c(s,\rho)&=\exp\{c\log L(s,\rho)+E_c(s)\},\\
 E_c(s)&=\sum_v\sum_{k\ge2}
   \frac{c^k-c}{k}\operatorname{tr}(A_v^k)(Nv)^{-ks}.
\end{align*}
The branch of $\log L$ is fixed by the Euler series in
$\operatorname{Re}s>1$.

For a full proof, expand $-\log\det(1-X)$ in powers of $X$ in the initial
half-plane.  The terms with $k=1$ cancel in
$\log D_c-c\log L$, giving the displayed $E_c$.  On a compact subset with
$\operatorname{Re}s\ge1/2+\varepsilon$, its terms in absolute value sum to
at most
\[
 2d\sum_v\sum_{k\ge2}(Nv)^{-k(1/2+\varepsilon)}
 \le \frac{2d}{1-2^{-1/2-\varepsilon}}
       \sum_v(Nv)^{-1-2\varepsilon}<\infty.
\]
Normal convergence, also uniform in $|c|\le1$, proves that $E_c$ is
holomorphic in $s$.  A nonvanishing holomorphic function on the simply
connected half-plane $H$ has a holomorphic logarithm, uniquely normalized
by its existing branch on the right.  The exponential above is therefore
holomorphic, nonzero, and agrees with the product where it converges.
This proves the continuation.

Under $M_H$, Brauer induction supplies the needed holomorphic nonvanishing
on $H$ for $\rho$ with no invariants: all Hecke factors have no zeros
there, and their orders at $1$ combine to
$-\dim V^{G_K}=0$.  This does not require Artin holomorphy on the boundary
line.  Under $A_D$ the same argument applies to sums of nonprincipal
Dirichlet characters over $\mathbb Q$.

The traces $\operatorname{tr}(A_v^k)$ are the local Adams-power characters;
the proof separates the first prime term from all higher Adams terms,
whose series converge in the larger half-plane.  This is a concrete
$\lambda$-ring link, rather than a claim that positivity of coefficients
follows from the formal operations.  For $c=1/m$ the phenomenon is closely
related to Gun--Hazra--Sahu, \emph{On holomorphy and non-vanishing of Artin
L-functions} (2025), Theorem 5.  The displayed formula covers all complex
$|c|\le1$ and makes its hypotheses explicit.

\subsection*{MT-L02. One irrational parameter detects the zero-free half-plane}
Fix an irrational real number $0<c<1$ and a nontrivial irreducible Artin
representation $\rho$.  Then the following are equivalent:
\begin{enumerate}
 \item $L(s,\rho)$ is holomorphic and nonzero on $H$;
 \item the Euler product $D_c(s,\rho)$ admits a single-valued meromorphic
       continuation to $H$.
\end{enumerate}
In the second assertion meromorphic is a global, single-valued condition;
a branch on a slit domain is not enough.

MT-L01 proves the forward implication.  Conversely $L$ is meromorphic by
Brauer induction and $E_c$ is holomorphic on $H$ by the same normally
convergent series without any hypothesis about the zeros of $L$.
Logarithmic differentiation in $\operatorname{Re}s>1$ gives
\[
 \frac{D_c'}{D_c}=c\frac{L'}L+E_c'.
\]
Both sides are meromorphic on connected $H$ if the asserted continuation
of $D_c$ exists, so the identity holds on $H$.  Taking residues at any
$s_0\in H$ gives
\[
 \operatorname{ord}_{s_0}D_c
      =c\operatorname{ord}_{s_0}L.
\]
Both orders are integers, and $c$ is irrational.  Both must be zero.
Thus $L$ is holomorphic and nonvanishing on $H$, proving the equivalence.
If instead $c=a/b$ in lowest terms, the same equality proves
$b\mid\operatorname{ord}_{s_0}L$ everywhere in $H$.

GRH in the meromorphic-divisor sense for the nontrivial irreducible Artin
factors is therefore equivalent, by the functional equation, to these
irrational-parameter continuation tests.  To include the trivial factor one
must separately add RH for $\zeta_K$; that input is already part of $M_H$.
This is an analytic multiplicity test.  It neither constructs the
continuation from arithmetic lattices nor removes poles on the critical
line.  With a trivial summand of multiplicity $r$, the known pole at $1$
contributes local monodromy $e^{-2\pi i c r}$; this is why the no-invariants
hypothesis was retained rather than suppressed.

\subsection*{MT-ART05. The exact induced-linear cone for \texorpdfstring{$A_5$}{A5}}
The coarse counterexample MT-ART04 can be sharpened to a degree-three
virtual character. More usefully, the same calculation gives explicit
lower bounds for any pole of an actual $A_5$ Artin function. These are
group-theoretic deductions, not realizations of the virtual examples as
arithmetic divisor orders.

Write the irreducible characters in the order
$1,\chi_3,\chi'_3,\chi_4,\chi_5$, and put
\[
 \Theta=a+b\chi_3+c\chi'_3+d\chi_4+e\chi_5,
 \qquad a,b,c,d,e\in\mathbb Z.
\]
Then the inequalities
$\langle\Theta,\operatorname{Ind}_H^{A_5}\eta\rangle\ge0$ for
every subgroup $H$ and every linear character $\eta$ are equivalent to
the following seven inequalities:
\[
 a\ge0,\quad e\ge0,\quad a+d\ge0,\quad b+c\ge0,\quad
 b+d+e\ge0,\quad c+d+e\ge0,\quad b+c+d\ge0.       \tag{MT-A5-cone}
\]
If $r=\Theta(1)=a+3b+3c+4d+5e$, then $d<0$ forces $r\ge3$,
and $b<0$ or $c<0$ forces $r\ge4$. These thresholds are sharp for
this abstract cone: respectively
\[
       1+\chi_3+\chi'_3-\chi_4,
       \qquad -\chi_3+\chi'_3+\chi_4               \tag{MT-A5-ex}
\]
have degrees three and four and satisfy every inequality.

\begin{proof}
We give the character and subgroup details so the finite assertion has
no unreported character-table oracle. The five conjugacy classes have
orders $1,2,3,5,5$ and sizes $1,15,20,12,12$. The natural four-dimensional
augmentation character has values $(4,0,1,-1,-1)$ and norm one, as in
MT-ART04. Its real representation is oriented because all permutations
are even. The Hodge-star decomposition of its exterior square has two
invariant real three-dimensional spaces. The exterior-square character
has values $(6,-2,0,1,1)$ by
\[
 \chi_{\wedge^2 V}(g)=\tfrac12(\chi_V(g)^2-\chi_V(g^2)).
\]
Its norm squared is two and it has no invariant vectors, so the two
three-dimensional summands are distinct irreducibles. Simplicity of
$A_5$ makes each faithful and its determinant trivial. An order-two or
order-three element in either corresponding $\mathrm{SO}_3$ representation
has trace $-1$ or $0$. At order five the two traces are
$\varphi=(1+\sqrt5)/2$ and $\bar\varphi=(1-\sqrt5)/2$; their sum is
one and each is the trace of a nontrivial order-five rotation.

There are six Sylow-five subgroups. Their normalizers have order ten
and are dihedral, since the centralizer of a five-cycle is its own
cyclic subgroup. The permutation action on the six subgroups has
values $(6,2,0,1,1)$: an order-five element fixes its unique Sylow
subgroup, an order-three element fixes none, and double counting gives
$6\cdot5/15=2$ fixed subgroups for an involution. Removing the constants
gives an irreducible character $\chi_5=(5,1,-1,0,0)$, of norm one.
The squared dimensions $1+9+9+16+25=60$ exhaust the character table.

Every proper subgroup is contained in an $A_4$, a $D_5$ of order ten,
or an $S_3$ of order six. Here is an elementary verification. Simplicity
excludes proper indices two, three and four. The remaining possible
proper subgroup orders divide sixty and belong to
$\{1,2,3,4,5,6,10,12\}$. There are no elements of order four, six or
ten. Thus subgroups of orders four, six and ten have types
$V_4,S_3,D_5$. The centralizer of an involution is its unique $V_4$,
whose normalizer is $A_4$. In a subgroup of order twelve, a normal
Sylow-three subgroup is impossible since the ambient normalizer of
$C_3$ has order six. There are therefore four Sylow-three subgroups,
contributing eight nonidentity elements. Three Sylow-two subgroups
would contribute nine further involutions, because distinct $V_4$'s
have disjoint nonidentity sets. Hence the Sylow-two subgroup is normal
and the order-twelve subgroup is its $A_4$ normalizer. The smaller
cyclic and Klein subgroups are contained in these listed groups.

Let $u,\bar u$ be the nontrivial linear characters of $A_4$,
$v$ its irreducible character of degree three, and let $\epsilon$
denote a sign character of either $D_5$ or $S_3$. The restriction
decompositions obtained from the displayed character values are
\[
\begin{array}{c|ccc}
 & A_4 & D_5 & S_3\\\hline
1&1&1&1\\
\chi_3&v&\epsilon+v_1&\epsilon+v_2\\
\chi'_3&v&\epsilon+v'_1&\epsilon+v_2\\
\chi_4&1+v&v_1+v'_1&1+\epsilon+v_2\\
\chi_5&u+\bar u+v&1+v_1+v'_1&1+2v_2
\end{array}                                                   \tag{MT-A5-res}
\]
Here $v_1,v'_1$ are the two dihedral characters of degree two and
$v_2$ is the standard $S_3$ character. The nonlinear characters in
these three groups are themselves induced from linear characters:
$v$ from a nontrivial character of $V_4$, $v_1,v'_1$ from $C_5$,
and $v_2$ from $C_3$. Transitivity of induction and Frobenius reciprocity
therefore show that all induced-linear tests imply nonnegative
multiplicities in every column of (MT-A5-res). Together with the
global trivial-character test $a\ge0$, these are exactly the seven
inequalities (MT-A5-cone), after deleting redundant inequalities.
Conversely those seven inequalities make every column an actual
character. Restricting further gives actual characters on all proper
subgroups, hence every required linear inner product is nonnegative.
The global linear test is again $a\ge0$. This proves the equivalence.

For the degree bounds, first put $d=-k<0$. Then $a\ge k$,
$b+c\ge k$, and $b,c\ge k-e$. If $e=0$, these imply
$r\ge k+6k-4k=3k\ge3$. If $e\ge1$, the first two inequalities give
$r\ge k+3k-4k+5e\ge5$. Next put $b=-k<0$. If $d\ge0$, we have
$c\ge k$, $d+e\ge k$, whence
$r\ge4d+5e\ge4k$. If $d=-j<0$, then
$a\ge j$, $c\ge k+j$, and $e\ge k+j$, giving
$r\ge j+3j-4j+5(k+j)\ge4k$. The argument for $c<0$ is identical.
Finally direct substitution verifies both examples (MT-A5-ex).
\end{proof}

\subsection*{MT-ART06. Pole multiplicity gates for actual \texorpdfstring{$A_5$}{A5} extensions}
Let $L/K$ be an $A_5$ extension and $s_0\ne1$. Its Heilbronn character
has coefficients equal to the local Artin orders. Hecke holomorphy and
induction give all the tests in MT-ART05, unconditionally. Hence every
Artin factor is holomorphic at a zero of $\zeta_L$ of order at most two.
A pole of the four-dimensional factor requires order at least three
and a zero of $\zeta_K$ at the same point. A pole of either
three-dimensional factor requires order at least four. Conditional on
$M_H$, all remaining nontrivial candidate poles lie on the critical line.

At the central point the lower bounds strengthen to six and eight,
respectively. To prove this final refinement, the permutation actions
of degrees five and six give
\[
 L(s,\chi_4)=\frac{\zeta_{L^{A_4}}(s)}{\zeta_K(s)},\qquad
 L(s,\chi_5)=\frac{\zeta_{L^{D_5}}(s)}{\zeta_K(s)}.
\]
The completed Dedekind zeta functional equation has sign $+1$, and
its completion factors are finite and nonzero at $1/2$. Consequently
the central orders $a,d,e$ are even integers, with $a,e\ge0$ by
(MT-A5-cone). If $d=-k<0$, then $k\ge2$; the preceding proof gives
$r\ge3k\ge6$ when $e=0$ and $r\ge5e\ge10$ otherwise.
If $b=-k<0$ and $d\ge0$, the nonnegative even integers $d,e$
satisfy $d+e\ge k\ge1$, so $r\ge4d+5e\ge8$.
If $d=-j<0$, then $j\ge2$ and the preceding bound
$r\ge5(k+j)$ is stronger. The same holds for $c<0$.
This proves all assertions; no assumption on simplicity of zeros has
been introduced. The degree-three and degree-four virtual examples
only show the limits of the indicated character tests away from the
additional central parity constraints.

These refinements put the residual Artin question into explicit local
inequalities on a finite Galois quotient, within the SGA 1 V \S7
framework. They sharpen ORDER item 42 and ZETA\_TWO ZT02 without
claiming general Artin holomorphy or constructing actual negative
Artin divisor orders.

\subsection*{MT-H01. A unitary Hecke family outside the algebraic-CM normalization}
Let $K=\mathbb Q(\sqrt5)$, $\epsilon=(1+\sqrt5)/2$, and let $\sigma$
be its nontrivial automorphism. Put $t_0=\pi/\log\epsilon$.
For each nonzero integer $k$ there is a conductor-one unitary Hecke
character defined on fractional ideals by
\[
       \eta_k((\alpha))=
       \left|\frac{\alpha}{\sigma(\alpha)}\right|^{ik t_0}.
                                                               \tag{MT-H1}
\]
It has infinite order, $\eta_k^\sigma=\eta_{-k}$, and is trivial
on the rational ideles. No positive power of $\eta_k$ becomes an
algebraic Hecke character after multiplying by a complex norm power
and a finite-order character. Thus the full unitary-Hecke input $M_H$
has an explicit scope beyond the algebraic Hecke characters used in
the potentially CM motive transfer.

\begin{proof}
The ring of integers is $\mathbb Z[\epsilon]$, with discriminant five.
Minkowski's ideal-class bound is $\sqrt5/2<2$, so every fractional
ideal is principal. An exact classical locator for the bound is
Milne, \emph{Algebraic Number Theory}, Theorem 4.3 (v3.08), PDF p.~72.
The units are $\{\pm\epsilon^j:j\in\mathbb Z\}$. For a direct
verification, a positive unit $u$ with $1<u<\epsilon$ would have
integer trace $u+u^{-1}$ in $(2,\sqrt5)$ if its norm were one, or
$u-u^{-1}$ in $(0,1)$ if its norm were minus one; neither interval
contains an integer. Multiplying an arbitrary positive unit by a
power of $\epsilon$ reduces it to $[1,\epsilon)$, so it is a power
of $\epsilon$, and signs give the assertion.

Multiplying $\alpha$ by a unit changes the logarithm of the absolute
ratio in (MT-H1) by $2j\log\epsilon$. Its phase therefore changes by
$2\pi i k j$, so (MT-H1) is well-defined and multiplicative. It is
unitary. The classical ideal/infinity-type description of Hecke
characters makes it an unramified finite-place character, with
archimedean exponents $\pm ikt_0$ (their simultaneous sign depends
on the ideal-to-idele convention). The unit compatibility just checked
is exactly what is required for this description. Rational ideles
have equal archimedean components and equal valuations at conjugate
places; their images are trivial. For a ramified rational place this
also follows here from the generator $\sqrt5$, whose absolute ratio
is one. Conjugating (MT-H1) inverts the ratio.

On the positive archimedean subgroup the exponents of any algebraic
Hecke character are integers, up to an overall convention sign.
A norm power adds the same complex number to the two exponents;
a finite-order character is trivial on this connected subgroup.
If a positive power $m$ of (MT-H1) had the asserted algebraic
normalization, subtraction of its two exponents would give
$2mikt_0$ equal to the difference of two integers. This is impossible.
The same nonconstant archimedean restriction proves infinite order.
\end{proof}

This does not assert a transcendence theorem for individual Hecke
values, or rule out every conceivable generalized motive. It precisely
separates the character class covered by CM01 from the broader unitary
character class quantified in Masullo's endpoint. A proposed transport
of the CM theta-value and algebraic-height steps must therefore address
this extra archimedean data; norm twisting alone cannot do so.

\subsection*{MT-H02. An explicit nonholomorphic \texorpdfstring{$\mathrm{GL}_2$}{GL2} family}
For $k\ge1$, the classical Hecke--Maa\ss\ correspondence associates to
$\eta_k$ a cuspidal dihedral Maa\ss\ newform $g_k$ of weight zero,
level five and nebentypus $\chi_5$. Its spectral parameter is $kt_0$,
its Laplace eigenvalue is $1/4+k^2t_0^2$, and
\[
       L(s,g_k)=L_K(s,\eta_k).
                                                               \tag{MT-H2}
\]
Consequently $M_H$ gives GRH for every member of this family. Its
entireness is already the classical Hecke theorem. The exact primary
locator for the correspondence, level, spectral parameter and coefficient
identity is Humphries--Khan, \emph{On the Random Wave Conjecture for
Dihedral Maa\ss\ Forms}, GAFA 30 (2020), Appendix A.1, printed pp.~106--107 (PDF pp.~73--74), the displayed
formula for $\psi((\alpha))$ and the subsequent formulas for
$t_g$ and $\lambda_{g_\psi}(n)$.

For clarity, the transfer proof is short and complete. MT-H01 gives
a non-norm character: its two archimedean exponents are unequal.
The cited correspondence is therefore cuspidal, and its formula
$\lambda_{g_k}(n)=\sum_{N\mathfrak a=n}\eta_k(\mathfrak a)$ identifies
the absolutely convergent Dirichlet series with (MT-H2), including all
finite primes. Analytic continuation gives the equality everywhere.
The character is unitary, so its nontrivial critical line is
$\operatorname{Re}s=1/2$, which is exactly the assertion supplied by
$M_H$. This proves the stated implication without a generic
non-dihedral Maa\ss\-form hypothesis.

These forms live on the analytic modular curve of level five and in
the automorphic-induction part of the adelic framework. They are an
explicit additional $\mathrm{GL}_2$ branch beside the holomorphic
CM and finite-image weight-one branches, not the missing general
non-CM $\mathrm{GL}_2$ theorem in ORDER item 42.

\subsection*{MT-H03. Exact tensor and symmetric-power products in that family}
Write $F_j(s)=L_K(s,\eta_j)$ for $j\ne0$ and
$F_0(s)=\zeta_K(s)=\zeta(s)L(s,\chi_5)$. Then $F_{-j}=F_j$.
For the actual local two-dimensional parameters $V_j$ of $g_j$, the
following finite Euler-product identities hold:
\[
 L(s,V_m\otimes V_n)=F_{m+n}(s)F_{m-n}(s),\qquad m,n\ge1,
                                                               \tag{MT-H3}
\]
and, for $n\ge0$ and $k\ge1$,
\[
 L(s,\operatorname{Sym}^n V_k)=
 \prod_{j=0}^{\lfloor(n-1)/2\rfloor}F_{(n-2j)k}(s)
 \begin{cases}\zeta(s),&n\text{ even},\\1,&n\text{ odd}.
 \end{cases}                                                   \tag{MT-H4}
\]
The product is empty when $n=0$. Conditional on $M_H$, all nontrivial
zeros of these products lie on $\operatorname{Re}s=1/2$. They are
entire except for a simple pole at one when $m=n$ in (MT-H3), or
when $n$ is even in (MT-H4). This is a covered arbitrarily large
degree family, not a use of an unknown general symmetric-power lift.

\begin{proof}
Conjugating ideals proves $F_{-j}=F_j$. At a split rational prime
the two local parameters of $V_k$ are $z^k,z^{-k}$, where
$z=\eta_1(\mathfrak p)$. The tensor weights are
$z^{m+n},z^{m-n},z^{-m+n},z^{-m-n}$, which gives (MT-H3).
The symmetric-power weights are $z^{(n-2j)k}$ for $0\le j\le n$;
pairing opposite weights leaves one middle weight $1$ precisely for
even $n$, proving (MT-H4) there.

At an inert prime the unique prime ideal is $(p)$ and has character
value one. The two induced Frobenius eigenvalues are $1,-1$ for
every nonzero $k$. Tensoring gives two of each sign. On the
symmetric power, the even-index monomials give eigenvalue $1$ and
the odd-index monomials give $-1$. If $n=2r$, these multiplicities
are $r+1,r$; if $n=2r+1$, they are $r+1,r+1$. These are exactly
the factors on the right sides. If $m=n$, the factor $F_0$ at an
inert prime is $(1-p^{-2s})^{-1}$, so the same check includes that
case.

The only ramified prime is five. The ideal $(\sqrt5)$ has
$\eta_k$-value one and the character is unramified at every finite
place. Thus its local character at that quadratic completion is
trivial, and its induced Weil parameter is $1\oplus\chi_{5,5}$.
The inertia-fixed dimension is one. A tensor product has two
trivial constituents and two $\chi_{5,5}$ constituents; a symmetric
power has respectively $r+1,r$ or $r+1,r+1$ of them. The latter
constituents have no inertia invariants. Hence the ramified Euler
factors also agree exactly, rather than only away from five.
Equivalently the proof is the weight-pair decomposition for diagonal
matrices $\operatorname{diag}(z,z^{-1})$ and an outside coset
acting by interchanging the two basis vectors.

Absolute convergence initially gives the identities for
$\operatorname{Re}s>1$. The right sides give their meromorphic
continuations. All $F_j$, $j\ne0$, are entire Hecke functions with
nonzero boundary value at one; the latter is Ramakrishnan--Valenza,
Theorem 7-28, printed p.~289. The same holds for $L(1,\chi_5)$.
Thus the displayed zeta factors give exactly the asserted simple
poles, with no cancellation at one. Finally $M_H$ locates all the
nontrivial zeros of each unitary Hecke factor. The identities contain
no negative exponents, so they introduce no hidden poles on the
critical line. This proves every assertion.
\end{proof}

The middle summand in (MT-H4) is the trivial character for this
specific family. In a differently normalized dihedral family it
depends on the central character; replacing it universally by either
$1$ or $\chi_K$ would be incorrect. This is the representation-ring
and lambda-operation part of the adelic framework, beside SGA 6 V
\S\S2,7, with the explicit analytic parameters of MT-H01 retained.

\subsection*{MT-ART07. The entire non-real $A_5$ Artin category}
Let $L/\mathbb Q$ be Galois with group $A_5$, and suppose $L$ is not
totally real. Every nontrivial irreducible Artin L-function for this
extension is entire, unconditionally. Under $M_H$ each satisfies GRH.
This closes the non-real $A_5$ specialization of the holomorphy gate;
MT-ART05/06 remain useful for totally real extensions and other base fields.

The holomorphy assertion is classical, not a new finding. An exact primary
locator is Calegari, \emph{The Artin conjecture for some $S_5$-extensions},
\S1, publisher-author PDF p.~2, the paragraph immediately before
Theorem 1.1: it states the stronger automorphy assertion for $A_5$ when
complex conjugation is nontrivial, citing Khare--Wintenberger and Kim's
icosahedral functorial constructions. The two conditions are equivalent
here because a Galois extension has a real embedding exactly when every
embedding is real, equivalently complex conjugation acts trivially.
We cite that classical result rather than repeat its proof.

Here is the complete GRH transfer and its tensor-category form.
For every finite-dimensional representation $V$ of $A_5$, put
$m=\dim V^{A_5}$. Complete reducibility gives
\[
 L(s,V)=\zeta(s)^m
          \prod_{\chi\ne1}L(s,\chi)^{\langle V,\chi\rangle},
 \qquad \langle V,\chi\rangle\in\mathbb Z_{\ge0}.
\]
The product after removing $\zeta^m$ is entire by the cited theorem.
For each irreducible factor, Brauer induction expresses its L-function
as a finite product of integer powers of finite-order Hecke L-functions.
At a point of the open critical strip off its middle line, $M_H$ makes
every such factor finite and nonzero. The product is consequently nonzero
there, including when some Brauer exponents are negative. This proves
the claimed GRH without making a positivity assumption on those exponents.

At $s=1$ every nontrivial irreducible factor is also nonzero. Indeed in a
Brauer expression $\chi=\sum_j a_j\operatorname{Ind}_{H_j}^{A_5}\eta_j$,
the order there is $-\sum_{\eta_j=1}a_j$, by the simple Dedekind pole and
Hecke boundary nonvanishing. Frobenius reciprocity identifies the sum with
$\langle\chi,1\rangle=0$. Thus $L(s,V)$ has a pole of order exactly $m$
at $1$, and has no other poles. Tensor products, duals, exterior and
symmetric powers, and every Schur functor produce representations of the
same finite group, so the same conclusions hold for them after removing
their explicitly determined invariant summands. This is a statement about
the whole Artin tensor category attached to this cover, not about a
general positive-weight motive. The companions in Grothendieck's language
are SGA~1, Expos\'e~V, \S7 (finite monodromy), and SGA~6, Expos\'e~V,
\S\S2,7 (representation-ring operations), together with AG-LAMBDA.

\subsection*{MT-ART08. A finite exceptional set is enough for a holomorphy transfer}
Here is a useful precise comparison lemma. Let $V$ be a self-dual,
finite-image complex Galois representation over a number field $K$.
Suppose an Euler product $F(s)$ converges in a right half-plane and
continues to an entire function, with the same finite local factors as $L(s,V)$ outside a
finite set $S$. Suppose also that, for each $v\in S$, the reciprocal of
the local factor of $F$ is a polynomial $Q_v((Nv)^{-s})$.
Then $L(s,V)$ is entire. No equality of ramified local parameters is
needed for this conclusion.

To prove this, write the Artin factor at $v$ as
$P_v((Nv)^{-s})^{-1}$. Every eigenvalue of Frobenius on the inertia
invariants is a root of unity, since $V$ has finite image. Hence the zeros
of $P_v((Nv)^{-s})$ all lie on $\operatorname{Re}s=0$.
Equality in a right half-plane, followed by meromorphic continuation,
gives
\[
 L(s,V)=F(s)\prod_{v\in S}
             \frac{Q_v((Nv)^{-s})}{P_v((Nv)^{-s})}.
\]
The right side is holomorphic on $\operatorname{Re}s>0$.
The classical completed Artin functional equation is
\[
 A^{s/2}\Gamma_V(s)L(s,V)
   =W A^{(1-s)/2}\Gamma_V(1-s)L(1-s,V),
\]
with $A>0$ and $W\ne0$; the gamma factor is a product of
$\Gamma_{\mathbb R}(s)$, $\Gamma_{\mathbb R}(s+1)$ and
$\Gamma_{\mathbb C}(s)$ with nonnegative integral multiplicities.
An exact locator for these formulas is Prasad--Yogananda,
\emph{A Report on Artin's holomorphy conjecture}, \S3.3,
PDF pp.~8--9, with Brauer induction in Theorem 3.3.
For $\operatorname{Re}s<1$, the numerator
$\Gamma_V(1-s)L(1-s,V)$ is holomorphic, since
$\operatorname{Re}(1-s)>0$. The reciprocal $1/\Gamma_V(s)$ is entire.
The functional equation therefore makes $L(s,V)$ holomorphic on this
second half-plane too. The two half-planes cover $\mathbb C$, proving
the assertion. For a non-self-dual $V$ the identical argument works if
the hypotheses are also supplied for $V^\vee$.

The finiteness of $S$ is an actual hypothesis. Agreement at a set of
places of density one does not give the finite product in the proof.
In particular this lemma cannot by itself upgrade the weak
five-dimensional correspondence in Calegari's Theorem 1.1.
In geometric terms this separates modification at finitely many closed
points of an arithmetic base from a density statement about infinitely
many Frobenius fibres.

\subsection*{MT-S501 and MT-S502. The $S_5$ transfer and its one-value gate}
Let $L/\mathbb Q$ be an $S_5$ extension, with complex conjugation of cycle
type $(12)(34)$. Assume $5$ is unramified and its Frobenius also has this
cycle type. Then the irreducible factors of dimensions $4$ and $6$ are
entire and, under $M_H$, satisfy GRH (MT-S501).

Let $E=L^{A_5}$, let $F\subset L$ have degree $6$ over $\mathbb Q$, and
put $H=EF$. Under $M_H$, the additional single condition
\[
                         \zeta_H(1/2)\ne0
\]
implies that every nontrivial irreducible factor is entire and satisfies
GRH (MT-S502). This is a proved reduction with the displayed extra input,
not a proof of that central nonvanishing.

The classical inputs are Calegari's Theorems 1.1 and 1.2,
publisher-author PDF pp.~2--3; the definition $H=EF$ is explicit in
Theorem 1.2 of arXiv:1112.1152v1, p.~2. The former supplies the
dimension-$4$ and dimension-$6$ holomorphy, and the latter supplies the
remaining case if $\zeta_H$ has no zero in the real interval $(0,1)$.
Under $M_H$ the only possible real zero in that interval is $1/2$.
Thus its interval condition is exactly equivalent to the displayed
single-value condition. Once holomorphy is available, the Brauer
zero-location argument of MT-ART07 proves GRH for every factor.
This is the full proof of both transfers. The one-dimensional factors
are Dirichlet functions and already fall under $A_D$.

The totally real quadratic field $E$ is the base for the odd projective
two-dimensional representations used in the cited theorem. Accordingly
this is a higher-rank consequence of a specific automorphic descent
theorem. It is not obtained by treating every $S_5$ factor as a single
Hecke character. In the Artin-motive programme, the remaining gate in
this family is one central Dedekind value. AG-SIGN controls a nonzero
leading coefficient after a central zero is removed; it does not prove
that the central order is zero.

\subsection*{MT-S503. One unconditional holomorphy example becomes a full GRH example}
Take the splitting field $L$ of
\[
                        x^5-x^3-x^2+x+1.
\]
Calegari, Theorem 5.1, publisher-author PDF p.~15, proves automorphy of
every representation of $\operatorname{Gal}(L/\mathbb Q)=S_5$;
the quintic field has discriminant $1609$. Consequently under $M_H$
every nontrivial irreducible Artin L-function of this explicit field
satisfies GRH, by the same signed Brauer argument. This example uses
that cited theorem, not the local-at-$5$ hypotheses of MT-S501, which
the source explicitly says this polynomial does not satisfy.
The historical bounded-height computation in Theorem 5.3 is part of
the classical proof's evidence; it is not a computation performed in
this run. Geometrically this supplies a concrete finite \(S_5\)-cover
of the arithmetic base whose complete Artin category is covered.

\subsection*{MT-M05. The stated Schauder quotient loses more than its zero evaluations see}
Lemma 9.3.2.5 on original PDF pp.~425--426 asserts a basis for the
quotient by the actual convolution ideal $\varphi*\mathcal S_E(G)$,
with coefficient functionals the Mellin jets at the zeros of
$\mathcal M\varphi$. We inspected the original p.~426 image: there
is no closure over this ideal. The assertion is false even when the
Mellin divisor consists of one simple zero. Thus the counterexample
does not rely on the zero-count defect in MT-M01.

Use $G=\mathbb R_{>0}$, trivial character and $u=\log t$; convolution
becomes ordinary additive convolution. Write
\[
 g_a(u)=(4\pi a)^{-1/2}e^{-u^2/(4a)},\qquad a>0.
\]
Choose $0<b<a$ and set $\varphi(e^u)=-g_a'(u)$,
$h(e^u)=-g_b'(u)$. Both functions belong to the exponential-Schwartz
space, and
\[
 \mathcal M\varphi(s)=s e^{as^2},\qquad
 \mathcal Mh(s)=s e^{bs^2}.
\]
The first function has just the simple zero $0$, lying in the interior
of the strip for $I=[-1,1]$. The only claimed coefficient functional is
$\delta_0$, and $\delta_0(h)=0$.
But $h\notin\varphi*\mathcal S_E$: if $h=\varphi*f$, multiplicativity
of the absolutely convergent Mellin transform gives
$\mathcal Mf(s)=e^{(b-a)s^2}$ for $s\ne0$, and then everywhere by
analyticity. On the imaginary axis this has modulus
$e^{(a-b)\tau^2}$. In contrast, $f(e^u)$ is integrable and therefore
$|\mathcal Mf(i\tau)|\le\int|f(e^u)|\,du$, a contradiction.
Thus the nonzero quotient class of $h$ has every asserted coefficient
zero, which is incompatible with uniqueness in a Schauder expansion.
The still simpler choice $\varphi=g_a$, $h=g_b$ has no Mellin zeros
and a nonzero quotient, so even the empty-divisor case fails.

\subsection*{MT-M06. The closed-ideal repair for every polynomial Gaussian divisor}
In logarithmic coordinates put
\[
 \mathcal E=\{f\in C^\infty(\mathbb R,\mathbb C):
   p_{N,j}(f):=\sup_x e^{N|x|}|f^{(j)}(x)|<\infty
   \text{ for every integer }N,j\ge0\}.
\]
This is $\mathcal S_E(\mathbb R_{>0})$ with its exponential-Schwartz
Fr\'echet topology. Polynomial factors in the usual Schwartz seminorms
are absorbed by increasing $N$. Let
\[
 P(s)=\prod_{\nu=1}^r(s-\rho_\nu)^{m_\nu},\quad
 d=\sum_\nu m_\nu,\qquad
 \varphi=P(-D)g_a,\quad D=d/dx.
\]
Distinct $\rho_\nu\in\mathbb C$, positive integral $m_\nu$, and
$a>0$ are arbitrary; $P=1,d=0$ is allowed. Then
\[
 \overline{\varphi*\mathcal E}
   =J_P:=\{h\in\mathcal E:
          (\mathcal Mh)^{(j)}(\rho_\nu)=0,
          \ 0\le j<m_\nu\},                         \tag{MT-M6}
\]
and the quotient by this \emph{closed} ideal is topologically
$\mathbb C^d$, with exactly the displayed Mellin jets as coordinates.
In particular its coordinate basis is absolute. This repairs the
finite-divisor case after making the closure operation explicit.

\begin{proof}
First, $g_a*\mathcal E$ is dense in $\mathcal E$. We give the full
argument, since density is distinct from the false surjectivity in
MT-M05. Convolution is continuous on $\mathcal E$: the inequality
$e^{N|x|}\le e^{N|x-y|}e^{N|y|}$ bounds each differentiated
convolution by an exponential supremum seminorm times an exponential
$L^1$ seminorm; the latter is bounded by a constant times
$p_{N+1,0}$.

Suppose a continuous linear functional $T$ vanishes on
$g_a*\mathcal E$. For every real $y$ and $\epsilon>0$ the function
$g_{a+\epsilon}(x-y)=g_a*g_\epsilon(\,\cdot-y)$ lies in this range.
It converges to $g_a(x-y)$ in every $p_{N,j}$ as
$\epsilon\downarrow0$, by uniform Gaussian estimates with the
parameter in a compact subinterval of $(0,\infty)$. Hence
$T(g_a(\,\cdot-y))=0$. The function
\[
 B(z)=T\bigl(e^{-x^2/(4a)}e^{zx/(2a)}\bigr)
\]
is entire: on every compact set of $z$, all its proposed complex
derivatives converge in every $p_{N,j}$ by Gaussian domination.
Completing the square gives $B(y)=0$ for real $y$, so the identity
theorem gives $B=0$.

For any Schwartz function $\psi$, Fourier inversion now gives
\[
 T(e^{-x^2/(4a)}\psi(x))
  =\frac1{2\pi}\int_{\mathbb R}\widehat\psi(\xi)
        B(2ai\xi)\,d\xi=0.
\]
This passage under $T$ is justified in the Fr\'echet space itself:
the $p_{N,j}$ seminorm of
$e^{-x^2/(4a)}e^{i\xi x}$ is at most
$C_{N,j,a}(1+|\xi|)^j$, whereas $\widehat\psi$ is rapidly
decreasing. For $h\in C_c^\infty$, take
$\psi=e^{x^2/(4a)}h$, which is again compactly supported and smooth.
Then $T(h)=0$. Finally $C_c^\infty$ is dense in $\mathcal E$:
multiply by a smooth cutoff equal to one on $[-R,R]$ and supported
in $[-2R,2R]$; the product rule and the bound from $p_{N+1,j}$ make
every $p_{N,j}$ of the discarded tail tend to zero. Thus $T=0$.
Hahn--Banach separation proves the asserted density.

Second, the range of $P(-D)$ on $\mathcal E$ is exactly $J_P$.
Integration by parts gives $\mathcal M(-Df)(s)=s\mathcal Mf(s)$,
so inclusion in $J_P$ follows. For the converse, if
$\mathcal Mh(\rho)=0$, define
\[
 f(x)=e^{-\rho x}\int_{-\infty}^x e^{\rho t}h(t)\,dt
     =-e^{-\rho x}\int_x^\infty e^{\rho t}h(t)\,dt.
\]
The equality is precisely the vanishing of the Mellin value.
The first expression controls the negative tail and the second the
positive tail. For any prescribed $N$, using an exponential bound
on $h$ with exponent larger than $N+|\operatorname{Re}\rho|+1$
proves exponential decay of $f$ on both tails; the equation
$f'=h-\rho f$ then gives the same bounds for every derivative.
Therefore $f\in\mathcal E$, and $(D+\rho)f=h$.
Equivalently $-f$ solves $(-D-\rho)(-f)=h$, and its Mellin transform
is $\mathcal Mh(s)/(s-\rho)$. Repeating this division through all
roots, including multiplicities, proves $P(-D)\mathcal E=J_P$.

The Mellin jet functionals are continuous, so $J_P$ is closed.
Since $g_a*\mathcal E$ is dense and $P(-D)$ is continuous,
$P(-D)(g_a*\mathcal E)$ is dense in $P(-D)\mathcal E=J_P$.
Differentiation commutes with convolution, identifying that range
with $\varphi*\mathcal E$. This proves (MT-M6).

For $d=0$ the quotient is zero and no interpolation is required.
For $d>0$, it remains to verify that no coordinate is missing. Given arbitrary
$d$ target jets, ordinary polynomial Hermite interpolation constructs
a polynomial $Q$ of degree below $d$ for which
$Q(s)e^{bs^2}$ has exactly those jets, with any fixed $b>0$.
Indeed multiplication by the nowhere-zero entire function $e^{bs^2}$
is invertible on each finite jet space, so interpolate the resulting
jets of $Q$. The function $Q(-D)g_b$ belongs to $\mathcal E$ and
realizes them. Choosing the targets to be the coordinate vectors
gives a continuous linear section from the finite-dimensional space
$\mathbb C^d$ to $\mathcal E$. The jet map consequently induces
a continuous isomorphism of the quotient with $\mathbb C^d$,
with continuous inverse. Finite-dimensional coordinate expansions are
absolute, completing the proof.
\end{proof}

For $P=s$, MT-M05 therefore exhibits a proper, nonclosed convolution
ideal whose closure is $\ker\delta_0$. The closed quotient really
does have the one coordinate the draft intended. For $P=1$ the
principal ideal is dense but proper, and the closed quotient is zero.
This is a concrete repair rather than a certification of the general
infinite-divisor statement. For arithmetic Hecke divisors the remaining
gate is controlled infinite interpolation and the claimed absolute
basis/nuclear estimates on the correct separated quotient. A local
$O(\log T)$ zero count alone does not prove those estimates.
In cote~1 trace language, Hausdorff completion and the identification
of the spectral quotient precede taking its nuclear trace.

\subsection*{MT-M07. An infinite closed quotient with an absolute Mellin basis}
The closed-ideal repair also admits a complete infinite-divisor example.
Fix $a,\ell>0$ in the logarithmic space $\mathcal E$ of MT-M06, and put
\[
 \varphi(x)=g_a(x-\ell)-g_a(x),\qquad
 \rho_n=\frac{2\pi i n}{\ell}\quad(n\in\mathbb Z).
\]
Then $\mathcal M\varphi(s)=e^{as^2}(e^{\ell s}-1)$, with precisely
the simple zeros $\rho_n$. Let
\[
 (\mathcal P_\ell f)(t)=\sum_{k\in\mathbb Z}f(t+k\ell).
\]
The quotient $\mathcal E/\overline{\varphi*\mathcal E}$ is
topologically $C^\infty(\mathbb R/\ell\mathbb Z)$ under
$\mathcal P_\ell$, and has an absolute Schauder basis whose coefficient
functionals are exactly $f\mapsto\mathcal Mf(\rho_n)$.

\begin{proof}
The periodization series and every derivative converge uniformly on a
fundamental interval, by exponential decay. This gives a continuous map
to the indicated smooth periodic space. Choose
$\chi\in C_c^\infty(\mathbb R)$ nonnegative and positive on
$[-\ell,\ell]$, and set
$\theta=\chi/\mathcal P_\ell\chi$. The denominator is positive,
smooth and periodic, so $\theta$ is smooth, compactly supported and
$\mathcal P_\ell\theta=1$. The map
$R(F)(x)=\theta(x)F(x)$ is a continuous section of periodization,
proving surjectivity.

Let $\Delta_\ell h(x)=h(x-\ell)-h(x)$. Telescoping shows that its
range lies in $\ker\mathcal P_\ell$. Conversely, if
$\mathcal P_\ell f=0$, define
\[
 h(x)=\sum_{k\ge1}f(x+k\ell)
     =-\sum_{k\le0}f(x+k\ell).
\]
These locally normally convergent series define a smooth function.
For $x\ge0$ use the first expression and for $x\le0$ the second.
For every positive integer $N$ and every $j$, summing the geometric
series of bounds $p_{N,j}(f)e^{-N|x+k\ell|}$ proves
$p_{N,j}(h)<\infty$; the $N=0$ case follows from $N=1$.
Thus $h\in\mathcal E$, and direct index shifting gives
$\Delta_\ell h=f$. The range of $\Delta_\ell$ is therefore the
closed space $\ker\mathcal P_\ell$.

By MT-M06, $g_a*\mathcal E$ is dense in $\mathcal E$.
Continuity of $\Delta_\ell$ and commutation with convolution imply
\[
 \overline{\varphi*\mathcal E}
 =\overline{\Delta_\ell(g_a*\mathcal E)}
 =\ker\mathcal P_\ell.
\]
The continuous section $R$ identifies the quotient topology with
the usual smooth periodic topology, rather than merely giving a
vector-space identification.

The Fourier expansion in the convention used here is
\[
 \mathcal P_\ell f(t)=\frac1\ell
       \sum_{n\in\mathbb Z}\mathcal Mf(\rho_n)e^{-\rho_n t}.
\]
Its coefficients are obtained by integrating over a fundamental
interval and summing translates, justified by absolute convergence.
They decrease faster than every power of $|n|$: integrate by parts
arbitrarily many times and use the integrability of every derivative
of $f$. The Fourier series consequently converges absolutely in
every smooth seminorm and equals $\mathcal P_\ell f$ by Fourier
uniqueness. Representatives for the basis are
\[
                     b_n(x)=\ell^{-1}\theta(x)e^{-\rho_n x}.
\]
Indeed $\mathcal Mb_n(\rho_m)=\delta_{m,n}$, by
$\mathcal P_\ell\theta=1$ and integration on one period. Moreover,
$p_{N,j}(b_n)\le C_{N,j}(1+|n|)^j$ since $\theta$ has compact
support. Hence the expansion
$\sum_n\mathcal Mf(\rho_n)b_n$ is absolutely convergent even in
$\mathcal E$, and represents the class of $f$ in the quotient.
The jets determine that class uniquely by the Fourier identification.
Ordering the indices as $0,1,-1,2,-2,\ldots$ gives an absolute
Schauder basis, as asserted.
\end{proof}

The actual ideal still cannot replace its closure: with $0<b<a$,
$h=\Delta_\ell g_b$ has all these Mellin values zero but is not
$\varphi*f$ for any $f\in\mathcal E$, by exactly the Fourier-growth
division argument of MT-M05. Thus this example tests an infinite
divisor with controlled spacing, not just the empty or finite case.

\subsection*{MT-M08. The corresponding trace and positive convolution formula}
For $f\in\mathcal E$, convolution by $f$ descends to the quotient of
MT-M07. On $L^2(\mathbb R/\ell\mathbb Z,dt)$ it is trace class and
\[
 \operatorname{Tr}C_f
   =\sum_{n\in\mathbb Z}\mathcal Mf(\rho_n)
   =\ell\sum_{k\in\mathbb Z}f(k\ell).                 \tag{MT-M8}
\]
All sums converge absolutely. If
$f=h*\widetilde h$ with $h\in\mathcal E$ and
$\widetilde h(x)=\overline{h(-x)}$, this operator is positive and
\[
                  \operatorname{Tr}C_f
                   =\sum_n|\mathcal Mh(\rho_n)|^2\ge0.
\]

To prove the assertions, first use absolute convergence and changes of
variables to periodize a convolution:
\[
 \mathcal P_\ell(f*g)(t)
  =\int_0^\ell(\mathcal P_\ell f)(t-u)
                     (\mathcal P_\ell g)(u)\,du.
\]
This proves descent and identifies the smooth kernel
$K_f(t,u)=(\mathcal P_\ell f)(t-u)$.
The orthonormal Fourier vectors
$\ell^{-1/2}e^{-\rho_n t}$ are eigenvectors with eigenvalues
$\mathcal Mf(\rho_n)$. These eigenvalues decrease faster than every
power, as proved in MT-M07, so their absolute sum is finite.
The diagonal expansion is thus a trace-class expansion and gives
the first trace sum. Evaluating the absolutely convergent Fourier
series of $\mathcal P_\ell f$ at zero gives the second equality
in (MT-M8), including its factor $\ell$. Finally the Mellin transform
of $\widetilde h$ at a purely imaginary argument is the conjugate
of that of $h$; the eigenvalues are therefore the displayed squares.

There is also a direct nuclear realization on the smooth quotient.
Explicitly, put $c_n(F)=\ell^{-1}\int_0^\ell F(t)e^{\rho_nt}\,dt$
and $e_n(t)=e^{-\rho_nt}$. The $c_n$ are an equicontinuous family,
bounded by the uniform norm. In
\[
 C_f=\sum_n(1+|n|)^{-2}\,
       c_n\otimes\bigl((1+|n|)^2\mathcal Mf(\rho_n)e_n\bigr),
\]
the scalars are summable and the vector family is bounded in every
$C^\infty$ seminorm. This is a nuclear decomposition. Fourier partial
sums converge uniformly on bounded sets of $C^\infty$ (apply a
uniform higher-derivative bound to the coefficient tail), so the
finite-rank trace limit is the same absolute sum above.

This supplies a fully checked spectral-quotient, basis and trace
model beside cote~1's nuclear trace language. It does not identify
an arithmetic prime distribution with this lattice or prove the
corresponding basis theorem for the actual Hecke zero set. In that
case the exact infinite interpolation and operator estimates remain
to be supplied on the correct closed quotient.

\subsection*{MT-M09. Even a bounded local zero count does not give the asserted point basis}
Passing to the closed ideal is necessary, but does not prove the
general basis claim. There exists $\varphi\in\mathcal E$ whose Mellin
transform has only simple zeros on the imaginary axis and only
$O(1)$ zeros in every unit-height interval, for which
\[
        Q=\mathcal E/\overline{\varphi*\mathcal E}
\]
has no Schauder basis with the individual zero evaluations as its
coefficient functionals. This holds for every enumeration of the
zeros, and therefore also rules out an absolute Schauder basis of
that specified kind.

\begin{proof}
Put $\epsilon_n=e^{-n^4}$ for $n\ge1$ and define
\[
 H(s)=\prod_{n=1}^{\infty}
       \left(1+\frac{s^2}{n^2}\right)
       \left(1+\frac{s^2}{(n+\epsilon_n)^2}\right),
 \qquad F(s)=e^{as^2}H(s),\quad a>0.
\]
The product converges uniformly on compact sets because the sum of
the absolute deviations of its factors from one does. Away from the
listed zeros the tail logarithms converge locally uniformly, so
there are no additional zeros. Its zeros are exactly
$\pm in,\ \pm i(n+\epsilon_n)$, all simple. Since
$0<\epsilon_n<1/2$, each unit-height interval contains a bounded
number of them.

For $r=|s|$ the product satisfies
\[
 \log|H(s)|\le2\sum_{n\ge1}\log(1+r^2/n^2)\le2\pi r.
\]
For the last bound compare the decreasing function of $x$ with its
integral on $(0,\infty)$; integration by parts gives
$\int_0^\infty\log(1+r^2/x^2)\,dx=\pi r$.
Consequently, on every strip $|\sigma|\le A$,
\[
 |F(\sigma+it)|
 \le \exp\{aA^2-at^2+2\pi\sqrt{A^2+t^2}\}.
\]
Define
\[
 \varphi(x)=\frac1{2\pi}\int_{\mathbb R}F(it)e^{-itx}\,dt.
\]
The Gaussian bound permits every derivative in $x$ under the integral.
Shifting the contour to $\operatorname{Re}s=\sigma$ gives, for
every real $\sigma$ and every $j\ge0$,
\[
 \varphi^{(j)}(x)
   =\frac{e^{-\sigma x}}{2\pi}
      \int_{\mathbb R}(-\sigma-it)^jF(\sigma+it)e^{-itx}\,dt.
\]
The horizontal sides of the contour tend to zero by the displayed
Gaussian estimate, uniformly for $\sigma$ in the finite interval
used in that shift. Choosing $\sigma=N+1$ for $x\ge0$ and
$\sigma=-N-1$ for $x\le0$ proves every seminorm $p_{N,j}$ is finite.
Thus $\varphi\in\mathcal E$. Cauchy estimates applied on slightly
larger strips show that $F(it)$ and all its derivatives are Schwartz.
Fourier inversion gives $\mathcal M\varphi(it)=F(it)$;
both sides are entire, so $\mathcal M\varphi=F$ everywhere.

Every evaluation at a listed zero annihilates the convolution ideal
and its closure, and therefore descends continuously to $Q$.
Let $p=p_{1,0}$ and let
$\overline p([f])=\inf_{j\in\overline{\varphi*\mathcal E}}p(f+j)$
be the corresponding quotient seminorm. For any $f\in\mathcal E$,
\[
 |(\delta_{in}-\delta_{i(n+\epsilon_n)})(f)|
  \le\epsilon_n\int_{\mathbb R}|x||f(x)|\,dx
  \le2\epsilon_n p(f).
\]
The first inequality follows from
$|e^{i\epsilon_n x}-1|\le\epsilon_n|x|$.
Because the difference functional annihilates the closed ideal,
the same bound holds with $p$ replaced by $\overline p$ on $Q$.

Suppose the claimed Schauder basis existed, and denote the basis
vector associated to $\delta_{in}$ by $b_n$. Biorthogonality of
basis vectors and coefficient functionals gives
$\delta_{in}(b_n)=1$ and
$\delta_{i(n+\epsilon_n)}(b_n)=0$. Hence
\[
                 \overline p(b_n)\ge(2\epsilon_n)^{-1}.
\]
Take the single test function $g_b\in\mathcal E$ for any fixed
$b>0$. Its coefficient at $\delta_{in}$ is $e^{-bn^2}$.
The corresponding terms of its proposed basis expansion satisfy
\[
 \overline p\bigl(e^{-bn^2}b_n\bigr)
        \ge\frac12 e^{n^4-bn^2}\longrightarrow\infty.
\]
Terms of a convergent series in a locally convex space must tend
to zero. In every enumeration these infinitely many distinct terms
occur arbitrarily far along the series, giving a contradiction.
This proves the assertion.
\end{proof}

This is a counterexample for general exponential-Schwartz Mellin
divisors, not a claim that such clustered zeros occur for an actual
Hecke L-function. It pinpoints why NT-Z10/MT-M04 cannot by themselves
establish the arithmetic Schauder step: even an $O(1)$ local count
does not control the norm of individual interpolation vectors.
The next possible repair must incorporate quantitative separation,
or group nearby evaluations into divided-difference blocks and prove
convergence of the resulting block expansion. Neither property is
supplied by a zero-location hypothesis alone in the present argument.

\subsection*{MT-M10. A point-jet basis forces a polynomial cluster-product bound}
The counterexample MT-M09 has a quantitative converse obstruction.
Let $\mathcal E$ be the exponential-Schwartz space of MT-M06,
let $J\subset\mathcal E$ be a closed linear subspace, and put
$Q=\mathcal E/J$. Fix a real $\sigma$ and a discrete set
$Z\subset\sigma+i\mathbb R$, with positive integer multiplicities
$m_\rho$. Suppose every functional
\[
 \delta_{\rho,k}(f)=\frac{(\mathcal Mf)^{(k)}(\rho)}{k!},
 \qquad 0\le k<m_\rho,
\]
annihilates $J$ and occurs as a coefficient functional of a
Schauder basis of $Q$ with continuous coordinate functionals.
Other basis coordinates are permitted.
Then there are constants $c>0$ and $A\ge0$, independent of
$\rho=\sigma+it$ and of the finite subset
$W\subset Z\setminus\{\rho\}$, such that
\[
       \prod_{\nu\in W}|t-\operatorname{Im}\nu|^{m_\nu}
                    \ge c(1+|t|)^{-A}.                \tag{MT-M10}
\]
In particular the distinct ordinates have a uniform polynomial
separation bound. The conclusion already follows from an ordinary
Schauder basis; absolute convergence is not needed.

\begin{proof}
The seminorms $p_{N,j}(f)=\sup_x e^{N|x|}|f^{(j)}(x)|$ make
$\mathcal E$ a Fr\'echet space: a Cauchy sequence converges uniformly,
with all derivatives, on every compact interval, and the weighted
uniform limits give the same smooth limit and all the seminorm
limits. The closed quotient $Q$ is Fr\'echet as well.
For a continuous seminorm $p$ on $\mathcal E$, write
\[
 \overline p([f])=\inf_{u\in J}p(f+u).
\]
The quotient seminorms of finite maxima of the $p_{N,j}$ generate
the quotient topology.

First fix any such $\overline p$. Let $S_n$ be the partial-sum
operators for the proposed basis, and let $R_n=S_n-S_{n-1}$ be
its rank-one coordinate projections, with $S_0=0$.
Convergence of every basis expansion makes
$\sup_n\overline p(R_nq)$ finite for each $q\in Q$.
For completeness, the uniform-boundedness step follows directly
from Baire's theorem. The closed sets
\[
 A_k=\{q:\sup_n\overline p(R_nq)\le k\},\qquad k=1,2,\ldots,
\]
cover $Q$. One contains $q_0+U$, where $U$ is a neighborhood of
zero, and $q_0$ itself may be chosen in that set. Subtraction
gives $\overline p(R_nu)\le2k$ for $u\in U$, uniformly in $n$.
Choose $U$ to contain a unit ball for a finite maximum of
quotient seminorms and rescale. This proves that for a continuous
quotient seminorm $\overline q$ and a constant $C$,
\[
             \overline p(R_nq)\le C\overline q(q)
                         \quad\hbox{for every }n,q.    \tag{1}
\]
The same conclusion holds when the seminorm on the right is
zero, by applying the bound to arbitrary scalar multiples.

Let $b_{\rho,0}$ be the basis vector with coefficient
$\delta_{\rho,0}$. For a fixed Gaussian $g_a$ with
$\mathcal Mg_a(s)=e^{as^2}$, use the test function
\[
                f_t(x)=e^{-\sigma x-itx}g_a(x).
\]
It satisfies $\delta_{\rho,0}(f_t)=1$.
Every fixed seminorm of $f_t$ is bounded by a fixed power of
$1+|t|$, by the Leibniz formula and Gaussian decay.
Apply (1) to $[f_t]$ and its coordinate projection. For each
fixed $\overline p$ there are constants $C_p,A_p$ such that
\[
       \overline p(b_{\rho,0})\le C_p(1+|t|)^{A_p}
                                      \quad(\rho\in Z). \tag{2}
\]
The constants do not depend on the enumeration position of
$b_{\rho,0}$.

We next give a uniform estimate for confluent divided differences.
For real nodes $t_0,\ldots,t_m$, repeated nodes allowed, the
divided difference of $u\mapsto e^{ixu}$ is
\[
 [t_0,\ldots,t_m]e^{ixu}
   =\int_{\Delta_m}(ix)^m
           e^{ix\sum_{j=0}^m\lambda_jt_j}\,d\lambda,
 \qquad
 \Delta_m=\{\lambda_j\ge0:\textstyle\sum\lambda_j=1\},
                                                        \tag{3}
\]
where $d\lambda$ has total mass $1/m!$.
One elementary verification suffices also for repeated nodes.
For distinct nodes, the divided difference of $u^k$ is zero
when $k<m$ and is the complete homogeneous polynomial
$h_{k-m}(t_0,\ldots,t_m)$ when $k\ge m$. This follows by
expanding the partial-fraction identity
\[
 \sum_{j=0}^m
 \frac{1}{(1-zt_j)\prod_{\ell\ne j}(t_j-t_\ell)}
          =\frac{z^m}{\prod_{j=0}^m(1-zt_j)}.
\]
On the simplex,
$\int_{\Delta_m}\lambda_0^{a_0}\cdots\lambda_m^{a_m}\,d\lambda
=\prod_j a_j!/(m+\sum a_j)!$; successive integration of
monomials proves this formula. Expanding the exponential and
using this integral gives exactly the same power series,
proving (3). Continuity as nodes coalesce gives the confluent
version, whose functionals are derivatives divided by factorials.
In particular the absolute value in (3) is at most $|x|^m/m!$.

Take $t_0=t$ and repeat each $\operatorname{Im}\nu$, $\nu\in W$,
exactly $m_\nu$ times among the other nodes. Apply the divided
difference to $u\mapsto\mathcal Mf(\sigma+iu)$.
All resulting point-jet functionals annihilate $J$.
Choose an integer $N\ge|\sigma|+1$ and put $p=p_{N,0}$.
Formula (3), Fubini's theorem and exponential decay give
\[
 \left|[t_0,\ldots,t_m]\mathcal Mf(\sigma+iu)\right|
 \le\int_{\mathbb R}\frac{|x|^m}{m!}e^{\sigma x}|f(x)|\,dx
 \le2p(f).                                             \tag{4}
\]
The last constant is independent of both the nodes and $m$,
since $\int_{\mathbb R}|x|^m e^{-|x|}\,dx=2m!$.
The same inequality holds on $Q$ with $\overline p$.

Biorthogonality says that on $b_{\rho,0}$ all the other
point jets in this divided difference vanish. The coefficient
of the simple node $t$ is
$\prod_{\nu\in W}(t-\operatorname{Im}\nu)^{-m_\nu}$.
For repeated other nodes this follows either from partial
fractions or by coalescing distinct nodes away from $t$.
Consequently (4) and (2) imply
\[
 \prod_{\nu\in W}|t-\operatorname{Im}\nu|^{-m_\nu}
       \le2\overline p(b_{\rho,0})
       \le2C_p(1+|t|)^{A_p},
\]
which is (MT-M10).
\end{proof}

For any fixed $0<r<1$, choosing all other zeros within distance
$r$ also bounds their total multiplicity by
$(A\log(1+|t|)-\log c)/|\log r|$.
The product condition is stronger than such a local count:
it controls the joint cost of arbitrarily close clusters.
Taking one other node recovers polynomial separation.
Taking the clusters in MT-M09 contradicts (MT-M10) directly.
This identifies an explicit necessary analytic input for
Masullo 9.3.2.5, PDF~425--426, even after the closed-ideal repair.
It neither asserts that arithmetic zeros violate that condition
nor proves that the condition suffices for a basis.

\subsection*{MT-M11. The same basis would bound arithmetic leading coefficients}
Fix a primitive unitary Hecke $L$-function $L(s,\eta)$ and assume
its nontrivial zeros lie on $\operatorname{Re}s=1/2$.
Suppose a closed quotient as in MT-M10 has all its nontrivial
zero jets as coefficient functionals, with their actual
multiplicities. Then there are constants $c,C,T_0>0$, depending
on the fixed function and the proposed basis, such that at every
nontrivial zero $\rho=1/2+iT$ with $|T|\ge T_0$ and order $m$,
\[
               \left|\frac{L^{(m)}(\rho,\eta)}{m!}\right|
                        \ge c(1+|T|)^{-C}.             \tag{MT-M11}
\]
In particular, at simple zeros the reciprocal of $|L'(\rho,\eta)|$
would have at most polynomial growth. The zero-location input is
$A_D$ for Dirichlet functions and $M_H$ for general Hecke functions;
the asserted quotient basis is an additional hypothesis.

\begin{proof}
We use the fixed-function analytic facts proved in NT-Z10:
polynomial growth on a fixed vertical strip away from its
finitely located poles, and $O(\log(2+|T|))$ zeros, with
multiplicity, in each bounded-height window.
The absolutely convergent Euler product gives
$|L(2+iT,\eta)|\ge c_\eta>0$: each local inverse factor
has modulus at least $(1+(\operatorname N\mathfrak p)^{-2})^{-1}$,
and the product of these lower bounds is positive.

Put $c_T=2+iT$. Choose $4\le R\le5$ so that the boundary of
$D=\{s:|s-c_T|<R\}$ contains no zero. At sufficiently large
$|T|$, it contains no pole or trivial zero either, because
the imaginary parts of their finitely many arithmetic
progressions parallel to the real axis are fixed by the infinity type.
Its zeros are therefore precisely the relevant critical-line
zeros in that disk, say $\rho_j$ with multiplicities $m_j$.
There are
$M=\sum_jm_j=O_\eta(\log(2+|T|))$ of them.
In the normalized coordinate $z=(s-c_T)/R$, put
$a_j=(\rho_j-c_T)/R$ and form the finite Blaschke product
\[
 B(s)=\prod_j
       \left(\frac{z-a_j}{1-\overline{a_j}z}\right)^{m_j}.
\]
Constant unimodular normalizations of the factors are
irrelevant. The function $G=L/B$ is holomorphic and zero-free
on the disk and its boundary. On the boundary $|B|=1$.
The polynomial growth bound and the maximum principle give
$|G(s)|\le M_T=C_1(1+|T|)^{C_2}$ throughout $D$, enlarging
$M_T$ if necessary to make it at least one.
At the center, $|B(c_T)|\le1$, so
$|G(c_T)|\ge c_\eta$.

The nonnegative harmonic function
$u(s)=\log M_T-\log|G(s)|$ satisfies the disk Harnack bound
\[
 u(\rho)\le
 \frac{R+|\rho-c_T|}{R-|\rho-c_T|}\,u(c_T)
 \le\frac{11}{5}\,u(c_T),
\]
because $|\rho-c_T|=3/2$ and $R\ge4$.
This elementary bound follows by comparing the maximum and
minimum of the Poisson kernel on each concentric boundary
circle and then letting that circle increase to radius $R$.
It implies $|G(\rho)|\ge c_2(1+|T|)^{-C_3}$.

Let $a=(\rho-c_T)/R$, so $|a|\le3/8$.
The leading coefficient of its own factor in $B$ contributes
$[R(1-|a|^2)]^{-m}\ge5^{-m}$.
Every other factor at $\rho$ has modulus
\[
 \frac{|a-a_j|}{|1-\overline{a_j}a|}
 \ge\frac{|\rho-\rho_j|}{10}.
\]
Thus, with $W$ the distinct other zeros in $D$,
\[
 \left|\lim_{s\to\rho}\frac{B(s)}{(s-\rho)^m}\right|
 \ge 10^{-M}\prod_{\rho_j\in W}|\rho-\rho_j|^{m_j}.
\]
All these points are on the same vertical line. Apply
(MT-M10) to this finite set $W$. Since
$10^{-M}\ge c_3(1+|T|)^{-C_4}$, the displayed expression
has a polynomial lower bound. Multiplying by the lower bound
for $|G(\rho)|$ proves (MT-M11).
\end{proof}

This is a necessary consequence of the proposed arithmetic
point-jet basis, not a derivative estimate obtained from GRH
alone. It converts the remaining interpolation problem into
a concrete additional test on arithmetic leading coefficients.
The work placement is the Masullo 9.3.2.5 repair, NT-Z10,
and the cote~1 trace/Fredholm programme reported in SA23
geo-H-lidskii/geo-H-genus. Positivity of the trace cannot
replace this topological prerequisite.

\subsection*{MT-M12. An exact simple-zero gate for the canonical arithmetic quotient}
Fix a primitive unitary Hecke character $\eta$ over a number field
$K$, and assume GRH for its $L$-function. Assume additionally that
all its nontrivial zeros are simple. Let
\[
 \Xi_\eta(s)=P_\eta(s)\Lambda(s,\eta)
\]
be the entire completed function with its possible poles removed.
For a non-pure-norm character take $P_\eta=1$. If
$L(s,\eta)=\zeta_K(s+i\tau)$, take
$P_\eta(s)=(s+i\tau)(s-1+i\tau)$.
The zeros of $\Xi_\eta$ are precisely the nontrivial zeros, denoted
$Z$. Put
\[
 \psi_\eta(x)=\frac1{2\pi}\int_{\mathbb R}\Xi_\eta(it)e^{-itx}\,dt,
 \qquad Q_\eta=\mathcal E/\overline{\psi_\eta*\mathcal E}.
\]
Then $\psi_\eta\in\mathcal E$, with
$\mathcal M\psi_\eta=\Xi_\eta$, and the following are equivalent:
\begin{enumerate}
\item There are $c,C>0$ such that
$|L'(\rho,\eta)|\ge c(1+|\operatorname{Im}\rho|)^{-C}$ for
every $\rho\in Z$.
\item $Q_\eta$ has an absolute Schauder basis with coefficient
functionals $[f]\mapsto\mathcal Mf(\rho)$, $\rho\in Z$.
\item $Q_\eta$ has a Schauder basis with these continuous
coefficient functionals.
\end{enumerate}
Under these conditions the closed convolution ideal is exactly
\[
 K_Z=\{f\in\mathcal E:\mathcal Mf(\rho)=0
                              \text{ for every }\rho\in Z\}. \tag{MT-M12a}
\]
Moreover $Q_\eta$ is topologically the rapidly decreasing sequence
space
\[
 s(Z)=\{(a_\rho):\sup_{\rho\in Z}
       (1+|\operatorname{Im}\rho|)^k|a_\rho|<\infty
                                  \text{ for every }k\ge0\}. \tag{MT-M12b}
\]
This equivalence uses $A_D$ for Dirichlet functions, or $M_H$
conditionally for general unitary Hecke functions. Simplicity is
an extra assumption throughout; condition (1) is the remaining
gate, not a consequence of those GRH inputs proved here.

\begin{proof}
\emph{Analytic preliminaries.}
Write $n=[K:\mathbb Q]$ and $\kappa=\pi n/4$.
The fixed-function polynomial strip bounds of NT-Z10 and
Stirling's formula for the $n$ factors $\Gamma_{\mathbb R}$
give, for every fixed $A>0$, constants $C_A,B_A$ such that
\[
 |\Xi_\eta(\sigma+it)|
       \le C_A(1+|t|)^{B_A}e^{-\kappa|t|}
                            \qquad(|\sigma|\le A).     \tag{5}
\]
On bounded $t$ the entire function is bounded on each compact
strip rectangle, so enlarging $C_A$ covers possible gamma-pole
cancellations there. Derivative bounds follow by Cauchy's formula
on a slightly larger strip. Thus $\Xi_\eta(it)$ is Schwartz.
The contour-shifting and Fourier-inversion argument of MT-M09,
using (5) in place of the Gaussian bound, shows that the displayed
inverse transform belongs to $\mathcal E$ and has Mellin transform
$\Xi_\eta$. More explicitly, shift to $\sigma=N+1$ on $x\ge0$
and to $\sigma=-N-1$ on $x\le0$; every integral with the extra
factor $(-\sigma-it)^j$ is absolutely convergent by (5).
This proves all $p_{N,j}$ bounds. The horizontal contour edges
tend to zero by the same estimate, and Fourier uniqueness then
proves the transform identity.

Assume condition (1). At a zero $\rho=1/2+iT$, differentiation
of the completed function gives
\[
 \Xi_\eta'(\rho)
  =P_\eta(\rho)q_\eta^{\rho/2}
       \prod_j\Gamma_{\mathbb R}(\rho+\mu_j)L'(\rho,\eta),
\]
where $q_\eta=|D_K|\operatorname N\mathfrak f_\eta$
is the finite conductor.
Stirling's formula and (1) imply, after enlarging constants to
cover the finitely many low zeros,
\[
        |\Xi_\eta'(\rho)|
           \ge c_1(1+|T|)^{-C_1}e^{-\kappa|T|}.         \tag{6}
\]

\emph{Polynomially bounded cardinal functions.}
Fix $a>0$. For each zero $\rho$ define the entire function
\[
 H_\rho(s)=
     \frac{\Xi_\eta(s)}
          {\Xi_\eta'(\rho)(s-\rho)}\,e^{a(s-\rho)^2}.    \tag{7}
\]
It takes the value one at $\rho$ and zero at every other zero.
For each fixed strip $|\operatorname{Re}s|\le A$, (5)--(6)
imply a bound
\[
 |H_\rho(\sigma+it)|
       \le C_A'(1+|T|)^{B_A'}
                              e^{-a(t-T)^2/2}.         \tag{8}
\]
Here are the details near the apparent denominator. If
$|s-\rho|\ge1$, use (5) and (6), together with
$\big||T|-|t|\big|\le|T-t|$ and
$1+|t|\le(1+|T|)(1+|t-T|)$.
The remaining fixed polynomial and
$e^{\kappa|t-T|}$ are absorbed by half the Gaussian.
If $|s-\rho|<1$, use
\[
 \frac{\Xi_\eta(s)}{s-\rho}
       =\int_0^1\Xi_\eta'(\rho+u(s-\rho))\,du
\]
and the Cauchy derivative bound from (5) on a slightly larger
strip. Since $|t-T|<1$, the same estimate (8) follows by
adjusting the constants.

Invert $H_\rho$ by Fourier transform and call the inverse
$h_\rho$. The same contour argument gives $h_\rho\in\mathcal E$,
$\mathcal Mh_\rho=H_\rho$, and, for each $N,j$,
\[
             p_{N,j}(h_\rho)
                    \le C_{N,j}(1+|T|)^{B_{N,j}}.      \tag{9}
\]
Indeed, on the shifted contour used for $x\ge0$ or $x\le0$,
integrating $(1+|t|)^j$ against the Gaussian in (8) costs
only another fixed power of $1+|T|$.

For any $f\in\mathcal E$, its values $\mathcal Mf(1/2+iT)$
decrease faster than every power of $|T|$. This follows by
integrating by parts arbitrarily many times in
$\int e^{x/2}f(x)e^{iTx}\,dx$; each resulting $L^1$ norm
is bounded by finitely many defining seminorms of $f$.
The local count NT-Z10 implies
$\sum_{\rho\in Z}(1+|\operatorname{Im}\rho|)^{-2}<\infty$.
Consequently
\[
              I(f)=\sum_{\rho\in Z}\mathcal Mf(\rho)h_\rho
                                                               \tag{10}
\]
converges absolutely in every seminorm of $\mathcal E$,
defines a continuous linear map, and has exactly the same
values as $f$ at all the zeros. Likewise, for every
$(a_\rho)\in s(Z)$ the series $\sum_\rho a_\rho h_\rho$
converges in $\mathcal E$ and depends continuously on that
sequence. These assertions use (9) and choose the coefficient
decay power separately for each target seminorm.

\emph{Division on vertical strips.}
We now show the ideal equality, which does not follow merely
from constructing the cardinal functions.
Condition (1) implies a uniform, fixed-strip minimum-modulus
estimate
\[
 |L(\sigma+it,\eta)|
   \ge c_A(1+|t|)^{-C_A}
            \min\{1,\operatorname{dist}(\sigma+it,Z)\},
 \qquad |\sigma|\le A,\quad |t|\ge T_A.                 \tag{11}
\]
Here is a complete proof. Put $r=A+2$ and center a disk at
$c_t=2+it$. Choose $R\in[4r+4,4r+5]$ with zero-free boundary.
At large $|t|$ there are no poles or trivial zeros in that disk.
Its $M=O_A(\log(2+|t|))$ zeros are all simple and on the
critical line. Remove them by the finite Blaschke product
$B$ for this disk, as in MT-M11, and put $G=L/B$.
Polynomial strip growth gives $|G|\le C(1+|t|)^C$ on the
whole disk, while the Euler lower bound at $c_t$ and Harnack's
inequality give
$|G|\ge c(1+|t|)^{-C}$ on the disk of radius $R/2$.
The Harnack factor on that smaller disk is at most three.
The point $s=\sigma+it$ is within radius $r<R/4$.

If the disk has no zeros, this already proves (11).
Otherwise choose a zero $\rho_0$ in it nearest to $s$.
If $|\rho_0-c_t|>R/2$, every zero in the disk is at
distance at least $R/4$ from $s$. Each Blaschke factor
therefore has modulus at least $1/8$, and
$|B(s)|\ge8^{-M}$, a polynomial lower bound.
If $|\rho_0-c_t|\le R/2$, simplicity gives
$L'(\rho_0)=G(\rho_0)B'(\rho_0)$.
In normalized coordinates $|a_0|\le1/2$, the derivative
of its own Blaschke factor is at most $4/(3R)$, and
every other factor at $\rho_0$ is at most
$2|\rho_0-\rho_j|/R$. Thus
\[
 |B'(\rho_0)|\le C_A^M
                     \prod_{\rho_j\ne\rho_0}|\rho_0-\rho_j|.
\]
The assumed lower bound for $L'$, the upper bound for $G$,
and $M=O_A(\log(2+|t|))$ imply
\[
          \prod_{\rho_j\ne\rho_0}|\rho_0-\rho_j|
                          \ge c(1+|t|)^{-C}.           \tag{12}
\]
Since $\rho_0$ is nearest to $s$,
$|\rho_0-\rho_j|\le2|s-\rho_j|$. Each denominator in a
Blaschke factor at $s$ is at most two. Hence
\[
 |B(s)|\ge(4R)^{-M}|s-\rho_0|
                       \prod_{\rho_j\ne\rho_0}|\rho_0-\rho_j|,
\]
and (12) proves (11). Zeros outside the large disk have
distance greater than one from $s$, so using the distance
to the full set $Z$ makes no difference in the minimum.
The estimate at a zero is interpreted as zero on both
sides and follows by continuity.

\emph{The closed ideal is the full zero kernel.}
The inclusion
$\overline{\psi_\eta*\mathcal E}\subset K_Z$ is immediate
from multiplication of Mellin transforms.
Take $f\in K_Z$ and write $F=\mathcal Mf$.
The quotient $F/\Xi_\eta$ is entire because all zeros
are simple and have been cancelled. On every fixed strip,
(11) and the rapid vertical decay of $F$ and $F'$ imply
rapid decay of $F/L$ for large $|t|$. If the distance to
$Z$ is at least one, divide the rapid bound for $F$
by (11). If it is less than one, choose a nearest zero
$\rho$ and use
$F(s)=\int_\rho^s F'(w)\,dw$ on the straight segment;
that segment lies in a slightly larger fixed strip and
a bounded-height window. The distance factor cancels
the corresponding factor in (11).

Inverting the gamma factors now costs at most
$C_A(1+|t|)^{C_A}e^{\kappa|t|}$ on the strip.
For any $\epsilon>0$ the entire function
\[
             J_\epsilon(s)=e^{\epsilon s^2}
                                  \frac{F(s)}{\Xi_\eta(s)}
\]
therefore decays faster than every power on every vertical
strip, in fact with a Gaussian bound. On compact rectangles
it is bounded by entireness. Contour shifting and Fourier
inversion give $j_\epsilon\in\mathcal E$ with this Mellin
transform. Uniqueness of the transform gives
\[
                  \psi_\eta*j_\epsilon=g_\epsilon*f.
\]
Finally $g_\epsilon*f\to f$ in $\mathcal E$ as
$\epsilon\downarrow0$. To verify this in each $p_{N,j}$,
write the difference as the integral of
$g_\epsilon(y)(f^{(j)}(x-y)-f^{(j)}(x))$.
For $|y|\le1$, the weighted mean-value bound is
$|y|e^{N|y|}p_{N,j+1}(f)$; outside any fixed
neighborhood of zero, use
$(e^{N|y|}+1)p_{N,j}(f)$ and the Gaussian tail.
Both integrated bounds tend to zero.
Thus $f$ belongs to the closed convolution ideal,
proving (MT-M12a).

The coordinate map $Q_\eta\to s(Z)$ is now injective.
It is continuous by the integration-by-parts bounds above,
which descend to quotient seminorms because all coordinates
annihilate the ideal. The continuous series map
$(a_\rho)\mapsto[\sum_\rho a_\rho h_\rho]$ is its inverse.
This proves (MT-M12b), and the usual coordinate vectors
give (2), with representatives $h_\rho$.
Statement (2) implies (3). Statement (3) implies (1)
by MT-M11 with $m=1$, extending its high-zero estimate
over the finitely many remaining zeros. This completes
the equivalence.
\end{proof}

Replacing $\psi_\eta$ by
$g_b*\psi_\eta$ for any $b>0$ leaves the closed ideal unchanged:
Gaussian convolution has dense range by MT-M06 and convolution
by $\psi_\eta$ is continuous.

The theorem settles a precise conditional interpolation and
spectral-synthesis interface, rather than identifying every
quotient occurring in Masullo's draft. For this canonical
completed-Hecke quotient with simple zeros, the Schauder and
absolute-Schauder existence gates are exactly the polynomial
reciprocal-derivative gate. The extension to arbitrary
multiplicities requires bounds for the whole local Hermite
interpolation data and is not asserted here.

\subsection*{MT-M13. Nuclear arithmetic convolution under the exact basis gate}
Under the equivalent conditions of MT-M12, convolution by
$f\in\mathcal E$ acts on $Q_\eta\simeq s(Z)$ as
\[
            (a_\rho)\longmapsto
                         (\mathcal Mf(\rho)a_\rho).
\]
It has an explicit nuclear representation on that Fr\'echet
space and a trace-class extension to $\ell^2(Z)$, with
\[
                   \operatorname{Tr}C_f
                         =\sum_{\rho\in Z}\mathcal Mf(\rho).
                                                               \tag{MT-M13}
\]
Here $\operatorname{Tr}$ denotes the trace of that Hilbert-space extension.
The sum converges absolutely. For
$h^\sharp(x)=e^{-x}\overline{h(-x)}$ and
$f=h*h^\sharp$, this Hilbert-space operator is positive and
\[
                   \operatorname{Tr}C_f
                         =\sum_{\rho\in Z}|\mathcal Mh(\rho)|^2
                         \ge0.
\]

\begin{proof}
Convolution descends because it preserves the closed
convolution ideal. Multiplication of Mellin transforms
gives the diagonal action. The eigenvalues decrease
faster than every power of the ordinate, so their
absolute sum is finite by NT-Z10. A diagonal operator
with these eigenvalues on the orthonormal coordinate
basis of $\ell^2(Z)$ is trace class and has the stated
trace.

To see nuclearity without an unspecified basis estimate,
write $b_\rho=[h_\rho]$ and
$\delta_\rho([g])=\mathcal Mg(\rho)$. The functionals
$\delta_\rho$ are equicontinuous: for a fixed
$N>1/2$, their absolute values are uniformly bounded
by a constant times the quotient seminorm of $p_{N,0}$.
The representation
\[
 C_f=\sum_{\rho\in Z}(1+|\operatorname{Im}\rho|)^{-2}
       \delta_\rho\otimes
       \bigl((1+|\operatorname{Im}\rho|)^2
                                  \mathcal Mf(\rho)b_\rho\bigr)
\]
has summable scalar coefficients, an equicontinuous
family of functionals, and a bounded family of vectors
in every quotient seminorm by (9) and rapid decay.
It is therefore an explicit nuclear representation.
Its diagonal sum is the same absolutely convergent
trace displayed above; the unambiguous Hilbert
trace needs no additional abstract trace convention.

The weighted involution preserves $\mathcal E$ and satisfies
\[
      \mathcal M h^\sharp(s)
                =\overline{\mathcal Mh(1-\overline s)}.
\]
At every $\rho$ on the critical line,
$1-\overline\rho=\rho$. Thus the diagonal eigenvalues
for $h*h^\sharp$ are the nonnegative squares stated,
which proves positivity and its trace formula.
\end{proof}

This is a rigorous operator model conditional on the exact
simple-zero basis gate. Its Hilbert completion is constructed
from the zero coordinates; it is not an independently
constructed geometric or prime-side operator proving GRH.
Its location in the programme is Masullo 9.3.2.5--9.3.3,
the Tate-adjoint involution, and cote~1's nuclear trace
framework as carried in SA23 geo-H-lidskii/geo-H-genus.

\subsection*{MT-M14. Bounded multiplicities and the full Hermite basis}
Retain the canonical completion and quotient of MT-M12, assume
GRH, and replace simplicity by the hypothesis
\[
                  1\le m_\rho\le m_0<\infty
                              \qquad(\rho\in Z),       \tag{13}
\]
where $m_\rho$ is the actual zero multiplicity.
Then the following are equivalent:
\begin{enumerate}
\item For some $c,C>0$,
$|L^{(m_\rho)}(\rho,\eta)/m_\rho!|
       \ge c(1+|\operatorname{Im}\rho|)^{-C}$ for all $\rho$.
\item The canonical quotient has an absolute Schauder basis
whose coefficient functionals are
$[f]\mapsto(\mathcal Mf)^{(k)}(\rho)/k!$ for all
$0\le k<m_\rho$.
\item It has an ordinary Schauder basis with precisely those
continuous coefficient functionals.
\end{enumerate}
Under these conditions the closed ideal is the full zero-jet
kernel, and the quotient is the rapid sequence space on
$\{(\rho,k):0\le k<m_\rho\}$, weighted by
$1+|\operatorname{Im}\rho|$. Assumption (13) and condition (1)
are additional arithmetic hypotheses, not consequences of GRH
established in this run.

\begin{proof}
The necessity of (1) from (3) is exactly MT-M11.
We give all changes needed for sufficiency; they also show why
the fixed bound $m_0$ is used.
Write $z=s-\rho$, $m=m_\rho$, and
\[
                  G_\rho(s)=\Xi_\eta(s)/(s-\rho)^m .
\]
The leading-coefficient hypothesis, multiplied by the gamma
and conductor factors, gives
\[
 |G_\rho(\rho)|\ge c(1+|T|)^{-C}e^{-\kappa|T|},
           \qquad T=\operatorname{Im}\rho.
\]
Cauchy's formula on a fixed larger disk around $\rho$ and
the strip bound (5) give
$|G_\rho^{(j)}(\rho)/j!|
 \le C(1+|T|)^C e^{-\kappa|T|}$ for $0\le j\le m_0$.
To see the division estimate directly, on a circle $|z|=2$
divide $\Xi_\eta(\rho+z)$ by $z^m$ and apply Cauchy's
formula inside; $m\le m_0$ keeps all constants fixed.

Fix $a>0$. Let $P_{\rho,k}(z)$ be the Taylor polynomial
through degree $m-1-k$ of the analytic germ
$e^{-az^2}/G_\rho(\rho+z)$.
The reciprocal-coefficient recursion starts with
$1/G_\rho(\rho)$ and at each subsequent degree divides
a finite sum by $G_\rho(\rho)$. The displayed upper and
lower bounds therefore bound every coefficient of
$P_{\rho,k}$ by
$C(1+|T|)^C e^{\kappa|T|}$, uniformly for
$0\le k<m\le m_0$. Only finitely many recursion steps
are involved. Define
\[
 H_{\rho,k}(s)=
      z^kP_{\rho,k}(z)e^{az^2}G_\rho(s).                \tag{14}
\]
Its Taylor expansion at $\rho$ is
$z^k+O(z^m)$, and it vanishes to the required multiplicity
at every other zero. On every fixed vertical strip its
bound has the same Gaussian form (8), with constants
depending also on $m_0$: away from $|z|<1$, use (5)
and the fixed polynomial degree in (14); inside,
use the Cauchy bound for $G_\rho$.
Inverse Mellin transforms $h_{\rho,k}\in\mathcal E$
therefore satisfy polynomial seminorm bounds (9)
uniformly over $k<m_\rho$.

The coefficients $(\mathcal Mf)^{(k)}(\rho)/k!$ are
uniformly rapidly decreasing in $|T|$ for
$k<m_0$. Indeed, they are Mellin transforms of the
fixed finite list $x^k f(x)/k!$, and the
integration-by-parts argument of MT-M12 applies to
each. The series with coefficients these jets and
vectors $h_{\rho,k}$ consequently converges absolutely
in $\mathcal E$, is continuous in $f$, and interpolates
all its jets. Conversely every rapid sequence of jets
is obtained from the corresponding convergent series.
Multiplicity is included in the local zero count,
so the same summability estimate applies.

It remains to verify the full ideal equality.
Repeat the disk argument proving (11), now with
Blaschke factors raised to the actual multiplicities.
The number of factors counted with multiplicity is still
$M=O_A(\log(2+|t|))$. If a nearest zero $\rho_0$
is in the inner half disk, the leading coefficient
identity is
\[
 \frac{L^{(m_{\rho_0})}(\rho_0)}{m_{\rho_0}!}
   =G(\rho_0)
       \lim_{s\to\rho_0}
                    \frac{B(s)}{(s-\rho_0)^{m_{\rho_0}}}.
\]
The upper bound for that leading Blaschke coefficient
is $C_A^M$ times
$\prod_{\rho_j\ne\rho_0}
          |\rho_0-\rho_j|^{m_{\rho_j}}$.
Condition (1) therefore supplies the polynomial lower
bound (12) for this multiplicity-weighted product.
Nearest-zero comparison gives
\[
 |L(s,\eta)|\ge c_A(1+|t|)^{-C_A}
                         |s-\rho_0|^{m_{\rho_0}}
\]
when the nearest distance is less than one.
When that distance is at least one, or the nearest
disk zero is outside the inner half disk, the same
argument gives a polynomial lower bound without
the distance factor. All constants are uniform
because $m_{\rho_0}\le m_0$ and $M=O_A(\log|t|)$.

If $F=\mathcal Mf$ vanishes to the required orders,
Taylor's formula at a nearest zero gives
\[
 F(s)=\frac{(s-\rho)^m}{(m-1)!}
       \int_0^1(1-u)^{m-1}
                         F^{(m)}(\rho+u(s-\rho))\,du .
\]
For $m\le m_0$ its derivative bound is uniformly
rapid in $|t|$. The distance factors cancel.
Thus $F/L$ is rapidly decreasing on each fixed
strip at large height. Exactly the Gaussian-regularized
division and approximate-identity proof in MT-M12
shows that $f$ lies in
$\overline{\psi_\eta*\mathcal E}$.
The reverse inclusion is immediate.
The interpolation series now identifies the quotient
topologically with the rapid jet-sequence space,
as in MT-M12. This proves (2), hence (3), and
completes the equivalence.
\end{proof}

No unspecified infinite Hermite estimate has been hidden in
the argument: the reciprocal-coefficient recursion has at most
$m_0$ steps. With unbounded multiplicities, the exponents of
the polynomial bounds produced by that recursion need not be
uniform. That case remains a separate interpolation gate.

\subsection*{MT-M15. Multiple zeros: positive traces and nilpotent operators}
Under the conditions of MT-M14, convolution has an explicit
nuclear representation on the rapid jet-sequence quotient
and a trace-class realization on the standard $\ell^2$
space of those jets, with
\[
          \operatorname{Tr}C_f
                    =\sum_{\rho\in Z}m_\rho\mathcal Mf(\rho).
                                                               \tag{MT-M15}
\]
For $f=h*h^\sharp$ this trace is
$\sum_\rho m_\rho|\mathcal Mh(\rho)|^2\ge0$.
However, if any $m_\rho>1$, there is no positive-definite
inner product on the full faithful jet quotient for which
every such adjoint-square convolution is a positive
operator. In particular, the trace positivity does not
upgrade to positivity of the entire nonsemisimple action.

\begin{proof}
At a zero $\rho$ let $a_j(f)=(\mathcal Mf)^{(j)}(\rho)/j!$.
The product rule in normalized jets is
\[
       a_k(f*g)=\sum_{j=0}^k a_j(f)a_{k-j}(g).
\]
Thus the block of $C_f$ at $\rho$ is a lower triangular
Toeplitz matrix with diagonal $a_0(f)$ and size $m_\rho$.
All its entries decrease rapidly with the ordinate,
uniformly in the at most $m_0$ indices.
The trace norm of a finite matrix is at most the sum
of the trace norms of its individual matrix units,
so the block trace norms are summable. This proves
the Hilbert trace-class assertion and its trace.

For nuclearity, use the rank-one expansion in these
matrix entries. The jet functionals are equicontinuous:
for a fixed $N>3/2$,
\[
 |(\mathcal Mg)^{(k)}(\rho)/k!|
 \le \frac1{k!}\int_{\mathbb R}|x|^k e^{x/2}|g(x)|\,dx
 \le2p_{N,0}(g),
\]
uniformly even before using $k<m_0$.
The interpolation vectors have polynomial seminorm
bounds, and the matrix entries have rapid decay.
Insert the summable scalar weight $(1+|T|)^{-2}$
as in MT-M13. With at most $m_0^2$ matrix units per
zero, this is an explicit nuclear representation.
The weighted involution has jet coefficients
\[
                   a_j(h^\sharp)=(-1)^j\overline{a_j(h)}.
\]
Its zeroth coefficient gives the asserted nonnegative
trace for an adjoint square.

Finally suppose $m=m_\rho>1$. The interpolating functions
in MT-M14 realize every local jet and zero jets elsewhere.
The local algebra is $\mathbb C[z]/(z^m)$, with
$z^\sharp=-z$ and conjugation on scalars.
Let $p$ denote its identity, extended by zero in every
other zero block, and let $n=z$ in this block and zero
elsewhere. Both are represented by actual functions
of $\mathcal E$, by (14). On the $m$-dimensional local
block, their convolution operators are $I$ and a
nonzero nilpotent $N$.
For every real $u$, the class $h=p+iu\,n$ satisfies
$h^\sharp=h$, and hence its adjoint square acts there as
\[
                         I+2iuN-u^2N^2.
\]
If all adjoint-square operators were positive in a
positive-definite inner product, these operators and
the corresponding operators with $-u$ would be
self-adjoint on the invariant local block.
Subtracting them for $u\ne0$ shows that $iN$ is
self-adjoint, equivalently $N^*=-N$.
A skew-adjoint operator on a finite-dimensional positive
inner-product space is normal and diagonalizable;
nilpotence forces all eigenvalues to vanish, so it is
zero. This contradicts the regular action of $z$
on $\mathbb C[z]/(z^m)$.
The argument is internal to a finite block and
does not depend on any choice of Hilbert completion.
\end{proof}

Passing to the semisimple quotient that retains only the
zeroth coordinate of each block removes the nilpotents.
Using $m_\rho$ copies of each scalar coordinate then
recovers the positive spectral trace with multiplicity.
It does not faithfully recover the discarded higher
jets. This is the precise distinction needed when
transferring Alpöge-style positivity to an orbital
jet quotient in Masullo 9.3.2.5--9.3.3, or to the
cote~1/SA23 trace programme. No assertion that the
draft explicitly claims a positive faithful Hilbert
realization of this entire quotient is made here.

\subsection*{MT-M16. The full arithmetic jet basis has one exact quantitative gate}
The bounded-multiplicity hypothesis in MT-M14 can be removed.
Fix a primitive unitary Hecke function, assume GRH, and use
the canonical $\Xi_\eta,\psi_\eta,Q_\eta$ of MT-M12.
Allow all actual multiplicities $m_\rho$, without a uniform
bound. Write $\rho=1/2+iT$ and set
\[
 P_\rho=
  \prod_{\substack{\nu\in Z,\ \nu\ne\rho\\
              |\operatorname{Im}\nu-T|\le1}}
              |T-\operatorname{Im}\nu|^{m_\nu},
\]
with empty product one. The following conditions are equivalent:
\begin{enumerate}
\item The first nonzero finite-$L$ Taylor coefficients satisfy
$|L^{(m_\rho)}(\rho,\eta)/m_\rho!|
                       \ge c(1+|T|)^{-C}$ uniformly in $\rho$.
\item The local products satisfy
$P_\rho\ge c(1+|T|)^{-C}$ uniformly in $\rho$.
\item The canonical closed convolution quotient has a
Schauder basis with the prescribed normalized Mellin
zero-jet coefficient functionals.
\item It has such an absolute Schauder basis, its closed
ideal is the full zero-jet kernel, and it is topologically
the rapid sequence space on
$\mathcal J=\{(\rho,k):0\le k<m_\rho\}$, weighted by $1+|T|$.
\end{enumerate}
Constants may differ between these equivalent conditions.
Only GRH for this fixed function is used analytically:
$A_D$ supplies it for Dirichlet functions and $M_H$
supplies it conditionally in general. The equivalent
quantitative conditions remain an extra gate; none is
being deduced here from zero location alone.

\begin{proof}
All constants below may depend on the fixed function and
on the indicated fixed strip. NT-Z10 gives both polynomial
strip growth of the finite $L$-function and a local count
including multiplicities. In particular
\[
       m_\rho=O(\log(2+|T|)),\qquad
       \sum_{\rho\in Z}m_\rho(1+|T|)^{-2}<\infty.       \tag{15}
\]
The finitely many low zeros can be included by adjusting
positive constants in every estimate.

\emph{Equivalence of the two quantitative conditions.}
Take the disk centered at $2+iT$ with a zero-free boundary
of radius $R\in[4,5]$. At large height it contains no
poles or trivial zeros. Its total zero multiplicity is
$M=O(\log(2+|T|))$. Remove its zeros by the finite
Blaschke product, with multiplicities, and write $L=GB$.
As proved in MT-M11, $G$ is bounded above by a polynomial
in $|T|$ and below by an inverse polynomial on the
inner half disk; this uses the Euler lower bound at
the center and Harnack's inequality.

The point $\rho$ lies in that inner disk. Its own leading
Blaschke coefficient is
$[R(1-|a|^2)]^{-m_\rho}$ with $a=-3/(2R)$.
Every other factor at $\rho$ is its distance from $\rho$
times a factor bounded above and below by fixed positive
constants. These statements follow directly from
$(z-a_j)/(1-\overline{a_j}z)$ and $|a|\le3/8$.
All other zeros of distance at most one are in the disk.
The distances of the remaining disk zeros are between
one and $2R$. Since there are $M$ factors with multiplicity,
all constant factors together cost at most a fixed power
of $2+|T|$. Therefore
\[
 c(1+|T|)^{-C}P_\rho
   \le\left|\frac{L^{(m_\rho)}(\rho,\eta)}{m_\rho!}\right|
   \le C(1+|T|)^C P_\rho.                              \tag{16}
\]
This proves (1) if and only if (2).

\emph{A radius adapted to the full multiplicity.}
Assume (2), and let $\delta_\rho$ be the lesser of one
and the distance to the nearest distinct zero.
If a distinct zero $\nu$ has this distance less than one,
the product $P_\nu$ contains the factor
$\delta_\rho^{m_\rho}$. Every other factor in $P_\nu$
is at most one. Hence
\[
 \delta_\rho^{m_\rho}\ge P_\nu
                         \ge c(1+|T|)^{-C}.
\]
When $\delta_\rho=1$ the same conclusion is immediate.
Put $r_\rho=\delta_\rho/4$. Equation (15) now gives
\[
                        r_\rho^{-m_\rho}
                                      \le C(1+|T|)^C. \tag{17}
\]
This is the quantitative replacement for the fixed number
of reciprocal-Taylor steps used in MT-M14.

\emph{Minimum modulus without bounded multiplicities.}
The disk proof of (11) works with all multiplicities once
(2) is available. We give the detail that makes the
constants independent of the largest multiplicity.
For a point $s=\sigma+it$ in a fixed strip, choose the
larger disk from MT-M12 and write again $L=GB$.
If that disk has no zeros, the lower bound for $G=L$
already gives the required polynomial lower bound.
If its nearest zero is outside the inner half disk,
every Blaschke factor at $s$ is at least $1/8$.
Otherwise let $\rho_0$ be a nearest disk zero in that
inner half disk. For every other zero,
$|\rho_0-\nu|\le2|s-\nu|$.
The product of the other distances, with multiplicity,
is at least $P_{\rho_0}$: omitted factors of distance
at most one only increase that product, while all
included distances greater than one are at least one.
The normalized Blaschke denominators and radii cost
only a constant to the total multiplicity
$M=O(\log(2+|t|))$. The lower bound for $G$ therefore
gives
\[
 |L(s,\eta)|\ge c_A(1+|t|)^{-C_A}
             \operatorname{dist}(s,Z)^{m_{\rho_0}}
                       \quad\hbox{if }\operatorname{dist}(s,Z)<1,
                                                               \tag{18}
\]
and a polynomial lower bound when that distance is
at least one. The choice of a nearest zero in (18)
is harmless in a tie. Zeros outside the larger disk
have distance greater than one from $s$.

\emph{Uniform Hermite interpolation.}
Put $m=m_\rho$, $z=s-\rho$, and
$G_\rho(s)=\Xi_\eta(s)/z^m$.
On the circle $|z|=r_\rho$, $\rho$ is the nearest zero.
Divide (18) by $|z|^m$ and use the fixed-strip lower
Stirling bound for the gamma factors. This gives
\[
             |G_\rho(s)|\ge
                 c(1+|T|)^{-C}e^{-\kappa|T|}
                            \quad(|z|=r_\rho).
\]
The function $G_\rho$ has no zero in that disk.
Consequently Cauchy's formula bounds the $j$th Taylor
coefficient of $e^{-az^2}/G_\rho(\rho+z)$ by
\[
       C(1+|T|)^C e^{\kappa|T|}r_\rho^{-j}.
                                                               \tag{19}
\]
For every $j<m$, (17) makes this a polynomial bound
times $e^{\kappa|T|}$, with constants independent of $j,m$.
Use exactly the Taylor truncation $P_{\rho,k}$ and
cardinal function (14) of MT-M14, now for all $k<m$.

Their fixed-strip bounds are uniform. When $|z|\ge1$,
the degree of $z^kP_{\rho,k}$ is at most $m-1$, so
\[
 |z^kP_{\rho,k}(z)/z^m|
       \le m\,C(1+|T|)^C e^{\kappa|T|}.
 \]
There is no growing positive power of $|z|$.
The factor $m$ is absorbed using (15), and the Gaussian
bound (8) follows from (5), just as before.
When $|z|<1$, the maximum principle applied to
$G_\rho$ on $|z|=2$ gives
$|G_\rho(s)|\le C(1+|T|)^C e^{-\kappa|T|}2^{-m}$.
The polynomial $z^kP_{\rho,k}$ is bounded by
$m$ times the coefficient bound (19), with (17).
This proves (8) on the remaining region as well,
after increasing a fixed constant.
Inverse transforms $h_{\rho,k}$ therefore satisfy
the polynomial seminorm bounds (9) uniformly over
every jet index $k<m_\rho$.

The normalized jets of any $f\in\mathcal E$ also have
uniform rapid height decay for \emph{all} $k\ge0$:
for each fixed $j$ and each fixed real-part strip,
\[
 \sup_{k\ge0}
 (1+|t|)^j\left|\frac{(\mathcal Mf)^{(k)}(\sigma+it)}{k!}\right|
                    \le C_{A,j}\,p_{N,j}'(f),          \tag{20}
\]
where $p_{N,j}'$ is a finite maximum of defining
seminorms, with sufficiently large fixed $N$.
Indeed, integrate by parts $j$ times in the transform
of $e^{\sigma x}x^k f(x)/k!$.
Every derivative of $x^k/k!$ is either zero or
$x^{k-\ell}/(k-\ell)!$; the remaining coefficients
depend only on $A,j$. For $N-A>1$, its weighted
$L^1$ norm is at most
$2/(N-A)^{k-\ell+1}\le2$, times the corresponding
seminorm of $f$. The region $|t|\le1$ follows from
the same integral bound without integration by parts.
This proves (20).

Equations (15), (20), and the uniform cardinal bounds
give absolute interpolation and its continuous inverse
for every rapid sequence on $\mathcal J$.
Finally, division proves the full ideal equality.
If $F=\mathcal Mf$ has all prescribed zero jets, its
Taylor remainder at a nearest zero is
\[
 F(s)=z^m\int_0^1m(1-u)^{m-1}
                  \frac{F^{(m)}(\rho+uz)}{m!}\,du .
 \]
The measure in this integral has total mass one.
Estimate (20) thus bounds the normalized remainder
uniformly in $m$, and its $|z|^m$ cancels the distance
factor in (18). The quotient $F/L$ is rapidly
decreasing on every fixed strip at large height.
The entire quotient $F/\Xi_\eta$, multiplied by
$e^{\epsilon s^2}$, is the Mellin transform of an
element of $\mathcal E$ for every $\epsilon>0$.
The Gaussian approximate-identity argument of
MT-M12 proves that $f$ lies in the closed ideal.
The reverse inclusion is immediate.
This yields the asserted topological rapid-sequence
identification and (4). Statement (4) implies (3),
and (3) implies (1) by MT-M11. All four conditions
are therefore equivalent.
\end{proof}

This removes both the simplicity assumption of MT-M12 and the
uniform multiplicity assumption of MT-M14. It does not prove the
remaining leading-coefficient or cluster-product bound.
The local $O(\log|T|)$ count is used together with that
quantitative condition; MT-M09 shows why the count alone
cannot replace it.

\subsection*{MT-M17. Nuclear trace with arbitrary zero multiplicities}
Under the equivalent conditions of MT-M16, every
$f\in\mathcal E$ acts by nuclear convolution on $Q_\eta$
and by a trace-class operator on $\ell^2(\mathcal J)$.
Its Hilbert trace is
\[
       \operatorname{Tr}C_f
                  =\sum_{\rho\in Z}m_\rho\mathcal Mf(\rho).
\]
For $f=h*h^\sharp$ this is the nonnegative sum
$\sum_\rho m_\rho|\mathcal Mh(\rho)|^2$.
The faithful-operator positivity obstruction of MT-M15
persists whenever any zero is multiple.

To prove these assertions without a uniform block-size
hypothesis, note that (15) bounds the total multiplicity
in each unit-height interval by $O(\log(2+|T|))$.
The number of Toeplitz matrix units in that interval is
therefore at most $O(\log^2(2+|T|))$.
Estimate (20) makes every matrix entry of $C_f$ uniformly
rapidly decreasing in height, independently of its jet
index. Thus the sum of the trace norms of all matrix
units is finite, giving the trace-class assertion and
the displayed diagonal trace. For nuclearity use the
equicontinuous normalized-jet functionals, the uniform
polynomial seminorm bounds for $h_{\rho,k}$, and the
scalar weight $(1+|T|)^{-2}$ for each matrix unit.
These weights are summable even with the
$O(\log^2(2+|T|))$ count, and the resulting vector
family is bounded in every seminorm by (20).
This is the same explicit nuclear decomposition as
MT-M15, now with all indices included.
The involution formula and the finite local nilpotent
argument in MT-M15 make no bounded-multiplicity
assumption, so they prove the remaining statements.

This is an exact reduction for the prescribed normalized-jet
coordinates. MT-M18 and MT-M19 below supply a different,
block-adapted absolute basis, spectral synthesis and nuclearity
without that quantitative condition. Identification with a
particular independently defined adelic quotient is a separate
map; MT-M21 begins the comparison with Meyer's construction.

\subsection*{MT-M18. Unconditional arithmetic block synthesis and an absolute adapted basis}

Fix a primitive unitary Hecke character \(\eta\) over a number field.
Use the entire completion \(\Xi=\Xi_\eta=P_\eta\Lambda(s,\eta)\)
and its inverse Mellin transform \(\psi=\psi_\eta\) from MT-M12,
but make no GRH or simplicity assumption.  The polynomial \(P_\eta\)
removes the possible pure-norm poles exactly as there.  Let \(Z\)
be the set of distinct zeros of \(\Xi\), with their actual
multiplicities \(m_\rho\).  These are the nontrivial zeros and
lie in \(0\le\operatorname{Re}s\le1\).
Put
\[
 \begin{gathered}
 \mathcal E=\{f\in C^\infty(\mathbb R,\mathbb C):
   p_{N,j}(f)=\sup_x e^{N|x|}|f^{(j)}(x)|<\infty
                    \text{ for all integers }N,j\ge0\},\\
 Q=\mathcal E/\overline{\psi*\mathcal E}.
 \end{gathered}
\]
All constants below concern this one fixed Hecke function; no
uniformity in its field, conductor, or infinity type is asserted.

There are real numbers
\[
 n+\tfrac14<h_n<n+\tfrac34\quad(n\in\mathbb Z)
\]
not equal to any zero ordinate, with the following properties.
Let \(\Omega_n\) be the interior of the rectangle
\([-1,2]+i[h_n,h_{n+1}]\), and orient its boundary
\(\Gamma_n\) positively.  Let \(A(\Omega_n)\) be the Banach
algebra of functions holomorphic in \(\Omega_n\) and continuous
on its closure, with the supremum norm.  Write
\[
 \begin{split}
 I_n&=\{F\in A(\Omega_n):
        F^{(j)}(\rho)=0\text{ for }\rho\in Z\cap\Omega_n,\
                                  0\le j<m_\rho\},\\
 J_n&=A(\Omega_n)/I_n,\qquad
 \|a\|_n=\inf_{F\bmod I_n=a}\sup_{\Gamma_n}|F|,\qquad
 d_n=\dim_{\mathbb C}J_n=\sum_{\rho\in Z\cap\Omega_n}m_\rho .
 \end{split}                                                    \tag{M18.1}
\]
An empty block has \(J_n=0\).  Then
\[
 d_n=O(\log(2+|n|)),\qquad
 s(J)=\{(a_n):q_k(a):=\sup_n(2+|n|)^k\|a_n\|_n<\infty
                                  \text{ for every }k\ge0\}      \tag{M18.2}
\]
is a Fr\'echet space, and the Mellin block restriction gives a
topological isomorphism
\[
 Q\simeq s(J).                                                    \tag{M18.3}
\]
In particular
\[
 \overline{\psi*\mathcal E}
 =\{f\in\mathcal E:(\mathcal Mf)^{(j)}(\rho)=0
                         \text{ for every }\rho,\ 0\le j<m_\rho\}.
                                                                  \tag{M18.4}
\]
The block expansion in \(Q\) converges absolutely in every
continuous seminorm.  Choosing an Auerbach basis in each
\((J_n,\|\cdot\|_n)\) gives an absolute Schauder basis of \(Q\)
with block-adapted continuous coefficient functionals.

Neither (M18.3) nor this last basis assertion prescribes the
individual zero jets as its coordinate functionals.  Thus the
result does not assert the quantitative condition of MT-M16.
The block partition and its adapted bases are choices; the
closed quotient in (M18.3) is the canonical one.

\begin{proof}
\emph{Analytic inputs and the transform topology.}
NT-Z10 proves unconditionally, including for a fixed
infinite-order unitary character, the polynomial upper bound
for the finite \(L(s,\eta)\) on every fixed vertical strip and
the bound \(O(\log(2+|T|))\) for the number of nontrivial
zeros, with multiplicities, in any fixed-height window.
It also gives
\[
                   |L(2+iT,\eta)|\ge\zeta_K(2)^{-1}.              \tag{M18.5}
\]
The standard completed Hecke functional equation locates all
the zeros of the entire completion in the fixed critical strip.
Exact classical source locators are Ramakrishnan--Valenza,
\emph{Fourier Analysis on Number Fields}, \S7.4,
Theorem 7-19, printed p.~279, and \S7.7, Theorem 7-28,
printed p.~289; the fixed-strip and local-count deductions
are proved in NT-Z10.  We use these unconditional facts,
not a zero-location consequence of either GRH input.

Put \(\kappa=\pi[K:\mathbb Q]/4\).  Stirling's formula and
the preceding strip bound give, for every fixed \(A>0\),
\[
 |\Xi(\sigma+it)|\le C_A(2+|t|)^{B_A}e^{-\kappa|t|}
                         \quad(|\sigma|\le A).                  \tag{M18.6}
\]
Cauchy's formula gives the same type of bound for \(\Xi'\)
on a smaller strip.  For large \(|t|\), the factor
\(\Xi(s)/L(s,\eta)\) is nonzero and has upper and lower
bounds by \(e^{-\kappa|t|}\) times fixed positive or negative
powers of \(2+|t|\), uniformly on each fixed strip.  Its
finitely located gamma exceptions are immaterial at large
height.

For clarity, the transform fact used below is elementary.
If an entire function \(H\) decreases faster than every
power of \(|t|\), uniformly on each strip \(|\sigma|\le A\),
its inverse Fourier integral on \(\sigma=0\) is the Mellin
transform inverse of an element of \(\mathcal E\).
For \(x\ge0\), shift the inverse contour to
\(\sigma=N+1\), and for \(x\le0\) shift it to
\(\sigma=-N-1\).  The horizontal edges tend to zero by the
uniform strip decay.  The \(j\)th derivative is bounded
by the integral of \(|s|^j|H(s)|\) on the shifted line,
times \(e^{-(N+1)|x|}\); these integrals are finite.
This proves all \(p_{N,j}\) bounds and continuity from the
corresponding strip seminorms.  Fourier inversion proves
the transform identity.  Conversely, integration by parts
in \(\mathcal Mf(\sigma+it)=\int e^{(\sigma+it)x}f(x)\,dx\)
shows rapid vertical decay, with each strip seminorm
bounded by finitely many \(p_{N,j}\).  Entireness follows
by exponential domination.  In particular (M18.6) proves
\(\psi\in\mathcal E\) and \(\mathcal M\psi=\Xi\), without GRH.

\emph{Selecting all horizontal boundaries at once.}
For \(n\in\mathbb Z\), let \(M_n\) count, with multiplicities,
the zeros whose ordinates lie in \([n-2,n+3]\).  Thus
\(M_n\le C\log(2+|n|)\).  On
\(I_n^*=(n+1/4,n+3/4)\), set
\[
 U_n(h)=\sum_{\substack{\rho\in Z\\
                    \operatorname{Im}\rho\in[n-2,n+3]}}
       m_\rho\log^+\frac1{|h-\operatorname{Im}\rho|}.
\]
Every summand has integral at most
\(\int_{-1}^1-\log|u|\,du=2\) over \(I_n^*\).
Averaging over its length \(1/2\) therefore permits a
choice \(h_n\in I_n^*\), avoiding every zero ordinate,
with \(U_n(h_n)\le 5M_n\); if \(M_n=0\), use its midpoint.
(The harmless slack from \(4M_n\) to \(5M_n\) gives a
positive-measure set when \(M_n>0\).)
Every zero with \(|h_n-\operatorname{Im}\rho|<1\)
is included in this sum.  Consequently
\[
 \prod_{\rho\in Z}
   \min\{1,|h_n-\operatorname{Im}\rho|\}^{m_\rho}
      \ge e^{-5M_n}\ge c(2+|n|)^{-C}.                            \tag{M18.7}
\]
Only finitely many factors in this product differ from one.
Also \(1/2<h_{n+1}-h_n<3/2\), so the rectangles have
uniformly bounded perimeter and cover all zero ordinates.

We claim, for every fixed \(A>0\), that
\[
 \min_{|\sigma|\le A}|\Xi(\sigma+i h_n)|
       \ge c_A(2+|n|)^{-C_A}e^{-\kappa|n|}.                      \tag{M18.8}
\]
Here is a complete minimum-modulus argument.  For large
\(|h_n|\), center a disk at \(c_n=2+i h_n\) and take a
zero-free boundary of radius
\(R\in[4(A+3)+4,4(A+3)+5]\).  Such a radius exists because
there are only finitely many zeros in that compact annulus.
At sufficiently large height the disk contains no finite-\(L\)
pole or trivial zero, since their possible ordinates are fixed
by the infinity type.  Its total zero multiplicity is
\(M=O_A(\log(2+|n|))\), by NT-Z10.  Remove its zeros by the
finite disk Blaschke product \(B\), with multiplicities,
and write \(L=GB\).  The factor \(G\) is holomorphic and
zero-free on the disk and its boundary.  Since \(|B|=1\)
on the boundary, (M18.6)'s finite-\(L\) antecedent and the
maximum principle give \(|G|\le C(2+|n|)^C\) there and
inside.  At the center, \(|B(c_n)|\le1\), so (M18.5)
gives \(|G(c_n)|\ge c>0\).
Apply Harnack's inequality to the nonnegative harmonic
function \(\log(C(2+|n|)^C)-\log|G|\).  On the inner
half disk its value is at most three times its value
at the center.  Thus
\[
                        |G(s)|\ge c(2+|n|)^{-C}
\]
on that inner disk.  This form of Harnack follows by
bounding the ratios of the Poisson kernels on concentric
circles; the ratio is \((R+r)/(R-r)\le3\) for \(r\le R/2\).

For \(s=\sigma+i h_n\), \(|\sigma|\le A\), the normalized
coordinate \(z=(s-c_n)/R\) has \(|z|<1/4\).
For a disk zero \(\rho\), with \(a_\rho=(\rho-c_n)/R\),
\[
 \left|\frac{z-a_\rho}{1-\overline{a_\rho}z}\right|
    \ge\frac{|s-\rho|}{2R}
    \ge\frac{\min\{1,|h_n-\operatorname{Im}\rho|\}}{2R}.
\]
Products of omitted factors in (M18.7) are at most one.
It follows that
\[
 |B(s)|\ge(2R)^{-M}
   \prod_{\rho\in Z}
      \min\{1,|h_n-\operatorname{Im}\rho|\}^{m_\rho}
             \ge c_A(2+|n|)^{-C_A}.
\]
Multiplication by \(G\) and by the lower Stirling bound
for \(\Xi/L\) proves (M18.8) at large height.
There are only finitely many remaining \(n\); each
horizontal segment is compact and zero-free for \(\Xi\),
so reducing \(c_A>0\) covers them.

The same disk proof, now centered at \(2+it\), gives a
polynomial lower bound for \(L(\sigma+it,\eta)\) on either
fixed vertical line \(\sigma=-A\) or \(\sigma=A\), for
every \(A\ge2\) and large \(|t|\).  Indeed their distances
from all disk zeros in the critical strip are at least one,
so every normalized Blaschke factor is at least \(1/(2R)\).
There is no logarithmic-distance product to control.
The lower Stirling bound and the compact low-height range
give the corresponding lower bound for \(\Xi\).
The same argument applies to the line \(\sigma=-1\)
and to \(\sigma=2\).  Combining it with (M18.8) yields
\[
 \inf_{w\in\Gamma_n}|\Xi(w)|
       \ge c(2+|n|)^{-C}e^{-\kappa|n|}.                          \tag{M18.9}
\]
All constants here are independent of separations between
zeros inside a rectangle.

\emph{The finite block algebras and an entire interpolation kernel.}
All zeros of a block are interior points.  Each of its
finitely many jet functionals is continuous on
\(A(\Omega_n)\), by Cauchy's formula on a disk contained
in \(\Omega_n\).  The map to these jets is onto, by
ordinary polynomial Hermite interpolation.  Its kernel
is \(I_n\), a closed ideal, proving that (M18.1) is a
genuine finite-dimensional quotient norm, with dimension
\(d_n\).  The maximum principle identifies the supremum
over the closure with that over \(\Gamma_n\).
The quotient norm is submultiplicative because the
supremum norm is.  The local zero count gives (M18.2)'s
dimension bound.

Fix \(a>0\).  For \(F\in A(\Omega_n)\) define
\[
 H_nF(z)=\frac1{2\pi i}\int_{\Gamma_n}
    F(w)\frac{\Xi(z)-\Xi(w)}{\Xi(w)(z-w)}
                      e^{a(z-w)^2}\,dw .                       \tag{M18.10}
\]
The apparent singularity at \(z=w\) is removable.
Consequently this is an entire function of \(z\).
If \(z\) is inside the rectangle, Cauchy's formula
for the second term gives
\[
 H_nF(z)=F(z)+\Xi(z)\frac1{2\pi i}
             \int_{\Gamma_n}
               \frac{F(w)e^{a(z-w)^2}}{\Xi(w)(z-w)}\,dw .
                                                                  \tag{M18.11}
\]
The integral on the right is holomorphic in the interior.
If \(z\) is outside the closed rectangle, that second
Cauchy integral is zero and \(H_nF(z)\) is divisible
locally by \(\Xi(z)\).
These statements follow also for continuous boundary
values by applying the formulas on interior rectangles
and taking limits.  Thus \(H_nF\) matches every
prescribed block jet of \(F\), and vanishes to the
full multiplicity at every zero outside the block.

Moreover, if all the block jets of \(F\) vanish, then
\(F/\Xi\) is holomorphic in \(\Omega_n\) and continuous
on its closure.  For exterior \(z\), both terms in
(M18.10) then integrate to zero by Cauchy's theorem.
Entireness implies \(H_nF=0\) everywhere.  It follows
that \(H_n\) depends only on \(a_n=F\bmod I_n\).
Denote this entire function by \(H_n a_n\).

For every fixed \(A>0\) there are constants \(C_A,B_A'>0\)
and \(c_a>0\) such that
\[
 |H_n a_n(\sigma+it)|
       \le C_A(2+|n|)^{B_A'}\|a_n\|_n
                         e^{-c_a(t-n)^2},
                 \qquad |\sigma|\le A.                         \tag{M18.12}
\]
To check the estimate before taking the quotient infimum,
bound \(|F|\) by its boundary supremum in (M18.10).
If \(|z-w|\ge1\), use (M18.6) and (M18.9) to bound
the difference quotient by
\[
 C_A(2+|n|)^C(2+|t|)^C
                    e^{\kappa\,\lvert t-\operatorname{Im}w\rvert}
       +1 .
\]
Here \(|\operatorname{Im}w-n|<2\), while
\[
 |e^{a(z-w)^2}|
       \le C_A e^{-a(t-\operatorname{Im}w)^2}.
\]
A fixed polynomial and the fixed linear exponential
are absorbed by a smaller Gaussian, and
\(2+|t|\le(2+|n|)(2+|t-n|)\).
If \(|z-w|<1\), the identity
\[
 \frac{\Xi(z)-\Xi(w)}{z-w}
      =\int_0^1\Xi'(w+u(z-w))\,du
\]
and the strip derivative bound give a polynomial in
\(2+|n|\) after division by (M18.9).
In this region \(|t-n|<3\), so it too satisfies (M18.12)
after a constant adjustment.  The perimeter is bounded
independently of \(n\).  Taking the infimum over all
representatives of \(a_n\) proves (M18.12).

The transform fact above now gives a unique
\(\mathcal T_n(a_n)\in\mathcal E\) whose Mellin transform is
\(H_n a_n\).  It is linear in \(a_n\), and for every
defining seminorm
\[
 p_{N,j}(\mathcal T_n(a_n))
       \le C_{N,j}(2+|n|)^{B_{N,j}}\|a_n\|_n .                  \tag{M18.13}
\]
Indeed on either shifted vertical line, the integral
of \((1+|t|)^j\) against (M18.12)'s Gaussian costs
only another fixed power of \(2+|n|\).
In particular no individual reciprocal derivative
or Hermite interpolation matrix norm has entered
these estimates.

\emph{Spectral synthesis for the closed ideal.}
Write \(J\) for the right side of (M18.4).
It is a closed subspace, since each Mellin jet is
continuous.  Multiplication of Mellin transforms
immediately gives
\(\overline{\psi*\mathcal E}\subseteq J\).
For the reverse inclusion take \(f\in J\) and
\(F=\mathcal Mf\).  The quotient \(F/\Xi\) is entire:
every zero, including its full multiplicity, has
been cancelled.

Fix \(A\ge2\).  Use the wider rectangles
\([-A,A]+i[h_n,h_{n+1}]\).  On their horizontal
edges (M18.8) applies; on their vertical edges
the vertical-line argument following (M18.8)
applies.  The rapid fixed-strip decay of \(F\)
and these lower bounds show, on the whole boundary,
\[
 \left|\frac{F(s)}{\Xi(s)}\right|
       \le C_{A,f}(2+|n|)^{C_A}e^{\kappa|n|}.
\]
The maximum principle for the entire quotient
extends this bound through the rectangle.
Since their heights have gaps between \(1/2\)
and \(3/2\), it follows that on every fixed strip
\[
 \left|\frac{F(\sigma+it)}{\Xi(\sigma+it)}\right|
       \le C_{A,f}(2+|t|)^{C_A}e^{\kappa|t|}.                  \tag{M18.14}
\]
For every \(\epsilon>0\), the entire function
\[
                   e^{\epsilon s^2}F(s)/\Xi(s)
\]
therefore has rapid decay on every fixed strip.
It is the Mellin transform of some
\(j_\epsilon\in\mathcal E\).
With \(g_\epsilon(x)=(4\pi\epsilon)^{-1/2}
e^{-x^2/(4\epsilon)}\), Fourier uniqueness gives
\[
                         \psi*j_\epsilon=g_\epsilon*f.
\]
The right side tends to \(f\) in \(\mathcal E\).
For completeness, on \(|y|\le1\) the weighted
mean-value bound for
\(f^{(j)}(x-y)-f^{(j)}(x)\) is
\(|y|e^{N|y|}p_{N,j+1}(f)\).
On the complement of any fixed neighborhood of
zero use
\((e^{N|y|}+1)p_{N,j}(f)\).
Integration against \(g_\epsilon(y)\) makes both
bounds tend to zero, by its first moment and
Gaussian tails.  This proves the reverse inclusion
and hence (M18.4), independently of the interpolation
construction.

\emph{The topological decomposition and the adapted basis.}
The restriction map
\[
 R:\mathcal E\longrightarrow s(J),\qquad
        (Rf)_n=(\mathcal Mf)|_{\Omega_n}\bmod I_n
\]
is continuous.  Indeed the quotient norm is at most
\(\sup_{\Gamma_n}|\mathcal Mf|\), and these boundary
suprema decrease faster than every power of \(|n|\),
with each such bound controlled by finitely many
seminorms of \(f\).  Define
\[
                    S(a)=\sum_{n\in\mathbb Z}\mathcal T_n(a_n).
\]
For each seminorm \(p_{N,j}\), (M18.13) and
\(q_{B_{N,j}+2}(a)\), with the exponent rounded up
if necessary, bound the sum of its term seminorms
by a constant times \(\sum_n(2+|n|)^{-2}\).
Thus this series converges absolutely in every
seminorm of \(\mathcal E\), and defines a continuous
linear map \(S:s(J)\to\mathcal E\).
Completeness of \(s(J)\) follows directly from the
complete finite-dimensional norms in its weighted
supremum seminorms; completeness of \(\mathcal E\)
follows from convergence of all weighted derivative
seminorms on compact intervals.

Every Mellin jet of \(S(a)\) may be computed
termwise by continuity.  The interpolation property
of (M18.11) gives \(RS=\mathrm{id}\).
Also \(\ker R=J=\overline{\psi*\mathcal E}\).
Therefore \(R\) induces the continuous bijection
(M18.3), with continuous inverse
\(a\mapsto[S(a)]\).
The block operators \(\mathcal T_n R_n\) on \(\mathcal E\)
are finite-rank continuous projections, mutually
annihilating for distinct \(n\).
Their absolute sum is the continuous projection
\(SR\); on \(Q\) their sum is the identity.

Finally every finite-dimensional complex normed
space has an Auerbach basis \(v_{n,1},\ldots,v_{n,d_n}\):
its vectors and dual coordinate functionals
\(\phi_{n,j}\) all have norm one.
Here is the elementary existence argument.  In
any fixed coordinates maximize the absolute
determinant of \(d_n\) vectors in the compact unit
ball.  A maximum is positive.  Replacing its
\(j\)th column by any unit-ball vector multiplies
the determinant by the corresponding dual
coordinate.  Maximality makes that coordinate
have modulus at most one; its value on its own
column is one.  This proves the claimed norms.

For \(a=(a_n)\in s(J)\), its resulting scalar
coefficients satisfy
\(|\phi_{n,j}(a_n)|\le\|a_n\|_n\).
Writing \(e_{n,j}\) for \(v_{n,j}\) in its block
and zero elsewhere, one has
\[
 \sum_{n,j}q_k\bigl(\phi_{n,j}(a_n)e_{n,j}\bigr)
 \le\sum_n d_n(2+|n|)^k\|a_n\|_n
 \le q_{k+3}(a)\sum_n d_n(2+|n|)^{-3}<\infty .
\]
This proves an absolute expansion in every
seminorm, in any enumeration of these vectors.
The coordinate functionals are continuous and
unique by the block isomorphism and the finite
basis in each block.  Transport through (M18.3)
proves the asserted absolute Schauder basis.
\end{proof}

The spectral block \(J_n\), as an algebra, is the
finite-length complex algebra
\[
 \mathbb C[s]\Big/
       \prod_{\rho\in Z\cap\Omega_n}(s-\rho)^{m_\rho}
 \simeq
 \prod_{\rho\in Z\cap\Omega_n}
                  \mathbb C[z]/(z^{m_\rho}).
\]
Its stable norm is the quotient holomorphic supremum
norm in (M18.1), not a norm obtained by declaring its
individual jet coordinates orthonormal.
Thus the construction is a finite-length spectral
divisor decomposition compatible with the Mellin
convolution algebra.  It directly refines the repair
of Masullo's Lemma 9.3.2.5, original PDF pp.~425--426:
closure of the convolution ideal is essential, and
the adapted basis has different coordinates from
that lemma's prescribed jets.  The inherited programme
locators are cote~1's nuclear-trace framework and
SA23 geo-H-lidskii/geo-H-genus, as carried in the
attached returns; no new inspection of original
cote scans is claimed here.  Nuclearity and the
trace action of these finite spectral algebras
are treated separately in MT-M19.

No first-priority claim is made for spectral synthesis in related
function spaces; the explicit block norms and bounds above supply
the precise topology and interpolation statement used here.



\subsection*{MT-M19. Unconditional nuclear convolution and its exact positivity criterion}
Use the canonical quotient, stable block algebras \(J_n\), and
topological isomorphism \(Q\simeq s(J)\) of MT-M18. No GRH,
simplicity, or reciprocal-derivative bound is assumed initially.
For every \(f\in\mathcal E\), convolution \(C_f\) is nuclear on
\(Q\). There is a Hilbert completion \(\mathcal H_B\) obtained
by making an Auerbach basis in each block orthonormal, on which
every \(C_f\) is trace class, with
\[
 \operatorname{Tr}_{\mathcal H_B} C_f
       =\sum_{\rho\in Z}m_\rho\mathcal Mf(\rho).       \tag{M19.1}
\]
The series is absolutely convergent and the trace is independent
of all block and Auerbach choices. Nuclearity is uniform when
\(f\) ranges over a bounded subset of \(\mathcal E\).

Define \(h^\sharp(x)=e^{-x}\overline{h(-x)}\). Then
\[
 \operatorname{Tr}_{\mathcal H_B} C_{h*h^\sharp}\ge0
          \quad\hbox{for every }h\in\mathcal E
 \quad\Longleftrightarrow\quad
 \text{GRH for }L(s,\eta).                              \tag{M19.2}
\]
Under GRH the trace equals
\(\sum_\rho m_\rho|\mathcal Mh(\rho)|^2\).
This concerns the trace; the chosen block Hilbert norm need
not make the operator \(C_{h*h^\sharp}\) positive.

\begin{proof}
Put \(F=\mathcal Mf\) and \(b_n(f)=\sup_{w\in\Gamma_n}|F(w)|\).
These numbers decrease faster than every power of \(2+|n|\),
with each decay bound controlled by finitely many defining
seminorms of \(f\). This follows from the integration-by-parts
estimate in MT-M18 on the one fixed strip containing all
\(\Gamma_n\). Convolution preserves the closed ideal: the
identity \(\mathcal M(f*g)=\mathcal Mf\,\mathcal Mg\) holds
by absolute integrability, and convolution is continuous on
\(\mathcal E\). Explicitly its seminorms are bounded by
\[
 p_{N,j}(f*g)
 \le p_{N,j}(f)\int_{\mathbb R}e^{N|y|}|g(y)|\,dy
 \le 2p_{N,j}(f)p_{N+1,0}(g).
\]
On \(J_n\), the induced map \(B_n(f)\) is multiplication by
the restriction of \(F\). The quotient Banach-algebra norm
therefore gives
\[
                      \|B_n(f)\|\le b_n(f).            \tag{M19.3}
\]

Let \(v_{n,j},\phi_{n,j}\), \(1\le j\le d_n\), be the
Auerbach vectors and dual functionals of MT-M18, all of norm
one. Regard each vector as supported in its own block and
write \(\ell_{n,j}(a)=\phi_{n,j}(a_n)\). The family
\(\ell_{n,j}\) is equicontinuous, since
\(|\ell_{n,j}(a)|\le q_0(a)\). Set
\[
 \lambda_{n,j}=(2+|n|)^{-2},\qquad
 u_{n,j}=(2+|n|)^2 B_n(f)v_{n,j}.
\]
The sum of \(|\lambda_{n,j}|\) is finite because
\(d_n=O(\log(2+|n|))\). For each \(k\),
\[
 q_k(u_{n,j})\le(2+|n|)^{k+2}b_n(f),
\]
so \(u_{n,j}\) is a bounded family in \(s(J)\). The series
\[
 C_f a=\sum_{n,j}\lambda_{n,j}\ell_{n,j}(a)u_{n,j}       \tag{M19.4}
\]
is thus an explicit nuclear decomposition, absolutely convergent
in every seminorm. Its value is the block multiplication map
because the finite Auerbach expansion in each block is exact.
The same estimates are uniform on a bounded set of \(f\)'s,
which proves the asserted uniformity.

Give the scalar Auerbach coordinates their ordinary square-sum
norm, and let \(\mathcal H_B\) be its completion. For each
nonempty block this Euclidean norm lies between
\(d_n^{-1/2}\|a_n\|_n\) and
\(d_n^{1/2}\|a_n\|_n\), since the vectors and coordinate
functionals have norm one. Consequently the inclusion
\(s(J)\to\mathcal H_B\) is continuous and injective, and
has dense image because it contains all finitely supported
block vectors. Every matrix entry of \(B_n(f)\) in these
orthonormal coordinates has modulus at most \(b_n(f)\),
by (M19.3). Expressing it as a sum of rank-one matrix units
therefore gives
\[
       \|B_n(f)\|_1\le d_n^2b_n(f),\qquad
       \sum_n\|B_n(f)\|_1<\infty.                     \tag{M19.5}
\]
The block direct sum is trace class and extends \(C_f\).

Algebraically, polynomial Hermite interpolation identifies
\(J_n\) with
\(\prod_{\rho\in Z\cap\Omega_n}\mathbb C[z]/(z^{m_\rho})\).
In the monomial basis of a local factor, multiplication by
\(F(\rho+z)\) is triangular with its constant term
\(F(\rho)\) on every diagonal entry. Its finite-dimensional
trace is \(m_\rho F(\rho)\). Trace is invariant under a
change of basis, so (M19.5) proves (M19.1). It also proves
the choice-independence stated there. Absolute convergence
follows directly from
\(\sum_n d_n b_n(f)<\infty\). All traces in the statement
are the ordinary Hilbert traces just constructed; no
unjustified trace on an arbitrary locally convex space is used.

Direct substitution gives
\[
 \mathcal Mh^\sharp(s)
       =\overline{\mathcal Mh(1-\bar s)}.              \tag{M19.6}
\]
The unitary Hecke functional equation implies that the zero
involution \(\rho\mapsto\rho^\dagger=1-\bar\rho\)
preserves \(Z\) and its multiplicities. Indeed it relates
\(\Lambda(s,\eta)\) to \(\Lambda(1-s,\bar\eta)\),
while complex conjugation relates the latter character to
\(\eta\). The pole-removal polynomial introduces no zeros
into \(Z\). Formula (M19.6) therefore preserves the full
zero-jet ideal, so \(\sharp\) descends to an antilinear
algebra involution of \(Q\). If GRH holds then \(\rho^\dagger=\rho\),
and (M19.1), (M19.6) give the stated nonnegative sum.

Conversely, suppose \(\rho\ne\rho^\dagger\) is a pair
off the critical line, of common multiplicity \(m\).
They have the same ordinate and hence lie in the same
block. In that block prescribe the constant jet 1 at
\(\rho\), the constant jet \(-1\) at \(\rho^\dagger\),
and the zero jet at every other zero; prescribe zero in
all other blocks. This is a finitely supported element
of \(s(J)\). The right inverse in MT-M18 realizes it as
\(\mathcal Mh\) for some \(h\in\mathcal E\). Formula
(M19.1) then has just two nonzero contributions for
\(h*h^\sharp\), both equal to \(-m\). Its trace is
\(-2m<0\), contradicting the left side of (M19.2).
This proves the equivalence.
\end{proof}

This is a precise zero-side Weil positivity criterion in
the canonical quotient, with an explicit unconditional
nuclear and trace-class realization. It supplies no proof
of positivity independent of the zeros. The trace is the
sum of traces on finite-length spectral algebras; multiple
zeros contribute their lengths. In the programme of
cote~1 and SA23 geo-H-lidskii/geo-H-genus, MT-M21--MT-M22 identify the specified Meyer adelic summation
quotient with this model, and MT-M24 identifies the actual
Masullo smooth quotient. MT-M25 treats the actual line Hilbert
completion and its lost jets. A cohomological realization and
a positivity argument defined independently of the zeros remain separate. Adopted \(A_D\)
supplies (M19.2)'s positive conclusion for Dirichlet
functions; \(M_H\) supplies it conditionally in general.

\subsection*{MT-M20. A weaker positive Hilbert completion and its exact obstruction}
Keep the unconditional canonical quotient of MT-M18.
If all nontrivial zeros are simple, individual Mellin
evaluation induces a continuous injective map with dense
image
\[
       Q\longrightarrow\ell^2(Z),\qquad
       [g]\longmapsto(\mathcal Mg(\rho))_{\rho\in Z}.   \tag{M20.1}
\]
It therefore defines a positive-definite continuous Hilbert
norm on \(Q\). Every convolution extends diagonally to a
trace-class operator on this completion. Under GRH,
\(C_{h^\sharp}=C_h^*\), and \(C_{h*h^\sharp}\) is a
positive operator. No reciprocal-derivative bound is
needed for this weaker Hilbert norm.

More precisely, the following are equivalent:
\begin{enumerate}
\item GRH holds for \(L(s,\eta)\), and every nontrivial
zero is simple.
\item There is a positive-definite continuous inner
product on \(Q\) such that every convolution extends to
a bounded operator on its Hilbert completion and every
\(C_{h*h^\sharp}\) is positive there.
\end{enumerate}
The inner product is required to be faithful on the
\emph{full jet quotient}. Semisimplifying its multiple-zero
local algebras changes this requirement.

\begin{proof}
Point evaluation at any zero in \(\Omega_n\) has norm at
most one on the quotient \(J_n\), by the maximum principle
and the definition of its quotient norm. Thus
\[
 \sum_{\rho\in Z}|\mathcal Mg(\rho)|^2
 \le\sum_n d_n\|(Rg)_n\|_n^2
 \le q_2(Rg)^2\sum_n d_n(2+|n|)^{-4}<\infty.           \tag{M20.2}
\]
This proves continuity. When all zeros are simple, zero
values mean all zero jets, and MT-M18 identifies the kernel
exactly with the closed ideal; hence the map is injective.
Every finitely supported zero-value vector lies in its
image, since it specifies elements of finitely many
\(J_n\)'s and the interpolation right inverse realizes
them. This proves density and (M20.1).

Convolution by \(f\) has diagonal entries
\(\mathcal Mf(\rho)\). Their absolute sum is finite by
the local zero count and rapid strip decay, so its extension
is trace class. Under GRH (M19.6) makes the entries for
\(h^\sharp\) the complex conjugates of those for \(h\).
The adjoint and positivity statements follow. This proves
the implication (1) to (2).

For the converse, finite-support elements of the block
algebras are actual elements of \(Q\), by MT-M18, and
convolution realizes their algebra multiplication. If
\(\rho\ne\rho^\dagger\), let \(e,e^\dagger\) be the
local idempotents equal to 1 in the full local algebras
at those two zeros and zero in every other local algebra.
The involution interchanges them. For \(h=e-e^\dagger\)
one has
\[
                  h*h^\sharp=-(e+e^\dagger).
\]
On the nonzero finite-dimensional invariant subspace
\((e+e^\dagger)Q\), convolution by this element is
\(-\mathrm{id}\). Its quadratic form is strictly negative
on every nonzero vector in any positive-definite inner
product. This contradicts (2), and proves GRH.

Now suppose a critical-line zero has multiplicity
\(m\ge2\). In its local algebra
\(\mathbb C[z]/(z^m)\), the involution is
\(z^\sharp=-z\), by (M19.6). Let \(e\) be that local
idempotent and \(n\) the element represented by \(z\),
both zero in all other local factors. They are realized
by elements of \(Q\). On the finite-dimensional invariant
space \(eQ\), write \(N=C_n\). It is a nonzero nilpotent
operator. For real \(u\), the element \(h_u=e+iun\)
satisfies \(h_u^\sharp=h_u\), and
\[
 C_{h_u*h_u^\sharp}|_{eQ}
                     =I+2iuN-u^2N^2.                \tag{M20.3}
\]
By (2) these operators are positive, in particular
self-adjoint, for both \(u=1\) and \(u=-1\).
Subtracting them shows that \(iN\) is self-adjoint
in the restricted positive-definite inner product.
A self-adjoint nilpotent on a finite-dimensional
inner-product space is zero: diagonalize the Hermitian
matrix, whose real eigenvalues must all vanish.
This contradicts \(N\ne0\). Thus all zeros are simple,
completing the converse.
\end{proof}

The Hilbert norm in (M20.1) can be strictly weaker than
the original Fr\'echet topology. Its existence therefore
does not make point evaluations a Schauder coordinate
system for that topology: MT-M16 still gives the exact
extra leading-coefficient condition for that assertion.
With multiple zeros, evaluation kills the higher-jet
nilpotent directions. A semisimple Hilbert model with
\(m_\rho\) copies of each scalar eigenvalue retains
the multiplicity-weighted trace, but is not a faithful
Hilbert completion of this full jet quotient. No
simplicity consequence of \(A_D\) or \(M_H\) is claimed.


\subsection*{MT-M21. Identification with Meyer's zeta-summation quotient over \(\mathbb Q\)}
For the Riemann zeta function, the canonical quotient in MT-M18
is naturally isomorphic, as a topological convolution module,
to the zero-space in Meyer's classical spectral realization.
This identification is unconditional and includes all zero
multiplicities.

Here are the precise source and conventions. In Ralf Meyer,
\emph{A spectral interpretation for the zeros of the Riemann
zeta function}, G\"ottingen Seminars 2004--05, pp.~117--137,
let \(H_- =\mathcal O(\mathbb R_+^\times)\) be the strong
Schwartz space of Section~2, let \(H_+\) be the even Schwartz
space on \(\mathbb R\), and let
\[
 H_\cap=\{u\in H_+:u(0)=\widehat u(0)=0\},\qquad
 (Zu)(x)=\sum_{k\ge1}u(kx).
\]
His Theorem~3.3, printed p.~124 (PDF p.~8), proves the
closed-image assertion for the zeta summation map; Section~4,
printed pp.~125--126, defines
\(H_-^0=H_-/ZH_\cap\).
Theorem~4.1, printed p.~126 (PDF p.~10), characterizes its
Mellin-transform kernel by two half-plane division conditions.
Under the logarithmic coordinate map
\[
       \mathcal L:\mathcal E\longrightarrow H_-,
       (\mathcal Lf)(x)=f(\log x),
\]
one has the exact equality of closed subspaces
\[
          \mathcal L\bigl(\overline{\psi_\zeta*\mathcal E}\bigr)
                         = ZH_\cap.                  \tag{M21.1}
\]
Consequently \(Q_\zeta\simeq H_-^0\), and the adapted
absolute basis, nuclear convolution, and trace constructions
of MT-M18--MT-M20 apply to this independently defined
zeta-summation quotient with their stated positivity hypotheses.

\begin{proof}
The map \(\mathcal L\) is a topological isomorphism:
the invariant derivatives \((x\,d/dx)^j\) become
\(d^j/dt^j\) for \(t=\log x\), and requiring rapid
decay after every real power weight \(x^\sigma\)
is equivalent to the defining seminorms
\(\sup_t e^{N|t|}|f^{(j)}(t)|\). Polynomial factors in
\(t\), if included in the Schwartz definition, are
absorbed by increasing \(N\). Its Mellin transform is
exactly \(\mathcal Mf(s)\), since \(d^\times x=dt\).
Multiplicative convolution becomes additive convolution.

Let \(F=\mathcal Mf\). Meyer's Theorem~4.1 says that
\(\mathcal Lf\in ZH_\cap\) precisely when
\[
 t\longmapsto F(\sigma+it)/\zeta(\sigma+it)
          \text{ is Schwartz for every }\sigma\ge\tfrac12,
                                                               \tag{M21.2}
\]
and
\[
 t\longmapsto F(\sigma+it)/\zeta(1-\sigma-it)
          \text{ is Schwartz for every }\sigma\le\tfrac12.
                                                               \tag{M21.3}
\]
The conditions include holomorphic removal of the corresponding
zeros on those lines. If they hold, every nontrivial zero
\(\rho\) of real part at least \(1/2\) is a zero of \(F\)
to at least its full multiplicity, by (M21.2). For zeros in
the other half-plane, the zeta functional equation says that
\(\zeta(1-s)\) has the same nontrivial zero divisor as
\(\zeta(s)\); condition (M21.3) gives the same conclusion.
Thus \(f\) belongs to the full zero-jet kernel in MT-M18.

Conversely, suppose \(F\) vanishes on this entire divisor.
The function \(F/\zeta\) is holomorphic on
\(\operatorname{Re}s>-2\). All nontrivial zero singularities
are removed, the pole of \(\zeta\) at 1 gives a removable
zero of its reciprocal, and its first trivial zero is at
\(-2\), outside this open half-plane.

Fix \(A\ge2\). Use rectangles
\([-1,A+1]+i[h_n,h_{n+1}]\), with the horizontal levels
constructed in MT-M18 for \(\zeta\). The finite-\(L\)
Blaschke estimates in that proof give a fixed inverse
polynomial lower bound for \(|\zeta|\) on the horizontal
edges at large height. The two vertical edges are a fixed
positive distance from every nontrivial zero, so the same
disk argument gives their polynomial lower bound. They
also avoid all trivial zeros. The finitely many low-height
boundaries cause no problem: the levels \(h_n\) are never
zero, and both vertical lines are zero-free.

For every \(N\), rapid fixed-strip decay of \(F\)
therefore bounds \(|F/\zeta|\) on these boundaries by
\(C_{A,N,f}(2+|n|)^{-N}\), after increasing the decay
power used for \(F\). The maximum principle for the
holomorphic quotient gives the same bound inside each
rectangle. Since adjacent heights have uniformly bounded
gaps, \(F/\zeta\) decreases rapidly on the whole strip
\([-1,A+1]\). Cauchy's formula on small fixed disks
inside that strip gives rapid decay of all its vertical
derivatives on \([1/2,A]\). This proves (M21.2).

Apply this same argument to \(F^\vee(s)=F(1-s)\).
It is an entire strip-Schwartz function and vanishes on
the full nontrivial divisor, because that divisor is
invariant under \(s\mapsto1-s\). The result for
\(F^\vee/\zeta\), followed by \(s\mapsto1-s\),
is exactly (M21.3). Meyer's image criterion now gives
\(\mathcal Lf\in ZH_\cap\).

We have proved that \(\mathcal L^{-1}(ZH_\cap)\)
is precisely the full zero-jet kernel. MT-M18 identifies
that kernel with \(\overline{\psi_\zeta*\mathcal E}\),
proving (M21.1). Both are closed subspaces of the
topologically identified test-function spaces, so their
quotients, with quotient topology, are naturally
topologically isomorphic. The logarithmic identification
commutes with scaling and integrated convolution, giving
the stated module assertion.
\end{proof}

This settles a specified arithmetic comparison map, beyond
a zero-built abstract model. It is a comparison with
Meyer's Poisson/zeta-summation construction, not an
identification with an unspecified adelic quotient in
Masullo's draft or with an \(\ell\)-adic cohomology group.
The classical spectral realization itself is Meyer's
result; the proof here supplies the explicit closed-ideal
and stable-block crosswalk used in this investigation.
The full number-field idele-class construction is located
in Meyer, \emph{On a representation of the idele class
group related to primes and zeros of L-functions},
Duke Math. J.~127 (2005), 519--595, author version
\path{math/0311468v3}, Theorem~5.8 (pp.~43--44) and
Theorem~5.11 (p.~47). A sectorwise topological comparison
with that construction requires its precise image
criterion and is not inferred merely from equal spectra.


\subsection*{MT-M22. The fixed Hecke sector of Meyer's adelic zero-space}

For every fixed primitive unitary Hecke \(L(s,\eta)\), the
canonical quotient \(Q_\eta\) of MT-M18 is isomorphic to the
corresponding fixed compact-character sector of Meyer's
idele-class zero-space.  The isomorphism is unconditional
and compatible with convolution, including all zero
multiplicities.  The sector is first given its quotient
Fr\'echet topology from \(H_-^\chi\); the proof also
identifies the quotient bornology by an explicit splitting.

We fix the character convention explicitly.  Write
\[
 L(s,\eta)=\prod_{\mathfrak p}
       (1-c_{\mathfrak p}(N\mathfrak p)^{-s})^{-1},
\]
where \(c_{\mathfrak p}=0\) at a ramified place.  Let
\(\beta\) be the associated unitary idele-class character
in Meyer's convention: for his \(S\)-prime element
\(p_v\in C_S\), at every unramified \(v\notin S\),
\[
                      \beta(p_v)^{-1}=c_{\mathfrak p_v}.
                                                               \tag{M22.1}
\]
Thus the component is parameterized by
\(\omega_s=\beta|\cdot|^s\).
This convention is dictated by the inverse in Meyer's
Euler factor, rather than by a choice of the name
\(\eta\) or \(\eta^{-1}\) for its idele realization.
For example, \(L(s,\eta)=\zeta_K(s+i\tau)\) corresponds
to \(\beta=|\cdot|^{i\tau}\).
If an idele character has instead been labeled by
its direct values \(c_{\mathfrak p_v}\) on Meyer's
\(p_v\), it must be inverted to obtain \(\beta\).

Let \(C=C_K\), \(C^1=\ker|\cdot|\), and choose a
continuous smooth splitting \(a:\mathbb R\to C\)
with \(|a(t)|=e^t\).
Normalize Haar measure on \(C^1\) to one, so
\(d^\times x=dc\,dt\) in \(x=ca(t)\).
Put \(\chi=\beta|_{C^1}\).
Use Meyer's spaces and embeddings
\[
 H_-=\mathcal S(C)_{\mathbb R},\qquad
 H_+=\mathcal S(\mathbb A_K)/K^\times,\qquad
 H_-^0=H_-/(H_+\cap H_-),
\]
where the intersection uses his embeddings into
\(\mathcal S(C)_{]1,\infty[}\oplus
  \mathcal S(C)_{]-\infty,0[}\).
Superscript \(\chi\) denotes the subspace on which
left translation by \(c\in C^1\) acts by \(\chi(c)\).
Then
\[
 (\mathcal L_\beta f)(x)=\beta(x)^{-1}f(\log|x|)
\]
induces an isomorphism
\[
              Q_\eta \ \simeq\ (H_-^0)^\chi.                  \tag{M22.2}
\]
It is topological for the specified quotient topology.
Norm translation by \(t\) corresponds to idele
translation by \(a(t)\), with the explicit unitary
factor \(\beta(a(t))\).

The primary source is Ralf Meyer, \emph{On a representation
of the idele class group related to primes and zeros of
\(L\)-functions}, Duke Math. J.~127 (2005), 519--595,
author version \path{math/0311468v3}.
The exact input is the weighted spaces in Definition 4.1
and (22), author pp.~24--25; the embeddings in Lemma 5.4,
pp.~40--41; and the image criterion in \S5.7, pp.~46--47,
following Lemma 5.10.  Theorem 5.1, pp.~38--39, supplies
the finite-level description used there.
Theorems 5.8, pp.~43--44, and 5.11, p.~47, are the
classical summability and spectral-multiplicity statements.
The comparison below uses the image criterion,
rather than inferring an isomorphism from equal spectra.

\begin{proof}
\emph{The exact test-function sector.}
A norm splitting exists by varying a positive component
at one archimedean place with the appropriate local
norm exponent.  In the resulting product
\(C=C^1\times a(\mathbb R)\), every \(\chi\)-homogeneous
function has the form \(\chi(c)^{-1}u(t)\).
Since \(\beta\) is unitary,
\(\beta(a(t))=e^{i\tau t}\) for a fixed real \(\tau\).
Multiplication by this fixed modulation preserves
every exponential-Schwartz seminorm, in both
directions, by the Leibniz formula.

Consequently \(\mathcal L_\beta\) is a topological
isomorphism from \(\mathcal E\) onto \(H_-^\chi\).
Indeed the norm-direction invariant derivative becomes
\(d/dt\); derivatives in the compact direction act
on its one fixed smooth character by constants.
All norm weights are \(e^{\alpha t}\), and the
requirement for every real \(\alpha\) is equivalent
to the seminorms \(e^{N|t|}|f^{(j)}(t)|\).
Polynomial factors in \(t\) are absorbed by increasing
\(N\).  The character is trivial on a suitable
compact subgroup with elementary compact quotient,
so these functions occupy one fixed Bruhat level.
There is no limit over varying compact characters.

Haar normalization as in Meyer's \S3.1, p.~16,
gives exactly
\[
 \widehat{\mathcal L_\beta f}(\omega_s)
       =\int_C\beta(x)^{-1}f(\log|x|)
                      \beta(x)|x|^s\,d^\times x
       =\mathcal Mf(s).                                        \tag{M22.3}
\]
The transform is zero on all other compact-character
components by orthogonality on \(C^1\).
A direct convolution calculation, using
\(\beta(y)^{-1}\beta(y^{-1}x)^{-1}=\beta(x)^{-1}\),
shows
\[
  (\mathcal L_\beta f)*(\mathcal L_\beta g)
                         =\mathcal L_\beta(f*g).
\]
Also
\[
 \lambda_b\mathcal L_\beta f
   =\beta(b)\mathcal L_\beta
         \bigl(f(\,\cdot-\log|b|)\bigr),\qquad b\in C.            \tag{M22.4}
\]
These identities specify the representation convention.

\emph{Meyer's image condition on this component.}
Choose once and for all a sufficiently large finite
set \(S\), containing the archimedean places, all
ramification of \(\beta\), and the places needed for
his normalized additive character and differential
idele.  His multiplier \(\mathcal L_S\) has transform
\[
 \widehat{\mathcal L_S}(\omega)
       =\prod_{v\notin S}(1-\omega(p_v)^{-1})^{-1}.
\]
By (M22.1), on \(\omega_s\) this is the partial finite
Hecke function
\[
 L_\eta^S(s)
       =L(s,\eta)D_{\eta,S}(s),\qquad
 D_{\eta,S}(s)=
        \prod_{v\in S_{\mathrm{fin}}}
             (1-c_{\mathfrak p_v}(N\mathfrak p_v)^{-s}).         \tag{M22.5}
\]
A ramified factor in \(D_{\eta,S}\) is one.
The reflected quasi-character is
\[
 \omega_s^*=|\cdot|\omega_s^{-1}
                    =\beta^{-1}|\cdot|^{1-s},
\]
and therefore
\[
              \widehat{\mathcal L_S}(\omega_s^*)
                       =L_{\bar\eta}^S(1-s).                   \tag{M22.6}
\]
Here \(\bar\eta\) denotes the dual unitary Hecke
character, with conjugate Euler coefficients.

For an entire strip-Schwartz function \(F(s)\)
on this component, the equivalent reduced
image criterion printed on Meyer p.~47 is
\[
 \begin{array}{ll}
 F(\alpha+it)/L_\eta^S(\alpha+it)
        \text{ is Schwartz in }t,&0<\alpha\le1,\\[2mm]
 F(it)/L_{\bar\eta}^S(1-it)
        \text{ is Schwartz in }t.&
 \end{array}                                                    \tag{M22.7}
\]
All apparent zero singularities must be removable.
The source's remaining conditions for
\(\alpha>1\) and \(\alpha<0\) follow from \(F\)
being the transform of \(H_-\).
Zero extension to other compact-character
components causes no additional Schwartz condition.
Thus (M22.7) characterizes membership of
\(\mathcal L_\beta f\) in \(\Sigma(H_+)\),
where \(F=\mathcal Mf\).

This is also the relevant intersection in
the quotient, rather than only its first-coordinate
image.  Meyer's Poisson formula gives
\[
 \Sigma u-J\Sigma\mathfrak F^*u
                         =c_0-c_1|\cdot|^{-1}.
\]
The second term lies in
\(\mathcal S(C)_{]-\infty,0[}\).
If \(\Sigma u\in H_-\), so does the right side.
Its decay as \(|x|\to0\) forces \(c_0=c_1=0\):
these distinct powers cannot decay after every
negative weight.  Hence
\(i_+u=i_-(\Sigma u)\).
Conversely, an element of the intersection has
its first coordinate in \(\Sigma(H_+)\cap H_-\).
Compact-character projection commutes with the
maps.  Therefore (M22.7) characterizes
\((H_+\cap H_-)^\chi\).

\emph{Removing the image criterion's growth gate.}
If (M22.7) holds, every zero of \(\Xi_\eta\) in
\(0<\operatorname{Re}s\le1\) is a zero of \(F\)
with at least its full multiplicity, by its first
condition.  The omitted finite Euler factors have
no zeros in that right half-plane, and the gamma
factors have neither zeros nor poles there.
The possible pure-norm pole is cancelled in
\(\Xi_\eta\) and is not a zero of it.
If a nontrivial zero is on \(\operatorname{Re}s=0\),
the functional equation identifies its multiplicity
with that of \(L_{\bar\eta}(1-s)\); the second
condition gives the same conclusion.
Thus membership implies all the zero-jet
vanishings of MT-M18.

Conversely suppose \(F=\mathcal Mf\) vanishes
on the full divisor of \(\Xi_\eta\).
Set \(G=F/\Xi_\eta\), which is entire.
The rectangle proof of (M18.14) gives,
for every \(A,N>0\), the sharpened estimate
\[
 \sup_{|\sigma|\le A}|G(\sigma+it)|
       \le C_{A,N,f}e^{\kappa|t|}(2+|t|)^{-N},
 \qquad \kappa=\pi[K:\mathbb Q]/4.                              \tag{M22.8}
\]
Indeed the boundary lower bound for \(\Xi_\eta\)
costs \(e^{\kappa|n|}(2+|n|)^{B_A}\).
Use the decay power \(N+B_A\), rounded up and
increased if needed, for \(F\) on that boundary.
The maximum principle and bounded block heights
give (M22.8), with no zero-spacing hypothesis.

Write the standard completion as
\[
 \Xi_\eta(s)=P_\eta(s)q_\eta^{s/2}
                  \Gamma_\eta(s)L(s,\eta).
\]
On each closed strip in \(\operatorname{Re}s>0\),
the unitary archimedean gamma product is
holomorphic and bounded above by
\(e^{-\kappa|t|}\) times a fixed power
of \(2+|t|\), by Stirling's formula.
Consequently
\[
        F(s)/L(s,\eta)
          =G(s)P_\eta(s)q_\eta^{s/2}\Gamma_\eta(s)
\]
is holomorphic and rapidly decreasing there,
by (M22.8).  The reciprocal of a finite-\(L\)
pole is a zero, so pure-norm poles cause
no exception.

The completed functional equation gives
\[
 \Xi_\eta(s)
   =\epsilon_\eta P_\eta(s)q_\eta^{(1-s)/2}
            \Gamma_{\bar\eta}(1-s)L(1-s,\bar\eta),
                    \qquad|\epsilon_\eta|=1.
\]
It follows that \(F(s)/L(1-s,\bar\eta)\)
is holomorphic and rapidly decreasing
on each closed strip in
\(\operatorname{Re}s<1\).
The reflected gamma arguments have positive
real part there.  Thus neither half-plane
division asks \(F\) to cancel trivial
finite-\(L\) zeros outside the appropriate half.

For \(\delta>0\), every nontrivial omitted
Euler factor satisfies
\[
 |1-c_{\mathfrak p}(N\mathfrak p)^{-s}|
                         \ge1-(N\mathfrak p)^{-\delta}
       \quad(\operatorname{Re}s\ge\delta).
\]
The ramified factors are one.
Thus \(D_{\eta,S}^{-1}\) is bounded on each
such right half-plane, as is
\(D_{\bar\eta,S}(1-s)^{-1}\) when
\(\operatorname{Re}s\le1-\delta\).
Their derivatives are bounded on smaller
fixed strips.  Therefore
\[
 \begin{split}
 F/L_\eta^S&\text{ is strip-Schwartz on }\operatorname{Re}s>0,\\
 F/L_{\bar\eta}^S(1-s)&\text{ is strip-Schwartz on }
                                      \operatorname{Re}s<1.
 \end{split}                                                    \tag{M22.9}
\]
Cauchy's formula on slightly larger strips
gives every vertical derivative estimate.
These statements imply (M22.7).
We have proved
\[
 \mathcal L_\beta^{-1}\bigl((H_+\cap H_-)^\chi\bigr)
   =\{f:(\mathcal Mf)^{(j)}(\rho)=0,\
          \Xi_\eta(\rho)=0,\ j<m_\rho\}
   =\overline{\psi_\eta*\mathcal E},                             \tag{M22.10}
\]
where the last equality is MT-M18.

\emph{The precise quotient topology and bounded lifting.}
Give \(H_-^\chi\simeq\mathcal E\) its
just-identified Fr\'echet topology and its
bounded-set bornology.  At this fixed character
these are precisely the source's weighted
Schwartz structure: it lies at one Bruhat
level, and the countable family of real
weight endpoints in (22) generates its
seminorms and bounded sets.
Here bounded and precompact sets coincide,
so the convention of using the precompact
bornology gives the same result.  Indeed,
for a bounded sequence in \(\mathcal E\),
the bounds on one additional derivative
give equicontinuity of every derivative
on each compact interval.  Diagonal
Arzel\`a--Ascoli gives a subsequence
converging with every derivative on every
compact interval.  For a prescribed
\(p_{N,j}\), the uniform \(p_{N+1,j}\)
bound makes the weighted tails outside
\([-R,R]\) at most \(Ce^{-R}\).
The subsequence is therefore Cauchy in
every defining seminorm and converges
in the complete space \(\mathcal E\).
This proves relative compactness of every
bounded set in its metrizable topology;
the converse is immediate.
The subspace in (M22.10) is closed.
Hence (M22.10) immediately gives a
topological isomorphism from \(Q_\eta\)
onto the sector quotient furnished with
the quotient topology from \(H_-^\chi\).

Here is why the original quotient bornology
agrees as well.  The maps \(i_+\) and \(i_-\)
are bornological embeddings by the cited
source.  Therefore the bornology of their
intersection equals the subspace bornology
inherited from \(H_-^\chi\): a subset of the
intersection is bounded in either embedded
space exactly when it is bounded in the
common ambient space.  No topology on
\(H_+^\chi\), or on its quotient, is being
assumed in this argument.

Compact-character averaging is a bounded
projection preserving \(H_+\) and \(H_-\).
At each finite level it is ordinary integration
over a compact group against its smooth
character, and these projections agree under
the level inclusions.
It preserves their intersection and splits
both spaces into its image and kernel.
The algebraic quotient and its quotient
bornology therefore split accordingly.
In particular \((H_-^0)^\chi\) is the
bornological quotient
\(H_-^\chi/(H_+\cap H_-)^\chi\);
representatives can be projected before
passing to the quotient.

Finally the continuous linear section
\(S:s(J)\to\mathcal E\) in MT-M18 supplies
an explicit lifting map for this quotient,
after transport by \(\mathcal L_\beta\).
Every bounded set in the quotient
Fr\'echet space lifts to a bounded set
through this section.  Conversely the
quotient map is continuous, so it takes
bounded sets to bounded sets.
The same two assertions hold for
precompact sets, using continuity of
the section and quotient map.
The quotient itself has bounded sets
precompact, since its continuous section
maps them to bounded, hence precompact,
subsets of \(\mathcal E\), and its
quotient map sends these back to the
original sets.
Thus the source's quotient bornology
is exactly the bornology of the specified
Fr\'echet quotient.  This is an explicit
splitting argument, not an inference
from equality of spectra.
Together with (M22.3)--(M22.4) and
multiplicativity, it proves all parts
of (M22.2).
\end{proof}

The stable blocks and nuclear traces in
MT-M18--MT-M19 therefore apply to this
independently defined adelic summation
quotient in each fixed Hecke sector.
All constants may depend on the character
and its infinity type; no uniform
decomposition over all characters is
asserted.  The trace positivity and
faithful-operator distinctions of
MT-M19--MT-M20 retain their exact GRH
and multiplicity hypotheses.
This identifies Meyer's specified
quotient, rather than an unspecified
quotient in Masullo's draft.
Its geometric location is the norm
direction of the idele-class action,
Tate's Poisson summation, and cote~1's
nuclear-trace framework as carried in
SA23 geo-H-lidskii/geo-H-genus.
The primary checks cover the named
normalization, image, and embedding
passages; no complete-paper audit
is claimed.


\subsection*{MT-M23. Every nonabelian Galois top field obstructs weak Mertens and faithful positivity}

Let $L/K$ be a nonabelian Galois extension of number fields. The
classical theorem of Hu--Kaneko--Martin--Schildkraut (HKMS),
\emph{Order of zeros of Dedekind zeta functions}, Theorem1.1,
author version2107.03269v2 p.2, proves unconditionally that
$\zeta_L$ has infinitely many nontrivial zeros of multiplicity
at least two. The complete first proof is \S3 p.5; a second
proof occupies \S4 pp.5--6. Their Theorem1.2, p.2 with proof
\S5 pp.6--8, gives infinitely many zeros of multiplicity at
least three whenever $\operatorname{Gal}(L/K)$ has an irreducible
representation of degree at least three. These are cited classical
results; no new multiplicity theorem is claimed here.

Set
\[
 M_L(x)=\sum_{N\mathfrak a\le x}\mu_L(\mathfrak a),\qquad
 I_L(X)=\int_1^X|M_L(x)/x|^2\,dx.
\]
The following consequences hold without either adopted GRH input:
\[
 \limsup_{X\to\infty}\frac{I_L(X)}{(\log X)^3}>0,
 \qquad
 \limsup_{X\to\infty}\frac{I_L(X)}{\log X}=+\infty.
 \tag{MT-M23a}
\]
If the group has an irreducible representation of degree at least
three, the exponent $3$ in the first assertion can be replaced by
$5$. The canonical full-jet quotient of $\zeta_L$ has no faithful
continuous Hilbert completion on which every convolution extends
boundedly and every $C_{h*h^\sharp}$ is a positive operator.
Nevertheless its nuclear trace construction and adapted absolute
basis exist unconditionally by MT-M18--MT-M19, and $M_H$ makes
the traces of these adjoint-square convolutions nonnegative.

\paragraph{Complete proof of the consequences.}
Choose a multiple zero supplied by HKMS, of actual multiplicity
$m\ge2$, or $m\ge3$ in the stronger case. If $\zeta_L$ has any
off-critical zero, the functional equation and conjugation give
a zero with real part greater than $1/2$. The last assertion of
NT-L16 then gives infinite limsup after division by every fixed
nonnegative power of $\log X$; all displayed conclusions follow.
Otherwise all nontrivial zeros are critical, and our chosen zero
$\rho$ satisfies the exact NT-L16 bound
\[
 \limsup_{X\to\infty}\frac{I_L(X)}{(\log X)^{2m-1}}
 \ge
 \frac{1}{(2m-1)((m-1)!)^2
       |\rho\,\zeta_L^{(m)}(\rho)/m!|^2}>0.
\]
When $2m-1$ equals the claimed exponent this is the first assertion.
When it is greater, take an unbounded sequence on which the
displayed ratio is bounded below by a fixed positive constant;
multiplying by the remaining positive power of $\log X$ gives
infinite limsup. The same argument with exponent one proves
the second assertion. This dichotomy uses no unproved zero
location. Finally MT-M20 proves that a faithful positive Hilbert
realization with the stated properties forces every nontrivial
zero to be simple. The actual multiple zeros contradict that
necessary condition. MT-M18--MT-M19 give the separate basis and
trace statements, and their GRH equivalence gives the conditional
trace positivity under $M_H$.

For the programme this is a finite-monodromy Artin-motive
obstruction to discarding nilpotent spectral blocks. It sharpens
the Kummer examples NT-L17 to all nonabelian Galois top fields.
The cote1/SA23 \path{geo-H-lidskii} trace route must retain the
difference between positivity of a scalar trace and positivity
of every operator in a faithful full-multiplicity Hilbert model.
HKMS's general multiplicity prediction for representations of
degree at least four remains their Conjecture6.2, not a theorem
used here. A non-Galois field is covered only when it is itself
the top field of a nonabelian Galois extension of some subfield.


\subsection*{MT-ART09. A smaller relative-zero gate for solvable Artin holomorphy}

Let $L/K$ be a finite solvable Galois extension, with group $G$,
and let $s_0\ne1$. Write
\[
 q(s_0)=\operatorname{ord}_{s=s_0}
             \frac{\zeta_L(s)}{\zeta_K(s)}\ge0.
\]
The nonnegativity follows from the classical Aramata--Brauer
theorem. When $|G|$ has at least two distinct prime divisors,
let $p_2$ be the second smallest and put
\[
                 B(G)=\max\{2,p_2-2\}.
\]
Then every Artin $L(s,\chi,L/K)$ is holomorphic at $s_0$ if
$q(s_0)\le B(G)$. This is the combination of Peng-Jie Wong,
\emph{Zeros of Dedekind zeta functions and holomorphy of Artin
$L$-functions}, Theorems1.3 and1.4, printed p.254, with their
complete proofs on pp.259 and260. The solvability hypothesis
is essential to the statements as cited here. In a $p$-group
all irreducible characters are monomial, so Artin holomorphy
already follows globally from Hecke holomorphy.

For $4\nmid |G|$ the sufficient bound $q(s_0)\le3$ holds,
without adding a solvability assumption: Wong Corollary5.3,
printed pp.259--260, includes the group-theoretic implication
that these groups are solvable. His Corollary5.1, p.259, is the
odd-degree special case. One may therefore replace $B(G)$ by
$\max\{B(G),3\}$ when applicable. The stronger Sylow-structure
variants in his Theorems5.2 and5.4 are separately located there;
no unstated prime-factor hypothesis from them is being used.

\paragraph{Complete transfer to the adopted Hecke hypothesis.}
Under $M_H$, a possible pole of an Artin function in
$0<\Re s<1$ must lie on $\Re s=1/2$: Brauer induction writes
that function as a finite quotient of Hecke functions, whose
possible nontrivial denominator zeros are all on that line.
This is the MT-ART01 confinement, and does not assert Artin
holomorphy. The contrapositive of the just-cited local theorem
now proves that a possible pole must satisfy
\[
 \Re s_0=\tfrac12,\qquad q(s_0)\ge B(G)+1.
 \tag{MT-ART09a}
\]
If $4\nmid |G|$, use the strengthened bound at least four.
Thus, if the relative zeta function has order at most $B(G)$
at every point of the critical line, all Artin functions are
holomorphic throughout the open critical strip. This last
zero-order condition is an extra input, not supplied by $M_H$.
For an unconditional local conclusion at a particular $s_0$,
only the corresponding relative-order inequality is needed.

This is a strictly finer criterion than measuring the order
of $\zeta_L$ alone: by the definition of a meromorphic order,
\[
 q(s_0)=\operatorname{ord}_{s_0}\zeta_L
                 -\operatorname{ord}_{s_0}\zeta_K.
\]
The base contribution is removed exactly, including when the
base has a multiple zero. It is not necessary to posit any
noncoincidence of zeros. The argument is a complete input
transfer to the cited classical local theorem, not a new proof
of that theorem or of general Artin holomorphy.

Geometrically, $\zeta_L/\zeta_K$ is the $L$-function of the
augmentation Artin motive attached to the finite etale cover
after deleting ramified primes, with the standard inertia
factors restored. The relative virtual permutation class is
$[\mathbb Q[G]]-[\mathbb Q]$ in its representation $K_0$.
This belongs to the cote1 trace and SGA5 finite-correspondence
framework already fixed in AG-C1, and sharpens the
solvable part of Q-ARTIN-POS beyond MT-ART02. It does not turn a
virtual-character inequality into coefficientwise positivity.


\subsection*{MT-M24. The specified Masullo smooth quotient is the Meyer quotient}
This is a comparison of the actual test spaces and summation maps. It does
not assume Masullo's global converse theorem, his spectral basis claims,
or $M$. Fix the primitive unitary Hecke function and the convention
$\beta(\mathfrak p_v)^{-1}=c_{\mathfrak p_v}$ from MT-M22; put
$C=C_K$, $C^1=\ker|\cdot|$, and $\chi=\beta|_{C^1}$.
Use the norm splitting and volume-one measure on $C^1$ chosen there.
Write $\mathcal E$ for the exponential Schwartz space of MT-M18,
$\mathcal I_\eta=\overline{\psi_\eta*\mathcal E}$ for its closed
completed-Hecke ideal, and $Q_\eta=\mathcal E/\mathcal I_\eta$.
For a real number $a$, let $M_a u(g)=|g|^a u(g)$, and define
\[
 \Theta_\chi=\Phi_{K,\chi}\bigl(\mathcal S_0(\mathbf A_K)\bigr),
 \qquad
 Q^{\mathrm{Mas}}_{\chi,a}
 =\mathcal S_{\mathrm E}(C)(\chi)\big/
       \overline{M_a\Theta_\chi}^{\,\mathcal S_{\mathrm E}(C)(\chi)}.
 \tag{M24.1}
\]
Here the symbols on the right have their literal meanings in Masullo:
\[
 \mathcal S_0(\mathbf A_K)
  =\{f\in\mathcal S(\mathbf A_K):f(0)=\widehat f(0)=0\},
 \qquad
 \Phi_K(f)(g)=\sum_{\xi\in K^\times}f(\xi g),
 \tag{M24.2}
\]
and $\Phi_{K,\chi}$ is compact-character averaging of $\Phi_K$.
The closure in (M24.1) does not enlarge $M_a\Theta_\chi$.
There is an explicit topological isomorphism
\[
 Q_\eta\xrightarrow{\ \sim\ }Q^{\mathrm{Mas}}_{\chi,a},
 \qquad
 [f]\longmapsto
 \bigl[g\longmapsto |g|^a\beta(g)^{-1}f(\log|g|)\bigr].
 \tag{M24.3}
\]
For $a=0$ this identifies the quotient in Masullo (7.1.4.7),
original PDF p.~327. For $a=1/2$ it identifies the centered quotient
in (1.8.3.1), original PDF p.~45. These two source displays use
different twists; (M24.3) keeps the twist explicit.

\begin{proof}
First verify the test-space topology rather than just its elements.
Masullo Definition D.3.3.1, original PDF p.~804, requires
$|\cdot|^b u\in\mathcal S(C)$ for every real $b$.
Definition D.3.3.3, p.~805, gives each fixed compact-invariance
level the projective topology of all these weighted Schwartz spaces,
and gives their union the locally convex inductive topology.
On the $\chi$-sector the map
\[
 \mathcal L_\beta f(g)=\beta(g)^{-1}f(\log|g|)
 \tag{M24.4}
\]
is a linear bijection from $\mathcal E$ onto that sector.
Indeed, the $\chi$-eigencondition for left translation says
$u(cg)=\chi(c)^{-1}u(g)$ for $c\in C^1$; its restriction to
the norm splitting therefore determines $u$. The unitary value
of $\beta$ on that splitting is $e^{i\tau x}$ for a fixed real
$\tau$, so removing it only changes derivatives by fixed finite
linear combinations.

Every function in this sector is invariant under $\ker\chi$.
The compact quotient $C^1/\ker\chi$ is a closed subgroup of the
circle, hence is either finite or a circle, and this is one permitted
elementary compact-invariance level. At this level the weighted
Schwartz seminorms, after $x=\log|g|$, are equivalent to
\[
 p_{N,j}(f)=\sup_x e^{N|x|}|f^{(j)}(x)|,
 \qquad N,j\ge0.
 \tag{M24.5}
\]
To see the equivalence directly, the weights $b=N$ and $b=-N$
bound (M24.5); conversely every polynomial factor in $x$ occurring
in a weighted Schwartz seminorm is absorbed by one larger
exponential weight. Derivatives on the compact factor multiply
this fixed character by constants. Inclusion of this level in the
inductive limit is continuous. The inverse of (M24.4) on the
sector is also continuous: compact-character averaging followed
by restriction to the norm splitting is continuous into
$\mathcal E$ on each elementary level, by these same seminorm
estimates. It is therefore continuous on the inductive limit.
This establishes the stated sector topology directly, without
using a topological claim about a different quotient.

Next identify the denominator as a set of actual functions.
Use Meyer's common ambient space and embeddings from MT-M22:
\[
 i_+(v)=(\Sigma v,J\Sigma\mathcal F_*v),
 \qquad i_-(u)=(u,u),
 \qquad H_+=\mathcal S(\mathbf A_K)/K^\times,
 \quad H_-=\mathcal S_{\mathrm E}(C).
\]
The summation maps agree term by term:
$\Sigma f(g)=\Phi_K(f)(g)$ by (M24.2).
Their Poisson identity is
\[
 J\Sigma\widehat f-\Sigma f
      =f(0)-|\cdot|^{-1}\widehat f(0).
 \tag{M24.6}
\]
For $f\in\mathcal S_0(\mathbf A_K)$ the two entries agree.
The first entry has all the right-half-plane Schwartz weights,
and the second all the left-half-plane weights, by the defining
Meyer embeddings. Their common value has all real weights:
between any two admitted weights the weighted derivative is
bounded by the two endpoint weighted derivatives, and polynomial
Schwartz factors are treated in the same way. Thus this common
value belongs to $H_-$, and
$\Phi_K(\mathcal S_0)\subset H_+\cap H_-$.

Conversely, if $(u,u)=i_+(v)$ lies in this intersection, choose
an adelic Schwartz representative $f$ of the coinvariant class
$v$. Equation (M24.6) then vanishes identically. The functions
$1$ and $|\cdot|^{-1}$ are linearly independent because the norm
map is onto $\mathbf R_{>0}$, so $f(0)=\widehat f(0)=0$.
Consequently
\[
       \Phi_K(\mathcal S_0(\mathbf A_K))=H_+\cap H_-.
 \tag{M24.7}
\]
This argument uses an actual representative of the coinvariant
class, not a claim that every completion class has an adelic
representative. $H_+$ here is Meyer's quotient of the adelic
Schwartz space.

Compact-character averaging preserves both embedded spaces
and their intersection. On functions, the volume-one projector is
\[
 P_\chi u(g)=\int_{C^1}u(cg)\chi(c)\,dc.
\]
Masullo (2.2.5.2), p.~95, is precisely this formula applied to
$\Phi_K f$, except that his chosen compact Haar measure can have
volume $c_K\ne1$. This multiplies the operator by the fixed
nonzero scalar $c_K$ and leaves its image unchanged. Hence
\[
       \Theta_\chi=(H_+\cap H_-)^\chi.
 \tag{M24.8}
\]
One can also verify surjectivity in (M24.8) directly: an element
of the intersection in the $\chi$-sector already has an
$\mathcal S_0$ representative by (M24.7), and applying
$\Phi_{K,\chi}$ to that representative returns $c_K$ times the
element. No choice of a compact subgroup lift in the idele group
is required.

MT-M22, whose image criterion and topology were proved using
Meyer's specified maps, now gives
\[
 \mathcal L_\beta^{-1}\Theta_\chi
 =\{f\in\mathcal E:(\mathcal Mf)^{(j)}(\rho)=0
         \text{ for all }0\le j<m_\rho\}
 =\mathcal I_\eta.
 \tag{M24.9}
\]
The middle space is closed, since each indicated evaluation is
continuous. Thus $\Theta_\chi$ is closed in the actual sector
topology just checked. Multiplication by $|\cdot|^a$ is an
automorphism of that topology, with inverse multiplication by
$|\cdot|^{-a}$; this follows either from (M24.5), using an
integer larger than $|a|$, or directly by translating the family
of all real weights. This proves the closure assertion and the
continuous two-sided quotient isomorphism (M24.3).
\end{proof}

The action is also explicit. For $b\in C$,
\[
 \lambda_b M_a\mathcal L_\beta f
  =\beta(b)|b|^{-a}M_a\mathcal L_\beta
       \bigl(f(\,\cdot-\log|b|)\bigr).
 \tag{M24.10}
\]
For an idele-class convolution test $k$, set
\[
 k_\beta(y)=\int_{C^1}k(c\tau(e^y))\beta(c\tau(e^y))\,dc.
\]
Changing variables in the Haar integral gives
\[
 C_k M_a\mathcal L_\beta f
   =M_a\mathcal L_\beta\bigl((e^{-a\cdot}k_\beta)*f\bigr).
 \tag{M24.11}
\]
In particular the spectral multiplier is evaluated at $\rho-a$,
with the compact-character factor $\beta$ retained. At $a=0$ this
is precisely the MT-M22 module comparison; at $a=1/2$ it is the
centered version, not an assertion that the same convolution
operator is unchanged by the twist.
An alternative Haar normalization rescales the scalar convolution
integral by the same fixed positive constant; it changes neither
the denominator nor the quotient topology.

\paragraph{Exact source scope and the remaining completion question.}
The original Masullo PDF pages used here are: the theta-sum
definition (2.2.3.1), pp.~81--82, and its underlying sum (2.2.2.1),
p.~77; the polar hyperplanes (2.2.4.23), p.~91;
character averaging (2.2.5.2), p.~95; the actual smooth quotient
(7.1.4.7), p.~327; its centered introductory version (1.8.3.1),
p.~45; and Definitions D.3.3.1/D.3.3.3, pp.~804--805.
The original p.~327 quotient formula was rendered and read;
the other defining formulas were checked against extraction
from the original PDF. The Meyer inputs are his summation map
and Lemma~5.3, pp.~39--40; the embeddings in Lemma~5.4,
pp.~40--41; and Poisson formula (31), p.~41, with the complete
provided arguments already read for MT-M22. No new proof of
the full Masullo draft is claimed.

Masullo Definition~10.1.1.1, p.~434, instead closes
$M_a\Theta_\chi$ in a weighted Hilbert norm, and
(10.2.1.1)--(10.2.1.2), p.~439, and (10.4.1.3), p.~456,
form the resulting Hilbert quotients. Those are different
topological completions of the identified smooth quotient.
Their kernel, retained jets and trace-class properties must be
proved in that norm; MT-M25 and MT-M26 address those maps.
The quotient identification itself is unconditional. Any
subsequent use of universal Hecke GRH is still the conditional
input $M$.

The geometric crosswalk is now an actual Tate--Poisson
summation quotient, rather than an unspecified space sharing
its zero multiplicities. It advances the inherited cote~1 /
SA23 geo-H-lidskii and geo-H-genus interfaces at the precise
test-space and closed-denominator level. It makes no claim
to construct an $\ell$-adic or motivic realization.

\subsection*{MT-M25. Exact line-completion kernel and the strict jet cutoff}
Fix the Hecke function, $\mathcal E$, $\Xi=\Xi_\eta$,
$\mathcal I=\mathcal I_\eta$ and $Q=\mathcal E/\mathcal I$
of MT-M18. For $s\ge0$ and $a\in\mathbf R$, put
\[
 H_s=L^2\bigl(\mathbf R,(1+x^2)^{s/2}\,dx\bigr),
 \qquad T_af(x)=e^{ax}f(x),
 \qquad J_{a,s}=\overline{T_a\mathcal I}^{\,H_s},
 \quad Z_{a,s}=H_s/J_{a,s}.
 \tag{M25.1}
\]
These are exactly the fixed-character versions of Masullo's
line norm (9.1.2.3), original PDF p.~399, his closure
(10.1.1.3), p.~434 with $I=\{0\}$, and his quotient
(10.2.1.2), p.~439, by MT-M24. The fixed unitary norm
modulation has absolute value one and is harmless for this norm.
Let $m_\rho$ be the multiplicity of a zero of $\Xi$, and set
\[
       q_s=\max\{0,\lceil(s-1)/2\rceil\}.
 \tag{M25.2}
\]
Then
\[
 J_{a,s}=
 \bigcap_{\substack{\rho=a+i\gamma:\Xi(\rho)=0\\
                    0\le k<\min(m_\rho,q_s)}}
 \ker\left[u\longmapsto\int_{\mathbf R}
                  u(x)x^k e^{i\gamma x}\,dx\right].
 \tag{M25.3}
\]
Every functional displayed here is continuous on $H_s$.
The natural map
\[
 q_{a,s}:Q\longrightarrow Z_{a,s},\qquad[f]\longmapsto[T_af]
 \tag{M25.4}
\]
is continuous with dense image, and its kernel consists exactly
of the classes whose jets $F^{(k)}(\rho)$, $F=\mathcal Mf$,
vanish at all the indices in (M25.3). Thus the line completion
retains only the zeros on $\Re\rho=a$ and only their derivatives
with the \emph{strict} inequality
\[
                         2k<s-1.
 \tag{M25.5}
\]
It is faithful on the full smooth quotient if and only if all
zeros lie on that line and $m_\rho\le q_s$ for every zero.
No GRH or simplicity is assumed in these assertions.

\begin{proof}
The continuous bilinear dual of $H_s$ can be written as
$H_{-s}$ under $\ell_g(u)=\int u(x)g(x)\,dx$.
This is just the Riesz representation theorem after multiplying
by $(1+x^2)^{s/4}$; complex conjugation can be absorbed into
the representing $g$.
For a polynomial $P$ of degree $k$ with nonzero leading
coefficient,
\[
 P(x)e^{i\gamma x}\in H_{-s}
 \quad\Longleftrightarrow\quad
 \int_{|x|\ge1}|x|^{2k-s}\,dx<\infty
 \quad\Longleftrightarrow\quad 2k<s-1.
 \tag{M25.6}
\]
This proves continuity of the displayed functionals and the
necessity of their strict numerical threshold. It also shows
that the number of admissible nonnegative integers $k$ is
exactly $q_s$, including $0\le s\le1$ when there are none.
Each such functional vanishes on $T_a\mathcal I$ by MT-M18's
full-jet ideal equality, hence on its Hilbert closure.

We prove there are no other annihilators. The elementary
inequality
\[
 1+(x+y)^2\le2(1+x^2)(1+y^2)
\]
implies, for either sign of $s$,
\[
 \|u(\,\cdot-y)\|_{H_{\pm s}}
       \le 2^{s/4}(1+y^2)^{s/4}\|u\|_{H_{\pm s}}.
 \tag{M25.7}
\]
For the negative sign use the same inequality with $x$ and
$x-y$ interchanged before taking reciprocals. Minkowski's
integral inequality therefore gives
\[
 \|k*u\|_{H_{\pm s}}
 \le2^{s/4}\|u\|_{H_{\pm s}}
       \int|k(y)|(1+y^2)^{s/4}\,dy.
 \tag{M25.8}
\]
In particular convolution by a Schwartz function is bounded
on both spaces. Smooth compactly supported functions are
dense in both weighted $L^2$ spaces: truncate first, and then
approximate on the remaining compact interval, where the
weight and its inverse are bounded. These facts also prove
strong continuity of translation, first on that dense subspace
and then by (M25.7).

Let $g\in H_{-s}$ annihilate $T_a\mathcal I$.
It defines a tempered distribution, since Cauchy--Schwarz
bounds $|\int ug|$ by the continuous Schwartz seminorm
$\|u\|_{H_s}\|g\|_{H_{-s}}$.
Use the Fourier convention
$\mathcal F_+u(t)=\int u(x)e^{itx}\,dx$ and its inverse
$(2\pi)^{-1}\int\phi(t)e^{-itx}\,dt$.
Define the tempered distribution $D$ by
$D(\phi)=\ell_g(\mathcal F_+^{-1}\phi)$.
Since $\psi_a=T_a\psi_\eta\in\mathcal E$ and
$\mathcal F_+\psi_a(t)=\Xi(a+it)$, we obtain
\[
                         \Xi(a+it)D=0.
 \tag{M25.9}
\]
Here is why the identity holds on every Schwartz test,
not just an entire Fourier transform. For $u\in\mathcal E$,
$\psi_a*u\in T_a\mathcal I$ because
$T_a(f_1*f_2)=T_af_1*T_af_2$ and $T_a\mathcal E=\mathcal E$.
Thus $\ell_g(\psi_a*u)=0$. The space $\mathcal E$ contains
$C_c^\infty$ and is dense in the usual Schwartz space;
convolution by $\psi_a$ is continuous there, and
$\ell_g$ is continuous there. This extends the equality to
all Schwartz $u$, which is precisely (M25.9).

The zeros of $t\mapsto\Xi(a+it)$ on the real line form a
locally finite set. Away from them, a compactly supported
smooth test can be divided by this nonvanishing multiplier,
so (M25.9) shows $\operatorname{supp}D$ is contained in that
set. Near a zero $\gamma$ of order $m_\rho$, multiplication
by a smooth invertible function reduces (M25.9) to
$(t-\gamma)^{m_\rho}D=0$. A distribution supported there
with this property depends only on the first $m_\rho$ jets
of a test: subtract its Taylor polynomial of degree
$m_\rho-1$ times a cutoff equal to one near $\gamma$;
the remainder is divisible by $(t-\gamma)^{m_\rho}$ there.
It is therefore a linear combination of
$\delta_\gamma^{(k)}$, $0\le k<m_\rho$.

Choose a smooth frequency cutoff supported near this single
zero and equal to one near it. Multiplying $D$ by that
cutoff corresponds to convolution of $g$ with a Schwartz
kernel, with an inessential reflection dictated by the
chosen Fourier convention. By (M25.8) the resulting
function still belongs to $H_{-s}$. Its inverse transform
is a polynomial times $e^{i\gamma x}$. Formula (M25.6)
therefore excludes every nonzero coefficient of degree
$k$ with $2k\ge s-1$. Locally only the functionals listed
in (M25.3) occur.

Finally these local pieces exhaust the annihilator in its
Hilbert norm. Choose $\vartheta\in C_c^\infty(\mathbf R)$
equal to one near zero and put
$\vartheta_R(t)=\vartheta(t/R)$, $R\ge1$.
The associated convolution kernels have the form
$Rk(Rx)$, possibly reflected, where $k$ is Schwartz and
$\int k=1$. Their weighted $L^1$ norms in (M25.8) are
uniformly bounded, because
\[
 \int R|k(Ry)|(1+y^2)^{s/4}\,dy
       =\int|k(v)|(1+v^2/R^2)^{s/4}\,dv
       \le\int|k(v)|(1+v^2)^{s/4}\,dv.
\]
Translation continuity and dominated convergence in the
same integral show $Rk(R\cdot)*g\to g$ in $H_{-s}$.
The distribution $\vartheta_RD$ is supported on finitely
many zeros, so each such convolution is a finite linear
combination of the permitted polynomial waves. Hence
the full annihilator of $T_a\mathcal I$ is the Hilbert
closure of the span of precisely the functionals in
(M25.3). Taking its annihilator proves (M25.3).

The inclusion $T_a:\mathcal E\to H_s$ is continuous,
since one exponential seminorm with weight greater than
$|a|$ absorbs the polynomial factor in (M25.1). Its range
is dense because $T_a\mathcal E=\mathcal E$ contains
$C_c^\infty$. It vanishes on $\mathcal I$ after passing
to the quotient, and so induces (M25.4), continuously
with dense image. The functional indexed by $(\rho,k)$
evaluates $T_af$ to
\[
       \int e^{ax}f(x)x^ke^{i\gamma x}\,dx
                    =(\mathcal Mf)^{(k)}(\rho).
\]
This gives its exact kernel. If all zero jets are retained,
MT-M18 makes that kernel zero. If any zero jet is omitted,
the finite-block interpolation of MT-M18 realizes a class
with just that jet nonzero and all other zero jets zero.
It is a nonzero class in the kernel of (M25.4). This proves
the stated necessary and sufficient condition for faithfulness.
\end{proof}

\paragraph{A literal boundary correction with an arithmetic witness.}
The original PDF (10.2.1.6), p.~440, displays
\[
 N_s(\rho)=\min\{N(\rho),[(s-1)/2+1]\}.
\]
At $s=3$ this allows the first derivative at any double zero.
But its proposed representing function is
$x e^{i\gamma x}$, whose squared $H_{-3}$ norm contains
\[
              \int_{|x|\ge1}\frac{x^2}{(1+x^2)^{3/2}}\,dx
                         =+\infty.
 \tag{M25.10}
\]
Thus the endpoint inclusion is false as written; the
correct retained multiplicity is
$\min\{m_\rho,q_s\}$. Both the source norm (9.1.2.3),
p.~399, and the displayed cutoff (10.2.1.6), p.~440, were
rendered and visually read from the original PDF.
This is not a hypothetical objection requiring an unknown
double zero of the Riemann zeta function. By NT-L17,
the actual degree-six field
$K=\mathbf Q(\sqrt[3]{2},\sqrt{-3})$ has nontrivial zeros
of $\zeta_K$ of order at least two. Choose such a zero
$\rho$, put $a=\Re\rho$, and use its trivial Hecke
character. Then (M25.10) is a concrete counterexample
to the stated membership at its permitted $a$ and $s=3$.
Under $M$, one may take $a=1/2$.

More generally, for every fixed finite $s$, choose an odd
prime $r$ with $r-1>q_s$. The Kummer field of NT-L17
has zeros of multiplicity at least $r-1$. Its line quotient
at $a=1/2$ cannot faithfully retain the full smooth
quotient: under $M$ it loses the higher jets, and without
$M$ either that loss occurs or an off-line zero is lost.
The strict cutoff is a property of this particular
Hilbert completion, not a failure of the unconditional
adapted-block synthesis in MT-M18.

This gives the precise analytic map at the inherited
cote~1 / SA23 geo-H-lidskii and geo-H-genus interface:
the Hilbert norm can discard part of the finite-length
spectral divisor. It does not prove or refute the
conditional endpoint $M$, and it invokes none of the
draft's unverified general basis or nuclearity arguments.

\subsection*{MT-ART10. An integral Galois-module class forces a central zero}

Let $N/\mathbb Q$ be a tamely ramified $Q_8$-extension and let
$\chi$ be the unique irreducible complex character of degree two.
Fr\"ohlich's classical theorem identifies the two signs
\[
 W(\chi)=+1\quad\Longleftrightarrow\quad
       \mathcal O_N\text{ is free over }\mathbb Z[Q_8].
 \tag{MT-ART10a}
\]
The exact primary statement is Fr\"ohlich, \emph{The Galois
module structure of algebraic integer rings in fields with
generalised quaternion group} (1974), Theorem1, printed p.82,
with the locally free class invariant defined on pp.81--82.
The original proof is his \emph{Artin-root numbers and normal
integral bases for quaternion fields}, Inventiones17(1972),
143--166. Its full proof was not retrieved here; (MT-ART10a)
is used as the classical theorem, not claimed as a fresh proof.
Martinet's \emph{Fonctions z\^eta au point $1/2$}, \S VI,
printed114--115, gives the equivalent normal-basis formulation.

Consequently, if $\mathcal O_N$ is not free then
\[
 L(\tfrac12,\chi)=0,
 \qquad \operatorname{ord}_{s=1/2}\zeta_N(s)\ge2.
 \tag{MT-ART10b}
\]
These assertions are unconditional and are compatible with
both adopted GRH inputs.

\paragraph{Complete short deduction.}
Choose a cyclic subgroup $C_4\subset Q_8$ and a faithful
degree-one character $\lambda$ of $C_4$. Its induction is
$\chi$: its values are $2$ at the identity, $-2$ at the central
involution, and zero elsewhere, which is the irreducible
degree-two character. Artin induction-invariance therefore
identifies $L(s,\chi)$ with a nontrivial finite-order Hecke
function over $N^{C_4}$, so it is entire. Since $\chi$ is
selfdual, its completed function satisfies
\[
                 \Lambda(s,\chi)=W(\chi)\Lambda(1-s,\chi).
\]
Nonfreeness and (MT-ART10a) give $W(\chi)=-1$. Evaluation at
$s=1/2$ forces $\Lambda(1/2,\chi)=0$. The archimedean factors
and conductor power are finite and nonzero there, so the
finite-part $L$-function also vanishes. If $K=N^{\{1,-1\}}$
is the biquadratic subfield, the regular-character identity
gives the all-prime factorization
\[
                         \zeta_N(s)=\zeta_K(s)L(s,\chi)^2.
\]
The first factor is holomorphic at $1/2$, so the order is at
least two. More precisely the order of $L(s,\chi)$ there is
positive and odd, by the same functional equation. This proves
(MT-ART10b) without a general Artin-holomorphy hypothesis.

\paragraph{A completely specified classical example.}
Martinet, \emph{Modules sur l'alg\`ebre du groupe quaternionien}
(1971), TheoremIV.1 printed406 and Example3 printed407, gives
\[
 K=\mathbb Q(\sqrt5,\sqrt{41}),\qquad
 M=\frac{5+\sqrt5}{2}\frac{41+\sqrt{205}}{2},\qquad
 N=K(\sqrt M).
 \tag{MT-ART10c}
\]
This is a totally real tame $Q_8$-extension, ramified only at
$5$ and $41$, with nonfree integer ring. In his trace criterion
the two values are $205$ and $(1+5+41+205)/4=63$; they differ
modulo four, so the nonfree case applies. The radicand was
checked directly in the original page407 image, where the
second square root is $\sqrt{5\cdot41}$.
Thus this particular field has $\zeta_N(1/2)=0$ and satisfies
(MT-ART10b). For a convenient polynomial description, multiplying
$X^2-M$ over the four conjugates of $M$ gives
\[
 X^8-205X^6+13940X^4-378225X^2+3404025.
 \tag{MT-ART10d}
\]
An exact rational coefficient calculation of this norm is
recorded separately; the Galois group, ramification and integral
module classification are the cited classical example.

The programme bridge here is precise: tameness makes
$\mathcal O_N$ a locally free integral Galois module, and its
class in the locally free class group of $\mathbb Z[Q_8]$
controls one functional-equation sign. Rationalizing erases
this integral class while retaining the usual regular
representation. This is an example of the integral $K_0$
information distinguished in AG-LEFSCH-LATTICE and
AG-LEFSCH-DEFECT, not a general formula for all arithmetic
$K$-groups. It also shows that a residual condition such as
the nonvanishing of a Dedekind zeta function at $1/2$ cannot
be discarded merely because $M_H$ holds. It is not an example
inside the particular $S_5$ family of MT-S502; that family's
central nonvanishing gate remains its own arithmetic question.


\subsection*{MT-M26. Trace class on the specified strip Hilbert quotient}
Fix one primitive unitary Hecke function and the notation of
MT-M18 and MT-M24. Let $s\ge0$, $a\in\mathbf R$, and let
$I=[u,v]$ be a compact interval with $u<v$. Define
\[
 \begin{split}
 w_{I,s}(x)&=\max\{e^{2ux},e^{2vx}\}(1+x^2)^{s/2},\\
 H_{I,s}&=L^2(\mathbf R,w_{I,s}(x)\,dx),\\
 J_{a,I,s}&=\overline{T_a\mathcal I_\eta}^{\,H_{I,s}},
 \qquad Z_{a,I,s}=H_{I,s}/J_{a,I,s},
 \qquad T_af(x)=e^{ax}f(x).
 \end{split}                                                    \tag{M26.1}
\]
Assume there is a closed strip $[c_0,c_1]$ containing the real
parts of all zeros of $\Xi_\eta$, with
\[
                 [c_0-a,c_1-a]\subset(u,v).
                                                               \tag{M26.2}
\]
Then the natural map $Q_\eta\to Z_{a,I,s}$ is continuous,
injective and has dense image. More precisely,
\[
 J_{a,I,s}
 =\{h\in H_{I,s}:(\mathcal Mh)^{(k)}(\rho-a)=0
       \text{ for all zeros }\rho\text{ and }0\le k<m_\rho\}.
                                                               \tag{M26.3}
\]
The Mellin transforms and derivatives in (M26.3) are defined
by absolutely convergent integrals inside the strip $I$.
For every $g\in\mathcal E$, convolution induces a trace-class
operator on this \emph{actual} Hilbert quotient, and
\[
       \operatorname{Tr}_{Z_{a,I,s}}C_g
          =\sum_{\Xi_\eta(\rho)=0}
                        m_\rho(\mathcal Mg)(\rho-a),
                                                               \tag{M26.4}
\]
with absolute convergence. The trace-class estimates are
uniform when $g$ ranges over a bounded subset of $\mathcal E$.

Unconditionally one may take $[c_0,c_1]=[0,1]$, so the result
holds whenever $[0,1]\subset a+\operatorname{int}I$.
Under GRH for this one function, one may take $c_0=c_1=1/2$.
In particular, under the conditional input $M$, the result
holds at $a=1/2$ for every $I=[-\alpha,\alpha]$ with
$\alpha>0$, for each fixed Hecke character separately.

\begin{proof}
We first specify the adapted blocks needed inside the available
strip. Choose fixed real numbers
\[
        a+u<b_-<c_0\le c_1<b_+<a+v.
\]
Use the selected horizontal heights $h_n$ from MT-M18 and
rectangles $\Omega_n=(b_-,b_+)\times(h_n,h_{n+1})$.
The proof of MT-M18 gives the same block theorem for these
vertical edges, as the following estimates make explicit.
The horizontal minimum-modulus estimate there was proved
on every fixed horizontal segment, so it applies on
$[b_-,b_+]$. On either vertical edge the distance from every
zero is at least
$\delta=\min(c_0-b_-,b_+-c_1)>0$.
In the Blaschke-product proof of MT-M18, each factor on these
edges is therefore at least
$\min(1,\delta)/(2R)$, where $R$ is the fixed disk radius
used there. There are $O(\log(2+|n|))$ such factors.
Their product has a polynomial lower bound. The same
zero-free factor estimate and the two-sided Stirling bound
then give
\[
 \inf_{w\in\partial\Omega_n}|\Xi_\eta(w)|
    \ge c(2+|n|)^{-B}e^{-\kappa|n|}.
                                                               \tag{M26.5}
\]
At bounded height the edges contain no zeros and compactness
supplies the bound. The fixed width and bounded height of the
rectangles give uniformly bounded perimeter. The entire
difference-kernel construction in MT-M18, its Gaussian estimate,
and its quotient-infimum argument consequently give finite
Banach algebras $J_n$ and lifts $\mathcal T_n:J_n\to\mathcal E$
with
\[
 \begin{gathered}
 d_n=\dim J_n=O(\log(2+|n|)),\qquad
 R_m\mathcal T_n=\begin{cases}\operatorname{id}_{J_n},&m=n,\\0,&m\ne n,
                         \end{cases}\\
 p_{N,j}(\mathcal T_n z)
       \le C_{N,j}(2+|n|)^{B_{N,j}}\|z\|_{J_n}.
 \end{gathered}                                                 \tag{M26.6}
\]
Here $J_n$ is the quotient of the holomorphic supremum-norm
algebra on the closed rectangle by its full block jets.
The absolute synthesis, closed-ideal equality and two-sided
block isomorphism of MT-M18 also hold: their proof uses
(M26.5), bounded perimeter and inclusion of every zero in
exactly one rectangle, all just verified. In particular the
lifts are not unstable individual Hermite interpolation
vectors. This establishes the precise version of MT-M18
being used below.

For $h\in H_{I,s}$ its Mellin integral is holomorphic on
$u<\Re z<v$. If $\sigma$ ranges over a compact subinterval
of $(u,v)$, then for every integer $k\ge0$,
\[
 \begin{split}
 |(\mathcal Mh)^{(k)}(\sigma+it)|
 &\le\|h\|_{H_{I,s}}
   \left(\int_{\mathbf R}
     \frac{|x|^{2k}e^{2\sigma x}}{w_{I,s}(x)}\,dx\right)^{1/2}\\
 &\le C_k\|h\|_{H_{I,s}},
 \end{split}                                                    \tag{M26.7}
\]
uniformly in $t$. The integral converges by the positive
distance to both endpoints, uniformly in that compact
subinterval. The same estimate justifies differentiation
under the integral on compact subsets of the strip.
Consequently
\[
 R_n^H(h)=\bigl[w\longmapsto(\mathcal Mh)(w-a)\bigr]\in J_n
 \quad\hbox{satisfies}\quad
       \|R_n^H(h)\|_{J_n}\le C\|h\|_{H_{I,s}},
                                                               \tag{M26.8}
\]
where $C$ is independent of $n$. This follows by bounding
the quotient norm by the supremum of this representative
on the contour. For $h=T_af$ the representative is
$\mathcal Mf(w)$, so $R_n^H$ vanishes on $T_a\mathcal I_\eta$,
and by continuity on its Hilbert closure. It induces a
bounded map $r_n:Z_{a,I,s}\to J_n$ with norm at most $C$.

Let $\pi:H_{I,s}\to Z_{a,I,s}$ be the quotient map, and put
$s_n=\pi T_a\mathcal T_n$. The map $T_a:\mathcal E\to H_{I,s}$
is continuous: a seminorm $p_{N,0}$ with
$N>|a|+\max(|u|,|v|)+1$ bounds its norm by an integrable
exponential times the fixed polynomial weight. Thus (M26.6)
implies, for some fixed constants $C,B$,
\[
 \|s_n\|_{J_n\to Z_{a,I,s}}\le C(2+|n|)^B,
 \qquad r_m s_n=\delta_{mn}\operatorname{id}_{J_n}.
                                                               \tag{M26.9}
\]
The last identity follows directly from the jets of the
lifts. In particular each block is actually realized in
the Hilbert quotient.

Convolution by $g\in\mathcal E$ is bounded on $H_{I,s}$.
Indeed, with $B_I=\max(|u|,|v|)$, the elementary polynomial
weight inequality used in MT-M25 gives
\[
 \|h(\,\cdot-y)\|_{H_{I,s}}
 \le2^{s/4}e^{B_I|y|}(1+y^2)^{s/4}\|h\|_{H_{I,s}}.
                                                               \tag{M26.10}
\]
Minkowski's inequality integrates this bound against
$|g(y)|$, and the resulting integral is finite.
The subspace $T_a\mathcal I_\eta$ is invariant under this
convolution because
\[
           g*T_af=T_a\bigl((e^{-a\cdot}g)*f\bigr)
                                                               \tag{M26.11}
\]
and $\mathcal I_\eta$ is a convolution ideal in $\mathcal E$.
Its Hilbert closure is therefore invariant as well. Hence
$C_g$ induces a bounded operator on $Z_{a,I,s}$.

Let $M_{g,n}$ be multiplication in $J_n$ by the block class
of $w\mapsto\mathcal Mg(w-a)$. Submultiplicativity of the
quotient algebra norm shows
\[
 \|M_{g,n}\|\le b_n(g):=
       \sup_{w\in\partial\Omega_n}|\mathcal Mg(w-a)|.
                                                               \tag{M26.12}
\]
For each $A>0$, $b_n(g)\le C_A(g)(2+|n|)^{-A}$,
with $C_A(g)$ bounded on bounded subsets of $\mathcal E$,
by its strip-Schwartz Mellin transform. These constants
are bounded by finitely many $\mathcal E$ seminorms:
integrating by parts sufficiently many times on the
fixed strip bounds the high ordinates, and direct
integration bounds the remaining compact set. The finite-rank
operators $s_nM_{g,n}r_n$ have rank at most $d_n$ and
satisfy the Hilbert trace-norm bound
\[
 \|s_nM_{g,n}r_n\|_1
       \le C d_n(2+|n|)^B b_n(g).
                                                               \tag{M26.13}
\]
Here we used the elementary inequality
$\|T\|_1\le\operatorname{rank}(T)\|T\|$ for a finite-rank
operator on a Hilbert space. Thus the series
\[
                     \sum_n s_nM_{g,n}r_n
                                                               \tag{M26.14}
\]
converges in trace norm, uniformly for $g$ in bounded test
sets. These same estimates show that
$g\mapsto C_g$ is continuous into the trace-class norm.

It equals the actual induced convolution. To verify this
without assuming spectral completeness in the Hilbert
space, first take a vector $\pi T_af$ with $f\in\mathcal E$.
Equation (M26.11) and the absolute block synthesis in the
\emph{already proved smooth quotient} give the identity
with (M26.14) on this vector. Applying $\pi T_a$ is
legitimate because it is continuous from that smooth
quotient into the Hilbert quotient; no injectivity has
been assumed. These vectors are dense: $T_a\mathcal E$
contains $C_c^\infty$, which is dense in the weighted
$L^2$ space by truncation and compact approximation.
Both operators are bounded, so the identity holds on
all of $Z_{a,I,s}$.

We can now prove the exact closure statement (M26.3).
One inclusion follows already from (M26.7). For the other,
if every zero jet of $h\in H_{I,s}$ vanishes, then
$r_n[ h ]=0$ for every $n$. The operator identity just
proved gives $C_g[h]=0$ for every $g\in\mathcal E$.
Take a smooth compactly supported approximate identity
$g_\epsilon$ of integral one. Translation is strongly
continuous on $H_{I,s}$: this holds first for compactly
supported smooth functions, then for all functions by
their density and (M26.10). That estimate is uniformly
bounded on small $y$, so
$g_\epsilon*h\to h$ in $H_{I,s}$.
Passing to the quotient gives $[h]=0$, proving (M26.3).
In particular $Q_\eta\to Z_{a,I,s}$ is injective, since
all zero jets of a smooth vector still determine its
class by MT-M18. Continuity and density were proved above.

Finally the finite-rank trace identity and (M26.9) give
\[
 \operatorname{Tr}(s_nM_{g,n}r_n)
       =\operatorname{Tr}_{J_n}(M_{g,n}r_ns_n)
       =\operatorname{Tr}_{J_n}M_{g,n}.
\]
In the local algebra at a zero $\rho$ of multiplicity
$m_\rho$, multiplication by a holomorphic function is
triangular in the powers of $w-\rho$, with that function's
value at $\rho$ repeated $m_\rho$ times on the diagonal.
This computes the trace as
$\sum_{\rho\in\Omega_n}m_\rho\mathcal Mg(\rho-a)$.
Summing is legitimate by (M26.13), and the scalar sum
is also absolutely convergent since $d_n=O(\log(2+|n|))$
and the Mellin transform is rapidly decreasing on the
fixed strip. This proves (M26.4).
\end{proof}

The injection from the smooth quotient has proper dense
image. Indeed $\Xi_\eta$ has infinitely many nontrivial
zeros, hence infinitely many nonzero $J_n$. The rapid
block space of MT-M18 is not normable: if one norm gave
its topology, it would be bounded above by a constant
times some $q_K$, while $q_{K+1}$ would be bounded by
a constant times that norm. A unit vector in a single
block at height $n$, scaled by $(2+|n|)^{-K}$, contradicts
these two inequalities as $|n|\to\infty$. Were the
continuous injection onto the Hilbert quotient surjective,
the Fr\'echet open mapping theorem would make the smooth
quotient normable. Thus this is a faithful Hilbert
completion, not an identification of its norm topology
with the stronger smooth topology.

At $a=1/2$ put $h^*(x)=\overline{h(-x)}$. Under GRH,
(M26.4) yields
\[
 \operatorname{Tr}C_{h*h^*}
      =\sum_\rho m_\rho|\mathcal Mh(\rho-1/2)|^2\ge0.
                                                               \tag{M26.15}
\]
On any strip satisfying the unconditional sufficient
condition $[0,1]\subset1/2+\operatorname{int}I$, positivity
of these traces for every $h\in\mathcal E$ is equivalent
to GRH for the fixed function. For if
$\rho^\dagger=1-\overline\rho\ne\rho$ is an off-line pair,
the unitary functional equation gives equal multiplicities.
MT-M18 finite-block interpolation, followed by $T_{1/2}$,
produces $h$ with Mellin values $1,-1$ at
$\rho-1/2,\rho^\dagger-1/2$, respectively, and zero at
all other zeros. Formula (M26.4) then gives
$\operatorname{Tr}C_{h*h^*}=-2m_\rho$.
No positivity of the operator $C_{h*h^*}$ itself is claimed;
the faithful full-jet obstruction of MT-M20 still applies.

\paragraph{Literal source interface and remaining scope.}
Masullo's norm (9.2.2.31), original PDF p.~416, becomes
(M26.1) after $x=\log|g|$ in the fixed compact-character
sector; the supremum over subintervals is attained at
the endpoints of the compact interval $I$. His closure
Definition~10.1.1.1, p.~434, becomes $J_{a,I,s}$ by
MT-M24. For $I=[-\alpha,\alpha]$ and $a>0$, his stated
parameter range, the resulting space
is his specified quotient (10.4.1.3), original PDF p.~456,
in that sector. The convolution convention is exactly
(M24.11); an idele-class test uses the scalar function
$k_\beta$, so its trace is evaluated at the correspondingly
shifted quasi-character. Compactly supported smooth tests
are included in $\mathcal E$.

This independently supplies fixed-character spectral
trace class on those literal quotients; it does not invoke
Theorem~10.3.2.1 or the draft's general Schauder claims.
All constants may depend on the fixed field and character.
A simultaneous summation over varying compact characters
would require additional uniform conductor and archimedean
parameter estimates; that is not inferred from a
fixed-character result. Neither the trace formula nor
the space identification proves $M$. The remaining
positivity input is still conditional.

In the inherited cote~1 / SA23 geo-H-lidskii and geo-H-genus
crosswalk, the map from the completed-Hecke spectral
divisor to a specified adelic Hilbert quotient is now
explicit and faithful whenever the containing strip has
interior room. The line quotient MT-M25 explains why
that interior condition cannot simply be discarded.

\subsection*{MT-M27. Full adelic trace class on a wide strip}
Let $K$ be a fixed number field, $C=C_K$, and
$C^1=\ker|\cdot|$. Choose a norm splitting and normalize
the compact Haar measure to have volume one. Put
\[
 \begin{split}
 E_K&=\mathcal S_{\mathrm E}(C),\qquad
 \Theta_K=\Phi_K(\mathcal S_0(\mathbf A_K)),\\
 H_{J,s}&=L^2\!\left(C,
       \max\{|g|^{2c},|g|^{2d}\}
       (1+|\log|g||^2)^{s/2}\,dg\right),\\
 \mathcal J_{J,s}&=\overline{\Theta_K}^{\,H_{J,s}},
 \qquad Z_{J,s}=H_{J,s}/\mathcal J_{J,s},
 \end{split}                                                    \tag{M27.1}
\]
where $s\ge0$ and $J=[c,d]$ satisfies $c<0<1<d$.
The topology on $E_K$ is the one specified in Masullo
Definitions D.3.3.1/D.3.3.3, original PDF pp.~804--805.
The Hilbert space and denominator in (M27.1) are his
specified norm and Hilbert closure, before the harmless
centering twist of MT-M24.

For each $\chi\in\widehat{C^1}$ choose a unitary extension
$\beta_\chi$ to $C$ and use the Hecke convention of MT-M22.
Let $\Xi_\chi$ have zeros $\rho$ with multiplicities
$m_{\chi,\rho}$. Define
\[
 \widehat f(\beta_\chi|\cdot|^z)
       =\int_C f(g)\beta_\chi(g)|g|^z\,dg.
                                                               \tag{M27.2}
\]
Then every $f\in\mathcal D(C)$ induces a trace-class
convolution operator on $Z_{J,s}$, and
\[
 \operatorname{Tr}_{Z_{J,s}}C_f
     =\sum_{\chi\in\widehat{C^1}}
        \sum_{\Xi_\chi(\rho)=0}
                   m_{\chi,\rho}
                     \widehat f(\beta_\chi|\cdot|^\rho).
                                                               \tag{M27.3}
\]
The series is absolutely convergent, uniformly on bounded
subsets of $\mathcal D(C)$. The trace norms are uniformly
bounded on those same sets. This includes the infinitely
many archimedean characters at any fixed finite level.
No GRH is required.

There are also exact kernel identities:
\[
 \begin{split}
 \overline{\Theta_K}^{\,E_K}
   &=\{u\in E_K:
      \partial_z^k\widehat u(\beta_\chi|\cdot|^z)|_{z=\rho}=0
            \text{ for all }\chi,\rho,0\le k<m_{\chi,\rho}\},\\
 \mathcal J_{J,s}
   &=\{u\in H_{J,s}:
      \partial_z^k\widehat u(\beta_\chi|\cdot|^z)|_{z=\rho}=0
            \text{ for all }\chi,\rho,0\le k<m_{\chi,\rho}\}.
 \end{split}                                                    \tag{M27.4}
\]
Thus the corrected full smooth quotient
$E_K/\overline{\Theta_K}^{E_K}$ maps continuously and
injectively, with proper dense image, into the specified
Hilbert quotient. We do not identify its stronger smooth
topology with the Hilbert norm topology.

\paragraph{Classical input and exact scope.}
The additional uniform input here is the published Meyer
Lemma~5.6, author PDF pp.~41--42, in
\emph{On a representation of the idele class group related
to primes and zeros of $L$-functions},
\url{https://arxiv.org/pdf/math/0311468} (v3).
For sufficiently large finite $S$, restricting its
stated $C_K$-valued map to $U_S$-invariants, as in
its proof, gives
\[
 (P_--P_+^S)C_f:
      L^2(C_S)_\alpha\oplus L^2(C_S)_{1-\alpha}
           \longrightarrow
          \mathcal S(C_S)_{>}\oplus\mathcal S(C_S)_{<}
                                                               \tag{M27.5}
\]
extends to a nuclear operator for $\alpha>1$,
uniformly for $f$ in bounded test sets.
Here $>$ means all weights $>1$, $<$ all weights $<0$,
$P_-$ is the diagonal projection, and $P_+^S$ has range
in the embedded adelic summation space $H_+^S$.
The supplied proof of this lemma was read in full,
including its Poisson and cutoff reduction; its cited
upstream local kernel lemmas are classical inputs, not
claimed newly fully audited. The rank-one meaning of
``nuclear'' is explicitly given in Meyer \S2.1, formula
(5), author pp.~6--7. We prove the transfer from (M27.5)
to the particular Hilbert quotient; it does not follow
merely from equality of spectral multiplicities.

\begin{proof}
\emph{Character projections and the two closed kernels.}
The compact abelian Fourier decomposition gives an
orthogonal Hilbert sum
\[
                     H_{J,s}=\mathop{\widehat\bigoplus}_\chi
                                      H_{J,s}(\chi).
                                                               \tag{M27.6}
\]
The weight depends only on the norm, so the compact
character projectors are orthogonal contractions.
The subspace $\Theta_K$ is invariant under these
projectors: by the global set identity (M24.7), it is
the intersection of the two embedded Meyer spaces,
both preserved by compact averaging. Its $\chi$ image
is exactly $\Theta_\chi$ of MT-M24.
Consequently $\mathcal J_{J,s}\cap H_{J,s}(\chi)$ is
the Hilbert closure of $\Theta_\chi$: apply the
contractive projector to a convergent approximating
sequence from $\Theta_K$. MT-M26 identifies this sector
closure with its full zero-jet kernel, because
$[0,1]\subset\operatorname{int}J$.
All these jet functionals are continuous on the full
Hilbert space: the compact projection has norm one,
and the Cauchy--Schwarz estimate (M26.7) applies in
that sector. Finite sums of character projections
converge strongly in (M27.6). It follows that vanishing
of all character/zero jets is equivalent to membership
in $\mathcal J_{J,s}$. This proves the second identity
in (M27.4).

For completeness the corresponding smooth closure
requires no uniform inverse estimates. An element
$u\in E_K$ lies at one elementary compact-invariance
level, with compact factor a finite abelian group
times a torus. At this level its seminorms are the
weighted suprema of all derivatives in the norm and
compact directions. Integration by parts on that
torus bounds every Fourier coefficient, in each
fixed norm-direction seminorm, by arbitrarily high
negative powers of its lattice frequency. For
example, apply $(1-\Delta)^N$ in the compact variable;
the character eigenvalue grows as $(1+|\nu|^2)^N$.
Taking $2N$ larger than the torus dimension makes
the Fourier tails summable, and increasing $N$
treats any specified compact derivative as well.
Thus finite character truncations converge to $u$
in every seminorm of that elementary level, hence
in $E_K$.
If all its zero jets vanish, each single-character
truncation belongs to the actual image $\Theta_\chi$
by MT-M24. Every finite truncation is therefore in
$\Theta_K$, and $u\in\overline{\Theta_K}^{E_K}$.
The reverse inclusion follows from continuity of
each zero-jet functional and its vanishing on
$\Theta_K$. This proves the first identity in
(M27.4). The natural map from the corrected smooth
quotient into the Hilbert quotient is consequently
injective and continuous. Its image is dense:
finite character sums times smooth compactly
supported functions of $\log|g|$ belong to $E_K$
and are dense in (M27.6). It is proper, since
surjectivity would, after a single compact-character
projection, contradict the proper-image assertion
of MT-M26 in that fixed sector.

\emph{A finite level for each test set, not a finite set of characters.}
Choose $\alpha>1$ so that
$[1-\alpha,\alpha]\subset(c,d)$.
A test $f\in\mathcal D(C)$ is invariant under
$U_S=\prod_{v\notin S}\mathcal O_v^\times$
for sufficiently large finite $S$ containing all
infinite places; $S$ can be enlarged to make the
$S$-ideal class group trivial. One common $S$ works
for a bounded set of tests. This follows directly
from the compact-invariance-level definition of
the test-function bornology, and is stated
explicitly in Meyer's proof of Theorem~5.8,
author p.~43. Enlarging $S$ preserves invariance.
For such $S$,
$C/U_S\simeq C_S=\mathbf A_S^\times/\mathcal O_{K,S}^\times$:
the trivial $S$-class group lets one adjust all
valuations outside $S$ by a principal idele; the
kernel is precisely the $S$-unit group.

Averaging over $U_S$ is an orthogonal contraction
on $H_{J,s}$, preserves $\mathcal J_{J,s}$, and
therefore gives an orthogonal projection on
$Z_{J,s}$. Since $f$ is $U_S$-invariant,
$C_f=C_fP_{U_S}=P_{U_S}C_f$.
It is enough to prove trace class on this invariant
Hilbert subspace, identified with functions on
$C_S$, and then extend it by zero on the orthogonal
complement. The invariant denominator is the
closure of $P_{U_S}\Theta_K$, by the same averaging
argument as above. No truncation of the infinitely
many compact Lie characters of $C_S^1$ is made.

\emph{An explicit bounded map from Meyer's ambient pair.}
Write
\[
 A_S=\mathcal S(C_S)_{>}\oplus\mathcal S(C_S)_{<},
 \qquad E_S=\mathcal S(C_S)_{\mathbf R},
 \qquad i_-(h)=(h,h).
\]
Choose a smooth norm cutoff $\phi$ equal to one
near zero and zero near infinity. Meyer's map
\[
        p_-(u_0,u_1)=(1-\phi)u_0+\phi u_1
                                                               \tag{M27.7}
\]
is bounded from $A_S$ to $E_S$.
This also follows immediately from weighted
seminorms: on the upper tail use a right weight
larger than the desired weight, on the lower
tail a smaller left weight, and treat derivatives
of the cutoff on a compact norm interval.
It satisfies $p_-i_-h=h$ and $P_-=i_-p_-$.

Put $\mathcal P=\operatorname{span}\{1,|\cdot|^{-1}\}$.
Poisson gives
\[
 X_S=H_+^S+i_-(E_S)
      =\{(u_0,u_1)\in A_S:u_0-u_1\in\mathcal P\}.
                                                               \tag{M27.8}
\]
Here is the explicit splitting needed for this identity
and for the subsequent bounded map. Choose two adelic
Schwartz functions, invariant outside $S$, whose pairs
$(f(0),\widehat f(0))$ are linearly independent.
They exist by taking fixed standard finite factors
and archimedean smooth functions: one vanishes at
zero but has nonzero integral, and another has
nonzero value at zero. Normalize their two pairs.
Their embedded images in $H_+^S$ give a linear map
$V:\mathcal P\to H_+^S$ such that
$dV=\operatorname{id}_{\mathcal P}$, where
$d(u_0,u_1)=u_0-u_1$. This uses exactly Poisson
(M24.6), including its two signs.
For any pair in $A_S$, evaluate its difference at
norms $1$ and $e$, and let $\ell(u_0,u_1)$ be the
unique element of $\mathcal P$ with those two
values. The matrix
$\left(\begin{smallmatrix}1&1\\1&e^{-1}\end{smallmatrix}\right)$
is invertible, so $\ell:A_S\to\mathcal P$ is a
bounded linear map and equals $d$ on $X_S$.
Subtracting $Vd$ makes the two entries equal;
their common function has all real Schwartz
weights. This proves (M27.8) directly.

Let $q:H_{J,s}^{U_S}\to Z_{J,s}^{U_S}$ be the
Hilbert quotient map. Inclusion $E_S\to H_{J,s}^{U_S}$
is bounded, by the exponential weights and compactness
of the norm-one factor. Therefore
\[
 B:A_S\longrightarrow Z_{J,s}^{U_S},\qquad
 B(y)=q\bigl(p_-(y-V\ell(y))\bigr)
                                                               \tag{M27.9}
\]
is bounded. It satisfies $B(i_-h)=q(h)$.
It vanishes on $H_+^S$: if $y\in H_+^S$, then
$y-V\ell(y)$ is still in $H_+^S$ and is diagonal;
its common function belongs to the theta image
by the Poisson intersection identity (M24.7),
so its Hilbert quotient class is zero.
Thus on $X_S$, (M27.9) is the actual canonical
quotient map, with an explicitly defined bounded
extension to the whole ambient pair. No assertion
about an unspecified quotient topology is used.

\emph{Transfer of the nuclear extension.}
There is a bounded diagonal inclusion
\[
 \delta:H_{J,s}^{U_S}\longrightarrow
       L^2(C_S)_\alpha\oplus L^2(C_S)_{1-\alpha},
       \qquad h\longmapsto(h,h).
                                                               \tag{M27.10}
\]
Indeed both endpoint exponential weights on the
right are bounded by the maximum of the weights
at $c,d$; the polynomial factor in (M27.1) is
at least one. The inclusion has norm at most
$\sqrt2$ for the chosen normalizations.
Compose (M27.5) with this inclusion and (M27.9):
\[
              A_f=B(P_--P_+^S)C_f\delta:
                     H_{J,s}^{U_S}\to Z_{J,s}^{U_S}.
                                                               \tag{M27.11}
\]
It is a nuclear map between Hilbert spaces,
uniformly for the bounded test sets under consideration.
To make the bornological-to-Hilbert step explicit,
Meyer's definition writes the middle map as
$\sum_j\lambda_j w_j\ell_j$, where
$\sum|\lambda_j|<\infty$, $(w_j)$ is bounded in
$A_S$ and $(\ell_j)$ is equibounded in the source
dual. The images $Bw_j$ are norm bounded.
The functionals $\ell_j\delta$ are uniformly
norm bounded. If a precompact rather than the
norm-bounded convention is used on the Hilbert
source, the same assertion follows from uniform
boundedness: equiboundedness on its compact
subsets implies pointwise boundedness of these
continuous functionals, hence a common operator
norm bound. A linear functional bounded on all
compact subsets of a Banach space is continuous,
since otherwise a sequence tending to zero
could have unbounded images. This also verifies
that no bornology convention inserts a hidden
topology assumption. The resulting rank-one
series has summable products of Hilbert norms.

On a smooth $h\in E_S$, the definitions give
$P_-C_fi_-h=i_-C_fh$, while
$P_+^SC_fi_-h\in H_+^S$ is killed by $B$.
Consequently $A_fh=q(C_fh)$ there.
Convolution by $f$ is bounded on $H_{J,s}$:
the radial weight has the translation bound
\[
 \|\lambda_b h\|_{H_{J,s}}
  \le2^{s/4}\max(|b|^{B_0},|b|^{-B_0})
       (1+|\log|b||^2)^{s/4}\|h\|_{H_{J,s}},
 \qquad B_0=\max(|c|,|d|),
\]
and $f$ is compactly supported and integrable.
It preserves the theta image and its closure,
so it induces a bounded operator $\overline C_f$
on the Hilbert quotient. Since $E_S$ is dense
in $H_{J,s}^{U_S}$, the preceding equality extends
to $A_f=\overline C_fq$ on the full Hilbert space.
The quotient has its isometric orthogonal
section $\sigma:Z_{J,s}^{U_S}\to H_{J,s}^{U_S}$.
Thus $\overline C_f=A_f\sigma$ is trace class.
The rank-one bounds just proved are uniform for
bounded test sets, and extension by zero to
$(Z_{J,s}^{U_S})^\perp$ preserves the trace norm.

\emph{The full trace and absolute character summation.}
All compact character projectors commute with
$C_f$, so (M27.6) descends to an orthogonal
sum of quotient sectors. Their operators are
exactly those of MT-M26, and their traces are
the inner sums in (M27.3). A trace-class operator
on an orthogonal sum has trace equal to the
sum of its sector traces, by computing in
orthonormal bases of the summands.
For clarity, absolute convergence of the full
\emph{zero} sum also follows directly from its
Hilbert trace norm. Take any finite collection
of the finite block subspaces of MT-M26 in
finitely many character sectors. Their direct
sum is finite dimensional and invariant under
$C_f$; its eigenvalues are exactly the indicated
Mellin values with their block multiplicities.
An orthonormal Schur basis for this invariant
subspace gives
\[
       \sum_{\text{these blocks}}m_{\chi,\rho}
            |\widehat f(\beta_\chi|\cdot|^\rho)|
                     \le\|C_f\|_1.
                                                               \tag{M27.12}
\]
Indeed choose diagonal phases so that the sum
of the absolute diagonal entries is a trace
against a finite-rank contraction; its absolute
value is at most the trace norm. Increasing the
finite set of blocks proves absolute convergence
of the whole spectral sum and the same bound.
Here are also the tails needed for uniform absolute
convergence, rather than only a uniform bound on the
total sum. A bounded test set lies in one space of
smooth functions with fixed compact support on an
elementary quotient. It is relatively compact there
by the derivative bounds and Arzel\`a--Ascoli.
The linear map from this Fr\'echet test space to
trace-class operators is continuous: its boundedness
on bounded sets, just proved, implies continuity
into a Banach space. Explicitly, failure of continuity
would give a sequence tending to zero in successively
more of the defining seminorms while its image norms
tend to infinity, contradicting bounded-set preservation.
The image of the test set is consequently relatively
compact in trace norm. If $P_F$ is a finite sum of
compact-character projections, then $P_F\to1$
strongly; hence
$\|T-P_FTP_F\|_1\to0$ uniformly for $T$ in this
relatively compact trace-class set. Pointwise this
follows by approximation by finite-rank operators;
uniformity follows from a finite trace-norm net and
the contraction property of the projections.
Applying (M27.12) on the complementary sectors bounds
all their absolute spectral terms by this tail norm.
For the remaining finitely many characters, the
fixed-sector estimates of MT-M26 give uniform
decay of the absolute sums outside finitely many
height blocks. The character projection of a
bounded test set is bounded in that sector's
$\mathcal E$, so those estimates apply. Combining
the two tails proves the asserted uniform absolute
convergence on bounded test sets.
This proves (M27.3).
\end{proof}

For the centered or otherwise twisted literal Masullo
quotient, put $J=a+I$. Multiplication by $|\cdot|^a$
is an isometry from (M27.1) to the Hilbert space with
base $I$ and denominator
$\overline{|\cdot|^a\Theta_K}$, and
$C_kM_a=M_aC_{|\cdot|^{-a}k}$.
Consequently, whenever $[0,1]\subset a+\operatorname{int}I$,
its full trace is
\[
       \sum_{\chi,\rho}m_{\chi,\rho}
                    \widehat k(\beta_\chi|\cdot|^{\rho-a}).
                                                               \tag{M27.13}
\]
For $a=1/2$ and $I=[-\alpha,\alpha]$, the unconditional
sufficient condition is $\alpha>1/2$.

At this centered wide-strip quotient, let
$h^*(g)=\overline{h(g^{-1})}$. Under $M$ every spectral
quasi-character in (M27.13) is unitary, and therefore
\[
       \operatorname{Tr}C_{h*h^*}
          =\sum_{\chi,\rho}m_{\chi,\rho}
               |\widehat h(\beta_\chi|\cdot|^{\rho-1/2})|^2\ge0.
                                                               \tag{M27.14}
\]
This uses $M$ as a condition; trace class alone does
not supply this positivity. Conversely, positivity
for every $h\in\mathcal D(C)$ would imply GRH for
every Hecke character of this fixed field: a single
off-line pair gives a negative fixed-character
Schwartz test by MT-M26; multiplying it by smooth
expanding norm cutoffs approximates it in that
sector's $\mathcal E$ topology by elements of
$\mathcal D(C)$, preserving a negative trace for
large cutoff. Continuity of the fixed-sector
trace-class map justifies this last step. No
positive-operator realization is inferred.

\paragraph{What remains after this transfer.}
This is the full adelic wide-strip result, not merely
a list of fixed-character trace formulas. The
varying-character summability is supplied by
Meyer's classical uniform nuclear extension, with
an explicit map into Masullo's specified Hilbert
quotient. The all-character test topology remains
the ordinary inductive union of finite invariance
levels; constants need not be uniform when that
level is allowed to vary without a boundedness
condition. The narrow centered intervals
$0<\alpha\le1/2$, even under $M$, do not embed
into the two Hilbert spaces in (M27.5).
MT-M26 treats each of their fixed-character
sectors; simultaneous trace-class control on
such a narrow interval remains a separate gate.

The original Masullo norm and closure locators
remain (9.2.2.31), p.~416; Definition~10.1.1.1,
p.~434; and (10.4.1.3), p.~456. The geometric
crosswalk to inherited cote~1 / SA23
geo-H-lidskii and geo-H-genus now includes the
actual full adelic Hilbert quotient on a wide
strip. Meyer's classical spectral realization
and uniform nuclearity retain their original
credit; this proof supplies the precise
comparison and completion transfer used here.

\subsection*{MT-ART11. The tame integral class is a central-parity invariant}
Let $N/K$ be a finite Galois extension of number fields, tamely
ramified at every finite place, with group $G$. All module classes
in this statement are over $\mathbb Z[G]$, after restriction of
scalars from $\mathcal O_K$. Put
\[
 U_{N/K}=[\mathcal O_N]-[\mathcal O_K[G]]
       \in\operatorname{Cl}(\mathbb Z[G]),\qquad
 m_\chi=\operatorname{ord}_{s=1/2}L(s,\chi,N/K)\in\mathbb Z.
                                                        \tag{ART11.1}
\]
Orders here are meromorphic orders, so a negative value denotes
a pole. Let $R_G^s$ be the virtual symplectic-character group and
$T(\chi)=\chi+\overline\chi$. The classical root-number map
\[
 t:\operatorname{Hom}_{\Omega_{\mathbb Q}}
       (R_G^s/T(R_G),\{\pm1\})
          \longrightarrow\operatorname{Cl}(\mathbb Z[G])
                                                        \tag{ART11.2}
\]
is additive after writing its sign-group law multiplicatively.
Then
\[
 U_{N/K}=t\bigl(\chi\longmapsto(-1)^{m_\chi}\bigr),
 \qquad 2U_{N/K}=0.                                    \tag{ART11.3}
\]
In particular,
\[
 \zeta_N(1/2)\ne0\ \Longrightarrow\ U_{N/K}=0,
 \qquad
 U_{N/K}\ne0\ \Longrightarrow\
        \operatorname{ord}_{s=1/2}\zeta_N(s)\ge2.          \tag{ART11.4}
\]
Neither implication uses Artin holomorphy or GRH. The first
means stable freeness as a $\mathbb Z[G]$-module of rank
$[K:\mathbb Q]$; it does not by itself assert a relative normal
integral basis over $\mathcal O_K[G]$.

\paragraph{Classical input.}
Taylor, \emph{On Fr\"ohlich's conjecture for rings of integers
of tame extensions}, Invent.\ Math.\ 63 (1981), 41--79,
Theorem~1, p.~41, gives $U=t(W)$ and $2U=0$.
Cassou-Nogu\`es--Taylor, Ann.\ Inst.\ Fourier 33(2) (1983),
1--17, (2.2)--(2.5), pp.~6--7, and Corollary~A, p.~10,
fix (ART11.2) and the relative normalization in (ART11.1).
The original Taylor theorem page was visually inspected;
the 1983 paper's full main proof, p.~11, was read with its
explicit upstream Gauss-sum theorems retained as citations.

\begin{proof}
The completed Artin function has a meromorphic functional
equation. If $\chi$ is symplectic it is self-dual, so
\[
 \Lambda(1/2+z,\chi)=W(\chi)\Lambda(1/2-z,\chi).
\]
Its archimedean gamma and conductor factors are holomorphic
and nonzero at $1/2$. Write its nonzero Laurent expansion as
$a z^{m_\chi}+O(z^{m_\chi+1})$, where $a\ne0$ and
$m_\chi$ may be negative. Comparing the two leading
coefficients gives $1=W(\chi)(-1)^{m_\chi}$.
The same calculation extends multiplicatively to virtual
symplectic characters. Thus the parity character in
(ART11.3) is exactly the classical root-number character;
in particular it has the Galois invariance and the
$T(R_G)$-quotient property needed in (ART11.2).
Taylor's theorem gives (ART11.3).

Here are full details of the nonvanishing step, so that no
unproved holomorphy assertion enters (ART11.4). Set
\[
             \Theta=\sum_{\chi\in\operatorname{Irr}(G)}
                         m_\chi\chi.
\]
Suppose $\zeta_N(1/2)\ne0$. For a cyclic subgroup
$C\le G$, put $F=N^C$. The abelian factorization is
\[
       \zeta_N(s)=\prod_{\eta\in\widehat C}
                       L(s,\eta,N/F).
\]
Every factor is holomorphic at $1/2$, including the trivial
factor $\zeta_F$. Their nonnegative orders sum to zero,
so each order is zero. Artin induction and Frobenius
reciprocity identify these orders with
$\langle\operatorname{Res}^G_C\Theta,\eta\rangle_C$.
Hence $\operatorname{Res}^G_C\Theta=0$ for every cyclic
$C$. Every $g\in G$ lies in $\langle g\rangle$, so
$\Theta(g)=0$ for every $g$, and irreducible-character
independence gives $m_\chi=0$ for every $\chi$.
The root-number character is therefore trivial and
$U_{N/K}=0$. This also proves that nonvanishing of
$\zeta_N$ excludes central Artin poles in this argument.

For the contrapositive, $\zeta_N$ is holomorphic at $1/2$.
Its completed functional equation has sign $+1$ and
nonzero archimedean factors there. A central zero
therefore has positive even order, proving the lower
bound two in (ART11.4).
\end{proof}

If $G$ has no irreducible symplectic characters, its
symplectic quotient $R_G^s/T(R_G)$ is zero and $U_{N/K}=0$.
For odd $|G|$, Taylor's Corollary, p.~41, with complete
proof p.~42, strengthens this to actual freeness over
$\mathbb Z[G]$, using the Eichler condition and Jacobinski
cancellation. This cancellation step is part of that
classical corollary. There is no general freeness claim
from a vanishing stable class in the preceding proof.

For use beside the positivity gate MT-ART02, there is the
following precise refinement. If all the central Artin
orders are nonnegative, then
\[
 \operatorname{ord}_{1/2}\zeta_N
   =\sum_{\chi\in\operatorname{Irr}(G)}\chi(1)m_\chi
   \ \ge\!
   \sum_{\substack{\chi\ {
m irreducible\ symplectic}\\W(\chi)=-1}}
                         \chi(1).                       \tag{ART11.5}
\]
Indeed each selected order is a nonnegative odd integer
and therefore at least one; every other term is
nonnegative. Thus an integral sign obstruction gives a
degree-weighted lower bound once the central Heilbronn
matrix is positive. Without that extra positivity, the
unconditional conclusion remains (ART11.4).

The placement is the locally free integral $K_0$ of the
finite Galois cover, in contrast with its regular rational
representation after the normal-basis theorem. It extends
the concrete integral distinction of MT-ART10 and
AG-LEFSCH-LATTICE/AG-LEFSCH-DEFECT. In the inherited
cote~16, scans~14--18, no.~7 integral-ring discussion,
it supplies an exact sign-valued invariant that survives
in the lattice but is absent from the rational module.
The associated Artin representation lies in the
SGA~6, V, \S2 and \S7, especially (7.1.1),
representation-$\lambda$-ring crosswalk fixed in AG-LAMBDA.
These are programme placements, not an assertion that
SGA~6 or the archive contains Taylor's arithmetic theorem.

\subsection*{MT-ART12. The hermitian trace class recovers every central parity}
Keep the tame extension and the right-module convention.
Equip $\mathcal O_N$ with its $G$-valued trace form
and the comparison module with its standard trace form:
\[
 \begin{split}
 h_N(m,n)&=\sum_{\gamma\in G}
                \operatorname{Tr}_{N/\mathbb Q}(m n^\gamma)
                                \gamma^{-1},\\
 h_0(m,n)&=\operatorname{Tr}_{K/\mathbb Q}(\overline m n),
       \qquad m,n\in K[G],\\
 X_{N/K}&=[\mathcal O_N,h_N]-[\mathcal O_K[G],h_0]
                   \in K_0H(\mathbb Z[G]).
 \end{split}                                             \tag{ART12.1}
\]
The bar in $h_0$ is the group-ring involution
$\gamma\mapsto\gamma^{-1}$, and the field trace is
applied coefficientwise. Here $K_0H$ is the Grothendieck
group for these locally free hermitian modules and
orthogonal sum. Its discriminant has the canonical
sign extraction $\theta$ of Cassou-Nogu\`es--Taylor,
(2.9), p.~8. Their main theorem, p.~10, and
Corollary~B, p.~11, give
\[
 \theta_\chi(X_{N/K})=W(\chi),\qquad
 \eta_\chi(X_{N/K})=
       \bigl(\operatorname{N}\mathfrak f(\chi)\bigr)^{1/2}
                                      W_\infty(\chi)
                                                        \tag{ART12.2}
\]
for every symplectic $\chi$. These exact normalizations
are (3.6)--(3.9), p.~9; the original page was inspected,
including the square root of the conductor and the bar
in $h_0$.

The consequence for the present spectral programme is
\[
       (-1)^{\operatorname{ord}_{1/2}L(s,\chi)}
                         =\theta_\chi(X_{N/K}).           \tag{ART12.3}
\]
Thus the hermitian trace class determines each
symplectic central parity individually, including
meromorphic orders, without GRH or Artin holomorphy.
The underlying locally free class is obtained by the
forgetful map and gives $t(W)$ as in MT-ART11; inversion
of $t$ is not needed in (ART12.3).

\begin{proof}
Apply the complete Laurent-series parity calculation
in MT-ART11 to each symplectic character and substitute
the classical first equality in (ART12.2). This proves
(ART12.3) with no positivity premise. If two tame
extensions over the same base have identified groups
and the same class $X$ under that identification, the
homomorphisms $\theta_\chi$ take the same values, so
all their symplectic central orders have the same
parity. Applying the forgetful map gives equal
underlying stable classes, by (ART11.3). Conversely,
this proof makes no injectivity assumption on $t$;
stable-class data alone have not been substituted
for the hermitian class.
\end{proof}

The definition of the reduced hermitian group in
the cited source uses extension to
$\overline{\mathbb Q}[G]$; the source's p.~5 and p.~10
were visually checked to retain that algebraic closure.
No isometry over $\mathbb Q[G]$ is silently inferred.
The construction supplies a precise hermitian refinement
of the integral $K_0$ comparison in MT-ART11 and of
the trace programme at cote~1. It gives root signs
and central parities, while nonnegative multiplicities
remain the distinct Heilbronn-positivity gate. The
extra conditional lower bound (ART11.5) states exactly
what that positivity would add.

\subsection*{MT-M28. Full adelic trace class on every positive-width critical strip}
Let \(K\) be a fixed number field and retain the actual summation
image and test-function conventions of MT-M24 and MT-M27. Choose
an archimedean norm splitting \(C_K=C_K^1\times\mathbf R_+^\times\)
and give \(C_K^1\) Haar volume one. For \(a\in\mathbf R\), \(s\ge0\)
and \(I=[u,v]\), \(u<v\), put
\[
 \begin{split}
 H_{I,s}&=L^2\!\left(C_K,
    \max\{|g|^{2u},|g|^{2v}\}
         (1+\log^2|g|)^{s/2}\,dg\right),\\
 Z_{a,I,s}&=H_{I,s}/
       \overline{|\cdot|^a\Theta_K}^{\,H_{I,s}},
 \qquad \Theta_K=\Phi_K(\mathcal S_0(\mathbf A_K)).
 \end{split}                                                    \tag{M28.1}
\]
Assume GRH for all unitary Hecke characters of this \(K\), as
supplied by the conditional input \(M\), and assume
\[
                         \tfrac12-a\in(u,v).                  \tag{M28.2}
\]
Every \(f\in\mathcal D(C_K)\) then induces a trace-class
convolution operator on the Hilbert quotient (M28.1). With the
character convention of MT-M22 and with each unitary extension
\(\beta_\chi\) chosen trivial on the norm splitting, its trace is
\[
 \operatorname{Tr}_{Z_{a,I,s}}C_f
   =\sum_{\chi\in\widehat{C_K^1}}\ \sum_{\Xi_\chi(\rho)=0}
          m_{\chi,\rho}\,
          \widehat f(\beta_\chi|\cdot|^{\,\rho-a}).
                                                               \tag{M28.3}
\]
The sum is absolutely convergent, uniformly on bounded test
sets, and the trace norms are bounded on those sets. At every
fixed compact-invariance and compact-support test level the
map \(f\mapsto C_f\) is continuous into trace norm.

The denominator in (M28.1) is exactly the kernel of all
character/zero jets
\[
 \partial_z^k\widehat h(\beta_\chi|\cdot|^z)
                  \big|_{z=\rho-a},\qquad 0\le k<m_{\chi,\rho}.
                                                               \tag{M28.4}
\]
Consequently the corrected full smooth quotient
\[
 |\cdot|^a\mathcal S_{\mathrm E}(C_K)\big/
           |\cdot|^a\overline{\Theta_K}^{\,\mathcal S_{\mathrm E}(C_K)}
                                                               \tag{M28.5}
\]
maps continuously and injectively into \(Z_{a,I,s}\), with
proper dense image. In particular, at \(a=1/2\), every
\(I=[-\alpha,\alpha]\) with \(\alpha>0\) is allowed.
The zero-width completion remains governed by MT-M25.

\paragraph{Classical input and scope of the uniform proof.}
The analytic input besides \(M\) is the usual Hecke
continuation and functional equation, with its finite-order
property, already used in NT-L03, NT-Z10 and NT-Z17.
The continuation/functional-equation locator is
Ramakrishnan--Valenza, \emph{Fourier Analysis on Number Fields},
Theorem~7-19, printed pp.~279--280; the real and complex
local factors are in \S7.1, printed pp.~244--245, with their
explicit integral calculations on pp.~250--252.
These local formulas and the entire supplied proof of
Theorem~7-19 were read, and the relevant local-factor pages
were inspected as original page images. We do not claim that
Theorem~7-19 separately states a uniform finite-order estimate:
only the classical finite-order property of each function is
used below, and all uniform estimates are proved here.
The source is
\url{https://webusers.imj-prg.fr/~julien.marche/Ref_GL_GT/Fourier_Analysis_on_Number_Fields.pdf}.
The exact summation-image identification is MT-M24, and
the already proved one-character Hilbert theorem is MT-M26.
Meyer's uniform nuclear lemma from MT-M27 is not needed in
the uniform estimates below.

\begin{proof}
\emph{1. Character parameters at one elementary test level.}
Fix a compact subgroup \(k\subset C_K^1\) such that
\(T=C_K^1/k\) is a compact abelian Lie group. Its dual is
\[
                  \Lambda=\widehat T\simeq\mathbf Z^d\times F,
                                                               \tag{M28.6}
\]
where \(F\) is finite. Every bounded subset of
\(\mathcal D(C_K)\) is contained in one such invariance
level, with a common compact support in \(x=\log|g|\)
and bounded derivatives of every order on \(T\times\mathbf R\).
This is the elementary-level definition of the Bruhat test
bornology/topology used in MT-M27. We prove uniform estimates
at this level, without replacing its infinitely many
archimedean characters by a finite set.

Choose generators of the free part in (M28.6), and extend
each character to \(C_K\) by making it trivial on the chosen
norm splitting. Then the extension at lattice index
\(\nu\in\mathbf Z^d\times F\) is the product of the fixed
generator extensions and one of finitely many torsion
extensions. Each fixed extension has finite conductor.
The conductor of a product divides the least common
multiple of the conductors of its factors. Hence all finite
conductors at this level divide one fixed ideal.

For the primitive Hecke function \(L_\nu\) in the convention
of MT-M22, write its completion as
\[
 \Xi_\nu(z)=Q_\nu^{z/2}
      \prod_{j=1}^{n_K}\Gamma_{\mathbf R}(z+\mu_{\nu,j})L_\nu(z),
 \qquad \Gamma_{\mathbf R}(z)=\pi^{-z/2}\Gamma(z/2),
                                                               \tag{M28.7}
\]
up to a nonzero scalar. Here \(n_K=[K:\mathbf Q]\).
For every \(\nu\ne0\) this completion is entire. The
single trivial compact character, with the chosen extension,
gives \(\zeta_K\); for it use the pole-removed completion
and the fixed-character result MT-M26 throughout. A pure
norm character trivial on the norm splitting is trivial,
so there are no other polar exceptions in this family.

The local formulas in the cited pages show that real
places contribute shifts \(\epsilon+i\tau\), with
\(\epsilon=0\) or \(1\), and complex places contribute
\[
      |m|/2+i\tau,\qquad |m|/2+1+i\tau.
 \]
Indeed
\(\Gamma_{\mathbf C}(z)=
\Gamma_{\mathbf R}(z)\Gamma_{\mathbf R}(z+1)\);
an alternative constant normalization cancels in all
ratios below. Angular integers and real norm parameters
of products grow at most linearly in the exponent vector.
Thus, writing \(H_\nu=2+|\nu|\) for any fixed lattice norm,
\[
 \begin{gathered}
  1\le Q_\nu\le Q_0,\qquad
  \operatorname{Re}\mu_{\nu,j}\in\tfrac12\mathbf Z_{\ge0},
  \qquad |\mu_{\nu,j}|\le C H_\nu .
 \end{gathered}                                                 \tag{M28.8}
\]
The constants depend only on the field and this test level.
They are independent of \(\nu\).

\emph{2. Uniform gamma ratios and finite-\(L\) bounds.}
We record the elementary gamma estimates needed to avoid
an uncanceled, possibly very large, completed-function
normalization. For \(\operatorname{Re}z\ge c>0\), the
digamma series gives
\[
 \psi(z)=-\gamma+
      \sum_{m\ge0}\left(\frac1{m+1}-\frac1{m+z}\right),
 \qquad
 |\psi(z)|\le C_c+2\log(2+|z|),
 \qquad |\operatorname{Im}\psi(z)|\le C_c.                    \tag{M28.9}
\]
Here is a direct bound. Split the series at
\(N=\lceil2|z|+2\rceil\). The finite part is bounded by
two harmonic sums, using \(|m+z|\ge m+c\); the tail is
bounded by
\(\sum_{m\ge N}|z-1|/((m+1)|m+z|)=O(1)\).
For \(z=x+iy\), the imaginary part is bounded by
\[
 |y|\sum_{m\ge0}\frac1{(m+x)^2+y^2}
 \le \frac1{2c}+\int_0^\infty
           \frac{|y|}{(t+x)^2+y^2}\,dt
 \le \frac1{2c}+\frac\pi2,
 \]
with the \(y=0\) case immediate.
Consequently a bounded change in real part, within a
right half-plane, changes \(|\Gamma_{\mathbf R}|\) by
at most a fixed power of \(2+|z|\), in either direction.
A change of imaginary part by \(h\) changes it by at
most \(e^{C_c|h|}\). These statements follow by
integrating the real or imaginary derivative of
\(\log|\Gamma_{\mathbf R}|\); the half-plane is free
of gamma poles and zeros.

Fix \(A>0\). The Euler product gives
\[
 \zeta_K(2)^{-1}\le |L_\nu(2+it)|\le\zeta_K(2),
 \qquad |L_\nu(A+2+it)|\le\zeta_K(2).
                                                               \tag{M28.10}
\]
The lower bound follows factor by factor from
\((1+x)^{-1}\ge1-x\), \(0\le x<1\), including the
omitted ramified factors. On \(\sigma=-A\), the
functional equation expresses \(L_\nu(\sigma+it)\)
as \(L_{\bar\nu}(1-\sigma-it)\), a bounded conductor
power and a product of ratios
\[
  \frac{\Gamma_{\mathbf R}(1-\sigma+
                   \operatorname{Re}\mu_{\nu,j}
                   +i(t+\operatorname{Im}\mu_{\nu,j}))}
       {\Gamma_{\mathbf R}(\sigma+\mu_{\nu,j}+it)}
                                                               \tag{M28.11}
\]
in absolute value. Conjugation was used in the
numerator; the root number has modulus one.
Use
\(\Gamma_{\mathbf R}(z+2)=
 (z/(2\pi))\Gamma_{\mathbf R}(z)\)
to move the denominator to a fixed right half-plane.
The reciprocal of the original denominator becomes
the reciprocal of that shifted gamma function times
a polynomial of fixed degree. This remains valid
at poles of the original denominator, by analytic
continuation. Now (M28.9) and (M28.8) bound every
ratio by a fixed power of \(H_\nu+|t|\).
It follows that both boundary lines of the strip
\([-A,A+2]\) satisfy such a polynomial bound.

For explicitness, the Phragm\'en--Lindel\"of step is
uniform despite using only finite order separately
for each \(\nu\). Divide \(L_\nu(z)\) by
\((H_\nu+C_A+z)^B\), with integer \(B\) sufficiently
large and \(C_A>A+2\). This denominator has no zero
in the strip, and on its boundaries its modulus
is comparable with \((H_\nu+|t|)^B\). The quotient
is uniformly bounded on both boundaries and has
finite order for each \(\nu\). Multiply it by
\[
          \exp\{-\epsilon\cos(\lambda(z-z_0))\},
 \]
where \(z_0\) is the real midpoint of the strip and
\(\lambda>0\) is small enough that
\(\cos(\lambda(\sigma-z_0))>c>0\) throughout it.
The multiplier has modulus at most one on the
vertical boundaries and tends to zero faster than
every finite-order growth on horizontal boundaries
as their heights tend to infinity. The maximum
principle on rectangles, followed first by their
height tending to infinity and then by
\(\epsilon\downarrow0\), gives the same uniform
bound throughout the strip. Thus, for every
fixed \(A\),
\[
       |L_\nu(\sigma+it)|
           \le C_A(H_\nu+|t|)^{B_A}
                 \quad (|\sigma|\le A,\ \nu\ne0).
                                                               \tag{M28.12}
\]
No GRH has been used in (M28.8)--(M28.12).

Jensen's formula, centered at \(2+iT\), now gives
for each fixed radius \(R\)
\[
  \#\{L_\nu\text{ zeros in }|z-(2+iT)|\le R\}
             \le C_R\log(H_\nu+|T|),                           \tag{M28.13}
\]
counting multiplicity. Indeed use a larger concentric
circle of radius at least \(2R\), (M28.12) for its
upper bound and the lower bound (M28.10) at the
center. A boundary avoiding zeros can be chosen in
a fixed bounded radius interval. The logarithmic
ratio in Jensen's inequality stays bounded below.
This counts trivial zeros as well as nontrivial
ones and is uniform at every height, not only
above a character-dependent exceptional range.

\emph{3. Uniform narrow block contours.}
Continue with \(\nu\ne0\); the single polar family
member is included using MT-M26 at the end of step~4.
Choose \(0<\delta<1/4\) so small that
\[
 [b_-,b_+]=[1/2-\delta,1/2+\delta]
                         \subset a+(u,v).
 \]
By the assumed GRH, all zeros of the entire
completion \(\Xi_\nu\) lie on \(\operatorname{Re}z=1/2\).
All other zeros of the finite \(L_\nu\) are the
trivial zeros canceling gamma poles, and have
nonpositive real part. Thus all nontrivial zeros
are separated by \(\delta\) from the two vertical
edges, and all trivial zeros are separated by at
least \(b_->1/4\) from every point of the narrow
closed strip.

Put \(X_{\nu,n}=H_\nu+|n|\). Count the nontrivial
zero ordinates in \([n-2,n+3]\), with multiplicity,
by \(M_{\nu,n}\). Equation (M28.13), with a fixed
sufficient radius, gives \(M_{\nu,n}\le C\log X_{\nu,n}\).
The elementary averaging argument of MT-M18 now
chooses
\[
        h_{\nu,n}\in(n+1/4,n+3/4)
 \]
avoiding every zero ordinate, such that
\[
 \prod_{\Xi_\nu(\rho)=0}
      \min(1,|h_{\nu,n}-\operatorname{Im}\rho|)^{m_\rho}
           \ge e^{-5M_{\nu,n}}\ge X_{\nu,n}^{-C}.             \tag{M28.14}
\]
For completeness, each function
\(\log^+(1/|h-\gamma|)\) has integral at most \(2\)
on an interval of length \(1/2\). Their sum
therefore has a value at most \(5M_{\nu,n}\);
the slack gives a positive-measure set when
\(M_{\nu,n}>0\). Every factor differing from one
comes from the specified ordinate window.

Let \(\Omega_{\nu,n}\) be the rectangle with
vertical edges \(b_\pm\) and horizontal edges
\(h_{\nu,n},h_{\nu,n+1}\); write
\(\Gamma_{\nu,n}=\partial\Omega_{\nu,n}\).
Its perimeter is uniformly bounded. The disk
Blaschke/Harnack proof of MT-M18, now using the
uniform bounds (M28.10)--(M28.13), gives
\[
        \inf_{w\in\Gamma_{\nu,n}}|L_\nu(w)|
                         \ge c X_{\nu,n}^{-B}.               \tag{M28.15}
\]
Here are the details needed for uniformity. Center
a disk at \(2+i\operatorname{Im}w\), with a
zero-free boundary of radius in a fixed interval
large enough that the narrow strip point lies in
the inner half disk. Remove all finite-\(L\)
zeros in this disk by its finite Blaschke product.
The remaining zero-free factor has modulus at
most \(C X_{\nu,n}^B\) on the boundary and at least
\(\zeta_K(2)^{-1}\) at the center. Harnack applied
to the logarithm of this upper bound divided by
its modulus gives a lower bound \(cX_{\nu,n}^{-B}\)
on the inner half disk. There are at most
\(C\log X_{\nu,n}\) Blaschke factors.

On a vertical edge every nontrivial factor has
distance at least \(\delta\) from its zero, and
every trivial factor distance at least \(b_-\).
Each normalized factor is bounded below by this
distance divided by twice the disk radius.
On a horizontal edge, use (M28.14) for the
nontrivial factors, and \(b_-\) for the trivial
factors. Omitting factors from the product in
(M28.14) only increases that product.
These lower bounds, each with \(O(\log X_{\nu,n})\)
factors, prove (M28.15).
The argument includes small heights and trivial
zeros whose ordinates grow with \(\nu\). It
does not absorb any such range into a
character-dependent compactness constant.

\emph{4. Uniform lifts with the completed gamma factors canceled.}
Let
\[
 J_{\nu,n}=A(\Omega_{\nu,n})/
       \{F:F^{(j)}(\rho)=0\
             \text{for every block zero and }j<m_\rho\},
                                                               \tag{M28.16}
\]
with the quotient of the supremum norm on the
closed rectangle. This is the same finite
Banach algebra as in MT-M18. Its dimension is
\[
                  d_{\nu,n}\le C\log X_{\nu,n}.                \tag{M28.17}
\]
If the block is empty its quotient is the zero
space. All assertions below then hold with
the corresponding maps zero.

Use the entire difference-kernel lift of MT-M18,
with its Gaussian parameter fixed to \(1\):
\[
 H_{\nu,n}F(z)=\frac1{2\pi i}
       \int_{\Gamma_{\nu,n}}
       F(w)\frac{\Xi_\nu(z)-\Xi_\nu(w)}
                    {\Xi_\nu(w)(z-w)}
                         e^{(z-w)^2}\,dw.                    \tag{M28.18}
\]
As proved there, it is entire, has the prescribed
jets in this block and zero full jets in all
other blocks, and depends only on the class of
\(F\) in \(J_{\nu,n}\). For each fixed \(\nu\),
the MT-M18 construction and the contour
adaptation in MT-M26 show that it is the
Mellin transform of an element of \(\mathcal E\).
We now prove a uniform estimate for that element.

For an integer \(N\ge0\), take the two inversion
lines
\[
                 \operatorname{Re}z=\pm(N+9/4).              \tag{M28.19}
\]
They lie outside all the rectangles. On them the
Cauchy integral of the \(-\Xi_\nu(w)\) term in
(M28.18) is zero. Only the ratio
\(\Xi_\nu(z)/\Xi_\nu(w)\) remains.
Move each numerator gamma factor
\(\Gamma_{\mathbf R}(z+\mu_{\nu,j})\) to a right
half-plane by a fixed even shift, if necessary.
Every recurrence denominator has real part in
\[
        \pm(N+9/4)+\tfrac12\mathbf Z+2\mathbf Z,
 \]
and therefore absolute value at least \(1/4\).
There is no gamma pole anywhere on either
entire inversion line, for any \(\nu\) or height.
The shifted numerator and the denominator
\(\Gamma_{\mathbf R}(w+\mu_{\nu,j})\) both lie
in a fixed right half-plane.

Use (M28.9) first to change the imaginary part,
then to change the real part by their bounded
difference. Uniformly for \(z=\pm(N+9/4)+it\)
and \(w\in\Gamma_{\nu,n}\), this gives
\[
 \left|
  \frac{\Gamma_{\mathbf R}(z+\mu_{\nu,j})}
       {\Gamma_{\mathbf R}(w+\mu_{\nu,j})}
 \right|
 \le C_N(H_\nu+|t|)^{B_N}
                      e^{C_N|t-\operatorname{Im}w|}.         \tag{M28.20}
\]
The fixed recurrence factors have been included
in \(C_N\). The finite conductor ratio is bounded
by (M28.8). The finite-\(L\) numerator is bounded
by (M28.12) and its denominator by (M28.15).
Thus (M28.7) gives the uniform completed ratio
\[
 \left|\frac{\Xi_\nu(z)}{\Xi_\nu(w)}\right|
 \le C_N X_{\nu,n}^{B_N}
            (H_\nu+|t|)^{B_N}
                      e^{C_N|t-\operatorname{Im}w|}.         \tag{M28.21}
\]
The normalization scalar in (M28.7) cancels
exactly. No bound for an uncanceled completed
function, and in particular no polynomial
bound for a potentially factorial-size gamma
normalization, is asserted or used.

The Gaussian in (M28.18) has modulus
\(\exp((\operatorname{Re}z-\operatorname{Re}w)^2
       -(t-\operatorname{Im}w)^2)\).
The real-part difference is bounded for fixed
\(N\), and
\(|\operatorname{Im}w-n|\le7/4\). The Gaussian
absorbs the exponential linear factor in
(M28.21) and every polynomial in \(|t-n|\).
Also \(H_\nu+|t|\le X_{\nu,n}(1+|t-n|)\).
The denominator \(|z-w|\) is bounded below
by a positive constant on the inversion lines.
Since the contour perimeter is bounded, we obtain
\[
 |H_{\nu,n}F(\pm(N+9/4)+it)|
 \le C_N X_{\nu,n}^{B_N}\|F\|_\infty
                              e^{-c(t-n)^2},                 \tag{M28.22}
\]
with \(c>0\) independent of \(\nu,n,N\), after
reducing it once.

Invert the Mellin transform on these lines.
For \(x\ge0\) use the positive line and for
\(x\le0\) the negative line. The \(j\)-th
derivative costs \(|z|^j\), whose Gaussian
integral is bounded by a constant times
\((2+|n|)^j\). Since the resulting function
is already in \(\mathcal E\) for each
individual \(\nu\), the contour shifts are
legitimate; (M28.22) now supplies their
uniform constants. Taking the infimum over
all representatives \(F\) of a fixed block
class yields linear lifts
\(\mathcal T_{\nu,n}:J_{\nu,n}\to\mathcal E\)
with
\[
 p_{N,j}(\mathcal T_{\nu,n}b)
       \le C_{N,j}X_{\nu,n}^{B_{N,j}}\|b\|_{J_{\nu,n}},
 \qquad R_{\nu,m}\mathcal T_{\nu,n}=\delta_{mn}\operatorname{id}.
                                                               \tag{M28.23}
\]
Here \(p_{N,j}(h)=\sup_x e^{N|x|}|h^{(j)}(x)|\).
The last identity means the identity in the
matching block and zero in the others.
The exceptional \(\nu=0\) has the corresponding
fixed-character polynomial bounds by MT-M26;
enlarging the constants once includes it.
These are stable quotient-block estimates,
with no separation condition on zeros inside
one rectangle.

\emph{5. The actual quotient and its full zero-jet kernel.}
The compact character projectors are orthogonal
contractions on \(H_{I,s}\). The denominator
of (M28.1) is invariant under them: MT-M24
identifies \(\Theta_K\) with the intersection
of the two summation spaces, and compact
averaging preserves that intersection. Applying
a projector to an approximating sequence
shows that its intersection with one character
sector is exactly the Hilbert closure of the
corresponding sector image.

Apply MT-M26 in each sector, using (M28.2).
It identifies that closure with all the
sector zero jets. Every such jet is a
continuous functional on the full Hilbert
space, by the compact projection and the
interior-strip Cauchy--Schwarz estimate.
Finite character truncations converge
strongly in the Hilbert direct sum.
Therefore the closed denominator is exactly
the simultaneous kernel (M28.4).

The full smooth closed-kernel identity of
MT-M27 was proved by compact Fourier
truncation and MT-M24, independently of
the Hilbert strip width. It continues to
apply here. It proves faithfulness and
continuity of (M28.5); compactly supported
smooth radial functions and finite
character sums prove density. If this
map were onto, applying any one compact
character projector would contradict
the proper-image assertion of MT-M26.
This proves the stated proper dense
smooth inclusion without equating
the smooth and Hilbert topologies.

\emph{6. Trace-norm summation over characters and blocks.}
Let \(\pi:H_{I,s}\to Z_{a,I,s}\) be the quotient map.
The block restriction map, after compact
character projection, is
\[
  r_{\nu,n}[h]
       =\left[w\longmapsto
          \widehat h(\beta_\nu|\cdot|^{\,w-a})\right]
           \in J_{\nu,n}.
 \]
The real parts \(w-a\) lie in a fixed compact
subinterval of \((u,v)\). Cauchy--Schwarz
against the radial weight, and the fact
that the compact character has modulus
one, therefore give
\[
                \|r_{\nu,n}\|\le C                         \tag{M28.24}
\]
uniformly in both indices. As in MT-M26,
this follows by bounding the quotient
norm by the supremum of the displayed
representative on the contour.

For \(b\in J_{\nu,n}\), define
\[
  s_{\nu,n}b
       =\pi\left[
          \beta_\nu(g)^{-1}|g|^a
               (\mathcal T_{\nu,n}b)(\log|g|)
              \right].
 \]
A single fixed seminorm \(p_{N,0}\), with
\(N>|a|+\max(|u|,|v|)+1\), bounds its
Hilbert norm. Equation (M28.23) gives
\[
 \|s_{\nu,n}\|\le C X_{\nu,n}^B,\qquad
 r_{\mu,m}s_{\nu,n}
      =\begin{cases}\operatorname{id},&(\mu,m)=(\nu,n),\\
                      0,&\text{otherwise}.
        \end{cases}                                          \tag{M28.25}
\]
No compact-direction derivatives occur
in this Hilbert norm, so no extra
character factor is concealed here.

For a test \(f\) invariant under \(k\), let
\[
 f_\nu(x)=\int_{C_K^1}
          f(c\,\tau(e^x))\,\beta_\nu(c)\,dc,\qquad
 b_{\nu,n}(f)=
       \sup_{w\in\Gamma_{\nu,n}}
                      |\mathcal M f_\nu(w-a)|.
 \]
For every \(A>0\),
\[
                   b_{\nu,n}(f)\le C_A(f)X_{\nu,n}^{-A},       \tag{M28.26}
\]
where \(C_A(f)\) is bounded by finitely
many seminorms at this test level and
is uniformly bounded on bounded test
sets. To see this directly, integration
by parts on the torus, using
\((1-\Delta_T)^r\), gives arbitrary
negative powers of \(1+|\nu|\).
Integration by parts in \(x\), on the
common compact support, gives arbitrary
negative powers of \(1+|\operatorname{Im}w|\),
uniformly for the real part in the
fixed contour interval. The derivatives
of \(e^{(\operatorname{Re}w-a)x}f\)
are bounded by finitely many test
seminorms. Since
\(|\operatorname{Im}w-n|\le7/4\),
combining these two arbitrary powers
gives (M28.26). The finite group in
(M28.6) only changes a constant.

Let \(M_{f,\nu,n}\) be multiplication
in \(J_{\nu,n}\) by
\([\mathcal Mf_\nu(w-a)]\). Its norm
is at most \(b_{\nu,n}(f)\), by the
submultiplicative quotient algebra
norm. The finite-rank operator
\[
                 A_{\nu,n}(f)
                    =s_{\nu,n}M_{f,\nu,n}r_{\nu,n}
 \]
has rank at most \(d_{\nu,n}\). Hence,
as an operator on the Hilbert quotient,
\[
 \|A_{\nu,n}(f)\|_1
       \le C\,d_{\nu,n}X_{\nu,n}^B b_{\nu,n}(f).
                                                               \tag{M28.27}
\]
This is the elementary Hilbert estimate
\(\|A\|_1\le\operatorname{rank}(A)\|A\|\);
the intervening finite-dimensional
spaces need not carry Hilbert norms.
Equations (M28.17) and (M28.26) show
that
\[
                  \sum_{\nu\in\Lambda}\sum_{n\in\mathbf Z}
                                      A_{\nu,n}(f)            \tag{M28.28}
\]
converges in trace norm. For example,
take \(A>B+d+3\) in (M28.26), since
the lattice plus the height index
has dimension \(d+1\). Its tails
converge uniformly on bounded test
sets, and the resulting trace norm
is bounded by finitely many test
seminorms.

This is the actual convolution, not
only a formal operator with the
same eigenvalues. Convolution by
\(f\) is bounded on \(H_{I,s}\):
the translation estimate (M26.10)
integrated over its compact support
proves this. It preserves the closed
denominator, since
\[
 C_f(|\cdot|^a h)
       =|\cdot|^a C_{|\cdot|^{-a}f}h
 \]
and the summation image is a
convolution submodule.
On a single-character smooth vector
the operator (M28.28) equals this
convolution by the already proved
one-character synthesis and dense
extension in MT-M26. On characters
not in \(\Lambda\), both operators
are zero: \(k\)-invariance of \(f\)
makes its Fourier coefficient zero.
Finite-character smooth vectors
are dense in the full Hilbert
quotient. Both operators are
bounded, so their equality extends
to that quotient. This completes
the full adelic trace-class proof.

\emph{7. Trace, uniform convergence, and the analytic scope.}
By (M28.25), the finite-rank trace
of \(A_{\nu,n}(f)\) is the trace
of multiplication by
\([\mathcal Mf_\nu(w-a)]\) in
\(J_{\nu,n}\). The local algebra
decomposes as the direct sum of
\(\mathbf C[z]/(z-\rho)^{m_\rho}\)
over its block zeros. Multiplication
is triangular in the power bases,
with diagonal value
\(\mathcal Mf_\nu(\rho-a)\) repeated
\(m_\rho\) times. Thus this trace
is the corresponding finite part
of (M28.3). Trace-norm convergence
allows summation of these traces.
The individual spectral terms
are absolutely summable as well:
\[
 \sum_{\rho\text{ in block}}
    m_\rho|\mathcal Mf_\nu(\rho-a)|
            \le d_{\nu,n}b_{\nu,n}(f),
 \]
by the maximum principle on the
block. Equations (M28.17) and
(M28.26) give uniform tails of
this double series directly.
This proves all the asserted
trace and continuity properties.

For a specified bounded test set
at the level \(k\), its trace-class
and trace-formula conclusions use
GRH only for the characters in
\(\widehat{C_K^1/k}\). Other
character sectors are annihilated
by these tests. The global
faithfulness statement (M28.4)
uses all character sectors.
Thus the full conditional input
\(M\) is sufficient, while the
operator conclusion at one level
has the stated smaller Hecke
family as its precise analytic
input.

At \(a=1/2\), define
\(h^*(g)=\overline{h(g^{-1})}\).
The group is abelian and
unimodular. Under the assumed
GRH, every
\(\beta_\chi|\cdot|^{\rho-1/2}\)
is unitary, and (M28.3) gives
\[
       \operatorname{Tr} C_{h*h^*}
          =\sum_{\chi,\rho}
                 m_{\chi,\rho}
                 |\widehat h(
                    \beta_\chi|\cdot|^{\,\rho-1/2})|^2
          \ge0.                                               \tag{M28.29}
\]
This is a trace identity. It
does not identify the group
involution with the operator
adjoint for the specified
Hilbert inner product.
All higher zero jets remain
present in (M28.4), consistently
with the obstruction in MT-M20
to a faithful positive
full-jet realization when
multiple zeros occur.
\end{proof}

\paragraph{Run and geometric crosswalk.}
This is a new uniform transfer proof in
\path{GRH_NIGHT/complexity}, following the fixed-character
MT-M18/MT-M24/MT-M26 and the full wide-strip MT-M27.
Its literal source targets are Masullo
Definition~9.2.2.31, original PDF p.~416, the Hilbert
closure in Definition~10.1.1.1, p.~434, and the
strip quotient in Definition~10.4.1.3, p.~456,
with the centering twist from MT-M24. It closes
the all-character positive-width strip gate
conditionally on \(M\); it is not a deduction
of \(M\). In the attached geometric programme,
it supplies an actual trace-class realization
for the idele/test-function quotient interface
of the \path{cote1}/\path{SA23geo-H-lidskii}
crosswalk. The positivity statement remains
the conditional trace statement (M28.29);
the stronger faithful positive-inner-product
gate is still exactly MT-M20.

\subsection*{MT-NB01. Quantitative approximation after the adopted Riemann hypothesis}

The qualitative Nyman--Beurling and B\'aez-Duarte criteria are
already background in \path{RETURN_ALPOGE_BRIDGES_cp02},
CC-P01--CC-P09. Here the additional issue is the rate, including
the contribution of multiple zeros. Put
\[
 d_N^2=\inf_{a_1,\ldots,a_N\in\mathbf C}
 \frac1{2\pi}\int_{\mathbf R}
 \left|1-\zeta(\tfrac12+it)
       \sum_{n\le N}a_n n^{-1/2-it}\right|^2
                 \frac{dt}{1/4+t^2}.
 \tag{NB01.1}
\]
The adopted input $A_D$ and the classical theorems below give,
for every $\epsilon>0$,
\[
 d_N^2\ll_\epsilon
 \frac{(\log\log N)^{5/2+\epsilon}}{\sqrt{\log N}},
 \qquad
 \liminf_{N\to\infty}d_N^2\log N
       \ge C_\zeta^{(2)}
       :=\sum_{\rho\ {
m distinct}}
                    \frac{m_\rho^2}{|\rho|^2}.
 \tag{NB01.2}
\]
The upper bound is Balazard--de Roton,
\emph{Sur un crit\`ere de B\'aez-Duarte pour l'hypoth\`ese
de Riemann}, Int. J. Number Theory 6(4) (2010), 883--903,
unnumbered main Theorem in \S1; author version
\path{arXiv:0812.1689v1}, p.~3. Its constructive choice
appears in \S2, Propositions~1--3, author pp.~4--5:
one can take
\[
 a_n=\mu(n)n^{-\varepsilon_N},\qquad
 \varepsilon_N=25(\log\log N)^{5/2+\epsilon}
                         (\log N)^{-1/2}
 \tag{NB01.3}
\]
for all sufficiently large $N$. The implied constants are
effective in the source's convention; no numerical constant
or numerical threshold is extracted here.
The lower bound is Burnol, \emph{A lower bound in an
approximation problem involving the zeros of the Riemann
zeta function}, Adv. Math. 170 (2002), 56--70,
Theorem~1.3, author version \path{math/0103058v2}, p.~4.
That theorem uses continuous dilations $\lambda\le\alpha\le1$;
the discrete space in (NB01.1) is a subspace at $\lambda=1/N$,
so its distance is at least the continuous distance.
The Mellin--Plancherel identification is stated and proved in
Balazard--de Roton, Proposition~1, author pp.~4--5.

The familiar sum over zeros with their ordinary multiplicities
is, under $A_D$,
\[
 C_\zeta^{(1)}:=\sum_\rho\frac{m_\rho}{|\rho|^2}
                    =2+\gamma-\log(4\pi).
 \tag{NB01.4}
\]
An exact locator for the identity in this normalization is
Balazard--de Roton, author p.~3, immediately before (3),
and Bettin--Conrey--Farmer, author p.~2.
Thus $C_\zeta^{(2)}\ge C_\zeta^{(1)}$, with equality precisely
when every nontrivial zero is simple. The adopted RH input
does not remove this distinction.

For comparison, Bettin--Conrey--Farmer,
\emph{An optimal choice of Dirichlet polynomials for the
Nyman--Beurling criterion}, Proc. Steklov Inst. Math.
280, Suppl.~2 (2013), S30--S36, Theorem~1,
\path{arXiv:1211.5191v1}, p.~2, proves the sharper conclusion
\[
 d_N^2\sim\frac{2+\gamma-\log(4\pi)}{\log N}
 \tag{NB01.5}
\]
under RH and the additional hypothesis
\[
 \sum_{|\operatorname{Im}\rho|\le T}
                    |\zeta'(\rho)|^{-2}
                  \ll T^{3/2-\delta}\quad\text{for some }\delta>0.
 \tag{NB01.6}
\]
The hypothesis includes simplicity. Their explicit polynomial is
$V_N(s)=\sum_{n\le N}(1-\log n/\log N)\mu(n)n^{-s}$.
The theorem gives this polynomial's error asymptotic; Burnol's
lower bound then gives (NB01.5) for the infimum. These are
classical results with their exact remaining inputs, not new
proofs of the cited analytic estimates.

\subsection*{MT-NB02. The multiplicity-sensitive Hecke approximation constant}

Let $K$ be a number field, and let $F(s)=L_K(s,\eta)$ be a
primitive unitary Hecke $L$-function. Allow $\eta=1$, but exclude
the nontrivial pure norm characters $|\cdot|^{it}$, whose pole
would lie away from $1$. For such a pure norm character the
following statement instead applies to its vertical translate
$\zeta_K(s)$. Write $F(s)=\sum_{n\ge1}a_n n^{-s}$ and define
\[
 \Psi_F(x)=\operatorname*{Res}_{s=1}
                   \frac{x^s F(s)}s-\sum_{n\le x}a_n,
 \qquad
 B_F^\lambda=\operatorname{span}\{
                  t\longmapsto\Psi_F(\alpha/t):
                               \lambda\le\alpha\le1\}.
 \tag{NB02.1}
\]
Assume GRH for this particular $F$, supplied by $A_D$ for a
primitive Dirichlet function and by the conditional input
$M_H$ in the general Hecke case. Then $\Psi_F(1/t)$ belongs
to $L^2((0,\infty),dt)$, the distance
\[
 D_F(\lambda)=\operatorname{dist}_{L^2(0,\infty)}
                      (\mathbf1_{(0,1)},B_F^\lambda)
 \tag{NB02.2}
\]
tends to zero, and
\[
 \liminf_{\lambda\downarrow0}D_F(\lambda)^2\log(1/\lambda)
    \ge C_F^{(2)}:=\sum_{\rho\ {
m distinct}}
                          \frac{m_\rho^2}{|\rho|^2}<\infty.
 \tag{NB02.3}
\]
The zeros in this formula are the nontrivial zeros of $F$,
all on the critical line under the stated input.
The precise classical theorem being specialized is
Anne de Roton, \emph{Une approche hilbertienne de
l'hypoth\`ese de Riemann g\'en\'eralis\'ee}, Bull. Soc.
Math. France 134(3) (2006), 417--445: Definition~1 p.~418,
Definition~2 and Theorem~1 p.~419, Theorem~2 p.~420.
The lower-limit argument is \S7, pp.~441--444, particularly
Propositions~7.1--7.2 and \S7.2. The original p.~420 image
and that complete final proof were read. Earlier technical
operator estimates remain cited prerequisites.

\begin{proof}[Complete proof of the stated Hecke specialization]
We check the source's Selberg-class hypotheses, rather than
assuming that an arbitrary meromorphic Artin quotient satisfies
them. Grouping the ideal Dirichlet series by absolute norm gives
the displayed coefficients. If $n_K=[K:\mathbf Q]$, their
absolute values are bounded by $d_{n_K}(n)$: at each rational
prime the local ideal-counting coefficients are bounded by those
of $(1-p^{-s})^{-n_K}$. Consequently $a_n\ll_\epsilon n^\epsilon$.
The unitary Euler factors have roots of modulus at most one.
For a rational prime $p$, the logarithmic coefficient at $p^r$
is a sum over prime ideals of residue degree $f\mid r$,
with each summand bounded by $f/r$. Since the sum of the
residue degrees is at most $n_K$, its absolute value is at
most $n_K/r$. All other logarithmic coefficients vanish.
This verifies the source's Euler-product bound with exponent
zero, strictly smaller than $1/2$.

The classical Hecke continuation, finite-order property and
functional equation supply the remaining axioms; see the exact
Ramakrishnan--Valenza Theorem~7-19 and local-factor locators
in MT-M28. The real gamma shifts have nonnegative real part,
and the complex shifts have real part $|m|/2\ge0$.
There is no pole for a non-pure-norm character. The allowed
trivial character gives $\zeta_K$, with its single pole at $1$.
Thus the permitted functions belong to the precise Selberg
class of de Roton's Definition~1. Her Theorem~1, with the
specified GRH, gives the $L^2$ assertion and convergence of
the distance defined in (NB02.1)--(NB02.2). Her Theorem~2
then gives the lower bound; squaring a nonnegative lower-limit
inequality yields (NB02.3).

For clarity, convergence of the multiplicity-square sum also
follows directly from the fixed-function local zero bound in
NT-Z10. In every unit height interval the sum of multiplicities
is $O_F(\log(2+|T|))$. The sum of their squares is at most
the square of that sum. The contribution of intervals of
height at least two is therefore bounded by a constant times
$\sum_{j\ge2}j^{-2}\log^2j$, which converges. The finitely many
bounded-height zeros cause no problem, since a critical-line
zero has $|\rho|\ge1/2$.
\end{proof}

\subsection*{MT-NB03. Integral Galois classes force an approximation penalty}

For a function in MT-NB02 set
\[
 C_F^{(1)}=\sum_\rho\frac{m_\rho}{|\rho|^2},\qquad
 \Delta_F=\sum_\rho\frac{m_\rho(m_\rho-1)}{|\rho|^2}.
 \tag{NB03.1}
\]
Then $C_F^{(2)}=C_F^{(1)}+\Delta_F$, and
$\Delta_F=0$ precisely when all its nontrivial zeros are simple.
For a tame finite Galois extension $N/K$ put
$U_{N/K}=[\mathcal O_N]-[\mathcal O_K[G]]$
in the locally free class group of $\mathbf Z[G]$, with the
restriction-of-scalars convention of MT-ART11. If
$U_{N/K}\ne0$, then, conditionally on GRH for $\zeta_N$,
\[
 \begin{split}
 \liminf_{\lambda\downarrow0}D_{\zeta_N}(\lambda)^2
                      \log(1/\lambda)&\ge C_{\zeta_N}^{(2)},\\
 C_{\zeta_N}^{(2)}&\ge16,\qquad
 C_{\zeta_N}^{(2)}\ge C_{\zeta_N}^{(1)}+8.
 \end{split}                                                   \tag{NB03.2}
\]
The same constants apply to the explicit tame quaternionic
field of MT-ART10, directly from its proved central order
$r\ge2$; this use does not identify nonfreeness with a nonzero
stable class. Independently of $U_{N/K}$, if $N/K$ is nonabelian
Galois, then $\Delta_{\zeta_N}>0$ under that same GRH input.
Thus a multiplicity-one constant is strictly too small in
all these nonabelian examples.

\begin{proof}
All the sums converge by the argument in MT-NB02. The identity
in (NB03.1) follows term by term from $m^2=m+m(m-1)$.
Every summand of $\Delta_F$ is nonnegative, and it vanishes
exactly when its positive integer multiplicity equals one.

MT-ART11 proves unconditionally, without Artin holomorphy,
that $U_{N/K}\ne0$ forces an even positive central order
$r=\operatorname{ord}_{s=1/2}\zeta_N(s)\ge2$. Its proof uses
Taylor's root-number class theorem and the cyclic-restriction
argument recorded there. The single central zero contributes
$4r^2\ge16$ to $C_{\zeta_N}^{(2)}$ and
$4r(r-1)\ge8$ to $\Delta_{\zeta_N}$. This proves (NB03.2)
by MT-NB02. The arithmetic central-zero assertion is
unconditional; GRH enters to obtain the approximation-space
conclusion and place every nontrivial zero in that formula.

For any nonabelian $N/K$, the classical
Hu--Kaneko--Martin--Schildkraut Theorem~1.1, already located
and used in MT-M23, gives infinitely many multiple nontrivial
zeros of $\zeta_N$. Under the present GRH each contributes a
strictly positive term to $\Delta_{\zeta_N}$; even one would
suffice. Hence the claimed strict inequality follows. No
uniform numerical lower bound is asserted in this last case.
\end{proof}

\subsection*{MT-NB04. A precise derivative exponent sufficient for the optimal rate}

Suppose in addition to $A_D$ that, for some $c>0$ and
$0\le\beta<1/4$, every nontrivial zero satisfies
\[
             |\zeta'(\rho)|\ge c(1+|\operatorname{Im}\rho|)^{-\beta}.
 \tag{NB04.1}
\]
Then all the zeros are simple, (NB01.6) holds for some
$\delta>0$, and the optimal asymptotic (NB01.5) follows.

\begin{proof}
The displayed positive lower bound implies simplicity. The
classical zero-count bound $N(T)=O(T\log(2+T))$, with the
same multiplicity convention as NT-Z10, gives for $T\ge2$
\[
 \sum_{|\operatorname{Im}\rho|\le T}|\zeta'(\rho)|^{-2}
                 \ll_{c,\beta}T^{1+2\beta}\log T.
\]
Take $\delta=1/4-\beta>0$. Since $\log T\ll_\delta T^\delta$,
the right side is $O(T^{5/4+\beta})=O(T^{3/2-\delta})$.
Bettin--Conrey--Farmer Theorem~1 gives the error asymptotic
for $V_N$, and Burnol gives the matching lower bound for
the infimum, as in MT-NB01. This proves the conclusion.
\end{proof}

This sufficient exponent condition is stronger than the
unspecified polynomial reciprocal-derivative bound forced by
the point-jet basis in MT-M11. It is not a consequence here
of $A_D$, of the repaired adapted-block basis, or of simple
zero locations alone. It makes the additional quantitative
condition needed by that particular route explicit.

\paragraph{Geometric and run placement of MT-NB01--04.}
The Hecke functions are the rank-one adelic objects of Tate's
thesis; for $\zeta_N$ the object is the pushforward of the unit
Artin motive along $\operatorname{Spec}N\to\operatorname{Spec}K$.
The functions $\Psi_F$ turn its coefficient-counting distribution
into a dilation subspace. The zero divisor, including its
orders, obstructs approximation by that subspace. In de Roton's
proof each order-$m$ zero supplies a full $m$-jet Gram block;
the inverse block's first diagonal entry is $m^2$ (\S7.2,
p.~444). This is the analytic reason that discarding higher
jets loses the constants in MT-NB03, consistent with MT-M20
and MT-M25.

For the archive programme, this is an approximation-theoretic
addition to the trace/Fredholm interface of cote~1, scans
79--80 and 111--112, inherited through
\path{OURS_INTAKE_RESEARCH_SA23_20261120.md} \S4.1,
\path{geo-H-lidskii}/\path{geo-H-genus}; the sheets themselves
are not cited as Nyman--Beurling theorems. MT-NB03 also uses
the integral $K_0$ comparison of MT-ART11 and
AG-LEFSCH-LATTICE/DEFECT, placed at cote~16 scans~14--18,
no.~7, and SGA6 Expos\'e~V \S2 and \S7 (7.1.1), with the
same inherited-locator qualification as MT-ART11. The precise
run extension is the quantitative completion of the qualitative
CC-P01--CC-P09 bridge, with the negative derivative moment
(NB01.6) retained as the remaining optimal-rate condition.

\subsection*{MT-M29. The determinant of the graded arithmetic spectral module}
Fix one unitary Hecke character and its pole-removed entire
completion \(\Xi\). Work first in the centered, actual wide-strip
Hilbert quotient of MT-M24/MT-M26, with \(a=1/2\) and
\([0,1]\subset1/2+\operatorname{int}I\). Its definition by adelic
summation and Hilbert closure is unconditional. Write \(U_y\)
for norm translation by \(\tau(e^y)\), and \(D\) for its
infinitesimal generator. Its dense smooth quotient is the
centered version of \(Q_\eta\).

There is a canonical algebraic submodule \(V_{\mathrm{fin}}\):
the vectors contained in finite-dimensional \(U_y\)-invariant
subspaces. It is dense in the smooth quotient and in the
specified Hilbert quotient. Its primary decomposition and
canonical filtration are
\[
 V_{\mathrm{fin}}=\bigoplus_\rho V_\rho,\qquad
 V_\rho\simeq\mathbf C[u]/(u^{m_\rho}),\qquad
 D|_{V_\rho}=(\rho-\tfrac12)+u,\qquad
 F^jV_\rho=(D-\rho+\tfrac12)^jV_\rho .
                                                               \tag{M29.1}
\]
The isomorphism uses the local Mellin coordinate
\(u=z-(\rho-1/2)\). Form the associated graded
\[
       V_{\mathrm{gr}}
          =\bigoplus_\rho\bigoplus_{j=0}^{m_\rho-1}
                         F^jV_\rho/F^{j+1}V_\rho.             \tag{M29.2}
\]
This construction is canonical and functorial at each finite
spectral stage for maps intertwining norm translation. It
retains \(m_\rho\) dimensions at \(\rho\); it replaces the
nilpotent extensions by scalar actions.

Assume now GRH for this function. It is supplied by \(A_D\)
for a Dirichlet function over \(\mathbf Q\), or by the
conditional input \(M\) for the general Hecke character.
Give the one-dimensional graded lines in (M29.2) the metric
in which the class of \(u^j\) has norm one, and make them
orthogonal. The resulting Hilbert completion
\(\mathcal H_{\mathrm{gr}}\) carries the unitary norm action
\[
        U_y^{\mathrm{gr}}=e^{iyT},\qquad
        T|_{(\rho,j)}=\gamma,\qquad \rho=\tfrac12+i\gamma.
                                                               \tag{M29.3}
\]
Every integrated \(\mathcal E\)-test operator is trace class,
with the same trace as on the actual wide-strip quotient:
\[
 \operatorname{Tr}_{\mathcal H_{\mathrm{gr}}}C_f^{\mathrm{gr}}
       =\sum_\rho m_\rho\,\mathcal Mf(\rho-\tfrac12)
       =\operatorname{Tr}_{Z_{1/2,I,s}} C_f .                  \tag{M29.4}
\]
The grading and this new Hilbert completion change the
representation when a multiple zero occurs; they do not
equip the original Jordan blocks with a positive invariant
metric.

Suppose in addition that
\[
 \Xi(\tfrac12+z)=\epsilon\Xi(\tfrac12-z),\qquad
       \epsilon\in\{1,-1\}.
                                                               \tag{M29.5}
\]
Delete the central spectral block, then take the
positive-frequency spectral subspace
\(\mathcal H_+=1_{(0,\infty)}(T)\mathcal H_{\mathrm{gr}}\).
On it,
\[
                          B=T_+^{-2}                         \tag{M29.6}
\]
is a bounded positive trace-class operator. If
\(r=\operatorname{ord}_{s=1/2}\Xi(s)\), set
\[
 H(t)=t^{-r}\Xi(\tfrac12+it),\qquad
              \mathfrak D(w)=H(\sqrt w)/H(0).
 \]
Then the operator from the actual arithmetic module has
\[
                       \det(1-wB)=\mathfrak D(w).             \tag{M29.7}
\]
This is the AG-DET operator, now identified as a functional
calculus of the norm generator on the graded finite
spectral module. Its numerical multiplicities are precisely
the original zero orders.

\begin{proof}
\emph{The finite spectral core comes from the actual quotient.}
MT-M18 and MT-M24 identify the smooth quotient with its full
zero jets. For a finite set \(S\) of zeros, there is a unique
class with any prescribed jets at \(S\) and zero jets at
every other zero: its existence follows from the finitely
many block lifts of MT-M18, and uniqueness from the exact
closed-ideal equality. These classes form a finite-dimensional
subalgebra \(V_S\), canonically isomorphic to
\[
                     \prod_{\rho\in S}
                              \mathbf C[u]/(u^{m_\rho}).
                                                               \tag{M29.8}
\]
It has a canonical unit \(e_S\), whose local class is one
at \(S\) and zero elsewhere. Multiplication by \(e_S\) is
the finite-jet quotient projector. Thus \(V_S\) is both
a subspace and a finite quotient of the actual smooth
arithmetic quotient; no choice of interpolation
representatives survives in (M29.8).

These finite projectors extend boundedly to the actual
wide-strip Hilbert quotient: its finite jet restrictions
are continuous by MT-M26, and the finitely many lifted
classes have finite Hilbert norm. This proves that \(V_S\)
is a genuine finite spectral submodule there as well.
The finite unions of the MT-M18 blocks exhaust all zeros,
and their partial sums converge in the smooth quotient.
Hence \(\bigcup_S V_S\) is dense in that quotient and,
by its dense inclusion, in the Hilbert quotient.

Norm translation acts in the Mellin transform by
multiplication by \(e^{zy}\) in the centered variable.
Consequently its action on the \(\rho\)-factor in (M29.8)
is multiplication by
\[
                  e^{(\rho-1/2+u)y}
                   =e^{(\rho-1/2)y}e^{uy}.
 \]
Differentiating this finite matrix exponential gives
the generator in (M29.1).

It remains to verify that this core is intrinsic, rather
than merely a specified collection of finite subspaces.
Let \(W\) be any finite-dimensional invariant subspace of
the Hilbert quotient. The continuous one-parameter group
on \(W\) is a matrix exponential; write \(D_W\) for its
matrix generator, and let \(P\) be its minimal polynomial.
Every finite block restriction intertwines the norm
group, and therefore intertwines its generator on \(W\).
On a local factor with \(P(\rho-1/2)\ne0\),
\[
                 P((\rho-1/2)+u)
 \]
is invertible, because its constant term is nonzero.
Thus the restriction of every vector of \(W\) to this
factor is zero. Only the finitely many zeros whose
centered parameters are roots of \(P\) can remain.
If \(S\) is that finite set, subtracting the finite
projector \(e_S\) kills all jets of a vector in \(W\).
The exact Hilbert kernel statement of MT-M26 forces
that difference to be zero. Therefore \(W\subset V_S\).
Conversely every \(V_S\) is finite dimensional and
invariant, so \(V_{\mathrm{fin}}=\bigcup_S V_S\).
This establishes the intrinsic description.

\emph{The filtration and its functoriality.}
For the primary summand \(V_\rho\), put
\(N_\rho=D-(\rho-1/2)\). Its image powers
\(N_\rho^j V_\rho\) are the filtration in (M29.1);
in the coordinate algebra they are \((u^j)\).
Any finite-stage intertwining map commutes with \(D\).
It preserves the primary decomposition and sends
\(N_\rho^j V_\rho\) into the corresponding image power
in its target. It therefore induces a map on every
graded quotient, compatibly with compositions.
The inclusions as \(S\) grows have the same property.
This proves the claimed algebraic functoriality.
No boundedness assertion for arbitrary infinite-stage
intertwiners is required.

On \(F^jV_\rho/F^{j+1}V_\rho\), the nilpotent part
\(N_\rho\) acts as zero, so \(U_y\) acts by
\(e^{(\rho-1/2)y}\). Each of these quotients has
dimension one, and there are exactly \(m_\rho\)
of them. The distinguished vector \([u^j]\)
uses the local unit in (M29.8) and the specified
global Mellin coordinate; it is independent of
any chosen theta-function representative.
Declaring these vectors orthonormal defines
the metric stated above.

\emph{The Hilbert action and its trace.}
Under GRH the centered parameter is \(i\gamma\).
Thus every norm-translation scalar on the associated
graded has modulus one. The Hilbert completion is
explicitly
\[
             \ell^2\{(\rho,j):0\le j<m_\rho\}.
 \]
The norm group acts by the scalars \(e^{i\gamma y}\).
It is strongly continuous, first on finite-support
vectors and then on all vectors by their density
and the unit norm of the action.
Its self-adjoint frequency generator is multiplication
by \(\gamma\), with domain
\[
       \left\{(a_{\rho,j}):
                    \sum_{\rho,j}\gamma^2|a_{\rho,j}|^2
                                      <\infty\right\}.
 \]
Self-adjointness follows directly: the adjoint of
this multiplication operator has exactly that
domain, since testing against the finite-support
coordinate vectors determines its coordinates.
This proves (M29.3).

For \(f\in\mathcal E\), the integrated action is
diagonal with entry \(\mathcal Mf(i\gamma)\)
on each of the \(m_\rho\) graded lines. The transform
decreases faster than every power of \(|\gamma|\).
The fixed-function local zero count gives
\(\sum_\rho m_\rho(2+|\gamma|)^{-2}<\infty\).
Hence this diagonal operator is trace class, with
trace and trace norm obtained by summing its
entries and their absolute values, respectively.
MT-M26 gives the identical trace on the actual
wide-strip Hilbert quotient. This proves (M29.4).
The proof also gives continuity into trace norm
for the \(\mathcal E\) test topology.

The central subspace has dimension \(r<\infty\).
After it is removed, positive zero ordinates
are bounded away from zero, since zeros of a
nonzero entire function have no finite accumulation
point. Thus \(T_+^{-2}\) is bounded and positive,
with eigenvalues \(\gamma^{-2}\), one for each of
the \(m_\rho\) graded lines at each positive
ordinate. Their sum converges by the same
zero count. This proves (M29.6), including trace
class and compactness. If there are no positive
zeros the space is zero and its determinant is one.

By (M29.5), the noncentral ordinates occur in
pairs \(\gamma,-\gamma\) with equal multiplicities.
AG-DET, whose complete factorization proof
applies to this entire order-at-most-one
completion, states precisely
\[
        \mathfrak D(w)=
          \prod_{\gamma>0}(1-w/\gamma^2)^{m_{1/2+i\gamma}}.
 \]
The spectrum of (M29.6), with its just-proved
graded multiplicities, gives this same locally
uniformly convergent Fredholm product.
This proves (M29.7) and the comparison with
AG-DET. It uses no assertion about individual
Hermite vectors forming a Schauder basis.
\end{proof}

\paragraph{The ordinary trace-GNS completion is different.}
The distinction is essential when zeros have multiplicity.
On the centered convolution algebra \(Q\), let
\[
       \tau(f)=\operatorname{Tr}_{Z_{1/2,I,s}}C_f,\qquad
       f^*(x)=\overline{f(-x)}
 \]
in the fixed-character radial notation. The Hecke
functional equation makes the centered ideal invariant
under this involution. Under the same GRH input,
\[
 \langle[f],[g]\rangle_\tau
          =\tau(f*g^*)
          =\sum_{\rho}m_\rho\,
                    \mathcal Mf(i\gamma)
                        \overline{\mathcal Mg(i\gamma)} .
                                                               \tag{M29.9}
\]
Its radical is exactly the ideal of classes whose
values, without their higher derivatives, vanish at
every zero. Indeed all weights in the sum for
\(f=g\) are positive, so zero norm is equivalent
to all those values vanishing. Finite interpolation
from (M29.8) gives every finitely supported list
of values. The completion of the quotient by
this radical is therefore
\[
       \mathcal H_{\mathrm{GNS}}
             \simeq\ell^2(\{\text{distinct }\rho\},m_\rho).
                                                               \tag{M29.10}
\]
The evaluation map is continuous from \(Q\):
rapid Mellin decay and the local zero count
bound the squared norm by finitely many
\(\mathcal E\) seminorms, and it vanishes on
the closed ideal. Thus the density assertion
in (M29.10) is a statement about the actual
arithmetic quotient, not only an abstract
diagonal model.

Multiplication by a test transform on (M29.10)
has one Hilbert dimension per distinct zero.
The normalized coordinate vector at \(\rho\)
is \(m_\rho^{-1/2}e_\rho\); its diagonal matrix
entry is still \(\mathcal Mf(i\gamma)\).
Consequently its ordinary Hilbert trace counts
each distinct zero once. The integer \(m_\rho\)
in (M29.9) is a weight in the spectral measure,
not a Hilbert spectral multiplicity.
It follows that the positive-frequency inverse
square on this GNS space has ordinary
Fredholm determinant
\[
                \prod_{\substack{\gamma>0\\\text{distinct}}}
                                      (1-w/\gamma^2).         \tag{M29.11}
\]
It agrees with (M29.7) if and only if every
positive zero is simple. This follows also
by comparing the orders of the isolated
zeros of the two entire determinants.
The associated graded construction (M29.2)
is what preserves the multiplicities required
by the arithmetic trace and the completed
function.

There are actual number-field examples in which
this distinction matters. For
\(E=\mathbf Q(\zeta_3,\sqrt[3]{2})\), NT-L17
proves unconditionally that \(\zeta_E\) has
infinitely many zeros of multiplicity at
least two. Under \(M\) these lie on the
critical line; infinitely many are
noncentral, and the selfdual symmetry
gives infinitely many positive ordinates.
Thus the ordinary GNS determinant (M29.11)
does not equal the pole-removed completed
\(\zeta_E\) determinant, whereas the
graded determinant (M29.7) does.
This example uses no simplicity conjecture.

\paragraph{The established family and the remaining geometric gate.}
For a selfdual primitive Dirichlet function,
including the separately pole-removed zeta
function, all of the positive conclusions
above use the designated theorem input
\(A_D\), without \(M\). A complex Dirichlet
character and its conjugate are handled
by the finite direct sum of their actual
arithmetic modules and the product of
their completions. Take the associated
graded in each summand, then the positive
frequency subspace of their direct sum.
The product completion is selfdual, its
zero multiplicities are the sums of the
summand multiplicities, and the resulting
Fredholm determinant is its AG-DET
determinant. Each of these assertions
follows from finite direct sums and the
already proved trace/product formulas.
The same construction applies under \(A_D\)
to a specified Hecke factor whose completed
function is a finite product of Dirichlet
completions with the same center, once
that classical factorization is supplied.
For general Hecke factors, and for their
selfdual finite products, the location
input remains conditional on \(M\).

This advances the comparison requested in
AG-Q08 and in the commentary after AG-DET:
the determinant operator is now obtained
from the canonical filtration of the
locally finite part of the actual
adelic summation quotient. Its positivity
still uses the stated zero-location
theorem. It is not an independently
prime-defined positive cohomology or
an identification with the
Gan--Gross--Prasad period determinants
in Alp\"oge's \emph{Remark}, steps~4
and~7 (PDF pp.~3--4). Those period and
integral-structure comparison maps
remain separate.

The precise Masullo source route is
different from an assertion of positivity
of the full Jordan representation.
Section~1.4.6, original PDF p.~30,
explicitly does not invoke Weil
positivity; \S1.4.8--1.4.10,
pp.~31--32, places the main burden
on computing the same geometric
trace intrinsically in arbitrarily
narrow strip completions.
Theorem~12.5.1.1, pp.~519--520,
equations (12.5.1.1)--(12.5.1.4),
asserts the comparison with the
zeros inside the chosen strip for
tests vanishing on \(C_K^1\).
Original pages~30,31,519,520
were rendered and read; the
complete displayed proof of
Theorem~12.5.1.1 was read as
a dependency map, not as an
audit of its cited upstream
kernel theorems. Equation
(10.3.1.17), p.~450, was
also inspected in the original
page image: its triangular
derivative block is consistent
with the local algebra above,
without supplying the missing
uniform orbital comparison.

MT-M28 proves the narrow-strip
trace class under \(M\). It
does not prove that geometric
comparison independently of
zero location. The remaining
gate is the zero-independent
geometric trace identity, or
an independent positive
cohomological/period comparison.
The new graded determinant
comparison closes the
particular module-identification
step and retains the full
numeric multiplicity, while
making the loss of nilpotent
extension data explicit.

\paragraph{Run locators.}
This new comparison is proved in
\path{GRH_NIGHT/complexity},
\path{work/root/graded_adelic_determinant.tex}.
It connects MT-M18/MT-M24/MT-M26
to AG-DET and to the inherited
\path{ALPOGE_BRIDGES} rows
BR-TR-01/02, the SA23
\path{geo-H-lidskii}/\path{geo-H-genus}
rows, and cote~1 scans~79--80
and~111--112 as recorded in
the supplied SA23 intake.
Those archive locators retain
their inherited status.

\subsection*{MT-M30. An ideal-valued splitting of restriction to norm one}
Let \(K\) be a fixed number field, choose the norm splitting
and Haar normalization of MT-M24, and put
\[
 E_K=\mathcal S_{\mathrm E}(C_K),\qquad
 \mathcal I_a=|\cdot|^a\overline{\Theta_K}^{\,E_K},\qquad
 E_K^0=\{f\in E_K:f|_{C_K^1}=0\},
                                                               \tag{M30.1}
\]
where \(a\in\mathbf R\) is fixed. Let
\(\mathcal S(C_K^1)\) denote the Bruhat smooth functions on the
compact group \(C_K^1\). Restriction to norm one is the
continuous map
\[
                 r:E_K\longrightarrow\mathcal S(C_K^1).
 \]
There is a continuous linear, compact-translation-equivariant
map
\[
         B_a:\mathcal S(C_K^1)\longrightarrow\mathcal I_a,
                   \qquad rB_a=\operatorname{id}.             \tag{M30.2}
\]
No GRH input is needed. In particular
\[
 P_a=\operatorname{id}-B_ar
 \quad\hbox{is a continuous projection onto }E_K^0,\qquad
                   P_af-f\in\mathcal I_a.                    \tag{M30.3}
\]
Inclusion consequently induces a topological isomorphism
\[
        E_K^0/(\mathcal I_a\cap E_K^0)
                          \ \simeq\ E_K/\mathcal I_a,         \tag{M30.4}
\]
whose inverse is induced by \(P_a\).
The target in (M30.2) is the \emph{closed} centered
theta ideal. Actual summation-image membership of
an infinite character sum is not asserted.

Write
\[
         \mathcal D^0(C_K)
              =\{f\in\mathcal D(C_K):f|_{C_K^1}=0\}.
 \]
This space is dense in \(E_K^0\).
Therefore two continuous linear functionals on \(E_K\)
that annihilate \(\mathcal I_a\) are equal everywhere
if and only if they agree on \(\mathcal D^0(C_K)\).
The map \(P_a\) need not preserve compact radial
support or positivity.

\begin{proof}
\emph{Uniform normalized elements of the ideal.}
At an elementary compact level, use the family of
characters \(\beta_\nu\), \(\nu\in\mathbf Z^d\times F\),
and the parameter \(H_\nu=2+|\nu|\) from step~1 of
MT-M28. Each extension is trivial on the fixed norm
splitting. For \(\nu\ne0\) use the entire completion
\[
 \Xi_\nu(z)=Q_\nu^{z/2}
        \prod_{j=1}^{n_K}\Gamma_{\mathbf R}(z+\mu_{\nu,j})
                         L_\nu(z).
 \]
For \(\nu=0\) use the fixed pole-removed completion
of \(\zeta_K\), for example its standard completion
multiplied by \(z(z-1)\). These choices have precisely
the nontrivial completed zero divisors used by MT-M24.
All satisfy \(\Xi_\nu(2)\ne0\), by the Euler product
and the finite nonzero gamma factors there.
Define
\[
          F_\nu(z)=
              \frac{\Xi_\nu(z)}{\Xi_\nu(2)}\,e^{(z-2)^2}.
                                                               \tag{M30.5}
\]
It is entire, vanishes to the full order at every
completed zero, and satisfies \(F_\nu(2)=1\).
Let \(\psi_\nu\) be its inverse Mellin transform.
For each fixed \(\nu\), the fixed-strip bounds and
Gaussian decay in MT-M18 show
\(\psi_\nu\in\mathcal E\).

We need, and now prove, bounds uniform at this
compact level:
\[
           p_{N,j}(\psi_\nu)
                        \le C_{N,j}H_\nu^{B_{N,j}}.            \tag{M30.6}
\]
The finite conductor is uniformly bounded and
\(\operatorname{Re}\mu_{\nu,j}\) is a nonnegative
half-integer with \(|\mu_{\nu,j}|\ll H_\nu\).
The Euler lower bound at \(2\) and the uniform
finite-\(L\) polynomial upper bound are precisely
(M28.10) and (M28.12). Their proofs were
unconditional. On either line
\(\operatorname{Re}z=\pm(N+9/4)\), the same
gamma-ratio proof as in (M28.20), now with the
single fixed denominator point \(w=2\), gives
\[
 \left|\frac{\Xi_\nu(z)}{\Xi_\nu(2)}\right|
       \le C_N(H_\nu+|\operatorname{Im}z|)^{B_N}
                              e^{C_N|\operatorname{Im}z|}.
                                                               \tag{M30.7}
\]
In detail, shift numerator gamma factors right
by a fixed even integer. Every recurrence
denominator has real part a quarter-integer
of absolute value at least \(1/4\). The
remaining gamma ratio is in a fixed right
half-plane. The bounds
\(|\operatorname{Im}\psi(z)|\le C_c\) and
\(|\psi(z)|\le C_c+2\log(2+|z|)\) there
give (M30.7) by integrating the logarithmic
derivative in imaginary and real directions.
This uses a canceled ratio throughout, so
no factorial-size completion normalization
remains. The single principal character
has its fixed polynomial pole-removal
factor and is included by enlarging
the constants once.

Multiplication by the Gaussian in (M30.5)
turns (M30.7) into
\[
 |F_\nu(\pm(N+9/4)+it)|
                 \le C_NH_\nu^{B_N}e^{-c t^2}.
 \]
Inverting on these lines, differentiating
under the absolutely convergent integrals,
and using the positive line for \(x\ge0\)
and negative line for \(x\le0\), proves
(M30.6), just as in (M28.22)--(M28.23).
This part of MT-M28 used only Euler
convergence, the classical functional
equation and per-function finite order.
Neither its GRH contour selection nor
its narrow-strip minimum modulus was
used here.

Set
\[
                 \varphi_\nu(x)=e^{ax}\psi_\nu(x).
 \]
Multiplication by \(e^{ax}\) is a continuous
automorphism of \(\mathcal E\), so it
preserves the polynomial-in-\(H_\nu\)
form of (M30.6), with changed indices
and constants. Its normalization is
explicitly
\[
                   \mathcal M\varphi_\nu(2-a)=1.              \tag{M30.8}
\]
The full-zero-jet identification in
MT-M24 shows that this is an element
of the centered \(\nu\)-sector
theta ideal \(\mathcal I_{a,\nu}\).
One can also see the ideal membership
before taking its closure: the
Gaussian factor in (M30.5) is itself
the Mellin transform of an element
of \(\mathcal E\), so \(\psi_\nu\)
lies in the principal convolution
ideal generated by the inverse
transform of \(\Xi_\nu\).

\emph{A polynomial lower bound for the normalizing norm.}
It is important not to divide by an
uncontrolled small number. Put \(t_0=2-a\)
and choose an integer \(L>|t_0|+1\).
By (M30.6),
\[
                  p_{L,0}(\varphi_\nu)\le C H_\nu^B.
 \]
For \(R\ge1\), the tail of the integral
in (M30.8) satisfies
\[
 \left|\int_{|x|>R}e^{t_0x}\varphi_\nu(x)\,dx\right|
       \le C' H_\nu^B e^{-(L-|t_0|)R}.
 \]
Choose \(R=C_0+C_1\log H_\nu\), with
fixed constants large enough that
this bound is at most \(1/2\).
The integral over \([-R,R]\) has
absolute value at least \(1/2\).
Cauchy--Schwarz gives
\[
 \frac12\le
       (2R)^{1/2}e^{|t_0|R}\|\varphi_\nu\|_2 .
 \]
Since \(R\ll1+\log H_\nu\), it follows
that, for fixed \(c,C>0\),
\[
                   \|\varphi_\nu\|_2^2\ge cH_\nu^{-C}.        \tag{M30.9}
\]
This argument also covers \(t_0=0\).

For a radial function write
\(\varphi^*(x)=\overline{\varphi(-x)}\), and define
\[
              b_\nu
                =\frac{\varphi_\nu*\varphi_\nu^*}
                             {\|\varphi_\nu\|_2^2}.           \tag{M30.10}
\]
Its first convolution factor is in
\(\mathcal I_{a,\nu}\), and the second
is in \(\mathcal E\). Thus
\[
          b_\nu\in\mathcal I_{a,\nu},
          \qquad b_\nu(0)=1,\qquad
          p_{N,j}(b_\nu)\le C_{N,j}H_\nu^{C_{N,j}}.            \tag{M30.11}
\]
The value at zero is exactly the
squared \(L^2\) norm in (M30.10).
For the seminorm bound, use
\[
 p_{N,j}(f*g)
   \le p_{N,j}(f)\int_{\mathbf R}e^{N|y|}|g(y)|\,dy
   \le 2p_{N,j}(f)p_{N+1,0}(g)
 \]
and then (M30.6), (M30.9).
No stability of \(\mathcal I_a\)
under the involution for arbitrary
\(a\) is needed for this membership:
only its first factor and its
convolution-ideal property are used.

\emph{Summing over the compact character family.}
For \(g\in\mathcal S(C_K^1)\) at the
chosen compact level, write
\[
       g(c)=\sum_{\nu\in\Lambda}g_\nu\beta_\nu(c)^{-1},
       \qquad
       g_\nu=\int_{C_K^1}g(c)\beta_\nu(c)\,dc .
 \]
The coefficients decrease faster
than every power of \(H_\nu\).
Define, in the norm-splitting coordinates,
\[
 (B_ag)(c\,\tau(e^x))
       =\sum_{\nu\in\Lambda}
                   g_\nu\beta_\nu(c)^{-1}b_\nu(x).
                                                               \tag{M30.12}
\]
Each compact derivative costs a
fixed power of \(H_\nu\), and
(M30.11) bounds every weighted
radial derivative by a fixed power.
The rapid Fourier coefficients
therefore make (M30.12) absolutely
convergent in every seminorm of
the elementary \(E_K\) level.
The required source coefficient
bound uses finitely many compact
derivatives, so \(B_a\) is
continuous at that level.

The functions \(b_\nu\) were
defined from the character,
its fixed norm extension and
(M30.5), independently of any
choice of generators used to
prove the bounds. Thus the
definitions agree when levels
are enlarged. The elementary
inductive-limit definition of
the Bruhat smooth topology
now gives the asserted global
continuous map.
Every finite partial sum lies
in \(\mathcal I_a\), and this
ideal is closed in \(E_K\);
hence the whole sum lies in it.
At \(x=0\), (M30.11) gives
\(rB_ag=g\), with convergence
in the compact smooth topology.
The character formula also
shows that \(B_a\) commutes
with compact translations.
This proves (M30.2).

Equations (M30.3) follow from
\(rB_a=\operatorname{id}\).
If \(f\in E_K^0\), then \(P_af=f\).
If \(f\in\mathcal I_a\), then
\(P_af\in\mathcal I_a\cap E_K^0\).
The maps induced by inclusion
and \(P_a\) on the two quotient
spaces are therefore inverse.
Both are continuous by the
defining quotient topologies,
which proves (M30.4).

\emph{Compact test density and extension of equalities.}
For \(f\in E_K^0\), choose smooth
radial cutoffs \(\kappa_R(x)\)
equal to one on \([-R,R]\)
and zero outside \([-2R,2R]\).
The functions \(\kappa_R f\)
are in \(\mathcal D^0(C_K)\):
their compact-variable level
is unchanged, their norm
support is compact, and
\(\kappa_R(0)f(c)=0\).
They converge to \(f\) in
every weighted derivative
seminorm. Indeed the
differences and the derivatives
of the cutoffs are supported
where \(|x|\ge R\); a bound
at one larger exponential
weight gives an additional
factor tending to zero.
This proves density.

Suppose \(T_1,T_2\) are
continuous on \(E_K\), vanish
on \(\mathcal I_a\), and
agree on \(\mathcal D^0(C_K)\).
By density they agree on
\(E_K^0\). For every \(f\in E_K\),
\[
            T_i(f)=T_i(P_af),
 \]
since \(f-P_af\in\mathcal I_a\).
They are therefore equal
on all of \(E_K\). The
reverse implication is immediate.
\end{proof}

\paragraph{An exact restricted-trace criterion.}
The splitting gives a useful precise
test of the residual source claim.
For each spectral pair \((\chi,\rho)\)
choose real weights
\[
                   0\le w_{\chi,\rho}\le m_{\chi,\rho}.
 \]
Define
\[
 \begin{split}
 \tau_{\mathrm{full}}(f)
    &=\sum_{\chi,\rho}m_{\chi,\rho}
              \widehat f(\beta_\chi|\cdot|^{\,\rho-a}),\\
 \tau_w(f)
    &=\sum_{\chi,\rho}w_{\chi,\rho}
              \widehat f(\beta_\chi|\cdot|^{\,\rho-a}).
 \end{split}                                                    \tag{M30.13}
\]
These are continuous on \(E_K\) and
annihilate \(\mathcal I_a\),
unconditionally. To verify convergence,
at each compact level the unconditional
Jensen estimate (M28.13) counts
the nontrivial zeros in each unit
height window by
\(O(\log(H_\nu+|n|))\).
Their real parts are in the fixed
strip \([0,1]\). Integration by
parts in the compact variables
and in the norm variable bounds
their Mellin values by arbitrary
negative powers of \(H_\nu+|n|\),
uniformly on bounded sets at that
level. In the norm variable the
integrals of derivatives are
controlled by finitely many
exponential-Schwartz seminorms;
compact radial support is not
needed. The same summation as
in MT-M28 proves absolute
convergence and continuity.
Vanishing on \(\mathcal I_a\)
is the full zero-jet kernel
identity from MT-M24/MT-M27.

Consequently
\[
 \left.\tau_{\mathrm{full}}\right|_{\mathcal D^0(C_K)}
       =\left.\tau_w\right|_{\mathcal D^0(C_K)}
       \quad\Longleftrightarrow\quad
             w_{\chi,\rho}=m_{\chi,\rho}
                     \text{ for every }(\chi,\rho).           \tag{M30.14}
\]
For the nontrivial implication,
the splitting extends the equality
to \(E_K\). Fix one spectral
pair and use the finite-jet
interpolation of MT-M18 in
that single character sector
to get transform value one
at that zero and zero at
all other zeros, with every
unwanted higher jet zero.
Evaluating the difference of
the two functionals gives
\(m_{\chi,\rho}-w_{\chi,\rho}\).
It must vanish. The reverse
implication follows directly
from the sums.
Equivalently, any omitted
zero or lowered multiplicity
has a witness already in
\(\mathcal D^0(C_K)\):
apply \(P_a\) to this
finite-spectral interpolant,
then approximate by the
compact radial cutoffs above.
Continuity keeps the nonzero
difference for all sufficiently
large cutoffs.

In particular, at \(a=1/2\)
choose \(w_{\chi,\rho}=m_{\chi,\rho}\)
inside a specified open
centered strip and zero outside.
Equality of the full and
strip-restricted functionals
on tests vanishing at norm
one is equivalent to the
absence of every omitted
zero. Requiring this for
all positive strip widths
is exactly the Hecke
zero-location statement.
The test restriction does
not weaken that criterion.
This equivalence supplies
no proof of either side.

\paragraph{What this changes in the source comparison.}
Masullo \S1.4.6, original PDF p.~30,
restricts the proposed geometric
trace calculation to
\(\varphi(C_K^1)=0\);
\S1.7.2, pp.~34--35,
notes that a literal convolution
square \(h*h^*\) satisfying
this condition is zero.
That last assertion is correct:
at the identity
\[
                     (h*h^*)(1)=\|h\|_2^2.
 \]
The new splitting shows that
it is not an obstruction
at the corrected arithmetic
quotient level. The class
of \(h*h^*\) has the
representative \(P_{1/2}(h*h^*)\)
vanishing on \(C_K^1\), and
every continuous quotient
trace takes the same value
on these two representatives.
The new representative need
not itself be a convolution
square or have any pointwise
positivity property.

Let \(W_{\mathrm{raw},1/2}\)
be the classical raw Weil
distribution in the centering
convention of MT-M24. Its
spectral version is
\[
 W_{\mathrm{spec}}(f)
   =\widehat f(|\cdot|^{1/2})
      +\widehat f(|\cdot|^{-1/2})
      -W_{\mathrm{raw},1/2}(f)
   =\tau_{\mathrm{full}}(f).
\]
An exact classical trace
locator is Meyer,
\emph{On a representation of
the idele class group related
to primes and zeros of
\(L\)-functions},
Theorem~5.8, author pp.~43--44,
with the spectral identification
of Theorem~5.11, p.~47.
The complete supplied proof
of Theorem~5.8 and its
normalization were read in
this run; its upstream
trace-theoretic lemmas
remain cited classical inputs.
Here its continuation to
\(E_K\) can be defined by
the absolutely convergent
full spectral sum (M30.13),
so no unproved geometric
kernel continuity is inserted.

Theorem~12.5.1.1 in the
Masullo draft, original
pp.~519--520, asserts that
the same geometric expression
equals the trace with only
the zeros inside the chosen
strip, on the restricted
compact tests. These original
pages and \S1.4.6 were
inspected as page images
for MT-M29. By (M30.14),
the remaining gate is
exactly that geometric
identity independently
of zero location, rather
than admissibility of
positivity tests or a
missing extension from
the restricted test space.
The uniform ideal-valued
splitting is unconditional;
the value of the full
functional on convolution
squares is nonnegative
under \(A_D\) in the
Dirichlet family and
under \(M\) in the
general Hecke family,
as in MT-M19/MT-M28.
No positivity of the
corrected representatives
is inferred.

\paragraph{Run and framework locators.}
This is a new proof in
\path{GRH_NIGHT/complexity},
\path{work/root/idele_restriction_splitting.tex}.
It advances the arithmetic
quotient comparison from
MT-M24/MT-M27/MT-M28 and
the determinant comparison
MT-M29. Its geometric
framework is the trace
functional on a quotient
by theta relations, beside
the inherited SA23
\path{geo-H-lidskii}
and \path{geo-H-genus}
rows, and cote~1
scans~79--80,111--112.
It supplies neither an
integral cycle comparison
nor an independent
Gan--Gross--Prasad
period positivity theorem.

\subsection*{MT-S504. The exact relative central condition in the $S_5$ transfer}

Retain the $S_5$ extension and local hypotheses of MT-S501:
complex conjugation and the unramified Frobenius at $5$ both
have cycle type $(12)(34)$. Write $E=L^{A_5}$, let $F$
be a degree-six subfield, and let $H=EF$. Let $\eta$ be the
quadratic character of $E/\mathbb Q$. Use Calegari's labeling
of the five-dimensional character $\rho_5$, whose values on
the classes $1,2A,2B,3A,4A,5A,6A$ are
$(5,1,1,-1,-1,0,1)$. Then the all-prime Artin factorization is
\[
       \zeta_H(s)=\zeta_F(s)L(s,\eta)L(s,\rho_5).
 \tag{S504.1}
\]
Under the conditional input $M_H$, the strictly more focused
sufficient condition
\[
 \operatorname{ord}_{1/2}\zeta_H
       \le\operatorname{ord}_{1/2}\zeta_F
                     +\operatorname{ord}_{1/2}L(s,\eta)
 \tag{S504.2}
\]
implies the same conclusion as MT-S502: every nontrivial
irreducible Artin factor is entire and satisfies GRH, and
the factors are automorphic by Calegari's argument.
In fact equality then holds in (S504.2). The hypothesis
allows zeros in the two other factors; it does not require
$\zeta_H(1/2)\ne0$.

\begin{proof}
Put $m_5=\operatorname{ord}_{1/2}L(s,\rho_5)$, interpreted
as the meromorphic order, so negative values are poles.
Equation (S504.1) shows that (S504.2) is exactly $m_5\le0$.
The factorization and the following precise classical inputs
are Calegari, \emph{The Artin conjecture for some
$S_5$-extensions}, \path{arXiv:1112.1152v1}: the comment
after Theorem~1.2, p.~2; Lemmas~4.8--4.9, pp.~10--11;
and the final proof of Theorem~1.2, p.~11. Their relevant
proofs have been read in full. The automorphic construction
in preceding sections remains a cited classical input.

That construction supplies a cuspidal automorphic
$\varpi$ on $\mathrm{GL}_5/\mathbb Q$, an exceptional set
$\Omega$ of primes, and a meromorphic comparison
\[
 \Lambda(s,\varpi)=\Lambda(s,\rho_5)\mu(s),\qquad
 \mu(s)=\mu_B(s)
       \prod_{p\in\Omega}
             \left(\frac{1+p^{-s}}{1-p^{-s}}\right)^2.
 \tag{S504.3}
\]
Here $\mu_B$ is holomorphic and nonzero on $(0,1)$.
Lemma~4.8 says that if $\Omega$ is infinite, $\mu$ has
a pole at a real point in $(0,1)$. Lemma~4.9 says that
finite $\Omega$ forces exact local equality everywhere,
including ramified and archimedean factors. The completed
automorphic numerator in (S504.3) is holomorphic on this
interval.

Brauer induction and $M_H$ imply that $L(s,\rho_5)$ is
holomorphic and nonzero at every real point of $(0,1)$
other than $1/2$: every Hecke factor in its signed
factorization has that property. Its archimedean factors
are finite and nonzero on the interval. Thus a real pole
of $\mu$ could only occur at $1/2$ and would require
$m_5>0$. Our hypothesis $m_5\le0$ rules that out;
a pole of the denominator would instead give a zero
of the comparison. Lemma~4.8 makes $\Omega$ finite,
and Lemma~4.9 proves $\Lambda(s,\varpi)=\Lambda(s,\rho_5)$.
The source's final argument also gives the exact match
for the quadratic twist $\rho_5\otimes\eta$.

The other irreducible dimensions $4$ and $6$ are already
automorphic by Theorem~1.1, as recorded in MT-S501;
the linear factors are classical Dirichlet functions.
This proves all the automorphy and holomorphy assertions.
Now holomorphy makes $m_5\ge0$, so $m_5=0$ and equality
holds in (S504.2). Finally signed Brauer factorization
and $M_H$ locate the nontrivial zeros of every entire
factor on the critical line, exactly as in MT-ART07.
\end{proof}

All $S_5$ irreducibles are realizable over $\mathbb R$;
Fr\"ohlich--Queyrut's real-representation theorem gives
root number $+1$. The meromorphic central-order parity
calculation of MT-ART11 therefore makes $m_5$ even.
Consequently failure of the sufficient condition leaves
only the gate
\[
 \operatorname{ord}_{1/2}\zeta_H
       -\operatorname{ord}_{1/2}\zeta_F
       -\operatorname{ord}_{1/2}L(s,\eta)
                       =m_5\ge2.
 \tag{S504.4}
\]
This is a necessary condition for this route to remain
unsettled, not a claim that a positive central order
prevents automorphy. It removes base-factor central zeros
from the test exactly; no independence or noncoincidence
of zeros is assumed.

Geometrically the virtual permutation identity in
(S504.1) isolates a single rank-five Artin motive from
the arithmetic finite covers. It sharpens the
MT-S502 one-value gate in the cote~1/SGA5 trace-divisor
programme, and the $S_5$ portion of Q-ARTIN-POS.
It leaves the actual relative central order as an
arithmetic input; neither designated GRH theorem
evaluates it.

\subsection*{MT-S505. Real-axis Artin holomorphy in the specified $S_5$ family}

Retain the extension and local hypotheses of MT-S501 and the
notation of MT-S504. Conditional on $M_H$, every irreducible
Artin factor of this extension is holomorphic at every real
point of $(0,1)$, including the centre. In particular
\[
 \operatorname{ord}_{1/2}\zeta_H
  \ge \operatorname{ord}_{1/2}\zeta_F
                  +\operatorname{ord}_{1/2}L(s,\eta).
 \tag{S505.1}
\]
The difference in (S505.1) is a nonnegative even integer.
If it is zero, every irreducible factor is automorphic
and satisfies Artin holomorphy and GRH, by MT-S504.
Thus no additional central nonvanishing hypothesis is
needed for the real-axis holomorphy assertion itself.

\begin{proof}
Write $m_5=\operatorname{ord}_{1/2}L(s,\rho_5)$,
allowing a negative order for a meromorphic pole.
If $m_5<0$, the sufficient condition $m_5\le0$
of MT-S504 holds. That theorem proves automorphy
and forces $m_5=0$, a contradiction. Hence $m_5\ge0$.
Meromorphicity then means that $L(s,\rho_5)$
is holomorphic at $s=1/2$. This argument does
not assume its central value is nonzero.

For completeness the other rank-five factor
is also holomorphic at the centre for an
unconditional elementary reason. Put
$J=\operatorname{Gal}(L/F)$, of order $20$.
Its Sylow $5$-subgroup $P$ is normal, since
its number divides $4$ and is $1$ modulo $5$.
The normalizer of $P$ in $S_5$ has order $20$;
thus $J$ is that normalizer. In the standard
affine action on $\mathbb F_5$, its elements
have cycle types
\[
       1\text{ once},\quad 5A\text{ four times},\quad
       2B\text{ five times},\quad 4A\text{ ten times}.
\]
Using Calegari's character table, arXiv:1112.1152v1,
p.~3, the average over these elements is $1$
for the trivial character and for $\rho_5\otimes\eta$,
and $0$ for $\eta,\rho_4,\rho_4\otimes\eta,\rho_5,\rho_6$.
For example the two rank-five averages are
\[
 \frac{5+4\cdot0+5\cdot1+10\cdot(-1)}{20}=0,
 \qquad
 \frac{5+4\cdot0+5\cdot1+10\cdot1}{20}=1.
\]
Frobenius reciprocity therefore proves the
permutation identity and its all-prime
Artin factorization
\[
 \operatorname{Ind}_{J}^{S_5}\mathbf1
          =\mathbf1\oplus(\rho_5\otimes\eta),\qquad
 \zeta_F(s)=\zeta_{\mathbb Q}(s)L(s,\rho_5\otimes\eta).
 \tag{S505.2}
\]
For $0<s<1$, the alternating series
\[
 A(s)=\sum_{n\ge1}(-1)^{n-1}n^{-s}
     =\sum_{j\ge1}\bigl((2j-1)^{-s}-(2j)^{-s}\bigr)
\]
converges and is strictly positive. Analytic continuation
of the elementary identity from $s>1$ gives
$A(s)=(1-2^{1-s})\zeta_{\mathbb Q}(s)$ on this interval.
The first factor is negative, so
$\zeta_{\mathbb Q}(s)<0$ throughout $(0,1)$.
Thus (S505.2) expresses the twisted rank-five
factor as a holomorphic quotient on the whole
interval, without any GRH assumption.

The factors of dimensions $4$ and $6$ are
entire by the cited Calegari Theorem~1.1,
and the linear factors are holomorphic
on $(0,1)$. Finally for $\rho_5$ at a
real point other than $1/2$, signed Brauer
induction and $M_H$ give a product and
quotient of nonzero holomorphic Hecke values,
so it is holomorphic there too.
This proves the real-axis assertion for
every irreducible factor.

Taking central orders in (S504.1) gives
(S505.1) with difference $m_5$. The real
realizability/root-number-$+1$ and
meromorphic functional-equation parity
argument already checked in MT-S504
makes $m_5$ even. The equality case is
exactly the sufficient condition of
MT-S504. All conclusions follow.
\end{proof}

This settles the real-axis portion of the
Artin pole gate for this precise $S_5$
family. Possible critical-line poles away
from the real axis are not ruled out in
the remaining $m_5\ge2$ case. Geometrically,
(S505.2) is the degree-six permutation
Artin motive of the finite cover, while
(S505.1) is an effectivity statement for
the central coefficient of the relative
rank-five divisor. It concerns that
coefficient, not the effectivity of the
whole Artin divisor required by
Q-ARTIN-POS and the cote~1/SGA5 trace programme.

\subsection*{MT-M31. Local Euler differences: actual range, Hilbert closure,
and the Schanuel dependency}

The distinction in this result is between the range of a bounded
operator and the closure of that range.  It occurs before any
GRH or positivity question.

\paragraph{The precise source assertion.}
Masullo, Proposition 11.2.0.1(2), original PDF pp.~484--485,
asserts that the extension of $\mathrm{PV}_{S_0}$ to the
zero-line Hilbert completion has closed range when
$s\geq 2|S_0|$.  Equation (11.2.0.3), p.~484, identifies
the closure of the smooth image with that actual operator range.
The proof on p.~484 gives, for $S_0=\{v\}$, the exact formula
\[
  P_v\Phi_F(f)=\Phi_F(P_vf),\qquad
  P_vf(x)=f(x)-f(\pi_v^{-1}x).
  \tag{M31.1}
\]
There is no additional norm-character multiplier in this formula.
On p.~485 the proof identifies the range with an intersection
of kernels of continuous coefficient functionals.  The argument
below proves that this identifies the \emph{closure}, and that
the actual range is not closed even for one prime over $\mathbf Q$.
The original images of both pages have been read, including
the multiplier and the range/closure equation.

\paragraph{General analytic formulation.}
For $s\geq0$ put
\[
 H_s=L^2(\mathbf R,w_s(x)\,dx),\qquad
 w_s(x)=(1+x^2)^{s/2},\qquad
 (U_hf)(x)=f(x-h).
 \tag{M31.2}
\]
This is precisely the radial form of the source norm
(9.1.2.3), original p.~399, with the $a=0$ completion
of (10.1.1.3), p.~434; see MT-M24--M25 for the actual
theta-image identification.  Fix a nonempty finite list
of $h_v>0$ and $a_v\in\mathbf C$ with $|a_v|=1$, and set
\[
 P=\prod_v(1-a_vU_{h_v}),\qquad
 D(z)=\prod_v(1-a_ve^{h_vz}).
 \tag{M31.3}
\]
Repeated factors are allowed.  Let $\mathcal T\subset\mathbf R$
be the locally finite set of $t$ with $D(it)=0$, and let
$m_t$ be its multiplicity.  Define
\[
 q_s=\max\{0,\lceil(s-1)/2\rceil\},\qquad
 J_{t,j}(f)=\int_{\mathbf R}f(x)x^je^{itx}\,dx.
 \tag{M31.4}
\]
The following assertions hold without any conjectural input:
\begin{enumerate}
\item $P$ is a bounded injective operator on $H_s$, and its
actual range is not closed, for every $s\geq0$.
\item Its exact closed range replacement is
\[
 \overline{P H_s}^{\,H_s}
 =\bigcap_{t\in\mathcal T}\ \bigcap_{0\leq j<\min(m_t,q_s)}
            \ker J_{t,j}.
 \tag{M31.5}
\]
Every displayed functional is continuous.  Thus for $0\leq s\leq1$
the actual range is dense and proper; for $s>1$ its closure
is proper and the range is a proper dense subspace of that closure.
\item If $\mathcal A$ is any dense linear subspace of $H_s$, then
\[
 \overline{P\mathcal A}^{\,H_s}
       =\overline{PH_s}^{\,H_s}.
 \tag{M31.6}
\]
In particular the corrected quotient is the Hilbert quotient by
(M31.5).  It is not the non-Hausdorff topological quotient
$H_s/PH_s$.
\end{enumerate}

\paragraph{Boundedness and injectivity.}
The inequality
\[
 1+(x+y)^2\leq2(1+x^2)(1+y^2)
 \tag{M31.7}
\]
shows that $\|U_h\|_{H_s\to H_s}\leq
2^{s/4}(1+h^2)^{s/4}$.  It gives the analogous estimate
on $H_{-s}=L^2(w_s^{-1}dx)$ as well.  Each factor and
their product are bounded.  Since $w_s\geq1$, $H_s$
embeds in ordinary $L^2(\mathbf R)$.  The ordinary Fourier
transform of $Pf$ is the transform of $f$ multiplied by
\[
            \prod_v(1-a_ve^{-ih_v\xi}).
\]
Its real zeros form a locally finite set, of Lebesgue measure
zero.  Consequently $Pf=0$ implies $f=0$ in $L^2$,
hence in $H_s$.

\paragraph{A quantitative failure of closed range.}
Choose one root $t_0$ of one of the factors of $D(it)$,
and choose nonzero $\psi\in C_c^\infty((1,2))$.
For $R\geq1$ put
\[
             f_R(x)=e^{-it_0x}\psi(x/R).
 \tag{M31.8}
\]
Changing variables gives constants $c,C>0$, independent of $R$,
such that
\[
      cR^{s+1}\leq\|f_R\|_{H_s}^2\leq CR^{s+1}.
 \tag{M31.9}
\]
For a factor with $a_ve^{it_0h_v}=1$,
\[
 (1-a_vU_{h_v})f_R(x)
   =e^{-it_0x}\{\psi(x/R)-\psi((x-h_v)/R)\}.
 \tag{M31.10}
\]
The mean-value bound, with support contained in an interval
of length $O(R)$ and $|x|=O(R)$, gives
\[
 \|(1-a_vU_{h_v})f_R\|_{H_s}
      \leq C_vR^{(s-1)/2}.
\]
The remaining finitely many factors are bounded independently
of $R$.  Since they commute,
\[
       \frac{\|Pf_R\|_{H_s}}{\|f_R\|_{H_s}}=O(R^{-1}).
 \tag{M31.11}
\]
An injective bounded operator between Hilbert spaces with
closed range has bounded inverse onto that range, by the
bounded-inverse theorem, and hence is bounded below.  Equation
(M31.11) contradicts this.  It proves the first assertion.

\paragraph{The exact closure.}
Use the continuous bilinear duality
\[
       H_s\times H_{-s}\longrightarrow\mathbf C,
       \qquad(f,g)\longmapsto\int f(x)g(x)\,dx.
 \tag{M31.12}
\]
Complex conjugation converts it to the usual Hilbert duality;
the bilinear version keeps the Mellin convention explicit.
The annihilator $\mathcal N$ of $PH_s$ consists precisely of
the $g\in H_{-s}$ satisfying
\[
          P^{\mathrm t}g=
             \prod_v(1-a_vU_{-h_v})g=0.
 \tag{M31.13}
\]
Each $g\in H_{-s}$ is a tempered distribution: Cauchy--Schwarz
bounds its pairing with a Schwartz test by
$\|g\|_{H_{-s}}\|\phi\|_{H_s}$.  With Fourier convention
$\widehat g(\xi)=\int g(x)e^{-i\xi x}dx$, (M31.13) is
\[
                         D(i\xi)\widehat g(\xi)=0.
 \tag{M31.14}
\]
It follows that $\widehat g$ is supported on $\mathcal T$.
Near an isolated $t\in\mathcal T$, a distribution killed
by the smooth function $D(i\xi)$ is a linear combination
of $\delta_t,\ldots,\delta_t^{(m_t-1)}$.  To verify the
order bound, divide by the nonvanishing smooth factor in
$D(i\xi)=(\xi-t)^{m_t}b(\xi)$ and use
$(\xi-t)\delta_t^{(j)}=-j\delta_t^{(j-1)}$.
The classical point-support distribution theorem,
\emph{Expository Notes on Distribution Theory}, AMSC/MATH~673,
Fall~2018, Theorem~2.1, p.~4 (University of Maryland),
then gives exactly this finite list.

Here localization is bounded in the required Hilbert norm.
For any Schwartz kernel $k$, (M31.7) and Minkowski give
\[
 \|k*g\|_{H_{-s}}
 \leq C_s\left(\int |k(y)|(1+|y|)^{s/2}dy\right)
                 \|g\|_{H_{-s}}.
 \tag{M31.15}
\]
Multiplying $\widehat g$ by a compactly supported smooth
function supported near one $t$ is therefore bounded on
$H_{-s}$.  The resulting inverse transform is a polynomial
times $e^{itx}$.  A nonzero polynomial of degree $j$
times this exponential belongs to $H_{-s}$ exactly when
\[
                              2j<s-1.
 \tag{M31.16}
\]
Indeed its squared absolute value is comparable to $|x|^{2j}$
for all sufficiently large $|x|$; the weighted integral
then has tail comparable to $\int_1^\infty x^{2j-s}dx$.
In particular equality in (M31.16) is not allowed.
This argument isolates each frequency first, so cancellation
among distinct frequencies cannot evade the cutoff.

For completeness the finite exponential-polynomial sums are
dense in $\mathcal N$ in its actual $H_{-s}$ norm.
Choose $\eta\in C_c^\infty(\mathbf R)$ equal to one near zero,
and put $\eta_R(\xi)=\eta(\xi/R)$, $R\geq1$.
Its inverse Fourier kernel has the form $k_R(x)=Rk(Rx)$,
with integral one.  Equation (M31.15) gives uniformly bounded
operators because
\[
 \int |k_R(y)|(1+|y|)^{s/2}dy
   =\int |k(u)|(1+|u|/R)^{s/2}du
   \leq\int |k(u)|(1+|u|)^{s/2}du.
\]
Translations are strongly continuous on $H_{-s}$, first on
$C_c^\infty$ and then by density and (M31.7).
Writing $k_R*g-g$ as the integral of
$k(u)(U_{u/R}g-g)$ proves convergence to zero in $H_{-s}$
by dominated convergence, using the same polynomial bound.
Each $\eta_R\widehat g$ has finite support in $\mathcal T$,
so its inverse transform is a finite sum of the terms
$x^je^{itx}$ with $j<m_t$ and (M31.16).
Conversely all these terms satisfy (M31.13), and their
closed span therefore equals $\mathcal N$.

Taking the annihilator once more proves (M31.5).
Cauchy--Schwarz proves continuity of every $J_{t,j}$ in that
formula using (M31.16).  For $s>1$, at least one root exists
and $J_{t,0}$ is a nonzero functional, so the closure is proper.
For $s\leq1$ no term survives and the closure is all of $H_s$.
Nonclosedness gives all remaining strictness assertions.
Finally boundedness of $P$ and density of $\mathcal A$ prove
(M31.6) directly by approximating each vector before applying $P$.

\paragraph{Application to the literal adelic assertion.}
It is enough to take $F=\mathbf Q$, $S_0=\{p\}$, and
the trivial character of the compact norm-one group.
Choose a norm splitting $C_{\mathbf Q}\simeq
C_{\mathbf Q}^1\times\mathbf R$ with radial coordinate
$x=\log|g|$, and normalize the compact Haar measure to one.
On this reducing sector (M31.1) is
$f(x)\mapsto f(x)-f(x+\log p)$.
Replacing $x$ by $-x$ gives exactly $1-U_{\log p}$ on
(M31.2).  Reflection is isometric because $w_s$ is even.

The ambient operator on the full adelic Hilbert space is
injective as well.  If $\Phi(g)=\Phi(\pi_p^{-1}g)$
almost everywhere, decompose $C_{\mathbf Q}$ into the
disjoint norm shells
$\{k\log p\leq\log|g|<(k+1)\log p\}$.
Translation by $\pi_p^{-1}$ maps successive shells onto
one another and preserves Haar measure.  The unweighted
$L^2$ mass is consequently equal on all shells.  Since
$H_s\subset L^2$ for $s\geq0$, it must be zero.
Thus the restriction to any closed invariant Hilbert subspace
is also injective.  If its range were closed, that restriction
would be bounded below on every sector it contains.

The actual zero-line completion of the trivial-character
theta image contains the whole radial $H_s$ sector.
This follows from MT-M24 and the exact MT-M25 kernel formula
at $a=0$, since the pole-removed completed zeta function has
no nontrivial zero on $\Re z=0$.  The latter is the classical
zero-free boundary theorem and functional equation, with no
RH premise.  Thus the normalized vectors from (M31.8)
belong to the actual completion, and (M31.11) contradicts
Proposition 11.2.0.1(2), including its stated range $s\geq2$.
This proves a counterexample in the specified space, not only
in an unrelated radial model.

By (M31.6), the left side of source (11.2.0.3) is precisely
the closure in (M31.5); its right side, read as the actual
extended-operator range, is a proper dense subspace of it.
Replacing the right side by its closure is the necessary
correction.  No conclusion about the truth of $M_H$ follows
from this counterexample.

\paragraph{What the local-pole input actually assumes.}
The following is a targeted dependency trace, not an audit
of the entire draft.
\begin{enumerate}
\item Theorem 7.1.3.2(1), original pp.~324--325, says that
two unramified finite places of distinct residue
characteristics have at most one common Euler pole.
If $s,s'$ were distinct common poles, then
$(s-s')\log q_v$ and $(s-s')\log q_w$ would both be
nonzero integral multiples of $2\pi i$, forcing
$\log q_v/\log q_w\in\mathbf Q$, impossible for powers
of distinct rational primes.  This part is unconditional.
\item Part (3b) of that theorem explicitly assumes Schanuel
for its sharper coincidence exclusion for a character
which is an algebraic character times a norm power.
The proof on p.~325 uses the four-exponentials consequence
of Schanuel.  Its weaker three-residue-characteristic
exclusion uses the unconditional six-exponentials theorem.
The source does not present the four-exponentials assertion
as an unconditional input in that argument.
\item For one rational prime $p$, the pole multiplicity of
$\prod_{v\mid p}L_v(s,\chi_v)$ is at most the number
of factors, hence at most $[F:\mathbf Q]$.
Each unramified Euler factor has simple poles, and ramified
factors contribute none.  This elementary bound does not
use Schanuel or coincidence control between different primes.
Definition 7.3.1.4, p.~343, and Theorem 7.3.1.6,
pp.~344--350, organize the construction by precisely these
$p$-primary multiplicities.  The supplied proof of the latter
was read as a dependency map; its invoked structure and
interpolation theorems are not certified here.
\item The source's zero-line coefficient norm is
$\int_{\mathbf R}|x|^{2j}(1+x^2)^{-s/2}dx$,
displayed in Proposition 11.2.0.4, p.~486.
The strict condition is (M31.16), not a conjectural
arithmetic statement.  If $j\leq [F:\mathbf Q]-1$
and $s\geq2[F:\mathbf Q]$, it holds strictly.
Thus the stated working degree $2[F:\mathbf Q]$
retains every $p$-primary jet without Schanuel.
The source's general boundary cutoff (11.2.0.6), p.~485,
has the same equality problem already corrected by MT-M25;
it does not affect this sufficient even-integer degree.
\item Theorem 11.2.2.1 and Remark 11.2.2.2,
pp.~490--491, use $s\geq2[F:\mathbf Q]$ and
$N(\varepsilon_p)\leq[F:\mathbf Q]$ to assert a
trace with full local multiplicity.  The statement and its
complete supplied proof contain no Schanuel premise.
Their Hilbert quotient and coefficient construction instead
depend on the actual/closed-range identification above.
The correct local quotient can be constructed using
(M31.5); the further circle model and local trace are
treated separately in MT-M32.
\end{enumerate}

\paragraph{The precise downstream gate.}
The dependency is explicit in the text, rather than inferred
from proximity of the statements.  Theorem 11.2.0.2,
pp.~485--486, uses the space asserted closed in
(11.2.0.3); the definition of the orbital Hilbert quotient
at pp.~489--490 and Theorem 11.2.2.1 then use this space.
Theorem 11.3.2.1, pp.~495--496, identifies a localized
geometric kernel with the local orbital action and invokes
the finite-rank trace approximation (11.2.2.12) from
Theorem 11.2.2.1.  Finally the proof of Theorem 12.5.1.1,
p.~520, explicitly invokes Theorem 11.3.2.1 for the
contributions away from the chosen finite set of places,
and Theorem 12.4.1.1 for its remaining local contribution.
The supplied statements and complete displayed proof of
11.3.2.1 were read for this trace.  The full proof of
12.4.1.1 is not claimed audited here.

Consequently the local closure correction preserves a
meaningful local orbital quotient, but does not by itself
prove the geometric-kernel identification in 11.3.2.1
or the narrow-strip trace identity of 12.5.1.1.
The Schanuel refinement belongs to a different branch:
the introductory Expectation 1.7.1, original pp.~37--38,
labels the sharper zero-width multiplicity consequences
as expectations not fully worked out.  No hidden Schanuel
assumption has been established in the positive-width
geometric trace route by this targeted dependency reading.
Its remaining issue is the stated analytic comparison,
not the elementary $p$-primary pole bound.

\subsection*{MT-M32. The corrected primary Euler quotient has a Riesz basis
and a local Poisson trace}

This result proves the local trace after the closure correction
of MT-M31.  It does not identify the global geometric kernel
with this local action.

\paragraph{Statement and conventions.}
Retain $H_s,w_s,U_h$ from MT-M31.  Fix $h>0$, positive
integers $d_1,\ldots,d_b$, and unit complex numbers
$a_1,\ldots,a_b$.  Define
\[
 Q(w)=\prod_{v=1}^{b}(1-a_vw^{d_v}),\qquad
 P=Q(U_h),\qquad
 Z_s=H_s/\overline{PH_s}^{\,H_s}.
 \tag{M32.1}
\]
Let $\lambda_1,\ldots,\lambda_L$ be the distinct roots
of $Q$, all on the unit circle, with multiplicities
$r_1,\ldots,r_L$.  Choose $0\leq\theta_\ell<2\pi$
such that $\lambda_\ell=e^{i\theta_\ell}$, and put
\[
 q_s=\max\{0,\lceil(s-1)/2\rceil\},\qquad
 r'_\ell=\min(r_\ell,q_s),\qquad
 t_{\ell,k}=\frac{\theta_\ell+2\pi k}{h}.
 \tag{M32.2}
\]
All roots and multiplicities here come from a finite Euler
polynomial; no Hecke GRH or Schanuel premise occurs.

The dual of $Z_s$, using the bilinear integral pairing of
MT-M31, has the Riesz basis
\[
 h^{-1/2}x^j e^{it_{\ell,k}x},
 \quad 1\leq\ell\leq L,\quad k\in\mathbf Z,
                  \quad0\leq j<r'_\ell.
 \tag{M32.3}
\]
When every $r'_\ell=0$, the quotient is zero and the basis
is empty.  The corresponding coefficient map on $Z_s$
is a bounded isomorphism onto the $\ell^2$ space of the
indices in (M32.3); its inverse supplies a biorthogonal
Riesz basis of the quotient.

For $f\in C_c^\infty(\mathbf R)$, or more generally for
an exponential-Schwartz test, let
\[
 C_f u(x)=\int_{\mathbf R}f(y)u(x-y)\,dy,
 \qquad M_f(z)=\int_{\mathbf R}f(y)e^{zy}\,dy.
 \tag{M32.4}
\]
The induced operator on $Z_s$ is trace class and satisfies
\[
 \operatorname{Tr}_{Z_s} C_f
      =\sum_{\ell=1}^{L}\sum_{k\in\mathbf Z}
                           r'_\ell M_f(it_{\ell,k})
      =h\sum_{n\in\mathbf Z}f(nh)
                          \sum_{\ell=1}^{L}r'_\ell\lambda_\ell^n.
 \tag{M32.5}
\]
Both expressions converge absolutely.  The map from tests
to trace-class operators is continuous in trace norm, with
uniform tails on bounded sets of tests at a fixed compact
support, and likewise for bounded sets in exponential
Schwartz space.

If $s>2\max_\ell r_\ell-1$, every local jet survives,
and the formula simplifies to
\[
 \operatorname{Tr}_{Z_s}C_f
       =\sum_{v=1}^{b}d_vh
                     \sum_{n\in\mathbf Z}a_v^{-n}f(nd_vh).
 \tag{M32.6}
\]
The strict condition is satisfied in particular by
$s\geq2b$.  For a primary number-field Euler product,
$h=\log p$, $d_v=f_v$ is the residue degree, and
$b\leq[F:\mathbf Q]$; hence the source's sufficient
degree $s\geq2[F:\mathbf Q]$ retains all multiplicities.

\paragraph{A finite matrix-weight model.}
If no index survives the cutoff, MT-M31 gives $Z_s=0$
and all the assertions are immediate.  Assume henceforth
that at least one index survives.
Let $\mathcal N_s\subset H_{-s}$ be the annihilator of
$PH_s$.  By the fully proved formula in MT-M31,
it is the closed span of the functions in (M32.3).
We construct a bounded isomorphism
\[
 \mathcal V_s:=\bigoplus_{\ell=1}^{L}
                     \bigoplus_{j=0}^{r'_\ell-1}L^2(0,h)
                  \longrightarrow\mathcal N_s.
 \tag{M32.7}
\]
For a coefficient vector $c=(c_{\ell,j})$, extend each
$c_{\ell,j}$ periodically with period $h$ and set
\[
 (Ac)(x)=\sum_{\ell,j}
       x^j e^{i\theta_\ell x/h}c_{\ell,j}(x\bmod h).
 \tag{M32.8}
\]
Endpoint representatives have measure zero and make no
difference to this definition.  Write $x=u+nh$ with
$0\leq u<h$, $n\in\mathbf Z$, and define
\[
 v_{\ell,j}(u,n)=(u+nh)^j
                         e^{i\theta_\ell(u/h+n)}.
 \tag{M32.9}
\]
The finite Hermitian matrix
\[
 G_{(\ell,j),(m,k)}(u)
   =\sum_{n\in\mathbf Z}
       \frac{\overline{v_{\ell,j}(u,n)}v_{m,k}(u,n)}
                  {(1+(u+nh)^2)^{s/2}}
 \tag{M32.10}
\]
converges absolutely and uniformly on $0\leq u\leq h$.
Indeed $j,k<q_s$ imply $j+k-s<-1$, so the tails
are bounded by a fixed summable power of $1+|n|$.
Thus $G(u)$ is continuous.  Tonelli, applied first to
finite coefficient sums and then to their squared absolute
values, gives the exact norm identity
\[
           \|Ac\|_{H_{-s}}^2
                =\int_0^h c(u)^*G(u)c(u)\,du.
 \tag{M32.11}
\]

The matrix $G(u)$ is positive definite for every fixed $u$.
If a vector $z$ has $z^*G(u)z=0$, then
\[
       \sum_{\ell,j}z_{\ell,j}
          e^{i\theta_\ell u/h}(u+nh)^j\lambda_\ell^n=0
                       \quad\text{for every }n\in\mathbf Z.
 \tag{M32.12}
\]
Exponential-polynomial sequences with distinct nonzero
$\lambda_\ell$ are linearly independent.
Here is a direct verification.  For one index $\ell$, apply
the finite-difference operator
$\prod_{m\ne\ell}(S-\lambda_m)^{r'_m}$, where
$(Sa)_n=a_{n+1}$, to (M32.12).  Every other component
vanishes.  If the $\ell$-component had a nonzero polynomial
of degree $j$, the resulting component would have the
same degree and a nonzero leading coefficient, multiplied
by $\prod_{m\ne\ell}(\lambda_\ell-\lambda_m)^{r'_m}$.
Applying $(S-\lambda_\ell)^j$ would then give a nonzero
constant multiple of $\lambda_\ell^n$, a contradiction.
Thus all the polynomials vanish.  The powers $(u+nh)^j$
are linearly independent polynomials in $n$, so $z=0$.

Continuity on the compact interval now gives constants
$0<c\leq C<\infty$, depending on this fixed local product
and $s$, with
\[
                  cI\leq G(u)\leq CI\quad(0\leq u\leq h).
 \tag{M32.13}
\]
Consequently (M32.8) is bounded and bounded below, and
its image is closed.  On each summand of $\mathcal V_s$,
the Fourier orthonormal basis
$h^{-1/2}e^{2\pi iku/h}$ maps to the corresponding
function in (M32.3).  Its image therefore contains a dense
subspace of $\mathcal N_s$ by MT-M31.  Conversely, finite
Fourier vectors are dense in $\mathcal V_s$ and their
images belong to the closed space $\mathcal N_s$, so
$A\mathcal V_s\subset\mathcal N_s$.  Its closed
image is exactly $\mathcal N_s$.  This proves (M32.7)
and the asserted Riesz basis, in the actual weighted norm.
No global Riesz assertion for a union of different rational
prime lattices is being made.

The bilinear integral pairing identifies $\mathcal N_s$
isometrically with the continuous dual of $Z_s$.
Taking the dual of the bounded isomorphism from the
standard $\ell^2$ basis onto (M32.3), and using reflexivity
of the Hilbert quotient, proves that the coefficient map
\[
 [u]\longmapsto
 \left(h^{-1/2}\int_{\mathbf R}u(x)x^je^{it_{\ell,k}x}\,dx
                                      \right)_{\ell,k,j}
 \tag{M32.14}
\]
is a bounded bijection with bounded inverse onto $\ell^2$.
This also proves the existence and uniform boundedness
of the biorthogonal quotient vectors.  It is stronger
than merely knowing that the individual coefficient
functionals are continuous.

\paragraph{The exact scalar quotient norm for one untwisted factor.}
For $P=1-U_h$ and $s>1$, the preceding model has a
particularly concrete primal form.  Put
\[
 W(u)=\sum_{n\in\mathbf Z}(1+(u+nh)^2)^{-s/2},
 \qquad
 (Rf)(u)=\sum_{n\in\mathbf Z}f(u+nh),\quad 0\leq u<h.
 \tag{M32.14a}
\]
The function $W$ is continuous and bounded above and
below by positive constants.  For almost every $u$,
Cauchy--Schwarz proves absolute convergence of $Rf(u)$
and
\[
 |Rf(u)|^2\leq W(u)
         \sum_{n\in\mathbf Z}|f(u+nh)|^2w_s(u+nh).
\]
Thus $R:H_s\to L^2((0,h),W(u)^{-1}du)$ is a contraction.
It is onto: a prescribed $F$ in the target has lift
\[
 f(u+nh)=\frac{F(u)}{W(u)w_s(u+nh)},
 \qquad
 \|f\|_{H_s}^2=\int_0^h\frac{|F(u)|^2}{W(u)}du.
 \tag{M32.14b}
\]
The coefficient $J_{2\pi k/h,0}(f)$ is the corresponding
Fourier coefficient of $Rf$ (with the positive-exponent
convention).  Completeness of Fourier series and MT-M31
give $\ker R=\overline{(1-U_h)H_s}$.  Therefore
\[
 H_s/\overline{(1-U_h)H_s}
        \simeq L^2((0,h),W(u)^{-1}du)
 \tag{M32.14c}
\]
is an isometry with its displayed quotient norm.
This is an explicit circle quotient even though the
actual difference-operator range is not closed.

\paragraph{The actual convolution and its trace norm.}
The polynomial-weight translation estimate of MT-M31
shows that $C_f$ is bounded on $H_s$ and $H_{-s}$.
It commutes with $P$ and preserves the closure of its range,
so it acts on $Z_s$.  On the dual functions one computes
directly
\[
 C_f^{\mathrm t}(x^je^{itx})
  =\int f(y)(x+y)^je^{it(x+y)}dy
  =\sum_{a=0}^{j}\binom ja M_f^{(a)}(it)
                                  x^{j-a}e^{itx}.
 \tag{M32.15}
\]
The same identity, paired with $u$, describes the quotient
operator in the coordinate map (M32.14).  It is the
direct sum over $(\ell,k)$ of triangular matrices of
size $r'_\ell$, with diagonal entry $M_f(it_{\ell,k})$.
Their sizes are bounded by the degree of $Q$, which is
fixed in this theorem.

For each integer $N\geq0$ and each derivative order
$0\leq a<\max r'_\ell$, integration by parts gives
\[
 |M_f^{(a)}(it)|\leq C_{N,a}(f)(1+|t|)^{-N}.
 \tag{M32.16}
\]
The constants are bounded by finitely many test seminorms:
on a fixed compact support they are controlled by finitely
many derivatives of $f$, and on exponential-Schwartz
space by finitely many of its defining seminorms.
For a finite $r\times r$ matrix its trace norm is at most
$r$ times its operator norm.  Hence the sum of the trace
norms of these blocks is finite, by (M32.16) with $N>1$.
Conjugating by the bounded isomorphism (M32.14) multiplies
the trace-norm estimate by at most the product of the two
isomorphism norms.  This proves trace class and its stated
continuity.

The block action was obtained on every coordinate of an
actual bounded convolution operator.  Thus the trace-class
sum equals that actual operator, first on finite coordinate
vectors and then by density.  Finite partial sums have
trace $\sum_{\ell,k}r'_\ell M_f(it_{\ell,k})$ over the
selected blocks.  Trace-norm convergence proves the first
identity in (M32.5).  The estimates also give uniform
tails on the stated bounded sets.  In particular no
unjustified passage from a Markushevich system to a
trace expansion is used.

\paragraph{Poisson summation with the normalization fixed.}
For a single $\theta$, periodize
$f(x)e^{i\theta x/h}$ with period $h$.
It is a smooth periodic function whose Fourier coefficients
are
\[
        h^{-1}M_f\bigl(i(\theta-2\pi k)/h\bigr).
\]
Its absolutely convergent Fourier series at zero gives
\[
 \sum_{k\in\mathbf Z}M_f\bigl(i(\theta+2\pi k)/h\bigr)
       =h\sum_{n\in\mathbf Z}f(nh)e^{i\theta n}.
 \tag{M32.17}
\]
The argument applies to exponential-Schwartz tests too,
since their periodization and all derivatives converge
uniformly.  Applying (M32.17) to each of the finitely
many roots proves the second equality in (M32.5).

If $r'_\ell=r_\ell$, count the roots with their
multiplicities in the original factorization of $Q$.
For a factor $1-a_vw^{d_v}$, its roots satisfy
$\lambda^{d_v}=a_v^{-1}$ and their power sum is
\[
 \sum_{\lambda^{d_v}=a_v^{-1}}\lambda^n
   =\begin{cases}
       d_va_v^{-n/d_v},&d_v\mid n,\\
       0,&d_v\nmid n.
     \end{cases}
 \tag{M32.18}
\]
This holds for positive, zero and negative integers $n$,
by the finite geometric-series identity on roots of unity.
Substituting in (M32.5) gives (M32.6).

\paragraph{Adelic normalization and geometric meaning.}
Here is an explicit conversion to a primary number-field
Euler product, with no hidden sign convention.
Choose a norm splitting
$C_F\simeq C_F^1\times\mathbf R$ and Haar measure
$dk\,dx$ with $\int_{C_F^1}dk=1$.
Fix a compact character $\beta$, extended to be trivial
on the chosen norm section, and write vectors in its
sector as $\Phi(k,x)=\beta(k)^{-1}u(x)$.
For a finite place $v\mid p$, write the uniformizer idele
as
\[
             \pi_v=(k_v,-f_v\log p).
\]
The literal operator in (M31.1) becomes
$u(x)-\beta(k_v)u(x+f_v\log p)$.
After the isometric reflection $x\mapsto-x$ this is
the factor in (M32.1) with
\[
              h=\log p,\quad d_v=f_v,
                         \quad a_v=\beta(k_v).
 \tag{M32.19}
\]
For a Bruhat test $\varphi\in\mathcal D(C_F)$, put
\[
 b_\beta(y)=\int_{C_F^1}\varphi(k,y)\beta(k)\,dk,
 \qquad
 \varphi_\beta(k,y)=\beta(k)^{-1}b_\beta(y).
 \tag{M32.20}
\]
The actual convolution on this sector is $C_{b_\beta}$;
after reflection its test is $f(y)=b_\beta(-y)$.
Consequently, when all primary jets survive, (M32.6) reads
\[
     \operatorname{Tr} A_\beta(\varphi)
       =\sum_{v\mid p}\log q_v
                          \sum_{n\in\mathbf Z}\varphi_\beta(\pi_v^n).
 \tag{M32.21}
\]
Only finitely many $n$ contribute for a compact radial
support.  If $\beta$ is unramified at $v$,
$\varphi_\beta$ is constant on the local unit group.
For any multiplicative Haar measure $\mu_v$, the
corresponding local expression is equivalently
\[
 \frac{\log q_v}{\mu_v(\mathcal O_{F_v}^{\times})}
          \int_{F_v^{\times}}\varphi_\beta(g_v)\,d\mu_v(g_v).
 \tag{M32.22}
\]
This specifies exactly where the local-unit volume
normalization enters.  The formulas do not silently
replace an integral by a sum while retaining its volume factor.

The finite matrix $G(u)$ is a concrete Hilbert model for
the finite local Euler polynomial $Q$, with polynomial
jet degrees retaining its nilpotent extensions.  Fourier
series on the norm circle supply the infinite lattice of
Euler poles.  The degree bound used for the local trace
therefore has a direct periodic-function explanation.
It supplies the corrected local quotient and trace intended
by Masullo Theorem 11.2.2.1, original pp.~490--491,
after MT-M31's closure correction and the strict jet cutoff.

The constants in (M32.13) are for a fixed prime product,
character and degree.  No uniform comparison as primes,
characters or global geometric truncations vary is asserted.
In particular this result does not show that the
localized kernel in Masullo Theorem 11.3.2.1,
pp.~495--496, acts on this corrected quotient in the
claimed way, nor that its nuclear limits commute with
the global shell degeneration.  That precise comparison
remains the subsequent geometric-trace gate.

\subsection*{MT-M33. Finite Euler quotients disappear in a separated strip}

This is an unconditional obstruction to a bounded module
identification, not an obstruction to every possible
cutoff or geometric-kernel calculation.

\paragraph{The actual norm and the centered Euler operator.}
For a compact nonempty interval $I=[u,v]$ and $s\geq0$ put
\[
 H_{I,s}=L^2\left(\mathbf R,
       \max(e^{2ux},e^{2vx})(1+x^2)^{s/2}\,dx\right),
 \qquad (U_yf)(x)=f(x-y).
 \tag{M33.1}
\]
The compact norm-one group may be included with its
volume-one Haar measure; the same estimates hold.
Formula (M33.1) is exactly the fixed-character source norm
(9.2.2.31), original PDF p.~416.  The supremum over
subintervals in that definition is attained at the
endpoints of this compact $I$.

Let $T_af(x)=e^{ax}f(x)$.  The source completion
(10.1.1.3), original p.~434, is the closure of
$T_a\Theta$ in (M33.1), where $\Theta$ is its smooth
theta image.  Both original pages were rendered and
read.  The equivariance convention stated just above
that definition agrees with
\[
 U_yT_a=T_a(e^{-ay}U_y),\qquad
            T_aU_yT_a^{-1}=e^{ay}U_y.
 \tag{M33.2}
\]
These identities fix the sign of the centering twist.

In a compact character sector, write a uniformizer idele
as $(k_v,-h_v)$, with $h_v=\log q_v>0$, and put
$b_v=\beta(k_v)$, so $|b_v|=1$.
The uncentered difference of MT-M31 is
\[
 P_{v,0}=1-b_vU_{-h_v}.
\]
Transporting this literal operator by $T_a$ gives
\[
 P_{v,a}=T_aP_{v,0}T_a^{-1}
                 =1-b_ve^{-ah_v}U_{-h_v}.
 \tag{M33.3}
\]
Its Mellin multiplier is
$1-b_ve^{-(a+z)h_v}$, whose zeros have real part $-a$.
The operator is not the unscaled uncentered difference
when $a\ne0$.

\paragraph{The inverse with a quantitative convergent bound.}
If $a+u>0$, (M33.3) is a bounded automorphism of
$H_{I,s}$, with inverse
\[
 P_{v,a}^{-1}
       =\sum_{n=0}^{\infty}b_v^ne^{-anh_v}U_{-nh_v},
 \tag{M33.4}
\]
where convergence is in operator norm.  More precisely,
\[
 \|P_{v,a}^{-1}\|
 \leq 2^{s/4}\sum_{n=0}^{\infty}
       (1+n^2h_v^2)^{s/4}e^{-(a+u)nh_v}<\infty.
 \tag{M33.5}
\]
If $a+v<0$, the inverse instead is
\[
 P_{v,a}^{-1}
      =-\sum_{n=1}^{\infty}b_v^{-n}e^{anh_v}U_{nh_v},
 \tag{M33.6}
\]
with bound obtained by replacing the final exponential
in (M33.5) by $e^{(a+v)nh_v}$ and starting at $n=1$.
Thus every finite product
$P_{S,a}=\prod_{v\in S}P_{v,a}$ is a bounded
automorphism whenever $0\notin a+I$.

\begin{proof}
For every $y\in\mathbf R$, the endpoint exponential
weight and the polynomial inequality from MT-M31 give
\[
 \|U_y\|_{H_{I,s}\to H_{I,s}}
 \leq2^{s/4}(1+y^2)^{s/4}
                      e^{\max(uy,vy)}.
 \tag{M33.7}
\]
Indeed, after the change of variables $x=t+y$, each
endpoint weight acquires its factor $e^{2uy}$ or
$e^{2vy}$, and
$1+(t+y)^2\leq2(1+t^2)(1+y^2)$.
For $B_v=b_ve^{-ah_v}U_{-h_v}$ this proves
\[
 \|B_v^n\|\leq2^{s/4}(1+n^2h_v^2)^{s/4}
                                      e^{-(a+u)nh_v}.
\]
If $a+u>0$, the series is summable and $B_v^n\to0$
in operator norm.  The telescoping identity
$(1-B_v)\sum_{n=0}^{N}B_v^n=1-B_v^{N+1}$,
and its commuting counterpart, prove (M33.4).
If $a+v<0$, apply the same argument to
$B_v^{-1}=b_v^{-1}e^{ah_v}U_{h_v}$; the identity
$(1-B_v)^{-1}=-\sum_{n\geq1}B_v^{-n}$ proves
(M33.6).  The factors and their inverses commute,
so their finite products are automorphisms too.
\end{proof}

\paragraph{All closed invariant subquotients have the same property.}
Let $Y_0\subset Y\subset H_{I,s}$ be closed subspaces
invariant under every $U_y$, $y\in\mathbf R$.
Then $P_{S,a}$ and its inverse preserve $Y_0$ and $Y$:
each term of (M33.4) or (M33.6) preserves them and
they are closed under the norm limit.
Consequently $P_{S,a}$ is an automorphism on $Y/Y_0$
with its quotient norm.  It is enough to assume
invariance under the particular shifts appearing in
the inverse series and their products.

In the actual arithmetic setting,
first work in one fixed compact character sector and put
\[
          J_{a,I,s}=\overline{T_a\Theta}^{\,H_{I,s}}
 \tag{M33.8}
\]
This space is invariant under all norm translations by (M33.2),
since $\Theta$ is translation invariant.
Therefore, without GRH,
\[
 P_{S,a}J_{a,I,s}=J_{a,I,s},\qquad
 \overline{T_aP_{S,0}\Theta}^{\,H_{I,s}}=J_{a,I,s}
                   \quad(0\notin a+I).
 \tag{M33.9}
\]
For the second equality, use
$T_aP_{S,0}=P_{S,a}T_a$, density in (M33.8), and
the fact that an automorphism maps a dense subspace
onto a dense subspace of its whole target.
In particular
\[
       J_{a,I,s}/\overline{T_aP_{S,0}\Theta}=0.
 \tag{M33.10}
\]
For $a=1/2$ and $I=[-\alpha,\alpha]$ this applies
to every $0\leq\alpha<1/2$, including the line
$\alpha=0$.  This concerns the finite Euler quotient,
not the quotient by the full global theta ideal:
the latter need not vanish.

The full compact-group assertion has the same proof:
replace each scalar $b_v$ by the corresponding compact
unitary translation.  Its integer powers have norm one,
so the bounds (M33.5)--(M33.7) are uniform in the compact
character.  For that version the closed subspaces must
also be invariant under these compact translations and
their inverses.  The full theta closure has this
invariance.  Norm-translation invariance alone is the
hypothesis for the fixed-character statement above.

\paragraph{The exact intertwining obstruction.}
Let $Z_{S,s}$ be the corrected nonzero local quotient of
MT-M32, in the original unreflected coordinates:
\[
       Z_{S,s}=H_s/\overline{P_{S,0}H_s}.
 \tag{M33.11}
\]
For instance one prime and $s>1$ gives the nonzero
scalar circle model (M32.14c).
Ordinary norm translations act on this quotient.
Give the same underlying Hilbert space its centered
action
\[
                   V_y=e^{-ay}U_y.
 \tag{M33.12}
\]
Then the polynomial in these operators
\[
       \prod_{v\in S}(1-b_ve^{-ah_v}V_{-h_v})
       =\prod_{v\in S}(1-b_vU_{-h_v})
       =0\quad\hbox{on }Z_{S,s}.
 \tag{M33.13}
\]
The last equality follows from its definition as a
quotient by the closure of the range, and does not
require the actual range to be closed.

Assume $0\notin a+I$, and let $W=Y/Y_0$ be any
closed invariant subquotient as above, with its
inherited norm-translation action.  Every bounded
linear map
\[
       T:(Z_{S,s},V)\longrightarrow W
       \quad\hbox{or}\quad
       T:W\longrightarrow(Z_{S,s},V)
 \tag{M33.14}
\]
which intertwines each norm translation is zero.
For the first direction, applying the same operator
polynomial to the intertwining relation gives
$P_{S,a}T=0$, while $P_{S,a}$ is injective on $W$.
For the second direction it gives $TP_{S,a}=0$,
while $P_{S,a}$ is surjective on $W$.
This proves both assertions.  The twist in (M33.12)
is essential: comparing the untwisted circle action
would compare a different polynomial and would not
prove this statement.

\paragraph{Location of the remaining source comparison.}
The source itself separates the spaces before attempting
the later comparison.  Theorem 11.1.2.1, p.~480,
is explicitly a theorem on the $0$-line.  The local
orbital quotient (11.2.2.1), p.~489, is also explicitly
defined with $a=0$, and Theorem 11.2.2.1,
pp.~490--491, gives its local trace.
In contrast, the absolute and invariant Paley--Wiener
constructions at pp.~487--489 allow a general center
$a$, while Theorem 11.3.2.1, pp.~495--496, sets
$a=1/2$ and permits $I=[-\alpha,\alpha]$ for
$\alpha\geq0$.  Its proof invokes the zero-line
carving theorem and the local orbital trace after
cutoffs and shell degeneration.  The exact norm and
centering convention at pp.~416 and 434 have been
visually checked for this comparison; the cited
transition statements were read in the supplied text.

Equations (M33.9)--(M33.14) show that the needed transition
cannot be a nonzero bounded identification of the
centered local circle module with a closed invariant
subquotient of the narrow Hilbert space.  Any valid
geometric calculation must instead justify the actual
non-equivariant cutoff operators and their limits.
Neither the local Poisson identity in MT-M32 nor the
conditional full narrow-strip spectral trace in MT-M28
supplies that justification.  No conclusion that the
Hecke-GRH endpoint is false is drawn here.

In the cote 1 trace programme (scans 79--80 and 111--112),
this locates the issue in transport from a local
closed-point Euler module to the global trace space.
The local Poisson calculation survives; the separated
strip changes the relevant translation spectrum, so
the geometric comparison needs its own cutoff-limit
argument.

\subsection*{MT-M34. Finite cutoffs erase a genuine Euler condition;
the exact cost of the escaping correction}

This result concerns the literal cutoff-strictness assertion in
Masullo's Theorem 11.1.2.1(2), original PDF p.~480.
It remains a counterexample after replacing every actual
Euler-difference range by its Hilbert closure as in MT-M31.
The uncut Euler quotient and its local trace survive, as
proved in MT-M32; the finite cutoff does not preserve that
quotient.

\paragraph{Statement.}
Fix $s>1$ and $h>0$, and retain the notation
\[
 w_s(x)=(1+x^2)^{s/2},\qquad
 H_s=L^2(\mathbf R,w_s(x)\,dx),\qquad
 U_hf(x)=f(x-h).
 \tag{M34.1}
\]
Let
\[
 V=\overline{(1-U_h)H_s}^{\,H_s},\qquad
 (Rf)(u)=\sum_{n\in\mathbf Z}f(u+nh),\quad0\le u<h.
 \tag{M34.2}
\]
The periodization is a bounded surjection onto $L^2(0,h)$,
and $V=\ker R$ is a proper closed subspace, by the exact
weighted quotient model (M32.14a)--(M32.14c).
If a measurable set $\Omega\subset\mathbf R$ has a
complement containing an interval of length $h$, up to
sets of measure zero, then its indicator multiplier satisfies
\[
               M_\Omega V=M_\Omega H_s.
 \tag{M34.3}
\]
This is equality of the actual images, not merely equality
of their closures.  In particular it holds for every $L>0$
and both
\[
 \Omega_{+,L}=[-L,\infty),\qquad
 \Omega_{0,L}=[-L,L].
 \tag{M34.4}
\]
The corresponding projections converge strongly to the
identity as $L\to\infty$, although $V\ne H_s$.

More generally, if $\rho\in L^\infty(\mathbf R)$ vanishes
almost everywhere on an interval of length $h$, then
\[
                   M_\rho V=M_\rho H_s.
 \tag{M34.5}
\]
This includes every bounded compactly supported cutoff.
In (M34.5) the common image need not be closed.

There is also an exact quantitative form of the obstruction.
Let $g\in H_s$ be compactly supported, put $F=Rg$, and
suppose $F\ne0$.  For either set in (M34.4), take $L$
large enough that $g$ is supported in $\Omega$.  Write
$E=\mathbf R\setminus\Omega$ and
\[
 W_E(u)=\sum_{\substack{n\in\mathbf Z\\u+nh\in E}}
                   \frac1{w_s(u+nh)}.
 \tag{M34.6}
\]
Then the minimum norm of a lift in the proper closed
subspace $V$ with prescribed cutoff $g$ is exactly
\[
 \min_{\substack{f\in V\\M_\Omega f=g}}\|f\|_{H_s}^2
       =\|g\|_{H_s}^2+
                      \int_0^h\frac{|F(u)|^2}{W_E(u)}\,du.
 \tag{M34.7}
\]
In particular, as $L\to\infty$,
\[
 \begin{split}
 \min_{\substack{f\in V\\M_{\Omega_{+,L}}f=g}}
                 \|f\|_{H_s}^2
  &=\|g\|_{H_s}^2+
        h(s-1)L^{s-1}\|F\|_{L^2(0,h)}^2(1+o(1)),\\
 \min_{\substack{f\in V\\M_{\Omega_{0,L}}f=g}}
                 \|f\|_{H_s}^2
  &=\|g\|_{H_s}^2+
       \frac{h(s-1)}2L^{s-1}\|F\|_{L^2(0,h)}^2(1+o(1)).
 \end{split}
 \tag{M34.8}
\]
Thus the norms themselves grow as $L^{(s-1)/2}$,
with leading constants $\sqrt{h(s-1)}\,\|F\|_2$
and $\sqrt{h(s-1)/2}\,\|F\|_2$, respectively.
For the source cutoff $L=\log\Lambda$ these are exact
leading constants in powers of $\log\Lambda$.

\paragraph{Repair outside the cutoff.}
For completeness, the scalar model of MT-M32 gives
\[
 W(u)=\sum_{n\in\mathbf Z}w_s(u+nh)^{-1},\qquad
 |Rf(u)|^2\le W(u)\sum_n|f(u+nh)|^2w_s(u+nh).
 \tag{M34.9}
\]
The positive continuous periodic function $W$ is bounded
above and below.  Therefore $R:H_s\to L^2(0,h)$ is
bounded.  The lift $f(u+nh)=F(u)/(W(u)w_s(u+nh))$
shows surjectivity.  The exact closure identity $V=\ker R$
is the already proved one in (M32.14c), including all
continuous Fourier coefficient constraints.  This also
shows that $V$ is proper.

Choose a half-open interval $I_0=[c,c+h)$ contained in
$\mathbf R\setminus\Omega$ up to a null set.  For almost
every $u\in[0,h)$ there is a unique integer $n(u)$ with
$u+n(u)h\in I_0$.  Define a lift $L_{I_0}$ by
\[
 (L_{I_0}F)(u+n(u)h)=F(u),\qquad
 L_{I_0}F=0\quad\hbox{off }I_0.
 \tag{M34.10}
\]
The piecewise translation $u\mapsto u+n(u)h$ preserves
Lebesgue measure.  Hence
\[
 \|L_{I_0}F\|_{H_s}^2
       \le\bigl(\sup_{x\in I_0}w_s(x)\bigr)\|F\|_2^2,
 \qquad R L_{I_0}F=F.
 \tag{M34.11}
\]
For a prescribed $g\in M_\Omega H_s$, the vector
\[
                        f=g-L_{I_0}Rg
 \tag{M34.12}
\]
belongs to $V$ and satisfies $M_\Omega f=g$.  This proves
(M34.3), and also supplies a bounded right inverse for each
fixed cutoff.  For (M34.5), start with any $f\in H_s$ and
replace it by $f-L_{I_0}Rf$ on an interval where $\rho=0$.
Its $M_\rho$ image is unchanged.  The reverse inclusions
in both image equalities follow from $V\subset H_s$.

For the two sets in (M34.4), the complements contain such
an interval for every finite $L$.  Their indicators increase
pointwise to one.  Dominated convergence applied to
$|f|^2w_s$ proves strong convergence of their multipliers
to the identity.  This proves all the qualitative assertions.

\paragraph{The minimum escaping correction.}
For a feasible $f$ in (M34.7), put $k=f-g$.  Then $k$ is
supported in $E$, $Rk=-F$, and the disjoint supports give
\[
                       \|f\|_{H_s}^2=\|g\|_{H_s}^2+
                                                    \|k\|_{H_s}^2.
 \tag{M34.13}
\]
For each fixed $L$, the measurable function $W_E$ is
positive and bounded above and below by positive constants
on $[0,h)$: the upper bound follows from $W_E\le W$;
for a lower bound, choose one full length-$h$ interval
in $E$ and use its bounded weight.  No continuity across
the endpoints modulo $h$ is needed.

Fiberwise Cauchy--Schwarz gives
\[
 |F(u)|^2\le W_E(u)
         \sum_{u+nh\in E}|k(u+nh)|^2w_s(u+nh).
 \tag{M34.14}
\]
After integration this is the lower bound in (M34.7).
Equality is attained by the measurable function
\[
 k(u+nh)=
 \begin{cases}
 -\dfrac{F(u)}{W_E(u)w_s(u+nh)},&u+nh\in E,\\[6pt]
 0,&u+nh\notin E.
 \end{cases}
 \tag{M34.15}
\]
The lower bound for $W_E$ proves that this function belongs
to $H_s$.  Its absolutely convergent periodization equals
$-F$ and its squared norm is the displayed integral in
(M34.7).  Equality in Cauchy--Schwarz on each fiber, or
strict convexity of the Hilbert norm, shows that this
minimum lift is unique up to equality almost everywhere.

Set $\varphi(t)=(1+t^2)^{-s/2}$ for $t\ge0$.
It is decreasing and
\[
 \int_L^\infty\varphi(t)\,dt
           =\frac{L^{1-s}}{s-1}(1+O(L^{-2})).
 \tag{M34.16}
\]
Indeed $\varphi(t)=t^{-s}(1+O(t^{-2}))$ uniformly for
$t\ge L\ge1$, and both powers can be integrated.
The decreasing-function comparison between a sum on a
lattice of step $h$ and its integral shows, uniformly in
the residue $u\in[0,h)$, that each one-sided tail satisfies
\[
 \sum_{u+nh>L}\varphi(u+nh)
   =\frac1h\int_L^\infty\varphi(t)\,dt+O(\varphi(L)),
 \tag{M34.17}
\]
with the same estimate for $u+nh<-L$ after reflection.
For clarity, the first point beyond $L$ lies between $L$
and $L+h$; comparing the decreasing sum with integrals
starting at that point incurs at most one summand, and
moving the integration endpoint back to $L$ incurs at
most $h\varphi(L)$.  These bounds are uniform in the
residue and also cover endpoint equalities.
Consequently
\[
 \begin{split}
 W_{(-\infty,-L)}(u)
    &=\frac{L^{1-s}}{h(s-1)}(1+o(1)),\\
 W_{\mathbf R\setminus[-L,L]}(u)
    &=\frac{2L^{1-s}}{h(s-1)}(1+o(1)),
 \end{split}
 \tag{M34.18}
\]
uniformly in $u$.  The uniform positive leading terms
allow reciprocation and integration against $|F|^2$.
Equations (M34.7) and (M34.18) prove (M34.8).

\paragraph{A precise missing bounded-lift hypothesis.}
Here is a useful general form of the failed inference.
Let $V$ be a closed subspace of a Hilbert space $H$, and
let orthogonal projections $E_n$ converge strongly to the
identity.  If a fixed $g\in H$ has lifts $v_n\in V$ with
\[
                 E_nv_n=E_ng,\qquad
                            \sup_n\|v_n\|<\infty,
 \tag{M34.19}
\]
then $g\in V$.  To prove this without a subsequence
argument, let $z\in V^\perp$.  Since
$\langle v_n-g,E_nz\rangle=0$,
\[
 |\langle g,z\rangle|
  =|\langle v_n-g,z\rangle|
  \le(\sup_n\|v_n\|+\|g\|)\,\|(1-E_n)z\|
          \longrightarrow0.
 \tag{M34.20}
\]
Thus $g\perp V^\perp$ and $g\in V$.  Equation (M34.8)
proves that this bounded-lift condition fails quantitatively
in the present cutoff problem whenever $Rg\ne0$.

\paragraph{Actual source specialization and exact refutation.}
Take $F=\mathbf Q$, the trivial continuous character of
$C_{\mathbf Q}$, $S=\{\infty\}$, and any rational prime
$p$.  Then $\operatorname{Pic}(\mathbf Z)=1$, the character
is unramified outside $S$, and the source index set
(11.1.2.1), p.~480, includes
\[
                         S_i=\varnothing,
                         \qquad S_{i'}=\{p\}.
 \tag{M34.21}
\]
These are distinct with $S_i\subset S_{i'}$; no place
of $S_i$ divides $p$, while a place of $S_{i'}$ does.
Choose any $s\ge2=2[\mathbf Q:\mathbf Q]$.
Thus every hypothesis of Theorem 11.1.2.1(2) is satisfied.

Normalize Haar measure on $C_{\mathbf Q}^1$ to have
mass one and use $x=\log|g|$ in its trivial sector.
The zero-line Hilbert space is exactly $H_s$.
Let $\Theta$ be the actual smooth theta image in this
sector.  MT-M24 and MT-M25 identify its closure with
$H_s$: the pole-removed completed rational zeta function
has no zero on $\Re z=0$, using the classical zero-free
boundary $\Re z=1$ and its functional equation.
This use of the boundary zero-free theorem requires no
GRH.  The source's one-prime operator is
$P_p=1-U_{-h}$ with $h=\log p$, as checked in MT-M31.
Since
\[
                    1-U_{-h}=-U_{-h}(1-U_h)
 \tag{M34.22}
\]
and $\ker R$ is translation invariant,
$\overline{P_pH_s}=V$; density and boundedness also give
$\overline{P_p\Theta}=V$.

Original Definition 9.3.3.3 and formulas
(9.3.3.4)--(9.3.3.5), p.~428, say that for $\Lambda>1$
the two cutoffs in Theorem 11.1.2.1 are exactly
$M_{\Omega_{+,L}}$ and $M_{\Omega_{0,L}}$ with
$L=\log\Lambda$.  The source defines its cutoff spaces
in (11.1.2.3), p.~480, as the Hilbert closures of the
cutoff smooth images.  For either multiplier $M_\Omega$,
boundedness and (M34.3) therefore imply
\[
 \overline{M_\Omega P_p\Theta}^{\,H_s}
       =\overline{M_\Omega V}^{\,H_s}
       =M_\Omega H_s
       =\overline{M_\Omega\Theta}^{\,H_s}.
 \tag{M34.23}
\]
The common subspace is closed because $M_\Omega$ is an
orthogonal projection.  Thus the source inclusion
(11.1.2.6) is equality for \emph{every} finite $\Lambda>1$,
for both stipulated types of cutoff.  In particular, for
the non-decreasing sequence $\Lambda_k=e^{k+1}$, no
subsequence has the strict inclusion asserted in
Theorem 11.1.2.1(2).  The uncut inclusion (11.1.2.5)
is nevertheless strict, since $V\ne H_s$.

The source's proof on p.~482, immediately after
(11.1.2.16), infers non-strictness of the uncut inclusion
from equality at each cutoff and strong convergence of
the cutoff multipliers to the identity.  Equations
(M34.3)--(M34.8) disprove exactly that inference and
identify the missing uniform control of its lifts.
The source later invokes Theorem 11.1.2.1 in the proof
of its geometric trace Theorem 11.3.2.1, pp.~495--496;
the main Theorem 12.5.1.1, p.~520, in turn invokes
Theorem 11.3.2.1.  This is a targeted dependency trace,
not a proof audit of the remaining kernel arguments.

\paragraph{Input and placement.}
The new cutoff and norm arguments above are elementary
and unconditional; the exact Euler closure and actual
theta-density identification are explicit antecedents
MT-M31--MT-M32 and MT-M24--MT-M25.  Original PDF
pp.~428, 480 and 482 were rendered and visually read
to check the definitions, full strictness hypotheses and
the disputed implication.  MT-M34 refutes this literal
intermediate assertion and proves its quantitative
replacement.  It does not assert that the conditional
Hecke-GRH endpoint $M_H$ is false, or that a different
geometric trace construction is impossible.

For cote 1, scans 79--80 and 111--112, this distinguishes
the trace of a genuine local Euler quotient from a
trace attributed to its finite radial cutoffs.  The
quotient supplies the local closed-point contribution;
the finite cutoff images alone carry no such quotient
condition.  A comparison of these traces therefore
requires quantitative control beyond strong convergence.

\subsection*{MT-M35. Finite Euler products become dense after cutoff;
prime exhaustion and cutoff do not commute}

MT-M34 proves actual cutoff-image equality for one Euler
factor.  The following different argument proves equality
of \emph{closed} cutoff images for every nonzero finite
translation polynomial, including Euler products with
incommensurate logarithms.  It does not assert actual-image
surjectivity for these general products.

\paragraph{Finite-translation statement.}
Let $s\ge0$, and set
\[
 H_s=L^2(\mathbf R,(1+x^2)^{s/2}\,dx),\qquad
 U_tf(x)=f(x-t).
 \tag{M35.1}
\]
Let $P$ be a nonzero finite translation polynomial.  Combine
coincident shifts first, delete zero coefficients, and write
\[
                P=\sum_{j=0}^{m}c_jU_{t_j},\qquad
          t_0<t_1<\cdots<t_m,\quad c_j\ne0.
 \tag{M35.2}
\]
Every term is bounded on $H_s$ and $H_{-s}$, by the
polynomial-weight translation bound (M25.7).
Let $\rho\in L^\infty(\mathbf R)$ have essential support
contained in some half-line, in either orientation.  Then
\[
 \overline{M_\rho P H_s}^{\,H_s}
                 =\overline{M_\rho H_s}^{\,H_s}.
 \tag{M35.3}
\]
If $\Theta\subset H_s$ is any dense linear subspace,
the left side is also $\overline{M_\rho P\Theta}$.
For an indicator $\rho=\mathbf1_\Omega$, with $\Omega$
essentially contained in a half-line, this says
\[
 \overline{M_\Omega P\Theta}^{\,H_s}
      =\overline{M_\Omega\overline{P H_s}}^{\,H_s}
      =M_\Omega H_s.
 \tag{M35.4}
\]
In particular it applies to both source cutoffs
$\Omega_{+,L}=[-L,\infty)$ and $\Omega_{0,L}=[-L,L]$.

\paragraph{Endpoint propagation for the weighted dual.}
We first prove the elementary unique-continuation fact
needed here.  Suppose a measurable locally integrable
function $g$ satisfies, almost everywhere,
\[
                 \sum_{j=0}^{m}d_jg(x-\tau_j)=0,
 \quad \tau_0<\cdots<\tau_m,\quad d_j\ne0,
 \tag{M35.5}
\]
and is supported in a half-line.  Then $g=0$ almost
everywhere.  The assertion for $m=0$ is immediate.
For $m\ge1$, translate the equation so that $\tau_0=0$,
and write $H=\tau_m>0$.

If $g=0$ almost everywhere on $(-\infty,A)$, let
$\varepsilon=\tau_1>0$.  For almost every
$x\in[A,A+\varepsilon)$ all the terms with $j\ge1$
evaluate $g$ below $A$.  Since $d_0\ne0$, equation
(M35.5) gives $g(x)=0$ there.  If by induction $g=0$
below $A+n\varepsilon$, the same equation on
$[A+n\varepsilon,A+(n+1)\varepsilon)$ has every
$x-\tau_j<A+n\varepsilon$ for $j\ge1$, and gives
zero on the next interval.  Countably many such steps
prove $g=0$ on the whole line.

If $g=0$ almost everywhere on $(B,\infty)$, put
$\varepsilon'=H-\tau_{m-1}>0$ and evaluate (M35.5)
at $x=y+H$.  It becomes
\[
 d_mg(y)=-\sum_{j<m}d_jg(y+H-\tau_j).
 \tag{M35.6}
\]
For almost every $y\in(B-\varepsilon',B]$, all
arguments on the right exceed $B$.  Induction in the
opposite direction, on consecutive intervals of length
$\varepsilon'$, proves the result.  Finite translations
of the null exceptional sets remain null at each step,
and their countable union is still null.  Combining the
coincident shifts before selecting the two endpoint gaps
is essential to the proof as stated.

We now use the continuous \emph{bilinear} pairing
\[
            \langle f,g\rangle=\int_{\mathbf R}f(x)g(x)\,dx
             \quad (f\in H_s,\ g\in H_{-s}).
 \tag{M35.7}
\]
It identifies $H_{-s}$ with the continuous linear dual
of $H_s$, with complex conjugation absorbed in the
representing function.  Consequently
\[
              P^{\mathrm t}=\sum_{j=0}^{m}c_jU_{-t_j},
              \qquad M_\rho^{\mathrm t}=M_\rho.
 \tag{M35.8}
\]
There is no coefficient conjugation in this convention.
All functions in $H_{-s}$ are locally integrable, since
the weight is bounded above and below on compact sets.

If $g\in H_{-s}$ annihilates $M_\rho PH_s$, then
$P^{\mathrm t}(\rho g)=0$ in $H_{-s}$, hence almost
everywhere.  The function $\rho g$ lies in $H_{-s}$ and
has essential support in a half-line.  Applying
(M35.5)--(M35.6), with the reflected and sorted shifts
from (M35.8), gives $\rho g=0$.  Thus $g$ annihilates
$M_\rho H_s$.  Conversely every annihilator of the
latter image plainly annihilates the former.  Equality
of continuous annihilators gives equality of the closed
linear subspaces, proving (M35.3).  This last assertion
also follows directly from Hilbert orthogonal
complements after converting the bilinear dual to the
usual conjugate-linear Riesz representation.

Boundedness of $M_\rho P$ and density of $\Theta$ prove
the dense-domain assertion.  For an indicator multiplier,
$M_\Omega$ is an orthogonal projection, so its image is
closed.  Boundedness again gives
$\overline{M_\Omega PH_s}=
\overline{M_\Omega\overline{PH_s}}$.
This proves (M35.4).

\paragraph{Finite Euler products and the cofinal cutoff intersection.}
For a finite set $T$ of rational primes, define
\[
                P_T=\prod_{p\in T}(1-U_{-\log p}),
                   \qquad P_\varnothing=1.
 \tag{M35.9}
\]
These are nonzero translation polynomials.  More generally
any product $\prod_v(1-a_vU_{h_v})$, where $h_v>0$ and
$a_v\ne0$, has nonzero constant and largest-shift
coefficients; combining coincident interior shifts cannot
cancel those endpoints.  Reflecting all shifts gives the
same conclusion for (M35.9).  No rational dependence or
independence of the shifts is used in (M35.3).

Let $\Theta$ be the actual trivial compact-character
theta image for $\mathbf Q$, and take $s>1$.
MT-M24--MT-M25, as used in MT-M34, give
$\overline\Theta^{\,H_s}=H_s$ unconditionally from the
classical zero-free boundary of the completed zeta function.
Put
\[
 V_T=\overline{P_T\Theta}^{\,H_s}
                 =\overline{P_TH_s}^{\,H_s}.
 \tag{M35.10}
\]
If $T\subset T'$, the factorization
$P_{T'}=P_TP_{T'\setminus T}$ gives $V_{T'}\subset V_T$.
For either source cutoff and every finite $T$, equation
(M35.4) gives
\[
                  \overline{M_\Omega P_T\Theta}^{\,H_s}
                                  =M_\Omega H_s.
 \tag{M35.11}
\]
Thus for any increasing cofinal sequence $T_k$ of finite
sets of rational primes,
\[
 \bigcap_{k\ge0}
       \overline{M_\Omega P_{T_k}\Theta}^{\,H_s}
                                   =M_\Omega H_s.
 \tag{M35.12}
\]
Every finite stage is the displayed full cutoff space;
there is no exchange of an image with an intersection
in this argument.  Any quotient between two of these
closed cutoff stages is zero.

\paragraph{The uncut cofinal intersection is zero.}
In contrast, for the same $s>1$,
\[
                         \bigcap_{k\ge0}V_{T_k}=\{0\}.
 \tag{M35.13}
\]
Indeed $\int_{\mathbf R}(1+x^2)^{-s/2}dx<\infty$,
so Cauchy--Schwarz gives a continuous embedding
$H_s\subset L^1(\mathbf R)$.  Each functional
\[
                  J_t(f)=\int_{\mathbf R}f(x)e^{itx}\,dx
 \tag{M35.14}
\]
is continuous on $H_s$, and $t\mapsto J_t(f)$ is
continuous by dominated convergence.  The norm-translation
convention gives
\[
       J_t(U_{-\log p}f)=e^{-it\log p}J_t(f).
 \tag{M35.15}
\]
Hence if $p\in T$, every $f\in V_T$ satisfies
$J_{2\pi n/\log p}(f)=0$ for all $n\in\mathbf Z$.
If $f$ belongs to the intersection in (M35.13), cofinality
makes these vanishings valid for every rational prime $p$.

The union of these frequency grids is dense in
$\mathbf R$.  For any fixed $t$, choose primes
$p_j\to\infty$ and choose an integer $n_j$ nearest to
$t\log p_j/(2\pi)$; then
\[
       \left|t-\frac{2\pi n_j}{\log p_j}\right|
                                  \le\frac\pi{\log p_j}\longrightarrow0.
 \tag{M35.16}
\]
Continuity therefore implies $J_t(f)=0$ for every real
$t$.  To recall the uniqueness argument in this setting,
Fourier inversion for a smooth compactly supported test
$\phi$ and Fubini give $\int f\phi=0$: the Fourier
transform of $\phi$ is integrable and $f\in L^1$.
Thus the distribution defined by $f$ vanishes, and
$f=0$ almost everywhere.  This proves (M35.13).

Equations (M35.12)--(M35.13) locate the precise failure
of a cutoff/prime-exhaustion interchange.  For every
fixed finite $L$, the cofinal cutoff intersection is
the full cutoff space; these spaces increase with $L$
and their union is dense in $H_s$.  The uncut cofinal
intersection is zero.  The result concerns subspaces
and their specified closures; it does not identify
the limits of any separately defined kernel operators.

\paragraph{Comparison with the literal Schwartz image in
Theorem 11.1.2.1(3).}
The source's original PDF p.~480 defines the cutoff
spaces in (11.1.2.3) exactly as the closures in (M35.11).
Its cofinal set-theoretic limit (11.1.2.4) is therefore
(M35.12), for $F=\mathbf Q$, $S=\{\infty\}$ and the
trivial character.  Its assertion (3), however, is an
identification with the closure of the specific Schwartz
image (11.1.2.7).  That separate assertion is \emph{not}
refuted by the present example: its stated right side is
also the full cutoff space, as we now check directly.

Every idele class of $\mathbf Q$ has a representative
with positive real component $t=|g|$ and all finite
components in $\mathbf Z_p^\times$: multiply by a rational
number to cancel the finitely many nonzero valuations,
and choose its sign to make the real component positive.
For $S=\{\infty\}$, the sum in (11.1.2.7), before
its cutoff, is then
\[
 \sum_{\xi\in\mathbf Q^\times}
   \left(f_\infty\otimes\prod_p\mathbf1_{\mathbf Z_p^\times}\right)
                   (\xi g)
                          =f_\infty(t)+f_\infty(-t).
 \tag{M35.17}
\]
Indeed only $\xi=1$ and $\xi=-1$ are units at every
finite prime.  This function is in the trivial compact
sector.  The source conditions in this case are
\[
 f_\infty\in\mathcal S(\mathbf R),\qquad
                  f_\infty(0)=0,\qquad
                  \widehat f_\infty(0)=\int_{\mathbf R}f_\infty=0.
 \tag{M35.18}
\]

Fix $L>0$ and take $u\in C_c^\infty(\Omega^\circ)$
for either $\Omega=\Omega_{+,L}$ or $\Omega_{0,L}$.
For $t\ne0$ define
\[
                    f_0(t)=\tfrac12u(\log|t|),
                    \qquad f_0(0)=0.
 \tag{M35.19}
\]
This is an even smooth compactly supported function
whose support is separated from zero, so
$f_0(e^x)+f_0(-e^x)=u(x)$ and $f_0(0)=0$.
Choose an even smooth compactly supported function $r$
with support in $0<|t|<e^{-L}$ and with $\int r\ne0$.
Then
\[
                    f_\infty=f_0-
                            \frac{\int f_0}{\int r}\,r
 \tag{M35.20}
\]
satisfies both conditions in (M35.18).  The repair term
in (M35.17) is supported in $x<-L$, outside either
cutoff.  Consequently the source's cutoff Schwartz
image contains $u$.  Smooth compactly supported
functions in $\Omega^\circ$ are dense in $M_\Omega H_s$:
truncate away from the endpoints and infinity, and then
use ordinary smooth approximation where the weight is
bounded above and below.  Thus the closure of the
right side in (11.1.2.7) is precisely $M_\Omega H_s$,
as claimed here.  No further refutation of part (3)
follows from this rational example.

\paragraph{Input and placement.}
The endpoint-propagation and cutoff-density theorems
are unconditional Hilbert-space deductions, with no
individual-zero or global-Riesz-basis assumption.
The actual theta-density input is the already proved
MT-M24--MT-M25.  Original PDF pp.~428, 480 and 482
were visually checked in MT-M34; the displayed formula
(11.1.2.7) and all its source hypotheses are retained
in the explicit comparison above.  MT-M35 strengthens
the corrected cutoff analysis without declaring the
conditional Hecke-GRH endpoint $M_H$ false or replacing
the remaining geometric-kernel comparison by a bounded
equivariant map.

\subsection*{MT-M36. The literal section average is zero;
normalized compact projections and genuine idelic cutoffs}

This is a targeted measure and normalization check of
Masullo's formula (11.3.2.2), original PDF p.~495.
The formula uses ordinary indicator functions and the
ordinary Haar integral on the idelic fundamental domain.
It does not display a Dirac distribution or a rescaled
limit of neighborhoods.  Under precisely that
interpretation its value is identically zero.
We prove this, and then give two legitimate replacements
with their exact scopes and trace properties.

\paragraph{The source's measure conventions.}
Write $\mathbb I_F=\mathrm{GL}_1(\mathbf A_F)$,
$\mathbb I_F^1=\ker|\cdot|_F$, and let
$\tau:\mathbf R_{>0}\to\mathbb I_F$ be a continuous
homomorphic section of the norm.  The source defines
\[
 E_F^1=\bigcup_{i=1}^{h_F}
              g_i(E_0\times\widehat{\mathcal O}_F^\times)
                         \subset\mathbb I_F^1
 \tag{M36.1}
\]
in Definition A.2.3.6, (A.2.3.14), original p.~549,
and explicitly gives this set the subspace topology.
It is not assigned the discrete topology.  Definition
A.3.2.1, pp.~559--560, uses Haar measures on ideles
and their class groups; the counting measure is on
$F^\times$.  The same distinction is explicit in
the compact-character integration of \S2.2.5.
The normalization (A.3.2.1), p.~560, gives volume one
on $C_{\mathbf Q}^1$.  The general measure discussion
in \S12.2, p.~504, is ordinary completed Radon/Lebesgue
integration of measurable functions.

The displayed two-variable function (11.3.2.2) is
therefore, with its suppressed Haar differential written,
\[
 B_0(x,y)=\sum_{\xi\in F^\times}\int_{E_F^1}
    \mathbf1_{\tau(\mathbf R_{>0})}(\xi gx)
    \mathbf1_{\tau(\mathbf R_{>0})}(\xi gy)\,d\mu_1(g).
 \tag{M36.2}
\]
The tensor symbol in the source just separates the two
named variables of these indicator functions.
For every $x,y\in\mathbb I_F$,
\[
                              B_0(x,y)=0.
 \tag{M36.3}
\]

\begin{proof}
Haar measure on $\mathbb I_F^1$ has no atoms.
Indeed this group contains an infinite compact subgroup:
embed $\mathcal O_{F_v}^\times$ at any one finite place
and use $1$ at every other place.  If one singleton had
positive Haar measure $c$, every singleton would have
the same measure, and $N$ distinct points in this compact
subgroup would have measure $Nc$ for every $N$.
This contradicts finiteness of a Radon measure on a
compact set.  Consequently every singleton in $E_F^1$
is null as well.  Equivalently, the quotient Haar
measure pulled back to a fundamental domain is atomless;
the discrete quotient by $F^\times$ does not change this
fact.  For $F=\mathbf Q$ this is particularly explicit:
$E_{\mathbf Q}^1=\{1\}\times\widehat{\mathbf Z}^\times$,
with its usual compact, atomless, volume-one Haar measure.

The section image is closed, hence Borel: it is the
fixed-point set of the continuous map
$z\mapsto\tau(|z|_F)$ on the Hausdorff group
$\mathbb I_F$.  Fix $\xi$ and $x$.  If
$\xi gx=\tau(t)$ with $g\in E_F^1$, the product formula
and $|g|_F=1$ force $t=|x|_F$.  Hence the only possible
value of $g$ is
\[
                       g=\xi^{-1}\tau(|x|_F)x^{-1}.
 \tag{M36.4}
\]
The first indicator in the integrand of (M36.2) is
supported on at most this one point.  Its product with
the second indicator has integral zero.  The field
$F$ is countable, so (M36.2) is a countable sum of zeros.
This proves (M36.3), without any interchange of
conditionally convergent expressions.
\end{proof}

\paragraph{The precise normalization failure.}
The source's indentation condition (12.2.1.2), original
p.~505, requires an orbit integral equal to one.
For a discrete subgroup $H$ this reads
$\sum_{\xi\in H}b(\xi x)=1$, with $b\ge0$.
Its two-variable lift is the averaged product
\[
                    B_b(x,y)=\sum_{\xi\in H}b(\xi x)b(\xi y),
 \tag{M36.5}
\]
as in (12.2.2.1), p.~508.  Nonnegativity and reindexing
give the necessary orbit identity
\[
 \sum_{\eta\in H}B_b(x,\eta x)
    =\left(\sum_{\xi\in H}b(\xi x)\right)^2=1.
 \tag{M36.6}
\]
Thus the zero function in (M36.3) cannot provide this
normalized lift.  The source's actual lifting rule
(12.2.2.2), p.~508, is multiplication of the pulled-back
kernel by the two-variable factor.  Inserting (M36.2)
into that rule produces the zero kernel.

This also gives a concrete nonzero test for the literal
interface in Theorem 11.3.2.1.  Take $F=\mathbf Q$,
the trivial character, $S=\{\infty\}$, $s=2$, and any
allowed interval $I$.  Fix a rational prime $p$ and
choose a nonnegative $\psi\in C_c^\infty((0,\infty))$
supported in a small neighborhood of $\log p$, containing
no other element of $\{n\log p:n\in\mathbf Z\}$, with
$\psi(\log p)>0$.  The test
\[
                             \varphi(g)=\psi(\log|g|)
 \tag{M36.7}
\]
is compactly supported on $C_{\mathbf Q}$, is invariant
under every finite local unit group, and satisfies all
the displayed ramification and support hypotheses.
For a local uniformizer $\pi_p$, whose idele norm is
$p^{-1}$, the right side of (11.3.2.1) is
\[
       \frac{\log p}{\mu_p(\mathbf Z_p^\times)}
                           \psi(\log p)>0,
 \tag{M36.8}
\]
because only the power $\pi_p^{-1}$ contributes.
Under the stated multiplication rule the left side
with (M36.2) is zero.  This refutes the literal
normalization/interface, not a reinterpretation using
a separately specified distributional regularization.

\paragraph{A normalized finite compact-level replacement.}
Let $K$ be an infinite compact metrizable profinite
abelian group with probability Haar measure; in the
actual rational example $K=C_{\mathbf Q}^1\simeq
\widehat{\mathbf Z}^\times$.  For an open subgroup
$U\subset K$, put $m_U=\mu(U)=1/[K:U]$ and define
\[
 a_U(k)=\frac{\mathbf1_U(k)}{m_U},\qquad
 B_U(k,l)=\int_Ka_U(rk)a_U(rl)\,d\mu(r).
 \tag{M36.9}
\]
Then
\[
 B_U(k,l)=\frac{\mathbf1_U(kl^{-1})}{m_U}
         =\sum_{\chi\in\widehat{K/U}}\chi(k)\overline{\chi(l)}.
 \tag{M36.10}
\]
The first equality follows because the two cosets of
$U$ in the integrand are either disjoint or identical;
in the latter case their intersection has measure $m_U$.
The second is the elementary character orthogonality
identity on the finite abelian group $K/U$.

Consequently the integral operator with kernel $B_U$
is exactly the orthogonal projection
\[
      (E_Uf)(k)=\frac1{m_U}\int_U f(ku)\,d\mu(u)
                    \quad\hbox{on }L^2(K).
 \tag{M36.11}
\]
Its range is the space of $U$-invariant functions,
and
\[
 \operatorname{rank}E_U=\operatorname{Tr}E_U=m_U^{-1},
       \quad\|E_U\|=1,\quad
       \|B_U\|_{L^2(K\times K)}^2=m_U^{-1}.
 \tag{M36.12}
\]
All these assertions follow either from (M36.10) or
from the orthonormal basis of normalized coset indicators.

If the neighborhood indicators are left unnormalized,
their averaged product is instead
\[
 \int_K\mathbf1_U(rk)\mathbf1_U(rl)\,d\mu(r)
       =m_U\mathbf1_U(kl^{-1})=m_U^2 B_U(k,l).
 \tag{M36.13}
\]
Its operator norm is $m_U^2$ and its trace is $m_U$.
Along a nested sequence of open subgroups $U_j$ forming
a neighborhood base at $1$, one has $m_{U_j}\to0$,
because their intersection is the Haar-null singleton.
Thus these unnormalized operators tend to zero even
in trace norm.  The normalized $E_{U_j}$ instead
converge strongly to the identity.  To prove the latter,
locally constant functions are dense in $L^2(K)$:
continuous functions are dense, and uniform continuity
approximates each by functions constant on the cosets
of a sufficiently small open subgroup.  Every such
function is fixed by all sufficiently late $E_{U_j}$,
and the projections have norm one.  Density proves
strong convergence.  There is no operator-norm
convergence to the identity, since each range is finite
dimensional whereas $L^2(K)$ is infinite dimensional.

These projections are a genuine normalized regularization
of the compact section.  Their diverging diagonal values
and Hilbert--Schmidt norms in (M36.12) also show why they
cannot be replaced by an ordinary point-indicator kernel.

\paragraph{Trace compatibility on the actual Hilbert modules.}
The same averaging defines orthogonal projections $E_U$
for any strongly continuous unitary representation of
$K$ on a Hilbert space $\mathcal H$.  They converge
strongly to the identity as $U$ shrinks to $1$:
\[
 \|E_Uv-v\|
       \le\sup_{u\in U}\|R(u)v-v\|\longrightarrow0.
 \tag{M36.14}
\]
This includes the compact action on the weighted idelic
Hilbert spaces and their closed invariant quotients:
the norm weight depends only on the radial variable.
If $T$ is trace class, then
\[
 \|E_UT-T\|_1\to0,\qquad
 \|TE_U-T\|_1\to0,\qquad
             \operatorname{Tr}(E_UT)\to\operatorname{Tr}(T).
 \tag{M36.15}
\]
For a complete proof, write a singular-value expansion
$T=\sum_ns_n u_n\otimes v_n^*$ with
$\sum_ns_n<\infty$ and unit vectors.  The first
trace-norm difference is at most
$\sum_ns_n\|(E_U-1)u_n\|$, which tends to zero by
(M36.14) and domination by $2s_n$.  For the second
use self-adjointness of $E_U$ and the same argument
on $v_n$.  Continuity of the trace gives the last claim.

There is exact stabilization for a test already at
finite compact level.  If a convolution test $\varphi$
is invariant under an open subgroup $U_0\subset K$,
its convolution operator $A_\varphi$ on any of these
modules satisfies
\[
        E_U A_\varphi=A_\varphi E_U=A_\varphi
                          \qquad(U\subset U_0).
 \tag{M36.16}
\]
Indeed translation of the test by any $u\in U_0$
does not change its convolution operator; integration
over $U$ gives both identities, using commutativity.
Thus if $A_\varphi$ is trace class its trace is preserved
exactly by all these finite-level projections.  On the
arithmetic quotients that trace-class hypothesis is
available unconditionally in the wide strips of MT-M27
and, with its explicit Hecke-GRH premise, in the narrow
positive-width strips of MT-M28.  Formula (M36.16) does
not supply that analytic hypothesis by itself.
For general number fields the compact group can have
connected archimedean factors; the finite-rank assertions
(M36.10)--(M36.12) here are expressly for the profinite
group, not for an untruncated archimedean character family.

\paragraph{A genuine discrete-idelic indentation and its orbital allocation.}
There is a different, also explicit, correction for
the discrete $\mathbf Q^\times$-normalization in
(M36.5)--(M36.6).  Use the standard section and set
\[
 D=\tau(\mathbf R_{>0})\,\widehat{\mathbf Z}^\times
       \subset\mathbb I_{\mathbf Q},\qquad b=\mathbf1_D.
 \tag{M36.17}
\]
The set $D$ is an open subgroup, hence also closed.
Every idele has a unique expression $q d$ with
$q\in\mathbf Q^\times$ and $d\in D$: cancel its
finite valuations by a rational number and then fix
the sign of its real component.  Uniqueness follows
from $\mathbf Q^\times\cap D=\{1\}$.
The map $D\to C_{\mathbf Q}$ is a homeomorphism;
openness follows from the quotient map and the
openness of $D$.  Therefore $b$ is continuous and
nonnegative, satisfies $\sum_qb(qx)=1$, and its support
meets the inverse image of any compact subset of
$C_{\mathbf Q}$ in a compact set.  These are exactly
the three indentation conditions of Lemma 12.2.1.1.
No auxiliary connected-domain assumption is needed
for this direct verification.

For $x\in D$, its two-variable lift has
\[
                   B_b(x,qx)=\mathbf1_{q=1}.
 \tag{M36.18}
\]
Thus it provides the normalized total orbital integral,
but places all of that diagonal mass in the identity
translate.  This allocation is not forced by the
normalization.  For a rational prime $p$ and
$0\le t\le1$, another valid indentation is
\[
                    b_t=(1-t)\mathbf1_D+t\mathbf1_{pD}.
 \tag{M36.19}
\]
Its support is a union of two clopen translates of
$D$, it has the same compactness property, and
$\sum_qb_t(qx)=1$.  Direct substitution, for $x\in D$,
gives
\[
 B_{b_t}(x,qx)=
 \begin{cases}
       (1-t)^2+t^2,&q=1,\\
       t(1-t),&q=p\text{ or }p^{-1},\\
       0,&\text{otherwise}.
 \end{cases}
 \tag{M36.20}
\]
The sum of these values is one.  Consequently, when
an integrable quotient diagonal is pulled back and
multiplied by $B_{b_t}$, its total orbital integral is
unchanged, while its individual rational-translate
contributions vary with $t$.  This follows directly
from (M36.20), since $x$ and $qx$ represent the same
class; no trace-limit interchange is used.

The normalized compact projection and the genuine
discrete-idelic indentation solve different problems.
Neither calculation identifies the source's desired
prime-by-prime geometric mass with the local Poisson
trace of MT-M32.  A correct comparison must specify
its normalized lift and prove the required orbital
allocation and limiting identity.  Positivity or the
conditional Hecke-GRH endpoint does not follow from
the measure repair alone.

\paragraph{Audit scope.}
Original PDF pp.~495, 505, 508, 549 and 560 were
rendered and visually read.  The definitions and
normalization text on pp.~504--506 and 559--560 were
also read from the supplied mathematical text.
This is not an audit of all the kernel theorems in
Chapter 12.  The literal failure (M36.3) and all the
replacement constructions above are unconditional.
The endpoint $M_H$ remains conditional throughout.

For cote 1, scans 79--80 and 111--112, the distinction
is between a normalized lift of a quotient trace and
its claimed local closed-point contributions.  The
former is available in the two explicit forms above;
the latter still requires a canonical allocation and
a proved geometric comparison.

\subsection*{MT-M37. Compatible nonnegative indentations cannot carry
the external prime shells}

The normalized replacements in MT-M36 do not by themselves
establish an orbital decomposition.  Here is a stronger
obstruction to one particular repair: a \emph{single}
nonnegative indentation normalized both globally and on
the initial local subgroup must vanish outside that subgroup.
This is exactly the simultaneous normalization offered in
Masullo's Proposition 12.2.2.4, original pp.~512--513.
An exact defect identity quantifies the obstruction
when the local normalization is relaxed.

\paragraph{The two normalizations force support.}
Let $G$ be an abelian group, $H\subset G$ a countable
subgroup, and $G_0\subset G$ a subgroup such that
\[
                 G=HG_0,\qquad H_0=H\cap G_0.
 \tag{M37.1}
\]
Suppose $b:G\to\mathbf R_{\ge0}$ satisfies
\[
 \sum_{h\in H}b(hx)=1\quad(x\in G),\qquad
 \sum_{h\in H_0}b(hx)=1\quad(x\in G_0).
 \tag{M37.2}
\]
Then
\[
                       b(y)=0\qquad(y\notin G_0).
 \tag{M37.3}
\]
For its averaged product
\[
                     B_b(x,y)=\sum_{h\in H}b(hx)b(hy),
 \tag{M37.4}
\]
one has
\[
                  B_b(x,y)=0\qquad(x^{-1}y\notin G_0).
 \tag{M37.5}
\]
In particular, if $G_0\subset G_k\subset G_{k+1}$ and
$\alpha_k=\mathbf1_{G_k}$, then
\[
       (\alpha_{k+1}(y)-\alpha_k(y))B_b(x,y)=0
                                      \quad(x\in G_0,\ y\in G).
 \tag{M37.6}
\]

\begin{proof}
For $x\in G_0$, subtract the second identity in (M37.2)
from the first.  Nonnegativity gives
\[
                       \sum_{h\in H\setminus H_0}b(hx)=0,
 \tag{M37.7}
\]
so every term is zero.  Any $y\notin G_0$ has the form
$hx$ with $h\in H$ and $x\in G_0$ by (M37.1).
That $h$ cannot belong to $H_0$, so (M37.7) proves
(M37.3).  If a term in (M37.4) is nonzero, both
$hx$ and $hy$ belong to $G_0$; hence $x^{-1}y\in G_0$.
This proves (M37.5).  The series (M37.4) is well-defined:
each factor is at most one by (M37.2), so its nonnegative
sum is at most $\sum_hb(hx)=1$.
Finally, for $x\in G_0$, either $y\notin G_0$ and
(M37.5) applies, or $y\in G_0\subset G_k$ and the
indicator difference is zero.  This proves (M37.6).
\end{proof}

The use of $hx$ instead of $h^{-1}x$ makes no difference
to the normalization, since inversion permutes each
subgroup.  The nonnegativity assumption is indispensable
to the termwise implication in (M37.7); it is part of
the source's definition of an indentation.

\paragraph{Exact defect and external orbital mass.}
Keep (M37.1) and the global normalization in (M37.2),
but do not impose the local normalization.  For $x\in G_0$
and a coset $c\in H/H_0$, define
\[
 A_c(x)=\sum_{h\in c}b(hx),\qquad
 \delta(x)=1-A_{H_0}(x).
 \tag{M37.8}
\]
Then $A_c\ge0$, $\sum_cA_c=1$, and $0\le\delta\le1$.
The total mass assigned to external translates is
\[
 \begin{split}
 W_{\mathrm{ext}}(x)
   &:=\sum_{\eta\in H\setminus H_0}B_b(x,\eta x)\\
   &=1-\sum_{c\in H/H_0}A_c(x)^2.
 \end{split}
 \tag{M37.9}
\]
In particular
\[
       2\delta(x)(1-\delta(x))
          \le W_{\mathrm{ext}}(x)
          \le2\delta(x)-\delta(x)^2\le2\delta(x).
 \tag{M37.10}
\]

\begin{proof}
Tonelli for nonnegative countable sums gives
\[
 \sum_{\eta\in H}B_b(x,\eta x)
    =\sum_{h\in H}\sum_{\eta\in H}b(hx)b(h\eta x)=1.
 \tag{M37.11}
\]
If $\eta$ ranges only over $H_0$, then $h\eta$ runs
over the same coset of $H_0$ as $h$.  Summing coset
by coset gives
\[
                   \sum_{\eta\in H_0}B_b(x,\eta x)
                              =\sum_c A_c(x)^2.
 \tag{M37.12}
\]
Subtracting proves (M37.9).  The identity coset has
mass $1-\delta$ and the other cosets have total mass
$\delta$, so
\[
       0\le\sum_{c\ne H_0}A_c^2\le\delta^2.
 \tag{M37.13}
\]
Substitution into (M37.9) gives both bounds in (M37.10).
\end{proof}

Here is the precise consequence for varying indentations.
Suppose $G$ carries a Haar measure, $b$ is measurable,
$D\subset G_0$ is a
measurable fundamental domain for $H$, and
$d\in L^1(D)$ is a prescribed quotient diagonal.
For any measurable selection of external translates,
its integral after multiplication by $B_b$ has absolute
value at most
\[
                    2\int_D\delta(x)|d(x)|\,d\mu_G(x).
 \tag{M37.14}
\]
Indeed the selected nonnegative weights are bounded by
$W_{\mathrm{ext}}$, and the quotient diagonal is the
same at $x$ and $\eta x$.  Equation (M37.10) proves
the bound; absolute convergence follows from
$W_{\mathrm{ext}}\le1$ and $d\in L^1$.
Thus if a sequence $b_j$ has defects $\delta_j\to0$
almost everywhere on $D$, every such external integral
tends to zero for the fixed diagonal, by dominated
convergence.  If the diagonals also vary, the same
conclusion follows either from
\[
       \|\delta_j\|_\infty\to0,\quad
                       \sup_j\|d_j\|_1<\infty,
 \tag{M37.15}
\]
or from $d_j\to d$ in $L^1$ and
$\delta_j\to0$ almost everywhere.  In the latter case
\[
 \int_D\delta_j|d_j|
       \le\|d_j-d\|_1+\int_D\delta_j|d|\longrightarrow0.
 \tag{M37.16}
\]
These are explicit sufficient conditions, not an
assertion that a varying source kernel meets them.

\paragraph{A trace-norm criterion for the varying diagonals.}
The second alternative has a useful operator formulation.
Fix a Hilbert realization $\mathcal H=L^2(D,\mu)$.
There is a canonical linear contraction
\[
 \mathcal D:\mathcal S_1(\mathcal H)\longrightarrow L^1(D,\mu),
 \qquad
 \operatorname{tr}(M_fT)=\int_D f\,\mathcal D(T)\,d\mu
                   \quad(f\in L^\infty(D,\mu)),
 \tag{M37.20}
\]
where $M_f$ denotes multiplication and $\mathcal S_1$ the
trace-class operators.  In particular
$\int_D\mathcal D(T)\,d\mu=\operatorname{tr}T$.
This is the trace diagonal; no pointwise restriction of an
arbitrary measurable two-variable kernel is being asserted.

\begin{proof}
Write a singular-value expansion
$T=\sum_n s_n\,u_n\otimes\overline{v_n}$, with
$\sum_n s_n=\|T\|_1$ and unit vectors $u_n,v_n$.
Set $\mathcal D(T)=\sum_n s_nu_n\overline{v_n}$.
Cauchy--Schwarz gives
$\|u_n\overline{v_n}\|_{L^1}\le1$, so the series converges
in $L^1$ with norm at most $\|T\|_1$.
For every bounded $f$, the rank-one trace formula and
trace-norm convergence give (M37.20).
Bounded measurable functions separate $L^1$ functions;
therefore (M37.20) makes the definition independent of
the expansion and proves linearity.  Applying the same
bound to $T_1-T_2$ proves contraction.
\end{proof}

Consequently, if $\|T_j-T\|_1\to0$ in this fixed realization,
$d_j=\mathcal D(T_j)$, and $\delta_j\to0$ almost everywhere,
every selected external orbital contribution in (M37.14)
tends to zero, with the explicit bound
\[
 2\|T_j-T\|_1+
       2\int_D\delta_j\,|\mathcal D(T)|\,d\mu
                         \longrightarrow0.
 \tag{M37.21}
\]
Indeed $\|d_j-\mathcal D(T)\|_{L^1}\le\|T_j-T\|_1$,
and the second term tends to zero by dominated convergence.
The assertion also applies to a fixed closed Hilbert subspace,
or to a closed subquotient, realized isometrically by $J$ in
this $L^2$ space: replace $T_j$ by $JT_jJ^*$, which preserves
trace norms and traces because it adds only zero singular
values.  If $w$ is measurable with $0<w<\infty$ almost
everywhere, a weighted $L^2(D,w\mu)$ realization can first
be conjugated by the unitary $Uu=\sqrt w\,u$ into $L^2(D,\mu)$.
The trace identity gives
$\mathcal D_\mu(UTU^{-1})=w\mathcal D_{w\mu}(T)$,
so (M37.21) uses the density for the fixed reference measure
after that conjugation.
This formulation requires that the proposed geometric diagonal
has actually been identified with the fixed trace diagonal.
Trace class alone does not identify it with a chosen version
of a measurable kernel on a null diagonal.

\paragraph{The source's compatible-indentation construction.}
Original Proposition 12.2.2.4, p.~512, assumes a second
countable locally compact abelian $G$, a discrete $H$
with counting measure, an increasing exhaustive family
of open subgroups $G_k$, and $H_k=H\cap G_k$.
Among its stated hypotheses are a fundamental domain
$D\subset G_0$ and the analogous fundamental-domain
property for each $G_k/H_k$.  In particular
$G=HD\subset HG_0$, which gives (M37.1).
Its asserted indentation $b$ is global and has
$\alpha_kb$ an $H_k$-indentation on $G_k$ for every
$k$.  At $k=0$ this is precisely the second
normalization in (M37.2).  The first is the global
indentation condition (12.2.1.2), p.~505.

Therefore every nonnegative indentation with the
asserted simultaneous normalization satisfies
(M37.3)--(M37.6).  This conclusion uses only the
two normalizations and $G=HG_0$; none of the additional
differentiation or domain-regularity hypotheses is
invoked.  The source's own construction on p.~513
explicitly extends an $H_0$-indentation on $G_0$
by zero to $G$, consistent with (M37.3).
In particular its tail integrand (12.2.2.26), evaluated
at $x\in D\subset G_0$ and $y=h^{-1}x$ outside $G_k$,
is already zero for each $k$.  The proposition's
zero-tail limit is not refuted by this observation.

\paragraph{The actual rational external-prime shells.}
Take
\[
 G=\mathbb I_{\mathbf Q},\quad H=\mathbf Q^\times,
 \quad G_0=\mathbf R^\times\times
                           \widehat{\mathbf Z}^\times,
 \quad H_0=\{1,-1\}.
 \tag{M37.17}
\]
Every idele can have its finite valuations canceled
by a rational number, so $G=HG_0$ directly.
Let $T_0=\varnothing$ and let $T_k$ exhaust all rational
primes, adding one prime at a time.  Define
\[
 G_k=\mathbf R^\times\times
         \prod_{p\in T_k}\mathbf Q_p^\times\times
         \prod_{p\notin T_k}\mathbf Z_p^\times.
 \tag{M37.18}
\]
These are open subgroups; adding one prime has quotient
$\mathbf Z$ by its valuation.  The standard section
$\tau(t)$ and every $g\in E_{\mathbf Q}^1=
\{1\}\times\widehat{\mathbf Z}^\times$ belong to $G_0$.
For \emph{any} single nonnegative $b$ satisfying the
global $\mathbf Q^\times$ normalization and the local
$\{1,-1\}$ normalization on $G_0$, (M37.6) gives
\[
 (\alpha_{k+1}-\alpha_k)(\xi g\tau(t))
                  B_b(\tau(t),\xi g\tau(t))=0
 \quad(\xi\in\mathbf Q^\times,\ g\in E_{\mathbf Q}^1,\ t>0).
 \tag{M37.19}
\]
This is pointwise, before integrating or taking any
kernel limit.  Multiplying by an additional radial
cutoff or by a pulled-back measurable kernel does
not change the vanishing wherever that kernel is
defined.  Thus no such single compatible indentation
can turn the source's external shell in (11.3.2.1)
into its nonzero local Poisson value, exhibited by
the explicit test (M36.7)--(M36.8).

The dependency distinction is precise.  The proof of
Theorem 11.3.2.1 on original p.~495 actually cites
Proposition 12.2.2.2 for its indented-kernel trace
interpretation, and on p.~496 claims an approximation
by the local orbital finite-rank kernels.  It does
\emph{not} there explicitly cite Proposition 12.2.2.4.
The latter is the later available construction of a
single indentation compatible with the whole local
filtration.  MT-M37 proves that substituting that
compatible construction for the invalid literal
section average of MT-M36 cannot supply the asserted
external-shell trace.  It does not invent a source
citation or refute the zero-tail statement itself.

\paragraph{What a varying-level repair must address.}
The family $b_t$ of MT-M36 drops the local normalization
at the initial subgroup and can redistribute mass to
the $p$ and $p^{-1}$ translates; its defect is exactly
$\delta=t$ on the standard domain, and it attains the
lower bound $2t(1-t)$ in (M37.10).  Thus there is no
obstruction here to all nonnegative global indentations.
The obstruction is to imposing the two normalizations
on the \emph{same} function.

A level-dependent family could instead normalize only
at its current subgroup and need not have zero defect
at the original $G_0$.  Such a replacement must specify
the resulting orbital weights and justify its kernel
and trace limits; (M37.14)--(M37.16) give concrete
tests for any claim that the initial defect disappears.
No boundedness, $L^1$ convergence, or invariant orbital
allocation is inferred from the existence of the
individual indentations.  The Hecke-GRH endpoint
$M_H$ remains conditional.

\paragraph{Audit scope.}
Original pp.~512--513 were freshly rendered and read
in full for the hypotheses and extension-by-zero
construction.  Original pp.~495, 505 and 508 were
visually checked in MT-M36, and the complete supplied
proof text on pp.~495--496 was read as a dependency
map.  The new support and defect statements above
are fully proved here and require no GRH input.


\subsection*{MT-S506. Dirichlet support of possible poles for every $S_5$ extension}

Let $L/\mathbb Q$ be a Galois extension with group $G=S_5$,
let $E=L^{A_5}$, and let $\eta$ be its quadratic character.
No restriction on complex conjugation or on ramification at $5$
is imposed in this subsection.  Use the irreducible representations
\[
 \begin{gathered}
 \mathbf1,\quad\eta,\quad V,\quad V\eta,\quad W,\quad W\eta,\quad X,\\
 V=\mathbf C^5/\mathbf C\mathbf1,\quad
 W=\operatorname{Sym}^2V-\mathbf1-V,\quad X=\bigwedge^2V.
 \end{gathered}
 \tag{S506.1}
\]
The notation for $W$ first denotes a virtual representation;
the character calculation below proves that it is an actual
irreducible representation of dimension $5$.  These are
Calegari's $\rho_4,\rho_5,\rho_6$, respectively, and $X\eta=X$.
All $L$-functions here include the finite factors at ramified primes.

For $z$ in the open critical strip $0<\Re z<1$, put
\[
 (a,b,c,d,x,y,e)=\operatorname{ord}_{s=z}
 (\zeta_{\mathbb Q},L_\eta,L_V,L_{V\eta},L_W,L_{W\eta},L_X).
 \tag{S506.2}
\]
An order is positive at a zero and negative at a pole.
Unconditionally $L_X$ is an entire Hecke $L$-function of a
nontrivial finite-order character over a degree-six field.
Consequently $a,b,e\ge0$, and the pole orders satisfy
\[
 \begin{aligned}
 \max(0,-c)&\le\min(a,e),&
 \max(0,-d)&\le\min(b,e),\\
 \max(0,-x)&\le b,&
 \max(0,-y)&\le a.
 \end{aligned}
 \tag{S506.3}
\]
In particular every possible pole of an irreducible Artin factor
in this strip lies in
\[
          Z(\zeta_{\mathbb Q})\ \cup\ Z(L_\eta).
 \tag{S506.4}
\]
Under the designated input $A_D$, every such pole therefore lies
on $\Re z=1/2$.  This conclusion uses only rational Dirichlet GRH.
It neither asserts that these poles exist nor proves their absence,
and it gives no general zero-location theorem for the other factors.

\begin{proof}
The complete character table, in the order (S506.1), is
\[
\begin{array}{c|r|rrrrrrr}
\text{cycle type}&\text{class size}&\mathbf1&\eta&V&V\eta&W&W\eta&X\\\hline
1^5&1&1&1&4&4&5&5&6\\
2\,1^3&10&1&-1&2&-2&1&-1&0\\
2^2\,1&15&1&1&0&0&1&1&-2\\
3\,1^2&20&1&1&1&1&-1&-1&0\\
3\,2&20&1&-1&-1&1&1&-1&0\\
4\,1&30&1&-1&0&0&-1&1&0\\
5&24&1&1&-1&-1&0&0&1
\end{array}
 \tag{S506.5}
\]
Indeed $\chi_V(g)=\#\operatorname{Fix}(g)-1$ and
\[
 \chi_W(g)=\frac{\chi_V(g)^2+\chi_V(g^2)}2-1-\chi_V(g),\qquad
 \chi_X(g)=\frac{\chi_V(g)^2-\chi_V(g^2)}2.
\]
These formulas give every displayed entry.  Weighted inner products
of the seven columns are $\delta_{ij}$: multiplication by the
class-size diagonal matrix gives $120I$ for the Gram matrix.
A virtual character of norm one is the positive or negative of
an irreducible character; its positive value at the identity fixes
the positive sign.  Thus all seven columns are actual irreducibles.
Their squared dimensions sum to $120$, so the list is complete.
This also proves $X\eta=X$ directly, since $X$ vanishes on odd
permutations.

Let $J=S_4$ be a point stabilizer and let
$F_{20}=\operatorname{AGL}_1(\mathbb F_5)\subset S_5$ be the
normalizer of a $5$-cycle.  Write an affine element as
$t\mapsto ut+v$, and let $\tau(u)$ be a primitive quartic character
of $\mathbb F_5^\times$, inflated to $F_{20}$.  The sign restricted
to $F_{20}$ is $\tau^2$.  Let
$T=S_2\times S_3$ be a two-subset stabilizer and let
$\alpha$ be sign on its $S_2$ factor and trivial on its $S_3$ factor.
Frobenius reciprocity and (S506.5) give
\[
\begin{array}{c|l}
\text{induced character}&\text{decomposition in }G\\\hline
\operatorname{Ind}_{J}^{G}\mathbf1&\mathbf1+V\\
\operatorname{Ind}_{J}^{G}(\eta|_J)&\eta+V\eta\\
\operatorname{Ind}_{F_{20}}^{G}\mathbf1&\mathbf1+W\eta\\
\operatorname{Ind}_{F_{20}}^{G}(\eta|_{F_{20}})&\eta+W\\
\operatorname{Ind}_{F_{20}}^{G}\tau&X\\
\operatorname{Ind}_{T}^{G}\alpha&V+X\\
\operatorname{Ind}_{T}^{G}(\alpha\eta|_T)&V\eta+X.
\end{array}
 \tag{S506.6}
\]
Here is an explicit check of the potentially sensitive quartic
identity.  The affine group has the identity, four nonidentity
translations of cycle type $5$, five elements of type $2^2\,1$,
and ten of type $4\,1$.  The first two sets have multiplier $1$,
the third has multiplier $-1$, and the last two multiplier classes
have $\tau$-values $i$ and $-i$.  For $X$ the weighted inner-product
numerator is $6+4+10=20$.  For every other column the numerator
is zero, because the two quartic multiplier classes cancel and
the remaining two multiplier contributions cancel.  Division by
$|F_{20}|=20$ proves the identity.  The other rows follow by the
same finite averaging formula
\[
 \bigl\langle\operatorname{Ind}_{H}^{G}\lambda,R\bigr\rangle_G
       =\frac1{|H|}\sum_{h\in H}\overline{\lambda(h)}\chi_R(h).
 \tag{S506.7}
\]
For the two $T$ rows, $\alpha$ and $\alpha\eta$ are respectively
sign on the $S_2$ and on the $S_3$ factor, which fixes the twist
convention.

For a linear character $\lambda$ of any subgroup $H\subset G$,
Artin induction and class field theory identify
$L(s,\operatorname{Ind}_H^G\lambda)$ with its Hecke $L$-function
over $L^H$.  It is holomorphic in the open critical strip,
including when $\lambda$ is trivial, in which case it is a
Dedekind zeta function.  For $\tau\ne\mathbf1$ it is entire.
Taking orders in (S506.6) yields
\[
 \begin{gathered}
 a+c\ge0,\quad b+d\ge0,\quad a+y\ge0,\quad b+x\ge0,\\
 e\ge0,\quad c+e\ge0,\quad d+e\ge0.
 \end{gathered}
\]
The orders exist as integers because Brauer induction gives
meromorphic continuation of every Artin factor.  Since $a,b,e$
are nonnegative, these inequalities prove (S506.3)--(S506.4).
The last assertion is precisely the Dirichlet zero-location
input $A_D$ applied to the trivial and the quadratic characters.
\end{proof}

The classical all-prime induction and multiplication identities
are Prasad--Yogananda, \emph{A Report on Artin's holomorphy
conjecture}, Section~2.1, Properties~1--2, PDF p.~3;
Hecke entireness is Section~3.1, pp.~4--5; meromorphicity is
Theorem~3.3, p.~5.  The character labels agree with Calegari,
\emph{The Artin conjecture for some $S_5$-extensions},
author-hosted publisher PDF p.~4, Section~2 character table.
The character calculations here are proved explicitly; the
classical analytic inputs are cited, not re-proved.

\subsection*{MT-S507. The complete monomial cone and the remaining five-dimensional gate}

Retain (S506.1)--(S506.2), and assume additionally that
$L_V$ and $L_{V\eta}$ are holomorphic at $z$, that is, $c,d\ge0$.
This is only a pointwise hypothesis for the following cone theorem.
It holds at every point for the Calegari family of MT-S501:
complex conjugation has cycle type $(12)(34)$, $5$ is unramified,
and $\operatorname{Frob}_5$ has the same cycle type.
Calegari's Theorem~1.1, author-hosted publisher PDF p.~2, supplies
this four-dimensional holomorphy unconditionally.  No GRH input
is needed for the cone theorem or its numerical bounds.

At $z$, all tests
\[
 \operatorname{ord}_z L\bigl(s,\operatorname{Ind}_H^G\lambda\bigr)
             \ge0\qquad(H\subset G,\ \dim\lambda=1)
 \tag{S507.1}
\]
are equivalent, after imposing $c,d\ge0$, to the following twelve
linear inequalities:
\[
 \begin{gathered}
      a,b,c,d,e\ge0,\\
 g_1=x+y\ge0,\quad g_2=b+x\ge0,\quad g_3=a+y\ge0,\\
 g_4=c+x+e\ge0,\quad g_5=d+y+e\ge0,\\
 g_6=a+c+x\ge0,\quad g_7=b+d+y\ge0.
 \end{gathered}
 \tag{S507.2}
\]
The assertion is an exact description of this cone, not a claim
that all of its integer points occur as arithmetic order vectors.
In particular $c,d\ge0$ are extra holomorphy inputs, not deductions
from the monomial tests alone.

Define the relative and total Dedekind orders
\[
 q=\operatorname{ord}_z\frac{\zeta_L}{\zeta_E}
       =4(c+d)+5(x+y)+6e,
 \qquad
 r=\operatorname{ord}_z\zeta_L=a+b+q.
 \tag{S507.3}
\]
If $L_W$ has a pole of order $v>0$ at $z$, then
\[
 \begin{gathered}
 x=-v,\quad y\ge v,\quad b\ge v,\quad c+e\ge v,
       \quad a+c\ge v,\\
 v\le\min\{b,c+e,a+c\},\qquad q\ge4v,\qquad r\ge5v.
 \end{gathered}
 \tag{S507.4}
\]
For a pole of $L_{W\eta}$, interchange
$(a,b),(c,d),(x,y)$.  Both five-dimensional factors cannot have
a pole at the same point.  In particular $q\le3$ implies that
all irreducible factors are holomorphic at $z$, and every possible
pole order is at most $\lfloor q/4\rfloor$.
In the branch $c=0$ of (S507.4),
\[
        a,b,e,y\ge v,\qquad q\ge6v,\qquad r\ge8v.
 \tag{S507.5}
\]
The analogous assertion with $d=0$ holds for $W\eta$.
All these numerical constants are sharp for the exact cone
(S507.2), including the extra branch condition in (S507.5).

Let $Z(f)$ denote its zero set within the open critical strip.
For this holomorphy family the corresponding precise necessary
simultaneous-zero conditions are
\[
\begin{aligned}
 \operatorname{Pole}(L_W)&\subseteq
 Z(L_\eta)\cap\bigl(Z(L_V)\cup
             (Z(\zeta_{\mathbb Q})\cap Z(L_X))\bigr),\\
 \operatorname{Pole}(L_{W\eta})&\subseteq
 Z(\zeta_{\mathbb Q})\cap\bigl(Z(L_{V\eta})\cup
             (Z(L_\eta)\cap Z(L_X))\bigr).
\end{aligned}
 \tag{S507.6}
\]
Under $A_D$ these possible poles lie on the critical line by
MT-S506.  Under $M_H$ and the Calegari local hypotheses they
must additionally be nonreal, by MT-S505.  Neither hypothesis
by itself supplies the missing noncoincidence statement in
(S507.6), and no simplicity of any zero is asserted here.

\begin{proof}[Proof of the exact cone]
Besides the rows already established in (S506.6), use the
subgroup $A_4$ fixing the fifth letter and a nontrivial
character $\nu$ of its quotient $A_4/V_4=C_3$; the two-subset
stabilizer $T$; and
\[
  D_8=\langle r_0,s_0\rangle\subset S_4,\qquad
  r_0=(1234),\quad s_0=(13),\qquad
  \lambda(r_0)=-1,\quad\lambda(s_0)=1.
\]
These values respect $r_0^4=s_0^2=1$ and
$s_0r_0s_0=r_0^{-1}$, so they define a linear character.
Formula (S506.7) and the table give
\[
 \begin{gathered}
 \operatorname{Ind}_{A_4}^{G}\nu=W+W\eta,\qquad
 \operatorname{Ind}_{D_8}^{G}\lambda=V+W+X,\\
 \operatorname{Ind}_{T}^{G}\mathbf1=\mathbf1+V+W.
 \end{gathered}
 \tag{S507.7}
\]
Twisting by $\eta$ gives the other two needed rows, since
$X\eta=X$.  For the $A_4$ identity, its identity and three
involutions contribute $5+3$ to each rank-five inner product;
the eight $3$-cycles split into two sets of four with character
values $\omega,\omega^2$, so their weighted contribution is
$-4(\omega+\omega^2)=4$.  Division by $12$ gives one.
The other inner products vanish.  For the $D_8$ identity,
its elements have cycle types $1$, two $4$-cycles, three double
transpositions, and two transpositions.  Their $\lambda$-weights
sum, within those types, to $1,-2,-1,2$, respectively.
Against the seven columns this yields
$(0,0,1,0,1,0,1)$.  For $T$, the permutation action on the ten
two-subsets has character $1+\chi_V+\chi_W$ by direct counting
of invariant two-subsets.  This proves every row of (S507.7).
Thus all seven $g_i$ are monomial tests.  The tests with
$H=G$ give $a,b\ge0$, and the quartic row gives $e\ge0$.
Together with the stated $c,d\ge0$ premise this proves the
necessity of (S507.2).

For completeness, the next two tables give every distinct induced
linear-character multiplicity vector, together with a nonnegative
integer decomposition of its order functional into (S507.2).
The columns labeled $a,\ldots,e$ are the coefficients of that
functional, not values of the order variables.  Thus each last
entry can be verified by substituting the definitions of $g_i$.
\[
\begin{array}{r|rrrrrrr|l}
 i&a&b&c&d&x&y&e&\text{nonnegative decomposition}\\\hline
1 & 0 & 0 & 0 & 0 & 0 & 0 & 1 & e \\
2 & 0 & 0 & 0 & 0 & 0 & 0 & 2 & 2e \\
3 & 0 & 0 & 0 & 0 & 1 & 1 & 0 & g_1 \\
4 & 0 & 0 & 0 & 1 & 0 & 0 & 1 & d+e \\
5 & 0 & 0 & 0 & 1 & 0 & 1 & 1 & g_5 \\
6 & 0 & 0 & 0 & 1 & 1 & 1 & 1 & g_1+d+e \\
7 & 0 & 0 & 1 & 0 & 0 & 0 & 1 & c+e \\
8 & 0 & 0 & 1 & 0 & 1 & 0 & 1 & g_4 \\
9 & 0 & 0 & 1 & 0 & 1 & 1 & 1 & g_1+c+e \\
10 & 0 & 0 & 1 & 1 & 0 & 0 & 2 & c+d+2e \\
11 & 0 & 0 & 1 & 1 & 1 & 1 & 1 & g_1+c+d+e \\
12 & 0 & 0 & 1 & 1 & 1 & 1 & 2 & g_4+g_5 \\
13 & 0 & 0 & 1 & 1 & 2 & 2 & 2 & g_1+g_4+g_5 \\
14 & 0 & 0 & 2 & 2 & 2 & 2 & 4 & 2g_4+2g_5 \\
15 & 0 & 1 & 0 & 0 & 0 & 0 & 0 & b \\
16 & 0 & 1 & 0 & 0 & 1 & 0 & 0 & g_2 \\
17 & 0 & 1 & 0 & 1 & 0 & 0 & 0 & b+d \\
18 & 0 & 1 & 0 & 1 & 0 & 1 & 0 & g_7 \\
19 & 0 & 1 & 0 & 1 & 1 & 1 & 0 & g_1+b+d \\
20 & 0 & 1 & 0 & 2 & 0 & 1 & 1 & g_7+d+e \\
21 & 0 & 1 & 0 & 2 & 1 & 2 & 1 & g_1+g_7+d+e \\
22 & 0 & 1 & 1 & 1 & 0 & 1 & 1 & g_7+c+e \\
\end{array}
\]

\[
\begin{array}{r|rrrrrrr|l}
 i&a&b&c&d&x&y&e&\text{nonnegative decomposition}\\\hline
23 & 0 & 1 & 1 & 1 & 2 & 1 & 1 & g_1+g_4+b+d \\
24 & 0 & 1 & 1 & 3 & 2 & 3 & 3 & g_1+g_4+g_5+g_7+d+e \\
25 & 1 & 0 & 0 & 0 & 0 & 0 & 0 & a \\
26 & 1 & 0 & 0 & 0 & 0 & 1 & 0 & g_3 \\
27 & 1 & 0 & 1 & 0 & 0 & 0 & 0 & a+c \\
28 & 1 & 0 & 1 & 0 & 1 & 0 & 0 & g_6 \\
29 & 1 & 0 & 1 & 0 & 1 & 1 & 0 & g_1+a+c \\
30 & 1 & 0 & 1 & 1 & 1 & 0 & 1 & g_6+d+e \\
31 & 1 & 0 & 1 & 1 & 1 & 2 & 1 & g_1+g_5+a+c \\
32 & 1 & 0 & 2 & 0 & 1 & 0 & 1 & g_6+c+e \\
33 & 1 & 0 & 2 & 0 & 2 & 1 & 1 & g_1+g_6+c+e \\
34 & 1 & 0 & 3 & 1 & 3 & 2 & 3 & g_1+g_4+g_5+g_6+c+e \\
35 & 1 & 1 & 0 & 0 & 0 & 0 & 0 & a+b \\
36 & 1 & 1 & 0 & 0 & 1 & 1 & 0 & g_2+g_3 \\
37 & 1 & 1 & 0 & 0 & 1 & 1 & 2 & g_2+g_3+2e \\
38 & 1 & 1 & 1 & 1 & 0 & 0 & 0 & a+b+c+d \\
39 & 1 & 1 & 1 & 1 & 1 & 1 & 0 & g_6+g_7 \\
40 & 1 & 1 & 1 & 1 & 2 & 2 & 0 & g_1+g_6+g_7 \\
41 & 1 & 1 & 2 & 2 & 1 & 1 & 2 & g_6+g_7+c+d+2e \\
42 & 1 & 1 & 2 & 2 & 3 & 3 & 2 & g_1+g_4+g_5+g_6+g_7 \\
43 & 1 & 1 & 4 & 4 & 5 & 5 & 6 & g_1+3g_4+3g_5+g_6+g_7 \\
\end{array}
\]

Here is an exhaustive finite algorithm for these tables, with no
numerical optimization premise.  List all $120$ permutations of
five letters and their exact composition table.  Begin with the
trivial subgroup; repeatedly adjoin every individual permutation
to every subgroup found, compute the resulting generated subgroup,
and discard duplicates.  This finds every subgroup: given a
subgroup $H$, successively adjoin elements of $H$ until its
finite generating list has been exhausted.  Each such step appears
in this search.  It yields $156$ subgroups and $19$ conjugacy
classes; conjugate subgroups give the same induced patterns.
For each representative, retain a generating list.

The exponent of $S_5$ divides $60$, so every linear character of
one of its subgroups has values among the $60$th roots of unity.
Assign to each generator all phases $j\in\mathbb Z/60$ allowed
by its order.  Propagate those phases through its Cayley graph,
reject inconsistent assignments, and check
$j(hh')=j(h)+j(h')$ for all pairs.  This both finds every linear
character and verifies that each retained assignment is one.
Compute (S506.7) in the exact ring
$\mathbb Q[Z]/(\Phi_{60}(Z))$, using
$\overline{\lambda(h)}=Z^{-j(h)}$; reduce every sum modulo
$\Phi_{60}$ and divide by $|H|$.  The resulting multiplicities
are the nonnegative integers listed above.  As checks, the number
of characters found for each $H$ is $|H/[H,H]|$, and the
weighted dimension of every row is $120/|H|$.
Discarding repeated vectors gives exactly the $43$ displayed rows.
All steps use finite permutations, integer arithmetic and exact
cyclotomic reduction.

The full executable enumeration and its output are
\path{work/checks/s5_hecke_cone.py} and
\path{s5_hecke_cone.json}; their exploratory floating-point
optimizer is not used in this proof.  The separate exact verifier
\path{work/checks/s5_cone_decomposition_verification.py}
and its JSON output reconstruct all $43$ rows as nonnegative
integer combinations of the twelve generators and verify the
sharpness witnesses below.  Its input SHA-256 is
\begin{center}\small\ttfamily
0b1ab9c1b47311477e3e113cced1e568a4c6fc417c8077875e87e7d8637ea917
\end{center}
The exhaustive permutation/character portion was re-executed
without its optimizer, and the displayed decompositions were
checked in integer arithmetic.  This is a transparent finite
computation with its algorithm and complete output supplied;
it is not a proof-assistant certification.

Every order vector satisfying (S507.2) pairs nonnegatively
with every displayed row.  Exhaustiveness therefore proves
sufficiency for (S507.1).  This establishes the exact cone.
\end{proof}

\begin{proof}[Proof of the pole bounds, sharpness and support condition]
The regular character of $G$ has multiplicities equal to
irreducible dimensions, whereas the character for the quadratic
field $E$ is $\mathbf1+\eta$.  All-prime Artin factorization
therefore gives (S507.3), including at ramified primes.
If $x=-v<0$, the inequalities $g_1,g_2,g_4,g_6\ge0$
give the first line of (S507.4).  Since $d,e,x+y\ge0$,
\[
 q=4(c+e)+4d+5(x+y)+2e\ge4(c+e)\ge4v,
 \qquad r=a+b+q\ge5v.
\]
If $c=0$, $g_4,g_6\ge0$ force $e,a\ge v$, so
$q\ge6e\ge6v$ and $r\ge a+b+q\ge8v$.
Twisting interchanges $a,b$, $c,d$ and $x,y$, giving every
symmetric assertion.  Since $g_1\ge0$, $x,y$ cannot both be
negative.  The other five irreducible factors are already
holomorphic at $z$ under the stated hypotheses.  A remaining
pole has integral order $v\ge1$ and $q\ge4v$; this proves both
the $q\le3$ holomorphy criterion and the floor bound.

The integer vector
\[
        (a,b,c,d,x,y,e)=(0,1,1,0,-1,1,0)
 \tag{S507.8}
\]
satisfies (S507.2): the seven $g_i$ equal
$(0,0,1,0,1,0,2)$ and the five required coordinates are
nonnegative.  It has a pole order $v=1$, $q=4$ and $r=5$.
Its sign twist provides the symmetric witness.  For the
$c=0$ branch, the vector
\[
        (a,b,c,d,x,y,e)=(1,1,0,0,-1,1,1)
 \tag{S507.9}
\]
has $g_i=(0,0,2,0,2,0,2)$, $q=6$ and $r=8$.
Multiplying these vectors by any positive integer gives every
pole order $v$.  By the exact cone, all $43$ monomial tests are
nonnegative at these witnesses.  This proves sharpness within
the specified positivity information.  It makes no claim that
these integer vectors arise from number fields or analytic
$L$-functions.

Finally $x<0$ forces $b>0$.  If $c>0$ the first branch of
(S507.6) holds.  If $c=0$, then $a,e>0$, which is exactly
its other branch.  Since $c\ge0$, these cases exhaust all
possibilities.  The sign twist proves the second inclusion.
Under $A_D$, (S506.4) confines the poles to the critical line.
Under $M_H$ and the Calegari family hypotheses, MT-S505 excludes
every real point of the open critical strip, including the
centre.  This leaves only the nonreal points described above.
\end{proof}

\paragraph{Input separation and the geometric gate.}
MT-S506 is a new reduction using classical character theory,
Artin induction and Hecke analyticity; only its last pole-location
assertion invokes $A_D$.  MT-S507 additionally uses the
four-dimensional holomorphy premise, supplied in the indicated
family by Calegari's theorem.  Its stronger exclusion of real
poles uses the separately proved, $M_H$-conditional MT-S505.
The finite cone does not establish noncoincidence, pole absence,
automorphy of the rank-five factors, or a new proof of GRH.

On the arithmetic base with the ramified primes removed, these
representations are the Artin motives of finite covers and their
direct summands.  Induction corresponds to finite pushforward
from a cover; (S507.1) records positivity of the zero orders
for all one-dimensional induced summands.  The witness (S507.8)
shows precisely why that information, even supplemented by the
rank-four holomorphy input, does not force the order divisor of
each irreducible summand to be effective.  The additional
all-prime induction identities are compatible with inertia
invariants at the removed primes.  Thus (S507.6), with the order
penalties (S507.4)--(S507.5), is a smaller explicit interface for
the Artin-positivity gate, rather than an identification of
zero-independent positive cohomology.  The attached run is
\path{GRH_NIGHT/complexity}; the inherited interfaces are
MT-ART07 and MT-S501/MT-S504/MT-S505, while the cone certificate
and the deductions in MT-S506/MT-S507 are new in this run.

\subsection*{MT-S508. One quadratic Hecke input and the exact automorphic central-order gate}

Retain the Calegari family of MT-S501: $L/\mathbb Q$ has group
$S_5$, complex conjugation has type $(12)(34)$, and the extension
is unramified at $5$ with Frobenius of the same type.  Let
$E=L^{A_5}$, let $F=L^{F_{20}}$ have degree $6$, and put $H=EF$.
Write $\eta$ for the quadratic character of $E/\mathbb Q$ and
$\eta_F$ for its restriction to $\operatorname{Gal}(L/F)$.
Then $H/F$ is quadratic, and its nontrivial Hecke character is
$\eta_F$.  Define the \emph{finite} $L$-functions
\[
 C(s)=L_{\mathbb Q}(s,\eta),\qquad
 B(s)=L_F(s,\eta_F)=\frac{\zeta_H(s)}{\zeta_F(s)}.
 \tag{S508.1}
\]
Both $B$ and $C$ are entire.  The all-prime induction identity
of MT-S506 gives
\[
             B(s)=C(s)L_W(s),\qquad
             \zeta_F(s)=\zeta_{\mathbb Q}(s)L_{W\eta}(s).
 \tag{S508.2}
\]
Let $\varpi$ be the particular cuspidal representation of
$\mathrm{GL}_5/\mathbb Q$ supplied by Calegari's construction,
and put $A(s)=L_f(s,\varpi)$.  This is its standard finite
$L$-function, which is entire.  Fix this comparison representation;
the statements below concern its exact transfer, not an arbitrary
choice of an automorphic representation of the same degree.
Here matching at every place means agreement of the standard
local $L$-factors, including the archimedean factor, in the
precise sense of the cited comparison theorem.

The only additional zero-location input needed here is
\[
  (\mathcal R_B)\qquad
  B(\sigma)\ne0\quad
  \bigl(0<\sigma<1,\ \sigma\ne\tfrac12,\ \sigma\in\mathbb R\bigr).
 \tag{S508.3}
\]
It concerns one quadratic Hecke character over one degree-six
field and only its real zeros.  GRH for this single Hecke
$L$-function implies $(\mathcal R_B)$; universal $M_H$ is a
sufficient input but is not required.  No Dirichlet GRH assumption
is needed for the central-order and automorphy assertions below.

Put
\[
 \beta=\operatorname{ord}_{1/2} B,\qquad
 b=\operatorname{ord}_{1/2} C,\qquad
 n=\operatorname{ord}_{1/2} A.
 \tag{S508.4}
\]
These are nonnegative integers.  Under $(\mathcal R_B)$ one has
\[
             \beta\ge b+n,
       \qquad \operatorname{ord}_{1/2}L_W=\beta-b\ge n\ge0.
 \tag{S508.5}
\]
Moreover the following are equivalent:
\[
 \begin{gathered}
       \beta\le b+n,\\
       \beta=b+n,\\
       \text{Calegari's exceptional set }\Omega\text{ is finite},\\
       \text{the specified }\varpi\text{ matches }W
             \text{ at every place}.
 \end{gathered}
 \tag{S508.6}
\]
Whenever these conditions hold, $\varpi\otimes\eta$ also matches
$W\eta$ at every place.  Thus every irreducible Artin factor of
$L/\mathbb Q$ is automorphic and every nontrivial one is entire.
If the comparison has not closed, its remaining gate is the
strict positive discrepancy
\[
                        \beta-b-n\ge1.
 \tag{S508.7}
\]
The older sufficient inequality $\beta\le b$ is a special case
of (S508.6).  The new condition retains possible central zeros
of the known automorphic numerator.

\paragraph{The exact classical comparison and its normalization.}
The classical construction and local comparison are Calegari,
\emph{The Artin conjecture for some $S_5$-extensions},
author-hosted publisher PDF: Theorem~1.1 p.~2, construction of
$\varpi$ p.~10, Theorem~4.7 p.~12, Lemma~4.8 pp.~13--14,
Landau Lemma~4.9 p.~14, Lemma~4.10 pp.~14--15, and the final
Theorem~1.2 proof p.~15.  This is the publisher numbering;
the last exact-comparison lemma is numbered 4.9 in the
arXiv v1 locators used in MT-S504.

We use the actual archimedean-completed quotient
\[
 R(s)=\frac{G_{\varpi}(s)A(s)}{G_W(s)L_W(s)},
 \tag{S508.8}
\]
where $G$ denotes the archimedean factor.  Including the usual
positive conductor powers would only multiply this by a
nowhere-zero entire function and would not change any order
or real-pole argument.  With the convention that a finite local
Euler factor is $P_p(p^{-s})^{-1}$, its correct factorization is
\[
 \begin{aligned}
 R(s)&=R_{\mathcal B}(s)
      \prod_{p\in\Omega}
             \left(\frac{1+p^{-s}}{1-p^{-s}}\right)^2,\\
 R_{\mathcal B}(s)&=
       \frac{G_{\varpi}(s)}{G_W(s)}
       \prod_{p\in\mathcal B}
             \frac{P_{W,p}(p^{-s})}{P_{\varpi,p}(p^{-s})}.
 \end{aligned}
 \tag{S508.9}
\]
Here $\mathcal B$ is the finite set of exceptional bad primes,
chosen disjoint from $\Omega$.  The source's displayed
$\mu_{\mathcal B}$ on publisher p.~13 reverses the two finite
polynomials relative to its subsequent reciprocal-Euler-factor
convention.  Equation (S508.9) fixes the direction directly from
the definition (S508.8); the source's holomorphic, nonzero
real-interval property is unaffected by this finite inversion.
In particular at $p\in\Omega$ its explicit polynomials are
\[
 P_{\varpi,p}(T)=(1+T^2)(1-T)^3,\qquad
 P_{W,p}(T)=(1+T^2)(1-T)(1+T)^2,
\]
which give exactly the displayed positive Euler quotient.

The cited comparison supplies two precise implications:
if $\Omega$ is infinite then $R$ has a pole at some real
$\sigma\in(0,1)$; if $\Omega$ is finite then the local factors
agree at all places, including the bad and archimedean ones.
The first uses its nonnegative Euler-series coefficients and
Landau's theorem; the second uses the cited finite-place
comparison with functional equations and twists.  The source's
final proof applies the same comparison to the quadratic twist.
We retain these classical results as inputs, not a new proof
of their automorphic construction.

The ratio $G_{\varpi}/G_W$ is holomorphic and nonzero on $(0,1)$.
For clarity, the Artin factor can be checked directly here.
Complex conjugation has trace $1$ on $W$, so its $+1,-1$
eigenspaces have dimensions $3,2$ and
\[
 G_W(s)=\Gamma_{\mathbb R}(s)^3\Gamma_{\mathbb R}(s+1)^2
       =\Gamma_{\mathbb R}(s)\Gamma_{\mathbb C}(s)^2.
 \tag{S508.10}
\]
Also $\eta$ is even, so $G_C(s)=\Gamma_{\mathbb R}(s)$.
The field $F$ has signature $(2,2)$, since the degree-six
permutation character $\mathbf1+W\eta$ takes value $2$ on
complex conjugation.  The real places split in $H/F$, because
$E$ is real.  Thus
$G_B(s)=\Gamma_{\mathbb R}(s)^2\Gamma_{\mathbb C}(s)^2
=G_C(s)G_W(s)$, consistently with (S508.2).
Here $\Gamma_{\mathbb R}(s)=\pi^{-s/2}\Gamma(s/2)$ and
$\Gamma_{\mathbb C}(s)=2(2\pi)^{-s}\Gamma(s)$.
All these factors are finite and nonzero at every real
$0<s<1$.  The source verifies that $G_{\varpi}$ likewise
uses standard real/complex gamma factors there.

\begin{proof}[Proof of the one-Hecke theorem]
First, the subgroup $F_{20}$ contains odd permutations, so
its intersection with $A_5$ has index two.  It follows that
$H=EF$ is the quadratic extension of $F$ cut out by the
nontrivial restriction $\eta_F$.  Artin induction, including
ramified factors, gives (S508.1)--(S508.2).

Substituting (S508.2) into the actual quotient yields
\[
          R(s)=U(s)\,\frac{A(s)C(s)}{B(s)},\qquad
          U(s)=\frac{G_{\varpi}(s)}{G_W(s)}.
 \tag{S508.11}
\]
For any real $0<\sigma<1$, the function $U$ is holomorphic
and nonzero.  Therefore
\[
 \operatorname{ord}_{\sigma}R
   =\operatorname{ord}_{\sigma}A
     +\operatorname{ord}_{\sigma}C
     -\operatorname{ord}_{\sigma}B.
 \tag{S508.12}
\]
This formula explicitly accounts for zeros of the automorphic
numerator; discarding them would give only the older sufficient
condition.

Under $(\mathcal R_B)$, the denominator in (S508.11) is nonzero
at every real point of $(0,1)$ other than the centre.
The other factors are holomorphic there, even if $C$ has a
zero.  Thus the only possible real pole of $R$ is at $1/2$,
where its order is $n+b-\beta$.
If $\Omega$ is infinite, the classical dichotomy forces a pole
there, hence $\beta>b+n$.  If $\Omega$ is finite, the exact
comparison gives $R\equiv1$ in the archimedean normalization
of (S508.8), hence $\beta=b+n$.  These two cases exhaust the
possibilities and prove (S508.5)--(S508.7).

The reverse implication in the last equivalence is also literal:
if all local factors of the specified comparison agree, then
no prime can have the two distinct polynomials displayed above,
so its exceptional set is empty.  Conversely, finite $\Omega$
forces exact local equality by the cited comparison.
For the quadratic twist, equality outside finitely many places
persists after twisting; the same exact-comparison argument,
as explicitly applied on source p.~15, identifies
$\varpi\otimes\eta$ with $W\eta$ at every place.  No second
Hecke zero-location hypothesis is needed for this automorphy
conclusion.  The dimensions $4$ and $6$ are already automorphic
in the specified family by Theorem~1.1, and the two linear
characters are classical.  The claimed holomorphy follows.
\end{proof}

\paragraph{The positive Euler quotient controls a whole half-plane.}
The same single real-zero hypothesis $(\mathcal R_B)$ implies
that $R$ is holomorphic and nonzero on $\Re s>1/2$.
In that half-plane it gives the exact divisor relation
\[
        \operatorname{ord}_z B
             =\operatorname{ord}_z A+\operatorname{ord}_z C.
 \tag{S508.13}
\]
In particular $L_W$ is holomorphic there and has the same
zero orders as $A$.  The poles of $L_W$ in the open critical
strip can only occur on its critical line, away from the real
axis.  This assertion requires neither $A_D$ nor full GRH for $B$.

\begin{proof}
Put
\[
 D_\Omega(s)=\prod_{p\in\Omega}
             \left(\frac{1+p^{-s}}{1-p^{-s}}\right)^2.
\]
This is the positive-coefficient Dirichlet series in the proof
of source Lemmas~4.8--4.9.  It is meromorphic on the plane
because $D_\Omega=R/R_{\mathcal B}$.  If $\Omega$ is infinite,
that proof identifies an abscissa of absolute convergence
$\gamma\in(0,1)$ and a pole at $\gamma$; the convergence at $1$
uses the stronger prime-sum bound proved in Theorem~4.7, not
Dirichlet density zero alone.  Since $R_{\mathcal B}$ is nonzero
and holomorphic at every real point of $(0,1)$, this is also a
pole of $R$.  Under $(\mathcal R_B)$, (S508.11) shows that this
pole can only be at $1/2$.  Hence $\gamma=1/2$.
If $\Omega$ is finite, the finite Euler product instead
converges for $\Re s>0$.

For every real $\sigma>1/2$ the nonnegative Dirichlet series
therefore converges, and its prime coefficients are $4$.
It follows that $\sum_{p\in\Omega}p^{-\sigma}<\infty$.
The logarithms of the Euler factors consequently converge
absolutely and locally uniformly on $\Re s>1/2$.
Each individual factor is nonzero in $\Re s>0$, so the
infinite product is holomorphic and nonzero on $\Re s>1/2$.
This also covers an empty or finite exceptional set.

The gamma quotient in $R_{\mathcal B}$ is holomorphic and
nonzero on all of $\Re s>0$: its standard gamma factors have
no poles or zeros there.  The finite polynomial factors are
also nonzero there, since their roots have modulus one and
$|p^{-s}|<1$.  Thus $R_{\mathcal B}$, and consequently $R$,
are holomorphic and nonzero on $\Re s>1/2$.
The same is true of $U$.  Taking orders in (S508.11) proves
(S508.13), and $L_W=UA/R$ proves the asserted holomorphy
and equality of zero orders.  The meromorphic Artin functional
equation for the selfdual $W$ reflects poles in
$0<\Re s<1/2$ to poles in $1/2<\Re s<1$;
its gamma and conductor factors are nonzero in this open strip.
Thus it excludes poles in the left half as well.  The centre
is already holomorphic by (S508.5).
\end{proof}

\paragraph{A single Hecke GRH also transfers to the known automorphic numerator.}
Assume now the full GRH for $B$, which implies
$(\mathcal R_B)$.  Then both $A=L_f(s,\varpi)$ and $C=L(s,\eta)$
satisfy GRH, whether or not the central equality (S508.6) holds.
Moreover the meromorphic divisor of $L_W=B/C$ in the open
critical strip is supported on the critical line.  This last
statement allows poles on that line; it does not prove Artin
holomorphy there.

\begin{proof}
For $1/2<\Re z<1$, GRH for $B$ says that
$\operatorname{ord}_z B=0$.  The two terms on the right of
(S508.13) are nonnegative, since $A$ and $C$ are holomorphic.
They are therefore both zero.  Their usual completed functional
equations then exclude zeros in $0<\Re z<1/2$ as well.
For the cuspidal unitary $\varpi$, the functional equation
pairs its standard $L$-function with that of its contragredient,
which is the complex conjugate under $s\mapsto\bar s$;
thus a zero is reflected across $\Re s=1/2$.
All gamma factors are nonzero in the open strip.  This proves
GRH for $A$ and $C$.  The factorization $L_W=B/C$ now proves
its critical-line divisor support.  If the central equality
also holds, the already proved exact transfer makes $L_W$
entire, so the divisor assertion is its ordinary GRH.
\end{proof}

This is specifically a \emph{one-$B$} transfer.  It is not a
claim that an $M_H$-conditional zero-location statement for
$\varpi$ was previously unavailable through its known base-change
product.  No GRH for $L_{W\eta},L_V,L_{V\eta}$ or $L_X$ follows
from the argument just given.  Once exact transfer is proved,
GRH for $\zeta_F$ separately gives GRH for $L_{W\eta}$ from
(S508.2), and universal $M_H$ is one sufficient way to supply
the other zero-location inputs.

\paragraph{The quadratic-twist expression for the same central gate.}
Put
\[
 A_\eta(s)=L_f(s,\varpi\otimes\eta),\qquad
 n_\eta=\operatorname{ord}_{1/2}A_\eta,
 \qquad f=\operatorname{ord}_{1/2}\zeta_F,
 \qquad k=\beta-b-n.
\]
Under $(\mathcal R_B)$ these nonnegative orders satisfy
\[
             k=n_\eta-f,\qquad 0\le k\le n_\eta.
 \tag{S508.14}
\]
Consequently $n_\eta\le f$ is equivalent to the same exact
transfer condition (S508.6).  In particular the additional
nonvanishing $A_\eta(1/2)\ne0$ suffices.  Neither that
nonvanishing nor this inequality is asserted unconditionally.

\begin{proof}
Every prime in $\Omega$ has Frobenius type $4\,1$, so
$\eta(p)=-1$.  Twisting its two displayed local polynomials
replaces $T$ by $-T$ and inverts their quotient.  Hence the
actual twisted completed ratio is
\[
 R_\eta(s)=\frac{G_{\varpi\eta}(s)A_\eta(s)}
                     {G_{W\eta}(s)L_{W\eta}(s)}
          =R_{\mathcal B,\eta}(s)D_\Omega(s)^{-1}.
 \tag{S508.15}
\]
The finite and archimedean factor $R_{\mathcal B,\eta}$ is
holomorphic and nonzero at every real $0<s<1$ before any
finiteness assertion about $\Omega$: $\eta$ restricts trivially
to $E$, so $(\varpi\otimes\eta)_E=\varpi_E$ has the same
known finite-image base change, and the local parameter argument
of source Lemma~4.8 still gives root-of-unity finite Euler
parameters after twisting.  In particular $\eta$ is even,
so it does not change the archimedean factors.  At the centre $R_{\mathcal B}$ and
$R_{\mathcal B,\eta}$ have order zero, and (S508.12) gives
$\operatorname{ord}_{1/2}D_\Omega=-k$.  Thus
$\operatorname{ord}_{1/2}R_\eta=k$.
On the other hand $\zeta_{\mathbb Q}(1/2)\ne0$, so (S508.2)
gives $\operatorname{ord}_{1/2}L_{W\eta}=f$.
Taking orders in the first quotient of (S508.15) gives
$k=n_\eta-f$.  Since $f\ge0$ and $k\ge0$ by (S508.5),
(S508.14) follows.  Equality $k=0$ is exactly (S508.6).
\end{proof}

\paragraph{Real-axis holomorphy needs only the one real-zero input.}
Under $(\mathcal R_B)$ every irreducible Artin factor in this
family is holomorphic at every real point of $(0,1)$.
For $W$, the half-plane argument and its functional equation
exclude every noncentral real pole, while (S508.5) handles
the centre.  For $W\eta$, (S508.2) and the elementary inequality
$\zeta_{\mathbb Q}(\sigma)<0$ for $0<\sigma<1$ prove
holomorphy on that real interval without any GRH input.
The other nontrivial factors are already entire by the
classical theorem.  Thus the earlier $A_D+(\mathcal R_B)$
corollary strengthens to $(\mathcal R_B)$ alone after the
positive Euler quotient is used.  This remains a conditional
statement about one Hecke function, not a verification of
$(\mathcal R_B)$.

\paragraph{The precise residual exceptional-prime law.}
If $k>0$, then $\Omega$ is infinite and, as $\epsilon\downarrow0$,
\[
 \sum_{p\in\Omega}p^{-1/2-\epsilon}
          =\frac{k}{4}\log\frac1\epsilon+O(1).
 \tag{S508.16}
\]
Thus an unclosed transfer forces an explicit positive logarithmic
singularity of the exceptional-prime harmonic sum, with the
same integer $k$ as the central excess.  This is an Abelian
statement at the Dirichlet-series boundary; no asymptotic
formula for the unweighted count of these primes is inferred.

\begin{proof}
By (S508.9), $D_\Omega$ has a pole of order $k$ at $1/2$.
It is positive for real $s>1/2$, so its leading Laurent
coefficient there is positive.  Consequently
\[
       \log D_\Omega(1/2+\epsilon)
             =k\log(1/\epsilon)+O(1).
\]
In the absolutely convergent half-plane the Euler logarithm is
\[
 \log D_\Omega(s)=4\sum_{p\in\Omega}p^{-s}+E(s),\qquad
 E(s)=4\sum_{p\in\Omega}\sum_{j\ge1}
              \frac{p^{-(2j+1)s}}{2j+1}.
 \tag{S508.17}
\]
The remainder is holomorphic for $\Re s>1/3$: on compact
subsets it is bounded termwise by a constant times
$\sum_p p^{-3\Re s}$, with locally uniform convergence.
In particular it is bounded near $s=1/2$.  Substituting the
real asymptotic for $\log D_\Omega$ into (S508.17) gives
(S508.16), with its exact coefficient $k/4$.
\end{proof}

\paragraph{Geometric and run attachment.}
The degree-six pushforward of the quadratic character of
$H/F$ is $\eta\oplus W$.  Equation (S508.12) isolates the
real divisor of the comparison between this Artin summand
and the known automorphic numerator.  The new gate is the
single central excess $\beta-b-n$, after the single real-zero
condition (S508.3); neither is evaluated by this proof.
This refines MT-S504/MT-S505 and the MT-S507 pole gate in
\path{GRH_NIGHT/complexity}.  The source comparison remains
a classical input, $A_D$ remains the designated Dirichlet
input, and $M_H$ remains conditional.  No universal Artin
holomorphy, functoriality, or zero-independent positive
cohomology statement is inferred.


% END root
\part{Unchecked claims and unresolved gates}
The complete proofs above establish only their stated implications.
The rows below are retained as research gates or as carried claims whose
complete source proofs were unavailable. They are not additional theorems.
The manuscript faults MT-M01 and MT-M02 have proved counterexamples;
the associated repairs do not certify every downstream operator or trace.

\paragraph{NT-Q06} Which pair-correlation and low-zero statistics are actually established?

\textbf{State:} proved/\allowbreak{}classical NT-Z01–\allowbreak{}NT-Z05 with support/\allowbreak{}moment restrictions; full GUE/\allowbreak{}correlations remain gate NT-Z06

\textbf{Source:} Montgomery1973 §1; HR2003 Thm3.\allowbreak{}1/\allowbreak{}5.\allowbreak{}1.

\paragraph{NT-Q07} Does rank-one Hecke GRH discharge the GL2 family hypothesis?

\textbf{State:} no such implication established; exact additional GL2-GRH gate NT-Z07 retained

\textbf{Source:} ILS2000 p56 and Thm1.\allowbreak{}1.

\paragraph{NT-Q10} Can carried class-group generator theorem be freshly source-inspected?

\textbf{State:} Bach1990 original remains inaccessible and carried honestly.\allowbreak{} The stronger modern generator theorem is freshly primary-checked in NT-C09/\allowbreak{}C10 (GM2025 exact constants/\allowbreak{}tables); BDDF and GM2025 finite certificates are freshly checked in NT-C11.\allowbreak{} Modern source gate closed without claiming recovery of Bach original.\allowbreak{}

\textbf{Source:} Bach1990 Thm4 pp375–\allowbreak{}376.

\paragraph{NT-Q19} Which prime-race statements follow from A alone and which still require ordinate independence?

\textbf{State:} NT-R01/\allowbreak{}R02 proved/\allowbreak{}classical under A; NT-R03/\allowbreak{}R04 retain explicit additional GSH gate

\textbf{Source:} Rubinstein–\allowbreak{}Sarnak1994 Thm1.\allowbreak{}1/\allowbreak{}1.\allowbreak{}2/\allowbreak{}1.\allowbreak{}4/\allowbreak{}1.\allowbreak{}5.

\paragraph{NT-Q20} Which multiplicity and individual gap/\allowbreak{}low-zero bounds follow, and does the cited theorem cover arbitrary infinity type?

\textbf{State:} NT-Z13–\allowbreak{}Z16 cover finite-order Hecke; full NT-Z17 symmetrization now proves conductor-uniform multiplicity/\allowbreak{}nearest-zero bounds for arbitrary non-pure-norm unitary infinity types, with fixed norm twists by translation.\allowbreak{} Individual unpaired S\ensuremath{\eta} argument bounds remain a separate gate.\allowbreak{} Central zeros are included.\allowbreak{}

\textbf{Source:} CCM2015 §4 axioms p6; Thm5p7; Thm7p13.

\paragraph{NT-Q30} Which concrete arithmetic condition supplies the canonical zero-basis gate rather than only GRH?

\textbf{State:} NT-L15 proves a fixed-Hecke weak mean-square hypothesis implies GRH, simplicity and exact reciprocal-derivative square summability; NT-Z19 deduces the absolute basis and positive trace realization.\allowbreak{} This is an additional hypothesis, not discharged by A/\allowbreak{}M.\allowbreak{} Whole proofs self-checked and root whole-peer passed; Hilbert-trace notation clarified.\allowbreak{}

\textbf{Source:} Titchmarsh§14.\allowbreak{}29pp376–\allowbreak{}380; MT-M16/\allowbreak{}M17.

\paragraph{NT-Q37} How far can CM Hasse deficits shrink absolutely while this GRH argument still gives an asymptotic?

\textbf{State:} For fixed \ensuremath{-}1/\allowbreak{}2<lambda<1/\allowbreak{}2, NT-E09 proves the exact potential-CM arcsine integral with x\textasciicircum{}(5/\allowbreak{}8+lambda/\allowbreak{}4)log²x error; its whole author/\allowbreak{}root readings passed.\allowbreak{} Established NT-E11 strengthens the all-CM-defined-over-base case, including nonabelian Frobenius restrictions, to the explicit optimized error delta[\ensuremath{\surd}xlogx L\_\allowbreak{}C+x\textasciicircum{}(5/\allowbreak{}8+lambda/\allowbreak{}4)\ensuremath{\surd}L\_\allowbreak{}C], hence x\textasciicircum{}(5/\allowbreak{}8+lambda/\allowbreak{}4)\ensuremath{\surd}logx for fixed data.\allowbreak{} E09 retains its valid older potential-CM bound here; no stronger descent theorem is silently substituted.\allowbreak{} The lambda=-1/\allowbreak{}2 norm-form gate remains, and the exact arcsine integral is essential for the stated power error when lambda>3/\allowbreak{}10.\allowbreak{}

\textbf{Source:} NT-E06; NT-E09.

\paragraph{NT-Q51} Can a cyclic class group be eliminated by deleting one explicitly bounded arithmetic point and obtain a Euclidean ring?

\textbf{State:} Yes under M:\allowbreak{} choose a degree-one generator prime within the C07 Hilbert-class-field bound.\allowbreak{} The full Pic quotient/\allowbreak{}localization and principal-class division proofs give a Euclidean one-point open.\allowbreak{} The needed subsequent Kummer RH family is for that chosen ring and is stated separately; cyclicity and existence versus computational complexity remain explicit.\allowbreak{} Whole author/\allowbreak{}geometry readings passed.\allowbreak{}

\textbf{Source:} NT-A21; NT-C07.

\paragraph{NT-Q53} Does M give effective non-CM residual open image, and what extra analytic family is used?

\textbf{State:} Yes over Q:\allowbreak{} every ell>964lograd(2N\_\allowbreak{}E)+5760 is surjective.\allowbreak{} Only the finite bounded-degree Dedekind family attached to relevant quadratic twists is used analytically.\allowbreak{} The modular-curve and local-inertia results are classical inputs.\allowbreak{} This is not the proposed uniform37 cutoff.\allowbreak{}

\textbf{Source:} NT-E15; Mayle–\allowbreak{}Wang2024Thm1p1401/\allowbreak{}§5.\allowbreak{}2pp1412–\allowbreak{}1413.

\paragraph{NT-Q56} Does the CM-recognition consequence also supply a near-quasilinear discriminant/\allowbreak{}endomorphism algorithm from GRH alone?

\textbf{State:} The complete dense-input E18 translation gives polynomial geometric CM decision only.\allowbreak{} It does not output D or endomorphisms.\allowbreak{} The cited faster CS Algorithm2 runtime retains heuristic(H) plus reference45§7.\allowbreak{}1 assumptions, and Enge floating-point exact-output complexity retains a rounding-error heuristic.\allowbreak{} Those specific gates are not silently removed by A or M; no impossibility claim about other algorithms.\allowbreak{}

\textbf{Source:} NT-E18; CSv4Theorem35p17; Enge2009Cor1.\allowbreak{}3.

\paragraph{AG-Q01} Which finite-type Z-scheme zeta functions have the adopted zero-support conclusion?

\textbf{State:} proved inherited exact Artin–\allowbreak{}Tate transfers; M adds identified potentiallyCM factors; residual nonCM continuation/\allowbreak{}RH and exceptional fibres remain explicit gates

\textbf{Source:} BRIDGES Q01; ZETA\_\allowbreak{}TWO ZT01–\allowbreak{}ZT07.

\paragraph{AG-Q08} Can the positive squared-variable determinant extend beyond \ensuremath{\zeta}?

\textbf{State:} Proved AG-DET/\allowbreak{}MOMENT/\allowbreak{}HANKEL and AG-SIGNED finite pole witness.\allowbreak{} MT-M29 now identifies the determinant operator with T\_\allowbreak{}+\textasciicircum{}\ensuremath{-}2 on the canonical image-power associated graded of the actual locally finite adelic module, retaining zero multiplicities.\allowbreak{} Independent positive prime/\allowbreak{}cohomological or GGP-period origin and arbitrary Artin moment positivity remain gates.\allowbreak{}

\textbf{Source:} ALPOGE\_\allowbreak{}BRIDGES BR-TR-01/\allowbreak{}02; SA23 geo-H-lidskii/\allowbreak{}geo-H-genus; cote1/\allowbreak{}79–\allowbreak{}80,111–\allowbreak{}112.

\paragraph{AG-Q11} Does finite-field purity imply openness of closure[pi0,H] in pi-prime under the original hypotheses?

\textbf{State:} literal ambientpi-prime refuted AG-C34-TYPE; intrinsic repair proved AG-C34-REPAIR/\allowbreak{}INDEX; with a specified section and finite derived length the raw action commutators are openH by AG-C34-SECTION; omitting solvability is refuted as an abstract inference by AG-C34-NONSOLV

\textbf{Source:} BRIDGES Q03; cote34/\allowbreak{}64 Rem1.\allowbreak{}17; definitions53–\allowbreak{}63.

\paragraph{AG-Q13} Do zero-location inputs construct integral Lefschetz correspondences?

\textbf{State:} partly resolved:\allowbreak{} AG-LEFSCH-CYCLES constructs actual inverse correspondences with the exact least denominator for product-polarized elliptic curves, including CM products; arbitrary varieties/\allowbreak{}polarizations and the particular theta-chain cycle comparison remain; AG-FOURIER-MIN gives the exact cohomological Fourier gate and AG-DP-EGS disproves its blanket very-general complex extension.\allowbreak{} AG-FOURIER-CM2 separately forces factor2 for full motivic Fourier on the actual CM curve; inverse-map identities remain valid.\allowbreak{}

\textbf{Source:} Hodge sketch§§3.\allowbreak{}2–\allowbreak{}3.\allowbreak{}3; cote16/\allowbreak{}14–\allowbreak{}18 no7,44–\allowbreak{}50.

\paragraph{AG-Q14} Does normalized L-sign prove arbitrary-cycle arithmetic Hodge index?

\textbf{State:} current gate:\allowbreak{} exact height identity for the given cycles and covered L-function; conditional sign deduction AG-HEIGHT proved

\textbf{Source:} Hodge sketch§3.\allowbreak{}5; cote44 Theorems4–\allowbreak{}5.

\paragraph{AG-Q16} Which finite/\allowbreak{}profinite GT gates move?

\textbf{State:} finite quotient detection/\allowbreak{}generation effective AG-FROB,CX-GEO-004/\allowbreak{}008; full profinite kernel/\allowbreak{}associator gate remains

\textbf{Source:} ALPOGE\_\allowbreak{}BRIDGES R-GTHAT; SA22 GT lines.

\paragraph{AG-Q17} Does global Hecke GRH give general weight-monodromy?

\textbf{State:} current gate:\allowbreak{} all-power trace bounds for finite exterior constructions and duals; prior reduction retained

\textbf{Source:} ALPOGE\_\allowbreak{}BRIDGES D-WM,BR-WM-02,CC-WM01–\allowbreak{}03.

\paragraph{AG-Q23} What exact extra positivity would eliminate poles after Hecke/\allowbreak{}Brauer divisor localization?

\textbf{State:} proved equivalence with all finite signed Hankel PSD and an explicit finite witness for any negative multiplicity; proving these inequalities from independent arithmetic data remains open

\textbf{Source:} AG-SIGNED; rootMT-ART02; cote1Fredholm framework.

\paragraph{AG-Q24} Does a specified arithmetic section and solvable H make its action-commutators open in H?

\textbf{State:} proved AG-C34-SECTION using all relevant finite covers; section existence not inferred

\textbf{Source:} Cote34PDF64/\allowbreak{}footer63 Rem1.\allowbreak{}17; AG-C34-SECTION.

\paragraph{AG-Q27} When does the Kummer class actually classify the punctured-normal-bundle fundamental-group extension?

\textbf{State:} classical comparison recovered under explicit char0/\allowbreak{}geoconnected/\allowbreak{}unibranch/\allowbreak{}pi2=0 hypotheses; actual normalizer map proved for the torsion-line-bundle family AG-C149-NORMAL, without pi2 for that group map; different families retain their comparison gate

\textbf{Source:} Cote149 originalPDF80–\allowbreak{}82/\allowbreak{}footer79–\allowbreak{}81, AG-C149-COMP.

\paragraph{AG-Q29} Does a Soulé-image assertion alone extend the regular-prime kernel-field proof to irregular primes?

\textbf{State:} exact freeness/\allowbreak{}generator-relation gate identified; under Greenberg canonical generators fail integrally despite Brown rational freeness; field-equality question remains unresolved AG-K8-TWIST now independently verifies the exactω\textasciicircum{}(\ensuremath{-}4) class-group component/\allowbreak{}H² twist5 vanishing; this is one indexed relation component, not the required complete freeness or generator presentation.\allowbreak{}

\textbf{Source:} LM6A44–\allowbreak{}A45; Fable§5; SharifiThm1.\allowbreak{}1/\allowbreak{}1.\allowbreak{}3.

\paragraph{AG-Q-CENTRAL-FUNCTOR} Does central divisor support construct a positive cohomological object with the actual Artin multiplicities?

\textbf{State:} Exact remaining gate AG-CENTRAL-FUNCTOR:\allowbreak{} all central Artin orders nonnegative, equivalently each full Heilbronn matrix positive.\allowbreak{} Under it the abstract star functor and dimensions are proved.\allowbreak{} Serre simplicity, geometric origin, integral lattice and leading-value regulator remain extra.\allowbreak{}

\textbf{Source:} Cote1; MT-ART02; Deninger2026 Thm2.\allowbreak{}15/\allowbreak{}2.\allowbreak{}16.

\paragraph{AG-Q-LEFSCH-CYCLES} Can the exact Wilson lattice denominator be attained by actual algebraic correspondences, with identities before discarding Chow torsion?

\textbf{State:} Yes for product-polarized elliptic curves:\allowbreak{} AG-LEFSCH-CYCLES constructs exact localized integral Chow inverses and proves optimality even among arbitrary algebraic cycles; no integral H1 basis or conjectural cycle-surjectivity is used.\allowbreak{} Other polarizations and the theta comparison retain their own gate.\allowbreak{}

\textbf{Source:} HodgeSketchRoute4; RETURN\_\allowbreak{}DP\_\allowbreak{}MULT DP15/\allowbreak{}Q4e; run20261112 H8; cote16scan14section7.

\paragraph{AG-Q-FOURIER-MIN} What exact cycle input permits an integral Fourier transform, and can a cohomological lift be silently promoted to rational Chow?

\textbf{State:} Classical gate recovered by AG-FOURIER-MIN:\allowbreak{} integral minimal curve class iff cohomological Fourier kernel/\allowbreak{}IHC1/\allowbreak{}algebraic divided-power closure.\allowbreak{} Weak Chow compatibility requires the stronger Chow-class lift.\allowbreak{} Full motivic compatibility is a further condition, explicitly obstructed in AG-FOURIER-CM2.\allowbreak{}

\textbf{Source:} HodgeSketchRoute4; BGF2023 Thm1.\allowbreak{}1/\allowbreak{}3.\allowbreak{}8, authorpp2,10–\allowbreak{}13.

\paragraph{AG-Q-FOURIER-ISOGENY} What quantitative integral obstruction remains when one attempts to transfer the Jacobian Fourier property through an isogeny after adding an abelian factor?

\textbf{State:} AG-FOURIER-ISOGENY proves degree d annihilates the entire one-cycle Hodge obstruction and kappa\_\allowbreak{}(g\ensuremath{-}1)|d; a gcd1 family suffices.\allowbreak{} The very-general EGS objects force every such degree even.\allowbreak{} Producing the required prime-to-obstruction isogenies in a specific CM family remains a gate.\allowbreak{}

\textbf{Source:} HodgeSketchRoutes2/\allowbreak{}4; BGF2023 Thm1.\allowbreak{}2/\allowbreak{}Prop4.\allowbreak{}5; EGS2026 Cor1.\allowbreak{}5p3/\allowbreak{}proofp81.

\paragraph{CX-Q04} What is the strongest available general finite-field factorization transfer?

\textbf{State:} proved conditional on M:\allowbreak{} CX-ALG-009; A\_\allowbreak{}D-only extension gate CX-ALG-010 (high valuation coefficient unit or alternative descent)

\textbf{Source:} CX\_\allowbreak{}GRH CX-N15/\allowbreak{}CX-U02.

\paragraph{CX-Q05} What does GRH give for class groups and what computation remains?

\textbf{State:} proved conditional generator transfer CX-ALG-011/\allowbreak{}013; full polynomial relation-search bound remains gate CX-ALG-012

\textbf{Source:} Bach1990T4; CX\_\allowbreak{}GRH CX-N07/\allowbreak{}N08.

\paragraph{CX-Q15} Do the cited GRH transfers close general derandomization or GCT hardness gates?

\textbf{State:} current gate:\allowbreak{} CX-LIM-001, uniform hitting sets or independent border-hardness theorem; no logical nonimplication claimed

\textbf{Source:} GCT V T4.\allowbreak{}7/\allowbreak{}L4.\allowbreak{}8; PIT/\allowbreak{}BPP questions.

\paragraph{CX-Q16} Does positivity/\allowbreak{}purity remove efficient multiplicity-construction gates?

\textbf{State:} specified earlier counterexamples carried CX-LIM-002; actual quotient ranks and uniform data construction remain

\textbf{Source:} ALPOGE\_\allowbreak{}GCT P-F7/\allowbreak{}P-F10 and PLETH PLE-F5/\allowbreak{}F7.

\paragraph{CX-Q32} What is the smallest sufficient family established here for the full factoring composition?

\textbf{State:} proved CX-ALG-025:\allowbreak{} pure Q(p\textasciicircum{}(1/\allowbreak{}r)) only for encountered prime r\ensuremath{\le}d with r\ensuremath{\nmid}(p\ensuremath{-}1)m and higher valuation; general A\_\allowbreak{}D construction of these prime-degree fields remains open

\textbf{Source:} CX\_\allowbreak{}GRH CX-N15/\allowbreak{}CX-U02; IKRS2011 Theorem5.\allowbreak{}8, Proposition4.\allowbreak{}6, Lemma4.\allowbreak{}2.

\paragraph{CX-Q33} Does exact extension construction share the same smaller analytic and unconditional promise gates?

\textbf{State:} proved CX-ALG-026 in poly(n log q), with n charged to output size; primes r|pm or valuation-one cases are unconditional

\textbf{Source:} Lenstra1991 Theorem9.\allowbreak{}1 pp344–\allowbreak{}346; CX-ALG-023/\allowbreak{}024.

\paragraph{CX-Q34} Can a structured degeneration avoid the unary Laurent-precision gate?

\textbf{State:} proved CX-GEO-010:\allowbreak{} complete torus-weight sign patterns have polynomial-bit integer representatives and a Boolean HN encoding

\textbf{Source:} CX-GEO-009; CX\_\allowbreak{}GRH CX-K02/\allowbreak{}CX-K05; ORDER28.

\paragraph{CX-Q36} Does the same encoding decide arbitrary orbit-closure membership?

\textbf{State:} proved only for closed target orbit (including nullcone); arbitrary-target one-cocharacter accessibility is not asserted, and succinct coordinate expansion remains a gate

\textbf{Source:} CX-GEO-011; Hilbert–\allowbreak{}Mumford closed-orbit theorem.

\paragraph{CX-Q37} Can a sparse degeneration yield AM without a unary Laurent-precision or syntactic-degree bound?

\textbf{State:} proved CX-GEO-012 when every intermediate circuit gate has unary/\allowbreak{}polynomial support width T; exponent classes and exact coefficient sums have one polynomial-size HN encoding

\textbf{Source:} CX-GEO-009/\allowbreak{}010; CX\_\allowbreak{}GRH CX-K02/\allowbreak{}CX-K05; ORDER28.

\paragraph{CX-Q43} Can whole zero blocks remove the individual derivative or separation gate for the canonical Xi closed quotient?

\textbf{State:} proved unconditionally MT-M18:\allowbreak{} stable quotient Banach-algebra block norms, all-jet closed ideal equality, rapid block isomorphism and absolute adapted Auerbach basis; prescribed individual jet coordinates retain MT-M16 gate

\textbf{Source:} Masullo9.\allowbreak{}3.\allowbreak{}2.\allowbreak{}5 PDF425–\allowbreak{}426; MT-M12/\allowbreak{}16; NT-Z10.

\paragraph{CX-Q46} Does the condition phi(C1)=0 obstruct comparison with all arithmetic quotient traces and positivity-test classes?

\textbf{State:} proved unconditionally MT-M30:\allowbreak{} continuous closed-ideal-valued restriction splitting, quotient isomorphism, compact restricted-test density, and exact omitted-spectrum witness; corrected tests need not retain positivity or support; zero-independent geometric narrow-strip identity remains the main gate

\textbf{Source:} Masullo§1.\allowbreak{}4.\allowbreak{}6p30;§1.\allowbreak{}7.\allowbreak{}2p35;(12.\allowbreak{}5.\allowbreak{}1.\allowbreak{}1)–\allowbreak{}(12.\allowbreak{}5.\allowbreak{}1.\allowbreak{}4)pp519–\allowbreak{}520; MT-M24/\allowbreak{}M27.

\paragraph{CX-Q52} Does the displayed norm-section average define a normalized indentation, and what valid replacement preserves traces?

\textbf{State:} MT-M36 proves literal average zero under the verified ordinary Haar convention.\allowbreak{} Normalized compact-open neighborhoods give exact finite-character projections, trace-norm convergence on trace-class modules and fixed-test stabilization; a genuine Q* indentation exists but its allocation among individual rational orbits varies with the cutoff while total trace is fixed.\allowbreak{} Prime-by-prime geometric allocation remains a distinct gate.\allowbreak{}

\textbf{Source:} Masullo11.\allowbreak{}3.\allowbreak{}2.\allowbreak{}2p495;A.\allowbreak{}2.\allowbreak{}3.\allowbreak{}14p549;A.\allowbreak{}3.\allowbreak{}2.\allowbreak{}1pp559–\allowbreak{}560;12.\allowbreak{}2.\allowbreak{}1.\allowbreak{}2p505;12.\allowbreak{}2.\allowbreak{}2.\allowbreak{}1–\allowbreak{}2p508.

\paragraph{CX-Q54} Can the remaining five-dimensional S5 Artin poles be excluded by an independently proved zero-noncoincidence or multiplicity bound beyond all monomial Hecke tests?

\textbf{State:} Full new proof complete:\allowbreak{} A\_\allowbreak{}D alone confines all S5/\allowbreak{}Q Artin poles in the open critical strip to \ensuremath{\zeta}Q/\allowbreak{}quadratic-Dirichlet zeros.\allowbreak{} With four-dimensional holomorphy, all43 monomial tests reduce to12 inequalities, forcing exact simultaneous-zero sets and sharp relative4v/\allowbreak{}total5v penalties (6v/\allowbreak{}8v when the companion four-dimensional order is0).\allowbreak{} UnderM\_\allowbreak{}H in the Calegari family S505 additionally excludes real poles.\allowbreak{} The finite cone has explicit negative-order witnesses, not asserted arithmetic; noncoincidence/\allowbreak{}pole absence remains open.\allowbreak{}

\textbf{Source:} MT-S506/\allowbreak{}S507; CalegariTheorem1.\allowbreak{}1; MT-S505; exact43-row S5 monomial certificate.

\paragraph{Q-ARTIN-POS} Can positivity remove critical-line Artin poles?

\textbf{State:} current general gate MT-ART02:\allowbreak{} positive semidefinite local Heilbronn matrix at multiple zeta zeros.\allowbreak{} Solvable branch narrowed by MT-ART09 to critical relative-zeta zeros of order\ensuremath{\ge}max(2,p2\ensuremath{-}2)+1, or\ensuremath{\ge}4 if4\ensuremath{\nmid}|G|.\allowbreak{} M\_\allowbreak{}H supplies location, not the remaining order bound.\allowbreak{} The specified S5 branch now has the strictly smaller MT-S504 relative central-order gate m5\ensuremath{\ge}2; central zeros inherited from \ensuremath{\zeta}\_\allowbreak{}F or L\ensuremath{\eta} no longer obstruct this sufficient route.\allowbreak{} MT-S505 additionally closes all real-axis poles in that family under M alone; only nonreal critical-line poles remain for that branch.\allowbreak{} MT-S506 proves A\_\allowbreak{}D-only critical-line pole support for every S5/\allowbreak{}Q extension.\allowbreak{} In the Calegari rank-four-holomorphic family, MT-S507 reduces the remaining rank-five poles to explicit simultaneous-zero sets with relative Dedekind order>=4 times the pole order (>=6 in the nonzero corresponding rank-four branch) and total order>=5 (respectively>=8).\allowbreak{} Relative order<=3 closes Artin holomorphy at the point.\allowbreak{} All43 monomial tests are exhausted; their formal sharp witnesses show the missing noncoincidence/\allowbreak{}order input is additional.\allowbreak{} MT-S508 weakens the real-axis exclusion input to A\_\allowbreak{}D plus the single quadratic-Hecke real-zero condition R\_\allowbreak{}B and refines the specified automorphic comparison gate to the excess beta-b-n, retaining central zeros of the known numerator.\allowbreak{} The expanded MT-S508 removes A\_\allowbreak{}D from that real-axis corollary:\allowbreak{} R\_\allowbreak{}B alone suffices, via positivity and Landau summability of the exceptional-prime quotient; a positive central excess also forces its exact square-root-boundary harmonic logarithm.\allowbreak{}

\textbf{Source:} ORDER42; Alpoge steps4,7.

\paragraph{Q-A5-LOCAL} What sharper local Artin-pole gates follow from all proper-subgroup Hecke tests for A5?

\textbf{State:} proved MT-ART05/\allowbreak{}06:\allowbreak{} exact cone, degree3/\allowbreak{}4 thresholds and central6/\allowbreak{}8 thresholds; actual divisor positivity remains an arithmetic gate

\textbf{Source:} MT-ART04; ORDER42; ZETA\_\allowbreak{}TWO ZT02.

\paragraph{Q-A5-NONREAL} Does the non-real A5-over-Q specialization retain an Artin holomorphy gate?

\textbf{State:} proved closed by classical automorphy; M transfer and tensor-category form MT-ART07; totally real case remains

\textbf{Source:} ORDER42; Calegari2013 §1 PDF2.

\paragraph{Q-S5-CENTRE} What is the exact remaining nonvanishing gate under M in the stated S5 family?

\textbf{State:} proved sharper exact comparison criterion MT-S508:\allowbreak{} only R\_\allowbreak{}B, the no-off-centre-real-zero condition for the single degree-six quadratic Hecke factor B, is needed.\allowbreak{} For beta=ord B, b=ord L\_\allowbreak{}eta and n=ord of the specified known cuspidal GL5 numerator at1/\allowbreak{}2, beta>=b+n; equality (equivalently beta<=b+n) holds iff the exceptional set is finite iff both rank-five factors match at every place.\allowbreak{} The remaining gate is beta-b-n>=1.\allowbreak{} Universal M supplies R\_\allowbreak{}B but is stronger; A\_\allowbreak{}D is unnecessary for this comparison theorem.\allowbreak{} MT-S504 is the n-discarding sufficient special case.\allowbreak{} No claim that failure of this specified comparison criterion rules out every possible automorphic realization.\allowbreak{} Expanded MT-S508 identifies the same excess k=n\_\allowbreak{}eta-ord(zeta\_\allowbreak{}F), so central nonvanishing of the known automorphic quadratic twist suffices.\allowbreak{} If k>0, the exceptional set has abscissa1/\allowbreak{}2 and sum p\textasciicircum{}(-1/\allowbreak{}2-epsilon)=(k/\allowbreak{}4)log(1/\allowbreak{}epsilon)+O(1).\allowbreak{} These are exact additional necessary features, not asserted nonvanishing or prime-count asymptotics.\allowbreak{}

\textbf{Source:} Calegari2013 Theorems1.\allowbreak{}1/\allowbreak{}1.\allowbreak{}2; ORDER42.

\paragraph{Q-MAS-INFINITE-BASIS} Can the required arithmetic infinite-divisor quotient be identified and given its asserted absolute basis and nuclear estimates?

\textbf{State:} proved unconditionally MT-M18/\allowbreak{}M19 adapted block synthesis, absolute basis and nuclear trace; MT-M21/\allowbreak{}M22 compare Meyer and MT-M24 identifies the actual Masullo smooth quotient.\allowbreak{} MT-M25 proves the exact line-Hilbert jet cutoff.\allowbreak{} MT-M26/\allowbreak{}27 prove fixed-character/\allowbreak{}full-adelic wide-strip trace class unconditionally; MT-M28 closes full-adelic every positive narrow strip under M.\allowbreak{} Prescribed individual-jet coordinates retain the MT-M16 quantitative coefficient gate; faithful positive full-jet action remains exactly GRH plus simple zeros (MT-M20/\allowbreak{}23).\allowbreak{} MT-M30 removes the restricted-test-domain gate; the independent geometric equality is Q-MAS-GEOMETRIC-NARROW.\allowbreak{}

\textbf{Source:} Masullo9.\allowbreak{}3.\allowbreak{}2.\allowbreak{}5; MT-M03/\allowbreak{}M04/\allowbreak{}M06; ORDER42.

\paragraph{Q-MAS-CANONICAL-SIMPLE} Can the canonical completed-Hecke quotient basis and its nuclear estimates be reduced to one explicit arithmetic condition?

\textbf{State:} proved MT-M12/\allowbreak{}M13 under GRH plus simple zeros:\allowbreak{} exact polynomial reciprocal-derivative gate; spectral synthesis, absolute basis and nuclear Hilbert trace follow when that gate holds

\textbf{Source:} Masullo9.\allowbreak{}3.\allowbreak{}2.\allowbreak{}5–\allowbreak{}9.\allowbreak{}3.\allowbreak{}3; NT-Z10; MT-M10/\allowbreak{}M11.

\paragraph{Q-MAS-CANONICAL-MULTIPLE} Does the exact canonical basis reduction extend beyond simple zeros?

\textbf{State:} proved MT-M16/\allowbreak{}M17 for every actual multiplicity pattern, removing MT-M14 uniform bound:\allowbreak{} one exact leading-coefficient/\allowbreak{}cluster-product gate gives the full basis, ideal and trace; no bounded-multiplicity assumption remains.\allowbreak{}

\textbf{Source:} Masullo9.\allowbreak{}3.\allowbreak{}2.\allowbreak{}5; MT-M11/\allowbreak{}M12.

\paragraph{Q-MAS-BLOCKS} Can grouping arbitrarily close arithmetic zeros remove the individual-coordinate basis and nuclearity gate?

\textbf{State:} proved unconditionally MT-M18/\allowbreak{}M19:\allowbreak{} stable finite Banach-algebra block norms, adapted absolute basis, exact spectral synthesis, explicit nuclear action and Hilbert trace; individual point-jet coordinates are a different assertion.\allowbreak{}

\textbf{Source:} Masullo9.\allowbreak{}3.\allowbreak{}2.\allowbreak{}5 originalPDF425–\allowbreak{}426; MT-M09/\allowbreak{}M16.

\paragraph{Q-ARTIN-RELATIVE-ORDER} Can the solvable local Artin gate remove zeros inherited from the base field?

\textbf{State:} proved MT-ART09:\allowbreak{} use q=ord(\ensuremath{\zeta}\_\allowbreak{}L/\allowbreak{}\ensuremath{\zeta}\_\allowbreak{}K), with thresholds max(2,p2\ensuremath{-}2) and at least3 when4\ensuremath{\nmid}|G|; possible poles under M\_\allowbreak{}H confined to the remaining higher-order relative zeros.\allowbreak{}

\textbf{Source:} Wong2019 Theorems1.\allowbreak{}3–\allowbreak{}1.\allowbreak{}4; Q-ARTIN-POS.

\paragraph{Q-CENTRAL-UNIVERSAL-NONVANISH} Could universal Hecke GRH by itself justify discarding a Dedekind central nonvanishing gate?

\textbf{State:} refuted as a universal nonvanishing assertion:\allowbreak{} explicit unconditional tame Q8 example in MT-ART10 has a central zero of order\ensuremath{\ge}2.\allowbreak{} Does not decide nonvanishing inside the specific S5 family.\allowbreak{}

\textbf{Source:} MT-S502 residual central gate; Fröhlich1974 Thm1; Martinet1971 Example3.

\paragraph{Q-NB-OPTIMAL-RATE} Does adopted RH give the conjectured optimal Nyman–\allowbreak{}Beurling rate and constant?

\textbf{State:} current gate:\allowbreak{} A gives the classical upper bound in MT-NB01 and Burnol gives the multiplicity-sensitive lower bound.\allowbreak{} The BCF route additionally needs \ensuremath{\Sigma}|\ensuremath{\zeta}\ensuremath{{}^{\prime}}|\textasciicircum{}\ensuremath{-}2\ensuremath{\ll}T\textasciicircum{}(3/\allowbreak{}2\ensuremath{-}\ensuremath{\delta}); MT-NB04 gives a sufficient all-zero derivative exponent\ensuremath{\beta}<1/\allowbreak{}4.\allowbreak{} Neither this moment nor simplicity is proved from A here.\allowbreak{}

\textbf{Source:} Balazard–\allowbreak{}deRoton2010mainTheorem; Bettin–\allowbreak{}Conrey–\allowbreak{}Farmer2013Theorem1; MT-M11.

\paragraph{Q-MAS-GEOMETRIC-NARROW} Does the intrinsic geometric trace remain the same on arbitrarily narrow completions without first assuming the Hecke zero location?

\textbf{State:} current exact gate after MT-M30:\allowbreak{} prove the source’s geometric equality to every narrow-strip trace on compact tests vanishing at norm one independently of zero location.\allowbreak{} MT-M30 unconditionally splits that test restriction through the closed theta ideal and proves any omitted zero has such a compact witness; thus no restricted-test-domain or positivity-test admissibility obstruction remains.\allowbreak{} MT-M27 is unconditional wide-strip trace, MT-M28 is narrow trace under M, MT-M29 is the conditional graded determinant comparison.\allowbreak{} None supplies the independent geometric equality.\allowbreak{} MT-M31 identifies and corrects the actual/\allowbreak{}closed-range defect in the local orbital quotient feeding11.\allowbreak{}3.\allowbreak{}2.\allowbreak{}1; this repair alone does not establish that geometric identification.\allowbreak{} MT-M32 supplies the corrected fixed-primary local Riesz/\allowbreak{}Poisson trace; comparison with the localized geometric kernel and uniform shell limits is still required.\allowbreak{} MT-M33 rules out nonzero bounded centered-module intertwiners into narrow closed invariant subquotients.\allowbreak{} MT-M34 refutes the literal finite-cutoff strictness used by11.\allowbreak{}3.\allowbreak{}2.\allowbreak{}1, even after closure repair:\allowbreak{} actual cutoff images are equal and their required lifts diverge at an exact rate.\allowbreak{} A replacement geometric-kernel/\allowbreak{}cutoff-limit construction is therefore required, rather than the source strong-convergence inference.\allowbreak{} MT-M35 proves collapse for every finite Euler product after cutoff and full cutoff cofinal intersection versus zero uncut intersection in the actual rational sector; an interchange of these constructions cannot supply the needed comparison.\allowbreak{} MT-M36 additionally refutes the literal section-average normalization in11.\allowbreak{}3.\allowbreak{}2.\allowbreak{}2 under its specified ordinary Haar measure.\allowbreak{} It constructs normalized projections and genuine idelic indentations, but exact local prime allocation and a replacement limiting kernel identity are still required.\allowbreak{} MT-M37 closes the single-compatible-indentation repair:\allowbreak{} global plus initial-local normalization forces support in G0 and kills every external prime shell.\allowbreak{} A varying-level replacement must specify its orbital weights and justify the trace limit; exact defect inequalities are now available.\allowbreak{}

\textbf{Source:} Masullo§1.\allowbreak{}4.\allowbreak{}6p30/\allowbreak{}§1.\allowbreak{}4.\allowbreak{}8–\allowbreak{}10pp31–\allowbreak{}32/\allowbreak{}Thm12.\allowbreak{}5.\allowbreak{}1.\allowbreak{}1pp519–\allowbreak{}520.

\paragraph{Q-S5-REAL-HOLOMORPHY} Does the specified S5 family retain a real-axis Artin pole gate under M when central nonvanishing is unknown?

\textbf{State:} proved MT-S505:\allowbreak{} every irreducible factor is holomorphic on real(0,1), and m5\ensuremath{\ge}0 even.\allowbreak{} A central pole would itself trigger full automorphy and contradict its pole.\allowbreak{} The twisted factor \ensuremath{\zeta}\_\allowbreak{}F/\allowbreak{}\ensuremath{\zeta}\_\allowbreak{}Q is holomorphic there unconditionally.\allowbreak{} Nonreal critical-line poles remain the residual branch when m5\ensuremath{\ge}2.\allowbreak{}

\textbf{Source:} Calegari2013Theorem1.\allowbreak{}1/\allowbreak{}Lemmas4.\allowbreak{}8–\allowbreak{}4.\allowbreak{}9; MT-S504.

\paragraph{Q-MAS-INDENTATION} Does the literal section average provide a normalized nonzero lift for the claimed geometric orbital trace?

\textbf{State:} refuted MT-M36 under the original ordinary Haar/\allowbreak{}indicator interpretation:\allowbreak{} the section average is zero for every number field.\allowbreak{} Normalized compact projections and genuine Q* indentations are explicitly constructed, with exact trace compatibility and demonstrated variability of individual orbital masses.\allowbreak{} Canonical prime allocation remains the separate comparison gate.\allowbreak{} MT-M37 closes the single-compatible-indentation repair:\allowbreak{} global plus initial-local normalization forces support in G0 and kills every external prime shell.\allowbreak{} A varying-level replacement must specify its orbital weights and justify the trace limit; exact defect inequalities are now available.\allowbreak{}

\textbf{Source:} Masullo11.\allowbreak{}3.\allowbreak{}2.\allowbreak{}2original495;12.\allowbreak{}2.\allowbreak{}1.\allowbreak{}2original505;12.\allowbreak{}2.\allowbreak{}2.\allowbreak{}1–\allowbreak{}2original508;A.\allowbreak{}2.\allowbreak{}3.\allowbreak{}14original549;A.\allowbreak{}3.\allowbreak{}2.\allowbreak{}1–\allowbreak{}2original560.

\paragraph{Q-S5-DIRICHLET-POLES} Does localizing possible S5 Artin poles in the critical strip require Hecke GRH over all fields?

\textbf{State:} proved MT-S506:\allowbreak{} A\_\allowbreak{}D suffices for every S5/\allowbreak{}Q extension, without the Calegari local hypotheses.\allowbreak{} All possible poles are supported on zeta\_\allowbreak{}Q and the quadratic Dirichlet zero sets; X is entire by explicit quartic induction.\allowbreak{} Absence of poles is a separate gate.\allowbreak{}

\textbf{Source:} MT-S501/\allowbreak{}505; PY07section2.\allowbreak{}1Properties1-2p3,section3.\allowbreak{}1pp4-5,Theorem3.\allowbreak{}3p5; S5 character table.

\paragraph{Q-S5-MONOMIAL-CONE} What do all induced one-dimensional Hecke positivity tests imply about the remaining rank-five S5 poles once the rank-four factors are holomorphic?

\textbf{State:} proved exact reduction MT-S507:\allowbreak{}43 distinct tests reduce to12 inequalities, force q>=4v/\allowbreak{}r>=5v and stated common zeros, with stronger q>=6v/\allowbreak{}r>=8v when the corresponding rank-four factor is nonzero.\allowbreak{} Four formal witnesses prove these cone constants optimal.\allowbreak{} Remaining gate:\allowbreak{} rule out the explicit nonreal simultaneous-zero configurations or their order thresholds; the formal witnesses are not realized arithmetic examples.\allowbreak{}

\textbf{Source:} MT-ART07; MT-S501/\allowbreak{}504/\allowbreak{}505; exact S5 subgroup-character certificate.

\paragraph{Q-S5-SINGLE-HECKE} Which individual Hecke input suffices for the S5 comparison, and can zeros of the automorphic numerator be retained in its central gate?

\textbf{State:} proved MT-S508:\allowbreak{} the one real-zero input R\_\allowbreak{}B suffices for the exact central criterion, absolute exceptional-prime Euler convergence on Re>1/\allowbreak{}2 and all real-axis Artin holomorphy, without A\_\allowbreak{}D.\allowbreak{} Full GRH\_\allowbreak{}B transfers to the known GL5 numerator and C before central equality; W has critical-line divisor support but can retain poles there.\allowbreak{} Exact transfer iff k=beta-b-n=n\_\allowbreak{}eta-ord(zeta\_\allowbreak{}F)=0.\allowbreak{} A positive k forces the exact exceptional-prime boundary logarithm(k/\allowbreak{}4)log(1/\allowbreak{}epsilon)+O(1).\allowbreak{} Other factors require their stated zero-location inputs.\allowbreak{}

\textbf{Source:} MT-S504/\allowbreak{}505; Calegari publisher Lemmas4.\allowbreak{}8-4.\allowbreak{}10pp13-15 and finalTheorem1.\allowbreak{}2proofp15.
\section*{Statement index}
\addcontentsline{toc}{section}{Statement index}
The complete statement, exact input, proof location and work placement for
every ID below is recorded in \path{CONSEQUENCES.tsv}. This index keeps
the paper navigable without duplicating the entire machine-readable table.
\paragraph{number theory} NT-P01, NT-P02, NT-P03, NT-P04, NT-P05, NT-P06, NT-P07, NT-P08, NT-P09, NT-P10, NT-C01, NT-C02, NT-C03, NT-C04, NT-C05, NT-H01, NT-H02, NT-H03, NT-H04, NT-A01, NT-A02, NT-A03, NT-F01, NT-F02, NT-F03, NT-F04, NT-L01, NT-L02, NT-L03, NT-L04, NT-L05, NT-L06, NT-L07, NT-L08, NT-Z01, NT-Z02, NT-Z03, NT-Z04, NT-Z05, NT-Z06, NT-Z07, NT-Z08, NT-Z09, NT-L09, NT-L10, NT-Z10, NT-P11, NT-P12, NT-Z11, NT-Z12, NT-H05, NT-H06, NT-H07, NT-H08, NT-R01, NT-R02, NT-R03, NT-R04, NT-Z13, NT-Z14, NT-Z15, NT-Z16, NT-P13, NT-P14, NT-P15, NT-P16, NT-L11, NT-L12, NT-H09, NT-H10, NT-Z17, NT-Z18, NT-C06, NT-C07, NT-C08, NT-L13, NT-L14, NT-H11, NT-L15, NT-Z19, NT-H12, NT-H13, NT-L16, NT-L17, NT-E01, NT-E02, NT-E03, NT-E04, NT-E05, NT-R05, NT-R06, NT-R07, NT-R08, NT-E06, NT-E07, NT-E08, NT-E09, NT-A04, NT-A05, NT-E10, NT-A06, NT-A07, NT-A08, NT-A09, NT-A10, NT-A11, NT-C09, NT-C10, NT-C11, NT-C12, NT-C13, NT-A12, NT-A13, NT-A14, NT-E11, NT-A15, NT-A16, NT-A17, NT-E12, NT-E13, NT-A18, NT-A19, NT-A20, NT-A21, NT-E14, NT-E15, NT-E16, NT-E17, NT-E18.
\paragraph{geometry programme} AG-AT, AG-ARTIN, AG-DED, AG-LOCAL, AG-MODEL, AG-CM, AG-WEIGHT, AG-SCHEMES, AG-NONCM, AG-NONCM-W0, AG-FROB, AG-LAMBDA, AG-SIGN, AG-HEIGHT, AG-SIGN-NORM, AG-CRITICAL, AG-DET, AG-MOMENT, AG-HANKEL, AG-EXT, AG-EXT-K, AG-C1, AG-C16, AG-C34-74, AG-C34-64, AG-C44, AG-C149, AG-NODE, AG-INERTIA, AG-GT, AG-PURITY, AG-WM, AG-STANDARD, AG-COVER, AG-KUMMER, AG-KUMMER-2, AG-CM-REC, AG-CM-PERIOD, AG-SIGNED, AG-C34-TYPE, AG-C34-REPAIR, AG-C34-INDEX, AG-C34-SECTION, AG-C34-NONSOLV, AG-C149-COMP, AG-C149-LIFT, AG-C149-ELLIPTIC, AG-LM6-REG, AG-LM6-IRREG, AG-LM6-NORM, AG-FG-EXP, AG-FG-DER, AG-FG-PPOWER, AG-FG-PRODUCT, AG-FG-THETA, AG-ARTHUR-ARCH, AG-ARTHUR-PARAM, AG-ARTHUR-MIN, AG-LEFSCH-LATTICE, AG-LEFSCH-SPLIT, AG-C149-NORMAL, AG-C149-FULLLIFT, AG-C149-CUSP, AG-LEFSCH-INDEX, AG-LEFSCH-DEFECT, AG-CYCLIC-DIV, AG-THETA-REG, AG-THETA-TORS, AG-THETA-HECKE, AG-THETA-CLASS, AG-THETA-RAY, AG-FG-POLE, AG-THETA-PCOND, AG-THETA-NONINT, AG-CM-INTEGRAL-MODEL, AG-CM-MODEL-DEGREE, AG-FG-FOURIER, AG-THETA-FOURIER, AG-SHIMURA-PERIOD, AG-SHIMURA-SLOPE, AG-SHIMURA-RAM, AG-DENINGER, AG-CENTRAL-STRUCT, AG-CENTRAL-FUNCTOR, AG-CENTRAL-REAL, AG-EXP-VANISH, AG-EXP-SEPARATE, AG-CENTRAL-Q8, AG-LOCAL-EXP, AG-LOCAL-EXCEPTION, AG-LOCAL-DET, AG-LOCAL-PURITY, AG-FROB-PROJECTOR, AG-FROB-ORDER, AG-FROB-COLLISION, AG-LOCAL-WD, AG-LOCAL-SPLICE, AG-FROB-NIL, AG-FROB-CHOW, AG-FROB-CYCLES, AG-FROB-Z, AG-LEFSCH-SMITH, AG-LEFSCH-CYCLES, AG-FOURIER-MIN, AG-DP-ALG, AG-DP-EGS, AG-FOURIER-ISOGENY, AG-FOURIER-CM2, AG-K8, AG-K8-TWIST.
\paragraph{complexity} CX-ALG-001, CX-ALG-002, CX-ALG-003, CX-ALG-004, CX-ALG-005, CX-ALG-006, CX-ALG-007, CX-ALG-008, CX-ALG-009, CX-ALG-010, CX-ALG-011, CX-ALG-012, CX-ALG-013, CX-ALG-014, CX-HN-001, CX-HN-002, CX-HN-003, CX-HN-004, CX-GEO-001, CX-GEO-002, CX-GEO-003, CX-GEO-004, CX-GEO-005, CX-GEO-006, CX-NN-001, CX-NN-002, CX-NN-003, CX-NN-004, CX-LIM-001, CX-LIM-002, CX-LIM-003, CX-GEO-007, CX-GEO-008, CX-ALG-015, CX-ALG-016, CX-ALG-017, CX-ALG-018, CX-GEO-009, CX-ALG-019, CX-ALG-020, CX-ALG-021, CX-ALG-022, CX-ALG-023, CX-ALG-024, CX-ALG-025, CX-ALG-026, CX-GEO-010, CX-GEO-011, CX-GEO-012, CX-GEO-013, CX-ALG-027, CX-ALG-028, CX-GEO-014, CX-GEO-015.
\paragraph{root} MT-A01, MT-M01, MT-M02, MT-M03, MT-M04, MT-ART01, MT-ART02, MT-ART03, MT-L01, MT-L02, MT-ART04, MT-Q01, MT-Q02, MT-ART05, MT-ART06, MT-H01, MT-H02, MT-H03, MT-ART07, MT-ART08, MT-S501, MT-S502, MT-S503, MT-M05, MT-M06, MT-M07, MT-M08, MT-M09, MT-M10, MT-M11, MT-M12, MT-M13, MT-M14, MT-M15, MT-M16, MT-M17, MT-M18, MT-M19, MT-M20, MT-M21, MT-M22, MT-M23, MT-ART09, MT-M24, MT-M25, MT-ART10, MT-M26, MT-M27, MT-ART11, MT-ART12, MT-M28, MT-NB01, MT-NB02, MT-NB03, MT-NB04, MT-M29, MT-M30, MT-S504, MT-S505, MT-M31, MT-M32, MT-M33, MT-M34, MT-M35, MT-M36, MT-M37, MT-S506, MT-S507, MT-S508.
\section*{Source registry and exact locators}
\addcontentsline{toc}{section}{Source registry and exact locators}
The inspection descriptions below record their actual scope. A located
classical theorem is not labelled as a fresh whole-proof audit unless
its event explicitly records that reading.

\paragraph{S-NT-LLS} Lamzouri–\allowbreak{}Li–\allowbreak{}Soundararajan, Conditional bounds for the least quadratic non-residue and related problems, corrected arXiv v3 (8 Dec 2016).

\textbf{Locator:} Theorem1.\allowbreak{}1 and Corollary1.\allowbreak{}1 p2; Corollary1.\allowbreak{}2 p3; Theorem1.\allowbreak{}5 pp3–\allowbreak{}4; Corollary1.\allowbreak{}3 and (1.\allowbreak{}3) p4.\allowbreak{} Author-page numbering.\allowbreak{}.

\textbf{Inspection:} Primary PDF text read for full statements, thresholds, U/\allowbreak{}V normalization and correction-version choice.\allowbreak{} These results were already background IA02; no fresh reproof of the long literature argument.\allowbreak{}

\url{https://arxiv.org/pdf/1309.3595v3}.

\paragraph{S-NT-GM} Grenié–\allowbreak{}Molteni, An explicit Chebotarev density theorem under GRH, J.\allowbreak{}Number Theory200(2019),441–\allowbreak{}485.

\textbf{Locator:} Theorem1.\allowbreak{}1 author p2; Corollary1.\allowbreak{}2 author p3; AppendixA TheoremA.\allowbreak{}1 author p26.\allowbreak{}.

\textbf{Inspection:} Original author PDF pp2–\allowbreak{}3 rendered and visually inspected for both explicit error polynomials and unramified convention.\allowbreak{} AppendixA finite-order Hecke zero-count statement read; that appendix is unconditional.\allowbreak{} The local O(log T) deduction uses consecutive heights.\allowbreak{} At03:\allowbreak{}27UTC original AppendixA p26 was also rendered and visually read:\allowbreak{} TheoremA.\allowbreak{}1 is unconditional for every primitive finite-order Hecke character, counting multiplicities and only0<Re\ensuremath{\rho}<1.\allowbreak{} Its TlogT main term proves infinitely many zeros for NT-L17.\allowbreak{} The long zero-count proof is cited, not independently recertified.\allowbreak{} At04:\allowbreak{}03 in this run, TheoremA.\allowbreak{}1 and Table1 text reread for the trivial-character uniform count used in NT-R08; no numerical constant newly asserted.\allowbreak{}

\url{https://sites.unimi.it/molteni/research/papers-pdf/44-molteni-An_explicit_Chebotarev_density_theorem_under_GRH.pdf}.

\paragraph{S-NT-BS} Bach–\allowbreak{}Sorenson, Explicit bounds for primes in residue classes, Math.\allowbreak{}Comp.\allowbreak{}65(1996),1717–\allowbreak{}1735.

\textbf{Locator:} §2 printed1720; Corollary3.\allowbreak{}3 printed1721; Theorem5.\allowbreak{}1 printed1729/\allowbreak{}PDF13.\allowbreak{}.

\textbf{Inspection:} Primary author-upload PDF opened in this agent.\allowbreak{} Geometry agent confirmed precise Thm5.\allowbreak{}1 statement; RH \ensuremath{\zeta}\_\allowbreak{}L and \ensuremath{\zeta}\_\allowbreak{}Q; degree-one prime; L\ensuremath{\ne}Q.\allowbreak{} Arbitrary extra-prime exclusion is not in this theorem.\allowbreak{}

\url{https://www.researchgate.net/profile/Jonathan-Sorenson-2/publication/220576951_Explicit_bounds_for_primes_in_residue_classes/links/55100ea70cf2ac2905afbceb/Explicit-bounds-for-primes-in-residue-classes.pdf}.

\paragraph{S-NT-SCHOENFELD} Schoenfeld, Sharper bounds for the Chebyshev functions \ensuremath{\theta}(x) and \ensuremath{\psi}(x).\allowbreak{}II, Math.\allowbreak{}Comp.\allowbreak{}30(1976),337–\allowbreak{}360.

\textbf{Locator:} Theorem10 equations(6.\allowbreak{}1)–\allowbreak{}(6.\allowbreak{}5), printed337; Corollary1 to Theorem10.\allowbreak{}.

\textbf{Inspection:} Primary indexed full statement of Theorem10 read, including 73.\allowbreak{}2/\allowbreak{}599 and 23·10\textasciicircum{}8.\allowbreak{} Direct AMS PDF403; mirror https:\allowbreak{}/\allowbreak{}/\allowbreak{}electronicsandbooks.\allowbreak{}com/\allowbreak{}edt/\allowbreak{}manual/\allowbreak{}Magazine/\allowbreak{}M/\allowbreak{}Mathematics\%20of\%20Computation/\allowbreak{}1960-2002/\allowbreak{}pdf/\allowbreak{}1976\_\allowbreak{}v030\_\allowbreak{}n134/\allowbreak{}2005976.\allowbreak{}pdf timed out via web.\allowbreak{} Corollary1 original page still to inspect; do not label full source proof checked.\allowbreak{}

\url{https://www.ams.org/mcom/1976-30-134/S0025-5718-1976-0457374-X/S0025-5718-1976-0457374-X.pdf}.

\paragraph{S-NT-DGM} Dudek–\allowbreak{}Grenié–\allowbreak{}Molteni, Primes in explicit short intervals on RH, Int.\allowbreak{}J.\allowbreak{}NumberTheory12(2016),1391–\allowbreak{}1407.

\textbf{Locator:} Theorem1.\allowbreak{}1 author p1.\allowbreak{}.

\textbf{Inspection:} Primary theorem statement, c=1/\allowbreak{}2+2/\allowbreak{}log x and x\ensuremath{\ge}2 read.\allowbreak{} Computational c=.\allowbreak{}55 refinement not used.\allowbreak{}

\url{https://sites.unimi.it/molteni/research/papers-pdf/36-molteni-Primes_in_explicit_short_intervals_on_RH.pdf}.

\paragraph{S-NT-GMP} Grenié–\allowbreak{}Molteni–\allowbreak{}Perelli, Primes and prime ideals in short intervals, arXiv:\allowbreak{}1602.\allowbreak{}02906v2.

\textbf{Locator:} Theorems1,2 and subsequent asymptotic statements, author p2; Proposition1 author p3.\allowbreak{}.

\textbf{Inspection:} Primary theorem statements and distinction Dirichlet-GRH versus RH \ensuremath{\zeta}\_\allowbreak{}K read.\allowbreak{} Absolute constants remain unspecified.\allowbreak{}

\url{https://arxiv.org/pdf/1602.02906}.

\paragraph{S-NT-SHOUP} Shoup, Searching for primitive roots in finite fields (1992).

\textbf{Locator:} Theorem3, author PDF p4.\allowbreak{}.

\textbf{Inspection:} Complexity agent downloaded original PDF with a standard browser user agent.\allowbreak{} This agent inspected the actual rendered p4 image at work/\allowbreak{}complexity/\allowbreak{}sources/\allowbreak{}Shoup\_\allowbreak{}Theorem3.\allowbreak{}png.\allowbreak{} Exact r\textasciicircum{}4(logr+1)\textasciicircum{}4log²p bound confirmed; no numerical implicit constant claimed.\allowbreak{}

\url{https://www.shoup.net/papers/primroots.pdf}.

\paragraph{S-NT-BACH} Bach, Explicit bounds for primality testing and related problems, Math.\allowbreak{}Comp.\allowbreak{}55(1990),355–\allowbreak{}380.

\textbf{Locator:} Theorem4 printed375 and corollary376, inherited exact locator.\allowbreak{}.

\textbf{Inspection:} Carried as background with no fault found; direct AMS requests inaccessible.\allowbreak{} Primary fresh inspection not claimed.\allowbreak{} Complexity agent owns the algorithmic application.\allowbreak{}

\url{https://www.ams.org/mcom/1990-55-191/S0025-5718-1990-1023756-8/S0025-5718-1990-1023756-8.pdf}.

\paragraph{S-NT-LENSTRA} Lenstra, On Artin’s conjecture and Euclid’s algorithm in global fields, Invent.\allowbreak{}Math.\allowbreak{}42(1977),201–\allowbreak{}224.

\textbf{Locator:} §2 pp202–\allowbreak{}206; (2.\allowbreak{}14),(2.\allowbreak{}15) pp205–\allowbreak{}206; Theorem(3.\allowbreak{}1)206; Theorem(4.\allowbreak{}1)207; Theorem(9.\allowbreak{}1),(9.\allowbreak{}2)218.\allowbreak{}.

\textbf{Inspection:} Primary definitions and theorem statements read; original p218 rendered and inspected, confirming \#S\ensuremath{\ge}2.\allowbreak{} Full elementary reduction in Theorem3.\allowbreak{}1/\allowbreak{}lemma3.\allowbreak{}2 read, but Cooke–\allowbreak{}Weinberger base theorem not independently reproved.\allowbreak{}

\url{https://pub.math.leidenuniv.nl/~lenstrahw/PUBLICATIONS/1977d/fulltext-12.pdf}.

\paragraph{S-NT-HOOLEY} Hooley, On Artin’s conjecture, J.\allowbreak{}reine angew.\allowbreak{}Math.\allowbreak{}225(1967),209–\allowbreak{}220.

\textbf{Locator:} Unnumbered final theorem pp219–\allowbreak{}220; original internal locator corroborated by later primary research.\allowbreak{}.

\textbf{Inspection:} Original publication metadata checked; original fulltext endpoint blocked.\allowbreak{} For the exact proof used here, cite inspected Lenstra Theorem3.\allowbreak{}1 instead.\allowbreak{} Original theorem internal-page check remains a source gate, not a mathematical gate on the classical corollary.\allowbreak{}

\url{https://doi.org/10.1515/crll.1967.225.209}.

\paragraph{S-NT-NEUKIRCH} Neukirch, Algebraic Number Theory (1999), ChapterVII §11.

\textbf{Locator:} VII(11.\allowbreak{}7) and proof printed533–\allowbreak{}534/\allowbreak{}PDF541–\allowbreak{}542; Cor(11.\allowbreak{}8), conductor–\allowbreak{}discriminant(11.\allowbreak{}9) printed534/\allowbreak{}PDF542; degree-one conductor identification(11.\allowbreak{}10) printed535/\allowbreak{}PDF543.\allowbreak{}.

\textbf{Inspection:} Full primary PDF retrieved locally from institutional host; web-open rejected its file size.\allowbreak{} Inspected original images of all three pages and read the entire supplied induction/\allowbreak{}conductor–\allowbreak{}discriminant proof, retaining its earlier local conductor and standard representation-theory inputs.\allowbreak{} Degree-one identification statement read.\allowbreak{} No claim to reprove all global class field theory.\allowbreak{}

\url{https://www.math.toronto.edu/~ila/Neukirch_Algebraic_number_theory.pdf}.

\paragraph{S-NT-CHANDEE} Chandee, Explicit upper bounds for L-functions on the critical line, Proc.\allowbreak{}Amer.\allowbreak{}Math.\allowbreak{}Soc.\allowbreak{}137(2009),4049–\allowbreak{}4063; arXiv0906.\allowbreak{}4177v1.

\textbf{Locator:} Axioms/\allowbreak{}conductor(2)–\allowbreak{}(6) authorp2; Cor1.\allowbreak{}1p3; Theorem2.\allowbreak{}1/\allowbreak{}logarithmic inequality proofpp4–\allowbreak{}6; Cor1.\allowbreak{}1 proof(14)–\allowbreak{}(15)p7; height-shift Remark4p4.\allowbreak{}.

\textbf{Inspection:} Original PDF downloaded; p2–\allowbreak{}3 rendered and visually inspected.\allowbreak{} Read complete Cor1.\allowbreak{}1 proof on p7, including the extra RH-zeta use for the explicit prime-power estimate.\allowbreak{} Confirmed no gamma-conjugate-pair restriction and retained the pole-free boundary assumption.\allowbreak{} Published numerical input in (15), its finite numerical check, and the entire main theorem proof are cited rather than independently recomputed/\allowbreak{}re-refereed.\allowbreak{} NT-L13 verifies the Hecke axioms, larger-conductor substitution before x is chosen, and exact epsilon/\allowbreak{}degree threshold in full.\allowbreak{}

\url{https://arxiv.org/pdf/0906.4177}.

\paragraph{S-NT-LUMLEY} Lumley, Explicit bounds for L-functions on the edge of the critical strip, J.\allowbreak{}NumberTheory188(2018),186–\allowbreak{}209; arXiv1804.\allowbreak{}09850v1.

\textbf{Locator:} Definitions§1.\allowbreak{}1 authorpp2–\allowbreak{}3; Thm2.\allowbreak{}1–\allowbreak{}2.\allowbreak{}2 and constants(2.\allowbreak{}1)–\allowbreak{}(2.\allowbreak{}3) pp3–\allowbreak{}4; Cor2.\allowbreak{}3p5; extra RH-zeta Lemmas3.\allowbreak{}6–\allowbreak{}3.\allowbreak{}7p10; proofs§4.\allowbreak{}1pp11–\allowbreak{}13 and §4.\allowbreak{}2pp14–\allowbreak{}16.\allowbreak{}.

\textbf{Inspection:} Primary PDF downloaded; original p3 axioms and p4 full theorem/\allowbreak{}constant images inspected.\allowbreak{} Read exact conductor convention, entire-completion requirement, complex infinity parameters, both thresholded inequalities and RH-zeta lemmas.\allowbreak{} J1 is bounded above, so NT-L14 explicitly substitutes its displayed majorant.\allowbreak{} Full main analytic proof and numerical optimization are cited, not independently re-refereed.\allowbreak{} NT-H11 Euler-factor/\allowbreak{}residue/\allowbreak{}conductor deduction is original and checked separately.\allowbreak{}

\url{https://arxiv.org/pdf/1804.09850}.

\paragraph{S-NT-TITCHMARSH} Titchmarsh, The Theory of the Riemann Zeta-function, 2nd ed.\allowbreak{}, revised by Heath-Brown (1986).

\textbf{Locator:} Lemma3.\allowbreak{}12 and proof printed61–\allowbreak{}62; Theorem14.\allowbreak{}2 and (14.\allowbreak{}2.\allowbreak{}5),(14.\allowbreak{}2.\allowbreak{}6) printed336–\allowbreak{}337; Theorem14.\allowbreak{}14(A); Theorem14.\allowbreak{}25(A)–\allowbreak{}(C) printed369–\allowbreak{}370.\allowbreak{} Weak Mertens §14.\allowbreak{}29:\allowbreak{} (14.\allowbreak{}29.\allowbreak{}1)/\allowbreak{}Theorem14.\allowbreak{}29(A) pp376–\allowbreak{}377; derivative bound(14.\allowbreak{}29.\allowbreak{}3)p377; Theorem14.\allowbreak{}29(B)/\allowbreak{}(14.\allowbreak{}29.\allowbreak{}4)p377, proofpp377–\allowbreak{}380.\allowbreak{} In the two-up scan, printed376–\allowbreak{}377 is PDF193.\allowbreak{}.

\textbf{Inspection:} Downloaded local PDF and text.\allowbreak{} Relevant theorem statements and the whole printed 14.\allowbreak{}25 argument read; Perron kernel locator checked for the new transfer proof.\allowbreak{} Not a check of the entire textbook.\allowbreak{} Two printed pages per PDF page.\allowbreak{} Full reference text excluded from deliverables.\allowbreak{} Original printed376–\allowbreak{}377 image inspected and whole Theorem14.\allowbreak{}29(A) proof read.\allowbreak{} Theorem14.\allowbreak{}29(B) statement and proof endpoint checked; its full textbook contour proof is not newly certified.\allowbreak{} NT-L15 instead supplies and checks a complete independent Laplace/\allowbreak{}Abel-Bessel proof for Hecke functions with an exact limsup constant.\allowbreak{}

\url{https://sites.math.rutgers.edu/~zeilberg/EM18/TitchmarshZeta.pdf}.

\paragraph{S-NT-SIMONIC} Simonič, Estimates for L-functions in the Critical Strip Under GRH with Effective Applications, Mediterr.\allowbreak{}J.\allowbreak{}Math.\allowbreak{}20(2023),87.

\textbf{Locator:} Theorem1 equations(4)–\allowbreak{}(6); Corollaries2–\allowbreak{}3; §2.\allowbreak{}1 SP axioms; Theorem2 equations(7),(8); proof §3.\allowbreak{}.

\textbf{Inspection:} Primary full-text theorem and axioms read.\allowbreak{} Fixed-character consequences derived in full in section.\allowbreak{}tex; effective Mertens constants checked in original HTML.\allowbreak{} No uniform hybrid-conductor Lindelöf estimate claimed.\allowbreak{}

\url{https://link.springer.com/article/10.1007/s00009-023-02289-2}.

\paragraph{S-NT-HARPER} Harper, Sharp conditional bounds for moments of the Riemann zeta function, arXiv:\allowbreak{}1305.\allowbreak{}4618.

\textbf{Locator:} Theorem1 authorp1; Theorem2 authorp15, proof§5 pp15–\allowbreak{}17, Propositions3/\allowbreak{}4.\allowbreak{}.

\textbf{Inspection:} Primary theorem statement read with k\ensuremath{\ge}0 fixed and implicit constant depending on k.\allowbreak{} Source introduction distinguishes order bounds from precise moments.\allowbreak{} No complete reproof of the paper.\allowbreak{} Theorem2 original p15 rendered and inspected for k\ensuremath{\ge}0, both signs of fundamental d, discriminant shell and exponent k(k+1)/\allowbreak{}2.\allowbreak{} New Markov deductions in NT-L12 are proved in full; the paper’s full moment argument is not independently recertified.\allowbreak{}

\url{https://arxiv.org/pdf/1305.4618}.

\paragraph{S-NT-MONTGOMERY} Montgomery, The pair correlation of zeros of the zeta function, Proc.\allowbreak{}Sympos.\allowbreak{}PureMath.\allowbreak{}24(1973),181–\allowbreak{}193.

\textbf{Locator:} §1 definition and unnumbered Theorem p181; Corollary2 p182; conjecture following restricted theorem.\allowbreak{}.

\textbf{Inspection:} Primary web PDF text read for normalization, \ensuremath{\alpha} range, uniformity and simplicity proportion; direct local download403.\allowbreak{} Full pair-correlation conjecture carefully separated.\allowbreak{}

\url{https://public.websites.umich.edu/hlm/paircor1.pdf}.

\paragraph{S-NT-HR} Hughes–\allowbreak{}Rudnick, Linear statistics of low-lying zeros of L-functions, Q.\allowbreak{}J.\allowbreak{}Math.\allowbreak{}54(2003),309–\allowbreak{}333.

\textbf{Locator:} v2:\allowbreak{} ordering p1; admissibility §2; Theorem3.\allowbreak{}1p5; Theorem4.\allowbreak{}1p7; Theorem5.\allowbreak{}1p9; Corollary8.\allowbreak{}2p20; Theorem8.\allowbreak{}3pp20–\allowbreak{}21.\allowbreak{}.

\textbf{Inspection:} Primary text read for unconditional complex-zero formulation, Fourier support, m<2/\allowbreak{}\ensuremath{\alpha} condition, and first NONNEGATIVE zero convention.\allowbreak{} Source statements checked, not full moment proof recertified.\allowbreak{} Definition2.\allowbreak{}1 and analytic-strip requirement on p4, plus Theorem3.\allowbreak{}1 closed Fourier support [-2,2] on p5, were reread for NT-Z18; the sinc-squared test and full positivity deduction are independently checked.\allowbreak{}

\url{https://arxiv.org/pdf/math/0208230}.

\paragraph{S-NT-ILS} Iwaniec–\allowbreak{}Luo–\allowbreak{}Sarnak, Low lying zeros of families of L-functions, Publ.\allowbreak{}Math.\allowbreak{}IHES91(2000),55–\allowbreak{}131.

\textbf{Locator:} Global RH assumption printed56; Theorem1.\allowbreak{}1 printed58–\allowbreak{}59; proof Theorem7.\allowbreak{}2.\allowbreak{}.

\textbf{Inspection:} Primary published PDF text inspected.\allowbreak{} The family GL2-RH hypothesis remains separate from auxiliary Dirichlet-GRH; not silently discharged by M.\allowbreak{}

\url{https://www.numdam.org/article/PMIHES_2000__91__55_0.pdf}.

\paragraph{S-NT-RT} Rodgers–\allowbreak{}Tao, The de Bruijn–\allowbreak{}Newman constant is non-negative, arXiv:\allowbreak{}1801.\allowbreak{}05914v5.

\textbf{Locator:} §1 definition/\allowbreak{}equivalence pp1–\allowbreak{}2; Theorem1p3.\allowbreak{}.

\textbf{Inspection:} Primary version5 dated3Jul2021; theorem and classical RH equivalence read.\allowbreak{} Combination with A gives Λ=0, with normalization retained.\allowbreak{}

\url{https://arxiv.org/pdf/1801.05914}.

\paragraph{S-NT-LAGARIAS} Lagarias, An elementary problem equivalent to the Riemann hypothesis, arXiv:\allowbreak{}math/\allowbreak{}0008177v2.

\textbf{Locator:} Theorem1.\allowbreak{}1 author p1.\allowbreak{}.

\textbf{Inspection:} Primary theorem statement read, including equality only at n=1; no reproduction of entire classical proof.\allowbreak{}

\url{https://arxiv.org/pdf/math/0008177}.

\paragraph{S-NT-MILNE} Milne, Class Field Theory v4.\allowbreak{}03 (2020).

\textbf{Locator:} ChapterVI Corollary2.\allowbreak{}12 and Example2.\allowbreak{}13, printed191/\allowbreak{}PDF200; V3.\allowbreak{}5–\allowbreak{}3.\allowbreak{}6 printed158/\allowbreak{}PDF167, Example3.\allowbreak{}9 printed159/\allowbreak{}PDF168; conductor–\allowbreak{}discriminant V3.\allowbreak{}27 printed165/\allowbreak{}PDF174.\allowbreak{}.

\textbf{Inspection:} Primary formula read for residue, full regulator, real-quadratic normalization, roots-of-unity factor.\allowbreak{} Source is author-hosted course notes.\allowbreak{} Ray-class existence/\allowbreak{}reciprocity and Hilbert-class examples read; the V3.\allowbreak{}27 proof is explicitly omitted there, so Neukirch is used for a primary supplied proof locator.\allowbreak{}

\url{https://www.jmilne.org/math/CourseNotes/CFT.pdf}.

\paragraph{S-NT-JOHNSTON} Johnston, Improving bounds on prime counting functions by partial verification of the Riemann hypothesis, RamanujanJ59(2022),1307–\allowbreak{}1321.

\textbf{Locator:} Theorem3.\allowbreak{}2 equations(3.\allowbreak{}3)–\allowbreak{}(3.\allowbreak{}6); §2 paragraph after(2.\allowbreak{}1), ordinary-count normalization.\allowbreak{}.

\textbf{Inspection:} Primary full-text theorem and normalization statement read.\allowbreak{} Under A the finite-height hypothesis can be met for every x, yielding \ensuremath{\psi} threshold59 and \ensuremath{\pi} threshold2657.\allowbreak{} This supplies an exact directly checked alternative proof locator to inaccessible Schoenfeld Corollary1.\allowbreak{}

\url{https://link.springer.com/article/10.1007/s11139-022-00616-x}.

\paragraph{S-NT-RV} Ramakrishnan–\allowbreak{}Valenza, Fourier Analysis on Number Fields, GTM186(1999).

\textbf{Locator:} §7.\allowbreak{}4 Theorem7-19 printed279; §7.\allowbreak{}7 Theorem7-28 printed289.\allowbreak{}.

\textbf{Inspection:} Original PDF downloaded and relevant full statements read directly in extracted text:\allowbreak{} completed unitary-Hecke meromorphy/\allowbreak{}functional equation and pure-norm pole exception; nonvanishing on Re s=1 for every unitary idele-class character.\allowbreak{} Web fetch exceeds size limit; local download succeeded.\allowbreak{} No page-image check or whole-book proof check claimed.\allowbreak{}

\url{https://webusers.imj-prg.fr/~julien.marche/Ref_GL_GT/Fourier_Analysis_on_Number_Fields.pdf}.

\paragraph{S-NT-FPS} Fiorilli–\allowbreak{}Parks–\allowbreak{}Södergren, Low-lying zeros of quadratic Dirichlet L-functions:\allowbreak{} Lower order terms for extended support, CompositioMath153(2017),1196–\allowbreak{}1216; arXiv1601.\allowbreak{}06833.

\textbf{Locator:} Definitions(1.\allowbreak{}1)–\allowbreak{}(1.\allowbreak{}2) and Theorem1.\allowbreak{}1 authorp2; Lemma2.\allowbreak{}5p7; Theorem3.\allowbreak{}5p12; final proof of Theorem1.\allowbreak{}1p14.\allowbreak{}.

\textbf{Inspection:} Primary PDF downloaded; exact family \ensuremath{\chi}\_\allowbreak{}\{8d\} for odd squarefree d of both signs, even nonnegative Schwartz weights, Fourier convention and support (\ensuremath{-}2,2) read.\allowbreak{} Checked Dirichlet-only GRH in the prime-sum step.\allowbreak{} Full literature proof not independently revalidated; the new nonvanishing deduction in section.\allowbreak{}tex is complete original prose.\allowbreak{}

\url{https://arxiv.org/pdf/1601.06833}.

\paragraph{S-NT-TV} Tsfasman–\allowbreak{}Vlăduţ, Infinite Global Fields and the Generalized Brauer–\allowbreak{}Siegel Theorem, MoscowMathJ2(2002),329–\allowbreak{}402; arXivmath/\allowbreak{}0205129v1.

\textbf{Locator:} Definition2.\allowbreak{}1 and Lemma2.\allowbreak{}2 authorp15; GRH Theorem3.\allowbreak{}1p19, proofpp19–\allowbreak{}21; GRH Theorem7.\allowbreak{}2p44, proofpp44–\allowbreak{}49; Theorem7.\allowbreak{}3p44 gives a distinct unconditional almost-normal-tower result.\allowbreak{}.

\textbf{Inspection:} Original PDF downloaded; p19 basic-inequality constants and p44 generalized Brauer–\allowbreak{}Siegel statement inspected as rendered images.\allowbreak{} Exact-family definition, Dedekind-only RH and absence of normality/\allowbreak{}tower assumptions in Thm7.\allowbreak{}2 checked.\allowbreak{} Full literature proofs not independently revalidated.\allowbreak{} Our derived splitting/\allowbreak{}tower arguments are complete in section.\allowbreak{}tex.\allowbreak{}

\url{https://arxiv.org/pdf/math/0205129}.

\paragraph{S-NT-RS} Rubinstein–\allowbreak{}Sarnak, Chebyshev’s Bias, ExperimentalMath3(1994),173–\allowbreak{}197.

\textbf{Locator:} Theorem1.\allowbreak{}1p174, proof§2.\allowbreak{}1; Theorem1.\allowbreak{}2p175, proof§2.\allowbreak{}2; GSH/\allowbreak{}formula(1.\allowbreak{}2)/\allowbreak{}Remark1.\allowbreak{}3p176; Theorems1.\allowbreak{}4/\allowbreak{}1.\allowbreak{}5p177, proofsProposition3.\allowbreak{}1/\allowbreak{}§3.\allowbreak{}2.\allowbreak{}.

\textbf{Inspection:} Author-hosted original PDF downloaded; text encoding is corrupt, so printedpp174–\allowbreak{}178 were rendered and inspected visually.\allowbreak{} Checked exact \ensuremath{\pi}(x) normalization, continuous-test versus ordering-indicator distinction, lower logarithmic density under GRH alone, GSH collection including central zeros, law-symmetry classification and fixed-r large-q quantifier.\allowbreak{} Full classical proofs not independently reread.\allowbreak{}

\url{https://www.math.uwaterloo.ca/~mrubinst/publications/Chebyshev.pdf}.

\paragraph{S-NT-CCM} Carneiro–\allowbreak{}Chandee–\allowbreak{}Milinovich, A note on the zeros of zeta and L-functions, Math.\allowbreak{}Z.\allowbreak{}281(2015),315–\allowbreak{}332; arXiv1503.\allowbreak{}00955v1.

\textbf{Locator:} Theorem1 authorp2, proof§3 pp3–\allowbreak{}5; §4 axioms p6; Theorem5 equations(4.\allowbreak{}4),(4.\allowbreak{}5) p7, proofpp7–\allowbreak{}12; Theorem7 and proofp13; Examples8–\allowbreak{}10p14.\allowbreak{}.

\textbf{Inspection:} Original PDF downloaded.\allowbreak{} Axioms p6 and full Theorem5 p7 rendered and visually inspected.\allowbreak{} Read Theorem1, exact conductor normalization and universal-constant assertion, Theorem7 and its entire p13 proof.\allowbreak{} Critical restriction:\allowbreak{} gamma parameters real or in conjugate pairs.\allowbreak{} Applied to finite-order Hecke characters with explicit axiom verification; did not assume this applies verbatim to arbitrary unitary infinity types.\allowbreak{} Full Theorem5 classical proof not independently recertified; NT-Z14 deduction is proved and checked in full.\allowbreak{} NT-Z17 subsequently verifies every axiom for the conjugate shifted product of any non-pure-norm unitary Hecke function; its degree2n and squared conductor preserve the exact universal constants.\allowbreak{} This is a new transfer proof, not a claim that an unpaired individual completion directly satisfies the source axioms.\allowbreak{}

\url{https://arxiv.org/pdf/1503.00955}.

\paragraph{S-NT-CMS} Carneiro–\allowbreak{}Milinovich–\allowbreak{}Soundararajan, Fourier optimization and prime gaps, Comment.\allowbreak{}Math.\allowbreak{}Helv.\allowbreak{}94(2019),533–\allowbreak{}568; arXiv1708.\allowbreak{}04122v2.

\textbf{Locator:} Theorem3, Corollary4, Theorem5 authorp5; proofs§5 and §6 pp21–\allowbreak{}26.\allowbreak{}.

\textbf{Inspection:} Downloaded original author v2.\allowbreak{} Page5 rendered and visually inspected for strict21/\allowbreak{}25 limsup, exact22/\allowbreak{}25 all-x interval, x\ensuremath{\ge}4 and p\_\allowbreak{}n>3.\allowbreak{} Full long Fourier/\allowbreak{}explicit-estimate proof not independently recertified; numerical routines not rerun.\allowbreak{}

\url{https://arxiv.org/pdf/1708.04122v2}.

\paragraph{S-NT-CMQR} Carneiro–\allowbreak{}Milinovich–\allowbreak{}Quesada-Herrera–\allowbreak{}Ramos, Fourier optimization, the least quadratic non-residue, and the least prime in an arithmetic progression; arXiv2404.\allowbreak{}08380v2 (2025), Math.\allowbreak{}Comp.\allowbreak{} DOI10.\allowbreak{}1090/\allowbreak{}mcom/\allowbreak{}4154.

\textbf{Locator:} Theorems1/\allowbreak{}2p2, Theorem3p3, definition(1.\allowbreak{}8)/\allowbreak{}Theorem6p5, Corollary7p6; proof§4pp13–\allowbreak{}17, §4.\allowbreak{}4pp16–\allowbreak{}17, §5pp18–\allowbreak{}22; numerical proof§6.\allowbreak{}.

\textbf{Inspection:} Downloaded original v2 (dated August2025).\allowbreak{} Theorem pages2,3,6 rendered and inspected.\allowbreak{} Checked exact least-prime-residue convention, character zeros excluded, fixed character order, 0.\allowbreak{}7615/\allowbreak{}0.\allowbreak{}6707/\allowbreak{}0.\allowbreak{}6131/\allowbreak{}0.\allowbreak{}5765 constants, \ensuremath{\ell}\ensuremath{\ge}6 formula and strict8/\allowbreak{}9 AP asymptotic; proof conclusion(5.\allowbreak{}15)–\allowbreak{}(5.\allowbreak{}16) checked for a-uniformity.\allowbreak{} The paper’s computer-assisted optimization proof is cited, not rerun.\allowbreak{}

\url{https://arxiv.org/pdf/2404.08380v2}.

\paragraph{S-NT-BFK} Bucur–\allowbreak{}Fité–\allowbreak{}Kedlaya, Effective Sato–\allowbreak{}Tate conjecture for abelian varieties and applications, JEMS26(2024),1713–\allowbreak{}1746; arXiv2002.\allowbreak{}08807v2.

\textbf{Locator:} Prop2.\allowbreak{}5(2.\allowbreak{}2)/\allowbreak{}(2.\allowbreak{}3), Remarks2.\allowbreak{}6–\allowbreak{}2.\allowbreak{}8 printed1720–\allowbreak{}1721 (authorpp6–\allowbreak{}7); §4.\allowbreak{}3, Lemma4.\allowbreak{}6, Prop4.\allowbreak{}7 and footnote2 printed1739–\allowbreak{}1740 (authorp20); Prop4.\allowbreak{}9 and Cor4.\allowbreak{}10 printed1741–\allowbreak{}1742 (authorpp21–\allowbreak{}22); Remark4.\allowbreak{}11 rational two-CM-field obstruction printed1743 (authorp22).\allowbreak{} Author URL https:\allowbreak{}/\allowbreak{}/\allowbreak{}arxiv.\allowbreak{}org/\allowbreak{}pdf/\allowbreak{}2002.\allowbreak{}08807.\allowbreak{}.

\textbf{Inspection:} Primary author PDF downloaded; original p6 and p20 rendered and visually read, including exact uniform weight dependence and explicit identification of the Hecke-power RH input.\allowbreak{} Entire §2 source discussion and entire Prop4.\allowbreak{}7 supplied proof read; Prop4.\allowbreak{}9 and Cor4.\allowbreak{}10 proofs read in extracted/\allowbreak{}published primary text.\allowbreak{} The long upstream Murty prime-sum proof remains a cited classical premise, not independently re-refereed.\allowbreak{} Our proof uses Prop2.\allowbreak{}5 plus original Erdős–\allowbreak{}Turán, not the general smoothing theorem.\allowbreak{} Standard Hecke completion degree [k:\allowbreak{}Q] is explicit in our specialization.\allowbreak{} No general non-CM analytic axioms discharged by M.\allowbreak{} At05:\allowbreak{}47UTC the entire short Remark4.\allowbreak{}11 obstruction was read in both author and publisher primary texts, confirming attribution to Andrew Sutherland; NT-E12/\allowbreak{}E13 now credit that earlier rational positive-extreme special case.\allowbreak{}

\url{https://ems.press/content/serial-article-files/47532?nt=1}.

\paragraph{S-NT-ET} Erdős–\allowbreak{}Turán, On a problem in the theory of uniform distribution.\allowbreak{}I, Indagationes Math.\allowbreak{}10(1948),370–\allowbreak{}378.

\textbf{Locator:} TheoremIII printed374, parallel Proceedings pagination1150, PDF5; proof continues in partII,406–\allowbreak{}413, https:\allowbreak{}/\allowbreak{}/\allowbreak{}www.\allowbreak{}renyi.\allowbreak{}hu/\allowbreak{}\textasciitilde{}p\_\allowbreak{}erdos/\allowbreak{}1948-03.\allowbreak{}pdf.\allowbreak{}.

\textbf{Inspection:} Original primary scans I and II downloaded.\allowbreak{} Original printed374 rendered and visually inspected in full, confirming arbitrary finite point sets, every positive integer truncation, numerical constant, 1/\allowbreak{}m harmonic weights and closed endpoints.\allowbreak{} Full classical Fourier-approximation proof in partII is cited, not independently reread.\allowbreak{} All endpoint and CM specializations are fully proved in our own text.\allowbreak{}

\url{https://users.renyi.hu/~p_erdos/1948-02.pdf}.

\paragraph{S-NT-JP} James–\allowbreak{}Pollack, Extremal primes for elliptic curves with complex multiplication, J.\allowbreak{}NumberTheory172(2017),383–\allowbreak{}391, with author erratum.

\textbf{Locator:} Theorem1 authorp2; comparison analytic Proposition2 authorp3.\allowbreak{} Erratum https:\allowbreak{}/\allowbreak{}/\allowbreak{}www.\allowbreak{}pollack-math.\allowbreak{}net/\allowbreak{}extremal-errata.\allowbreak{}pdf, single page.\allowbreak{}.

\textbf{Inspection:} Primary author Theorem1 and whole separate erratum read.\allowbreak{} Theorem states the unconditional Q asymptotic for each extremal sign with constant2/\allowbreak{}(3\ensuremath{\pi}).\allowbreak{} Erratum explicitly replaces the narrow-sector input after a gap in the earlier Maknys argument.\allowbreak{} This run does not independently certify that revised analytic input or re-referee the full unconditional proof.\allowbreak{} Source is comparison only; NT-E08 has its own conditional proof from a different cited premise.\allowbreak{}

\url{https://www.pollack-math.net/extremal.pdf}.

\paragraph{S-NT-RAJAN} C.\allowbreak{} S.\allowbreak{} Rajan, Distribution of values of Hecke characters of infinite order, Acta Arithmetica85(1998),279–\allowbreak{}291.

\textbf{Locator:} §1 character/\allowbreak{}gamma/\allowbreak{}conductor conventions and unit-interval zero count(2) printed281–\allowbreak{}283; Theorem1 and normalized-angle convention printed286; conditional Proposition4 and entire supplied proof outline printed289–\allowbreak{}290.\allowbreak{} DOI10.\allowbreak{}4064/\allowbreak{}aa-85-3-279-291.\allowbreak{}.

\textbf{Inspection:} Primary publisher PDF full indexed text recovered through search reference after direct fetch403.\allowbreak{} Read §1 analytic conventions/\allowbreak{}conductor-power behavior and exact local-zero-count statement; Theorem1 full character restriction; complete Prop4 statement and supplied proof outline.\allowbreak{} Normalized arg is modulo1; Dirichlet means finite-order Hecke, and equivalence is modulo norm twists.\allowbreak{} Long earlier explicit-formula arguments and cited Murty proof are not independently re-refereed.\allowbreak{} Local PDF downloads failed and screenshot tool returned only placeholders, so no fresh image inspection claimed.\allowbreak{} NT-A04/\allowbreak{}05 new deductions are fully written/\allowbreak{}read.\allowbreak{}

\url{https://matwbn.icm.edu.pl/ksiazki/aa/aa85/aa8537.pdf}.

\paragraph{S-NT-ANS} Akbary–\allowbreak{}Ng–\allowbreak{}Shahabi, Limiting distributions of the classical error terms of prime number theory, arXiv1306.\allowbreak{}1657v1 (2013).

\textbf{Locator:} Introduction authorp6, discussion immediately after(1.\allowbreak{}15), cites Titchmarsh p377(14.\allowbreak{}29.\allowbreak{}4) for weak-Mertens reciprocal derivative control; their stronger limiting-distribution moment assumptions are separate.\allowbreak{}.

\textbf{Inspection:} Primary author PDF text read for exact weak-Mertens versus RH/\allowbreak{}simplicity distinction.\allowbreak{} Used as corroboration, not as source of a new theorem.\allowbreak{} No full read of its limiting-distribution proof claimed.\allowbreak{}

\url{https://www.cs.uleth.ca/~nathanng/RESEARCH/AkbaryNgShahabi.pdf}.

\paragraph{S-NT-EV} Ellenberg–\allowbreak{}Venkatesh, Reflection principles and bounds for class group torsion, IMRN2007 rnm002, author preprint.

\textbf{Locator:} Author preprint Lemma3 and complete proofp6; conditional Proposition1 and complete short proofp7.\allowbreak{} Height input Lemma2p5; Arakelov definitionspp4–\allowbreak{}5.\allowbreak{} Published numbering differs from the author preprint.\allowbreak{}.

\textbf{Inspection:} Downloaded original 12-page author PDF from identical author mirror https:\allowbreak{}/\allowbreak{}/\allowbreak{}virtualmath1.\allowbreak{}stanford.\allowbreak{}edu/\allowbreak{}\textasciitilde{}akshay/\allowbreak{}research/\allowbreak{}sch.\allowbreak{}pdf after first host timed out.\allowbreak{} Rendered p6 and p7 and visually read exact \ensuremath{\delta}<1/\allowbreak{}[2\ensuremath{\ell}(n\ensuremath{-}1)], non-extended-prime condition, torsion bound and full supplied Lemma3/\allowbreak{}Proposition1 proofs.\allowbreak{} Earlier Arakelov/\allowbreak{}height machinery is cited, not globally recertified.\allowbreak{} NT-H12 supplies and checks the independent Dedekind-only degree-one-prime reduction in full.\allowbreak{}

\url{https://math.stanford.edu/~akshay/research/sch.pdf}.

\paragraph{S-NT-STACKS-KUMMER} Stacks Project, §59.\allowbreak{}28 Kummer theory.

\textbf{Locator:} Lemma59.\allowbreak{}28.\allowbreak{}1 and subsequent long exact étale cohomology sequence; Lemma59.\allowbreak{}28.\allowbreak{}5, Tag040Q, line bundles with power trivialization.\allowbreak{}.

\textbf{Inspection:} Read the complete Lemma59.\allowbreak{}28.\allowbreak{}1 proof and following cohomology sequence in primary HTML.\allowbreak{} Read Lemma59.\allowbreak{}28.\allowbreak{}5 statement and distinction H¹ versus Picard torsion; full descent proof of59.\allowbreak{}28.\allowbreak{}5 not independently recertified.\allowbreak{} Invertibility of \ensuremath{\ell} on U is explicit.\allowbreak{}

\url{https://stacks.math.columbia.edu/tag/03PK}.

\paragraph{S-NT-MILNE-ANT} Milne, Algebraic Number Theory, version3.\allowbreak{}08 (19July2020).

\textbf{Locator:} Remark4.\allowbreak{}11 and infinite-place footnote printed72/\allowbreak{}PDF74; ordinary unit Thm5.\allowbreak{}1printed85/\allowbreak{}PDF87; S-unit Thm5.\allowbreak{}11 and complete proof printed90/\allowbreak{}PDF92; Hermite Thm8.\allowbreak{}43printed151–\allowbreak{}152/\allowbreak{}PDF153–\allowbreak{}154.\allowbreak{} Minkowski Thm4.\allowbreak{}3printed70 and Thm4.\allowbreak{}9proof72.\allowbreak{} Theorem8.\allowbreak{}42printed151 and Theorem8.\allowbreak{}43printed151–\allowbreak{}152, complete supplied finiteness proofs, freshly reread for NT-E14–\allowbreak{}16.\allowbreak{}.

\textbf{Inspection:} Primary text read for exact ordinary Hilbert-class-field real splitting convention, ordinary/\allowbreak{}S-unit ranks, torsion roots of unity and Hermite finiteness.\allowbreak{} Whole S-unit extension proof read.\allowbreak{} The underlying full unit theorem and global class field theory are cited rather than independently rederived.\allowbreak{} Freshly read Thm4.\allowbreak{}3 and the entire short Thm4.\allowbreak{}9 proof for NT-R08; full Minkowski lattice proof is a cited classical premise.\allowbreak{} At06:\allowbreak{}00–\allowbreak{}06:\allowbreak{}01UTC, official v3.\allowbreak{}08 text freshly read for both complete finiteness arguments, with the local-field finiteness and Minkowski inputs remaining classical.\allowbreak{} The finite-family conclusion for fixed bounded degree/\allowbreak{}unramified set is explicit.\allowbreak{}

\url{https://www.jmilne.org/math/CourseNotes/ANT.pdf}.

\paragraph{S-NT-CS} Campagna–\allowbreak{}Stevenhagen, Cyclic reduction densities for elliptic curves, Res.\allowbreak{}NumberTheory9:\allowbreak{}61(2023); author arXiv2001.\allowbreak{}00028v2(2022).

\textbf{Locator:} Authorv2 Thm1.\allowbreak{}1p3; §2 Lemma2.\allowbreak{}1/\allowbreak{}Cor2.\allowbreak{}2/\allowbreak{}Props2.\allowbreak{}3–\allowbreak{}2.\allowbreak{}4pp6–\allowbreak{}9; factorization proofpp9–\allowbreak{}10; §5/\allowbreak{}Thm5.\allowbreak{}1pp14–\allowbreak{}16.\allowbreak{} Published HTML same theorem numbering.\allowbreak{}.

\textbf{Inspection:} Primary published HTML and downloaded authorv2 read for all CM/\allowbreak{}nonCM/\allowbreak{}base-field hypotheses, density series, exact critical primes, finite entanglement and vanishing.\allowbreak{} Original authorp3 andp8 rendered and visually read.\allowbreak{} Entire supplied §2 argument read; underlying Serre/\allowbreak{}open-image and CM degree-growth theorems remain cited, not independently re-refereed.\allowbreak{} Thm5.\allowbreak{}1 statement and complete displayed construction read in published HTML.\allowbreak{} PDFdownload https:\allowbreak{}/\allowbreak{}/\allowbreak{}arxiv.\allowbreak{}org/\allowbreak{}pdf/\allowbreak{}2001.\allowbreak{}00028; institutional publishedPDF403.\allowbreak{}

\url{https://link.springer.com/article/10.1007/s40993-023-00463-9}.

\paragraph{S-NT-MURTY-ARTIN} M.\allowbreak{} Ram Murty, On Artin’s conjecture, J.\allowbreak{} Number Theory 16 (1983), 147–\allowbreak{}168.

\textbf{Locator:} Theorem1 printed153/\allowbreak{}PDF7, entire supplied proof153–\allowbreak{}155/\allowbreak{}PDF7–\allowbreak{}9; general-family definitions150–\allowbreak{}153.\allowbreak{}.

\textbf{Inspection:} Author-hosted original scan downloaded.\allowbreak{} Text extraction empty, so original pages148–\allowbreak{}155 rendered and visually read.\allowbreak{} Exact Theorem1 explicitly assumes RH for each squarefree-compositum Dedekind zeta; entire theorem proof153–\allowbreak{}155 read and checked for the analytic-family input.\allowbreak{} Earlier lattice machinery and cited effective Chebotarev remain classical premises.\allowbreak{}

\url{https://mrammurty-ai.github.io/artinconj.pdf}.

\paragraph{S-NT-ZYWINA-LT} Zywina, Bounds for the Lang–\allowbreak{}Trotter conjectures, author2015 preprint arXiv1508.\allowbreak{}07682.

\textbf{Locator:} Thm1.\allowbreak{}2p1; Thm1.\allowbreak{}3/\allowbreak{}Cor1.\allowbreak{}4p2; Dedekind-GRH and AHC definitionsp3; Thm2.\allowbreak{}3p5; mixed representations§3.\allowbreak{}2–\allowbreak{}3.\allowbreak{}3pp7–\allowbreak{}10; entire§4pp10–\allowbreak{}11; §5 algebraic reductionpp11–\allowbreak{}13 and conditional proof§5.\allowbreak{}1pp13–\allowbreak{}14; Cor1.\allowbreak{}4proofp15.\allowbreak{}.

\textbf{Inspection:} Original author PDF downloaded; theorem p1 rendered and visually checked.\allowbreak{} Read full §4 proof, §3.\allowbreak{}3 joint-representation construction, complete algebraic and conditional §5.\allowbreak{}1 proof and entire Cor1.\allowbreak{}4 proof.\allowbreak{} General analytic Chebotarev proof§7 is cited, not independently re-refereed.\allowbreak{} Recorded abelian quotients mean no additional AHC in these elliptic applications.\allowbreak{} Source p1 has a trace-sign typo in point-count formula; our a\_\allowbreak{}p uses the standard characteristic-polynomial convention and bounds are uniform under a\ensuremath{\mapsto}\ensuremath{-}a.\allowbreak{}

\url{https://pi.math.cornell.edu/~zywina/papers/LTMurty.pdf}.

\paragraph{S-NT-MMS} Murty–\allowbreak{}Murty–\allowbreak{}Saradha, Modular Forms and the Chebotarev Density Theorem, Amer.\allowbreak{}J.\allowbreak{}Math110(1988),253–\allowbreak{}281.

\textbf{Locator:} Lemma4.\allowbreak{}4 printed269–\allowbreak{}270/\allowbreak{}PDF18–\allowbreak{}19, complete proofp270; quadratic-field splitting use of Proposition3.\allowbreak{}2printed261.\allowbreak{}.

\textbf{Inspection:} Author-hosted original journal scan downloaded and extracted text read for Lemma4.\allowbreak{}4 and its entire proof.\allowbreak{} Exact prime-selection input is RH for quadratic Dedekind zeta functions and \ensuremath{\zeta}\_\allowbreak{}Q; no general Artin holomorphy used in this lemma.\allowbreak{} Full paper theorem and source page images are not newly recertified.\allowbreak{}

\url{https://mrammurty-ai.github.io/Murty_ModularFormsChebotarevDensity.pdf}.

\paragraph{S-NT-FJ} Fiorilli–\allowbreak{}Jouve, Distribution of Frobenius elements in families of Galois extensions, J.\allowbreak{} Inst.\allowbreak{} Math.\allowbreak{} Jussieu 23 (2024), 1169–\allowbreak{}1258; author arXiv2001.\allowbreak{}05428v2.

\textbf{Locator:} Definitions(9),(10),(13),(15),(16) authorpp9–\allowbreak{}11; cyclic induction Lemmata3.\allowbreak{}14–\allowbreak{}3.\allowbreak{}15p30; Cor3.\allowbreak{}17pp31–\allowbreak{}32; Proposition3.\allowbreak{}18pp32–\allowbreak{}34 and entire supplied proof.\allowbreak{} Published DOI10.\allowbreak{}1017/\allowbreak{}S1474748023000154.\allowbreak{} Comparison:\allowbreak{} Prop4.\allowbreak{}6(87)authorp40 and its no-AHC variance proofp41.\allowbreak{}.

\textbf{Inspection:} Read the exact assumptions§1.\allowbreak{}2pp5–\allowbreak{}7, definitionspp9–\allowbreak{}11, all cyclic-induction proofs and the full Prop3.\allowbreak{}18 proofpp30–\allowbreak{}34 in the primary author PDF.\allowbreak{} Rendered and visually inspected original p32 statement and mean/\allowbreak{}variance formulas.\allowbreak{} The non-AHC branch reduces to cyclic Hecke functions; RH \ensuremath{\zeta}\_\allowbreak{}L supplies the needed zeros.\allowbreak{} Source analytic approximation cites Ng/\allowbreak{}ANS/\allowbreak{}Devin; their upstream proofs were not independently rebuilt.\allowbreak{} Our joint-law and finite-étale deductions are fully written and read.\allowbreak{} No Gaussian/\allowbreak{}LI claim; coefficients retain multiplicities.\allowbreak{} Read the entire proof of inequality(87)pp40–\allowbreak{}41 and image-checked Theorem2.\allowbreak{}1p11.\allowbreak{} The latter author-version equation(17) lacks the square on the L1 norm; Prop4.\allowbreak{}6 and its proof have the homogeneous squared norm.\allowbreak{} NT-R08 is independently derived from GM and does not use the misprinted scaling.\allowbreak{}

\url{https://arxiv.org/pdf/2001.05428}.

\paragraph{S-NT-RW} Rudnick–\allowbreak{}Waxman, Angles of Gaussian primes, arXiv 1705.\allowbreak{}07498 v2 (29 September 2018).

\textbf{Locator:} Theorem1.\allowbreak{}1authorp3; smooth definitions§3.\allowbreak{}1pp10–\allowbreak{}11; Lemma3.\allowbreak{}1pp11–\allowbreak{}12; Theorem3.\allowbreak{}2p12 and complete proof§4.\allowbreak{}4p20; Hecke family/\allowbreak{}completion§4.\allowbreak{}1p13; Cor4.\allowbreak{}6and proofpp16–\allowbreak{}17; Lemmas4.\allowbreak{}7–\allowbreak{}4.\allowbreak{}8pp17–\allowbreak{}20.\allowbreak{}.

\textbf{Inspection:} Downloaded original v2 and visually inspected authorp12 andp20.\allowbreak{} Read complete mean proof, Cor4.\allowbreak{}6 proof, Lemmas4.\allowbreak{}7–\allowbreak{}4.\allowbreak{}8 and the entire supplied Theorem3.\allowbreak{}2 deduction, tracking the nonzero Hecke-power RH family and the normalized Haar measure.\allowbreak{} The upstream general local-zero-count/\allowbreak{}explicit-formula machinery remains a cited premise; no full-paper re-referee claim.\allowbreak{} New deductions use only K\ensuremath{\le}X\ensuremath{\le}K²; numerical implied constant not extracted.\allowbreak{} The conjectural variance asymptotic is not used.\allowbreak{}

\url{https://arxiv.org/pdf/1705.07498}.

\paragraph{S-NT-IK} Iwaniec–\allowbreak{}Kowalski, Analytic Number Theory, AMS Colloquium Publications53 (2004).

\textbf{Locator:} §5.\allowbreak{}1printed93–\allowbreak{}95:\allowbreak{} analytic axioms, real-or-conjugate-paired gamma parameters, analytic conductor(5.\allowbreak{}6)–\allowbreak{}(5.\allowbreak{}7); Theorem5.\allowbreak{}15printed114, equation(5.\allowbreak{}56), with its preceding proof sentence using(5.\allowbreak{}53)pp112–\allowbreak{}113 and local zero count(5.\allowbreak{}26).\allowbreak{}.

\textbf{Inspection:} Read the primary book text in the publicly indexed scan at https:\allowbreak{}/\allowbreak{}/\allowbreak{}www.\allowbreak{}scribd.\allowbreak{}com/\allowbreak{}document/\allowbreak{}807332382/\allowbreak{}Colloquium-Publications-53-Henryk-Iwaniec-Emmanuel-Kowalski-Analytic-Number-Theory-American-Mathematical-Society-2004 :\allowbreak{} entire §5.\allowbreak{}1 definition and the complete supplied short Theorem5.\allowbreak{}15 proof sentence.\allowbreak{} The earlier explicit-formula development/\allowbreak{}exercise is a cited premise, not independently rederived.\allowbreak{} No page-image inspection or full-book proof audit claimed.\allowbreak{} Checked author correction list at https:\allowbreak{}/\allowbreak{}/\allowbreak{}people.\allowbreak{}math.\allowbreak{}ethz.\allowbreak{}ch/\allowbreak{}\textasciitilde{}kowalski/\allowbreak{}corrections-ant.\allowbreak{}pdf ; no Theorem5.\allowbreak{}15 correction listed there.\allowbreak{} Original publisher DOI identifies the book; direct AMS access returned403.\allowbreak{} Uniform degree/\allowbreak{}conductor normalization additionally cross-checked against de Boer2022 thesis §4.\allowbreak{}3.\allowbreak{}3pp143–\allowbreak{}147, Prop4.\allowbreak{}6/\allowbreak{}Lemma4.\allowbreak{}8 and equation(4.\allowbreak{}48), https:\allowbreak{}/\allowbreak{}/\allowbreak{}ir.\allowbreak{}cwi.\allowbreak{}nl/\allowbreak{}pub/\allowbreak{}32136/\allowbreak{}32136D.\allowbreak{}pdf; that special Arakelov family is not substituted for arbitrary infinity types.\allowbreak{} No numerical implied constant extracted.\allowbreak{}

\url{https://doi.org/10.1090/coll/053}.

\paragraph{S-NT-GM2025-CG} Grenié–\allowbreak{}Molteni, Breaking the 4 barrier for the bound of a generating set of the class group, Math.\allowbreak{} Comp.\allowbreak{}94 (2025),3145–\allowbreak{}3163.

\textbf{Locator:} Author version dated1August2025:\allowbreak{} Theorems1–\allowbreak{}3/\allowbreak{}Table1–\allowbreak{}2p2; Theorem4p3; explicitformula(3)p3; §2pp3–\allowbreak{}7; Proposition1pp8–\allowbreak{}9; Proposition2pp9–\allowbreak{}10; §4proofThms1/\allowbreak{}2pp10–\allowbreak{}12; §5proofThms3/\allowbreak{}4pp12–\allowbreak{}13; Theorem5(20)andcompleteproofpp13–\allowbreak{}14; Corollaries1–\allowbreak{}2pp14–\allowbreak{}15.\allowbreak{} DOI10.\allowbreak{}1090/\allowbreak{}mcom/\allowbreak{}4114.\allowbreak{}.

\textbf{Inspection:} Downloaded original author PDF.\allowbreak{} Actual p2,p3,p13 images inspected; complete analytic derivations §2–\allowbreak{}§5 and Theorem5 proof read.\allowbreak{} Explicitly retained RH \ensuremath{\zeta}\_\allowbreak{}K and class characters from zero sums and RH \ensuremath{\zeta}\_\allowbreak{}Q in Proposition1.\allowbreak{} Threshold computations, Odlyzko bounds and the zeta-zero table inputs are cited published premises, not numerically rerun.\allowbreak{} Introductory T\ensuremath{{}^{\prime}} wording is interpreted by the precise §5 conclusion that the generated subgroup contains Cl².\allowbreak{} Our certificate states support in logarithmic coordinates and counts every contributing prime power.\allowbreak{}

\url{https://sites.unimi.it/molteni/research/papers-pdf/53-molteni-Breaking%20the%204%20barrier%20for%20the%20bound%20of%20a%20generating%20set%20of%20the%20class%20group.pdf}.

\paragraph{S-NT-BDDF} Belabas–\allowbreak{}Diaz y Diaz–\allowbreak{}Friedman, Small generators of the ideal class group, Math.\allowbreak{} Comp.\allowbreak{}77 (2008),1185–\allowbreak{}1197; author preprint5December2005.

\textbf{Locator:} Author PDF15pp:\allowbreak{} Theorem1.\allowbreak{}1p3/\allowbreak{}(2), Theorem2.\allowbreak{}1pp5–\allowbreak{}6/\allowbreak{}(5) and complete proof, Corollary2.\allowbreak{}2pp6–\allowbreak{}7/\allowbreak{}(10)–\allowbreak{}(13); input representation§3pp7–\allowbreak{}8.\allowbreak{}.

\textbf{Inspection:} Downloaded author primary PDF and inspected actual p3 and p7 images.\allowbreak{} Read complete Theorem2.\allowbreak{}1 proof and Corollary2.\allowbreak{}2 proof; exact constants \ensuremath{\pi}²/\allowbreak{}2 and4 Catalan and strict prime cutoff retained.\allowbreak{} Only nontrivial ideal-class Hecke RH is assumed in this certificate.\allowbreak{} Historical computational-performance observations are not promoted to present algorithmic guarantees; no experiments rerun.\allowbreak{}

\url{https://www.math.u-bordeaux.fr/~kbelabas/research/OnBach.pdf}.

\paragraph{S-NT-GM2018-CG} Grenié–\allowbreak{}Molteni, Explicit bounds for generators of the class group, Math.\allowbreak{} Comp.\allowbreak{}87 (2018),2483–\allowbreak{}2511.

\textbf{Locator:} arXivv5(12April2017):\allowbreak{} introductionp1 includes RH \ensuremath{\zeta}\_\allowbreak{}K; Theorem2.\allowbreak{}4and explicitformula(2.\allowbreak{}2)authorp3; Proposition3.\allowbreak{}1pp4–\allowbreak{}5; Theorem3.\allowbreak{}5(3.\allowbreak{}9)–\allowbreak{}(3.\allowbreak{}10)andproofauthorp6.\allowbreak{} DOI10.\allowbreak{}1090/\allowbreak{}mcom/\allowbreak{}3281.\allowbreak{}.

\textbf{Inspection:} Freshly read exact small-discriminant premise Theorem3.\allowbreak{}5 and its supplied proof, with §2 class-character/\allowbreak{}Dedekind RH scope.\allowbreak{} This is an upstream premise of GM2025, not a separately re-audited numerical bound; no original page-image or whole-paper audit claimed.\allowbreak{}

\url{https://arxiv.org/pdf/1607.02430}.

\paragraph{S-NT-DP} Dick–\allowbreak{}Pillichshammer, Digital Nets and Sequences, author preprint of Cambridge University Press (2010) monograph.

\textbf{Locator:} Author preprint printed/\allowbreak{}PDFp72 Definition3.\allowbreak{}13 and explicit multiset convention; Theorem3.\allowbreak{}26 printed/\allowbreak{}PDFp84, Erdős–\allowbreak{}Turán–\allowbreak{}Koksma inequality.\allowbreak{} Its displayed proof reference is Drmota–\allowbreak{}Tichy, Sequences, Discrepancies and Applications (1997), §1.\allowbreak{}2.\allowbreak{}2; publisher DOI10.\allowbreak{}1007/\allowbreak{}BFb0093404.\allowbreak{} These are author-preprint page numbers, not asserted to match the published book pagination.\allowbreak{}.

\textbf{Inspection:} Downloaded the author-hosted627-page PDF, read the full discrepancy definition/\allowbreak{}multiset convention onp72 and the exact Theorem3.\allowbreak{}26 statement onp84.\allowbreak{} Locally rendered and actually viewed p84:\allowbreak{} factor(3/\allowbreak{}2)\textasciicircum{}r, cutoff2/\allowbreak{}(H+1), frequency cube and product max(1,|m\_\allowbreak{}i|) weights checked.\allowbreak{} The cited Drmota–\allowbreak{}Tichy proof is a classical premise, not freshly read; no whole-book proof audit claimed.\allowbreak{} New arithmetic tensor-family, shrinking-box, least-prime and monomial-obstruction deductions are written and whole-read separately.\allowbreak{}

\url{https://web.maths.unsw.edu.au/~josefdick/preprints/DP_book_preprint.pdf}.

\paragraph{S-NT-LENSTRA-IDEAL} H.\allowbreak{} W.\allowbreak{} Lenstra Jr.\allowbreak{}, Euclidean ideal classes, Astérisque 61 (1979), 121–\allowbreak{}131.

\textbf{Locator:} Definition(1.\allowbreak{}1)/\allowbreak{}(1.\allowbreak{}2)printed123; Lemma(1.\allowbreak{}5), Theorem(1.\allowbreak{}6) and fixed-function uniqueness printed124–\allowbreak{}125; Proposition(2.\allowbreak{}2)printed127/\allowbreak{}PDF8.\allowbreak{} Exact input:\allowbreak{} |S|\ensuremath{\ge}2 and RH \ensuremath{\zeta}\_\allowbreak{}\{K(\ensuremath{\zeta}\_\allowbreak{}m,(O\_\allowbreak{}\{K,S\}×)\textasciicircum{}(1/\allowbreak{}m))\} for every positive squarefree m.\allowbreak{}.

\textbf{Inspection:} Downloaded the original12-page Numdam PDF.\allowbreak{} Printedp127 locally rendered and actually viewed:\allowbreak{} confirmed \ensuremath{\ge}2 against misleading searchable OCR >2, exact squarefree unit-Kummer family and conclusion.\allowbreak{} Read all elementary proofs of Lemma1.\allowbreak{}5/\allowbreak{}Theorem1.\allowbreak{}6 and the fixed-function uniqueness paragraph onpp124–\allowbreak{}125, plus the precise definition.\allowbreak{} Proposition2.\allowbreak{}2 cites methods of Lenstra1977 rather than printing the analytic sufficiency proof; that classical input is not claimed freshly reproved or fully audited here.\allowbreak{} Complete new elementary reformulation and bounded one-prime localization deduction are written separately.\allowbreak{}

\url{https://www.numdam.org/item/AST_1979__61__121_0.pdf}.

\paragraph{S-NT-MAYLE-WANG} Jacob Mayle and Tian Wang, On the effective version of Serre’s open image theorem, Bull.\allowbreak{} London Math.\allowbreak{} Soc.\allowbreak{} 56 (2024), 1399–\allowbreak{}1416.

\textbf{Locator:} Theorems1–\allowbreak{}2 printed1401; Corollary6 and proof1404–\allowbreak{}1405; discriminant Lemma7/\allowbreak{}proof1405; Proposition12/\allowbreak{}full proof1407–\allowbreak{}1408; entire Theorem2 proof1408–\allowbreak{}1410, equations13–\allowbreak{}16; Lemmas14–\allowbreak{}15 and entire Theorem1 proof1411–\allowbreak{}1413.\allowbreak{} Publisher PDF mirror https:\allowbreak{}/\allowbreak{}/\allowbreak{}pure.\allowbreak{}mpg.\allowbreak{}de/\allowbreak{}rest/\allowbreak{}items/\allowbreak{}item\_\allowbreak{}3579985\_\allowbreak{}3/\allowbreak{}component/\allowbreak{}file\_\allowbreak{}3598865/\allowbreak{}content ; author arXiv2109.\allowbreak{}08656v3(25January2024), Theorems1–\allowbreak{}2p2.\allowbreak{} NOS Theorem8p1406 and Corollary9pp1406–\allowbreak{}1407, Remark11p1407; Faltings criterion(11)p1407.\allowbreak{}.

\textbf{Inspection:} Obtained the original18-page publisher PDF; printed1400 and1401 locally rendered and viewed, confirming the Dedekind-GRH convention and exact964/\allowbreak{}5760 and482/\allowbreak{}2880 constants.\allowbreak{} Read all supplied mathematical proofs in §3,§4 and§5.\allowbreak{}1–\allowbreak{}5.\allowbreak{}2, including the actual A/\allowbreak{}lA finite quotient, degree192→384, quadratic avoidance, discriminant calculation and twist argument.\allowbreak{} Exact sufficient RH family identified from these proofs.\allowbreak{} Upstream Bach–\allowbreak{}Sorenson Table3 numerical computations, Faltings/\allowbreak{}semisimplicity, local inertia and modular-curve exclusion theorems remain cited classical premises, not newly re-refereed.\allowbreak{} Full new all-CM input-bit-model decision proof is written separately.\allowbreak{} These exact NOS statements and the full supplied Corollary9 proof were reread at06:\allowbreak{}00UTC; Theorem8 points to Serre–\allowbreak{}Tate1968Theorem1.\allowbreak{} Faltings original Corollary2 was subsequently opened directly as S-NT-FALTINGS.\allowbreak{}

\url{https://londmathsoc.onlinelibrary.wiley.com/doi/full/10.1112/blms.13002}.

\paragraph{S-NT-FALTINGS} G.\allowbreak{} Faltings, Endlichkeitssätze für abelsche Varietäten über Zahlkörpern, Invent.\allowbreak{} Math.\allowbreak{}73 (1983),349–\allowbreak{}366.

\textbf{Locator:} Sätze3–\allowbreak{}4printed360; Korollar1 and Korollar2(i)–\allowbreak{}(iv), with supplied short proof, printed361.\allowbreak{} DOI10.\allowbreak{}1007/\allowbreak{}BF01388432.\allowbreak{} Standard corrected theorem is used; the1984erratum is Invent.\allowbreak{}Math.\allowbreak{}75,p381.\allowbreak{}.

\textbf{Inspection:} Primary original paper text opened and read for exact semisimplicity/\allowbreak{}Hom statements and the entire Corollary2 statement plus its short deduction.\allowbreak{} The equivalence with equality of local Euler factors at all but finitely many places directly supplies an odd good trace witness, so no extra semisimple-character citation is required.\allowbreak{} This is not a fresh audit of the long Faltings proof or its upstream finiteness arguments; the accepted classical isogeny theorem is the input.\allowbreak{}

\url{https://math.uchicago.edu/~drinfeld/Deligne%27s_conjecture_Manin_conf/Faltings_argument/Faltings.pdf}.

\paragraph{S-NT-CS} John E.\allowbreak{} Cremona and Andrew V.\allowbreak{} Sutherland, Computing the endomorphism ring of an elliptic curve over a number field, Contemporary Mathematics796 (2024),75–\allowbreak{}102.

\textbf{Locator:} Authorv4(15December2023),21pp:\allowbreak{} Theorems1–\allowbreak{}2p2; equation2.\allowbreak{}1p3; equation2.\allowbreak{}2/\allowbreak{}Proposition13/\allowbreak{}Corollary14p8; Proposition16p9; Proposition19 and full proofpp9–\allowbreak{}10; Algorithm1p13, Lemma23pp13–\allowbreak{}14, Theorems24/\allowbreak{}27 and full supplied proofs p14; Algorithm2/\allowbreak{}Theorem30 and full proofp15; heuristic(H)/\allowbreak{}Theorem35 and full supplied proofp17.\allowbreak{} Publisher DOI10.\allowbreak{}1090/\allowbreak{}conm/\allowbreak{}796/\allowbreak{}15998.\allowbreak{}.

\textbf{Inspection:} Downloaded original authorv4.\allowbreak{} Actual p9 image inspected for Proposition19 constants and Proposition16.\allowbreak{} Read supplied algebraic derivations in §§2.\allowbreak{}2–\allowbreak{}2.\allowbreak{}7 and complete Algorithm1 correctness/\allowbreak{}complexity proofs, plus Algorithm2 termination proof and entire Theorem35 deduction with its explicit extra heuristics.\allowbreak{} Upstream K-curve/\allowbreak{}volcano theory, numerical class-number tabulations, modular-polynomial routine, and computational subroutine analyses remain cited classical inputs.\allowbreak{} NT-E17 replaces the least-ordinary-prime choice by a biquadratic split-prime proof, avoiding Lemma23 finite-list dependence, and proves only polynomial complexity directly.\allowbreak{} No experiment or upstream numerical optimization rerun.\allowbreak{}

\url{https://arxiv.org/pdf/2301.11169v4}.

\paragraph{S-NT-LT} Alessandro Languasco and Timothy S.\allowbreak{} Trudgian, Uniform effective estimates for |L(1,chi)|, J.\allowbreak{} Number Theory236 (2022),245–\allowbreak{}260.

\textbf{Locator:} Authorv1(17November2020):\allowbreak{} Theorem1p1, equations2–\allowbreak{}3; Corollary1 and class-number formula p2; explicit q>=5 fundamental quadratic-character scope.\allowbreak{}.

\textbf{Inspection:} Downloaded original author PDF.\allowbreak{} Read exact Theorem1 and entire Corollary1 statement with its supplied class-number-formula deduction.\allowbreak{} This identifies the quadratic-Dirichlet analytic input to CS Proposition19.\allowbreak{} The full analytic proof and small-conductor numerical verification are cited published premises, not newly re-proved or rerun.\allowbreak{}

\url{https://arxiv.org/pdf/2011.08348}.

\paragraph{S-NT-ENGE} Andreas Enge, The complexity of class polynomial computation via floating point approximations, Math.\allowbreak{} Comp.\allowbreak{}78 (2009),1089–\allowbreak{}1107.

\textbf{Locator:} Authorv2(25July2008):\allowbreak{} abstractp1; Theorem1.\allowbreak{}1, Theorem1.\allowbreak{}2 and heuristic Corollary1.\allowbreak{}3pp2–\allowbreak{}3; DOI10.\allowbreak{}1090/\allowbreak{}S0025-5718-08-02200-X.\allowbreak{}.

\textbf{Inspection:} Downloaded originalv2 and read the exact theorem/\allowbreak{}corollary statements and intervening explanation.\allowbreak{} The near-output-linear exact class-polynomial conclusion explicitly assumes rounding errors do not spoil the computation; the height upper bound itself is rigorous.\allowbreak{} Recorded solely to preserve this additional hypothesis in comparisons.\allowbreak{} Neither NT-E17 nor NT-E18 uses the floating-point algorithm.\allowbreak{} No whole-paper proof or numerical-error audit claimed.\allowbreak{}

\url{https://arxiv.org/pdf/cs/0601104v2}.

\paragraph{S-NT-BCRS} Ian F.\allowbreak{} Blake, Janos A.\allowbreak{} Csirik, Michael Rubinstein and Gadiel Seroussi, On the Computation of Modular Polynomials for Elliptic Curves (1999).

\textbf{Locator:} Author12-page PDF:\allowbreak{} §2pp4–\allowbreak{}7, definitions2–\allowbreak{}4, exact computation steps1–\allowbreak{}3p6, CRT/\allowbreak{}bit-complexity discussionp7 and coefficient bound(9)p7.\allowbreak{}.

\textbf{Inspection:} Read the entire supplied §2 algorithm and analysis from the primary author PDF, including exact integer-power-series arithmetic, finite truncation, CRT reconstruction and polynomial coefficient sizes.\allowbreak{} Upstream modular invariance, coefficient-height theorem and fast power-series routines are cited premises, not independently re-proved; no performance experiments rerun.\allowbreak{} This gives a direct source for the polynomial Phi\_\allowbreak{}ell construction used in E17.\allowbreak{}

\url{https://www.math.uwaterloo.ca/~mrubinst/publications/phi.pdf}.

\paragraph{S-AG-ZETA} Supplied ZETA\_\allowbreak{}TWO return.

\textbf{Locator:} RETURN\_\allowbreak{}ZETA\_\allowbreak{}TWO\_\allowbreak{}final\_\allowbreak{}IMPLICATIONS.\allowbreak{}md §§1–\allowbreak{}4, ZT01–\allowbreak{}ZT08, CM01, SV01–\allowbreak{}SV02; matching LEDGER and QUESTIONS.

\textbf{Inspection:} Full implication text read and relevant reductions rechecked algebraically; original earlier full proof file not included.\allowbreak{} Earlier findings retain their IDs and are not claimed new.\allowbreak{}


\paragraph{S-AG-BRIDGES} Supplied ALPOGE\_\allowbreak{}BRIDGES and original BRIDGES returns.

\textbf{Locator:} RETURN\_\allowbreak{}ALPOGE\_\allowbreak{}BRIDGES\_\allowbreak{}cp02\_\allowbreak{}IMPLICATIONS.\allowbreak{}md and QUESTIONS:\allowbreak{} BR-TR-01/\allowbreak{}02, BR-MIX-01, BR-WM-02, AQ-FINITE, R-GTHAT; OURS\_\allowbreak{}BRIDGES\_\allowbreak{}BRIDGES\_\allowbreak{}20261120.\allowbreak{}tex F1–\allowbreak{}F4,Q01–\allowbreak{}Q08.

\textbf{Inspection:} Implication/\allowbreak{}ledger/\allowbreak{}question records read; reused background is distinguished from the original proof package, which is not attached.\allowbreak{}


\paragraph{S-AG-BS} Bach–\allowbreak{}Sorenson, Explicit bounds for primes in residue classes, Math.\allowbreak{}Comp.\allowbreak{}65(1996),1717–\allowbreak{}1735.

\textbf{Locator:} Theorem5.\allowbreak{}1 printed1729/\allowbreak{}PDF13; §2 printed1720 (RH \ensuremath{\zeta}E,\ensuremath{\zeta}Q); Corollary3.\allowbreak{}3 printed1721; Theorem5.\allowbreak{}3 printed1731–\allowbreak{}1732/\allowbreak{}PDF15–\allowbreak{}16, with Lemma5.\allowbreak{}2p1731, gives p<2(q logq)\textasciicircum{}2 for coprime residue classes.

\textbf{Inspection:} Original author-upload PDF theorem and hypothesis text read; no reproof of full classical explicit estimate.\allowbreak{} Freshly read full displayed Theorem5.\allowbreak{}3 proof and Lemma5.\allowbreak{}2, and inspected original imagep1731 from work/\allowbreak{}complexity/\allowbreak{}sources/\allowbreak{}BachSorenson1996.\allowbreak{}pdf.\allowbreak{} The authors’ numerical program is a cited classical input, not freshly certified.\allowbreak{} The printed last summation sentence of Lemma5.\allowbreak{}2 does not follow just from its displayed discriminant inequality; for the prime moduli used here the correction is proved directly in notes.\allowbreak{}md, since only the principal character is imprimitive.\allowbreak{}

\url{https://www.researchgate.net/profile/Jonathan-Sorenson-2/publication/220576951_Explicit_bounds_for_primes_in_residue_classes/links/55100ea70cf2ac2905afbceb/Explicit-bounds-for-primes-in-residue-classes.pdf}.

\paragraph{S-AG-DG} Deligne–\allowbreak{}Goncharov, Groupes fondamentaux motiviques de Tate mixte.

\textbf{Locator:} arXiv math/\allowbreak{}0302267v2,27Feb2003; (2.\allowbreak{}1.\allowbreak{}3), Propositions2.\allowbreak{}2–\allowbreak{}2.\allowbreak{}3, §2.\allowbreak{}4, printed14/\allowbreak{}PDF16.

\textbf{Inspection:} Primary formula and surrounding hypotheses read, distinguishing rational Ext and free Lie presentation from canonical generators and integral torsion.\allowbreak{}

\url{https://arxiv.org/pdf/math/0302267}.

\paragraph{S-AG-MCMULLEN} C.\allowbreak{}McMullen, Advanced Complex Analysis, Harvard Math213a course notes.

\textbf{Locator:} Version16Nov2026; Theorem3.\allowbreak{}16 printed92/\allowbreak{}PDF93; Theorems3.\allowbreak{}10–\allowbreak{}3.\allowbreak{}12 zero exponent.

\textbf{Inspection:} Primary author notes used for exact classical Hadamard locator; statement and applicability read.\allowbreak{} The theorem is classical, not presented as a2026 innovation.\allowbreak{}

\url{https://people.math.harvard.edu/~ctm/home/text/class/harvard/101/22/html/www/math213a/home/course/course.pdf}.

\paragraph{S-AG-KUMMER} Stacks Project, Kummer theory.

\textbf{Locator:} Lemma59.\allowbreak{}28.\allowbreak{}1 and immediately following long exact cohomology sequence; fppf variant Lemma59.\allowbreak{}28.\allowbreak{}3 kept separate.

\textbf{Inspection:} Entire proof of Lemma59.\allowbreak{}28.\allowbreak{}1 and displayed sequence read; ell-invertibility and inverse-limit transition conventions checked.\allowbreak{}

\url{https://stacks.math.columbia.edu/tag/03PK}.

\paragraph{S-AG-MILNE-ANT} J.\allowbreak{}S.\allowbreak{}Milne, Algebraic Number Theory, v3.\allowbreak{}08,19July2020.

\textbf{Locator:} Theorem4.\allowbreak{}3 PDF72 (Minkowski bound); Theorem4.\allowbreak{}4 class group finiteness.

\textbf{Inspection:} Exact bound and hypotheses read; explicit -35 ideal-class example proved in this run, not copied from source.\allowbreak{}

\url{https://www.jmilne.org/math/CourseNotes/ANT.pdf}.

\paragraph{S-AG-149} Alexander Grothendieck, cote149, Autour de la Longue Marche à travers la théorie de Galois.

\textbf{Locator:} 136PDFpages including cover.\allowbreak{} PDF80/\allowbreak{}footer79 eqs13–\allowbreak{}17; PDF81/\allowbreak{}footer80 eqs18–\allowbreak{}19; PDF82/\allowbreak{}footer81 eqs20–\allowbreak{}21; PDF85/\allowbreak{}footer84 eq22.\allowbreak{} Images PDF79–\allowbreak{}86 inspected.\allowbreak{}.

\textbf{Inspection:} Original official PDF downloaded and selected original page images read.\allowbreak{} The content is normalizer/\allowbreak{}inertia extensions, normal line bundles and Kummer cohomology.\allowbreak{} ORDER31 association to LM6 A44–\allowbreak{}A45 is separate inherited evidence; no K8 theorem attributed to sheets.\allowbreak{}

\url{https://grothendieck.umontpellier.fr/149.pdf}.

\paragraph{S-AG-MASULLO} Supplied Masullo The Music of the Primes draft.

\textbf{Locator:} Original PDF425 Lemma9.\allowbreak{}3.\allowbreak{}2.\allowbreak{}4;426 Theorem9.\allowbreak{}3.\allowbreak{}2.\allowbreak{}6 eq9.\allowbreak{}3.\allowbreak{}2.\allowbreak{}10;804 DefinitionD.\allowbreak{}3.\allowbreak{}3.\allowbreak{}1;11 norm/\allowbreak{}SE definition.

\textbf{Inspection:} Original images425,426,804 inspected, norm text11 read.\allowbreak{} Independent whole-proof checks MT-M01–\allowbreak{}03 recorded in events.\allowbreak{}tsv.\allowbreak{} No claim that M itself is disproved by auxiliary counterexamples.\allowbreak{}


\paragraph{S-AG-HODGE} Supplied Hodge bridge sketch.

\textbf{Locator:} OURS\_\allowbreak{}HODGE\_\allowbreak{}BRIDGE\_\allowbreak{}SKETCH\_\allowbreak{}20261121.\allowbreak{}md §§1–\allowbreak{}5, archive map§5; cote16 scans14–\allowbreak{}18,44–\allowbreak{}50; cote44 scans2,4,6,8–\allowbreak{}9; cote47 scans3,5–\allowbreak{}9; Route1.\allowbreak{}4/\allowbreak{}Suggestion6, Route4 and cote16scans14–\allowbreak{}18no7; RETURN\_\allowbreak{}DP\_\allowbreak{}MULT\_\allowbreak{}cp01\_\allowbreak{}IMPLICATIONS.\allowbreak{}md §6 DP15/\allowbreak{}Q4e; war-game20261112 H8.

\textbf{Inspection:} Relevant arguments and reported archive locators read.\allowbreak{} Offset in cote16 scan convention retained.\allowbreak{} Separate2026CM-Hodge claims not independently certified from this sketch.\allowbreak{} Arthur and integral routes reread; DP15 returned result reused with original ID; full H8 proof not supplied and H8(iv) not newly certified.\allowbreak{}


\paragraph{S-AG-SA} Supplied SA22 and SA23 intakes.

\textbf{Locator:} OURS\_\allowbreak{}INTAKE\_\allowbreak{}RESEARCH\_\allowbreak{}SA22\_\allowbreak{}20261120.\allowbreak{}md §5.\allowbreak{}1 A6-node-character/\allowbreak{}A:\allowbreak{}inertia-maximality,§5.\allowbreak{}3Soulé/\allowbreak{}GT,§5.\allowbreak{}4–\allowbreak{}5.\allowbreak{}8; SA23§4.\allowbreak{}1 geo-H-lidskii/\allowbreak{}geo-H-genus.

\textbf{Inspection:} Relevant statements and intake check records read; original proofs/\allowbreak{}scans absent except149 separately retrieved.\allowbreak{} Cote1 scans79–\allowbreak{}80,111–\allowbreak{}112 are inherited locators.\allowbreak{}


\paragraph{S-AG-ANDERSON} G.\allowbreak{}W.\allowbreak{}Anderson, Cyclotomy and an extension of the Taniyama group, Compositio Mathematica57(1986),153–\allowbreak{}217.

\textbf{Locator:} §5.\allowbreak{}3 printed178/\allowbreak{}PDF27; Theorem6 §5.\allowbreak{}6 printed179–\allowbreak{}180/\allowbreak{}PDF28–\allowbreak{}29; Cor5.\allowbreak{}7.\allowbreak{}2/\allowbreak{}Rmk5.\allowbreak{}7.\allowbreak{}3/\allowbreak{}Cor5.\allowbreak{}7.\allowbreak{}4 printed181/\allowbreak{}PDF30; Prop5.\allowbreak{}7.\allowbreak{}5 printed181–\allowbreak{}182/\allowbreak{}PDF30–\allowbreak{}31; §5.\allowbreak{}8 Theorem7 printed182/\allowbreak{}PDF31.

\textbf{Inspection:} Original primary PDF downloaded; images27–\allowbreak{}31 and complete proofs Theorem6, Prop5.\allowbreak{}7.\allowbreak{}5 and Theorem7 read.\allowbreak{} Deeper cited Serre/\allowbreak{}Langlands/\allowbreak{}Deligne/\allowbreak{}Blasius results are classical inputs, not re-proved.\allowbreak{} Printed Rmk5.\allowbreak{}7.\allowbreak{}3 refers to Cor5.\allowbreak{}8.\allowbreak{}2, an internal numbering error noted explicitly.\allowbreak{}

\url{https://www.numdam.org/item/CM_1986__57_2_153_0.pdf}.

\paragraph{S-AG-FRAMEWORK} Classical geometric framework locators carried in ZETA\_\allowbreak{}TWO.

\textbf{Locator:} SGA4 XVII§§5.\allowbreak{}1–\allowbreak{}5.\allowbreak{}3; SGA5 XV§3 no2 unnumbered theorem eq1; EGA IV3(8.\allowbreak{}8.\allowbreak{}2); WeilI Thm1.\allowbreak{}6; WeilII Thm3.\allowbreak{}3.\allowbreak{}1; SGA6 V§§2,7(7.\allowbreak{}1.\allowbreak{}1); Deligne1979§§1.\allowbreak{}1–\allowbreak{}1.\allowbreak{}3,(3.\allowbreak{}1.\allowbreak{}2).

\textbf{Inspection:} Carried exact locators, not a claimed new full reading of the original SGA/\allowbreak{}EGA/\allowbreak{}Deligne works.\allowbreak{}


\paragraph{S-AG-34} Alexander Grothendieck, cote34, ExposéII.

\textbf{Locator:} 210PDFpages includingcover; Thm1.\allowbreak{}1 PDF53/\allowbreak{}footer52; struckRem1.\allowbreak{}17 PDF64/\allowbreak{}footer63; Thm1.\allowbreak{}1 proof §§2–\allowbreak{}3 PDF70–\allowbreak{}78/\allowbreak{}footer69–\allowbreak{}77; (2.\allowbreak{}5.\allowbreak{}5) PDF74/\allowbreak{}footer73.

\textbf{Inspection:} Official84,513,436bytePDF downloaded; pageimages53–\allowbreak{}78 read, with enlarged crops53/\allowbreak{}64.\allowbreak{} Exact normal/\allowbreak{}locallyfinite/\allowbreak{}geoconnected hypotheses and pro-ell solvable quotient wording recovered.\allowbreak{} Full theorem1.\allowbreak{}1 proof route §§2–\allowbreak{}3 read and geometric reductions checked; foundational SGA citations not separately re-proved.\allowbreak{} Webopen failed because filelarge; original local download succeeded.\allowbreak{}

\url{https://grothendieck.umontpellier.fr/34.pdf}.

\paragraph{S-AG-RV} D.\allowbreak{}Ramakrishnan and R.\allowbreak{}J.\allowbreak{}Valenza, Fourier Analysis on Number Fields(1999).

\textbf{Locator:} §7.\allowbreak{}4 Theorem7-19 printed279; §7.\allowbreak{}7 Theorem7-28 printed289, whole proof289–\allowbreak{}292.

\textbf{Inspection:} Primary PDF downloaded by number\_\allowbreak{}theory; theorem/\allowbreak{}hypothesis text and complete boundary proof289–\allowbreak{}292 directly read here.\allowbreak{} Boundary includes all unitary characters and norm twists; completed norm-character pole exceptions retained.\allowbreak{} Full proof of Theorem7-19 not separately re-refereed.\allowbreak{}

\url{https://webusers.imj-prg.fr/~julien.marche/Ref_GL_GT/Fourier_Analysis_on_Number_Fields.pdf}.

\paragraph{S-AG-MILNE-JAV} J.\allowbreak{}S.\allowbreak{}Milne, Jacobian Varieties and Abelian Varieties, revised notes of1986 chapters.

\textbf{Locator:} JacobianVarieties version12June2021 Prop9.\allowbreak{}1 pp30–\allowbreak{}32; AbelianVarieties version2Jan2022 Remarks15.\allowbreak{}3/\allowbreak{}15.\allowbreak{}5 pp28–\allowbreak{}29.

\textbf{Inspection:} Primary statements and comparison proof read; the Galois-compatible Tate-module/\allowbreak{}abelianized-fundamental-group identification is used, not an asserted cycle comparison.\allowbreak{}

\url{https://www.jmilne.org/math/xnotes/JVs.pdf}.

\paragraph{S-AG-BES} N.\allowbreak{}Borne, M.\allowbreak{}Emsalem and J.\allowbreak{}Stix, Lifting Galois sections along torsors, Math.\allowbreak{}Nachr.\allowbreak{}288(2015),1939–\allowbreak{}1951.

\textbf{Locator:} arXiv1301.\allowbreak{}4429v2:\allowbreak{} algebraicK(pi,1) curves §3.\allowbreak{}1PDF4; Prop3.\allowbreak{}1/\allowbreak{}3.\allowbreak{}4PDF5–\allowbreak{}6; continuous-cohomology Def4.\allowbreak{}1/\allowbreak{}Prop4.\allowbreak{}2PDF7; Prop4.\allowbreak{}4PDF8–\allowbreak{}9; Lemmas4.\allowbreak{}5–\allowbreak{}4.\allowbreak{}6PDF9.

\textbf{Inspection:} PrimaryPDF4–\allowbreak{}9 fully read, including whole proofs of exactness, kernel action, continuous-cohomology limits, class/\allowbreak{}transgression comparison and section criterion; images5,8,9 inspected.\allowbreak{} Bibliographic journal details confirmed from publisher.\allowbreak{} Foundational étale-homotopy and Jannsen spectral-sequence references are invoked, not re-proved.\allowbreak{}

\url{https://arxiv.org/pdf/1301.4429}.

\paragraph{S-AG-SHARIFI} R.\allowbreak{}T.\allowbreak{}Sharifi, Relationships between conjectures on the structure of pro-p Galois groups unramified outside p, Proc.\allowbreak{}Sympos.\allowbreak{}PureMath.\allowbreak{}70(2002),275–\allowbreak{}284.

\textbf{Locator:} 13-page author version:\allowbreak{} definitionsPDF1–\allowbreak{}2; Thm1.\allowbreak{}1PDF2/\allowbreak{}proof8–\allowbreak{}9; Thm1.\allowbreak{}3PDF3/\allowbreak{}proof10; Eq2.\allowbreak{}10PDF7; full supporting §§2–\allowbreak{}3PDF3–\allowbreak{}10; Thm4.\allowbreak{}1PDF10–\allowbreak{}11.

\textbf{Inspection:} Primary author text read through the complete §§2–\allowbreak{}3 proofs, including maximal-image element construction, pro-p group lemma, regular-prime theorem and Greenberg obstruction.\allowbreak{} Exact domain and factorial normalization checked.\allowbreak{} Underlying Ihara, Iwasawa and Brown inputs cited rather than independently re-proved.\allowbreak{} Screenshot tool returned text markers only; no visual-inspection claim.\allowbreak{}

\url{https://www.math.ucla.edu/~sharifi/publ/galstr.pdf}.

\paragraph{S-AG-HOSHI} Y.\allowbreak{}Hoshi, On the Kernels of the Pro-l Outer Galois Representations Associated to Hyperbolic Curves over Number Fields, revised March2014 manuscript.

\textbf{Locator:} Definition1.\allowbreak{}8, Remark1.\allowbreak{}8.\allowbreak{}1 and Theorem1.\allowbreak{}9PDF11; references[4]/\allowbreak{}[22]PDF27.

\textbf{Inspection:} Primary exact field definition and regular-prime statement read, including its proof citation to Brown and Sharifi.\allowbreak{} Other Hoshi theorems not being certified here.\allowbreak{}

\url{https://www.kurims.kyoto-u.ac.jp/~yuichiro/rims1782revised.pdf}.

\paragraph{S-AG-BROWN} F.\allowbreak{}Brown, Mixed Tate motives over Z, Ann.\allowbreak{}Math.\allowbreak{}175(2012),949–\allowbreak{}976.

\textbf{Locator:} IntroductionConjecture1printed949/\allowbreak{}PDF1; Theorem1.\allowbreak{}1 and implication to Conjecture1 printed950/\allowbreak{}PDF2.

\textbf{Inspection:} Primary exact statement, category and Deligne–\allowbreak{}Ihara consequence in introduction read.\allowbreak{} The main 28-page proof was not re-refereed; treated as a classical cited theorem.\allowbreak{}

\url{https://annals.math.princeton.edu/wp-content/uploads/annals-v175-n2-p10-p.pdf}.

\paragraph{S-AG-ALPOGE-FABLE} Supplied Alpoge Remark and Fable programme explanation.

\textbf{Locator:} ALPOGE\_\allowbreak{}remark.\allowbreak{}pdf step2PDF3, step17/\allowbreak{}fn6PDF4; OURS\_\allowbreak{}FABLE\_\allowbreak{}PROGRAM\_\allowbreak{}SUGGESTIONS\_\allowbreak{}20261120\_\allowbreak{}EXPLAINED.\allowbreak{}md §§2,5,6; ALPOGE\_\allowbreak{}LEAN\_\allowbreak{}MAP\_\allowbreak{}20261120.\allowbreak{}md step2 row.

\textbf{Inspection:} Local notePDF3–\allowbreak{}4 text and exact step2 formula read; Fable§2 full inherited proof sketch read.\allowbreak{} New AG-FG proofs replace its special-normal-form passage by a direct recursion.\allowbreak{} Added archive locators are explicitly inherited from Fable, not visually rechecked.\allowbreak{} No exact Lean lemma claimed beyond supplied low-confidence map.\allowbreak{}

\url{https://alpo.ge/remark.pdf}.

\paragraph{S-AG-ARTHUR} J.\allowbreak{}Arthur, Unipotent automorphic representations:\allowbreak{} conjectures, Astérisque 171–\allowbreak{}172(1989), 13–\allowbreak{}71.

\textbf{Locator:} §5 printed29ff local cohomological hypotheses; §9(9.\allowbreak{}1)printed58; Lefschetz actionprinted60–\allowbreak{}61; Prop9.\allowbreak{}1printed62–\allowbreak{}66/\allowbreak{}ClayPDF50–\allowbreak{}54.

\textbf{Inspection:} Entire Prop9.\allowbreak{}1 proof and surrounding §9 construction read as primary text; §5 initial hypotheses checked.\allowbreak{} Vogan–\allowbreak{}Zuckerman Prop6.\allowbreak{}19 remains cited, not independently re-proved.\allowbreak{} No blanket claim that global Conj8.\allowbreak{}1 remains open or is established in every intended case.\allowbreak{}

\url{https://www.claymath.org/library/cw/arthur/pdf/34.pdf}.

\paragraph{S-AG-GGP} W.\allowbreak{}T.\allowbreak{}Gan,B.\allowbreak{}H.\allowbreak{}Gross,D.\allowbreak{}Prasad, Branching laws for classical groups:\allowbreak{} the non-tempered case, arXiv1911.\allowbreak{}02783v2(17Aug2020).

\textbf{Locator:} §2PDF8, separate Deligne/\allowbreak{}ArthurSL2 and A-parameter→L-parameter; introductory conventionPDF2.

\textbf{Inspection:} Full relevant definitions read; local factor calculation proved separately.\allowbreak{} No general non-tempered branching conjecture is claimed a theorem.\allowbreak{}

\url{https://arxiv.org/pdf/1911.02783}.

\paragraph{S-AG-DELIGNE-LOCAL} P.\allowbreak{}Deligne, Les constantes des équations fonctionnelles des fonctions L, LNM349(1973),501–\allowbreak{}597.

\textbf{Locator:} §8.\allowbreak{}4.\allowbreak{}1 WD definition, IASPDF68; §8.\allowbreak{}12 local factors printed572/\allowbreak{}IASPDF72.

\textbf{Inspection:} Relevant primary textual extraction consulted, kernel-of-monodromy/\allowbreak{}inertia definition crosschecked with Wedhorn(3.\allowbreak{}2.\allowbreak{}1).\allowbreak{} Web screenshot emitted only a marker, so no image-check claim.\allowbreak{}

\url{https://publications.ias.edu/sites/default/files/Number20.pdf}.

\paragraph{S-AG-WEDHORN} T.\allowbreak{}Wedhorn, The local Langlands correspondence for GL(n) over p-adic fields, arXivmath/\allowbreak{}0011210v2(25Nov2000).

\textbf{Locator:} §3.\allowbreak{}1(3.\allowbreak{}1.\allowbreak{}10)PDF36 WD/\allowbreak{}DeligneSL2; §3.\allowbreak{}2(3.\allowbreak{}2.\allowbreak{}1)PDF37 determinant on (kerN)\textasciicircum{}I and geometric Frobenius.

\textbf{Inspection:} Relevant author exposition read, local formula directly used and new diagonal/\allowbreak{}monodromy calculations checked.\allowbreak{} No new proof of local Langlands claimed.\allowbreak{}

\url{https://arxiv.org/pdf/math/0011210}.

\paragraph{S-AG-STIX-CUSP} Jakob Stix, On cuspidal sections of algebraic fundamental groups, author manuscript dated24Feb2012.

\textbf{Locator:} IntroductionPDF1 char0/\allowbreak{}full-profinite scope; §2PDF4 Kummer convention; §3Lemma14/\allowbreak{}Prop11 PDF7; Prop16(1)–\allowbreak{}(2) PDF8.

\textbf{Inspection:} Read complete Lemma14 and Proposition16 proofs, inertia/\allowbreak{}decomposition definitions, Proposition11 root splitting and Kummer convention.\allowbreak{} Underlying Tamagawa torsion-freeness input cited, not separately reproved.\allowbreak{} Exact theorem used for B=P1\ensuremath{-}\{0,1,\ensuremath{\infty}\}, extended to nI by a new argument.\allowbreak{}

\url{https://www.math.uni-frankfurt.de/~stix/research/preprints/STIXcuspsec20120802homepage.pdf}.

\paragraph{S-AG-KUNNETH} SGA1, Revêtements étales et groupe fondamental, ExposéX Corollaires1.\allowbreak{}7–\allowbreak{}1.\allowbreak{}8, recomposed arXivmath/\allowbreak{}0206203v2.

\textbf{Locator:} ExposéX Cor1.\allowbreak{}7 recomposedpp203–\allowbreak{}204, PDF219–\allowbreak{}220; Cor1.\allowbreak{}8 and proof recomposedpp204–\allowbreak{}205, PDF220–\allowbreak{}221; original marginalpp265–\allowbreak{}266.

\textbf{Inspection:} Entire product corollary proof and following algebraically-closed-base-change proof read.\allowbreak{} Proper-factor condition retained.\allowbreak{} Arithmetic fibre-product identification proved separately from geometric exact sequences; no assertion of unrestricted positive-characteristic Künneth.\allowbreak{}

\url{https://arxiv.org/pdf/math/0206203}.

\paragraph{S-AG-PF26-GEOMETRY} OpenAI, Deterministic Polynomial Factorization over Prime Fields,4Oct2026.

\textbf{Locator:} §4pp13–\allowbreak{}19; Lemma4.\allowbreak{}1p15, Lemma4.\allowbreak{}2pp15–\allowbreak{}16, Lemmas4.\allowbreak{}3–\allowbreak{}4.\allowbreak{}4p17, Cor4.\allowbreak{}5p18, Prop4.\allowbreak{}6pp18–\allowbreak{}19.

\textbf{Inspection:} Entire §4 source proofs independently read/\allowbreak{}rederived from local original PDF text; local SHA256 recomputed.\allowbreak{} Official repository metadata freshly opened.\allowbreak{} Geometry review excludes norm-solving implementation, companion analytic claims and entire factoring theorem.\allowbreak{}

\url{https://github.com/openai/math/blob/main/preprints/Deterministic-Polynomial-Factorization-over-Prime-Fields-October-4-2026/Deterministic-Polynomial-Factorization-over-Prime-Fields.pdf}.

\paragraph{S-AG-BK-FOURIER} Kenichi Bannai and Shinichi Kobayashi, Integral structures on p-adic Fourier theory, Ann.\allowbreak{} Inst.\allowbreak{} Fourier66(2016),521–\allowbreak{}550.

\textbf{Locator:} §2 printed526–\allowbreak{}527/\allowbreak{}PDF7–\allowbreak{}8:\allowbreak{} Fourier period and valuation; Proposition3.\allowbreak{}7 printed533–\allowbreak{}535/\allowbreak{}PDF14–\allowbreak{}16; Proposition4.\allowbreak{}1 printed536–\allowbreak{}537/\allowbreak{}PDF17–\allowbreak{}18; Theorem4.\allowbreak{}3 printed538/\allowbreak{}PDF19; §5 printed543–\allowbreak{}549/\allowbreak{}PDF24–\allowbreak{}30, especially Proposition5.\allowbreak{}3(i),(ii) and proof pp544–\allowbreak{}545, Lemma5.\allowbreak{}4 and Proposition5.\allowbreak{}5 pp546–\allowbreak{}547; Theorem5.\allowbreak{}8 p549.

\textbf{Inspection:} Primary whole §5 read.\allowbreak{} The new Fourier transfer uses Proposition5.\allowbreak{}3(i),(ii), whose entire displayed proof was checked, and the full torsion-average/\allowbreak{}moment argument of Lemma5.\allowbreak{}4 and Proposition5.\allowbreak{}5.\allowbreak{} Also read the complete proof of Proposition3.\allowbreak{}7 and Proposition4.\allowbreak{}1, and checked Theorem4.\allowbreak{}3 as their reformulation.\allowbreak{} Exact period valuation in §2 located; underlying Tate/\allowbreak{}Lubin–\allowbreak{}Tate period theorem remains cited.\allowbreak{} No independent full-paper rereferee claim.\allowbreak{} The later sigma-radius theorem has not been read and is not imported.\allowbreak{}

\url{https://www.numdam.org/item/10.5802/aif.3018.pdf}.

\paragraph{S-AG-VELU} Jacques Vélu, Isogénies entre courbes elliptiques, C.\allowbreak{} R.\allowbreak{} Acad.\allowbreak{} Sci.\allowbreak{} Paris273(1971), A238–\allowbreak{}A241; Alexandru Ghitza English translation.

\textbf{Locator:} §2 formula7 on translationPDF2; §3 formulas8–\allowbreak{}11 and their derivation PDF3–\allowbreak{}4; Remark2 on normalized differential and Remark3 on descent PDF4.

\textbf{Inspection:} Complete relevant quotient construction and formula derivation in §§2–\allowbreak{}3 read in the author-source translation, including normalization and descent remarks.\allowbreak{} Used integral coefficient formulas and their reduction; did not rely on visibly corrupted text-layer derivative/\allowbreak{}addition formulas.\allowbreak{} Original French page images were not inspected.\allowbreak{}

\url{https://aghitza.github.io/publications/translation-velu.pdf}.

\paragraph{S-AG-BK-THETA} Kenichi Bannai and Shinichi Kobayashi, Algebraic theta functions and the p-adic interpolation of Eisenstein–\allowbreak{}Kronecker numbers, arXiv:\allowbreak{}math/\allowbreak{}0610163v4 (11Dec2007); Duke Math.\allowbreak{}J.\allowbreak{}153(2010),229–\allowbreak{}295.

\textbf{Locator:} Definition1.\allowbreak{}8/\allowbreak{}Example1.\allowbreak{}9, PDF10–\allowbreak{}11, exact reduced-theta/\allowbreak{}sigma convention; §2.\allowbreak{}3 ordinarity condition(28) PDF27; Lemma2.\allowbreak{}13/\allowbreak{}Propositions2.\allowbreak{}14–\allowbreak{}2.\allowbreak{}15 PDF28–\allowbreak{}29; Corollaries2.\allowbreak{}16–\allowbreak{}2.\allowbreak{}17 PDF29–\allowbreak{}30.

\textbf{Inspection:} Read entire reduced-theta construction in Example1.\allowbreak{}9 and exact normalization; read displayed complete Lemma2.\allowbreak{}13/\allowbreak{}Propositions2.\allowbreak{}14–\allowbreak{}2.\allowbreak{}15 proof arguments and ordinary hypotheses.\allowbreak{} Earlier algebraic-theta inputs of those propositions were cited, not all independently reread.\allowbreak{} Supersingular radius experiments were not used as theorems.\allowbreak{} New coefficient counterexample is independently derived from the curve.\allowbreak{}

\url{https://arxiv.org/pdf/math/0610163}.

\paragraph{S-AG-KRY-CM} Stephen S.\allowbreak{} Kudla, Michael Rapoport and Tonghai Yang, A peculiar modular form of weight one, arXiv:\allowbreak{}math/\allowbreak{}9808143v1 (31Aug1998).

\textbf{Locator:} §5 beginsPDF23 with arbitrary imaginary quadratic field; Prop5.\allowbreak{}1 and full proof/\allowbreak{}Cases1–\allowbreak{}2 PDF23–\allowbreak{}25; Prop5.\allowbreak{}3/\allowbreak{}Cor5.\allowbreak{}4 PDF25; automorphism identity(5.\allowbreak{}2)PDF24.

\textbf{Inspection:} Read complete Prop5.\allowbreak{}1 argument including both deformation cases and coarse-moduli discussion, exact arbitrary-K scope, automorphism formula and class-field identification.\allowbreak{} Serre–\allowbreak{}Tate/\allowbreak{}Drinfeld/\allowbreak{}Gross inputs are cited by the source and were not independently re-proved.\allowbreak{} Added full level3 and field-degree/\allowbreak{}differential normalization argument is new and checked separately.\allowbreak{}

\url{https://arxiv.org/pdf/math/9808143}.

\paragraph{S-AG-AI-CM} Fabrizio Andreatta and Adrian Iovita, Katz type p-adic L-functions for primes p non-split in the CM field, Comment.\allowbreak{} Math.\allowbreak{} Helv.\allowbreak{}99(2024),641–\allowbreak{}716.

\textbf{Locator:} Published §2.\allowbreak{}4 pp654–\allowbreak{}655; Cor4.\allowbreak{}23 and Def4.\allowbreak{}24 pp681–\allowbreak{}682; §5.\allowbreak{}1–\allowbreak{}5.\allowbreak{}2 pp682–\allowbreak{}684; Prop5.\allowbreak{}6 pp687–\allowbreak{}688, Def5.\allowbreak{}7/\allowbreak{}Rem5.\allowbreak{}8 p688, Prop5.\allowbreak{}9 and complete proof pp689–\allowbreak{}690.\allowbreak{} Author arXiv1905.\allowbreak{}00792v3:\allowbreak{} §2.\allowbreak{}1–\allowbreak{}2.\allowbreak{}2 pp7–\allowbreak{}9 (infinity-type and avatar conventions), §2.\allowbreak{}4 pp12, Lem3.\allowbreak{}1 pp15–\allowbreak{}16, Lem5.\allowbreak{}1 p33, Prop5.\allowbreak{}6 pp36–\allowbreak{}37, Prop5.\allowbreak{}9 pp37–\allowbreak{}38.\allowbreak{} Ramified extension:\allowbreak{} published Lem3.\allowbreak{}2 pp659–\allowbreak{}660, §6 pp691–\allowbreak{}694, Lem6.\allowbreak{}1 p692, Prop6.\allowbreak{}4 pp693–\allowbreak{}694 and Prop6.\allowbreak{}6 p694; arXivv3 Lem3.\allowbreak{}2 p16 and complete §6 pp39–\allowbreak{}41.\allowbreak{}.

\textbf{Inspection:} Read the complete relevant proofs of Lem3.\allowbreak{}1, Lem5.\allowbreak{}1, Prop5.\allowbreak{}6 and Prop5.\allowbreak{}9, and the exact §2.\allowbreak{}1–\allowbreak{}2.\allowbreak{}4 conventions and Def4.\allowbreak{}24 thresholds; published Prop5.\allowbreak{}9 and operator normalization Cor4.\allowbreak{}23 cross-checked.\allowbreak{} Detailed overconvergent Gauss–\allowbreak{}Manin iteration and external Katz/\allowbreak{}Waldspurger theorems remain cited classical inputs, not independently rerefereed.\allowbreak{} Retained F[p]=F\ensuremath{-}VU F and the adelic class-sum denominator from the precise definitions despite abbreviated inconsistent intermediate displays.\allowbreak{} Further full reading of Lem3.\allowbreak{}2 and all §6, including Lem6.\allowbreak{}1 proof and Prop6.\allowbreak{}6 proof, with published exact locators checked.\allowbreak{} The ramified different valuation and n\ensuremath{\ge}2 threshold are retained; generic-fibre transversality and even prime-ideal conductor exponent are explicit in the transfer.\allowbreak{}

\url{https://ems.press/content/serial-article-files/48931?nt=1}.

\paragraph{S-AG-DENINGER-CENTRAL} Christopher Deninger, Is there a Birch and Swinnerton-Dyer conjecture for Dedekind zeta functions?, arXiv2504.\allowbreak{}15767v3 (1Apr2026); J.\allowbreak{}Number Theory286(2026),47–\allowbreak{}66.

\textbf{Locator:} Serre central-order hypothesis pp1–\allowbreak{}2; speculative cohomology §2 pp2–\allowbreak{}9; Thm2.\allowbreak{}15/\allowbreak{}2.\allowbreak{}16 p10, complete proofs pp10–\allowbreak{}14; exponential-motive discussion pp15–\allowbreak{}16; Fresan–\allowbreak{}Jossen Thm4.\allowbreak{}1 and complete reproduced proof pp16–\allowbreak{}17.\allowbreak{} All locators are author arXivv3 pages, not asserted publisher page mapping.\allowbreak{} The fixed-part input is Fresan–\allowbreak{}Jossen Exponential motives Thm6.\allowbreak{}5.\allowbreak{}1 in the version cited as source[8], quoted on Deningerp17.\allowbreak{}.

\textbf{Inspection:} Read all18pages.\allowbreak{} Abstract theorem is unconditional; Serre enters the analytic dimension comparison.\allowbreak{} Reconstructed arbitrary non-Galois descent by the invariant-space averaging projector, rather than using the abbreviated Pi\_\allowbreak{}K direct-sum notation outside its Galois scope.\allowbreak{} Made G-equivariance explicit in the compatible-star uniqueness claim.\allowbreak{} Read full reproduced FJ extension proof; the quoted fixed-part theorem is a cited input, not independently rerefereed in the full book.\allowbreak{}

\url{https://arxiv.org/pdf/2504.15767v3}.

\paragraph{S-AG-FQ-REAL} Jacques Queyrut, Sur l’équation fonctionnelle des séries L d’Artin pour des caractères réels, Mémoires de la SMF37(1974),133–\allowbreak{}135.

\textbf{Locator:} Unnumbered theorem p133 for virtual representations realizable over R; dihedral reduction p134; rank-one quadratic-relative root-number formula p135.\allowbreak{} Joint full theorem reference Frohlich–\allowbreak{}Queyrut, Invent.\allowbreak{}Math.\allowbreak{}20(1973),125–\allowbreak{}138.\allowbreak{}.

\textbf{Inspection:} Downloaded public primary PDF in memory and read entire three-page presentation.\allowbreak{} Exact theorem scope and reduction checked; final local root-number identity is quoted there rather than fully derived, so no independent complete rereferee of the original1973 proof is claimed.\allowbreak{}

\url{https://www.numdam.org/item/MSMF_1974__37__133_0.pdf}.

\paragraph{S-AG-BK-LOCAL} Spencer Bloch and Kazuya Kato, L-functions and Tamagawa numbers of motives, The Grothendieck Festschrift I (1990),333–\allowbreak{}400.

\textbf{Locator:} Local groups§3.\allowbreak{}7 pp353–\allowbreak{}354; Prop3.\allowbreak{}8 and complete duality proof pp354–\allowbreak{}359; Cor3.\allowbreak{}8.\allowbreak{}4 exact sequences pp355–\allowbreak{}356; Def3.\allowbreak{}10 p359; formal-group and abelian-variety examples pp359–\allowbreak{}361.\allowbreak{} PDF page=printed page\ensuremath{-}332.\allowbreak{}.

\textbf{Inspection:} Read the entire relevant primary pp353–\allowbreak{}361, including the complete displayed Prop3.\allowbreak{}8 proof, exact finite/\allowbreak{}exponential sequences and examples.\allowbreak{} Earlier period-ring/\allowbreak{}Tate inputs are retained as classical cited facts.\allowbreak{} Screenshot tool returned textual image markers, so no new visual page-inspection claim.\allowbreak{}

\url{https://math.stanford.edu/~conrad/BSDseminar/refs/BKTamagawa.pdf}.

\paragraph{S-AG-CF} Pierre Colmez and Jean-Marc Fontaine, Construction des représentations p-adiques semi-stables, Invent.\allowbreak{}Math.\allowbreak{}140(2000),1–\allowbreak{}43.

\textbf{Locator:} TheoremA authorPDFp2:\allowbreak{} every weakly admissible filtered(phi,N)-module over the specified complete discretely valued field is admissible; tensor equivalence described immediately above.\allowbreak{} Author file37pages, not publisher pagination.\allowbreak{}.

\textbf{Inspection:} Read exact TheoremA, base-field hypotheses and categorical interpretation in author introduction pp2–\allowbreak{}4.\allowbreak{} The full37-page proof is a classical located input and was not reread in this check.\allowbreak{} New2-dimensional weak-admissibility calculation is proved in full separately.\allowbreak{}

\url{https://webusers.imj-prg.fr/~pierre.colmez/CF.pdf}.

\paragraph{S-AG-BKNO} Ashay Burungale, Shinichi Kobayashi, Kentaro Nakamura and Kazuto Ota, A local sign decomposition for symplectic self-dual Galois representations of rank two, arXiv2508.\allowbreak{}17776v1(25Aug2025).

\textbf{Locator:} Coefficient rings and selfduality§2.\allowbreak{}2p12; genericity Def2.\allowbreak{}11p16; Thm3.\allowbreak{}1 and Cor3.\allowbreak{}2p19; fundamental exact sequence(4.\allowbreak{}1)p23; deRham determinant comparison Prop4.\allowbreak{}17pp28–\allowbreak{}29; specialization definition§5.\allowbreak{}1.\allowbreak{}1p35.\allowbreak{}.

\textbf{Inspection:} Read all §§2–\allowbreak{}4pp12–\allowbreak{}34, introductory scope and exact specialization definitionp35.\allowbreak{} Genericity over a coefficient field imposes H0=H2=0 only; no purity or phi=1 exclusion found.\allowbreak{} Cor3.\allowbreak{}2 has the explicit AG-LOCAL-EXCEPTION counterexample.\allowbreak{} Prop4.\allowbreak{}17 uses an inverse exponential unavailable there.\allowbreak{} Main sign decomposition is not declared false; upstream epsilon-isomorphism/\allowbreak{}Colmez theorems remain cited inputs, not fully rerefereed.\allowbreak{} AG-LOCAL-SPLICE now proves the exceptional de Rham determinant comparison by preserving the full dual exact complex; this does not restore the Cor3.\allowbreak{}2 nonvanishing claim.\allowbreak{}

\url{https://arxiv.org/pdf/2508.17776v1}.

\paragraph{S-AG-BMS} Bhargav Bhatt, Matthew Morrow and Peter Scholze, Integral p-adic Hodge theory, Publ.\allowbreak{}Math.\allowbreak{}IHES128(2018),219–\allowbreak{}397.

\textbf{Locator:} Author version8Jan2019:\allowbreak{} Thm1.\allowbreak{}1(i)p2, repeated Thm14.\allowbreak{}6(i)pp121–\allowbreak{}122; proofp122 reduces to Thm14.\allowbreak{}5(i),Prop13.\allowbreak{}21 and filtered comparisons.\allowbreak{}.

\textbf{Inspection:} Read exact proper-smooth-formal hypotheses, full Thm14.\allowbreak{}6 statement and its complete displayed deduction.\allowbreak{} The substantial upstream comparison theorems are located classical inputs; no full124-page proof reread claimed.\allowbreak{}

\url{https://www.math.uni-bonn.de/people/scholze/integralpadicHodge.pdf}.

\paragraph{S-AG-KATZ-MESSING} Nicholas M.\allowbreak{}Katz and William Messing, Some consequences of the Riemann hypothesis for varieties over finite fields, Invent.\allowbreak{}Math.\allowbreak{}23(1974),73–\allowbreak{}77.

\textbf{Locator:} Theorem1pp74–\allowbreak{}75 equality of Frobenius polynomials for reasonable Weil/\allowbreak{}crystalline cohomology; Cor1pp75–\allowbreak{}76 purity and hard Lefschetz; Theorem2pp76–\allowbreak{}77 rational algebraic Frobenius Kunneth projectors.\allowbreak{} Original scan has coverPDF1 then printed73 atPDF2.\allowbreak{}.

\textbf{Inspection:} Read all five primary printed pages through in-memory OCR of the original author-hosted scan, including whole Theorem1, Cor1 and Theorem2 proofs.\allowbreak{} Deligne purity/\allowbreak{}monodromy and Berthelot weak-Lefschetz inputs remain explicitly classical.\allowbreak{} No image-inspection claim for the web screenshot marker.\allowbreak{}

\url{https://web.math.princeton.edu/~nmk/old/katzmessing.pdf}.

\paragraph{S-AG-NAKAMURA-EPS} Kentaro Nakamura, A generalization of Kato’s local epsilon-conjecture for (phi,Gamma)-modules over the Robba ring, Algebra \& Number Theory 11 (2017), 319–\allowbreak{}404.

\textbf{Locator:} Author v3(15Feb2015), §3.\allowbreak{}1 pp28–\allowbreak{}30 graded lines, exact-sequence orientation(17), acyclic trivialization(19), duality(21); §3.\allowbreak{}3 pp34–\allowbreak{}40, actual finite/\allowbreak{}dual splice(22)–\allowbreak{}(24)pp36–\allowbreak{}37, Lemma3.\allowbreak{}4pp37–\allowbreak{}38, Remark3.\allowbreak{}6p38, Lemma3.\allowbreak{}7pp38–\allowbreak{}40.\allowbreak{} WD construction/\allowbreak{}unramified epsilon normalizationp35.\allowbreak{}.

\textbf{Inspection:} Read all of §3.\allowbreak{}1pp28–\allowbreak{}30 and §3.\allowbreak{}3pp34–\allowbreak{}40, including complete Lemma3.\allowbreak{}4 and Lemma3.\allowbreak{}7 proofs.\allowbreak{} Re-read the full determinant conventions and(22)–\allowbreak{}(24) while proving AG-LOCAL-SPLICE, checking actual Tate-dual Frobenius and WD rather than underlying-Weil epsilon.\allowbreak{} Upstream local epsilon construction and the full paper are not re-audited.\allowbreak{}

\url{https://arxiv.org/pdf/1305.0880v3}.

\paragraph{S-AG-BERGER-MON} Laurent Berger, Équations différentielles p-adiques et (phi,N)-modules filtrés, Astérisque 319 (2008), 13–\allowbreak{}38.

\textbf{Locator:} §I.\allowbreak{}1 printedp16 /\allowbreak{}PDF5:\allowbreak{} filtered module and N phi=p phi N conventions; §III.\allowbreak{}2 printedpp26–\allowbreak{}27 /\allowbreak{}PDF15–\allowbreak{}16:\allowbreak{} differential monodromy input and TheoremsIII.\allowbreak{}2.\allowbreak{}3–\allowbreak{}III.\allowbreak{}2.\allowbreak{}4; §V.\allowbreak{}1 printedp31 /\allowbreak{}PDF20 states de Rham iff potentially semistable for representations.\allowbreak{}.

\textbf{Inspection:} Read the exact definitions, relevant §III.\allowbreak{}2 statements and complete displayed proof of III.\allowbreak{}2.\allowbreak{}3, its immediate categorical consequenceIII.\allowbreak{}2.\allowbreak{}4, and the representation statement§V.\allowbreak{}1p31.\allowbreak{} Classical p-adic monodromy theorem and its upstream Andre/\allowbreak{}Kedlaya/\allowbreak{}Mebkhout proofs remain located inputs; no whole-source proof audit claimed.\allowbreak{}

\url{https://www.numdam.org/item/AST_2008__319__13_0.pdf}.

\paragraph{S-AG-ANDRE-KIMURA} Yves André, Motifs de dimension finie, Séminaire Bourbaki no.\allowbreak{}929 (2003–\allowbreak{}2004), Astérisque 299 (2005), 115–\allowbreak{}145.

\textbf{Locator:} §2.\allowbreak{}1pp118–\allowbreak{}120 rational Chow motives, transpose graph and super realization; Proposition2.\allowbreak{}1pp120–\allowbreak{}121 and Theorem2.\allowbreak{}8p124 abelian-type finiteness; Lemma3.\allowbreak{}13/\allowbreak{}Theorem3.\allowbreak{}14pp130–\allowbreak{}131 nilpotence; Cor3.\allowbreak{}19–\allowbreak{}3.\allowbreak{}22p132 dimension and categorical Cayley–\allowbreak{}Hamilton; §4.\allowbreak{}2p138 natural tensor Frobenius over finite fields.\allowbreak{}.

\textbf{Inspection:} Read complete printedpp118–\allowbreak{}124,130–\allowbreak{}133,138–\allowbreak{}139, including full displayed nilpotence/\allowbreak{}dimension/\allowbreak{}Cayley–\allowbreak{}Hamilton deductions, abelian-type deduction and finite-field Frobenius scope.\allowbreak{} Upstream general tensor-category and Shermenev/\allowbreak{}O’Sullivan results remain classical inputs, not a full original-source audit.\allowbreak{} AG-FROB-CHOW explicitly uses the monic reciprocal normalization and rederives its short partial-trace identity; coefficient field Q is distinct from geometric base Fq.\allowbreak{}

\url{https://www.numdam.org/item/SB_2003-2004__46__115_0.pdf}.

\paragraph{S-AG-WILSON-CHAINS} E.\allowbreak{} Ghorbani, G.\allowbreak{} B.\allowbreak{} Khosrovshahi, Ch.\allowbreak{} Maysoori and M.\allowbreak{} Mohammad-Noori, Inclusion Matrices and Chains, J.\allowbreak{} Combinatorial Theory A 115 (2008), 878–\allowbreak{}887; primary proof of Wilson’s 1990 diagonal-form theorem.

\textbf{Locator:} Author v1 (20Sep2007), pp3–\allowbreak{}6 rank-chain definitions; §4p7 unimodular conventions and equations(5)–\allowbreak{}(6); Theorem1 and whole inductionpp8–\allowbreak{}9; Corollary3 and whole proofp10.\allowbreak{} Original Wilson:\allowbreak{} European J.\allowbreak{} Combinatorics11(1990),609–\allowbreak{}615, DOI10.\allowbreak{}1016/\allowbreak{}S0195-6698(13)80046-7.\allowbreak{}.

\textbf{Inspection:} Read all authorpp1–\allowbreak{}11 in original-PDF text, including complete rank-chain construction and full Theorem1/\allowbreak{}Cor3 proofs.\allowbreak{} Confirmed integer unimodular equivalence, inclusive t\ensuremath{\le}k\ensuremath{\le}v\ensuremath{-}t range, diagonal entries and multiplicities.\allowbreak{} The 1990 original was bibliographically located but not obtained or claimed freshly read; the primary 2007 proof is the inspected input.\allowbreak{} No whole §5 audit or screenshot-inspection claim.\allowbreak{}

\url{https://arxiv.org/pdf/0709.3144}.

\paragraph{S-AG-BGF} Thorsten Beckmann and Olivier de Gaay Fortman, Integral Fourier transforms and the integral Hodge conjecture for one-cycles on abelian varieties, Compositio Mathematica159(2023),1188–\allowbreak{}1213.

\textbf{Locator:} Author final25-page version:\allowbreak{} Thm1.\allowbreak{}1p2; Def3.\allowbreak{}1 and Lem3.\allowbreak{}2pp6–\allowbreak{}7; Lem3.\allowbreak{}3–\allowbreak{}3.\allowbreak{}6pp7–\allowbreak{}10; Thm3.\allowbreak{}8pp10–\allowbreak{}12; Prop3.\allowbreak{}11/\allowbreak{}3.\allowbreak{}12pp12–\allowbreak{}13; full Thm1.\allowbreak{}1 proofp13; Cor4.\allowbreak{}1/\allowbreak{}Thm1.\allowbreak{}2p14.\allowbreak{} Distinguish rational-Chow lift from cohomological lift.\allowbreak{}.

\textbf{Inspection:} Read all relevant §3pp5–\allowbreak{}13 and §4.\allowbreak{}1pp13–\allowbreak{}15 in original PDF text, including full displayed Fourier identities, equivalence proofs and Jacobian-product deduction.\allowbreak{} The MP divided-power input was separately read in its original full proof.\allowbreak{} Beauville/\allowbreak{}Mukai/\allowbreak{}Orlov and Lefschetz(1,1) remain explicit classical upstream inputs; no whole §4.\allowbreak{}2–\allowbreak{}§6 audit claimed.\allowbreak{}

\url{https://www.math.ens.psl.eu/~odegaay/IHC-fourier-final.pdf}.

\paragraph{S-AG-EGS} Philip Engel, Olivier de Gaay Fortman and Stefan Schreieder, Matroids and the integral Hodge conjecture for abelian varieties, arXiv2507.\allowbreak{}15704v3 (27Mar2026; manuscript dated31Mar2026).

\textbf{Locator:} Thm1.\allowbreak{}1 and Cor1.\allowbreak{}2p2; Cor1.\allowbreak{}5p3; Thm1.\allowbreak{}6/\allowbreak{}1.\allowbreak{}8/\allowbreak{}1.\allowbreak{}9pp4–\allowbreak{}6; final proof chain§8.\allowbreak{}1–\allowbreak{}8.\allowbreak{}3pp77–\allowbreak{}80; full Cor1.\allowbreak{}5 proofp81.\allowbreak{} Thm1.\allowbreak{}1/\allowbreak{}Cor1.\allowbreak{}2 apply to very general complex ppav, not a supplied CM or number-field point.\allowbreak{}.

\textbf{Inspection:} Read entire introductionpp1–\allowbreak{}8, complete displayed final reductionspp77–\allowbreak{}80 and Cor1.\allowbreak{}5 proofp81, including exact dependence on Thms4.\allowbreak{}10/\allowbreak{}5.\allowbreak{}11/\allowbreak{}6.\allowbreak{}1/\allowbreak{}7.\allowbreak{}1 and the stated rank computations.\allowbreak{} Imported primary theorem, not a new audit of the complete85-page degeneration/\allowbreak{}combinatorial proof or Sage certificates; no theorem overQbar inferred.\allowbreak{}

\url{https://arxiv.org/pdf/2507.15704v3}.

\paragraph{S-AG-MP} Ben Moonen and Alexander Polishchuk, Divided powers in Chow rings and integral Fourier transforms, Advances in Mathematics224(2010),2216–\allowbreak{}2236.

\textbf{Locator:} Authorv1(25Apr2009),20pages:\allowbreak{} Thm1.\allowbreak{}6pp6–\allowbreak{}7/\allowbreak{}Cor1.\allowbreak{}7p7; Prop2.\allowbreak{}5pp10–\allowbreak{}11 and Rem2.\allowbreak{}6p11; integral motivic Def3.\allowbreak{}4p12; Lem3.\allowbreak{}6pp13–\allowbreak{}14; full Thm3.\allowbreak{}9pp15–\allowbreak{}17; Thm3.\allowbreak{}11pp17–\allowbreak{}19, exact d·modified-diagonal criterion.\allowbreak{}.

\textbf{Inspection:} Read all20pages in original PDF text.\allowbreak{} Targeted complete proof checks cover symmetric-power rational-equivalence descent and Thm1.\allowbreak{}6/\allowbreak{}Cor1.\allowbreak{}7, Prop2.\allowbreak{}5/\allowbreak{}Rem2.\allowbreak{}6, the motivic composition conventions and complete Thm3.\allowbreak{}9/\allowbreak{}3.\allowbreak{}11 proofs.\allowbreak{} Classical Fulton intersection theory and Beauville/\allowbreak{}Roitman torsion description remain cited upstream inputs.\allowbreak{} Explicit CM detector and F0 square were independently rederived in AG-FOURIER-CM2.\allowbreak{}

\url{https://arxiv.org/pdf/0904.3995}.

\paragraph{S-AG-K8-JOINT} Mathieu Dutour Sikirić, Philippe Elbaz-Vincent, Alexander Kupers and Jacques Martinet, Voronoi complexes in higher dimensions, cohomology of GL\_\allowbreak{}N(Z) for N\ensuremath{\ge}8 and triviality of K8(Z), JIMJ25(2026),565–\allowbreak{}590.

\textbf{Locator:} Author arXiv1910.\allowbreak{}11598v1(25Oct2019),22pp:\allowbreak{} Prop6.\allowbreak{}1p18, Prop6.\allowbreak{}3/\allowbreak{}Thm6.\allowbreak{}4p19, Cor7.\allowbreak{}1p20.\allowbreak{} Journalmetadata confirms online24Nov2025, issueJan2026; author-page locators are not journal-page locators.\allowbreak{}.

\textbf{Inspection:} Read complete final reduction proofs Props6.\allowbreak{}1/\allowbreak{}6.\allowbreak{}3 and Thm6.\allowbreak{}4 plus §7 on authorpp18–\allowbreak{}20 and introduction.\allowbreak{} Bibliographic journal record verified directly.\allowbreak{} PublishedPDF endpoint returnedHTML, so no claim of inspecting final26-page journalPDF.\allowbreak{} Upstream Voronoi enumeration/\allowbreak{}Arlettaz theorem remain cited inputs; explicit2-primary step independently supplied from WeibelVI9.\allowbreak{}8.\allowbreak{}

\url{https://jamartin.perso.math.cnrs.fr/Publications/DEKM_arXiv.pdf}.

\paragraph{S-AG-K8-SHORT} Alexander Kupers, A short proof that K8(Z)=0, author note dated21October2019.

\textbf{Locator:} Entire4-page authornote; TheoremAp1, rank-filtration table/\allowbreak{}Lemmas2.\allowbreak{}1–\allowbreak{}2.\allowbreak{}2p2, complete proof of Lemma2.\allowbreak{}2/\allowbreak{}Lemma2.\allowbreak{}5/\allowbreak{}TheoremApp3–\allowbreak{}4.\allowbreak{}.

\textbf{Inspection:} Read all4pages and checked the complete argument relative to its stated earlier homology inputs; the Serre transgression from(9,0)to(0,8) isd9 in the usual homological convention, correcting printedd8 without changing the proof.\allowbreak{} Did not freshly audit prior Voronoi/\allowbreak{}Church–\allowbreak{}Putman computations.\allowbreak{} Webopen403, but originalauthorPDF successfully read through in-memory download.\allowbreak{}

\url{https://www.utsc.utoronto.ca/people/kupers/wp-content/uploads/sites/50/2021/01/k8zshorter.pdf}.

\paragraph{S-AG-WEIBEL-K} Charles A.\allowbreak{} Weibel, The K-book, ChapterVI:\allowbreak{} The higher K-theory of fields, author chapter PDF.

\textbf{Locator:} 66-page chapterPDF:\allowbreak{} §8 Theorem8.\allowbreak{}2p47/\allowbreak{}Corollary8.\allowbreak{}3p48, Brauer sequence8.\allowbreak{}1.\allowbreak{}1p47/\allowbreak{}Example8.\allowbreak{}3.\allowbreak{}2p48; Corollary9.\allowbreak{}8p55; Theorem10.\allowbreak{}1 and low-degree tablep60; Proposition10.\allowbreak{}5p63; Theorem10.\allowbreak{}9p65.\allowbreak{} These are chapterpage numbers.\allowbreak{}.

\textbf{Inspection:} Read complete indicated source proofs:\allowbreak{}8.\allowbreak{}2/\allowbreak{}8.\allowbreak{}3,9.\allowbreak{}8,10.\allowbreak{}5 and10.\allowbreak{}9; low-degree table and exact Brauer formula checked.\allowbreak{} Upstream norm-residue/\allowbreak{}motivic spectral sequence and global classfield theory remain classical inputs, not freshly audited.\allowbreak{} New full K8-TWIST transfer checked separately.\allowbreak{}

\url{https://sites.math.rutgers.edu/~weibel/Kbook/Kbook.VI.pdf}.

\paragraph{CX-S-KOIRAN} Pascal Koiran, Hilbert’s Nullstellensatz is in the Polynomial Hierarchy, DIMACS 96-27.

\textbf{Locator:} Section 2 Theorem 2 pp.\allowbreak{}5–\allowbreak{}6; Section 5 Theorem 9 and Corollary 1 printed p.\allowbreak{}12 (PDF13); 8 August 1997 erratum p.\allowbreak{}1.

\textbf{Inspection:} Downloaded both author PS files, converted to PDF, read theorem statements, protocol proof, number-field input, entire erratum.\allowbreak{}

\url{https://perso.ens-lyon.fr/pascal.koiran/Publis/DIMACS.96-27.ps}.

\paragraph{CX-S-BACH} Eric Bach, Explicit bounds for primality testing and related problems, Math.\allowbreak{} Comp.\allowbreak{} 55 (1990), 355–\allowbreak{}380.

\textbf{Locator:} Section 6 Theorems 2–\allowbreak{}3 p.\allowbreak{}372, Miller consequence p.\allowbreak{}373; Section 7 Theorem 4 p.\allowbreak{}375, corollary p.\allowbreak{}376.

\textbf{Inspection:} Original AMS PDF returned 403.\allowbreak{} Exact locators and statements carried from CX\_\allowbreak{}GRH/\allowbreak{}ALPOGE\_\allowbreak{}GCT; not represented as freshly checked original proof.\allowbreak{}

\url{https://doi.org/10.1090/S0025-5718-1990-1023756-8}.

\paragraph{CX-S-SCHOOF} René Schoof, Four primality testing algorithms, MSRI Publications 44 (2008).

\textbf{Locator:} Theorem 2.\allowbreak{}2 pp.\allowbreak{}105–\allowbreak{}106, including full proof specifying quadratic Dirichlet characters.

\textbf{Inspection:} Read statement and full proof via primary PDF text.\allowbreak{} PDF text loses minus signs; transfer uses standard Miller test and the theorem’s stated cutoff.\allowbreak{}

\url{https://pub.math.leidenuniv.nl/~stevenhagenp/ANTproc/05rene.pdf}.

\paragraph{CX-S-AKS} Agrawal, Kayal and Saxena, PRIMES is in P, Ann.\allowbreak{} Math.\allowbreak{} 160 (2004), 781–\allowbreak{}793.

\textbf{Locator:} Theorems 4.\allowbreak{}1 and 5.\allowbreak{}1; 2019 erratum Ann.\allowbreak{} Math.\allowbreak{} 189, 317–\allowbreak{}318.

\textbf{Inspection:} Primary article and erratum publication located; used only for unconditional P classification.\allowbreak{}

\url{https://annals.math.princeton.edu/wp-content/uploads/annals-v160-n2-p12.pdf}.

\paragraph{CX-S-SHOUP} Victor Shoup, Searching for Primitive Roots in Finite Fields, Math.\allowbreak{} Comp.\allowbreak{} 58 (1992), 369–\allowbreak{}380.

\textbf{Locator:} Author version Theorem 3 p.\allowbreak{}4, Section 6 p.\allowbreak{}13; Proposition 9 p.\allowbreak{}10.

\textbf{Inspection:} Downloaded author PDF and visually inspected p.\allowbreak{}4:\allowbreak{} Theorem3 is O(r\textasciicircum{}4(log r+1)\textasciicircum{}4(log p)\textasciicircum{}2), r=omega(p\ensuremath{-}1).\allowbreak{} Text encoding is corrupted; image is authoritative.\allowbreak{} No full original proof check claimed.\allowbreak{}

\url{https://shoup.net/papers/primroots.pdf}.

\paragraph{CX-S-HAC} Menezes, van Oorschot and Vanstone, Handbook of Applied Cryptography (1997).

\textbf{Locator:} Algorithm 3.\allowbreak{}34 and Fact 3.\allowbreak{}35 p.\allowbreak{}100; Fact 3.\allowbreak{}46 (composite square-root/\allowbreak{}factoring relation).

\textbf{Inspection:} Primary author-hosted chapter, exact square-root algorithm and its random step checked.\allowbreak{}

\url{https://cacr.uwaterloo.ca/hac/about/chap3.pdf}.

\paragraph{CX-S-GCTV} Ketan Mulmuley, Geometric Complexity Theory V:\allowbreak{} Efficient Algorithms for Noether Normalization, arXiv:\allowbreak{}1209.\allowbreak{}5993v5.

\textbf{Locator:} Theorem 2.\allowbreak{}10 p.\allowbreak{}25; Theorem 3.\allowbreak{}4 pp.\allowbreak{}27–\allowbreak{}28; Theorem 4.\allowbreak{}1 p.\allowbreak{}30; Theorem 4.\allowbreak{}7/\allowbreak{}Lemma 4.\allowbreak{}8 p.\allowbreak{}33; Theorem 4.\allowbreak{}10 p.\allowbreak{}34; Theorem9.\allowbreak{}8 p68 and following Remark3 p69.

\textbf{Inspection:} Primary PDF downloaded and relevant statements/\allowbreak{}proof of3.\allowbreak{}4 read; output and equation-access qualifications checked.\allowbreak{} Later complete Theorem9.\allowbreak{}8 proof and following Remark3 pp68–\allowbreak{}69 read for the orbit-intersection interface; its explicit-quotient/\allowbreak{}fixed-rank hypotheses retained.\allowbreak{}

\url{https://arxiv.org/pdf/1209.5993}.

\paragraph{CX-S-FS} Michael Forbes and Amir Shpilka, A PSPACE Construction of a Hitting Set for the Closure of Small Algebraic Circuits, ECCC TR17-163.

\textbf{Locator:} Theorem 7.\allowbreak{}1 p.\allowbreak{}21; introduction Theorems 1.\allowbreak{}2–\allowbreak{}1.\allowbreak{}3.

\textbf{Inspection:} Primary theorem and unary-encoding condition checked.\allowbreak{}

\url{https://eccc.weizmann.ac.il/report/2017/163/download/}.

\paragraph{CX-S-GSS} Guo, Saxena and Sinhababu, Algebraic Dependencies and PSPACE Algorithms in Approximative Complexity over Any Field, Theory of Computing 15(16) (2019).

\textbf{Locator:} Theorem1.\allowbreak{}3 p5; Lemma4.\allowbreak{}1/\allowbreak{}Theorem4.\allowbreak{}2 pp15–\allowbreak{}16; Theorems1.\allowbreak{}5–\allowbreak{}1.\allowbreak{}6 p6.

\textbf{Inspection:} Primary journal PDF freshly retrieved via theoryofcomputing.\allowbreak{}org; checked APS statement/\allowbreak{}closure characterization and precise separating-e.\allowbreak{}s.\allowbreak{}o.\allowbreak{}p.\allowbreak{} statements in Theorems1.\allowbreak{}5–\allowbreak{}1.\allowbreak{}6.\allowbreak{}

\url{https://theoryofcomputing.org/articles/v015a016/v015a016.pdf}.

\paragraph{CX-S-IKRS} Ivanyos, Karpinski, Rónyai and Saxena, Trading GRH for algebra:\allowbreak{} algorithms for factoring polynomials and related structures.

\textbf{Locator:} 2008 author version Main Theorem p3, Lemma4.\allowbreak{}2 p14, Proposition4.\allowbreak{}6 pp16–\allowbreak{}17; published authorPDF Theorems1.\allowbreak{}1/\allowbreak{}1.\allowbreak{}3 p3, Lemma2.\allowbreak{}3 pp9–\allowbreak{}10, Lemma4.\allowbreak{}2 p16, Proposition4.\allowbreak{}6 pp18–\allowbreak{}19, Theorem5.\allowbreak{}8 pp28–\allowbreak{}29.

\textbf{Inspection:} Published primary statements freshly checked.\allowbreak{} Complete proofs of Lemma2.\allowbreak{}3, Lemma4.\allowbreak{}2 and Proposition4.\allowbreak{}6 read for the new composition; Theorem5.\allowbreak{}8 exact statement/\allowbreak{}interface and runtime checked, not a fresh whole-paper proof audit.\allowbreak{} Older-version locators retained separately.\allowbreak{}

\url{https://www.cse.iitk.ac.in/users/nitin/papers/nogrh.pdf}.

\paragraph{CX-S-BS} Eric Bach and Jonathan Sorenson, Explicit Bounds for Primes in Residue Classes, Math.\allowbreak{} Comp.\allowbreak{}65 (1996),1717–\allowbreak{}1735.

\textbf{Locator:} Section2 p.\allowbreak{}1720; Theorem5.\allowbreak{}1 p.\allowbreak{}1729.

\textbf{Inspection:} Primary author-uploaded PDF retrieved; statement and precise analytic premise read; geometry agent independently verified same theorem.\allowbreak{}

\url{https://www.researchgate.net/profile/Jonathan-Sorenson-2/publication/220576951_Explicit_bounds_for_primes_in_residue_classes/links/55100ea70cf2ac2905afbceb/Explicit-bounds-for-primes-in-residue-classes.pdf}.

\paragraph{CX-S-MILLER} Gary L.\allowbreak{} Miller, Riemann’s Hypothesis and Tests for Primality, JCSS13(3) (1976),300–\allowbreak{}317.

\textbf{Locator:} Lemma5 and Theorem3 p310; Theorem4 pp311–\allowbreak{}312.

\textbf{Inspection:} Author-hosted primary PDF downloaded; p310 visually inspected and entire Lemma5 proof read.\allowbreak{} \ensuremath{\lambda}\ensuremath{{}^{\prime}} multiple and polynomial bit-growth conditions confirmed; explicit-input version proved here.\allowbreak{}

\url{https://www.cs.cmu.edu/~glmiller/Publications/Papers/Mi76.pdf}.

\paragraph{CX-S-BIMS} Bhargava, Ivanyos, Mittal and Saxena, Irreducibility and deterministic r-th root finding over finite fields, ISSAC2017.

\textbf{Locator:} Corollary1.\allowbreak{}3 authorPDF p2; Theorem3.\allowbreak{}1 and full proof p3; Corollary3.\allowbreak{}2 p4.

\textbf{Inspection:} Statement and complete Theorem3.\allowbreak{}1 proof read; original p3 rendered and inspected.\allowbreak{} Literal nonvanishing step fails for the explicit F7 cubic example CX-ALG-016.\allowbreak{} Corrected basis-search proof CX-ALG-017; even-degree quadratic case independently re-proved CX-ALG-018.\allowbreak{} Corollary conclusions not refuted.\allowbreak{}

\url{https://www.cse.iitk.ac.in/users/nitin/papers/irred-and-rth-root.pdf}.

\paragraph{CX-S-BE} Jean-François Biasse and Muhammed Rashad Erukulangara, A proof of the conjectured run time of the Hafner–\allowbreak{}McCurley class group algorithm, Advances in Mathematics of Communications17(6) (2023),1388–\allowbreak{}1405.

\textbf{Locator:} Author PDF MainTheorem and Conjecture1 p4; Algorithm4 p15; Proposition6 pp15–\allowbreak{}16; Corollary1 p16.

\textbf{Inspection:} Direct PDF request returned503 but web source recovered through indexed reference.\allowbreak{} Main theorem, explicit Hecke-GRH definition, complete Algorithm4, Proposition6 and Corollary1 with its proof read.\allowbreak{} Entire 18-page proof not claimed re-audited; prior CX-N12 detailed transfer retained.\allowbreak{}

\url{https://par.nsf.gov/servlets/purl/10324557}.

\paragraph{CX-S-JMV} David Jao, Stephen D.\allowbreak{} Miller and Ramarathnam Venkatesan, Expander graphs based on GRH with an application to elliptic curve cryptography, J.\allowbreak{} Number Theory129 (2009),1491–\allowbreak{}1504.

\textbf{Locator:} Author arXiv0811.\allowbreak{}0647v2:\allowbreak{} Theorem1.\allowbreak{}1 p3; Remark1.\allowbreak{}2(a)/\allowbreak{}Corollary1.\allowbreak{}3 p4; Theorem1.\allowbreak{}5 p5; Theorem1.\allowbreak{}6 pp6–\allowbreak{}7; Lemma2.\allowbreak{}1 pp7–\allowbreak{}8; proof of Theorem1.\allowbreak{}1 p9 and Cor1.\allowbreak{}3 p10; Proposition3.\allowbreak{}1 p11, Theorem3.\allowbreak{}3 p12.

\textbf{Inspection:} Primary author-hosted PDF read for all listed statements.\allowbreak{} Full proofs of Theorem1.\allowbreak{}1, Cor1.\allowbreak{}3 and Lemma2.\allowbreak{}1 read; verified estimates are character-by-character, enabling the genus quotient refinement.\allowbreak{} Isogeny statements and supplied-solver scope checked; no full underlying CM/\allowbreak{}Deuring proof audit claimed.\allowbreak{}

\url{https://web.williams.edu/Mathematics/sjmiller/public_html/crypto/handouts/JaoMillerVenkatesan_ExpanderGraphsGRHEllCurveCrypto.pdf}.

\paragraph{CX-S-L91} Hendrik W.\allowbreak{} Lenstra Jr.\allowbreak{}, Finding isomorphisms between finite fields, Mathematics of Computation56 (1991),329–\allowbreak{}347.

\textbf{Locator:} Theorems1.\allowbreak{}1/\allowbreak{}1.\allowbreak{}2 p330; subfields §2p332; Theorem3.\allowbreak{}1 pp334–\allowbreak{}336; Proposition4.\allowbreak{}2 p337; Theorems5.\allowbreak{}1/\allowbreak{}5.\allowbreak{}2 pp338–\allowbreak{}340; Theorem7.\allowbreak{}1 p342; Theorem9.\allowbreak{}1 pp344–\allowbreak{}346, Lemma9.\allowbreak{}2 p345.

\textbf{Inspection:} Author-hosted primary PDF downloaded.\allowbreak{} Complete normal-basis proof, prime-degree Teichmüller construction and converse, Artin–\allowbreak{}Schreier proof and final field-construction proof read.\allowbreak{} pp338–\allowbreak{}339 rendered and visually read to recover formula exponents lost in text extraction.\allowbreak{} New application and rewritten full proof in CX-ALG-023/\allowbreak{}026; no claim of new priority for these classical algorithms.\allowbreak{}

\url{https://pub.math.leidenuniv.nl/~lenstrahw/PUBLICATIONS/1991b/art.pdf}.

\paragraph{CX-S-AL86} Leonard M.\allowbreak{} Adleman and Hendrik W.\allowbreak{} Lenstra Jr.\allowbreak{}, Finding irreducible polynomials over finite fields, STOC1986,350–\allowbreak{}355.

\textbf{Locator:} Theorem1 p350; AlgorithmA pp350–\allowbreak{}351; Proposition3 and polynomial-time analysis p351; Lemma6 p352 (Gaussian periods under inertness).

\textbf{Inspection:} Fresh primary author PDF read through AlgorithmA and its complete correctness proof pp350–\allowbreak{}352.\allowbreak{} For prime degree r\ensuremath{\ne}p only steps2–\allowbreak{}3 and Lemma6 are needed; least-prime existence is replaced by the exact CX-ALG-024 input and bound.\allowbreak{} No extra radical or Hecke oracle is hidden in the period construction.\allowbreak{}

\url{https://pub.math.leidenuniv.nl/~lenstrahw/PUBLICATIONS/1986a/art.pdf}.

\paragraph{CX-S-BHMR} Michael Bate, Sebastian Herpel, Benjamin Martin and Gerhard Röhrle, Cocharacter-closure and the rational Hilbert–\allowbreak{}Mumford Theorem, arXiv1411.\allowbreak{}7849v5 (2016).

\textbf{Locator:} Introduction p1 (unique closed orbit and classical Hilbert–\allowbreak{}Mumford formulation); Theorem4.\allowbreak{}3 with complete proof pp12–\allowbreak{}13; Corollary4.\allowbreak{}5 pp13–\allowbreak{}14.

\textbf{Inspection:} Primary theorem statement and complete Theorem4.\allowbreak{}3 proof read.\allowbreak{} Only its algebraically closed-field specialization is used.\allowbreak{} Kempf1978 Theorem1.\allowbreak{}4 is identified by the source, but the original published proof was not retrieved; no original-proof check claimed.\allowbreak{} Elementary bounded-weight and HN reductions are proved independently in CX-GEO-010/\allowbreak{}011.\allowbreak{}

\url{https://arxiv.org/pdf/1411.7849v5}.

\paragraph{CX-S-PF26} OpenAI, Deterministic Polynomial Factorization over Prime Fields, 4 October2026.

\textbf{Locator:} Theorem1.\allowbreak{}2p3; complete algebraic proof §§2–\allowbreak{}10pp6–\allowbreak{}44; Lemma2.\allowbreak{}2pp8–\allowbreak{}9; Proposition2.\allowbreak{}3pp9–\allowbreak{}10; §4pp13–\allowbreak{}19; Theorem6.\allowbreak{}1 with Lemmas6.\allowbreak{}2–\allowbreak{}6.\allowbreak{}5pp24–\allowbreak{}32; Proposition7.\allowbreak{}1pp33–\allowbreak{}34; Proposition8.\allowbreak{}1pp35–\allowbreak{}37; Lemma9.\allowbreak{}1p38; Proposition10.\allowbreak{}1pp41–\allowbreak{}44.

\textbf{Inspection:} Complete invoked algebraic proof §§2–\allowbreak{}10 read from primary mirror and local PDF, with full post-reading reasoning checks.\allowbreak{} Root independently read complete §6 and §§2.\allowbreak{}1–\allowbreak{}2.\allowbreak{}3; geometry independently read complete §4.\allowbreak{} No gap found in these readings.\allowbreak{} Recent algebraic assertions remain explicit extra input PF26; source companion analytic theorem not used or audited here; no formal certification or implementation.\allowbreak{} Tang2611.\allowbreak{}00135v1 led to discovery but its newer analytic inputs are not adopted.\allowbreak{}

\url{https://github.com/openai/math/blob/main/preprints/Deterministic-Polynomial-Factorization-over-Prime-Fields-October-4-2026/Deterministic-Polynomial-Factorization-over-Prime-Fields.pdf}.

\paragraph{CX-S-MEYER-IDELE} Ralf Meyer, On a representation of the idele class group related to primes and zeros of L-functions, Duke Math.\allowbreak{}J.\allowbreak{}127(2005),519–\allowbreak{}595; arXiv:\allowbreak{}math/\allowbreak{}0311468v3.

\textbf{Locator:} §3.\allowbreak{}1 Haar normalization p16; Definition4.\allowbreak{}1/\allowbreak{}(22)pp24–\allowbreak{}25; (24)p28; Lemma4.\allowbreak{}20p36; Theorem5.\allowbreak{}1pp38–\allowbreak{}39; Lemma5.\allowbreak{}4pp40–\allowbreak{}41; Theorem5.\allowbreak{}8pp43–\allowbreak{}44; Lemma5.\allowbreak{}10p45 and exact image criterionpp46–\allowbreak{}47; Theorem5.\allowbreak{}11p47.\allowbreak{} Author page equals PDF page.\allowbreak{}.

\textbf{Inspection:} Read complete provided Lemma4.\allowbreak{}20, Theorem5.\allowbreak{}1, §5.\allowbreak{}3, Theorem5.\allowbreak{}8 and §5.\allowbreak{}7 arguments, plus exact weighted-space and Haar conventions.\allowbreak{} Rendered and visually read original pp46–\allowbreak{}47, verifying inverse Euler factor and reduced image criterion at Re0.\allowbreak{} Earlier homological/\allowbreak{}local trace inputs are classical cited dependencies; no complete-paper audit claimed.\allowbreak{} New fixed-sector transfer MT-M22 is proved in full separately.\allowbreak{}

\url{https://arxiv.org/pdf/math/0311468}.

\paragraph{CX-S-MASULLO-QUOTIENT} Masullo, The Music of the Primes, supplied 2026-11-21 draft:\allowbreak{} specified quotient/\allowbreak{}test-space definitions.

\textbf{Locator:} ["(2.\allowbreak{}2.\allowbreak{}3.\allowbreak{}1)pp81–\allowbreak{}82", "(2.\allowbreak{}2.\allowbreak{}4.\allowbreak{}23)p91", "(2.\allowbreak{}2.\allowbreak{}5.\allowbreak{}2)p95", "(7.\allowbreak{}1.\allowbreak{}4.\allowbreak{}7)p327", "(1.\allowbreak{}8.\allowbreak{}3.\allowbreak{}1)p45", "DefinitionD.\allowbreak{}3.\allowbreak{}3.\allowbreak{}1p804", "DefinitionD.\allowbreak{}3.\allowbreak{}3.\allowbreak{}3p805", "(9.\allowbreak{}1.\allowbreak{}2.\allowbreak{}3)p399", "(9.\allowbreak{}2.\allowbreak{}2.\allowbreak{}31)p416", "Definition10.\allowbreak{}1.\allowbreak{}1.\allowbreak{}1p434", "(10.\allowbreak{}2.\allowbreak{}1.\allowbreak{}1)–\allowbreak{}(10.\allowbreak{}2.\allowbreak{}1.\allowbreak{}2)p439", "(10.\allowbreak{}2.\allowbreak{}1.\allowbreak{}6)p440", "(10.\allowbreak{}4.\allowbreak{}1.\allowbreak{}3)p456", "§1.\allowbreak{}4.\allowbreak{}6p30", "§1.\allowbreak{}4.\allowbreak{}8–\allowbreak{}1.\allowbreak{}4.\allowbreak{}10pp31–\allowbreak{}32", "(10.\allowbreak{}3.\allowbreak{}1.\allowbreak{}17)p450", "Theorem12.\allowbreak{}5.\allowbreak{}1.\allowbreak{}1pp519–\allowbreak{}520", "§1.\allowbreak{}7.\allowbreak{}2/\allowbreak{}(1.\allowbreak{}7.\allowbreak{}2.\allowbreak{}1)p35", "Theorem7.\allowbreak{}1.\allowbreak{}3.\allowbreak{}2pp324–\allowbreak{}325", "Definition7.\allowbreak{}3.\allowbreak{}1.\allowbreak{}4p343", "Theorem7.\allowbreak{}3.\allowbreak{}1.\allowbreak{}6pp344–\allowbreak{}350", "Proposition11.\allowbreak{}2.\allowbreak{}0.\allowbreak{}1/\allowbreak{}(11.\allowbreak{}2.\allowbreak{}0.\allowbreak{}3)pp484–\allowbreak{}485", "Theorem11.\allowbreak{}2.\allowbreak{}0.\allowbreak{}2pp485–\allowbreak{}486", "Proposition11.\allowbreak{}2.\allowbreak{}0.\allowbreak{}4p486", "Theorem11.\allowbreak{}2.\allowbreak{}2.\allowbreak{}1/\allowbreak{}Remark11.\allowbreak{}2.\allowbreak{}2.\allowbreak{}2pp490–\allowbreak{}491", "Theorem11.\allowbreak{}3.\allowbreak{}2.\allowbreak{}1pp495–\allowbreak{}496", "Expectation1.\allowbreak{}7.\allowbreak{}1pp37–\allowbreak{}38", "Definition9.\allowbreak{}3.\allowbreak{}3.\allowbreak{}3/\allowbreak{}(9.\allowbreak{}3.\allowbreak{}3.\allowbreak{}4)–\allowbreak{}(9.\allowbreak{}3.\allowbreak{}3.\allowbreak{}5)p428", "Theorem11.\allowbreak{}1.\allowbreak{}2.\allowbreak{}1(2)/\allowbreak{}(11.\allowbreak{}1.\allowbreak{}2.\allowbreak{}3)–\allowbreak{}(11.\allowbreak{}1.\allowbreak{}2.\allowbreak{}6)p480", "proof(11.\allowbreak{}1.\allowbreak{}2.\allowbreak{}16)and following strong-limit inference p482", "(11.\allowbreak{}3.\allowbreak{}2.\allowbreak{}2)p495", "(A.\allowbreak{}2.\allowbreak{}3.\allowbreak{}14)p549 subspace topology", "DefinitionA.\allowbreak{}3.\allowbreak{}2.\allowbreak{}1/\allowbreak{}(A.\allowbreak{}3.\allowbreak{}2.\allowbreak{}1)pp559–\allowbreak{}560", "Lemma12.\allowbreak{}2.\allowbreak{}1.\allowbreak{}1/\allowbreak{}(12.\allowbreak{}2.\allowbreak{}1.\allowbreak{}2)pp504–\allowbreak{}505", "Definition12.\allowbreak{}2.\allowbreak{}1.\allowbreak{}3p506", "Proposition12.\allowbreak{}2.\allowbreak{}2.\allowbreak{}2/\allowbreak{}(12.\allowbreak{}2.\allowbreak{}2.\allowbreak{}1)–\allowbreak{}(12.\allowbreak{}2.\allowbreak{}2.\allowbreak{}2)p508", "Proposition12.\allowbreak{}2.\allowbreak{}2.\allowbreak{}4/\allowbreak{}(12.\allowbreak{}2.\allowbreak{}2.\allowbreak{}26)–\allowbreak{}(12.\allowbreak{}2.\allowbreak{}2.\allowbreak{}31)pp512–\allowbreak{}513"].

\textbf{Inspection:} 


\paragraph{CX-S-RV} Ramakrishnan–\allowbreak{}Valenza, Fourier Analysis on Number Fields, GTM186(1999).

\textbf{Locator:} ["§7.\allowbreak{}1 local factors printed244–\allowbreak{}245 (PDF264–\allowbreak{}265)", "explicit real/\allowbreak{}complex local integral calculations printed250–\allowbreak{}252 (PDF270–\allowbreak{}272)", "Theorem7-19 entire statement and supplied proof printed279–\allowbreak{}280 (PDF299–\allowbreak{}300)"].

\textbf{Inspection:} Fresh full reading of supplied Theorem7-19 proof and local-factor definitions/\allowbreak{}calculations; original printed244,245,252 rendered and visually checked.\allowbreak{} The finite-order property is classical per-function background as in NT-L03/\allowbreak{}NT-Z17; no separate uniform finite-order theorem attributed to RV7-19.\allowbreak{} M28 proves all required uniform strip, gamma-ratio and zero-block bounds directly.\allowbreak{} No whole-book audit.\allowbreak{}

\url{https://webusers.imj-prg.fr/~julien.marche/Ref_GL_GT/Fourier_Analysis_on_Number_Fields.pdf}.

\paragraph{CX-S-ALPOGE-NOTE} Levent Alpöge, Remark, supplied four-page note.

\textbf{Locator:} ["step4 PDF3:\allowbreak{} determinant positivity", "step7 PDF4:\allowbreak{} determinant/\allowbreak{}Gan–\allowbreak{}Gross–\allowbreak{}Prasad periods", "step8 PDF4:\allowbreak{} iterated theta lifts"].

\textbf{Inspection:} Read all four pages in supplied PDF text extraction, including exact steps4,7,8.\allowbreak{} A\_\allowbreak{}D remains the user-designated Dirichlet-GRH theorem input; no audit of the Lean archive, complete internal proof, or GGP determinant comparison is claimed.\allowbreak{} M29 distinguishes its quotient-derived spectral determinant from these period mechanisms.\allowbreak{}

\url{https://alpo.ge/remark.pdf}.

\paragraph{CX-S-PY07} Dipendra Prasad and C.\allowbreak{} S.\allowbreak{} Yogananda, A Report on Artin’s holomorphy conjecture (23 February 2007).

\textbf{Locator:} ["§2.\allowbreak{}1 Properties1–\allowbreak{}2 PDF3:\allowbreak{} all-prime Artin multiplication and induction", "§3.\allowbreak{}1 PDF4–\allowbreak{}5:\allowbreak{} Artin reciprocity and Hecke entireness", "Theorem3.\allowbreak{}3 PDF5:\allowbreak{} Brauer meromorphicity"].

\textbf{Inspection:} The exact cited primary statements were freshly read.\allowbreak{} Classical proofs/\allowbreak{}dependencies are credited as inputs, not re-audited.\allowbreak{} Older historical status discussions are not used for present scope.\allowbreak{}

\url{https://www.math.iitb.ac.in/~dprasad/artin.pdf}.

\paragraph{CX-S-CALEGARI13} Frank Calegari, The Artin conjecture for some S5-extensions, Math.\allowbreak{} Ann.\allowbreak{}356(2013),191–\allowbreak{}207.

\textbf{Locator:} ["Theorem1.\allowbreak{}1 author-hosted publisher PDF2:\allowbreak{} exact local hypotheses and four-dimensional automorphy", "§2 character table author-hosted publisher PDF4", "Construction of varpi publisher PDF10; Theorem4.\allowbreak{}7 PDF12", "Lemma4.\allowbreak{}8 PDF13–\allowbreak{}14; Landau Lemma4.\allowbreak{}9 PDF14; exact-comparison Lemma4.\allowbreak{}10 PDF14–\allowbreak{}15", "Final Theorem1.\allowbreak{}2 proof PDF15:\allowbreak{} simultaneous exact transfer for quadratic twist"].

\textbf{Inspection:} The exact theorem statement and complete character table were freshly read; the automorphic theorem is cited rather than re-proved.\allowbreak{} MT-S505 is retained as a separately checked antecedent for the M\_\allowbreak{}H-conditional exclusion of real poles.\allowbreak{}

\url{https://math.uchicago.edu/~fcale/papers/S5.pdf}.

\paragraph{A-SOURCE} Levent Alpoge, Remark.

\textbf{Locator:} PDF pp3–\allowbreak{}4 steps2–\allowbreak{}20; ALPOGE\_\allowbreak{}Main.\allowbreak{}lean five declarations; ALPOGE\_\allowbreak{}Statements.\allowbreak{}lean definitions.

\textbf{Inspection:} User-supplied PDF text pp1–\allowbreak{}4 read; declarations and definitions read; no independent Lean kernel run.\allowbreak{} Cited theorem input A\_\allowbreak{}D as instructed.\allowbreak{}

\url{https://alpo.ge/remark.pdf}.

\paragraph{M-SOURCE} A.\allowbreak{}M.\allowbreak{} Masullo, The Music of the Primes, supplied 940-page draft.

\textbf{Locator:} §1.\allowbreak{}2 pp12–\allowbreak{}19; Lemma9.\allowbreak{}3.\allowbreak{}2.\allowbreak{}4 p425; Lemma9.\allowbreak{}3.\allowbreak{}2.\allowbreak{}5 pp425–\allowbreak{}426; Theorem9.\allowbreak{}3.\allowbreak{}2.\allowbreak{}6 p426; DefD.\allowbreak{}3.\allowbreak{}3.\allowbreak{}1 p804; Theorem12.\allowbreak{}5.\allowbreak{}1.\allowbreak{}1 pp519–\allowbreak{}520;13.\allowbreak{}3.\allowbreak{}1.\allowbreak{}1 and placeholder13.\allowbreak{}3.\allowbreak{}1.\allowbreak{}2 p522.

\textbf{Inspection:} Original PDF425,426 images inspected; source definitions11,804 read.\allowbreak{} Claims M remain conditional, independent intermediate assertions not inferred from endpoint.\allowbreak{} Original p426 re-inspected for Lemma9.\allowbreak{}3.\allowbreak{}2.\allowbreak{}5 actual ideal versus closure; MT-M05 refutes its quotient claim independently of the local zero-count issue.\allowbreak{}


\paragraph{KW09} Chandrashekhar Khare and Jean-Pierre Wintenberger, Serre's modularity conjecture (I).

\textbf{Locator:} Theorem10.\allowbreak{}1(ii), author PDF20; Cor10.\allowbreak{}2(ii), author PDF21.

\textbf{Inspection:} Primary source theorem and Corollary read, including A5 scope and Artin-L-function discussion.\allowbreak{}

\url{https://www.math.ucla.edu/~shekhar/papers/results.pdf}.

\paragraph{STARK74} Harold M.\allowbreak{} Stark, Some effective cases of the Brauer-Siegel theorem, Invent.\allowbreak{}Math.\allowbreak{}23(1974),135–\allowbreak{}152.

\textbf{Locator:} Theorem3,p144.

\textbf{Inspection:} Exact source locator corroborated by primary later paper Gun-Hazra-Sahu2025 introduction; original volume web fetch timed out.\allowbreak{} Classical theorem statement used; no claim to have read original proof.\allowbreak{}

\url{https://gdz.sub.uni-goettingen.de/download/pdf/PPN356556735_0023/PPN356556735_0023.pdf}.

\paragraph{GHS25} Sanoli Gun, Suhita Hazra and Dhananjaya Sahu, On holomorphy and non-vanishing of Artin L-functions, Monatshefte fur Mathematik206(2025),853–\allowbreak{}883.

\textbf{Locator:} Theorem5 and Corollary6; §6; introduction for Stark statement.

\textbf{Inspection:} Primary statements and proof mechanism of integer-root Euler deformation checked.\allowbreak{} MT-L01/\allowbreak{}L02 gives own proof for general |c|\ensuremath{\le}1 and irrational c.\allowbreak{}

\url{https://link.springer.com/article/10.1007/s00605-025-02060-7}.

\paragraph{GM19-A1} Loic Grenie and Giuseppe Molteni, An explicit Chebotarev density theorem under GRH.

\textbf{Locator:} AppendixA TheoremA.\allowbreak{}1, author PDF26.

\textbf{Inspection:} Number\_\allowbreak{}theory worker verified primary finite-order Hecke zero count; root uses the result as a cited classical input and subtracts consecutive-height counts.\allowbreak{}

\url{https://sites.unimi.it/molteni/research/papers-pdf/44-molteni-An_explicit_Chebotarev_density_theorem_under_GRH.pdf}.

\paragraph{HK20-A1} Peter Humphries and Rizwanur Khan, On the Random Wave Conjecture for Dihedral Maaß Forms, GAFA30(2020),34–\allowbreak{}125.

\textbf{Locator:} AppendixA.\allowbreak{}1 printed106–\allowbreak{}107/\allowbreak{}PDF73–\allowbreak{}74:\allowbreak{} character/\allowbreak{}infinity-type formula, correspondence, spectral parameter, coefficient sum and all finite-place Satake cases; LemmaA.\allowbreak{}1 printed107.

\textbf{Inspection:} Primary publisher HTML and PDF appendix read directly; exact character formula and normalization checked.\allowbreak{} Classical correspondence cited, not re-proved; explicit MT-H01/\allowbreak{}03 proved separately.\allowbreak{}

\url{https://link.springer.com/article/10.1007/s00039-020-00526-4}.

\paragraph{CALEGARI13} Frank Calegari, The Artin conjecture for some S5-extensions, Math.\allowbreak{}Ann.\allowbreak{}356(2013),191–\allowbreak{}207.

\textbf{Locator:} §1 publisher-author PDF2 paragraph before Theorem1.\allowbreak{}1; Theorems1.\allowbreak{}1/\allowbreak{}1.\allowbreak{}2 PDF2–\allowbreak{}3; Theorem5.\allowbreak{}1 PDF15 and Theorem5.\allowbreak{}3 PDF16; arXiv1112.\allowbreak{}1152v1 Theorem1.\allowbreak{}2 p2 for H=EF; MT-S504 additionally uses arXiv:\allowbreak{}1112.\allowbreak{}1152v1 comment after Theorem1.\allowbreak{}2p2, Lemmas4.\allowbreak{}8–\allowbreak{}4.\allowbreak{}9pp10–\allowbreak{}11, final Theorem1.\allowbreak{}2 proofp11; MT-S506/\allowbreak{}507 character labels:\allowbreak{} section2tablep4; pointwise holomorphy premise:\allowbreak{} Theorem1.\allowbreak{}1p2; MT-S508 uses publisher Theorem4.\allowbreak{}7p12,Lemma4.\allowbreak{}8pp13-14,LandauLemma4.\allowbreak{}9p14,Lemma4.\allowbreak{}10pp14-15 and finalTheorem1.\allowbreak{}2proofp15.

\textbf{Inspection:} Primary author-hosted publisher PDF statements and §1 read; arXiv v1 read to recover explicitly the H=EF definition missing from extracted publisher statement.\allowbreak{} Known automorphy cited, not re-proved; computations not re-run.\allowbreak{} Root freshly read those exact v1 statements and the full final three proofs; initial automorphic construction retains classical credit.\allowbreak{} Root read the full publisher comparison/\allowbreak{}Landau/\allowbreak{}final proofs for MT-S508.\allowbreak{} The finite bad-prime polynomial ratio on p13 is corrected using the explicitly defined completed quotient; its real-interval nonvanishing is unchanged.\allowbreak{}

\url{https://math.uchicago.edu/~fcale/papers/S5.pdf}.

\paragraph{PY07} Dipendra Prasad and C.\allowbreak{}S.\allowbreak{}Yogananda, A Report on Artin’s holomorphy conjecture, February23 2007.

\textbf{Locator:} Theorem3.\allowbreak{}3 PDF5; §3.\allowbreak{}3 PDF8–\allowbreak{}9 completed functional equation and gamma factors; MT-S506/\allowbreak{}507 also use section2.\allowbreak{}1Properties1-2p3 and section3.\allowbreak{}1pp4-5.

\textbf{Inspection:} Primary author-hosted exposition read for exact finite Euler definition, Brauer theorem and functional-equation formulas.\allowbreak{} Its pre-2009 status discussion is not used for current modularity scope.\allowbreak{} Root additionally read the exact all-prime induction/\allowbreak{}multiplication and Hecke analyticity statements for MT-S506/\allowbreak{}507.\allowbreak{}

\url{https://www.math.iitb.ac.in/~dprasad/artin.pdf}.

\paragraph{MEYER05-Q} Ralf Meyer, A spectral interpretation for the zeros of the Riemann zeta function, Göttingen Seminars2004–\allowbreak{}05,117–\allowbreak{}137.

\textbf{Locator:} Section2 strong Schwartz definitions; Theorem3.\allowbreak{}3 printed124/\allowbreak{}PDF8; Section4 quotient definitions printed125; Theorem4.\allowbreak{}1 and complete proof printed126/\allowbreak{}PDF10; Corollary4.\allowbreak{}2 printed127; Proposition5.\allowbreak{}6 printed131–\allowbreak{}132.

\textbf{Inspection:} Primary author/\allowbreak{}seminar PDF downloaded; complete Theorem3.\allowbreak{}3 and Theorem4.\allowbreak{}1 proofs read, with their stated classical prerequisites.\allowbreak{} Proposition5.\allowbreak{}6 statement located; no whole-paper certification claimed.\allowbreak{} MT-M21 proves the additional exact closed-ideal comparison.\allowbreak{}

\url{https://math.nyu.edu/~tschinke/WS04/pdf/meyer.pdf}.

\paragraph{MEYER05-ID} Ralf Meyer, On a representation of the idele class group related to primes and zeros of L-functions, Duke Math.\allowbreak{}J.\allowbreak{}127(2005),519–\allowbreak{}595.

\textbf{Locator:} Author version math/\allowbreak{}0311468v3, Definition4.\allowbreak{}1/\allowbreak{}(22)pp24–\allowbreak{}25; Theorem5.\allowbreak{}1pp38–\allowbreak{}39; Lemma5.\allowbreak{}4pp40–\allowbreak{}41; Theorem5.\allowbreak{}8pp43–\allowbreak{}44; Section5.\allowbreak{}7pp45–\allowbreak{}47, including Lemma5.\allowbreak{}10 and Theorem5.\allowbreak{}11; Lemma5.\allowbreak{}5 and Lemma5.\allowbreak{}6 pp41–\allowbreak{}42; §2.\allowbreak{}1 formula(5) pp6–\allowbreak{}7.

\textbf{Inspection:} Primary arXivv3 retrieved.\allowbreak{} Root and complexity read the full Section5.\allowbreak{}7 image argument, Lemma5.\allowbreak{}4 embedding proof, and Theorem5.\allowbreak{}11 short proof, with scope-specific checks of definitions, levels and character convention.\allowbreak{} Complexity inspected source images46–\allowbreak{}47.\allowbreak{} MT-M22 has complete author and root proof readings.\allowbreak{} No complete-paper audit claimed.\allowbreak{} Root and complexity subsequently read the complete Lemma5.\allowbreak{}6 proof and exact nuclear-extension scope for MT-M27; the upstream local kernel results retain classical credit.\allowbreak{}

\url{https://arxiv.org/pdf/math/0311468}.

\paragraph{HKMS22} Daniel Hu, Ikuya Kaneko, Spencer Martin and Carl Schildkraut, Order of zeros of Dedekind zeta functions, Proc.\allowbreak{}AMS150(2022),5111–\allowbreak{}5120, DOI10.\allowbreak{}1090/\allowbreak{}proc/\allowbreak{}16041.

\textbf{Locator:} Author2107.\allowbreak{}03269v2 Theorems1.\allowbreak{}1–\allowbreak{}1.\allowbreak{}2 p2; complete proofs Sections3–\allowbreak{}5 pp5–\allowbreak{}8; general degree>=4 statement is Conjecture6.\allowbreak{}2.

\textbf{Inspection:} Root read the entire nine-page author paper, including both proofs of Theorem1.\allowbreak{}1 and the full Theorem1.\allowbreak{}2 proof.\allowbreak{} Cited group-theory and Stark prerequisites were not all independently reproved.\allowbreak{} Source multiplicity result is classical.\allowbreak{}

\url{https://arxiv.org/pdf/2107.03269}.

\paragraph{WONG19} Peng-Jie Wong, Zeros of Dedekind zeta functions and holomorphy of Artin L-functions, J.\allowbreak{}Ramanujan Math.\allowbreak{}Soc.\allowbreak{}34(2019),253–\allowbreak{}261.

\textbf{Locator:} Theorems1.\allowbreak{}3–\allowbreak{}1.\allowbreak{}4 p254, complete proofs259–\allowbreak{}260; Section4 lemmata256–\allowbreak{}258; Corollaries5.\allowbreak{}1 and5.\allowbreak{}3 pp259–\allowbreak{}260; Theorems5.\allowbreak{}2/\allowbreak{}5.\allowbreak{}4 additional Sylow variants.

\textbf{Inspection:} Root read the complete mathematical text pp253–\allowbreak{}260 and reference list p261, including both main proofs and all Section4 lemmata.\allowbreak{} Exact solvability hypothesis retained from theorem statements, despite broader abstract wording; classical theorem cited, not independently reproved.\allowbreak{}

\url{https://www-math.nsysu.edu.tw/~pjwong/stuff/LowOrderZero.pdf}.

\paragraph{FROHLICH74} Albrecht Fröhlich, The Galois module structure of algebraic integer rings in fields with generalised quaternion group, Mém.\allowbreak{}SMF37(1974),81–\allowbreak{}86.

\textbf{Locator:} Locally free class invariant pp81–\allowbreak{}82; Theorem1 p82.\allowbreak{} Original proof:\allowbreak{} Artin-root numbers and normal integral bases for quaternion fields, Invent.\allowbreak{}Math.\allowbreak{}17(1972),143–\allowbreak{}166, DOI10.\allowbreak{}1007/\allowbreak{}BF01418937.

\textbf{Inspection:} Root read the invariant definition and full Theorem1 statement and surrounding explanation pp81–\allowbreak{}82.\allowbreak{} Original1972 full proof was not accessible; classical theorem used by exact locator, not certified as reread in full.\allowbreak{}

\url{https://www.numdam.org/article/MSMF_1974__37__81_0.pdf}.

\paragraph{MARTINET71} Jacques Martinet, Modules sur l’algèbre du groupe quaternionien, Ann.\allowbreak{}Sci.\allowbreak{}ENS4(1971),399–\allowbreak{}408.

\textbf{Locator:} TheoremIV.\allowbreak{}1 and proof p406; Example3 p407 gives the tame totally real nonfree Q8 field with primes5,41.

\textbf{Inspection:} Root read TheoremIV.\allowbreak{}1 complete displayed proof and Examples1–\allowbreak{}3; original page407 rendered and visually read to verify sqrt(5*41) in the radicand.\allowbreak{} Upstream integral-module classification theorems not all reproved.\allowbreak{}

\url{https://jamartin.perso.math.cnrs.fr/AlgNT/AnnENS1971.pdf}.

\paragraph{MARTINET72} Jacques Martinet, Fonctions zêta au point1/\allowbreak{}2, Séminaire de théorie des nombres de Grenoble1(1971–\allowbreak{}1972),108–\allowbreak{}116.

\textbf{Locator:} SectionV proposition and complete proof pp113–\allowbreak{}114; SectionVI quaternion root-number theorem p114 and corollary p115.

\textbf{Inspection:} Root read complete SectionV proof and SectionVI theorem/\allowbreak{}corollary; corroborates the exact Fröhlich theorem and monomial quaternion factorization.\allowbreak{} No wholesale current-status reliance on this historical survey.\allowbreak{}

\url{https://numdam.org/item/STNG_1971-1972__1__108_0.pdf}.

\paragraph{TAYLOR81} Martin J.\allowbreak{} Taylor, On Fröhlich’s conjecture for rings of integers of tame extensions, Invent.\allowbreak{}Math.\allowbreak{}63(1981),41–\allowbreak{}79, DOI10.\allowbreak{}1007/\allowbreak{}BF01389193.

\textbf{Locator:} Theorem1 and Corollary p41; complete odd-degree/\allowbreak{}cancellation proof p42; definition of the root-number class in §2.

\textbf{Inspection:} Retrieved public GDZ IIIF manifest and original scanned pages41–\allowbreak{}42 (volume scan47–\allowbreak{}48); visually read full theorem and corollary, its entire short proof and introduction.\allowbreak{} The long proof of Theorem1 was not reread; CNT1983 main proof and exact upstream references were read separately.\allowbreak{}

\url{https://gdz.sub.uni-goettingen.de/id/PPN356556735_0063}.

\paragraph{CNT83} Philippe Cassou-Noguès and Martin J.\allowbreak{} Taylor, Constante de l’équation fonctionnelle de la fonction L d’Artin d’une représentation symplectique et modérée, Ann.\allowbreak{}Inst.\allowbreak{}Fourier33(2)(1983),1–\allowbreak{}17, DOI10.\allowbreak{}5802/\allowbreak{}aif.\allowbreak{}913.

\textbf{Locator:} §1 hermitian K0/\allowbreak{}discriminant pp3–\allowbreak{}5; (2.\allowbreak{}2)–\allowbreak{}(2.\allowbreak{}10) and Proposition2.\allowbreak{}6 pp6–\allowbreak{}8; (3.\allowbreak{}6)–\allowbreak{}(3.\allowbreak{}9)p9; main Theorem/\allowbreak{}CorollaryA p10, CorollaryB and complete main proofp11; local comparison §4pp12–\allowbreak{}17.

\textbf{Inspection:} Root read the complete retrieved mathematical text, including Proposition2.\allowbreak{}6 and the entire main proofp11 and local comparison.\allowbreak{} Originalpp5,9,10 rendered and visually checked for Qbar scalar extension, square-root conductor factor and the group-ring involution in the standard trace form.\allowbreak{} Upstream Gauss-sum/\allowbreak{}determinant theorems remain classical cited inputs, not independently reproduced.\allowbreak{}

\url{https://www.numdam.org/item/10.5802/aif.913.pdf}.

\paragraph{BURNOL02} Jean-François Burnol, A lower bound in an approximation problem involving the zeros of the Riemann zeta function, Adv.\allowbreak{}Math.\allowbreak{}170(2002),56–\allowbreak{}70, DOI10.\allowbreak{}1006/\allowbreak{}aima.\allowbreak{}2001.\allowbreak{}2066.

\textbf{Locator:} Theorem1.\allowbreak{}3 authorp4; preceding definitions of B\ensuremath{\lambda} and D(\ensuremath{\lambda}), authorpp2–\allowbreak{}4; arXiv:\allowbreak{}math/\allowbreak{}0103058v2.

\textbf{Inspection:} Root read exact definitions and full Theorem1.\allowbreak{}3 statement, with distinction between continuous dilations and discrete approximation.\allowbreak{} The complete long source proof was not reread; lower bound remains cited classical input.\allowbreak{}

\url{https://arxiv.org/pdf/math/0103058}.

\paragraph{BDR10} Michel Balazard and Anne de Roton, Sur un critère de Báez-Duarte pour l’hypothèse de Riemann, Int.\allowbreak{}J.\allowbreak{}NumberTheory6(4)(2010),883–\allowbreak{}903, DOI10.\allowbreak{}1142/\allowbreak{}S1793042110003307.

\textbf{Locator:} §1 unnumbered main Theorem authorp3; §2 Propositions1–\allowbreak{}3 and explicit epsilon\_\allowbreak{}N choice pp4–\allowbreak{}5; arXiv:\allowbreak{}0812.\allowbreak{}1689v1.

\textbf{Inspection:} Root read exact main statement, effective-constant convention, full Mellin–\allowbreak{}Plancherel Proposition1 proof, Propositions2–\allowbreak{}3 and explicit damping parameter25.\allowbreak{} Later long estimates remain classical source dependencies, not a new whole-source audit.\allowbreak{}

\url{https://arxiv.org/pdf/0812.1689}.

\paragraph{DEROTON06} Anne de Roton, Une approche hilbertienne de l’hypothèse de Riemann généralisée, Bull.\allowbreak{}Soc.\allowbreak{}Math.\allowbreak{}France134(3)(2006),417–\allowbreak{}445, DOI10.\allowbreak{}24033/\allowbreak{}bsmf.\allowbreak{}2516.

\textbf{Locator:} Definition1p418; Definition2/\allowbreak{}Theorem1p419; Theorem2p420; §7pp441–\allowbreak{}444, Propositions7.\allowbreak{}1–\allowbreak{}7.\allowbreak{}2 and §7.\allowbreak{}2.

\textbf{Inspection:} Root read precise Selberg axioms, complementary-function/\allowbreak{}dilation definitions and both theorem statements.\allowbreak{} Originalp420 image checked.\allowbreak{} Entire final §7 proof read, including all Gram-limit and inner-product calculations; earlier §§2–\allowbreak{}6 operator estimates are cited prerequisites.\allowbreak{} MT-NB02 separately verifies the Hecke specialization.\allowbreak{}

\url{https://www.numdam.org/item/10.24033/bsmf.2516.pdf}.

\paragraph{BCF13} Sandro Bettin, J.\allowbreak{}Brian Conrey and David W.\allowbreak{}Farmer, An optimal choice of Dirichlet polynomials for the Nyman–\allowbreak{}Beurling criterion, Proc.\allowbreak{}SteklovInst.\allowbreak{}Math.\allowbreak{}280Suppl2(2013),S30–\allowbreak{}S36.

\textbf{Locator:} Theorem1 and explicit V\_\allowbreak{}N authorp2; Lemmas2–\allowbreak{}3 pp5–\allowbreak{}7; complete Theorem1 proof pp7–\allowbreak{}9; arXiv:\allowbreak{}1211.\allowbreak{}5191v1.

\textbf{Inspection:} Root read exact Theorem1 and its simplicity requirement, then the entire final proofpp7–\allowbreak{}9 and displayed Lemma3 proof.\allowbreak{} Main result is cited classical input; no claim of fully auditing all earlier lemmas or proving the derivative-moment hypothesis.\allowbreak{}

\url{https://arxiv.org/pdf/1211.5191}.

\paragraph{DIST18} Matei Machedon, Expository Notes on Distribution Theory, AMSC/\allowbreak{}MATH673, Fall2018, University of Maryland.

\textbf{Locator:} §2 Theorem2.\allowbreak{}1p4:\allowbreak{} distributions supported at one point; full Taylor-remainder proof.

\textbf{Inspection:} Root read precise theorem and complete short proof; used only the classical finite delta-derivative representation.\allowbreak{} Weighted cutoff and annihilator estimates in MT-M31 are separately proved here.\allowbreak{}

\url{https://www.math.umd.edu/~matei/distributions13-14.pdf}.

\end{document}
